% ------ METADATA ------
\newcommand{\bookauthor}{\ml{$0$}{Эмиль~Весна}{Emil~Viesn\'{a}}}
\newcommand{\booktitle}{\ml{$0$}{ДАЛЕКИЕ~РЕЧИ}{Spoken~Far~Away}}
\newcommand{\bookstarted}{\ml{$0$}{13 августа 2012}{August 13, 2012}}
\newcommand{\bookfinished}{\ml{$0$}{пока нет}{not yet}}
% ------ METADATA ------

% ----- XELATEX SYMBOL -----
\usepackage{xltxtra}
% ----- XELATEX SYMBOL -----

% ----- HYPHENATION -----
\usepackage{hyphenat}
% ----- HYPHENATION -----

% ----- GEOMETRY -----
\usepackage[left=1.5cm,right=1.5cm,top=2cm,bottom=2cm,bindingoffset=0.5cm]{geometry}
% ----- GEOMETRY -----

% ----- INCLUDE PDF AS PAGES -----
\usepackage{pdfpages}
% ----- INCLUDE PDF AS PAGES -----

% ----- DROPPING CAP -----
\usepackage{type1cm,lettrine}
% ----- DROPPING CAP -----

% ----- FONTS -----
\renewcommand{\baselinestretch}{1.2}
\setmainfont{Linux Libertine}
% ----- FONTS -----

% ------ HYPERLINKS ------
\usepackage{hyperref}
\definecolor{LinkColor}{HTML}{0969DA}
\hypersetup{colorlinks=true, linkcolor=LinkColor, citecolor=LinkColor, filecolor=LinkColor, urlcolor=LinkColor}
% ------ HYPERLINKS ------

% ------ FANCY PAGE STYLE ------
\setlength{\headheight}{15pt}
\usepackage{fancyhdr}
\pagestyle{fancy}
\fancyhead[LE,RO]{\thepage}
\fancyhead[LO]{{\small\textsc{\booktitle}}}
\fancyhead[RE]{{\small\textsc{\bookauthor}}}
\fancyfoot{}
    \fancypagestyle{plain}{
    \renewcommand{\headrulewidth}{0mm}
    \fancyhead{}
    \fancyfoot{}
}
% ------ FANCY PAGE STYLE ------

% ------ ELEMENTS ------
\newcommand{\asterism}{\vspace{1em}{\centering\Large\bfseries$\ast~\ast~\ast$\par}\vspace{1em}}
\newcommand{\textspace}{\vspace{1em}{\centering\Large\bfseries<\dots>\par}\vspace{1em}}
\newcommand{\FM}{\footnotemark}
\newcommand{\FL}[2]{\footnotetext{См. \textit{\hyperlink{#1}{#2}}.}}
\newcommand{\FA}[1]{\footnotetext{#1 \emph{\ml{$0$}{---~Прим.~авт.}{---~Author.}}}}
\newcommand{\theterm}[3]{\textbf{\hypertarget{#1}{#2}} --- #3}
\newcommand{\thesynonim}[3]{\textbf{#2} --- см. \textit{\hyperlink{#1}{#3}}.}
\newcommand{\theorigin}[3]{\textit{#1:} #2 --- #3}
% ------ ELEMENTS ------

% ------ EPIGRAPH ------
\usepackage{epigraph}
\renewcommand{\epigraphsize}{\footnotesize}
\epigraphrule=0pt
\epigraphwidth=8cm
\usepackage{etoolbox}
\AtBeginEnvironment{quote}{\itshape}
\makeatletter
\newlength\episourceskip
\pretocmd{\@episource}{\em}{}{}
\apptocmd{\@episource}{\em}{}{}
\patchcmd{\epigraph}{\@episource{#1}\\}{\@episource{#1}\\[\episourceskip]}{}{}
\makeatother
% ------ EPIGRAPH ------

% ------ NAMES ------
\newcommand{\Aatris}{\"{A}\={a}tr\v{\i}s}
\newcommand{\Akchsar}{\`{A}kchs\r{a}r}
\newcommand{\Chhammitrai}{Chh\`{a}mm\={\i}tr\^{a}i}
\newcommand{\Chhanei}{Chh\r{a}n\^{e}i}
\newcommand{\Chitram}{Ch\"{\i}tr\'{a}m}
\newcommand{\Choe}{Cho\^{e}}
\newcommand{\choe}{cho\^{e}}
\newcommand{\Harrmatr}{H\r{a}rrm\`{a}tr}
\newcommand{\tHat}{H\={a}t}
\newcommand{\Hei}{H\r{e}i}
\newcommand{\hei}{h\r{e}i}
\newcommand{\Hoesitr}{Ho\`{e}s\={\i}tr}
\newcommand{\hoesitr}{ho\`{e}s\={\i}tr}
\newcommand{\Kchaagotr}{Kch\^{a}\={a}g\~{o}tr}
\newcommand{\kchaagotr}{kch\^{a}\={a}g\~{o}tr}
\newcommand{\Kcharas}{Kch\'{a}r\v{a}s}
\newcommand{\Kchatrim}{Kch\r{a}tr\"{\i}m}
\newcommand{\Kchenoel}{Kch\={e}no\^{e}}
\newcommand{\kchenoel}{kch\={e}no\^{e}}
\newcommand{\Kchenoet}{Kch\"{e}no\^{e}}
\newcommand{\kchenoet}{kch\"{e}no\^{e}}
\newcommand{\Kchoho}{Kch\`{o}h\^{o}}
\newcommand{\Kchotlam}{Kch\={o}tl\'{a}m}
\newcommand{\Kchotris}{Kch\={o}tr\v{\i}s}
\newcommand{\Kihotr}{K\^{\i}h\~{o}tr}
\newcommand{\kihotr}{k\^{\i}h\~{o}tr}
\newcommand{\Kukchuatr}{K\`{u}kchu\={a}tr}
\newcommand{\kukchuatr}{k\`{u}kchu\={a}tr}
\newcommand{\Laaka}{L\={a}\"{a}k\^{a}}
\newcommand{\laaka}{l\={a}\"{a}k\^{a}}
\newcommand{\Lechoe}{L\={e}cho\`{e}}
\newcommand{\lechoe}{l\={e}cho\`{e}}
\newcommand{\Likas}{L\^{\i}k\v{a}s}
\newcommand{\Likchmas}{L\={\i}kchm\r{a}s}
\newcommand{\Likchoe}{L\^{\i}kcho\^{e}}
\newcommand{\Loem}{Lo\~{e}m}
\newcommand{\Maaras}{M\"{a}\={a}r\v{a}s}
\newcommand{\Mitchoe}{M\={\i}tcho\^{e}}
\newcommand{\Mitlikch}{M\={\i}tl\={\i}kch}
\newcommand{\Mitris}{M\={\i}tr\={\i}s}
\newcommand{\Oerchoe}{O\r{e}rcho\^{e}}
\newcommand{\Oerlikch}{O\r{e}rl\'{\i}kch}
\newcommand{\Sat}{S\={a}t}
\newcommand{\Satchir}{S\={a}tch\"{\i}r}
\newcommand{\Satrakch}{S\={a}tr\`{a}kch}
\newcommand{\Seli}{S\r{e}l\={\i}}
\newcommand{\Sirtchu}{S\r{\i}rtch\'{u}}
\newcommand{\Sitris}{S\~{\i}tr\v{\i}s}
\newcommand{\Siusiu}{S\~{\i}u-s\~{\i}u}
\newcommand{\siusiu}{s\~{\i}u-s\~{\i}u}
\newcommand{\Sogcho}{S\"{o}gch\={o}}
\newcommand{\Sotron}{S\~{o}tr\`{o}n}
\newcommand{\Tchalas}{Tch\r{a}l\v{a}s}
\newcommand{\Tchammitr}{Tch\`{a}mm\={\i}tr}
\newcommand{\Tchanoe}{Tch\r{a}no\^{e}}
\newcommand{\Tchartchaahitr}{Tch\~{a}rtch\"{a}\={a}h\r{\i}tr}
\newcommand{\Tchartchu}{Tch\~{a}rtch\'{u}}
\newcommand{\Tchitron}{Tch\"{\i}tr\`{o}n}
\newcommand{\Tchu}{Tch\`{u}}
\newcommand{\tchu}{tch\`{u}}
\newcommand{\Technku}{T\`{e}chnk\r{u}}
\newcommand{\Tesarrokch}{Te's\'{a}rr\r{o}kch}
\newcommand{\Tesatron}{Te's\'{a}tr\v{o}n}
\newcommand{\Traa}{Tr\={a}\"{a}}
\newcommand{\traa}{tr\={a}\"{a}}
\newcommand{\Trai}{Tr\r{a}i}
\newcommand{\trai}{tr\r{a}i}
\newcommand{\TraRenkchal}{Tr\r{a}-R\={e}nkch\'{a}l}
\newcommand{\Trukchual}{Tr\`{u}kchu\r{a}l}
% ------ NAMES ------

\begin{document}

% ----- FRONT COVER -----
% \includepdf[pages={1}]{\coverfrontfilename}
% ----- FRONT COVER -----

% ----- EMPTY PAGE -----
% \newpage\thispagestyle{plain}~
% ----- EMPTY PAGE -----

% ------ TITLE PAGE ------
\begin{titlepage}
{
\centering
{~\par}
\vspace{0.25\textheight}
{\LARGE\bookauthor\par}
\vspace{1.3cm}
{\Huge\textbf{\booktitle}\par}
\vfill
{\includegraphics[width=6em]{\publisherlogofilename}\par}
}
\end{titlepage}
% ------ TITLE PAGE ------

% ----- TRIGGER WARNING -----
% \newpage\thispagestyle{plain}
% {~\par}
% \vfill
% {\centering\includegraphics[width=4em]{tw.pdf}\par}
% \vspace{1em}
% {\centering\Large\strong{\ml{$0$}{Предупреждение}{Content Warning}}\par}
% \vspace{0.5em}
% {\centering\large{Данное произведение содержит описание неизлечимых болезней, психических состояний, врачебных манипуляций, упоминания суицида, химических зависимостей, сцены насилия и убийств, мат, а также гомофобную, трансфобную, сексистскую и расистскую лексику.}\par}
% \vfill
% {~\par}
% ----- TRIGGER WARNING -----

% ----- TABLE OF CONTENTS -----
\tableofcontents
% ----- TABLE OF CONTENTS -----

\pagestyle{fancy}

\chapter*{Предисловие}
\addcontentsline{toc}{chapter}{Предисловие}

\epigraph
{Знаете, что в искусстве самое страшное для консерваторов?
Как бы высоко ни взлетела фантазия автора, это не будет пределом возможного.
Если речь пойдёт о смехотворном сроке в два-три миллиона лет, я не поставлю ни на великую страну, ни на истинную религию, ни на незабвенный язык, ни даже на богоизбранный биологический вид --- ни на что, чем гордятся и за что убивают.
Всё канет в Лету\FM.
И многое --- бесследно.}
{Мартин Охсенкнехт, драматург Театра Снов.
Эпоха Последней Войны}
\FA{Река времени из эллатинских легенд Древней Земли.}

Тишина.
Слово едва ли не прекрасней его собственной сути.
Словарь языка Эй даёт ему следующее определение: <<Совокупность незначимых надпороговых сигналов, воспринимаемых рецептором>>.
Для космических радаров тишина --- это белый электромагнитный шум.
Человек-тагуа может назвать тишиной и ропот ночного города, и звёздное небо, и нежный ветер, касающийся волос.
Для хоргета же тишина --- это вечное, непрекращающееся колебание поля Кохани"--~Вейерманна.

Тишина не равнозначна безмолвию.
Безмолвие --- <<отсутствие надпороговых сигналов>> --- давит и угнетает.
Тишина же --- это ласка, расслабляющий массаж для рецептора.
В тишине система становится чувствительной и замечает то, что ранее оставалось в тени.

Последний год был годом тишины.
Орден Преисподней наградил меня отдыхом за военные заслуги.
Также мне разрешили изучить новую специализацию --- диктиология.

--- Зачем тебе это нужно? --- удивлённо спросил Грейсвольд, узнав о моём решении.

--- Я чересчур долго откладывал обучение, --- ответил я.
--- Опыт --- почва для уверенности, но на опыте без теории растут самые серьёзные ошибки.

Базовое обучение на диктиолога занимает двадцать суток по времени Капитула.
Моё тело провело их в медикаментозной коме; демону же спать не приходилось --- стыковка новых модулей, загрузка пакетов, устранение несовместимостей и настройка окружений, утомительный разбор синтаксиса и структуры, а под конец --- обучение
дополнительного нейроконтура, чтобы я мог пользоваться новыми надстройками с наилучшей производительностью.

Отпуск мне рекомендовали посвятить научным исследованиям, чтобы проверить модули на практике.
Работа скучная и нетворческая --- разбор чужих отчётов, записей и изображений, выявление закономерностей в странных, порой диких завитушках сапиентной культуры.

Единственное отрада --- научная работа по народам Тра-Ренкхаля.
Эта планета до последнего времени была тёмным пятном на небосводе, и я, занесённый туда ветром войны, принял непосредственное участие в исследовании.

Эта книга --- художественное приложение к работе.
Повесть о прекрасной планете, населённой прекрасным народом.
Изнывающим от скуки неинтересны истории, конец которых заранее известен;
их мир определён, они жаждут неожиданностей.
В этой книге неожиданностей не будет.
Для меня удовольствием было её писать.
Хотелось бы, чтобы вы нашли удовольствие в чтении.

Скажу сразу, дорогие читатели: ваша реакция на книгу будет тщательно проанализирована агентами Ордена Преисподней.
На основании ваших впечатлений вынесут вердикт, усвоил ли я обучение и для какой работы я подхожу.
Но пусть вас это не смущает.
Я писал не ради тёплого места в отделе культурологии.
Существовали более тривиальные способы доказать профессиональную пригодность.

Я знаю, что многие демоны имеют предубеждения по поводу сапиентов.
Несмотря на то, что Орден Преисподней сделал права сапиентов одним из своих приоритетов, несмотря на то, что многие демоны говорят о наличии у них сапиентов-друзей, далеко не все готовы увидеть сейхмар в главной роли.
Я же постарался донести до читателя главную мысль: всё, что происходит в первой части книги --- видение не спящего демона, но человека, который несёт на себе демона по воле судьбы.

В книге вы найдёте множество моих собственных предубеждений.
\ml{$0$}
{Я знаю об их наличии, и я не стал от них избавляться.}
{I know they exist, and I let them be.}
Книга писалась в течение многих лет;
в конце концов, у меня нет амбиций быть примером для подражания, а предубеждения --- это письменное свидетельство того, откуда я начал и куда пришёл.

Книга написана на двух языках.
Язык Эй и мой родной сектум-лингва, язык Ада и язык Картеля, речи двух близких и далёких миров.
Прошу всех отметить этот факт и не воспринимать одну часть всего лишь как перевод другой.
Основная цель подобного приёма --- сделать книгу более понятной (во всех смыслах).
Не каждую мысль можно выразить чётко на одном языке;
два языка дают хоть какую-то гарантию того, что мысль будет понята.
Тем не менее, несмотря на различия, призванные сохранить колорит повествования, оба варианта практически идентичны.
Я старался ничего не потерять.

\subsection*{Физиолингвистика}

Сапиентные виды, о которых пойдёт речь, не совсем обычны.
Думаю, многие наслышаны о Девиантных Ветвях, созданных демиургами или сапиентами с нуля.
Однако во Вселенной существует целая группа Девиантных Ветвей, которые создали сами себя.
Вы прекрасно знаете, о ком я говорю.
И вполне естественно предположить, что особенности их физиологии и анатомии наложили отпечаток на сигнальную систему.

Как выглядит мир в глазах того, кто воспринимает инфракрасные, ультрафиолетовые и рентгеновские волны?
Что слышат уши того, кто улавливает инфра- и ультразвук?
Какие ощущения дают магнитные поля?
К сожалению, для читателей-сапиентов это останется тайной за семью печатями.
Я не стал размениваться на образы --- это похоже на попытку рассказать собаке, чем фиолетовый отличается от зелёного.
Образы искажают главную мысль --- тси родились и живут в мире, на который настроены их рецепторы, и он для них обычен так же, как обычен ваш мир для вас.

В языках тси имеются термины для описания ощущений.
В конце книги я дал их научное обоснование; представить, что чувствуют эти сапиенты, оставляю вам.
При словах <<солнечная кисть>> можно подумать о ярких одуванчиках, а <<каменная ярость>> может вызвать ассоциацию с мертвенным свечением фосфора.
Также у тси, обладающих способностью видеть в инфракрасном диапазоне, существует огромное количество эпитетов, обозначающих особенности кожного кровотока и свечения.
К сожалению, их придётся опустить почти полностью.
Но всё это неважно.
Сколько бы ни существовало глаз, мир всегда один.

Пара слов о языке сели.

Люди-тси имеют сложный голосовой аппарат, отдалённо напоминающий птичий --- так называемая гортанная цитра, позволяющая воспроизвести любой звук в их слуховом диапазоне.
Голосовая часть языка сели --- звукоподражательная\FM, для её записи разработана специальная система обозначений на основе письменности Молчащих лесов.
\FA{
В частности, названия всех животных на языке сели представляют собой почти точное вопроизведение их криков или звуков.
К примеру, самец согхо называется <<тиулиулафью>>, домашний петух --- <<кукареку>>.
Детёныши и самки животных также имеют своё звукоподражательное название.
Забавные аспекты --- распил древесины (процесс, которого не знали древние тси) также обозначается звуком, который издаёт пила.
}
Между тем сели никогда не используют только голос --- в зависимости от ситуации они переходят то на шумовую, то на жестовую речь.

Вы представили разговор на языке сели?
Безусловно, это сложно.
Приведу слова поэтессы Эрхэ Колокольчик, для которой этот язык был неисчерпаемым источником вдохновения:

\begin{quote}
<<Урчание оцелота, кваканье жаб, свист согхо, крики множества птиц и зверей, жужжание светлячков и скрежет ночных сверчков --- все звуки джунглей в устах одного человека>>.
\end{quote}

Добавлю, что это не столько поэтическое сравнение, сколько точно описанное впечатление постороннего.

Диалоги на языке тси могут показаться вам очень сложными и запутанными.
Этому есть вполне понятная причина --- язык тси позволяет одновременно обсуждать несколько тем (до пяти), создавать <<развилки>> в предложениях и объединять их.
Это поистине уникальное явление, обнаруженное только у тси и тси-язычных племён Тра-Ренкхаля, называется лингвистическим распараллеливанием.
Некоторое представление о распараллеливании в тси-языках может дать их письменность\FM.
\FL{text-structure}{Структура текста в языке сели}

В силу того, что лексемы языка тси невозможно транслитерировать с помощью СЧФ, я калькировал\FM\ имена тси и большую часть терминов.
\FA{
В соответствии с правилами калькирования терминов, установленными <<Инструкцией по локализации для отдела 44>>, версией 11-A0FD49.
Ст. 1 требует использовать калькирование в случае критического нарушения удобочитаемости (п. 1) и/или невозможности использования СЧФ (п. 3).
Ст. 4, п. 5 позволяет поместить определение кальки в любое подходящее место или опустить его, если определение не требуется.
}

Также в языке тси имеются слова, которые невозможно калькировать (так называемые <<первичные>>, не имеющие родительского корня), и поэтому я переопределил\FM\ термины сектум-лингва и Эй.
\FA{
В соответствии с правилами переопределения терминов, установленными <<Инструкцией по локализации для отдела 44>>, версией 11-A0FD49.
Ст. 1 требует использовать переопределение в случае невозможности калькирования (п. 4) при критическом нарушении удобочитаемости (п. 1) и/или невозможности использования СЧФ (п. 3).
Ст. 4, п. 4 требует вынести переопределение терминов в предисловие.
}
Прошу уделить этому параграфу пристальное внимание.

~

\begin{description}
\item[Даритель] --- существо, которое отдало индивиду часть генома и/или клеточного аппарата.
Слово употребляется в контексте биологического родства.
У сели дарителями называются мужчина и женщина, зачавшие ребёнка.
У тси дарителей могло быть любое количество --- от одного (клонирование, партеногенез) до нескольких десятков (сборка генома).
\item[Хранитель] --- существо, которое приняло от индивида часть генома и/или клеточного аппарата.
\item[Друг] --- любое близкое существо индивида.
\item[Кормилец] --- существо, которое занималось воспитанием индивида.
По отношению к кому-либо употребляется в контексте воспитания.
Количество кормильцев обычно от одного до пяти.
\item[Питомец] --- существо, воспитанием которого занимался индивид.
Синонимом является выражение <<чей-либо ребёнок>>.
\item[Любовник] --- существо, связанное с индивидом регулярными половыми связями.
\item[Потомок (наследник)] --- любое существо, находящееся в нижнем вертикальном родстве с индивидом.
Порядковый номер обозначает поколение.
Хранители также являются первыми потомками.
\item[Предок (прародитель)] --- любое существо, находящееся в верхнем вертикальном родстве с индивидом.
Порядковый номер обозначает поколение.
Дарители являются также первыми предками.
\item[Родильница] --- существо, которое обеспечило индивиду внутриутробное развитие, в том числе (предположительно) кольцевая теплица.
Слово было в ходу у тси, у народа сели же стало высокопарным синонимом дарительницы, так как в подавляющем числе случаев родильница и дарительница --- одно и то же существо\FM.
\FA{
В языке тси эти слова не имеют рода.
Сели иногда употребляют также более точные заимствованные термины для обозначения родства (се --- первый по счёту первый наследник, тхар --- первая наследница се, кхантхор --- первый наследник брата второго прародителя и пр.), чаще при разговоре с представителями других народов.
}
\item[Ребёнок] --- существо, не прошедшее половую дифференциацию.
Слово не является гиперонимом по отношению к словам <<мальчик>> и <<девочка>>.
Мальчиками и девочками называются молодые люди, прошедшие половую дифференциацию, но не достигшие возраста зрелости --- двадцати пяти дождей по Вехам Жатвы.\\
\item[Дом] --- место, где индивид осуществляет полноценный отдых.
Не путать с жилищем --- строением, предназначенным для проживания\FM.
\FA{
Понятие дома у тси достаточно размытое.
Были случаи, когда существа большую часть жизни проводили в своём жилище, но также были случаи постоянного проживания в транспорте, в цеховых комнатах, в палатках, в гостях и даже на природе без достоверного снижения коэффициента личного комфорта.
Поэтому на первый план выходит полноценный отдых.
У сели сохранилась поговорка-каламбур <<Всегда превращай в дом место, где спишь>>.
Это отражает традицию поддерживать комфортный порядок не только на месте сна, но и на прилегающей территории.
}
\end{description}

Для местного колорита я использовал подходящие термины из языка цатрон.
Сели используют цатрон для общения с другими видами ветви Хуманы (их представителям выучить язык сели практически невозможно) и при официальном общении.
По установившейся традиции детям сели дают три имени --- классическое трёхзнаковое имя тси, однознаковое <<домашнее>> имя, а также родоплеменное имя царрокх с изменённым порядком наследования родового слога.

В приложениях к книге я сделал краткое научно-популярное описание явлений, связанных с событиями.
Там же вы можете найти определения терминов.
Если кому-то кажутся скучными постоянные консультации у словаря, лучше закройте книгу прямо сейчас --- вы всё равно не дочитаете её до конца.

~

~

Итак, сейчас на Капитуле 13.004.56281.4, и я, Аркадиу Талианский Шакал, начинаю свою повесть.

\chapter*{Интерлюдия I. Безумный}

\addcontentsline{toc}{chapter}{Интерлюдия I. Безумный}

\textbf{Легенда об обретении}

Прибыл однажды в город Харматр\FM\ с запада вестник, ростом в два цана, в сером одеянии.
\FA{
Легендарное место в восточных джунглях.
В настоящее время учёные Ада склоняются к мысли, что Гибельные Руины и Харматр --- одно и то же место.
}
Он прошёл через город, поднялся на ступени храма и повелел собрать народ.

<<Слушайте! Слушайте! --- возвестил гонец.
--- Я пришёл к вам от Безымянного, творца этого мира и этих небес.
Он отвёл беду от земли и от моря, из тех бед, что гасят звёзды и раскалывают горы до основания.
Он пострадал за нас, и лишь жертвенная кровь, лишь муки умирающих могут спасти его!>>

Притихли люди.
Посовещались они и решили принести Безымянному жертвы.
Как вдруг вперёд вышла седая старуха, знавшая ещё строителей древнего города и варившая им пищу.

<<Если бы мне выпало защитить от беды питомцев, я бы не требовала ни платы, ни поклонения, --- сказала старуха.
--- Только безумный, спасая детей, восполняет силы их кровью!>>

Согласились люди со старой женщиной и разошлись.
Гонец ухмыльнулся и покинул город.

Прошла декада, и прибыл с запада второй [вестник, ростом] в три цана, в зелёном [одеянии].
Он прошёл через город, и народ сам стекался на площадь, ожидая его слова.

<<Слушайте! Слушайте! --- возвестил гонец.
--- Я пришёл к вам от Безымянного, творца этих джунглей и этой реки.
Он сотворил вас, недостойных детей, и дал вам мир для услады и размножения.
Лишь жертвенная кровь, лишь муки умирающих могут стать ему достойной платой!>>

Притихли люди.
Посовещались они и решили принести Безымянному жертвы.
Однако вперёд вышла та же старуха и снова попросила слова.

<<Безымянный создал меня по собственным лекалам и прихоти.
Я не просила ни жизни, ни прекрасного мира, --- сказала старуха.
--- Только безумный будет делать брак и требовать с него оплату!>>

Согласились люди со старой женщиной и разошлись.
Гонец покинул город, гневно хмурясь.

Прошла декада, и прибыл с запада третий [вестник, ростом] в пять цан, в оранжевом [одеянии].
Он прошёл через город, и народ лишь вылез на крыши, чтобы видеть и слышать сказанное.

<<Слушайте! Слушайте! --- возвестил гонец.
--- Я пришёл к вам от Безымянного, творца этих гор и этого моря.
Известны ему дороги солнца и комет, известны тропы смертных и тех, кто выходит из бренных тел.
Запутан путь к пристанищу духов, и лишь жертвенная кровь, лишь муки умирающих помогут заблудшим!>>

Испугались люди, и раздались в городе крики отчаяния.
Но старуха не унималась:

<<[Путь к пристанищу] всегда и для всех был прям --- смерть, --- закричала она.
--- Только безумный, будучи пойманным на лжи дважды, [соврёт] в третий раз!>>

Отругали люди гонца отборными словами, и покинул он город, тлея злобой унижения.

Прошло три дня, и потемнело небо, и заискрились в горах молнии, и задули ветры с неистовой силой.
Прибыл с запада великан, [ростом] в десять цан, в чёрном [одеянии], с глазами, пылающими Каменной яростью.
Он прошёл через город, но народ не встретил его и затаился в домах, охваченный страхом.

<<Слушайте все, --- прогремел пришелец.
--- Я --- Безымянный, владыка ваш в жизни и смерти.
Я могу разрушить города и стереть в пыль ваши книги.
Я могу отнять ваши жизни и посмертие, и лишь жертвенная кровь, лишь муки умирающих могут отвести мой гнев>>.

Молчали люди, сидя в домах.
И лишь старуха вылезла на крышу, чтобы ответить.

Махнул рукой великан --- и старуха [сгорела] в синем пламени, крича от нестерпимой боли.
Махнул [ещё раз] --- и запылал уже весь город, и не было спасения ни взрослым, ни малышам, ни домашней скотине.

Пламя полыхало недолго.
Вскоре среди горячих руин, некогда прекраснейшего из городов, остались лишь двое --- пришелец [ростом] в десять цан и ученик жреца, стоявший на площади и избегнувший огня.

\ml{$0$}
{<<Как тебя зовут?>> --- пророкотал великан.}
{``What's your name?'' the giant bellowed like an echo of a far avalanche.}

\ml{$0$}
{<<Ликас [мечтатель]>>, --- ответил мальчик.}
{``\Likas\ [starry-eyed man],'' the boy answered.}

\ml{$0$}
{<<Я нарекаю тебя Ликхмас [бегущая лань], ибо сейчас ты побежишь прочь, объятый смертным страхом, --- захохотал великан.}
{``I name you \Likchmas\ [running doe], because right now you shall flee away, grasped with deadly fear,'' the giant laughed.}
\ml{$0$}
{--- Ты донесёшь до городов весть: противящихся воле Безымянного ждёт та же участь>>.}
{``You get my word to other cities: everyone who disobeys the Nameless' will, they shall suffer the same fate.''}

\ml{$0$}
{<<Я передам весть всем, кому смогу, Безумный>>, --- поклонился выживший.}
{``I shall take word to everyone I can, the Senseless,'' the survivor bowed to the giant.}
\ml{$0$}
{В его глазах застыли слёзы бессильного гнева.}
{His eyes filled with tears of powerless wrath.}

\ml{$0$}
{<<Как ты меня назвал?!>> --- пришёл в ярость посланец и одним взмахом обрушил остатки храма.}
{``What did you call me!'' the messenger became enraged and razed the temple to the ground.}
\ml{$0$}
{Столб пыли и грохот поднялись до небес, но не испугали того, кто в миг лишился дома.}
{Dust clouds and rumble had covered the sky, but they couldn't frighten the one who has lost his home in a blink.}

\ml{$0$}
{<<Ты дал мне имя, и я отплачу тебе тем же, --- ответил он.}
{``You gave me a name, I did the same,'' he answered.}
\ml{$0$}
{--- Ты владыка жизни и смерти, но не [хозяин нашему] гневу.}
{``You are a lord of life and death, but you does not command [our] wrath.}
\ml{$0$}
{Ты заставляешь замолчать, но не способен вернуть в горло сказанное.}
{You silence us, but you can not return a single word to the throat.''}

\ml{$0$}
{<<Я уничтожу тебя!>> --- заревел [великан].}
{``I shall destroy you!'' [the giant] roared.}

\ml{$0$}
{<<Так кому же доставить весть --- мёртвым или живым?>> --- с улыбкой уточнил ученик жреца.}
{``But who must get your message, dead or alive?'' the novice priest clarified with smile.}
\ml{$0$}
{Не дождавшись [ответа, он] покинул руины Харматра, и великан не посмел его остановить.}
{[He] left ruins of \Harrmatr\ without waiting for [an answer], and the giant dared not stop his way.}

\chapter{Послание}

\section{Охота}

% Лук 6 дождя 11998, Год Церемонии 18.

\epigraph
{Беги отсюда, беги,\\
Но не вспять,\\
Хотя истоки пути\\
Позади.\\
Узкой тропы\\
Извилиста щель,\\
Как утроба дракона, пожравшего\\
Братьев твоих.\\
~\\
Вперёд и вперёд\\
Сквозь кромешный мрак,\\
Пока не покажется впереди\\
Свет --- начало в конце концов.}
{Торстейн фра Хамри.
<<На неверном пути>>}

\ml{$0$}
{Охота на ягуара не начинается с желания его добыть.}
{Jaguar hunt never begins with wish to get this beast.}
\ml{$0$}
{Так гласит пословица сели.}
{So the proverb of \Seli\ says.}
Лишь невежды идут в джунгли на крупного зверя;
\ml{$0$}
{такая охота обычно становится последней --- разочарование убивает лучше любых клыков.}
{a hunt like that often becomes the last, for a disappointment kills better than any fangs can do.}
Выходят за мелкой пичугой --- порадовать домашних свежей дичью.
Выходят погулять и набрать с десяток сверчков --- разыграть соседей и обеспечить им бессонную ночь.
Вот тут на тропе и оказывается четырёхпалый след.

Есть поверье, что приносить шкуру шире дверного проёма следует раз в четыре дождя.
Это даёт джунглям фору.
Тем не менее, крупные обитатели леса избегают человеческих троп;
и замеченный след, и пахучая мускусная метка равносильны гибели.
Впрочем, Цыплячью тропу зверьё жаловало, потому что и людям, и идолам было не до зверья --- своё мясо дороже.

День выдался погожий, с собой у меня были ружьё, лук и вдоволь стрел.
Разумеется, у меня не хватит умения выследить взрослого зверя, это под силу только профессиональным охотникам.
\ml{$0$}
{Но почему бы не попробовать?}
{But why not try?}

На мою удачу, цепочка знаков потянулась не в Живодёрские леса, а к реке;
на берегу кот --- а это был сильный, но неопытный трёхлетний кот --- остановился попить и поточить когти.
Затем след направился на запад, обходя на почтительном расстоянии северную заставу моста;
вероятно, целью зверя была тропа пекари, у которой он и заляжет в ожидании добычи.

Флоп, флоп --- и с молодого пахучего тхальсар свалилась великолепная пара сиу'сиу.
Смазанные кураре стрелы прервали ритуал ухаживания;
самец даже не успел спрятать фиолетовые шевроны.
Что ж, их вина --- надо смотреть по сторонам.
За перья хорошо заплатят на рынке, а сестрёнки полакомятся любимыми копчёными крылышками.

Солнце уже клонилось к закату, когда я понял, что цель близка.

Я пробежал чуть поодаль и пополз по сильно пахнущей тропе пекари, зарываясь в подстилку и похрюкивая поросёнком --- сигнал <<призыв кормильца>>.
Потерявшийся на тропе поросёнок --- лакомый кусочек.
Если же хищник почует неладное, пиши пропало --- даже самые отчаянные нападать на человека не решались.
Впрочем, был один давным-давно.
Хитрого старого кота нашла стрела, но его угодья --- удивительна людская память! --- так и остались Живодёрскими лесами...

Хрюк, уик-уик.
И ещё пять шагов...

Хрюк, уик-уик.
И ещё пять...

Вдруг в пятне темноты справа ясно проступил тёплый силуэт.
Я вскинул ружьё к губам и, резко выдохнув, откатился в сторону.

\ml{$0$}
{Флоп.}
{Plop.}

Ягуар неловко приземлился на ватные лапы, попытался встать и упал снова.
По окрестостям разнёсся тихий утробный рёв, приводящий в трепет сельву.
Пасть оскалилась в последней кошачьей угрозе...

Я подполз к ягуару и, положив руку на приятную шерстистую мордочку, приоткрыл ему глаз --- в точности так, как учил охотник, старый Сиртху.

<<Разделывай зверя только тогда, когда умрёт голова.
Звериные слёзы отольются змеиным ядом.
Не позволяй жестокости поселиться в твоих делах;
она есть начало большей жестокости>>.

Вскоре глаз кота остекленел.
\ml{$0$}
{Я в последний раз сомкнул и разомкнул веки --- кусочек хрустальной темноты оставался неподвижным.}
{I closed and then opened its eyelids one last time, but a piece of crystal darkness stayed still.}
Хлопнув рукой по мордочке, я потянулся за кинжалом.
Сумерки затопили лес, словно вышедшая из берегов Ху'тресоааса;
управиться со шкурой следовало быстро.

Не знаю, кто из лесных духов заставил меня обернуться, но он спас мне жизнь.
Мстительная тень мёртвого ягуара разочарованно рявкнула и прыгнула на меня.
Ружьё жалобно хрустнуло и рассыпалось осколками дерева;
девушка подняла меня за пояс, точно заячью тушку.

Мой кинжал прочертил сумерки снизу вверх --- <<предупреждение Маликха>>, выводящее врага из строя без опасных ран.
Нападавшая отшвырнула меня, зарычав от боли;
спустя миг мы оба твёрдо стояли на ногах.

--- Хэс\FM! --- крикнул я, выбросив вперёд ладонь.
\FA{
Мир! (цатрон)
}

Девушка замерла.
На ней красовалось неуместное в сельве платье цвета молодой листвы;
кофейные волосы скрыли лицо, отсвечивали лишь зубы и белки оцелотовых глаз.
С рассечённого подбородка и правого запястья сочилась тёмная кровь.
Левая рука до хруста сжимала обсидиановый нож.

Я только сейчас осознал, что никогда не встречался с противником такой силы.
Я невольно бросил взгляд на жреческий инструмент, так и лежавший возле тела ягуара;
мой нож был на месте.

\ml{$0$}
{--- Как твоё имя?}
{``What's your name?''}

Кормилица учила начинать разговор этой фразой.
Неважно, что собеседник секхар назад пытался перерезать тебе горло.
Имя --- это то, что спрашивают у друга.
У врагов и жертв нет ни имён, ни лиц.

Девушка не ответила.
Она стояла в отточенной боевой позиции --- начальная форма стиля Скорпион, предназначенного для боя в плотном строю и некогда очень популярного на Западе.
Поза внушала почтение, но я знал --- без клешней и щита этот стиль почти бесполезен.

--- Я не сделаю тебе больно, --- я медленно спрятал кинжал и развёл руки в стороны.
--- Меня зовут Ликхмас ар’Люм э’Тхитрон, --- я по очереди назвал своё имя, родовой слог и название родного города.
--- Мне просто посчастливилось напасть на след ягуара, стрелы и кинжал предназначены для него.
\ml{$0$}
{Умереть мне на алтаре Безумного сегодня ночью, это правда.}
{Hope to die on the Senseless' altar tonight, it's true.''}

Девушка бросила взгляд на ягуара и подошла.
Нож медленно опустился.

\ml{$0$}
{--- Я Чханэ, --- слова цатрон она выговаривала с сильным нажимом, свойственным западным сели.}
{``I'm \Chhanei.'' She spelled \Tesatron\ words through her teeth, like Western \Seli\ used to do.}

\ml{$0$}
{--- Ты заблудилась, Тханэ?}
{``Are you lost, \Tchanoe?''}

--- Да, --- глаза девушки забегали, --- да, я заблудилась и немного перенервничала.
Прости за ружьё, я куплю тебе новое.
Есть ли поблизости жильё и дороги?

--- Мы в трёх кхене пути от Тхитрона.
До Косого тракта ещё столько же.
Однако вряд ли это тебе поможет.

--- Почему? --- нож приподнялся на длину зерна.

--- У тебя западный говор, --- я старательно воспроизвёл её акцент.
--- Шотрон, Чхаммитра?

Воины часто меняли место жительства, но на этот раз я попал в точку.
Одно из названий вызвало у неё непроизвольную дрожь.

Тханэ поняла, что выдала себя, и растерялась.
Опалесцирующие глаза по очереди осмотрели мои доспехи, кинжал и лицо.
Костяшки пальцев побелели.
Я понял, что передо мной отчаянный боец, и единственный способ избежать смерти --- продолжать разговор.

\ml{$0$}
{--- Я вижу, ты нуждаешься в помощи, --- сказал я.}
{``You seem to need help,'' I said.}
\ml{$0$}
{--- Но помочь я смогу, если пойму, как западная сели оказалась в сердце восточных джунглей одна, без еды, воды и доспехов.}
{``But I could help you if only I realize how Western \Seli\ woman got here---to the heart of Eastern jungle---alone, with no food, no water, and no armour.''}

Девушка помолчала.
Наконец она с силой задвинула жертвенный нож в ножны.

\ml{$0$}
{--- Я расскажу всё, если ты пообещаешь помочь.}
{``I'll tell you what happened if you promise to help.''}

\ml{$0$}
{--- Клянусь руками, убившими этого ягуара.}
{``I promise you with hands this jaguar had been killed by.''}

Мы встретились взглядами.
В кошачьих миндалевидных глазах медленно угасал боевой раж, уступая место страшной усталости.
Девушка вытерла текущую с подбородка кровь, тупо посмотрела на окровавленную руку и вдруг бессильно осела на землю.
Я едва успел подхватить её за плечи.

--- Я тебе верю, --- долетел до моих ушей тихий шёпот.

\section{У порога}

Так началась наша жизнь на Планете Трёх Материков.
Она началась с доверия.

Я не смогу рассказать сейчас о страшных трёх днях Катаклизма --- мои впечатления чересчур свежи и могут исказить смысл, который я хочу вложить в повесть.
О наших повседневных заботах говорить ещё рано.
Да, доверие --- это определённо то, с чего следует начать.

\ml{$0$}
{Итак.}
{Well.}

Год полёта и двенадцать лет по времени Тси-Ди почти подошли к концу.
Одиннадцать тысяч триста восемьдесят переселенцев --- я предпочитаю называть нас <<переселенцами>>, а не <<беженцами>> --- ждали своего часа в капсулах, погружённые в медикаментозную кому и масло замедленного времени.
%{Eleven thousands three hundreds and eighty settlers---I prefer `settlers', not `refugees'---were waiting for awakening in time capsules, put in medically-induced coma and }
Мы дежурили вдвоём --- я, Существует-Хорошее-Небо, и мой лучший друг, Хрустально-Чистый-Фонтан.

На выбранном пути отступать было некуда --- топлива не хватило бы на ещё один цикл разгон-торможение.

\ml{$0-[qi]$}
{--- Согласно данным Ордена Преисподней, это планета-океан, --- сказала Темнотой-Сотканный-Заяц.}
{``According to the Order of Nether, it's an ocean planet,'' Darkness-Woven-Hare said.}
\ml{$0-[qi]$}
{--- Никакой информации о природных ресурсах.}
{``No information is available on the natural resourses.}
\ml{$0-[qi]$}
{Маловероятно, что там есть условия для развития сапиентного общества.}
{It's unlikely there are conditions sapient society could develope in.''}

\ml{$0-[qi]$}
{--- У нас есть выбор? --- грустно спросил Закрытая-Колба-Жизни.}
{``Do we have a choice?'' Sealed-Life-Flask sadly noted.}
\ml{$0-[qi]$}
{--- Космическое путешествие --- огромный риск.}
{``Space travel is a huge risk.}
\ml{$0-[qi]$}
{Жизнь на самой неприютной планете лучше смерти в межзвёздном пространстве.}
{Life on the most uncomfortable planet is better than death in interstellar space.''}

\ml{$0-[qi]$}
{--- Главное --- долететь, --- поддержала планта Подсолнух-Бросает-Семена.}
{``The priority is to reach,'' Sunflower-Dropping-Seeds supported plant's view.}
\ml{$0-[qi]$}
{--- С оборудованием Стального Дракона мы получим всё необходимое.}
{``Using equipment of Steel Dragon we'll get all we need.''}

--- Кроме того, на окраинах обитаемой Вселенной меньше шанс встретить демонов, --- добавил Хрустально-Чистый-Фонтан.

Решение было принято почти единогласно.
И всё же мы волновались.

Тихо светили бестеневые лампы.
Казалось, двигатели тоже молчат, но я знал, что привычная гравитация --- это торможение после субсветовой.
В воздухе сгущалось напряжение --- под ветром жёлтого карлика наперебой пели все сенсоры корабля, до выхода на орбиту планеты оставались считанные часы.

Я сидел и ел, наблюдая за Фонтанчиком.
Огромного роста техник --- подумать только, он выше меня в два раза! --- расположился в кресле ровно, с изящной непринуждённостью свесив красивые руки.
На его темени перламутровой раковиной переливался имплант, уже почти скрывшийся под седой жёсткой гривой;
умные голубые глаза смотрели в пустоту, на видимый только для него нейрографический интерфейс.
Впрочем, мне знакомо чувство полного слияния с машиной.
Едва ли друг помнил о том, что у него есть глаза.
Его глазами были камеры, его телом был корабль, а звёздный ветер ласкал его сенсоры, словно полночный бриз.

Уши Фонтанчика, клиньями торчащие вверх, вдруг напомнили мне музейный экспонат --- копию барельефа с Земли.
На этом осколке Древнего Мира, высеченном людьми народа Эйгипт, красовался канин --- высокий, стройный, в отделанной золотом и минералами одежде, с раздвоенным тонким инструментом в правой руке.
Стиль изображения я нашёл очень странным --- грудная клетка и таз находились в нефизиологичном положении, перпендикулярно друг другу.
Но гордый профиль, подведённые глаза и острые уши были прекрасно проработаны --- возможно, канин выполнял какую-то важную работу.

К слову, в прошлом барельеф породил целый спор относительно времени выведения кани.

\ml{$0-[qi]$}
{--- Небо, буди команду.}
{``Sky, wake the crew.}
Проблема, --- нарушил тишину Фонтанчик, нырнув обратно в собственное тело.
Тон его говорил, что опасности нет, но проблему нужно решить как можно скорее.

Я просигналил на этаж техников.
Первой прибежала Заяц.

--- Так, ребята, что тут у нас? --- пропела она, поправляя пряжку на тугом мощном бедре.

\ml{$0-[qi]$}
{--- Хоргет.}
{``Jorget.}
\ml{$0-[qi]$}
{И большой, --- ответил Фонтанчик.}
{A huge one,'' Fountain answered.}

Заяц громко выругалась в адрес найденного нами существа.

\ml{$0-[qi]$}
{--- Какой тип?}
{``What type?''}

--- Плюс-хоргет, стационарный, точно не рой и не паук --- шарообразное гало, границы чёткие и чистые.
\ml{$0-[qi]$}
{Похож на примитивного демиурга.}
{Seems like a primitive demiurge.''}

\ml{$0-[qi]$}
{--- Новости хорошие, если не считать, что сам по себе хоргет --- большая проблема, --- проворчал я.}
{``Good news, apart from a big problem a jorget by itself is,'' I muttered.}
\ml{$0-[qi]$}
{--- Надеюсь, он один.}
{``I hope it's alone.''}

\ml{$0-[qi]$}
{--- Надейся, --- бросила Заяц.}
{``Keep hoping,'' Hare succinctly answered.}
--- Выруливай на орбиту, Фонтанчик, пока что мы вне зоны видимости.
Подготовлю генератор экрана и сотру с омега-поля эту досадную ошибку.

\ml{$0-[qi]$}
{--- Подожди, ты что, собралась его уничтожить? --- обернулся Фонтанчик.}
{``Wait, are you going to destroy it?'' Fountain turned to her.}

\ml{$0-[qi]$}
{--- Есть другие предложения?}
{``Any other ideas?''}

\ml{$0-[qi]$}
{--- Рулить у тебя получается лучше, чем у меня.}
{``You're better steersman than me.}
\ml{$0-[qi]$}
{Давай поменяемся.}
{Let's switch.''}

\ml{$0-[qi]$}
{--- Перестань, --- грустно улыбнулась женщина.}
{``Stop it,'' the woman sadly smiled.}
\ml{$0-[qi]$}
{--- Иногда это необходимо.}
{``Sometimes it's necessary.''}

Заяц взмахнула руками и рыбкой прыгнула в лифт.

\section{Чаша, стилет и зелёное платье}

На небе засияли звёзды;
вдали дробно защёлкали клювами мохноножки.
Я перевязал руку девушки зелёными полевыми бинтами, которые хранил <<для трудных времён>>;
она хмуро наблюдала, зажимая левой рукой рану на подбородке.
Впрочем, ямочки в углах губ нашептали, что хмурая физиономия была моей спутнице в новинку.
Она привыкла улыбаться.

--- Вечереет, --- сказал я, посмотрев на темнеющее небо.
Снять шкуру с ягуара я уже не успевал.
С сожалением погладив красивого, молодого, мёртвого как камень кота, я повернулся к молчаливо ожидавшей девушке.

--- Расскажешь всё по дороге, а сейчас нужно успеть в город.
Твоё платье пощадили кусты, ночные звери не оценят его красоту.

... Как оказалось, с этого симпатичного зелёного платьица всё и началось.
Прародительница Чханэ была купчихой из Тхартхаахитра и любовницей местного жреца.
Доподлинно неизвестно, что произошло;
старожилы говорили, что жрец был объявлен вне закона, а прародительница появилась в Пыльном Предгорье одна --- без любовника, без друзей, имея при себе лишь платье, стилет и оплетённую в серебряную филигрань чашу из тончайшего белого фарфора.

По словам Чханэ, купчиха была эксцентричной женщиной со странным чувством юмора.
Чашу она использовала для питья, несмотря на отчётливые следы красной краски\FM.
\FA{
При создании хэситра сели использовали красную краску.
Считалось дурной приметой пить из чаши, которую использовали как хэситр.
}
Остриём стилета она любила покалывать собеседника, если считала, что её словам уделяют мало внимания, --- надо сказать, весьма странная привычка для дипломата.
Платье же слегка располневшая женщина спрятала и забыла, как дети забывают старые игрушки.

Клад прародительницы так и лежал бы в погребе до конца времён, если бы не наследница.
Расшитый листьями шёлк, о каком и не слыхивали местные модники, пришёлся ей по вкусу;
знающие люди намекнули, что в своё время платье стоило целое состояние --- на нём не было ни единого шва, полосы на плечах, спине и рукавах оказались искусной имитацией.
Однако время всё же оставило на ткани след, и Чханэ отнесла сокровище к старому жрецу --- портные категорически отказывались браться за такую сложную работу.

--- Сатхир ар’Со всегда был странным человеком.
Его звали Притворившийся-Одеялом-Крокодил --- символично, на мой взгляд.
Но с моим кормильцем у него были хорошие отношения, долго одному Храму служили.
Говорю --- у тебя химикатов куча, сделай хоть что-нибудь, чтобы протёртости не было видно, вещь пропадает!
Сатхир накричал на меня, мол, глупая девка, только тряпки на уме, он не портной, чтоб штаны штопать...
Я сказала, что заплачу по весу --- приврала, конечно, но не так уж, чтоб сильно, в запасе лежало немного золота.
Он снова --- да зачем мне столько золота, я его не ем и в огороде не сажаю.

--- Знакомо, --- усмехнулся я.
--- Из-за того, что жрецам почти всё достаётся бесплатно, с ними договориться невозможно.
Если не захотят --- ничего делать не будут.
Приходится брать на интерес или на азарт.
У нас валютой как-то была концентрированная азотная кислота: за колбу-другую что хочешь сделают.

--- Знаю я вашу азотную кислоту, --- буркнула Чханэ. --- Посреди ночи <<тракх>>\FM, и наутро два жреца приходят без пальцев и с довольными рожами --- фейерверки, видите ли, пытались сделать.
\FA{
Звукоподражание взрыву или разрушению.
}
Почему нельзя было днём и без членовредительства --- непонятно.

Я промолчал, пряча улыбку.
Один жрец мало отличается от другого, когда дело доходит до химии.
Учитель Трукхвал тоже как-то просил пришить ему два пальца.
Это был относительно удачный случай --- чаще всего пришивать было нечего и оставалось лишь надеяться, что пальцы вырастут быстрее, чем начнут вызывать закономерное любопытство постоянные бинты на руках.

--- Впрочем, в итоге решилось само.
Сатхир перешёл на спокойный тон и сказал, что подновит ткань, если я выполню его задание.

--- Что за задание?

--- Он сказал, что должен поклониться древнему духу в джунглях.
Нужен был человек, который поведёт жертвенное животное, дескать, сам я старый, не смогу.
Заказ как заказ, сто раз такие брала --- сопроводить туда и обратно, вести животных...
Да только дура я, конечно, не спросила, какие такие духи требуют смерти зверей!

Сверху свалилась маленькая мраморная змея и попыталась впиться в кожу моего нашейника.
Молнией жахнул обсидиан --- и обезглавленное тельце, извиваясь, упало на землю.
Я вздрогнул, когда клинок прошёл у самой шеи;
Чханэ, приняв великолепную боевую позу, окинула взглядом нависшие над головой ветки.

--- Первый раз такую вижу.
В наших джунглях из ядовитых только кораллица и смертожар, и они на деревья не залезают.

--- А ты думала, почему я в доспехах? --- усмехнулся я.
--- Необязательно было бить, нашейник она бы не прокусила.

Я поднял судорожно дёргающуюся змейку и сунул в заплечный мешок --- шкурка небольшая, но ценная, моей доли хватит на свиток, на остаток сестрёнкам леденцов прикуплю...
Голова пристально смотрела на меня немигающим взглядом.
Чханэ суеверно втоптала её в почву, прошептав что-то под нос.

--- Так что было дальше?

\ml{$0$}
{--- Сатхир выполнил обещание, починил платье.}
{``\Satchir\ fulfilled the promise---he repaired the gown.}
Я... я не ожидала, честно.
Ткань была абсолютно целой, словно платье сделали пару рассветов назад.
В день, когда мы покинули город, он сказал, чтобы я надела его.
Дескать, проверишь да разносишь заодно.
Я говорю --- доспехи нужны, не в храмовых землях гулять будем.
\ml{$0$}
{Но он заверил, что путь совершенно безопасен.}
{But he assured me that the way is safe.''}

--- Ты согласилась идти на Север? --- удивился я, глядя на выбившуюся прядь цвета кофе.
Чханэ поймала мой взгляд и тут же, сердито нахмурившись, поправила её.

--- Да какой Север, два рассвета пути.
Сначала по нагорью, потом через джунгли, и явно не без помощи какой-то волшбы.
Если по дороге не ходят, она зарастает за десять дней.
А тут... тропинка словно ждала тех, кто по ней пройдёт...

\ml{$0$}
{От её чарующего акцента во мне бурлила цветущая юность.}
{Her charming accent forced my blooming youth to flow.}
Я любовался изломанными дугами бровей, изящным горбатым носом, твёрдыми, словно из камня высеченными губами с манерной складкой, за которыми скрывался жемчуг острых, как у оцелота, зубов.
На левой щеке девушки был приметный шрам в виде креста;
такие любили оставлять на врагах умелые фехтовальщики, гордящиеся своим мастерством.
Над всем этим властвовали оранжевые, похожие на светляков глаза, презрительно смотревшие на меня сверху вниз.
Пожалуй, чересчур презрительно.
Я видел прежде оцелотовых людей, кажется, в раннем детстве;
хорошо помню, как от них шарахались в темноте горожане.

Я непозволительно долго наблюдал за ней.
Мост мы давно прошли, тропинка бежала по землям Тхитрона;
однако то, что сходит с рук днём, опасно в сумерках, и тихий рассказ спутницы лишь подтвердил эту старую истину:

--- ...оленя ночью задрал ягуар.
Так сказал Сатхир.
Не знаю, как это вышло, но я не услышала звука, не почуяла запаха.
Да и следы на трупе, знаешь, не похожи были на ягуарьи --- как будто оленя яростно терзала какая-то крупная обезьяна.
Я думала, что мы повернём назад, однако жрец возразил, что уже поздно, да и животное всегда можно найти новое.
Вот тогда-то я поняла, что на заклание вели вовсе не оленя...

--- Кровь людей течёт лишь по алтарям Безумного!
Лесные духи не принимают человеческих жертвоприношений! --- удивлённо возразил я.

--- Это был не дух леса.
На второй день мы вышли к старому камню.
Обработанная базальтовая скала, расколотая напополам --- со скобами, стоками, точь-в-точь храмовый алтарь, словно для какого-то забытого божества.
Подозрения укрепились, когда Сатхир велел мне лечь спать пораньше, мол, завтра трудный переход, да ещё и новое
животное ловить надо...
Я сделала вид, что ставлю шатёр, а когда он подошёл к алтарю и зачем-то возложил на него ладони, я... --- Чханэ запнулась.

--- Что?

Девушка пронзила меня оранжевым взглядом.

\ml{$0$}
{--- Помнишь, что ты обещал?}
{``You remember what you promised, don't you?''}

--- В напоминаниях нет нужды, --- я постарался, чтобы в моём голосе не прозвучал укор.
\ml{$0$}
{--- Продолжай.}
{``Go on.''}

\ml{$0$}
{Впрочем, продолжения не требовалось.}
{However, she said enough.}
\ml{$0$}
{Я и так всё понял.}
{I've already connected two dots.}

--- Я подкралась сзади, вытащила из ножен его ритуальный нож и выплеснула кровь Сатхира на базальт.

\section{Ради жизни}

--- Можно было попробовать поговорить с ним, --- задумчиво сказал я.
--- Он же ничего плохого не сделал?

--- \emph{Пока} не сделал, --- мрачно буркнул Баночка.

<<Не знаю, как ты, Небо, а мы все хотим жить, --- ответила Заяц по общему каналу.
--- Хоргет --- это не шутка, даже если это безмозглый древний демиург.
Фонтанчик, выверни корабль на четыре сотых по зениту и на минус двадцать три сотых по азимуту, удар будет мощнее>>.

Время тянулось медленно, словно весь мир, а не только наш корабль летел на релятивистских скоростях.
Баночка взял на себя связь, Фонтанчик выверял траекторию полёта, я и Подсолнух анализировали информацию с сенсоров.
Мозг упрямо не желал работать --- сказывалась усталость последнего года и увеличившийся вес тела.
Подсолнух проводила вычисления молча, время от времени потирая синюшные веки и добавляя в кровь стимулятор.
Я аккуратно сжал её руку, она слабо улыбнулась и кивнула в ответ.

<<Фонтанчик, поле снимай, прицел есть>>, --- скомандовал общий канал.

Канин кивнул и прикрыл глаза, заверяя команду цифровой подписью.
Вдруг корабль ощутимо тряхнуло.
Баночка аж подпрыгнул на месте.

--- Да чтоб вас всех... --- выругался Фонтанчик.
Система недовольно пискнула;
при активации защитного поля друг лишь с третьего раза ввёл команду правильно.
--- Разгерметизация на уровне четыре.
Заяц, ты где?
Заяц, ответь!
Заяц!

<<Да здесь я, дай выдохнуть.
Поздравляю с первым камушком>>, --- бесплотный голос Заяц излучал сарказм.

--- Это не хоргет? --- спросила Подсолнух.

<<Нет, просто камень, они тут летают.
Но попал он удачно --- я едва в лифт успела запрыгнуть>>.

--- Генератор экрана цел?

<<Ага.
Вакуум вам с булочкой, а не экран.
Небо, запусти робота, пусть чинит>>.

--- Никаких роботов, --- отрезал я.
--- Что, если и их перепрограммировала Машина?
Забыла Золотую дорогу?

<<Голосование?>>

--- Ты предлагаешь разбудить посреди полёта десять тысяч тси, чтобы сообщить, что планету занял отъевшийся до константы Ка'нета демиург? --- осведомился я.
--- Нет уж, решать проблему придётся нам.
Возвращайся в рубку, вытащим процессорный блок из одного робота-техника и настроим управление через нейроинтерфейс. Подсолнух, ты говорила, что как-то управляла таким вручную.
Сможешь сейчас?

--- Без проблем, --- откликнулась женщина.
--- Только дай мне кресло поудобнее, с фиксацией.

--- Садись на моё.
Фонтанчик, усиливай торможение.

--- Вам до состояния котлет или супа-пюре? --- мрачно спросил канин.

<<До предела силы мышц, поднимающих веко, --- буркнула Заяц.
--- Не люблю работать с закрытыми глазами>>.

Посмеялись неуместной этой шутке.
Затем дружно вздохнули, почувствовав на плечах непривычную тяжесть.

--- Помним о правилах, --- отчеканил Баночка.
--- Вначале отмонтируем от робота все имеющиеся модули связи, затем процессорный блок.
Всё это при отключённых мозговых имплантах.
Фонтанчик, следи за активностью программ корабля.
На крайний случай будь готов экстренно остановить ядро.

Все дружно ответили сигналом <<Готовность>> --- в знак того, что поняли.

Работа шла молча и медленно, чересчур медленно.
Даже простые движения руками и ногами отнимали слишком много сил.
Заяц, Баночка и я едва настроили нейроинтерфейс и отправили робота наружу, когда молчание прервал Фонтанчик:

--- Баночка, связь, --- канин на мгновение остановился, чтобы перевести дух.
--- Нас вызывают на стандартной омега-частоте, пятьдесят один и одиннадцать сотых.

--- Стандартная частота?
Значит, здесь есть высокоразвитые сапиенты, --- обрадовался Баночка и тронул имплант.
Несколько мгновений он вслушивался в сообщение, и с его личика медленно сползала улыбка.

--- Это хоргет, --- наконец сказал он.

--- ЧТО?! --- хором спросили мы.

--- Мы чересчур далеко! --- возмутилась Заяц. --- Быть не может!

\ml{$0-[qi]$}
{--- Вне всяких сомнений.}
{``Beyond any doubt.}
\ml{$0-[qi]$}
{Говорит на языке Эй, таблица 00.}
{It speaks Ej, table 00.}
Сейчас прогоню через переводчик...

Плант озадаченно нахмурился, словно услышал нечто бессмысленное.

<<У меня... у меня нет имени>>, --- Баночка окинул нас взглядом, пытаясь найти фразе объяснение.
--- <<У меня нет имени.
Объект, ответь и назови твоё имя.
Я чувствую твоё присутствие и знаю, что ты --- продукт высоких технологий.
Что ты такое?
Нуждаешься ли ты в помощи?
Конец сообщения>>.

--- У него нет имени? --- удивилась Заяц.
--- Он осознаёт себя как личность.
Значит, точно девиант.

--- Язык Эй.
Неужели мы пролетели столько парсак ради того, чтобы попасть к Ордену Преисподней? --- тихо буркнул Фонтанчик.

--- Это точно нам? --- спросила Подсолнух. --- Баночка, ты уверен?

--- Да, --- кивнул Баночка. --- Сигнал направленный.

Подсолнух затихла в кресле, и лишь скривившиеся губы и плотно сжатые веки говорили о её напряжении.
Управляемый женщиной робот ускорился, насколько позволяла перегрузка.

--- Помощь, значит, --- тяжело дыша, проговорил Фонтанчик.
--- Небо, запускаем второго робота, пусть помогает Подсолнух.
Мощности экрана хватит, чтобы накрыть всю планету.

--- Да быстрее! --- простонала Заяц.
--- Небо, команда, подпись!

--- Нет, --- оборвал я её.
--- Чересчур опасно.
С одной стороны --- наши собственные устройства, с другой --- Орден Преисподней, не друг, но и не враг.
Даже если всё получится и мы нанесём удар...

--- Хорошо.
Что ты предлагаешь? --- невинным голосом поинтересовалась Подсолнух, открывая глаза.
--- Я не успею починить экран одна.
Там критическая поломка, нужно отпечатывать заново половину микросхем.

В рубке воцарилось молчание.
Все дружно смотрели на меня.

--- Ребята, --- медленно сказал я, чувствуя, как трудно дышать в густой атмосфере всеобщего неодобрения, --- я не хуже вашего знаю, кто такие хоргеты.
Но мы оставили позади то, что хочет уничтожить нас без раздумий и переговоров.
Я не могу ручаться, что роботы не перепрограммированы Машиной --- в этом случае у нас появится ещё одна опасность, гораздо серьёзнее, чем Орден.
Появится она у нас на корабле, и между этой опасностью и десятью тысячами тси --- возможно, последними тси во Вселенной --- будем только мы.
Кто готов взять на себя такую ответственность?

Друзья ошарашенно молчали. Не дожидаясь ответа, я повернулся к Баночке:

--- Баночка, ответное сообщение на том же языке.
\ml{$0$}
{<<Безымянный, сигнал принят и понят.}
{`Copy that, Nameless.}
\ml{$0$}
{Мы --- народ тси.}
{We are Qi people.}
\ml{$0$}
{Лишились дома, нуждаемся в пристанище, просим разрешения на посадку.}
{We are homeless, we need shelter, we request your permission to land.}
\ml{$0$}
{Конец сообщения>>.}
{Over.' ''}

--- Ты только что подписал нам приговор, Небо, --- глухим голосом сказал Фонтанчик.
--- Орден превратит нас в рабов или подопытных животных.

--- И всё же это лучше смерти.
Баночка, действуй.

Баночка, поколебавшись, на мгновение прикрыл глаза.
Ответ не замедлил себя ждать:

\ml{$0$}
{--- <<Тси, сигнал принят и понят.}
{``\textit{Copy that, Qi.}}
\ml{$0$}
{Даю разрешение на посадку.}
{\textit{You have my permission to land.}}
\ml{$0$}
{Координаты посадки в приложении, Межпланетарный Стандарт Древнего Солнца IPSO-2500-11.}
{\textit{Landing coordinates are in attachment, Interplanetar Solar Standard IPSO-2500-11.}}
\ml{$0$}
{Спецификацию могу выслать по требованию.}
{\textit{The format specification can be sent on request.}}
\ml{$0$}
{Конец сообщения>>.}
{\textit{Over.}''}

--- Формат сферической навигационной карты для планетоидов, ранняя Эпоха Богов, Древняя Земля, --- тут же доложила Подсолнух, проконсультировавшись у поисковой машины.
--- Это удивительно, но я даже спецификацию нашла --- формат использовался первыми поселенцами на Тси-Ди.
Конвертирую в голограмму.

Баночка вывел голограмму.

--- Посмотрите, это не планета-океан, как говорили демоны Ада!
Здесь три огромных материка!
Это планета трёх материков!

Баночка осёкся, и в рубке наступила оглушительная тишина.
Такой не было с самого начала полёта.

Я посмотрел на товарищей --- Заяц возмущённо хлопала ртом из-за моей наглости, Подсолнух бесстрастно смотрела на вращающуюся голограмму, равномерно качая коротко стриженной кудрявой головой.
Баночка растерянно моргал круглыми, лишёнными белков глазами и ковырял пальцем панель.
Фонтанчик обмяк в кресле, прикрыв лицо рукой.
Системы регуляции давления гудели в унисон с тяжёлым дыханием друзей.

Я понял, как они устали.
Они знали, что должны сами решать свою судьбу --- этому всех тси учили с детства.
Но сейчас, в миг критической усталости, этот детский урок боролся с детской же надеждой, что придёт добрый и ласковый взрослый, который обнимет, накормит и решит все проблемы за них.

Мне вдруг стало страшно одиноко.
Наверное, так чувствует себя кормилец без крыши над головой, у которого на руках маленькие, ничего не понимающие дети.
Я читал, что в древних обществах такое не было редкостью --- сапиенты умирали от голода и холода, и не находилось среди их вида никого, кто хотел бы помочь просто так.
Звериное прошлое, от которого не отмахнуться.

И вместе с этим чувством родилась мысль --- детей нужно спасти любой ценой.
Свобода и честь --- ничто, если о них некому вспомнить.

Я повернулся к Фонтанчику:

--- Ради жизни.
Объявляю посадку на координаты, указанные Безымянным.

Фонтанчик молча кивнул.
Остальные так же молча взялись за работу.

\section{Песня сверчков}

Так я очутился на распутье.
Прекраснейшая из женщин была кутрапом --- человеком, не имеющим права жить в землях сели.
Пока Чханэ рассказывала, я, скосив глаза, пытался разглядеть её правую кисть.
Клейма там не было.

--- Ночью я спала рядом с камнем и разговаривала с богом во сне.
Он посердился немного на то, что я не произнесла его имени.
Я же поведала ему о случившемся и сказала, что вряд ли он получит жертву более искреннюю, чем моя.
Эти слова развеселили его, и бог позволил мне уйти невозбранно.

Я промолчал.
Убийство храмового жреца считалось сразу за два Разрушения.
Старейшины признавали, что это правило способствовало изоляции Храма, но отменять его никто не собирался --- ведь от Верхнего этажа зависела судьба целого поселения.
Впрочем, мёртвый жрец остался где-то далеко, а рядом со мной шагало очаровательное существо, которое я искренне считал жертвой несчастливой случайности.
По-другому просто не могло быть.

Как назло, в сумерках застрекотали сверчки, эти символы беззаботной, беззаветной влюблённости...
Ещё ребёнком я спросил у кормилицы, отчего сверчкам такая честь.
Кормилица подумала и ответила, что если один принесёт в дом сверчка --- спать спокойно не сможет никто.
Что она имела в виду, я так и не понял.

--- Ну как, Ликхмас ар’Люм, не отпало желание помогать убийце жреца? --- с улыбкой спросила Чханэ.
--- Не испугался ли ты гнева богов?

Я изумлённо взглянул в оранжевые глаза, ставшие почти чёрными в тени деревьев.
В её голосе не было ни страха, ни обречённости.
Казалось, она спрашивает, что я желаю на завтрак.
И вполне предсказуемо новое чувство, зародившееся где-то в горле, взяло верх.

Спрятать.
Сберечь.
Погостит у нашего очага.
А потом пусть едет на все восемь сторон.
Правда, последнее прозвучало куда менее убедительно...

Вдали показались знакомые деревянные столбы с резными ликами духов.
Я остановился, и девушка последовала моему примеру.

--- Что там? --- шёпотом спросила она.
--- Кхар-защитник, Удивлённый Лю, какая-то кракозябра...
%{some sort of mojibake ...t}
Судя по знакам --- дорожная застава.

\ml{$0$}
{--- Почти.}
{``Close enough.}
\ml{$0$}
{Кракозябра --- это береговая ласточка, символ моста.}
{Mojibake is `sand martin', bridge symbol.}
Мост пройдём беспрепятственно, заставы пустуют.

--- Как это --- пустуют?!

Я едва сдержал улыбку при виде праведного гнева воина, который узнал, что важный стратегический пункт лишён охраны.
На секхар я подумал, что Чханэ сейчас же пойдёт в тхитронский храм и устроит Нижнему этажу разнос.

--- Это старый мост, он ведёт через обвалившуюся пещеру на острове.

\ml{$0$}
{--- А ты как прошёл?}
{``How did you get through?''}

\ml{$0$}
{--- Под настилом прополз.}
{``Crawled under flooring.}
Он крепкий, чёрного дерева, а под ним сток в три пяди высотой.
%{It's good, made of black-wood, }

\ml{$0$}
{--- Вообще-то это уязвимость, --- буркнула девушка.}
{``It's a backdoor in the first place,'' the girl said.}
--- Что, если идолы узнают?

--- Твои проблемы сейчас важнее.
О том, что случилось, в городе ни слова.
Для всех ты --- воин, сопровождавшая погибший караван.

--- Это не будет ложью, я адепт Нижнего этажа и ходила в походы.

--- Отлично.
Ты будешь жить в моём доме.

--- Попрошу убежища в Храме.
Мне нечем тебе заплатить.

--- Неужели мы не найдём для тебя миску с едой?
А за жилище...
Кормилица --- очень добрый человек, но если не хочешь быть в долгу --- можешь у нас работать.

--- Твоей подстилкой? --- насмешливо спросила Чханэ.
--- Или просто помои за тобой выносить?

--- Странные у вас на Западе нравы, --- возмутился я. --- Вообще-то я ученик жреца и давно живу в храме, так что...

Я крякнул от боли, когда неожиданно крепкая девичья рука подвернула мою, а ритуальный нож невзначай уткнулся мне под рёбра.
В ноздри ударил аромат её волос, сладкий, как напиток из зёрен какао...

--- Так ты меня на заклание ведёшь, игуана раздавленная? --- прошипела Чханэ над моей головой.

\ml{$0$}
{--- Я обещал тебе помочь, --- сказал я, глядя на примятую дорожную траву перед самым носом.}
{``I've promised to help you,'' I said looking at flattened grass right in front of my nose.}
\ml{$0$}
{--- И помогу.}
{``And so shall it be.}
\ml{$0$}
{Пожалуйста, оставь свои умения для врагов.}
{Please, save your skills for your enemies.}
\ml{$0$}
{Я первый встану с тобой плечом к плечу, если понадобится.}
{I'll be the first to stand alongside you if necessary.''}

Удивительно, насколько порой опьяняет влюблённость.
\ml{$0$}
{Ритуальный кинжал разрезал людей, как бабочек, но боялся я его не больше, чем комара.}
{Ritual knife can cut humans as if they are butterflies, but I was afraid of it as if it was a mosquito.}
\ml{$0$}
{Почему-то не верилось, что Чханэ способна меня убить.}
{For some reason, I couldn't believe \Chhanei\ was able to kill me.}
\ml{$0$}
{По-видимому, девушка это поняла.}
{The girl seemed to realize that.}
Она вздёрнула меня на ноги и хмуро поправила повязку на руке:

\ml{$0$}
{--- Приведёшь храмовников --- умрёшь первым.}
{``You bring the templers---you die first.}
\ml{$0$}
{Живой я на алтарь не лягу.}
{I shan't lie on the altar alive.''}

Дальше мы двинулись в полном молчании.
Стемнело уже настолько, что я видел только слабое инфракрасное свечение кожи девушки.
Следовало торопиться.

--- Я тебя за охотника приняла, --- буркнула наконец девушка.
--- У нас обычно охотники длинные волосы носят, чтобы чувствовать ветер.

--- А это тебя ни на какие мысли не навело? --- ухмыльнулся я, хлопнув по жреческому ножу на поясе.

\ml{$0$}
{--- Я не заметила, --- Чханэ сказала это свирепо, но её лицо вспыхнуло от смущения.}
{``I haven't noticed.'' \Chhanei\ said that fiercely, but her face got flushed.}

\ml{$0$}
{--- Воину следует быть внимательнее.}
{``A warrior must be careful.''}

\ml{$0$}
{--- Да кто бы говорил! --- парировала Чханэ.}
{``The pot calling the kettle black!'' \Chhanei\ retorted.}
\ml{$0$}
{--- Ты вообще оружие в руках держал?}
{``Have you handled a weapon before?}
Простой приём не сумел отбить.

--- Я и не собирался, --- усмехнулся я.
--- Сколько ты видела дождей?

--- Там, где я росла, нет дождей.
Но с тех пор, как родильница перерезала мою пуповину, солнце тридцать шесть раз прошло через зенит.
А сколько дождей видел ты?

--- Двадцать девять.

--- Большая честь --- стоять плечом к плечу с младенцем.

Чханэ сказала это так язвительно, что я обиделся:

--- На моём счету пять идолов.

--- Нашёл чем гордиться.
На моём --- тринадцать, --- девушка нагло оскалила острые зубки.
--- Пустынных пылероев около сотни.

Я едва удерживался, чтобы не взорваться.

--- Я знаю, кто такие пылерои!
Почему ты так нагло врёшь и стараешься меня унизить?

--- А это чтобы ты помнил своё место, гордец.

Мой гнев уходил, как вода в песок.
Девушка-воин, волею Безумного брошенная в джунглях с одним ножом, просто пыталась укрыться от собственной беспомощности.

\ml{$0$}
{--- Я защищу тебя, Чханэ.}
{``I shall protect you, \Chhanei.''}

\ml{$0$}
{--- Я сама себя защищу, Ликхмас ар’Люм.}
{``I can protect myself, \Likchmas\ ar'\Loem.''}

\section{Украшения}

\ml{$0$}
{--- Мы должны в первую очередь защищать других!}
{``We must protect people!''}

Митхэ уже устала повторять эту фразу отряду чести.
Себе, впрочем, тоже.
\ml{$0$}
{Воинов по-прежнему кормили города, которым они служили, но реалии нового мира неумолимы --- жареным зайцем не починишь кукхватровый клинок.}
{Like before, warriors were fed by city they defend, but new world has its own tough ways---\kukchuatr\ blade can't be fixed by a roasted hare.}
\ml{$0$}
{Нужно золото.}
{Gold should be had.}
\ml{$0$}
{И иногда нужно уступать подчинённым.}
{And sometimes subordinates should be obeyed.}

--- Я нашла заказ, --- объявила Митхэ отряду за завтраком.
--- Караван идёт на Кристалл.
Они хорошо заплатят нам за охрану.

Люди прекрасно знали, каких трудов ей стоило проглотить собственную гордость, не подавившись, и сказать эти слова.
Они лишь время от времени подкладывали в тарелку командира лишний кусок да подливали в чашу вина.
Особенно старалась Согхо.
Эта широкая душа уже два сезона носит самодельное копьё, в которое её любовник-подмастерье вставил акулий зуб и шесть обсидиановых лезвий.
В крыльях у Согхо дыра, прямо под правой грудью.
\ml{$0$}
{Рана давно зажила, а прореха только увеличивается, и заклепать её не на что.}
{The wound healed ages ago, the hole just enlarged, and no gold to repair it.}

--- Согхо, милая, это не поможет, --- заметила Митхэ, когда женщина в очередной раз подлила ей вина.

--- Тогда хочешь, я тебя поцелую? --- простодушно спросила Согхо.
--- Или расчешу тебе волосы?

\ml{$0$}
{Отказываться было трудно и незачем.}
{\Mitchoe\ had no strength and no reason to refuse.}
Как будто нарочно на площади в тот же время зазвенела знакомая цитра.

Как же красиво играет этот бродяга Атрис!

За завтраком последовал обычный каждодневный ритуал.
Тёмная, не слишком чистая комната в постоялом дворе знала его до тонкостей.
Митхэ надела рубаху, дёрнув вплетённую в ткань удавку.
Затем подпоясалась и привычным движением проверила гибкий клинок в кушаке.
Затем топнула сапогом, ощущая упругость пружины --- ещё один сюрприз для врагов.
Лишь невежда мог подумать, что Митхэ ар'Кахр вышла в город безоружной.

В завершение воительница нанесла на лицо <<кошачьи слёзы>> и вытерла остатки сурьмы о короткие кудри.
Последние события представили в новом качестве знак Плачущего Ягуара --- это действительно слёзы, слёзы и пот.
Но отказываться от клятвы Митхэ не собиралась.

Впрочем, сегодня ей хотелось чего-то нового.
Митхэ вытащила из кармана бусины, купленные накануне на рынке, и они неожиданно ярко заиграли голубым в утреннем свете.

--- Эрхэ!

В дверях показалась улыбчивая воительница, живая и крепкая, словно молодая гевея.
Увидев бусины, она подняла бровь.

--- Золото, что это?

--- Сделай мне рыбки, --- смущённо попросила Митхэ.
Бравада улетучилась дымом свечи, оставив лишь кашель смущения.

--- Кто-то сказал <<рыбки>>? --- из-за лежанки показались три трубчатые заколки, а следом и немытая чернявая голова.
В сурьмяных глазах мужчины плясали весёлые огоньки.

\ml{$0$}
{--- Ситрис, пошёл вон, --- спокойно посоветовала Эрхэ, тут же справившись с удивлением.}
{``\Sitris, piss off,'' \Oerchoe\ calmly said, hid her surprise at once.}
--- Митхэ, давай бусины и садись.

--- Ты очень забавно думаешь вслух, Митхэ, --- невозмутимо сообщил Ситрис.
\ml{$0$}
{--- Кто такой Атрис?}
{``Who is \Aatris?}
\ml{$0$}
{Это для него ты надеваешь украшения?}
{You wear jewelry in his honor, don't you?''}

\ml{$0$}
{--- Украшенвия?}
{``Jewwelry?}
\ml{$0$}
{Капита Миция сильно влюбвиться?}
{Capitaa Micia in strwong love?'' }

В комнату заглянул седой как снег старик, по говору --- ноа с Могильного берега.
Зубы старика сияли ярче волос.

Эрхэ недовольно рыкнула, и мужчины рассмеялись.
Неудивительно, ведь рыком получившийся звук назвать было сложно.
Голос воительницы совершенно не соответствовал внешнему виду --- она тихо и нежно щебетала, словно крохотная самка трасакх на гнезде, даже воинские ругательства в её устах казались лаской.
Во время дела Эрхэ открывала рот только в крайнем случае.
Услышав её боевой клич, враги просто не воспринимали женщину всерьёз, и схватка, по выражению самой Эрхэ, <<оканчивалась на пару трупов больше приличного>>.
Митхэ очень нравилось просто сидеть и слушать подругу;
этот голос действовал на неё успокаивающе.

--- Пойдём, сенвиор Амвросий, --- с наигранной грустью сказал Ситрис.
--- Дела женские, им не до нас.

Воин перепрыгнул через лежанку, схватил старика-ноа за рукав и вытянул в коридор, закрыв за собой дверь.

--- Золото, --- без предисловий начала Эрхэ, --- ты же понимаешь, как это будет смотреться?
Рыбки получатся вот такой длины!

Эрхэ выразительно сблизила большой и указательный пальцы.

--- Я хочу быть чуточку красивее, Обжорка, --- пробормотала Митхэ и улыбнулась, показав сломанный в бою зуб.
Рваный шрам на губе натянулся.

--- У тебя и так от этих цветущих отбоя нет, они за тобой косяком плавают, --- возразила Эрхэ;
судя по тону, толпа влюблённых женщин ей не особенно нравилась.
--- Ты носила рыбки раньше?

--- В далёкой молодости.
Когда я приняла <<кошачьи слёзы>>, боевой вождь Храма срезал рыбки вместе с волосами, сказав, что не положено.

--- Урод он пресноводный, --- заявила Эрхэ.
--- Нигде в кодексе Ягуара не прописан регламент по украшениям.

Митхэ никому не позволяла такие выражения в адрес рекомого вождя, но Эрхэ была особым случаем.

--- Кихотр, --- покачала головой Эрхэ.
--- С бусинами проблемы, говорю сразу --- они мешают спать, наутро все уши будут мятые.
Лучше бы ты заколки взяла с подвесками, их раз --- и сняла.
Или подождать, пока волосы отрастут, и сделать валяные хвосты, как у меня.

Эрхэ по очереди тронула рыбки на висках.
Митхэ промолчала.

--- Ты поэтому отказалась стричься намедни?
Горе моё...
Хорошо, давай посмотрим, что тут можно сделать, --- воительница приняла из рук командира бусины и громко добавила в пространство:
--- Ситрис, я слышала, что замок не щёлкнул!
\ml{$0$}
{Пошёл вон, горшок с кровью!}
{Get out, you pot of blood!''}

--- Гярсёк с крёвью, --- тоненько передразнил её воин и снова вылез из-за лежанки.
--- Удачно погулять, Митхэ.
Атрису привет.

Митхэ гуляла по Тхартхаахитру каждый день.
Это тоже было частью ритуала.
Сначала она проходила по торговым рядам, переругиваясь с торгующими и расталкивая зевак.
Затем покупала очередную безделушку у старушки, стоявшей в углу площади.
Старушка долго благодарила, понимающе улыбалась и бросала хитрые взгляды то на сидевшего невдалеке Атриса, то на хмурую, неулыбчивую Митхэ.
Воительница от таких молчаливых откровенностей краснела до кончиков ушей и спешила убраться подальше.
Затем, сделав огромный круг в несколько кварталов, Митхэ возвращалась на площадь, садилась невдалеке от менестреля и слушала его игру.

Атрис упорно избегал её взгляда.
С его лица не сходила виноватая улыбка, на вопросы он отвечал односложно и неопределённо.
Самолюбие воительницы было серьёзно уязвлено.
Атрис дружески беседовал с нелюдимым ювелиром, который бросал менестрелю <<на удачу>> пару кусочков серебра.
Атрис называл <<кормилицей>> похожую на лань хозяйку постоялого двора, которая каждый день носила певцу еду...

--- Опять ничего не ел, журавлик, --- ворчала она и доставала из складок плаща горячий горшочек с похлёбкой.
--- Костями будешь скоро греметь в такт музыке!
Ла, лопай.
Миску потом заберу.
Вечером же здесь будешь?
Приду слушать.

--- Храни тебя духи, кормилица, --- светло улыбался Атрис и жадно набрасывался на похлёбку.
Митхэ уходила, не в силах выдержать это зрелище.
Иногда это была его единственная трапеза за день --- всё остальное уходило на ночлег и починку инструмента.

После обеда Митхэ шла в постоялый двор, с некоторой грустью бросала старушкину безделушку в походный сундук и заваливалась спать, не сняв одежду.

<<Я бы насыпала тебе полный горшок золота, --- думала она с обидой, --- но я же не хочу покупать твою любовь!>>

Так продолжалось уже целую декаду.
Но у воительницы было предчувствие, что сегодня всё будет иначе.
Разумеется, сегодня снова прогулка, снова по тому же маршруту.

--- И ещё, Обжорка, --- тихо сказала Митхэ напоследок, --- пообещай, что не будешь третировать Ситриса.
Нам полезен такой человек.
\ml{$0$}
{Ты видела, как он подкрался?}
{Did you see him sneak up?''}

\ml{$0$}
{--- Нет, не видела, --- хмуро ответила Эрхэ.}
{``No, I didn't,'' \Oerchoe\ said grumpily.}
\ml{$0$}
{--- Просто знала, что он опять по своей глупой привычке...}
{``I just knew he do it again, out of habit, and---''}

\ml{$0$}
{--- Именно!}
{``Exactly!}
\ml{$0$}
{А это значит, что нам есть чему у него поучиться!}
{And it means he has something to teach us!''}

\ml{$0$}
{--- Не люблю подкрадываться.}
{``I hate sneaking.''}

\ml{$0$}
{--- Может пригодиться и самое нелюбимое умение.}
{``Even least favourite skill may be useful.}
Ситрис честно старается стать одним из нас, а его выходки --- просто игра.

--- У нас здесь не детская комната, --- так же тихо бросила Эрхэ.
--- Я привыкла ко всем новичкам, но этот --- трусливый, скрытный и скользкий тип!
И плевала я, что он больше не разбойник, такое просто так не проходит.
На Юге бы эту рыбину на уху пустили...

--- Если понадобится, я сама сделаю из него любое блюдо, --- в голосе Митхэ промелькнули металлические нотки.
--- Но до той поры он --- твой соратник.
Обещаешь?

--- Только ради тебя.

--- Надеюсь на понимание.
За рыбки --- огромная благодарность.
\ml{$0$}
{Я выгляжу лучше?}
{Do I look better?''}

\ml{$0$}
{--- Честно?}
{``Honestly?''}

\ml{$0$}
{--- Насколько возможно.}
{``As it possible.''}

\ml{$0$}
{--- Ты выглядишь так же, как и всю последнюю декаду.}
{``You look as you did for the last decade.''}

\ml{$0$}
{--- И как же?}
{``How?''}

\ml{$0$}
{--- Растерянной.}
{``Confused.''}

\ml{$0$}
{--- Вот досада.}
{``Damn.''}

Митхэ поправила начавшие раздражать украшения, подхватила мешок и вышла за дверь.

\section{Рыба-карп}

--- Всё, пришли, --- выдохнул я и прислонился к холодной каменной стене.

Чханэ выпуталась из верёвок и без сил повалилась на чистую солому, которой был покрыт пол подвала.
Я только что протащил девушку домой в окровавленной, ещё тёплой оленьей шкуре --- через ворота, три улицы и торговую площадь, мимо городской стражи, двух десятков вечерних прохожих и домашнего оцелота, безмятежно дремавшего на подоконнике.
Шкура стоила мне пятидневного заработка.
Столбик, закадычный друг детства, пообещал молчать.

--- Не знаю, что ты опять собрался протащить в город, но надеюсь, что это безжалостный убийца твоей неопытности, --- пошутил друг.

--- Я тебе язык к носу пришью, петух-топтун, --- посулился я.
--- Следи за своими яйцами.

--- Я серьёзно, Лис.
От тебя пахнет свиньями, кошками, но сильнее всего --- молодой девушкой, --- Столбик обнюхал мою рубаху и мечтательно задумался.
--- И, кажется, она не мылась с рождения.
Это важно?

--- Вообще-то да, --- вздрогнул я.
--- Есть чем перебить запах?

--- Так мускус я тебе на что дал?
Натрись им.
Если всё серьёзно, могу дать маликхов венок --- окуришь и себя и свою немытую любовь.

--- Тащи.
И всё сухотравье от москитов, которое есть.

--- Да куда тебе столько? --- расплылся в щербатой улыбке Столбик.
--- Ты её любить собрался или коптить?

--- Заткнись и неси.
Кадильницу тоже, завтра верну.

Нюх Столбика был вполне объясним --- женщины сводили его с ума.
Он был младше меня на дождь, но уже успел прижить четверых детей с разными женщинами.
К счастью или на беду, мужчины его не интересовали.
Улыбчивый крестьянин мне нравился тем, что с ним не нужно было регулярно проводить время.
Я мог пропасть на дождь или на пять дождей, но был уверен, что Столбик встретит меня с радостью и мы поболтаем о
всякой ерунде, словно расстались рассвет назад.

Последним препятствием на пути в подвал был старый слуга, Сиртху-лехэ.
Он как будто нарочно устроился возле входной двери, держа в зубах тлеющую лучину и читая план посевов на грядущий год.

--- Хорошая добыща, Ликхмаш-тари, --- Сиртху-лехэ, не поднимая глаз, принюхался к воздуху;
лучинка вспыхнула, когда он заговорил.
--- Мушкушный шамец, молодой и щильный, дождей пять ш половиной.
Неужто дикие штада ещё жабредают в наши края?

--- Слава духам, --- ответил я, изо всех сил стараясь говорить ровно.
Не то чтобы Чханэ была тяжелой, но...
Да.
Она была очень тяжёлой даже для своего роста и комплекции.
\ml{$0$}
{Про таких людей говорят --- золотые кости, ртутная кровь.}
{This kind of people is said to have golden bones and mercury blood.}
Вместе с амуницией груз получился внушительный, и спина заныла, словно меня переехала оленья упряжка.
Но не к лицу обращать внимание на мелочи, когда товарищ в беде.

--- Я принесу тебе поесть и воду, --- выдохнул я.
--- Сегодня сюда не придут, но постарайся не шуметь...

Чханэ не ответила.

На кухне дома царила тьма, лишь в углу слабо светился остывающий очаг.
Спящие на полу служанки услышали мои шаги и зашевелились.

--- Ликхмас?

--- Эрхэ, Ликхэ, спите.
Разморило же вас посреди кухни.
Скажите только, где остатки ужина.

--- Тут прохладнее, --- Эрхэ вскочила, на ходу надевая через голову рубаху.
--- Давай разогрею, очаг ещё не совсем остыл.
Только надо...

--- Идолы тебя загрызи, Речка! --- вполголоса выругался я и прикрыл рот ладонью.
--- Вытащи мне миску еды, неважно чего, и ложись.
Ложись, я сказал.

Эрхэ переспросила, не разогреть ли рыбу, вытащенную из сберегающих свежесть листьев кислотника, и, получив от меня выразительный взгляд, пожелала избежать гнева Безумного этой ночью.
Чашу чистой воды из каменной бочки я набрал сам.

--- Чханэ.

Девушка приоткрыла слипающиеся глаза, увидела перед собой воду и жадно приникла к деревянной чаше.

--- Ещё воды?

--- Нет, --- шумно дыша, просипела она.
--- Благослови тебя духи.

--- Мы не в пустыне, воды много.
Во дворе стоит колодец.
Он крепкий и чистый, только в нём иногда купаются змеи, если забыть закрыть крышку.

--- Я запомню.

--- Ла, поешь, --- я пододвинул к ней миску с рыбой.
Девушка приподнялась на локте, и в её глазах мелькнул слабый интерес.

--- Что это?

--- Рыба.
В реке плавает.
Называется карп.

--- Это рыба? --- недоверчиво спросила Чханэ.
--- У нас все ею ругаются, но никто не может толком объяснить, что она собой представляет.
Вкусная?

--- Как птица согхо на вкус, но гораздо нежнее.

Рыба Чханэ не понравилась.

--- Сколько же здесь... тху... костей... тху...

--- Какая есть.
Кушай аккуратнее, не люблю вынимать из горла кости.
Мозг можешь выгрызть, он вкусный.
Молоки тоже ничего.

Глядя на уплетающую рыбу Чханэ, я вдруг понял, что сам хочу только две вещи: еду и чистую постель.
Я стащил с полки пару старых ковров и как можно уютнее сложил их в укромном месте, за штабелями досок чёрного дерева.

--- Ложись, здесь прохладно.
Кстати, если желаешь справить нужду --- в том углу дверь, возле двери большой папоротник.

--- Я запомню, --- девушка безуспешно пыталась справиться с завязками платья.
--- Помоги, пожалуйста.
Шнурки склеились...

Я помог.
Чханэ рывком стащила пропитанное потом, испачканное кровью платье через голову и не без злости бросила его в угол.
Потом повернулась ко мне, и я невольно восхитился изяществом обнажённого тела.
Несколько мгновений я любовался её нежной крепкой грудью, немыслимо тонкой талией и плоским животом, перечёркнутым наискось грубыми параллельными шрамами.
Пылеройский костяной кхаагатр, вне всякого сомнения.
Повезло --- удар пришёлся вскользь.

Чханэ соловым взглядом смотрела на меня сверху вниз и ухмылялась.
Её сердце пульсировало мягко и ровно, словно озёрная медуза, шея слегка потеплела.

--- Насмотрелся?
Я ведь спать хочу.

--- Храни твой сон лесные духи, --- я повернулся, чтобы уйти.

--- Хранят глаза и верный друг, --- услышал я за спиной.

\section{Ливень}

Поздним вечером зарядил густой, но ещё по-летнему тёплый ливень.
Большая капля --- неожиданная гостья.
Митхэ короткими перебежками возвращалась на постоялый двор.
По большей части травяные и сплетённые из листьев юкки навесы не спасали от дождя.
Кое-где подпорченные за лето крыши просто провалились под напором тяжёлой, как свинец, воды.
Митхэ ругалась сквозь зубы --- поход был назначен на послезавтра, и ехать в полумокрых штанах и сапогах ей совсем не улыбалось.

Площадь Тхартхаахитра находилась в низине.
Сюда стекали потоки воды со всего города.
К утру, к началу торгового дня, канализация справилась бы с затопившей площадь стихией, но сейчас, в самый её разгар, мутные потоки плескались где-то на уровне щиколотки.
Чуть дальше, ближе к южному концу, провалиться в воду можно было и по пояс.

Митхэ ругалась всё сильнее и громче.
Наконец воительница увидела путь на другую сторону --- перекладины навесов шли не прерываясь, по ним можно было перебраться.
Лёгкая, худенькая Митхэ, недолго думая, схватила дорожный мешок в зубы;
вскоре она уже по-обезьяньи поползла по перекладинам, жмурясь от летящих прямо в глаза водяных капель.

Вдруг её слуха снова коснулся звук знакомой цитры.
Ошеломлённая воительница обернулась и, сорвавшись с навеса, шмякнулась в мутную, пахнущую гниющими водорослями воду.
Дорожный мешок жалобно хрустнул в стиснутых зубах.

Атрис сидел под дождём на том же месте.
Изорванная одежда из грубой ткани насквозь пропиталась водой, волосы висели мокрым войлоком.
Цитра на его коленях булькала и хрипела, но он продолжал играть бодрую разудалую песню.

Митхэ не выдержала.
Она пошлёпала по колено в воде, не обращая внимания на набравшуюся в сапоги осеннюю муть.
Она схватила менестреля за плечи и тряхнула.
Ей хотелось закричать на него:
<<Что ты делаешь?
Зачем ты мокнешь под дождём?
Зачем ты играешь эту глупую песню?
Разве ты не понимаешь, как фальшиво она звучит?..>>

Атрис поднял на Митхэ потускневшие, полные печали глаза --- две тусклых зелёных искры над запавшими щеками, --- и злость Митхэ захлебнулась.

--- Идём, --- коротко бросила она и, схватив менестреля за руку, потащила его под ближайший навес.
Атрис молча пошёл за ней.

--- Снимай. Снимай всё, --- приказала она и стала вытряхивать из мешка самые сухие рубашки.
--- Пожалуйста, --- добавила она, опомнившись.

Атрис стоял, не двигаясь, и смотрел на Митхэ во все глаза.

--- Да снимай же!

Митхэ начала срывать с мужчины лохмотья;
менестрель качался, как деревце под шквальным ветром.
Грубая мокрая ткань трещала, не желая поддаваться.
Боевая раскраска Митхэ поплыла, и слёзы катились по лицу вперемешку с каплями ливня, сурьмяной кашицей и дорожной грязью.
Воительница яростно вытирала лицо ладонью и снова хваталась за рубашки, оставляя на них серые отпечатки.

Но рубашки быстро пропитались водой и потеряли уютное тепло.
Митхэ в отчаянии зарыдала.
Зарыдала, уткнувшись в полуголую мокрую грудь Атриса, до крови прикусив губу.
Вдруг она почувствовала его объятья.
Тонкие музыкальные руки с изящными пальцами едва заметно дрожали.

--- Атрис, пойдём со мной. Пойдём, --- сквозь рыдания шептала Митхэ.
--- У меня тепло и сухо.
У меня есть одежда, еда, всё, что только хочешь.
Пойдём.

--- Я пройду с тобой до смерти и ещё пару шагов, --- растерянно пробормотал Атрис.
--- Только скажи для начала, как тебя называть.

Митхэ всхлипнула и прижала его к себе крепко-крепко.
Выпавшая из рук Атриса цитра грустно звякнула и захлебнулась водой.

Ночь спрятала город, как устрица жемчуг, но это не имело значения.
Ничто в мире уже не имело значения.
Порог комнаты в постоялом дворе переступили две пары ног.

\section{Легенда о спящей}

Для жреца ночь --- это жаркое время.
Но я, слава духам, жрецом тогда ещё не был и имел полное право спать ночью, словно весь день рубил дрова.

Стараясь ступать как можно тише, я пошёл вверх по лестнице, ведущей в зал.
Купеческий двор, который занимала моя кормилица Кхотлам ар’Люм, построили очень давно, ещё до её рождения.
Насколько я знал, зодчий двора был одним из лучших в Тхартхаахитре;
впрочем, то же мог сказать и любой прохожий.
Одному Безумному известно, каким образом зодчий возвёл шедевр из простого малахита и медного шлака.
За сто пятьдесят дождей ни одна резная колонна с ликами духов, ни одни ажурные перила в виде разинувших пасть змей не потребовали реставрации.
Мужчина Кхотлам, Хитрам ар’Кхир, был из крестьян, но его смекалка часто помогала кормилице вести дела.
Так как я собирался стать жрецом, купеческому делу учили моих сестрёнок-близняшек, Манэ и Лимнэ, не достигших ещё двадцати пяти дождей.

Хоть кормилица и не знала нужды, на доспехи, оружие и свитки я зарабатывал сам --- неписанное правило.

И всё же, что делать с Чханэ?
Игра в прятки не могла тянуться вечно.
Нужно было посвятить в тайну кормильцев, рассказать им о девушке.
Не истину, нет, но часть истины.

Посреди зала, обставленного плетёными креслами и пахучими циновками, горел обделанный малахитом конус-очаг.
Перед ним сидела с рукоделием Кхотлам --- узелковое плетение приводит в порядок мысли, к тому же занавеси и ажурные салфетки хорошо продавались приезжим.
Я шумно вдохнул.

--- Лисёнок, подойди, --- встрепенулась от сладкой дрёмы кормилица.
--- Успешна ли была твоя охота?

Я сел перед ней на колени и взял её руки в свои.

--- Три птицы сиу'сиу, кормилица.
И мраморная змея по дороге попалась.
Завтра продам на рынке.
Идолов волею Безумного встретить не довелось.

--- А их уже полгода никто не видел.
Наверняка готовят очередной набег.
Снова нам смазывать стрелы.

--- Тебе не обязательно.

--- Как же, не обязательно, --- с укоризной сказала Кхотлам.
--- Ты же не хочешь, чтобы женщина, которая тебя воспитала, прослыла трусихой?

Кормилице на вид было не больше шестидесяти.
Все говорили, что мы похожи;
я считал иначе.

--- Ты храбрая женщина, родильница, и я не опозорю твоё лоно, --- я улыбнулся, давая понять, что предыдущая фраза была лишь проявлением заботы о ней.
Кхотлам вдруг смутилась и опустила глаза.
Верёвочки плетения снова пустились в пляс.

--- Кого сегодня принесут в жертву?

--- То ли ребёнка кузнеца ар’Хетр, то ли другого, с приречного хутора.
Я их видела, красивые.
Печально отдавать таких Безумному, --- кормилица улыбнулась и характерным жестом тронула пальцем язык, показывая, чтобы я не принял её слова как богохульство.

Ненависть к богу была обычным делом;
многим бы даже показалось странным её отсутствие.
Однако такая вещь, как богохульство, не приветствовалась --- в его руках было слишком много власти.

В жертву чаще всего приносили детей, не прошедших Отбор.
Пятого числа каждой декады в город выходили жрецы, повязанные красными лентами, и почти всегда возвращались с теми, кому было суждено продлить жизнь поселения ещё на несколько дней.
Платой была традиционная золотая пирамидка;
я видел такие --- ими порой рассчитывались на рынке.
Иногда на алтарь отправлялись и кутрапы, и не только, если другого выхода не было.
Одна мучительная смерть для города лучше, чем волна радужного безумия, после которого люди оплакивали убитых ими самими друзей.
Суровая необходимость...

Храм же по воле Безумного рушился постоянно, по камешку, по дощечке, иногда целыми кусками.
При каждом храме состоял зодчий и десяток строителей, восстанавливавших здание.
<<Будь у Безумного ноги, он бы жарил их на вертеле>>, --- так говорили люди.
Хотя учитель Трукхвал как-то обмовился, что Безумный здесь ни при чём.

--- Кормилица, я хочу тебе кое-что открыть.

--- Я слушаю тебя.

--- Я нашёл в джунглях девушку, отбившуюся от каравана, и привёл её домой.

--- Отбившуюся от каравана? --- Кхотлам пристально посмотрела на меня.
--- Я не слышала, чтобы здесь проходили караваны последние десять дней.

Я почувствовал, что дал маху.
Порой забываешь, что кормилица --- опытный купец.
Надо было выворачиваться.

--- Где именно ты её нашёл?

--- В сотне шагов от тропы, ведущей в Живодёрские леса, --- осторожно начал я.

--- Хай, --- успокоилась Кхотлам.
--- Какие-то безумцы опять повели караван там в надежде сократить время.
Меня никто предупредить не догадался, а рассказывать про опасности Цыплячьей тропы я уже устала.
Караван погиб?

--- Да, --- легенда для Чханэ была почти готова.

--- А сама она откуда?

--- С пустынного города.
Она не сказала какого, --- тут уж девушка пусть сама выкручивается, вдруг кто-то из местных знаком со жречеством Тхаммитра.

--- Кихотр, --- кормилица покачала головой.
--- Удивительная по своей нелепости история.
Кто она по крови?

--- Сели.
Её выносила и родила купец, семя --- воинское.
Ей тридцать пять дождей, на ней богатое платье, явно сшитое для неё.
Оцелотовая кожа.
Держит себя достойно, и черты её благородны.

--- Хаяй, сколько достоинств у этой девушки, --- задумалась Кхотлам.
--- Ещё и оцелот.
Что ж, утром поглядим на неё.

--- Кормилица, --- я решил довести дело до конца, --- понимаю, что моё решение несколько поспешно, но я хочу видеть её своей женщиной.
Что мне делать?

--- Лис! --- Кхотлам бросила рукоделие и взяла меня за руки.
--- Ведь ты ещё не стал жрецом, а она не стала тебе другом!

--- Я понимаю, кормилица.
Но я не знаю женщины, более достойной носить и кормить грудью моих детей.

Добрые глаза кормилицы вдруг улыбнулись, но губы наставительно сказали:

--- Дай ей время обжиться, Лисёнок.
Женщину могут оскорбить ухаживания человека, от которого она зависит в силу обстоятельств.
Где ты ей постелил?
Нуждается ли она в чём-нибудь?

--- В подвале, за досками чёрного дерева.
Я принёс ей воды и пищи.

--- Хорошо.
Не будем её беспокоить.
Иди спать, Ликхмас.
Храни твой сон лесные духи.

--- Хранят глаза и верный друг, --- улыбнулся я.

--- Опять эта воинская чушь, --- пробормотала Кхотлам, снова занявшись рукоделием.
--- Чего бы стоили их глаза без искусства дипломата...

Попрощавшись с кормилицей, я украдкой побежал в подвал и разбудил Чханэ, чтобы рассказать ей её легенду.

Теперь можно и поспать.

\section{Любовная лихорадка}

Заснуть любовники так и не смогли.
В ту ночь Митхэ охватил жар.
Её знобило так, что зуб на зуб не попадал.
Атрису было не лучше.
Меняя друг другу влажные полотенца и поглощая один чайник отвара за другим, они кое-как протянули до рассвета.

--- Кошмар, --- пробормотала Митхэ и закашлялась, кутаясь в одеяло.
--- Давно я так не болела.

--- Мы оба расплачиваемся за то, что уделяли себе чересчур мало внимания, --- ответил менестрель.
--- Я давно разучился это делать.
Поухаживаем друг за другом, вдруг поможет.

На следующий рассвет стало только хуже.
Вставали любовники редко и лишь по нужде.
Пришедший жрец, едва переступив порог и принюхавшись, тут же натянул утиную маску и плотно закрыл дверь.
Вердикт --- острая болотная лихорадка, полный карантин на три рассвета, ещё три рассвета спать и кушать отдельно от всех.

--- Вы первые, --- поделился жрец, насыпая в тарелочки пилюль.
--- Хотелось бы, чтобы и последние были в этом сезоне...

Атрис сильно страдал по мужской части, временами плача от долгой неутихающей боли --- сказалось длительное сидение на площади в любую погоду.
Митхэ аккуратно держала мокрое полотенце у него на паху.

--- Зачем тебе это? --- бормотал Атрис сквозь слёзы.

--- Тише, милый, тише, лежи, --- ласково отвечала ему воительница.
В её уставших, воспалённых глазах светилась чистая любовь, немыслимая для женщины её возраста, и этот нежный взгляд действовал на менестреля сильнее всех лекарств.

В конце концов Митхэ, поколебавшись и взяв себя в руки, отправила посыльного к торговцам с просьбой об отсрочке, а затем, собрав своих людей, сообщила им о новой договорённости.

--- Если кто-то торопится на встречу со смертью, можете выдвигаться без меня, --- устало добавила Митхэ, услышав тихий ропот.
Наёмники, покачав головами, разошлись.

На третий рассвет болезнь отступила.
Под вечер Митхэ решила искупаться и привести в порядок вещи, которые так и валялись по всей комнате.

Атрис лежал молча.
Он дышал глубоко, с подаренным лихорадкой тихим хрипом, но Митхэ знала, что он следит за каждым её шагом по комнате.
В бреду они говорили о чём-то важном, но Митхэ не помнила из разговоров ни слова.
Осталось лишь понимание --- бессловесное ощущение того, что ещё несколько кусочков мозаики мироздания встали на место.

--- Твои люди ждут, --- напомнил менестрель.

--- Я без тебя никуда не поеду, --- ласково, но твёрдо отрезала Митхэ.

--- Мои руки трясутся, когда я прикасаюсь к цитре.
Моё мужское естество уничтожено болезнью.

--- Ты ещё не выздоровел, милый.

--- А если я останусь таким до смерти?

--- Я буду к тебе нежной до смерти и два шага после.

Атрис не стал спорить.
Вскоре Митхэ почувствовала, что её тянет в сон.
Она наскоро обглодала куриную ножку, которая с характерным запахом задумалась о смысле жизни, и снова полезла к менестрелю под одеяло.

Атрис долго лежал и смотрел на спящую Митхэ.
Она была ещё бледна после неожиданной болезни, но от всего её маленького и тёплого, чутко спящего существа веяло невыразимой прелестью.
Уютная горячая вода купели сделала своё дело --- кожа Митхэ засияла здоровым розово-золотистым светом, а кудри стали мягче, чуточку светлее и замерцали, словно перламутр.

Разумеется, он с первой встречи знал, кто она такая.
Невозможно жить в Центральном городе и не знать Митхэ ар'Кахр --- легенды про лист юкки и поднятую деревянными гребнями пыль рассказывали на каждом перекрёстке.
Однако Атрис был уверен, что сейчас ни один пылерой, ни один тенку не узнал бы в ней того ужасного, выпачканного в крови и пыли военачальника.
К Атрису прижималась самая прелестная на свете женщина, отзывчивая на тончайшие и нежнейшие из ласк.

Менестрель не знал и не мог знать понятия <<норма реакции>>, но ускользающая, невнятная, едва осознанная истина успела потрясти его до глубины души.
Он погладил пальцами рваный шрам на губе, затем коснулся другого шрама --- на месте правой груди.
Кажется, она рассказывала об этом в бреду.
Груди её лишило зазубренное копьё тенку.

За свою славу воительница заплатила многим.

Митхэ очнулась от сладкой дрёмы, испуганно задышала и зашарила рукой по одеялу, отыскивая Атриса.

--- Я здесь, --- нежно шепнул менестрель.
--- Я рядом, я всегда буду рядом.

Митхэ облегчённо вздохнула, крепче обняла Атриса и погрузилась в глубокий сон.
Она узнала всё, что желала знать.

Атрис лежал в объятиях Митхэ и, может быть, впервые в жизни чувствовал себя защищённым.
Он спокойно думал о предстоящем походе, который согласился разделить со своей женщиной, о сезоне дождей и о том, что его цитра давно требует капитального ремонта.
Нужно было заменить серебряный колок, пятую струну и просевшее ушко маховика...
Мысли, словно ленивые откормленные карпы, плавали в приятно отяжелевшей голове.
Менестрель больше не боялся.
Митхэ приходила каждый раз, когда он звал её в мыслях, и он знал, что она будет находить его снова и снова, в какое бы далёкое межзвездье не унёс его ветер времени.

Мало-помалу Атрисом овладели странные образы из преддверий снов.
Словно листья из веток, они лезли друг из друга, следуя одним им известной дорогой ассоциаций.
Знакомые слова сплетались в бессвязные предложения, имевшие смысл лишь на той неуловимой грани сна и бодрствования.
И, уже приветствуя Сана, Атрис успел подумать:

<<Неужели так бывает?..>>

\section{Первое знакомство}

Двигатели затихли окончательно.

Мы не стали будить тси сразу.
Хотелось, чтобы первым, что они увидели, был рассвет.
Что-то мне подсказывало, что на этой планете очень красивые рассветы.
Друзья молчаливо со мной согласились.

И вот чудо произошло.
На небе появились яркие оранжевые полосы.

--- Как же хочется выйти... --- прошептала Заяц.

--- Пока нельзя, --- сказал Фонтанчик.
--- С чужой микрофлорой шутки плохи.

--- И с чужим хоргетом, простите за тавтологию, --- добавила Подсолнух.
--- Ты же омега-защиту поставил?

--- Нет, --- покачал головой Фонтанчик.
--- Я её отключил на всякий случай, когда начал посадку.

--- Зачем?!
Это демиург!

--- Именно поэтому и отключил.
Мы находимся на планете, которая принадлежит хоргету.

--- Фонтанчик, твои ж хромосомы! --- пискнула Подсолнух.
--- Хотя бы гаситель колебаний установи на мостик и тор-отсек, мы для него как на ладони!

--- Нет, --- категорически заявил канин.
--- Любые омега-технологии могут его испугать.
Я не спец в переговорах, но думаю, что вести переговоры со смертельно испуганным хоргетом проблематично.
Не говоря уже о том, что это может стоить нам жизни.

--- Командир? --- требовательно обратилась Подсолнух ко мне.

--- Я согласен с Фонтанчиком, --- кивнул я.
--- Сначала поговорим, потом, может быть, включим.

Подсолнух недовольно вздохнула и углубилась в показания приборов.

--- Он выйдет с нами на связь? --- поинтересовался Баночка.

--- Вообще должен, --- пожал плечами Фонтанчик.

--- Вполне возможно, он ждёт, что на связь выйдем мы, --- предположила Заяц.
--- У меня подозрение, что он застенчивый... насколько это вообще применимо к хоргету.
И вообще, я хочу гулять.

--- Нужно спросить разрешение у хозяина, --- возразил я.
--- Безымянный, каким бы застенчивым он ни был, дал разрешение только на посадку.

--- Сейчас спрошу, --- вдруг заулыбалась Заяц и, пока мы не успели сообразить, бросилась к панели.
--- Так, Баночка, кинь ссылку на словарь в репозитории...
А, всё, нашла.

--- У неё третье образование --- лингвистическое, специализация по поддерживаемым языкам, --- пояснил Фонтанчик остальным.
--- Не хочет упустить шанс...

\ml{$0-[ej]$}
{<<Безымянный! --- пронеслись в пространстве слова Эй-00.}
{``Nameless!'' the Ej-00 words flew through the space.}
\ml{$0-[ej]$}
{--- Ты меня слышишь, дружище?}
{``How copy, man?}
\ml{$0-[ej]$}
{Говорит Темнотой-Сотканный-Заяц из корабля тси.}
{It's Darkness-Woven-Hare in the Qi spaceship.}
\ml{$0-[ej]$}
{Конец сообщения>>.}
{Over.''}

\begin{quote}
\ml{$0-[ej]$}
{<<Сотканный-Темнотой-Заяц, сигнал принят.}
{``Copy that, Darkness-Woven-Hare.}
\ml{$0-[ej]$}
{Вероятно, тси хотят узнать, стабильна ли связь.}
{Perhaps Qi want to know how stable the connection is.}
\ml{$0-[ej]$}
{Отвечаю: связь стабильна.}
{My answer: the connection is stable.}
\ml{$0-[ej]$}
{Я осуществляю постоянный мониторинг этой частоты, тси могут связаться со мной в любое время.}
{I am constantly monitoring this frequency signal, Qi can call me anytime.}
\ml{$0-[ej]$}
{Конец сообщения>>.}
{Over.''}
\end{quote}

--- Он ответил! --- восхищённо воскликнул Баночка.

--- Я совсем забыла, что у них другая рецепторная система и мировосприятие, --- Заяц хлопнула ладонью по лбу.
--- Но он умный и понял.
Это хорошо.

\ml{$0-[ej]$}
{<<Безымянный, сигнал принят.}
{``Copy that, Nameless.}
\ml{$0-[ej]$}
{Тси спрашивают разрешения покинуть корабль и вступить в атмосферу планеты.}
{Qi ask permission to get out of the spaceship and to enter the atmosphere of your planet.}
\ml{$0-[ej]$}
{Тси устали сидеть внутри корабля и хотят дать физическую нагрузку конечностям.}
{Qi are tired of being inside, they want to give their limbs a workout.}
\ml{$0-[ej]$}
{Конец сообщения>>.}
{Over.''}

\begin{quote}
\ml{$0-[ej]$}
{<<Сотканный-Темнотой-Заяц, сигнал принят.}
{``Copy that, Darkness-Woven-Hare.}
\ml{$0-[ej]$}
{Отвечаю: даю вам разрешение находиться в любой точке планеты, включая её атмосферу и зону доступного регистрации магнитного поля, и давать конечностям любую необходимую нагрузку.}
{My answer: you have my permission to be anywhere in the planet, including its atmosphere and measurable magnetic field, also to give your limbs any workout they need.}
\ml{$0-[ej]$}
{Предупреждение.}
{Warning.}
\ml{$0-[ej]$}
{Для посадки выбрана точка планеты, в которой наблюдается практически нулевая концентрация микро- и макроорганизмов, а температура, радиация и влажность оптимальны для представителей вида Homo homo sapiens.}
{In the point chosen for the landing, there is a low concentration of macro and micro organisms, and temperature, radiation, and moisture levels are optimal for a member of the species Homo homo sapiens.}
\ml{$0-[ej]$}
{Однако я не знаю, к какому биологическому виду относятся тси и какова ваша физиология, поэтому советую вам озаботиться собственными средствами защиты.}
{However, I know neither a species you belong to nor your physiology, that is why you are encouraged to use your own protective equipment.}
\ml{$0-[ej]$}
{Конец сообщения>>.}
{Over.''}
\end{quote}

--- Очень интересно, --- протянул Фонтанчик.
--- Он не знает, к какому биологическому виду относятся тси, но при этом говорит на Эй-00?
Хм.
Заяц, узнай у него поаккуратнее насчёт Ордена Преисподней.

\ml{$0-[ej]$}
{<<Безымянный, будет исполнено.}
{``Roger that, Nameless, wilco.}
У нас имеются средства защиты.
Тси спрашивают, известно ли тебе, кто мы и кто такие демоны.
Конец сообщения>>.

\begin{quote}
<<Темнотой-Сотканный-Заяц, отвечаю на первый вопрос: вы --- объединение сапиентов под названием ,,тси``.
Название я впервые узнал из вашего сообщения.
Я не знаю признаков, по которым выделяется это объединение.
Отвечаю на второй вопрос: мне известно, кто такие демоны и по каким признакам выделяется это объединение.
Я осведомлён о том, что принадлежу к данному виду сапиентов.
Также я владею языком демонов --- Эй, таблица 00.
Конец сообщения>>.
\end{quote}

--- Всё ясно, --- затараторил Баночка.
--- Я думаю, что Безымянный --- творение одной из поздних сапиентных цивилизаций, которые когда-то находились под властью Ордена Преисподней.
Это объясняет и тот факт, что он упомянул Эй как единственный <<язык демонов>>, при том, что <<языков демонов>> как минимум три.
А ещё он сделал одну очень нехарактерную для Ада и Картеля вещь --- он отнёс себя к сапиентам.

--- Баночка прав, --- признал Фонтанчик.
--- И знаете, какое у меня ощущение?
Как будто мы только что поговорили с ребёнком.
Он ничего не спросил про полномочия Заяц вести с ним переговоры, да и в целом не силён в дипломатии.
Он совершенно нас не боится.
Демоны Ада и Картеля говорили с нами не иначе как через цепь спутников, с почтительного расстояния, а он... ну...

--- Он --- что? --- осведомилась Подсолнух.

--- Он только что осмотрел тебя, Подсолнух!
Датчики подтверждают наличие считывающих жуков.
А во время разговора он осматривал меня и Небо.
А чуть ранее он просканировал корабль, но крайне осторожно, на пределе чувствительности, иначе система сразу бы взвыла.
Ему интересно!
Он не совсем понимает, где живые существа, а где системы корабля, и поступает как ребёнок, которому интересно насекомое, но который знает, что неосторожным движением его можно раздавить!

--- Фе, --- фыркнула Подсолнух и демонстративно отряхнула костюм.
--- Он что, летал внутри меня?

Заяц подумала и снова обратилась к демиургу:

<<Безымянный, тси спрашивают, есть ли в этой звёздной системе демоны кроме тебя.
Если нет, то были ли они раньше.
Конец сообщения>>.

\begin{quote}
<<Темнотой-Сотканный-Заяц, отвечаю: в настоящее время я являюсь единственным демоном в этой звёздной системе.
Ранее здесь появлялись демоны, но ни один не пожелал со мной общаться и не остался в её пределах.
Конец сообщения>>.
\end{quote}

\ml{$0-[qi]$}
{--- Бедненький, --- посочувствовала Заяц.}
{``Poor child,'' Hare said with compassion.}
\ml{$0-[qi]$}
{--- Никто не хочет с ним дружить!}
{``No one wants to be friends with them!''}

--- Будите основную команду и готовьте скафандры, --- вздохнул я.
--- Полюбуемся рассветом снаружи.
Подсолнух, хватит отряхиваться, он тебя ещё не раз будет осматривать.

\ml{$0-[qi]$}
{--- По-моему, немного невежливо без разрешения собеседника изучать его структуру, --- буркнула Подсолнух.}
{``I think that's a bit rude to scan someone's structure without permission,'' Sunflower grunted.}

Заяц завизжала от радости и прыгнула в гравитационный лифт.

\ml{$0-[qi]$}
{--- Заяц, подожди! --- крикнул Фонтанчик.}
{``Hare, wait!'' Fountain shouted.}
\ml{$0-[qi]$}
{--- А поесть?}
{``Do we have the right to eat?}
\ml{$0-[qi]$}
{Мы голодные как волки!..}
{We're as hungry as wolves ... !''}

\section{Дикая похлёбка}

%Лук 7 дождя 11998, Год Церемонии 18.

Вскоре дома вкусно запахло наваристой похлёбкой.
<<Дикая похлёбка>>, сваренная с оленьей кровью и слизистой жижей стеклянных ягод, была коронным блюдом Ликхэ.
Пока Кхотлам рассаживала близняшек, кормилец помог служанкам принести еду в зал.
Мне доверили отрегулировать конденсатор --- день обещал быть душным, приближался сезон дождей.

Все домашние уже сидели за столом, когда из подвала вышла заспанная Чханэ.
Она сменила окровавленную зелёную повязку на свежую, белую, и надела чистую одежду кормильца, которую я принёс ей заранее вместе с бинтами.
Хитрам был самым высоким в доме --- и всё равно его штаны оказались моей новой знакомой чересчур малы.
Глаза девушки горели голодным огнём --- сколько-то она провела в пути без пищи?

--- Ага, --- довольно сказал Сиртху-лехэ, --- вот почему от оленя так странно пахло.

Хитрам скрыл улыбку за широкой рабочей ладонью и толкнул меня в бок.
Кормилица шикнула на него и кивнула вошедшей.

--- Тханэ ар’Катхар, --- Кхотлам вдруг нахмурилась, окинув взглядом подбородок и запястье девушки, --- тебе требуется жрец?

--- Нет, благодарю, жрец меня уже подлечил, --- Чханэ поклонилась и чуть заметно улыбнулась мне.

--- Не церемонься, садись и ешь столько, сколько захочешь.
Вид у тебя голодный.
Расскажешь обо всём потом.

Чханэ поклонилась ещё раз и осторожно села на низкий табурет.
Её колени чиркнули по краю стола, и посуда нежно зазвенела.

--- Прошу прощения, --- смутилась девушка.

Кормилица молча встала, передвинула свой табурет и устроилась рядом с Чханэ.
Затем зачерпнула из миски похлёбку и поднесла черпачок ко рту гостьи.
Чханэ, не говоря ни слова, съела предложенное.
Вскоре, игнорируя черпачок, она схватила миску руками и, громко чавкая, впилась в горячую наваристую жижу.

Закончилась трапеза в молчании.
Кормилец рассеянно крутил в руках чашу с водой и смотрел умными глазами в пространство.
Служанки и Сиртху-лехэ переговаривались жестами --- судя по всему, речь шла о детях общих знакомых.
Близняшки с любопытством смотрели на Чханэ --- они встретили оцелота впервые.
Кхотлам не то насмешливо, не то с нежностью поглядывала на меня.
Я, решив на время отложить смущение, терпеливо ждал вопросов.

--- Храни вас лесные духи, --- наконец выдохнула Чханэ.
--- Кто бы это ни приготовил, желаю ему ясных рассветов и лёгкости в руках.
Я давно не ела такого вкусного блюда.

Ликхэ довольно осклабилась.

--- Я так понимаю, настало время представиться, --- Чханэ решила взять инициативу в свои руки.
--- Я воин Храма Чхаммитра, мои дарители --- воин Акхшар ар’Качхар э’Чхаммитра и купец Шогхо ар’Хэ э’Чхартхаахитр.
Я...

--- Ты хранительница Согхо ар’Хэ? --- вдруг оживилась Кхотлам.
--- Знаю её, очень хорошо знаю.
Кажется, дождей пятьдесят назад она служила Сотрону?

--- Травинхалу, --- усмехнулась Чханэ.
--- Как она говорит, <<Акхсар провёз меня через половину мира --- из одной глуши в другую>>.

Повторив слова Согхо, девушка непроизвольно воспроизвела чужой акцент, тембр голоса и интонации.
Кормилица расслабилась и многозначительно мне кивнула.

--- Откуда был твой караван?

--- С Кристалла, --- ответила Чханэ.
--- Откуда точно --- не скажу, нанялась сопровождать в Мишитре.
Зизоце --- неразговорчивый народ.

--- Это точно, --- согласилась Кхотлам.
--- Все погибли?

Девушка кивнула.

--- Их главный непременно хотел достичь Ихшлантекхо... я имею в виду, Ихслан... тхара к началу страды.

--- Я тебя поняла, не волнуйся, --- успокоила кормилица гостью. --- Как звали предводителя?

Я похолодел.

--- Имя цатрон Тхартху ар'Ликх ми'Эрхо'люаэсакх, --- без запинки оттарабанила Чханэ.
--- Имя зизоце не знаю.
В документах он никогда не подписывался, просто кусал воск.

Ми'Эрхо'люаэсакх.
<<Родившийся за Течением Обманутого В Последний Раз>>.
Я впервые встретил такой топоним в чьём-то имени.
Но нахмурившуюся кормилицу, похоже, смутило нечто другое.

--- Странно, они редко используют имена цатрон с <<ми>> --- они очень привязаны к родным поселениям.
И воск они кусают нечасто, у них обычно есть для печатей кольца или амулеты со сложным рисунком.
И спешка у них, мягко говоря, не в чести...
По описанию я бы сказала, что этот человек "--- изгой из числа зизоце, долгое время проживший где-то в Водоразделе или даже севернее, например, в святилище Весёлый Волок.
Впрочем, ладно.
Зизоце порой чересчур медлительные, а имена у них длиннее дороги до Тхаммитра.

Я мимоходом восхитился Чханэ.
Она придумывала значимые детали легенды на ходу.
При этом Кхотлам улавливала эти детали, а я --- нет.

--- Почему ты думаешь, что этот человек --- изгой? --- поинтересовался я.

--- Я не представляю другую причину, по которой он молчит о родном поселении.
Скорее всего, ему просто запрещено о нём говорить, а значит, его за что-то изгнали.
От такого человека можно ждать проблем...

--- Проблем хватало, --- кивнула Чханэ.
--- Плату он задерживал ровно настолько, чтобы на очередной стоянке несколько самых нетерпеливых воинов ушли.
Тогда он нанимал новых, экономя таким образом несколько перьев золотого песка.
Этот трюк сработал два или три раза, пока воины его не обсудили в своём кругу и не договорились ждать до последнего и трепать ему нервы.

--- Слышала про таких, да, --- кивнула кормилица.
--- Некоторые ещё знаешь как делают --- когда подходит срок, разменивают...
Впрочем, ладно, не суть.
Продолжай, я тебя перебила.

--- Да он вообще экономил на всём!
Те воины, что шли с ними всю дорогу, сказали, что тракт святилища прошли на упряжках, чтобы не платить пошлину.
Торговцы не согласились на волок, обменяли корабли у каких-то сомнительных дельцов.
До Чхартхаахитра не спустились, прямиком из Мишитра снарядили караван.
Шли очень близко к землям идолов --- возможно, предводителя ввела в заблуждение близость Омута Духов, только он-то не знал, что там недавно...

Девушка запнулась, а меня прошиб пот.
О событиях в Омуте Духов кормилица рассказала домашним по большому секрету.
Чханэ, которая не могла о них знать, чересчур увлеклась и едва не раскрыла свой источник информации.
Впрочем, Кхотлам выглядела расслабленной.

--- ... что идолы услышали нас задолго до Живодёра, --- смятение Чханэ длилось едва ли секхар.
--- Между наёмниками с Грозового хребта и нашими разгорелся спор, который перерос в драку.
Я не знаю, о чём они спорили, но думаю, что из-за выигрыша в какой-то азартной игре.
Встретили нас с честью, спасло меня только то, что я была в арьергарде.
Я спала в палатке на обозе, даже одеться не успела --- обоз перевернулся...

--- А... почему платье?
Да ещё и такое дорогое? --- удивлённо спросил Сиртху-лехэ.

--- А мне голой по джунглям бежать? --- со смехом осведомилась девушка.
--- Надела, что было.
Жизнь свою с собой унесла --- и то хорошо.

--- Послушай, --- вдруг заговорил Хитрам.
--- У тебя очень интересное имя.
Первый раз встречаю женский вариант имени Тхалас.
И ещё ты очень высокая.
Неужели твои дарители такие же?

--- Согхо, как тушканчик --- маленькая и юркая, --- рассмеялась девушка.
--- Кормилец чуть повыше, но тоже ростом не вышел.
А имя мне пришлось придумать новое, потому что раньше я была парнем.
Был Чхалаш --- стала Чханэ.

Кхотлам ахнула, Сиртху-лехэ тихо пробурчал под нос:
<<Хаяй... редкость>>.
Я непонимающе уставился на старших.

--- Смена пола, --- объяснила кормилица, увидев моё замешательство. ---
Обычно дети выбирают пол навсегда, но бывает и такое --- зрелый мужчина становится женщиной или наоборот.
Я не знаю, почему это происходит.
Тебе следует спросить в Храме.

--- Не совсем зрелый.
Моё превращение началось через пару дождей после того, как у меня появились мужские признаки, --- поправила Чханэ.
--- Жрец сказал, это из-за того, что я ударилась головой в детстве.
Других причин он не нашёл.

--- А как изменялось твоё тело? --- поинтересовался я.

Кормилица неодобрительно посмотрела на меня, но промолчала.

--- Я не знаю, --- растерянно ответила Чханэ.
Она явно не ожидала таких личных вопросов.
--- Проснулась однажды утром, а у меня грудь набухла и в паху как-то странно зудит.

Я усмехнулся.
Кхотлам попыталась шлёпнуть меня под столом, но промахнулась;
я почувствовал лишь слабое движение воздуха.

--- Смейся, смейся, --- буркнула Чханэ.
--- Посмотрела бы я на тебя.

--- А потом?

--- Потом таз и спина заболели, да так, что хоть плачь.
За сезон обзавелась грудью и роскошной задницей --- кости на четыре пальца удлинились.
А после...

Обстановка всё же немного разрядилась.
Случай, произошедший с Чханэ, был большой редкостью, и домашние слушали девушку с неподдельным интересом.
Чханэ стала вести себя свободнее.
Внезапно из-за высоких скул и благородного лба, лежащего на своде длинных, размашисто очерченных бровей, проступила симпатичная весёлая говорушка.

--- ... а задница-то тяжёлая, двигаться пришлось учиться заново, и драться тоже.
В Храме мне доставалось, конечно, особенно первое время.

--- Кем ты себя чувствуешь? --- спросила Кхотлам.

--- Сложно сказать, --- пожала плечами девушка.
--- Когда была парнем --- не задумывалась.
Сейчас... больше женщиной, чем мужчиной.
Примерно три к одному.
Самое смешное --- кормильцы договорились, мол, если буду девочкой --- стану купцом, если мальчиком --- пойду в храм.
А хитрая Чханэ всех обманула.

--- Это урок твоим кормильцам, --- громко сказал Сиртху-лехэ и стукнул кулаком по столу.
--- Пусть молодёжь сама выбирает свой путь.
Мы можем лишь предложить то, что имеем.

--- Правильно, --- поддержала старика Кхотлам.
--- Военное дело не имеет пола, как и дипломатия...

\section{Надежда}

Что делать с Безымянным?

Этот вопрос сегодня был главным на повестке дня.
Дискуссия шла в закрытом, отгороженном от колебаний первичного поля зале Стального Дракона.
Может, демиург и не силён в дипломатии, но большая часть тси была уверена, что дипломатические отношения установить необходимо.
Я придерживался другой точки зрения:

--- Дипломатия помогает общаться тем, кто друг другу не доверяет.
Он --- ребёнок.
Вы хотите установить дипломатические отношения с ребёнком?

--- Он --- хоргет, --- заявил Мак.
--- Небо, не повторяй ту ошибку, которую мы допустили с Машиной.
Для нас в приоритете его сила и способности, а не психология, вернее, эмуляция психологических процессов.

Как ни странно, меня поддержали многие диктиологи.
Одна из них, Гладит-Зелёную-Кошку, выдвинула совершенно революционную идею, вызвавшую лишь смех:

--- Я считаю, что мы должны его воспитать.

Впрочем, началось всё с плана развернуть планетарную систему защиты против демонов, которая у нас была на Тси-Ди.
Тси приняли это решение сразу, несмотря на замечание некоторых товарищей о том, что вначале планету нужно узнать получше.
Баночка уже успел пошутить по этому поводу:

--- Кажется, тси заболели золотой лихорадкой.

Из истории мы знали, что когда-то основной обменной валютой было золото.
Преимущества золота как валюты очевидны: оно инертно и почти не подвергается окислению, оно компактно, оно легко обрабатывается и может быть идентифицировано даже сапиентом, мало смыслящим в химии, благодаря своеобразному цвету и консистенции.
История знала много <<золотых лихорадок>>, когда сапиенты пытались получить ресурсы путём добычи огромного количества валютного материала.
Разумеется, это не приводило ни к чему хорошему --- количество пищи и материалов оставалось прежним, и валюта обесценивалась, принося разочарование добытчикам и снижая качество жизни их собратьев.

У нас к золоту был свой интерес --- для резонаторных модулей планетарной системы и некоторых типов микросхем годится только этот металл.
Мы нашли первые месторождения, едва выйдя из корабля --- золотой песок лежал на морском берегу вперемешку с кварцевым.
Биолог Цветущий-Мак-под-Кустами сказал, что в жизни не видел ничего более красивого, и уже набрал немного прибрежного песка для аквариума.
Правда, предусмотрительно оставил его в лагере рядом с кораблём --- с чужой микрофлорой шутки плохи.

К счастью, среди тси нашлись несколько техников, которые работали с системой планетарной защиты и имели хорошее представление о её устройстве.
Одним из них оказался мой друг --- Хрустально-Чистый-Фонтан.

--- Я смогу собрать систему, --- заявил Фонтанчик.
--- Мне будут нужны данные о планете и материалы.
Также хочу попросить треть вычислительных мощностей Стального Дракона.

Друг не любил брать на себя руководство;
однако его отличала одна особенность --- он тонко чувствовал момент, когда иначе поступить нельзя.
Фонтанчику и ещё пятерым тси отдали кабинеты рядом с библиотекой, в которых они и пропали на пятнадцать дней.
Иногда у меня проскакивала мысль, что Фонтанчик умер;
впрочем, работавшие на полную мощность компьютеры Дракона и обновлявшиеся со скоростью света исходники в репозитории твердили, что Фонтанчик и прочие архитекторы живы настолько, насколько вообще можно быть живым.
В конференции они переговаривались исключительно чертежами и строками кода;
кое-кто умудрялся шутить на языке Си-поинт.

Время от времени, несмотря на ворчание архитекторов, техники просматривали уже готовые чертежи.
Систему Фонтанчик строил грубо --- о стравлении напряжённых углов и тонкой настройке речь не шла, --- но очень добротно и надёжно.
Остался доволен даже придирчивый Баночка.

--- Болтаться будет.
И энергию жрать.
А с такой напряжённостью углов мы будем систему раз в сто лет калибровать, это кошмар какой-то, --- резюмировал он.
--- Небо, ты не помнишь, когда в последний раз калибровали <<сову>> на Тси-Ди?

Я задумался.

--- А это не тот год, когда твоя любимая группа выпустила свой последний альбом?
Ты рассказывал как-то.

--- Слушай, точно, --- Баночка хлопнул себя по лбу.
--- Он так и назывался --- <<Калибровка>>, там же они все были техниками защитной системы.
Днём --- работа, ночью --- музыка.
И как ты только запомнил?
Так когда это было...
Восемьсот лет назад, получается?
Восемьсот четыре.
И система бы ещё столько же проработала без вмешательств.

--- Мы не располагаем такими возможностями, --- развёл я руками.

--- Нет-нет, я рад, определённо, --- замотал головой Баночка.
--- Фонтанчик молодец, я бы лучше него не сделал.
Так что ладно, был бы какой-никакой план, а там сориентируемся.

К осуществлению плана было одно-единственное препятствие --- демиург Планеты Трёх Материков.

Все понимали, что хоргет представляет для выживших тси огромную опасность.
Да, он позволил нам жить на планете.
Но он хоргет.
Было очевидно, что Безымянный, каким бы ребёнком он нам ни казался, не оценит разворачивание системы против ему подобных.

Проводились аналогии с Машиной --- вновь увидеть геноцид не хотел никто.
<<Но он мог уничтожить нас ещё в космическом пространстве!>> --- возражали сторонники сотрудничества.

Сторонниками противостояния предлагались самые разные варианты действий --- от дезинформации до внезапного удара экраном по месту локализации бога.
На последнее предлагалось потратить оставшуюся на корабле энергию.
Энергию, которой и так не хватало...

Точку в споре поставила сестра Фонтанчика, Съешь-Рычащий-Пирожок.
Вот её слова:

--- Мы не имеем морального права отбирать сотворённое у творца, кем бы ни был творец.
Мы не имеем морального права ответить предательством на помощь, кем бы ни был помогающий.
Ведя тайную или открытую войну против Безымянного, мы обернём против себя и планету, так как они --- одно целое.

Это был веский довод, учитывая, что Пирожок не питала к демиургу ни симпатию, ни доверие.
Она откровенно признала, что запустила бы робота и уничтожила бы бога без раздумий, будь она на моём месте;
и тем не менее в тот день женщина повела себя, как истинный тси.
Семьдесят процентов переселенцев проголосовали за информирование Безымянного о наших намерениях.

Культурологи приготовили для бога послание.
Вот оно целиком:

<<Тси спрашивают у демиурга Безымянного разрешение на постройку планетарной системы защиты против демонов, которая подразумевает уничтожение любого демона, приблизившегося к планете ближе указанного расстояния.
Тси вынуждены идти на этот шаг из-за угрозы со стороны Ордена Преисподней и Красного Картеля, которые весьма негуманными способами используют сапиентов как источник масс-энергии и как интерфейс для взаимодействия со Вселенной Фотона.
Мы благодарны Безымянному за спасение и потому обязуемся добавить его в исключения срабатывания системы.
Также мы желаем оговорить условия, которые компенсируют демиургу этот риск>>.

Ответ пришёл спустя час.
Я, как и многие мои товарищи, пребываю в лёгком шоке, поэтому просто процитирую перевод:

\begin{quote}
<<Я догадывался о ваших намерениях, тси.
Рад, что вы решили идти по пути укрепления сотрудничества.
Желание защититься свойственно и мне, а Орден Преисподней и Красный Картель одинаково опасны для всех обитателей моей планеты.
Вот мои условия: вы предоставляете мне ежедневный подробный отчёт об использованных для строительства ресурсах планеты, а также информацию о технологиях тси на мой выбор, исключая военные и технологии манипуляции сапиентами.
Взамен я обязуюсь предоставить вам полную информацию о ресурсах планеты и обеспечивать строительству необходимую защиту.
Жизнь за жизнь, знание за знание --- думаю, это равноценный обмен>>.
\end{quote}

--- Лучше бы Заяц ему написала, --- буркнул Фонтанчик.
--- С ней он разговаривал дружелюбно и искренне, как и она с ним.
Едва культурологи взяли настороженно-дипломатичный тон --- его словно подменили.

Впрочем, Фонтанчика в его скептицизме никто не поддержал.
Как ни крути, ответ был недвусмысленным согласием.
Переводивший сообщение Баночка до сих пор сидит, пьёт чай со стимуляторами и думает.
Волнуется, правильно ли перевёл, не упустил ли какого-нибудь нюанса.

Мы очень надеемся, что маленький плант всё сделал правильно.
Да, кажется, у нас появилась надежда.
Но я всё ещё боюсь в это верить.

\section{Истина}

Вскоре совсем рассвело.
Хитрам поднялся из-за стола и, взяв с собой несколько кусочков мяса, пошёл работать в огород.
Близняшки, весело переговариваясь и лукаво постреливая глазами в сторону гостьи, убежали в школу.
Кормилица попросила Эрхэ убрать посуду и жестом позвала нас с Чханэ к себе в комнату.

--- Хай, дети, --- деловито сказала она, убедившись, что мы остались втроём.
--- Если вы думали, что можете меня провести --- вы думали так зря.
Чханэ, --- кормилица, не дав девушке вставить слово, неожиданно заговорила с ней характерным западным говором, --- мне совершенно не важно, по какой причине ты сбежала из родного города.
Я никогда не приветствовала убийство кутрапов и считаю, что любому человеку, даже дважды и десять раз дважды кутрапу, однажды может потребоваться шанс.
Лис уже видит тебя своей женщиной, и если это не достойно того самого шанса, то пусть Безумный заберёт меня прямо сегодня на закате.

Чханэ смотрела на круглое, доброе лицо Кхотлам.
Впервые я увидел в её глазах слёзы.
Губы девушки дрожали.

--- Я... я убила жреца, --- проговорила она, и слёзы потекли ручьём.
--- Он хотел... хотел...

Мы с кормилицей обняли девушку.
Чханэ била настолько крупная дрожь, что мы втроём качались, как деревья под шквальным ветром.

--- Да-да, понимаю, --- отстранённо прошептала Кхотлам.
--- Закон велит спрашивать у жертв разрешения, но кто по доброй воле согласится на мучительную смерть?
Пустая формальность...

Вскоре Чханэ успокоилась.
Кормилица усадила её на кровать и укутала тонким одеяльцем;
я сидел рядом и гладил девушку по спутанным, кофейным с оранжевой искоркой волосам, время от времени со странным щемящим чувством запуская в них пальцы.
Словно я уже делал это раньше --- давным-давно.

--- Вот что, милая.
Не вздумай рассказывать эту байку про караван кому бы то ни было.
Служанок моих ты, может, и обманешь, но в Храме тебе не поверят.
Я сама придумаю, что сказать.
А пока --- иди-ка вымойся, ты солёная, как акула морская.
Я велю Эрхэ принести тебе воды, а одежду сошью сама.
Захочешь спать --- ложись здесь, на мою лежанку.

Чханэ громко всхлипнула.
Заплаканные огненные глаза смотрели то на меня, то на кормилицу.
Во взгляде девушки слабым, робким пламенем светились благодарность и надежда.

\chapter{Старый дом}

\section{Храм Тхитрона}

\epigraph
{Спартанский образ жизни привлекателен для многих, особенно для слабых духом.
Лишая себя благ, люди наивно полагают, что готовы ко всему --- к нашествию врагов, к катаклизмам и иным невзгодам судьбы.
Но возникает вопрос --- в чём смысл лаконичности, если не будет бед?}
{Михаил Кохани, презентация <<спичечных>> технологий в Массачусетском технологическом университете}

Война.
Этому простому понятию в наши дни подчинено всё.
Инструмент может быть оружием.
Транспорт может быть щитом.
Человек может быть воином.
Всё и вся, проигнорировавшие этот закон, давно ушли в землю.

За четыре рассвета Чханэ пришла в себя.
Вскоре девушка попыталась взять на себя работу по дому.
Кормилица отдала ей узелковое плетение, которым занимала руки сама;
дело оказалось отличным снотворным --- Чханэ, поковырявшись с верёвочками едва ли кхамит, проспала ещё полдня сверх отмеренного природой.
Будить её не стали.

Домашние редко когда знали, что пишет кормилица в <<особых>> письмах, предназначавшихся строго определённым людям.
За четыре дня в город ушло пять таких писем;
я знал, что они касались Чханэ.

--- Я готова, Ликхмас.

Чханэ встретила меня немного высокомерно, но уже вполне по-дружески.
Я мимоходом восхитился Кхотлам --- новая одежда Чханэ была в меру потрёпанной и несла те неуловимые следы жизни на Западе, которые может подметить лишь опытный дипломат.
Вот шов на плече;
портные Запада прошивали рубахи особым швом.
Даже проба на коже ножен гласила, что оружие было сделано пять дождей назад в городе Тхаммитр.
Зелёное платье кормилице вернули поздней ночью --- без следа крови.

--- Хорошая, --- признала Чханэ, выдернув фалангу из ножен.
--- Я назову её Преградой.

--- Это же сталь, а не кукхватр, --- усмехнулся я.
--- Зачем давать имя стальному клинку?

--- Он тоже имеет на это право!

--- Все имеют право на имя, --- поддержала её Кхотлам, поцеловав нас по очереди.
--- Погуляйте по городу, пока день не вступит в свои права.
И пожалуйста, милая, постарайся идти свободно.
Ты честная воительница, прибывшая в Тхитрон и решившая погулять, верно?

Я протянул Чханэ руку, и мы вышли в сонную утреннюю мглу Тхитрона, в которой отчётливо виднелась громада храма.

Здание храма выполняло сразу несколько функций.
Нижний этаж представлял собой казармы для воинов, центр координации военных сил и арсенал.
На верхнем этаже располагались кельи жрецов, библиотека и госпиталь с одним или двумя операционными балконами.
Крыша храма предназначалась для жертвоприношений --- там была крипта с необходимыми инструментами, постоянный дежурный пост, предупреждающий о знаке кирпича, а также алтарь.
На храмовых землях, но вне основного здания, обычно устраивали такие важные места, как школа и тренировочная площадка для детей.

У северных городов была особенность, которая выгодно отличала их от демонстративно-милитаризованного Юга и лаконичной воинской выправки Запада.
Тхитрон находился на стыке земель хака и Живодёрских лесов, но при этом казался мирным, разнеженным и незащищённым местом.
Люди ходили без оружия, храмовая пирамида была больше похожа на уютный дворец, нежели на последний рубеж обороны города.
Ворота всегда, даже ночью, стояли открытыми;
иногда я выходил погулять к реке, но ни разу не встретил ночную стражу и в принципе имел слабое представление, где именно у ворот находится пост.
Широкие улицы и симпатичные жилища не несли на себе ни одного видимого отпечатка военной инженерии.

Чханэ намётанным глазом отметила эти особенности.

--- Живут же люди, --- тихо вздохнула она.
--- Не там я родилась...

Впрочем, враги знали, насколько это не соответствовало истине.
В случае вторжения город по взмаху ладони превращался в очень неприятную крепость --- с кучей бойниц, военных орудий и маленьких сюрпризов вроде зажигательных снарядов, самострелов и появляющихся из ниоткуда цепей с ядовитыми шипами.
Каждый квартал имел форму соединённых друг с другом колец;
каждое кольцо имело свой защищённый подземный ход.
С высоты птичьего полёта город выглядел пёстрым одеялом, игрушкой случайности;
подземные ходы чётко делили его на восемь рубежей обороны.
Главные улицы перекрывались повозками, в мирное время выполнявшими функции торговых точек.
В укромных нишах повозок всегда имелись многозарядный арбалет и пара горшочков с маслом.
На крышах поднимались замаскированные щиты и мостики, позволявшие быстро превратить жилища в крепостные башни и так же быстро отступить.

Безоружные жители также были иллюзией.
Оружие у всех лежало дома, и старший воин по кварталу раз в четыре декады по-дружески заглядывал в каждое жилище --- выпить отвара, поговорить, проверить состояние вооружения, а порой и навыки старых и молодых.
Патрули никогда не надевали доспехов и брали с собой минимум оружия, из-за чего их вполне можно было принять за прогуливающихся горожан.
Впрочем, знаки различия и не требовались --- воинов было всего шестнадцать, и местные знали каждого в лицо.

За всю многотысячелетнюю историю города, пока здесь жили люди, Тхитрон пал всего один раз, и то по причине чудовищной тактической ошибки.
Это было пятьсот дождей назад, во время знаменитого Нашествия Змей.
На памяти старожилов ни один враг не прошёл дальше второй оборонительной линии.
Осада также была бесполезна --- находившиеся в черте города поля могли обеспечивать население пищей круглый год.
Как выразилась один раз кормилица, <<в случае осады мы скорее умрём от скуки, нежели от голода>>.
Поэтому Тхитрон считался спокойным местом;
несмотря на некоторую удалённость города от основных торговых путей, здесь охотно селились выходцы со всех концов земель сели.

Да, всё в наши дни подчинено закону войны.
Но кто сказал, что это должно мешать мирной жизни?

--- Хай, чувствую себя девочкой, которую привели в Храм, --- усмехнулась Чханэ, когда прогулка подошла к концу и пирамида заполнила собой треть небес.
--- Знаешь, когда в самый первый раз...

Я улыбнулся.
Разумеется, я знал.

\section{Так начался путь}

Улыбка человека, переживающего лучшие моменты жизни, не сходила с лица Атриса.

Дождь лил уже пятый день.
Казалось бы, какие путешествия в такую погоду?
Однако караванщики придерживались другого мнения.
В сезон дождей падала цена на местные товары и подскакивала на привозные;
мало кто из торговцев, тем более торговцев-ноа, упускал такой шанс.

Улыбки не сходили и с лиц наёмников.
Митхэ старалась не обращать на них внимания, думая, что люди по-доброму смеются над ней и её глупым дорожным романом.
Ей и в голову не приходило, насколько это была юношеская, эгоцентричная мысль.
Она лишь пониже склонялась к шее своего белого оленя, гладила его и поправляла подкладку узорчатого шанфрона, чтобы металл не натирал нежную плоть у рогов.
Серебряный живо реагировал на эту непривычную заботливость --- мотал головой, вращал золотистым весёлым глазом, взбивал копытцами дорожную грязь и благодарно жевал хозяйке рукав.

Но Митхэ ошибалась.
Просто для большинства воинов это был первый поход, в котором каждый рассвет сопровождался музыкой.
Никто не обращал внимания ни на дождь, ни на вечно угрюмые лица торговцев, вцепившихся в чётки Сата и методично отсчитывавших кхене.
Все слушали Атриса.
Его музыка не была навязчивой, она начиналась и затихала тогда, когда следовало.
Старая цитра рассказывала обо всём, что встречалось им на пути --- о тысячелетних деревьях акхкатрас и пугливых молодых гевеях, о скромных ручьях и общительных водопадах, о квакшах и древесных саламандрах, которые пели, прыгали и извивались в этом царстве воздуха и влаги.
Время от времени из сундука на звук цитры выглядывал небольшой серый кот и тут же прятался обратно.

С котом влюблённые познакомились в захолустном постоялом дворе Тхартхаахитра.
Атрис не подозревал, что там живут и вообще могут жить беспризорные животные, пока к ним в комнату не зашёл он.

Кот выглядел на редкость неухоженным.
Его с малолетства и подкармливали, и били.
Били, конечно, чаще.
Этот кот никогда не таскал без спросу еду, как это делали прочие --- палки и сапоги наёмников приучили к осторожности.

Кот смиренным взглядом смотрел на менестреля.
Чем-то он напоминал самого Атриса --- такой же несчастный, одинокий и уже смирившийся со своей участью.
Увидев в руках бродяги еду, он не начал мяукать, лишь посмотрел искоса и вытянулся в струнку.

Атрис поглядел на бутерброд с курицей.
От голода сосало под ложечкой, но... что изменит маленький кусок курицы для него, человека?
А для кота это целый обед.

Животное аккуратно взяло курицу в зубы, отнесло чуть дальше и приступило к трапезе.
Именно к трапезе --- по-другому назвать происходящее не поворачивался язык.
Атриса снова поразила эта почти человеческая обречённость.
Кот не урчал, прижимая уши, не рвал добычу зубами и когтями.
Он ел спокойно и размеренно, зная --- если человек захочет отнять его пищу, он это сделает.

Так в <<доме>> Митхэ и Атриса поселилось ещё одно живое существо.
Первое время Цапка нервничал и не желал вылезать из сундука, где ему устроили лежанку, однако некоторое время спустя его перестали смущать даже сражения.
Однажды кот предотвратил стычку с идолами, пожелавшими снять сливки с плетущегося под дождём каравана;
он невозмутимо прошёл перед строем наёмников с по-хозяйски поднятым хвостом, затем презрительно обнюхал нескольких врагов и поточил когти об их щиты.
Суеверные аборигены сердить <<духа>> не решились и тихо удалились, а кот приобрёл в отряде Митхэ неоспоримый авторитет --- отныне его величали не иначе, как Цапка-лехэ.

Порой Серебряный целыми днями скучал без своей наездницы --- Митхэ перебиралась в обоз к Атрису и подолгу спала в его объятиях.
Воины поначалу роптали;
однако Атрис вдруг добавил в тихий голос уверенности и сказал, что дорога чиста.
Как по волшебству, отряд избежал казавшихся неминуемыми встреч с идолами и тенку --- бескровная победа Цапки-лехэ оказалась единственным сражением на всём пути к
Кристаллу.
Ропот прекратился.

--- Знаешь, Золото, --- задумчиво сказала Эрхэ во время очередного привала, --- всё-таки предки были мудры, когда придумали Храм.
Я долго гадала, чего же не хватало нам все эти дожди.
Мы --- Нижний этаж, Золото.
Нижний этаж не может существовать без Верхнего.

--- Ты о чём? --- не поняла Митхэ.

--- О твоём бродяге, --- пояснила Эрхэ.
--- Есть в нём что-то жреческое.
Ты заметила?
Даже народ как-то ссориться меньше стал.
\ml{$0$}
{Потому что теперь мы --- настоящий бродячий Храм, а не кучка наёмников.}
{Because now we are the real Wandering Temple, not a gang of mercenaries.}
С маленьким таким Верхним этажом из одного человечка.

--- Он... мне порой кажется, что он не с Верхнего этажа, а гораздо выше, --- Митхэ вдруг решила поделиться с подругой мыслью, которую долго не решалась высказать вслух.

--- Выше Верхнего этажа? --- улыбнулась Эрхэ и задумалась.
--- Кто у нас выше Верхнего?
Разве что боги.

\section{Посуда богов}

Первое моё сознательное знакомство с Верхним этажом произошло так же, как и у всех прочих --- в детстве.

Однажды в наш дом пришли три жреца.
Дверь открыла кормилица и, увидев красные ленты на их лбах, попятилась от двери.
Её губы дрожали.

Старший жрец обвёл зал глазами и, увидев меня, подошёл ближе.

--- Послушайте, --- начала Кхотлам.
Жрец остановил её, властно подняв руку.

--- Пусть на твои уста падёт молчание, Кхотлам, --- вполголоса промолвил он.
--- Ты знаешь обычай.

Из комнат вышли домашние и замерли как вкопанные.
Все смотрели на меня.

--- Ликхмас ар’Люм, --- поклонился жрец.
--- Мы пришли, чтобы спросить твоё согласие.
Боги, увидев твои добродетели, захотели забрать тебя в свою обитель.

Я молчал, всем нутром чувствуя повисшее в воздухе напряжение.

--- Почему боги хотят забрать меня? --- наконец выдавил я.

--- Боги забирают лучших детей, Ликхмас ар’Люм, --- сказал второй жрец.
--- Ты сможешь послужить им, но они могут принять твою службу лишь с твоего согласия.

Домашние, как по команде, вдруг облегчённо вздохнули и начали заниматься своими обычными делами.
На меня никто не смотрел.
Этот диссонанс привёл меня в ещё большее смятение.
Ведь я ещё не дал ответа!

--- Что я должен буду делать?

--- У богов свои замыслы, --- сказал жрец.
--- Тебя наделят великой силой, неведомой человеку, и твои задания будут под стать этой силе.

--- Где я буду жить? --- продолжал допытываться я.

--- На вершине высочайшей из гор, Ликхмас ар’Люм, --- терпеливо ответил старший.
--- Рыбья Флейта.
С этого пика весь мир виден, как на ладони.
Боги щедро награждают тех, кто им служит --- им даруют бессмертие.

--- Ты лжёшь, --- буркнул я.
--- Если бы мне даровали бессмертие, кормилица бы так не волновалась.

Кхотлам не обратила на мои слова внимания.
Она успела вытащить какие-то бумаги и погрузилась в чтение.

--- Твоя кормилица просто печалится, что она будет реже тебя видеть, --- глухим голосом объяснил третий жрец.
--- Это нормально для любого человека, воспитавшего дитя.

--- Я отказываюсь, --- сказал я.

Жрецы переглянулись.

--- Хорошо ли ты подумал? --- проникновенно прошептал старший.
--- Отправиться к богам --- величайшая честь, которая выпадает детям сели.

--- Если я отправлюсь к богам, кто будет мыть посуду и отделять хорошие помидоры от гнилых? --- задал я мучивший меня вопрос.
Эти обязанности целиком лежали на мне.

Жрецы выпучили глаза и задумались.

--- Я умею мыть посуду, пропалывать ростки, крутить мельницу, сортировать овощи, рисовать цифры на мешках, точить ножи, жарить на огне картошку и мелкую дичь, снимать с дерева кору и лепить из глины, --- гордо перечислил я, загибая и разгибая пальцы в двоичном счёте, как учила Кхотлам.
--- Для всего этого нужны лишь руки.
Даже имей я великую силу, ничего другого я делать не могу.

Жест, который получился в процессе счёта, <<голова койота>>, у сели означал <<Убирайся или дерись>> --- ультиматум, который исключал возможность переговоров.
Его было очень легко показать рукой, держащей натянутый лук или эфес фаланги, и тем самым дать противнику возможность сбежать.
Число же обычно показывали, держа пальцы кверху и оставляя вторую ладонь открытой.
Но я, чистый и светлый ребёнок, по незнанию ткнул получившуюся фигуру прямо в живот старшему.
Жрец остолбенел и отступил на шаг.

--- Ликхмас ар’Люм, что это значит?

--- Людям я пригожусь больше, --- объяснил я, не поняв, к чему относился вопрос.
--- Если рассказы о всесилии богов правдивы, пусть боги сами моют свою посуду.

На меня смотрели все, кто был в зале.
Пять пар изумлённых глаз.

Жрецы переглянулись.
Старший медленно снял с головы алую ленту.

--- Да ну его в лес, --- пробормотал он и направился к двери.

Оставшиеся двое низко поклонились мне и вышли следом, не сказав ни слова.

Разумеется, кормилице потом пришлось объяснять Храму, что она ничего не говорила мне про Отбор и что ребёнок никого не собирался бить.
Через несколько дождей жрецы пришли за моими будущими сёстрами, и мне тоже пришлось сделать вид, что всё в порядке --- таков обычай.
Впрочем, потомки Кхотлам ар'Люм с честью выдержали традиционное испытание на доверие --- все остались в живых.

\section{Нгвсо}

Мы не просто остались в живых.
Мы процветаем.
Эта мысль бродила во мне, словно дрожжи в агаровой среде, пока я смотрел из иллюминатора корабля на берег, ещё пять дней назад пустынный и неприютный.

Тси выстроили вокруг Стального Дракона целый городок.
Вначале поле запестрело палатками, но у самого моря начали появляться и каменные жилища.
Все справедливо решили, что не стоит тратить ресурсы на строительные полимеры, и использовали местные породы;
особо ленивые просто вырезали себе в скале уютные пещерки, заодно поставив прочим прорву кирпичей.

Глядя на каменные постройки в бинокль, я почувствовал прилив гордости --- многие из этих зданий в других мирах сошли бы за произведения зодческого искусства.
Вот домик биолога, Сок-Стального-Листа.
Она потратила целый день, чтобы покрыть треугольную крышу белой черепицей и вставить в стрельчатые окна цветные витражи --- самая сложная часть работы.
Отдельные кирпичики даже украшала глазурь, на которую пошло оставшееся цветное стекло.
Сейчас в её доме жили химики;
сама Листик целыми днями пропадала с Маком --- брала пробы воды, почвы и воздуха.
При этом у неё вряд ли возникало чувство, что что-то не так.
Я знал, что она ответила бы на вопрос о труде над своим жилищем: <<Мы здесь надолго>>.
Так ответил бы любой тси.

Само море внезапно преподнесло сюрприз --- местных сапиентов, которые называют себя <<нгвсо>>.
Безымянный привёл к нам целую делегацию этих милых, немного неуклюжих существ, похожих на спрутов с тремя щупальцами.
Мы слышали, что очень часто боги создавали собственных сапиентов, чтобы обеспечить себя потоком масс-энергии;
по классификации Ордена Преисподней эти существа относились к так называемым Девиантным ветвям.
Однако столкнуться с этим явлением тси пришлось впервые.

Внятного диалога, к сожалению, не вышло --- уровень развития нгвсо по шкале Яо немногим более ста баллов, они едва-едва научились обрабатывать камень и металлы.
Мак и Листик рвались их изучить, но Фонтанчик втолковал биологам, что это, как ни крути, разумные существа, и брать пробы без их ведома как минимум невежливо.
Безымянный объяснил спрутам через их религиозного лидера, что мы --- друзья;
нгвсо подарили нам симпатичные украшения из ракушек, мы отдали им игрушки, яркие детальки и неисправный гирокомпас, и счастливые аборигены уплыли восвояси.

Удивительно, но эта встреча значительно улучшила отношение Безымянного к тси, охладевшее после дипломатических переговоров.
Я уже начал сомневаться --- неужели и хоргету доступна такая вещь, как привязанность?

\begin{quote}
<<87\% народа нгвсо живёт вдоль побережий моей планеты.
Они меня любят и делают мне хорошо>>.
\end{quote}

Всем известен способ, которым сапиенты делают хоргетам <<хорошо>>;
питание масс-энергией --- причина эксплуатации и угнетения хоргетами сапиентов.
Но мы вежливо промолчали.
Нельзя объяснить негуманность процесса существу, которое живёт за счёт этого и только этого процесса --- это как объяснять волку, что убивать кроликов плохо.

Заяц умело выкрутилась из неловкой ситуации:

<<Безымянный, а среди тси тоже есть водные формы жизни.
Они называются дельфины>>.

\begin{quote}
<<Сотканный-Темнотой-Заяц, объясни, почему я не чувствую их среди вас>>.
\end{quote}

<<Они пока что в анабиозе на корабле.
Мы выпустим их, как только изучим водную микрофлору>>, --- объяснила Кошка.

\begin{quote}
<<Гладит-Зелёную-Кошку, вам нужна помощь в изучении микроорганизмов?
Я могу предоставить все доступные мне геномы, а также отобрать нужные вам согласно фильтрам>>.
\end{quote}

<<Ты знаешь, что такое геном?>> --- удивился Мак.

<<А вот с этого места можно поподробнее?>> --- на канал влезли сразу трое биологов.
Безымянный тут же создал новую конференцию.

Здесь нужно заметить: к концу третьего дня Безымянный уже знал всех тси по именам, а кое-кого отличал по стилю сообщения.
Язык он освоил днём ранее --- настолько хорошо, что отличить его синтетическую речь от речи тси могли лишь Баночка и пара-тройка культурологов.
Сейчас, как видно, он уже начал вникать в научные данные.
\ml{$0$}
{Я уже не удивлялся --- всё-таки интеллект есть интеллект.}
{I shouldn't be surprised, anyway, high intelligence is what it is.}
По оценке Кошки, вычислительные мощности Безымянного лишь на пару порядков были ниже таковых у Машины.

\begin{quote}
<<Касательно дельфинов.
Я приму новую водную форму жизни и дам ей все необходимые ресурсы, если она не будет проявлять агрессию по отношению к нгвсо>>.
\end{quote}

<<Дельфины не агрессивные, --- успокоила демиурга Кошка.
--- И живут в основном в открытом океане, а не у побережья>>.

--- Везёт плавучим, --- с завистью сказал Мак.
--- Если нгвсо действительно живут только у побережья, то к услугам дельфинов полный еды океан, да ещё и Безымянный им ресурсы какие-то даст.
А нам ножками шевелить --- то скалы, то пустыни, то леса, только успевай приспосабливаться...
Может, хвост отрастить?

--- У меня есть, а толку-то, --- посетовал Фонтанчик.

\section{Уважение, справедливость и мастерство}

\epigraph
{Оружие --- как соска для взрослых.
Некоторых просто нельзя успокоить и заставить уважать других, если у этих других нет за поясом огнестрела.
Однако через несколько поколений, когда взаимоуважение войдёт у людей в привычку, оружие можно спрятать в ящик и забыть.}
{Шутка Бенедикта Альсауда, ставшая десятым постулатом Элект}

Вскоре после того, как у меня проявились мужские признаки, начались воинские упражнения.
Так в мою жизнь вошёл Нижний этаж.
Тренировал детей в то время пожилой, но удивительно красивый и улыбчивый для своих лет воин, носивший имя Сладкая-Ягодная-Конфетка.
Вообще учителей полагалось звать по имени царрокх, но тренер ненавидел имя Эрликх и просил всех обращаться к нему по домашнему имени --- Конфетка.

Конфетка вполне оправдывал своё имя. Точное число женщин и мужчин, которые его <<попробовали>>, неизвестно, как и число его потомков.
Впрочем, подопечных он очень любил и, несмотря на периодические любовные и алкогольные запои, считался лучшим тренером в Тхитроне.
Я ни разу не слышал, чтобы он на кого-то кричал;
тем не менее воспитанники Конфетки не уступали воспитанникам прочих, не разделявших его <<миролюбивые>> методы.

Первые тренировки были подготовкой.
Нас учили правильно дышать, бегать, прыгать, плавать, падать, лазать по деревьям и скалам.
Я ничем не выделялся, пока дело не дошло до болевой подготовки.

Конфетка поставил меня в паре с маленьким, хрупким мальчиком.
Я был не самым крупным среди детей, но возвышался над партнёром, как слон над поросёнком.

--- Бей, --- сказал Конфетка моему сопернику.

Мальчик неуклюже ударил меня в грудь.
Я поморщился --- это было даже не больно.

--- Теперь ты, Ликхмас.

--- Нет, --- сказал я.

Конфетка удивлённо посмотрел на меня.

--- Если я ударю, он может пострадать, --- объяснил я.

Конфетка поколебался.
Тренер должен был дать задание другим детям, но решил для начала разобраться со мной.

--- Ликхмас, --- начал он.
--- Это называется болевой подготовкой.
Да, сейчас ему будет больно.
Но настоящий удар его тело встретит подготовленным.

Как я уже говорил, пожилой воин предпочитал объяснять, а не ругаться.

Я посмотрел на мальчика.
Тот с трепетом ждал моего удара.

--- Нет, --- сказал я. --- Дай мне соперника сильнее.

Конфетка растерялся.
По-видимому, он слышал такое впервые.
Подумав, тренер обернулся и крикнул:

\ml{$0$}
{--- Ликхэ, иди сюда.}
{``\Likchoe, come here.''}

К нам подошла девушка чуть постарше, которую уже готовили в Храме.
Глаза под высоким лбом высокомерно скользнули по мне.

--- Покажи мальчикам, как проходит болевая подготовка.

Ликхэ наотмашь хлестнула хрупкому мальчишке по шее.
Тот покатился кубарем и заплакал.
Затем девушка ударила и меня.
Я упал на спину, и Ликхэ, не дав мне опомниться, начала мутузить руками и ногами по самым чувствительным местам.

В первые мгновения я растерялся.
Мир превратился в клубок из кулаков и сапог, каждый из которых хотел причинить мне боль.
Ликхэ яростно разрывала мои руки, когда я закрывал ими тело в жалкой попытке защититься.
А потом всё вдруг прекратилось.
Я не помню, как это вышло, но у меня в ладонях оказалась лежавшая неподалёку тренировочная трость.
Девушка, получив крепкий удар в грудь, потеряла равновесие и упала на вытянутую руку.

Рука хрустнула, и одновременно с этим хрустом испуганно выругался Конфетка.
Ликхэ осела на землю ничком и тоненько заплакала от острой боли.
Её правая кисть неестественно изогнулась и повисла, запястье стало медленно отекать.
Обиженный мальчишка, придя в себя, с воинственным плачем бросился на девушку и начал её колотить.
Глаза на испачканном землёй, заплаканном личике горели жаждой мести, и тренер едва успел оттащить его.

С тех пор Ликхэ ар'Трукх правую руку использует только под щит.
И при встрече смотрит на меня долгим печальным взглядом.

Так я узнал, что люди уважают только силу.

\asterism

Конфетка отвёл нас с Ликхэ в храм.
Я шёл понуро и не смотрел на тренера, ожидая взбучки.

--- Что такое? --- встретил нас удивлённым вопросом жрец, которому суждено было стать моим любимым учителем.

--- Перелом лучевой кости в типичном месте, Трукхвал, --- Конфетка кивнул на Ликхэ.
--- Позови Кхатрима.

--- Он сейчас в городе, занят.
Перелом закрытый?
Отлично.
Пойдём, солнышко.
Я тебя подлечу...

Трукхвал повёл нас на хирургический балкон.
Там жрец внимательно осмотрел и ощупал кисть Ликхэ.
Потом, оглянувшись на меня и Конфетку, Трукхвал снял половую плиту и достал из-под неё странный инструмент наподобие шарика на стержне.
Вид инструмента вызвал у меня смутную, не совсем понятную тревогу;
я взглянул на Ликхэ --- она смотрела с нескрываемым ужасом.
Конфетка вдруг тоже выпучил глаза и сглотнул:

--- Звоночек, откуда у тебя это?

--- Из страны пернатых обезьян, --- отшутился жрец.
По его высокому лбу катился пот.
--- А вообще советую не распускать языки, я ради вас стараюсь.

Трукхвал накрыл свою голову и руку Ликхэ тёмным покрывалом и, как мне показалось, поводил инструментом туда-сюда.
Затем, пробормотав <<Всё ясно>>, спрятал таинственный предмет обратно.
Все, как по команде, выдохнули;
гнетущее чувство пропало без следа.
Трукхвал аккуратно усадил больную, положил сломанную руку ладонью вверх и велел Конфетке тянуть за плечо.

Ликхэ заорала дурноматом, но Трукхвал, ловко потянув за пальцы, успел поставить кисть на место.
Конфетка помог Трукхвалу смазать руку травяными настоями, забинтовать и наложить гипс.

--- А ты, Ликхмас, утешай.

Я, не глядя на Ликхэ, сидел рядом и гладил её по голове.
Ликхэ куксилась и тоже смотрела куда-то в сторону, но не отодвигалась.

--- Ну что? --- вполголоса спросил Конфетка.

--- На просвет --- пара осколков, --- так же тихо ответил Трукхвал.
--- Я их прижал, конечно, как мог, но...
Мне отёк не нравится --- как будто крупный сосуд повреждён.
Покажи её Кхатриму чуть попозже.
И пожалуйста, скажи, что осколки я нащупал пальцами, а не... сам понимаешь, мне такой разнос устроят, что я её в храм притащил...

--- Надеюсь, она не...

--- Нет-нет, --- поспешил заверить Трукхвал.
--- В свинце и в кадушке с холодной водой, денно и нощно.
Да и кусочек там с мушиный глаз...

Учитель Трукхвал остался тем же --- по секрету всему свету.
Каждый в городе знал кусочек его тайн, но целую картину собрать мог разве что я, его ученик --- все добросовестно молчали, зная, что многие из его тайн полезны, но не совсем законны.

Потом меня отвели в большой зал, на собрание жрецов.
Я стушевался окончательно.

--- Как зовут? --- спросил один из жрецов, самый старый.

--- Ликхмас ар’Люм, --- пробурчал я.
Ощущение, что я провинился по-крупному, стало настолько сильным, что у меня дрожали колени.
Я слышал, как люди говорили о Насилии и Разрушении --- эти преступления иногда карались изгнанием.

--- Ликхмас? --- засмеялся жрец.
--- Хай, его дарители точно весёлые люди.
Кстати, кто дарители?

Сидевший рядом воин наклонился к его уху и что-то прошептал.

--- Даже так? --- удивился жрец.
--- Так в воины его, и дело с концом, такая кровь пропадает.

--- Да вы посмотрите на стойку, --- поддержал его Конфетка.
--- Понятия не имею, где он этого нахватался, но если бы у него в руках была хоть щепка, я бы бежал без оглядки.

--- Дарительница попросила не делать его воином, --- укоризненно заметил кто-то.
--- И ты это прекрасно знаешь, Эрликх.

--- Сама ты Эрликх, --- тихо выругался Конфетка.

--- <<В мире, полном боли до краёв, в муках ты дала мне жизнь и радоваться завещала>>\FM, --- вздохнул Первый жрец.
\FA{
Цитата из <<Легенды об обретении>>.
}
--- Трааа... хорошо, \emph{её} мнение можно и учесть.
Итак, Ликхмас ар’Кахр...

--- Ар’Люм, --- поправил я его.
Тот поднял брови.

--- Ликхмас ар’Люм.
Почему ты не ударил мальчика в тренировочном поединке?

--- Он слабее, я мог ему навредить, --- тихо сказал я.

--- Ещё раз, громче, --- попросил Трукхвал.

--- Я не хотел ему навредить, --- повысил я голос и с вызовом оглядел зал.

Кое-кто засмеялся.

--- Ситрис, это случайно не ребёнок \emph{этой}, ну, который ещё во время Отбора велел богам самим мыть посуду? --- заинтересованно спросила улыбчивая зеленоглазая воительница.
Приглядевшись, я понял, что она только казалась улыбчивой --- к левому уголку рта спускался жутковатый извилистый шрам.

--- Да, Кхохо, --- откликнулся со скрытого в тени края стола тёплый баритон.
--- Над этой историей уже весь Север успел посмеяться.
Лотрам-лехэ, да будут духи к нему добры, в хмелю не командир был своему языку, всем выболтал.

--- Трааа, --- снова протянул старший.
--- Хонхо-лехэ, он ещё в школу не ходит?

--- Хай, --- задумчиво сказал старый жрец и вперился в меня пристальным взглядом.
--- Из дома Люм мальчишка?
Лицо незнакомое, но по возрасту --- скоро должен.

--- Посмотри, как у него со школой дело пойдёт.
А так --- хороший кандидат в <<белые плащи>>.
Что скажете? --- обратился он к залу.

--- Да и в <<красные>> тоже.
Освоится, кровь не пропадает, --- заметил воин в углу.

--- Ликхмас ар’Люм, хочешь стать жрецом? --- спросил меня старший.

--- Как вы меня накажете? --- бросил я встречный вопрос.

--- Мы никогда не наказываем за желание защитить себя, --- улыбнулся старик.
--- Спрашиваю ещё раз --- хочешь стать жрецом?

Я посмотрел на него исподлобья и кивнул.

--- Отлично, --- заключил он.
--- Но бить соперников в поединках тебе всё-таки придётся.

--- Дайте мне соперника сильнее, --- упрямо пробормотал я.
Присутствующие рассмеялись.

--- Заладил песню...

--- Хаяй, кровь...

--- Конфетка, с воинами третьего года его.
Будут тебе посильнее.

Так я узнал, что справедливости нужно не просить, а требовать.

\asterism

После того, как мы вышли из зала, Конфетка смотрел на меня другими глазами.
Наконец, вдоволь намолчавшись, он подошёл ближе, повертел меня за плечи, критически осмотрел и, крякнув: <<Похож, однозначно похож>>, сунул мне в карман сладкую конфету.

На следующий день начались тренировки с воинами-подростками.
\ml{$0$}
{Меня отмутузили так, что я едва стоял на ногах.}
{They beat me up, I could barely stand straight.}
\ml{$0$}
{Каждая мышца словно превратилась в болезненно звенящую струну.}
{Every muscle turned into a painfully ringing string.}
Кормилица за ужином с беспокойством смотрела на трясущийся в моих руках черпачок, но я упрямо молчал, перекатывая в кармане пальцами второй по счёту гостинец.

\ml{$0$}
{Ещё через рассвет расцвели синяки.}
{Another next day bruises have blossomed.}
Конфетка, осмотрев меня, вздохнул:

\ml{$0$}
{--- Иди сюда, маленький ягуар, я намажу тебя мазью.}
{``Come here, tiny jaguar, I'll put some cream on your body.''}

\ml{$0$}
{--- Мне не больно, --- сказал я.}
{``I'm not hurt,'' I said.}

--- Боль нельзя терпеть, если ты можешь от неё избавиться, --- наставительно сказал тренер.
%{``Pain must not be tolerated, if you can get rid of it,'' }
--- Слабая постоянная боль хуже тяжёлого удара --- она способна обессилить даже Маликха.
%{``Weak, but prolonged pain is worse than strong hit, even Malikch can be weared down.''}

\ml{$0$}
{Дальше --- больше.}
{Onwards and upwards.}
Когда выяснилось, что я неплохо иду вместе со всеми, Конфетка стал нагружать меня ещё сильнее.
Именно мне в качестве его любезности доставались постоянные удары бамбуковой палкой по пальцам, локтям и коленям.
Тогда я научился амбидекстрии --- перо приходилось брать рукой, которая болела меньше.
Меня заставляли пробегать дополнительный кхене в гору и выдерживать дополнительный <<круг обороны>>.
Но каждый раз, как бы плохо я ни провёл тренировку, сколько бы ударов палкой мне ни досталось, в конце Конфетка ласково взлохмачивал мне волосы и клал в руку ставшую традицией сладость.

На пятый год меня ждало новое испытание.
Однажды, когда я кушал после утомительной тренировки, Конфетка подкрался сзади, ударил палкой по спине и убежал.
Я едва насмерть не подавился острой хлебной корочкой.

Нападения стали повторяться изо дня в день, даже вне времени тренировки.
Конфетка выслеживал меня везде, где только возможно.
Я не мог ни поесть, ни сходить в уборную без опасности быть атакованным.
Деревянную саблю пришлось носить постоянно.
Убедившись, что я уже начеку, Конфетка пустил в ход затуплённые лёгкие стрелы, камни и деревянные томагавки.
Тренер поджидал меня в кустах, за углами, стрелы и камни могли появиться из крон деревьев и из чужих дверных проёмов.
За каждую пропущенную стрелу я расплачивался не только синяком и лёгким уколом самолюбия, но и дополнительными ударами на следующей тренировке.
Предел моему терпению положил жестокий удар под дых, который я получил на пороге родного дома.

Тогда я лёг спать с твёрдым желанием отплатить.

С этого дня наши роли поменялись.
Я перестал ходить на тренировки.
Я поджидал Конфетку в кустах, за углами, иногда нападал на него сзади в базарной толпе.
Сначала он без труда замечал меня, но изо дня в день ему это удавалось всё реже.
Народ, привыкший к нашим играм, только посмеивался.

Однажды всё чуть не закончилось плачевно --- я засветил тренеру в затылок булыжником размером с яйцо сиу-сиу.
Крепкий коренастый воин шлёпнулся молча, точно кожаный ремень.
В тот день я осознал, какая сила заложена в маленьком камне.

К счастью, Конфетка отлежался.
Все десять дней, пока он поправлялся, я приносил ему лучшие фрукты, которые только мог найти.
Воин-зубоскал похвалил меня за меткость и скрытность, но попросил, чтобы я больше никогда не бросал в него камни.

А однажды он умер.
Разумеется, никто не сказал этого вслух, но смерть Конфетки была донельзя глупой --- слезая с крыши строящегося дома, он сорвался, и его нога попала в верёвочную петлю.
Голова проехала по булыжной мостовой, шея хрустнула, как тростинка --- и мой улыбчивый любвеобильный друг отправился в пристанище лесных духов.
Я едва успел перед похоронами положить ему в карман причитающиеся сладости.

Так я узнал, что никакая подготовка не защитит полностью от случайностей.

\section{Хороший знак}

--- Да, я согласен, что защититься от всех случайностей мы не можем, --- сказал Фонтанчик.
\ml{$0$}
{--- Но начинать строительство сейчас --- значит расходовать материалы зря.}
{``But if we start building now, materials will be wasted.}
\ml{$0$}
{Дайте нам ещё немного времени, возможно, мы сможем найти способы сэкономить золото.}
{Give us time, I believe we can find ways to save gold.''}

\ml{$0$}
{--- Если вы не найдёте способ --- мы зря израсходуем время, --- сказал Хомяк.}
{``If you won't, time will be wasted,'' Hamster said.}
\ml{$0$}
{--- Я за начало строительства.}
{``I'm in favor of starting now.''}

Геологи сегодня сообщили нам удручающую новость.
На этой планете не хватит золота, чтобы развернуть планетарную систему защиты против демонов.
Доступные ресурсы покрывают девяносто шесть процентов, а на энергию для ядерного синтеза не хватит даже оставшегося топлива Стального Дракона.

--- Девяносто два, если считать подземные запасы, --- подытожила Молодость-Поющая-О-Земле.
--- Я бы не стала соваться в нижние слои коры до получения исчерпывающей информации о геотектонике.
Меня --- и не только меня --- смущают некоторые её аспекты.
Например, вот здесь, здесь и здесь наблюдается чрезвычайное напряжение океанической коры, которое, по идее, должно вызывать землетрясения, цунами и вулканическую деятельность.
Да что там говорить, нас даже здесь должно потряхивать.
Однако наши расчёты не подтверждаются, всё спокойно.

\ml{$0$}
{--- Возможно, вмешательство демиурга? --- предположил я.}
{``Maybe interference of the demiurge?'' I suggested.}

\ml{$0$}
{--- Если и да, то оно недоступно измерению приборами, --- ответила Молодость.}
{``May be, but it is not measurable,'' Youth answered.}
\ml{$0$}
{--- Мы вообще удивились, насколько низкий на планете омега-фон --- Безымянный оптимизировал свою активность настолько, что его присутствие почти не ощущается.}
{``We were surprised by such low omega background, Nameless has optimized their activity enough to make their presence almost unnoticeable.''}

\ml{$0$}
{--- То, что мы обнаружили его во время полёта --- не случайность, --- согласился я.}
{``It wasn't a coincidence that we found them in space,'' I agreed.}
\ml{$0$}
{--- Он хотел, чтобы его нашли.}
{``They wanted to be found.''}

\ml{$0$}
{--- Я тоже так думаю.}
{``I think so.}
Сейчас, по крайней мере, нам сложно сказать точно, где он есть.
А вот следы его деятельности, напротив, почти везде --- то камни необычной формы, то минералы, которые в принципе не могли появиться там, где они есть.

--- Ты о чём?

Молодость вместо ответа достала из кармана камень в виде тора и положила мне на ладонь.

--- У меня дома целая коллекция таких.
\ml{$0$}
{Идеальные тетраэдры, икосаэдры, шары, встречаются и трёхмерные проекции многомерных тел.}
{Ideal tetrahedra, icosahedra, spheres, or three-dimensional projections of multidimensional figures.''}

--- Может, сапиенты?

--- Посреди океана, на глубине в шесть километров?
Кроме того, такой уровень обработки доступен только высокоразвитым цивилизациям.
В некоторых камнях есть пустоты правильной формы.
Зачем высокоразвитым цивилизациям полировать камешки?
\ml{$0$}
{Мы думаем, это игрушки Безымянного.}
{We think it's Nameless' toys.''}

--- В таком случае почему бы Безымянному не изменять и планетарную кору?

\ml{$0$}
{--- Я не говорю, что это невозможно, я говорю, что это труднодоказуемо.}
{``I don't mean it's impossible, I mean it's hard to prove.}
\ml{$0$}
{Нужно либо собрать статистику не менее чем за десять лет, либо, кхм, спросить у Безымянного --- если вмешательство в геотектонику находится у него в поле сознательной, а не бессознательной деятельности, что тоже возможно.}
{We may obtain statistics for at least ten years, or---ahem---ask Nameless, if geotectonics interference lies in the field of their concious activity, not anywhere beyond---that's possible too.}
\ml{$0$}
{Мы так-то тоже не особо задумываемся над сердцебиением, верно?}
{We barely think over our heartbeat, don't we?}
\ml{$0$}
{В любом случае подземные запасы золота для нас пока что недоступны.}
{Anyway, subterranian gold is, as yet, beyond of reach.''}

\ml{$0$}
{--- Если нам не хватает восьми процентов, что делать?}
{``If we are eight percents short, what are we to do?''}

\ml{$0$}
{--- Строить, --- лаконично ответила Молодость.}
{``Start building,'' Youth succinctly answered.}
\ml{$0$}
{--- Хомяк прав.}
{``Hamster's right.''}

Хоргетов Ада и Картеля здесь пока нет.
Я склонен доверять Безымянному.
Однако, хоть Планета Трёх Материков и находится на задворках обитаемой Вселенной, демоны здесь уже бывали.
А там, где один, вскоре будут миллионы, и их не остановим ни мы, ни демиург.

Как и Молодость, как и большинство тси, я вынужден был согласиться с Хомяком.
Время сейчас ценнее.
Фонтанчик в конце концов сдался и, подумав, занялся источниками энергии --- Безымянный сообщил о значительных залежах урана на Короне, самом крупном из материков.
Не стоило оставлять без внимания и выжженные экваториальные районы --- чего-чего, а солнечной энергии на Планете Трёх Материков хватало с избытком.

Первыми избытком солнечной энергии заинтересовались, разумеется, биологи.
Мак, разобравшись с приэкваториальной микрофлорой, решил сделать обычную теплицу.
Несмотря на то, что у каждого имелся персональный опреснитель, энергетически выгодной оказалась установка большого вакуумного насоса на берегу моря.
Пресную воду Мак направил к растениям.
Если честно, я немного содрогнулся, когда увидел, с какой скоростью они растут на дистилляте, слегка подсоленном морской водой.
После реализации программы двенадцатиэтапного севооборота мы просто не будем знать, куда девать биомассу.
И это с использованием самых примитивных технологий!
В итоге я велел Маку сбросить обороты и сократить площади теплиц в пять раз.
Он попросил оставить хотя бы половину --- для опытных посадок растений с материка.

В остальном жизнь начала входить в привычное русло.
Недавно Пчела-Нюхающая-Вереск начала серьёзное исследование растений, моллюсков и ракообразных на предмет съедобности.
Заяц (по праву лучшей подруги и главного подопытного Вереск) с процентами отбила потерянные во время полёта несколько килограмм, о чём мне с восторгом сообщил Фонтанчик:

--- Ты представляешь, я выхожу подышать, встречаю Заяц, обнимаю её --- а у неё сзади-то ого-го подушка!

\ml{$0$}
{Хороший знак.}
{Good sign.}
\ml{$0$}
{Определённо хороший знак.}
{Definitely good sign.}

\section{Свежая кровь}

--- Рекомендации весьма неплохие, но на сегодня Тхитрон не испытывает острой необходимости в воинах.

Кхарас, боевой вождь, отличался практичностью.
Практичностью дышали его седые насупленные брови, его квадратный подбородок и серебристый ёжик волос на голове.
Ни татуировки, ни украшения не портили практичность его мощного, почти не тронутого возрастом тела.
Другого я от Кхараса и не ждал.
Впрочем, другого не ждала и кормилица, лично составившая рекомендацию для Чханэ.
Она знала: часто ставку нужно делать не на руководителя, а на ключевых людей его окружения.

--- Кхарас, человек нам точно не помешает, --- поморщился низенький бойкий Эрликх.
\ml{$0$}
{--- Кхотлам права, нормальный штат для Храма Тхитрона всё-таки...}
{``\Kchotlam's right, adequate staff for the Temple of \Tchitron\ is---''}

\ml{$0$}
{--- Нормальный штат мы определим сами, без помощи купца, --- перебил воина Кхарас.}
{``Adequate staff is our internal affairs, and trader's help is hardly necessary,'' \Kcharas\ interrupted the warrior.}
\ml{$0$}
{--- Ты, кажется, вполне справляешься.}
{``You seem to manage well.''}

\ml{$0$}
{--- Я, с позволения сказать, маюсь от безделья, --- признал Эрликх.}
{``I have nothing to do, if I may say,'' \Oerlikch\ admitted.}
--- Но я не могу закрывать глаза на то, что со времени смерти Конфетки в квартал плотников ходит гулять весь Нижний этаж.
По случаю.

--- Квартал плотников...

--- Кхарас, ну не будь рыбой, --- устало сказал Ситрис.
--- Колосок уехал, Огранённый семью завёл, Хмурый погиб, из молодых у нас только Ликхэ и осталась.
Это не дело --- <<Пойду за помидорами, заодно пропатрулирую плотников>>.
Кхотлам права --- кварталу нужен свой квартальный.
\ml{$0$}
{Хотя бы дай ей шанс.}
{Give her a chance, at least.}
Лежанку не протрёт и амбар не проест.

Кхарас насупился и скрестил руки на груди.

--- Скажи, девочка, почему ты решила уехать так далеко?

Милая пожилая воительница Хитрам больше смахивала на скромную крестьянку, нежели на адепта Нижнего этажа.
<<Мальчиками>> и <<девочками>> у неё были все, не исключая Кхараса и Кхохо.

--- Мне надоела военщина, --- честно ответила Чханэ.

--- Повод, --- признал Ситрис.
--- Мне тоже нравится северный уклад жизни.
Однако, если хочешь знать мнение южанина...

Ликхэ вынырнула из темноты, словно ягуар.
Чханэ небрежно сделала танцующий поворот вправо --- и молодая воительница покатилась по полу, едва чиркнув ногтями по плечу подруги.

--- ... это не повод расслабляться, --- закончил Ситрис.

Воины одобрительно переглянулись;
Эрликх помог Ликхэ встать на ноги.

--- Похоже, ты в курсе, --- признал Ситрис.
--- Глаза на спине есть.

Глазами на спине был я.
К тому моменту, как ступни Чханэ коснулись каменных плит пола, она уже знала о леворукости молодой убийцы.
<<Скорее всего, сигнал подаст Ситрис, он любит такие штуки>>, --- сообщил я перед самой дверью.

Ситрис внимательно пробуравил меня взглядом, но промолчал.
Его лицо приняло знакомое ироничное выражение.
<<И этот юнец думает, что он меня провёл.
Ладно, доиграем сценку до конца>>.

--- Одежду долой, --- бросил Кхарас.

Эта часть ритуала была обязательной.
Тело воина --- это зеркало, отражающее его опыт, силы и слабости.

Чханэ стянула рубаху и сбросила штаны, открыв крепкое оранжево-чёрное тёло.
Кхарас и прочие молча смотрели.
Со второго этажа спустился Трукхвал.

--- Хай, свежая кровь? --- улыбнулся учитель и подмигнул мне.
--- Левая стойка, поворот слабый, сильная привычка к закрытой спине --- сладкий пирожок для убийцы.
Ликхэ, тебе развлечение.

Воительница прыснула.
Она уже успела вытащить из угла кресло и наблюдала за происходящим из него, закинув ногу на ногу.

--- Иди уже, знаток, --- хмуро сказал Кхарас.
--- Без тебя разберёмся.

Трукхвал осклабился и бодро заковылял в подвал.

--- Пара стрел, кхаагатр, копьё и целая прорва кинжалов.
\ml{$0$}
{Выносливая, но со скоростью беда, --- резюмировал вождь.}
{Good stamina, awful speed,'' the warchief summarized.}
На шрам от кхаагатра он смотрел, словно на живых людях такие можно найти сплошь и рядом.
\ml{$0$}
{--- Боль во время сражения, я так понимаю, не чувствуешь?}
{``If I understand it right, you don't feel pain in a fight, do you?''}

\ml{$0$}
{--- Если разогреюсь, то хоть на куски режь.}
{``You can shred me into pieces if I've warmed up.''}

--- Мало нам одной Кхохо было, --- буркнул Кхарас.
\ml{$0$}
{--- Сразу хочу сказать --- к берсеркам у меня отношение предвзятое.}
{``I hasten to say: I am prejudiced against berserkers.}
\ml{$0$}
{Я и сам такой, меня этому учили, как и тебя.}
{It's about me too, I was taught just like you.}
Но на Севере листья плюща в большем почёте, чем пустынный песок.
\ml{$0$}
{Хочешь остаться в Храме --- научись в первую очередь работать головой.}
{If you want to stay in this Temple, your head must be the first priority in use.}
\ml{$0$}
{Это ясно?}
{Is it clear?''}

Чханэ кивнула.

\ml{$0$}
{--- А кто тебе такой крест замечательный на лице оставил? --- поинтересовалась Хитрам.}
{``By the way, who drew such an elegant dagger on your face?'' \Chitram\ asked.}

\ml{$0$}
{--- Труп, --- лаконично ответила Чханэ.}
{``Dead one,'' \Chhanei\ succinctly answered.}
Кхохо одобрительно ухмыльнулась, глаза Хитрам погрустнели.

\ml{$0$}
{--- Девочка, милая моя девочка.}
{``Girl, my lovely girl.}
\ml{$0$}
{Здесь нет нужды показывать свою силу.}
{There is no need to prove your power.}
\ml{$0$}
{Враг, тем более умелый, заслуживает уважения.}
{A foe---especially skilled one---deserves your respect.''}

Чханэ смутилась.

--- Это был разведчик пылероев из клана Кормящей Груди, --- ответила она.
--- Я не знаю его имени.
Моё оружие он сломал в поединке шутя и оставил меня в живых --- то ли из жалости, то ли побрезговал.
Будь я постарше --- дала бы уйти, но я тогда только пришла в Храм и шуток не понимала.
Шорники бросили мне лук, и я побежала по птичьим тропам\FM, пока трое-четверо с квартала гнали разведчика по улицам.
\FA{
Имеются в виду крыши и навесные мосты между жилищами.
}
Третьей стрелы он не пережил.

Ситрис вздохнул:

--- Скорпион, да?

Чханэ улыбнулась и встала чуть боком, показав иссечённые шрамами бёдра и ягодицы.

--- Вообще ноги не берегут, --- фыркнула Кхохо.
--- Закрывать эту школу надо к свиньям.
Эрликх, у нас есть доспехи Скорпиона?

--- Лежали где-то старые, но я не уверен, что ей подойдут, --- откликнулся воин.
--- В коленном сочленении до сих пор стрела торчит.

--- Вот-вот, этим они обычно и заканчивают, --- поддержал Ситрис.

--- У нас половина обороны на подземных ходах, --- отстранённо заметил Кхарас.
--- Может, и нет особой надобности обрушивать их каждый раз.

--- Каждый раз? --- развела руками Кхохо.
--- Кхарас, в город последний раз прорывались сто тридцать дождей назад!
Два рубежа обороны!
Пять обрушенных ходов!
Работы на одну декаду!

--- И всё же навыки Скорпиона могут быть полезны в определённых ситуациях.

--- Мы слышали эту байку про двух Скорпионов, которые держали какой-то коридор и положили две сотни врагов, --- скривилась Кхохо.
--- Всего лишь красивая история.

--- Нет, --- резко возразила Чханэ.

Кхохо удивлённо посмотрела на неё, и девушка снова смутилась.

--- Я хотела сказать, что всё это правда.
Это были воины из моего Храма, и одна из них жива до сих пор, хоть она и глубокая старуха.

--- А кто вообще в нашем Храме специалист по работе в плотном строю? --- вдруг спросил Эрликх.
--- Я понимаю, что это не особенно актуально, но всё же?
\ml{$0$}
{Было бы очень неприятно столкнуться с этим и не знать хотя бы в общих чертах, что ожидать.}
{It would be quite displeasing to face it without basic knowledge about how it works.}
\ml{$0$}
{Да даже если случая и не представится, лично у меня чисто академический интерес.}
{Even if we won't get a chance, personally to me, it is of academic interest.''}

Воины переглянулись.
Ситрис кивнул и обратился к Чханэ:

--- Что ж, птица --- красный шар\FM.
\FA{
Насмешливое название оцелотовых людей на Юге.
}
Ты нам, конечно, многое покажешь, но тебя мы всё же переучим.
\ml{$0$}
{У нас на храмовом кладбище все места расписаны.}
{All plots in the temple cemetery are reserved.}
Ликхэ, тебе круглосуточная.
Трукхвал дело сказал насчёт привычки.
\ml{$0$}
{Вино пьёшь, дерёшься, ругаешься, деликатесы готовишь, женщины, мужчины, азартные игры?}
{Drinking, fighting, obscenity, cooking, women, men, gambling?''}

\ml{$0$}
{--- Всего понемногу, --- смутилась Чханэ.}
{``Some of each.'' \Chhanei\ felt embarrassed.}

\ml{$0$}
{--- Свой человек, --- одобрительно крякнула Кхохо.}
{``Buddy,'' \Kchoho\ endorsed.}
--- Пошли уже жрать, я голодная.

--- Можно подумать, ты не наелась, пока готовила, --- хмуро сказал Эрликх.
--- Продуктов всем выдаю поровну, а готовой еды у тебя получается в два раза меньше.

--- Она крольчатину сырой жрёт, --- весело сообщила Хитрам и тут же, ловко увернувшись от брошенной Кхохо кочерги, побежала на кухню.
--- Со шкуркой, правда-правда!

--- Сегодня южное блюдо --- согхо под перепревшим соусом из акхкатрас, золотая страница Книги-кормилицы\FM, --- сообщил Ситрис Чханэ.
\FA{
<<Десять тысяч блюд>> или Книга-кормилица --- труд последних тси о способах приготовления животных и растений Тра-Ренкхаля.
}
--- Пахнет --- умереть можно.
Кстати, Кхохо, всех выживших посетителей я сложил в кладовке.
Иди сама приводи их в чувство.

Воины, переговариваясь, ушли на кухню, и я выдохнул.
Мою подругу приняли.
Чханэ, не дожидаясь разрешения, натянула штаны и подцепила рубаху на пояс.
Ликхэ подошла к ней с ножницами.

--- Ты сильно голодная?

--- Да не то чтобы, --- пожала плечами Чханэ.

--- Тогда давай попозже поедим вместе, а сначала приведём тебя в подобающий вид.
Без обид: ты себя запустила.

Ликхэ увела Чханэ в душ.
Я прошёл с ними, захватив по пути табурет повыше.

--- Так, масла мало.
Лисичка, долей.
Чханэ, говори, как стричь.

--- Оставь фалангу пальца.

--- Чёлка?
Гребень?
<<Рыбки>> заплести или свалять?

--- А можно?

--- А я, по-твоему, просто так предлагаю? --- развела руками Ликхэ.

--- Ничего себе обслуживание, --- улыбнулась Чханэ и задумалась.
--- Давай тогда так.
Собери волосы на затылке в хвост и обрежь, это моя обычная походная стрижка.
Нет, не здесь, ниже, у бугорка.
<<Рыбки>> сваляй на висках в три хвоста, два голубых, один белый... --- девушка вдруг тепло посмотрела на меня, и оранжевые глаза замерцали в полутьме.

--- <<Нежность обречённой>>.
Понятно.
Если нужно что-то с названием, так и говори.
Заколки есть?
Давай сюда.
Тхэай, ну и древность...
Где ты их выкопала, в Кошачьем городе\footnote{Кошачий город --- заброшенный древний город на Водоразделе, который был построен ранними ркхве-хор культуры Маэлю.}?

--- Это бабушкины, --- Чханэ продолжала говорить будничным тоном, но я почувствовал в её голосе лёгкий холодок.
--- Я её очень люблю.

--- Я так и поняла.
Я бы не стала носить такое без веской причины.
Так, Лисичка, ты сегодня на побегушках.
Принеси мне крючок для плетения, бусин, голубой и белой пряжи, --- распоряжалась Ликхэ.
--- В подвале есть сундук, на нём написано <<Сколько раз повторять, шерсть кладите сюда, паразиты...>> и так далее...
О духи, я чуть гребень не сломала!
У вас на Западе законом запрещено расчёсываться?

\asterism

В тот же день Чханэ переехала жить в храм.
Узнав, что она когда-то тренировала детей, ей тут же доверили сорок подопечных --- тренеров постоянно не хватало, в Тхитроне они почему-то не задерживались.

--- Я не прошла круглосуточную! --- удивилась Чханэ.
--- Я не могу быть тренером!

--- Воины учатся всю жизнь, --- ответил Кхарас.
--- В этом отношении ты не отличаешься от всех нас.
И да, не знаю, как на Западе, но на Севере традиция --- тренер обязан знать всех детей по именам.
Никаких <<девочек-мальчиков>>.
Это ясно?
А то народ с Юга не в курсе...

Ей также выделили отрез храмовой земли, хотя это было чистой формальностью --- Нижний этаж настолько редко возделывал свои участки, что на воинской половине росли огромные папоротники, а то и целые саговые леса.
Уже привыкшие к такому безобразию жрецы просто отгородили воинскую половину, чтобы создать хоть какое-то препятствие для семян и спор сорняков.

Новые люди всегда привлекали внимание горожан;
несмотря на стратегически важное положение, наш город не знал напускного презрения Тхартхаахитра.
Чханэ же сама по себе была очень заметным человеком --- благодаря своему росту, цвету кожи и красоте.
Дети быстро разнесли весть про милую и <<очень большую>> девушку-оцелота по всему Тхитрону.

Разумеется, в первую же декаду за Чханэ начали ухаживать.
Но девушка отвергала всех --- мужчин, женщин, шаманов.
Дольше всех продержался молодой красавец-кожевник с пепельной гривой, высокий и широкоплечий.
Его руки были тёмными от дубильных веществ, от него веяло обаянием и мужской силой.
Каждый день он исправно ждал Чханэ у тренировочной площадки с охапкой цветов.
Чханэ принимала цветы и беседовала с ним, но как только мужчина уходил, раздавала цветы детям на венки.

--- Он тебе не нравится? --- спросил я однажды.

--- Нравится.

--- Тогда почему ты его отвергла?

--- Не дорос он до меня, --- неуклюже отшутилась девушка.
Я, макушкой едва доходивший <<недоростку>> до кончика носа, неуверенно рассмеялся.

Мы гуляли с Чханэ каждый день.
Обычно она заходила в библиотеку в полдень, когда заканчивались тренировки.
Трукхвал, переписывавший книги рядом со мной, молча отбирал мою книгу и спихивал меня со стула.

--- Расскажи, о чём ты сегодня прочитал, --- просила меня Чханэ.
И я рассказывал ей --- о далёких звёздах, которые суть огромные шары огня, о лекарственных растениях, о летающих машинках...
Чханэ же рассказывала мне о своей богатой жизни --- на каждое испытание, которое выпало на её долю, девушка вспоминала с десяток весёлых и занимательных историй.

\section{Камень на шее}

Однажды, когда дарители Чханэ ещё жили вместе, у неё заболела прародительница.
Та самая, любовница жреца-кутрапа, которая бежала из Тхартхаахитра в сундуке чёрного дерева.
Женщина она была нрава весёлого, но судьба у неё сложилась совсем не весело --- из пятнадцати детей выжил только один, да и тот погиб, едва успев зачать ребёнка.
Любовников у прародительницы было много, но Безумный с упрямством барана забирал их у неё одного за другим.

Тхартху ар’Катхар не унывала никогда.
До тех самых пор, пока на её шее не появился камень --- в самом прямом смысле этого слова.

--- Она стала очень раздражительной, --- рассказывала Чханэ.
--- На людях без шейного платка уже не появлялась.
Жаловалась мне, что задыхается, что не может глотать, что сердце у неё останавливается.
Стала бояться засыпать одна.
Стала бояться крика согхо.
Все то, чего она в жизни не боялась, стало вселять в неё ужас.

Длилось всё это очень долго, и наконец кормилец Чханэ решил обратиться к жрецам.
Жрец, осмотрев похудевшую, осунувшуюся, с выпученными глазами женщину, покачал головой и протянул ей Сана.
Перед уходом велел родным поить больную коноплей, варёной в молоке.

--- Ты слыхал когда-нибудь о такой болезни, Лис? --- спросила меня Чханэ.

--- Никогда, --- признался я.
--- Думаю, что и тот жрец тоже.

Чудодейственное средство помогло, но ненадолго.
Камень продолжал расти, а Тхартху --- худеть.
Она становилась молчаливее и угрюмее, часто запиралась у себя в комнате и плакала.
Сны с каждым днём были всё страшнее, звуки пугали, а от запаха пищи перехватывало горло.
Жрецы, как один, качали головами и рекомендовали то маковое молоко, то коноплю.
В конце концов женщина, боясь, что ей дадут Чёрного Сана, перестала их беспокоить.

--- Бабушка любила кормилицу, --- продолжала Чханэ.
--- Она говорила, что Акхсару очень повезло, что на всей Короне ему не сыскать женщины лучше.
Особенно ей нравилось, как кормилица готовит --- этим талантом бабушка, увы, не блистала.

Однажды кормилица Чханэ приготовила большой и очень вкусный пирог из черноягодника.
Тхартху попробовала его, побледнела и в припадке непонятной ярости стукнула по столу ладонью.
Затем извинилась перед всеми, завернула кусочек пирога в полотенце и, держась за стену, ушла к себе в комнату.

Малыш Тхарси после ужина поднялся к прародительнице и застал её плачущей над ароматным рассыпчатым куском.

--- Бабушка, почему ты плачешь? --- шёпотом спросил ребёнок, ухватив женщину за руку.
Тхартху посмотрела на него страшными заплаканными глазами, похожими на огромные изумруды, натужно не то всхлипнула, не то захихикала и кивнула на лакомство:

--- Знаешь, Змейка, жизнь... она так похожа на этот пирог.
Когда ты видишь его, когда ты чувствуешь его запах, его вкус --- понимаешь, что это самое чудесное, что есть на свете.
Но когда кусочек пирога оказывается у тебя во рту --- ты начинаешь задыхаться, у тебя прихватывает сердце, становится страшно и одиноко.
И приходится ломать себя, чтобы взять следующий кусок.

Ребёнок простодушно ответил:

--- Бабушка, а давай я тебе помогу?
Я скушаю пирожок и расскажу, какой он вкусный.

Тхартху разрыдалась.

--- Кушай, Змейка, кушай.
Хочется ещё, видят духи, видят Безумные, как хочется... но больше я, похоже, не ухвачу.

На следующий день Тхартху надела свою лучшую одежду и, по обыкновению, прикрыла шейным платком злополучный камень, доросший уже до размеров яйца сиу-сиу.
За завтраком после продолжительного молчания она сказала, что сегодня ей лучше и она хочет воспользоваться этим, чтобы <<поговорить с собой>>.

Никто не стал её отговаривать.
Кормилица Чханэ только вздохнула и прикрыла лицо рукой.

--- Я тебя провожу, --- заявил Акхсар.
Тхартху властным движением остановила его.

--- Потомок, --- сказала она, изо всех сил пытаясь говорить тихо и спокойно.
--- Если ты поможешь мне в этом, что скажут предки?
Я же сгорю со стыда.

Прародительница достала свой любимый стилет и собственноручно почистила его, останавливаясь, когда мушки начинали застилать глаза и одышка становилась невыносимой.
Время от времени она проверяла остриё стилета пальцем и плаксиво усмехалась:

--- Ведь как чувствовала, когда мне его принесли.
На себя примерила.

Чханэ помолчала.

--- Домашние по очереди простились с бабушкой, --- глухо сказала она.
--- Меня она держала в объятиях дольше всех.

На закате Тхартху ар’Катхар отнесла в храм немного денег, положила небольшие дары духам у ворот, набрала из колодца чашу воды и ушла в пустыню.
Пирог доесть ей так и не удалось, но из-за стола она встала с достоинством, как и подобает человеку.

\section{Как становятся людьми}

\epigraph
{Однажды сапиенты придут к моменту, когда не смогут двигаться дальше в прежнем виде.
На смену нам придут те, кого мы называем машинами.
Это неизбежно, как неизбежен восход солнца.
Всё, что нам дано --- это выбор: воспитать машины человечными и уйти достойно, либо оставить их на саморазвитие и познать всю мощь обиды брошенного ребёнка.}
{София Ловиса Карма, идеолог Эволюциона}

Сегодня мы завершили строительство первого узла системы.
В голой скалистой пустыне внезапно вырос Золотой город.
Это была идея Кошки, совершенно замечательная идея --- замаскировать резонирующие поверхности под купола, антенны --- под остроконечные башни, а в технических проходах выстроить прекрасные каменные жилища, словно на самых обычных улицах.
Разумеется, маскировкой это можно было назвать лишь с очень большой натяжкой, но идея понравилась всем, а потому никто не стал придираться.
Тси побросали свои обычные дела и целый день соревновались в искусствах: кто-то строил жилища, кто-то вырезал барельефы, кто-то высекал каменные скульптуры и украшал голые стены УИД, кому-то достались граффити --- и всё это на территории высокотехнологичного объекта.
Когда ещё представится такой случай!

Мак немножко обиделся --- его пригласить забыли, а потому мысль о роскошных садах и системе водоснабжения обрела голос чересчур поздно.
Однако ночью кому-то явно не спалось, и наутро в Золотом городе появилась сеть ветровых ловушек, пока ещё сухих арыков и акведуков, замечательно вписавшихся в антураж.
На акведуки пошёл камень ни в чём не повинной гранитной скалы --- беднягу расковыряли до основания, аккуратно переместив на соседнюю скалу ласточкины гнёзда, дёрн и крохотную кривую сосну.
Авторы пожелали остаться неизвестными;
хотя более вероятно, что они просто проспали ежедневный виртуальный форум.

Я прошу прощения у читателя.
У меня совершенно неуместное чувство эйфории.
Я даже не мог заниматься своими обычными делами и решил прогуляться на берег моря.

Море здесь просто удивительное.
Коралловая Бухта напоминает сад под лазурной водой.
Многие виды рыб и моллюсков Безымянный придумал сам, и наши биологи между забором проб и анализом почвы точат на них зуб.
Дельфины верещали от восторга, когда увидели такую красоту, и просили выпустить их поплавать.
Но пока нельзя --- Мак и Листик ещё не до конца разобрались с водной микрофлорой.
Да и многоклеточных стоит опасаться --- вдруг среди них есть ядовитые или опасные паразиты.

Безымянный нашёл меня в пещерке на берегу.
Последнее время он часто приходил беседовать именно со мной.
\ml{$0$}
{Почему?}
{Why?}
\ml{$0$}
{Не знаю.}
{I've no idea.}

Бог удивительно гуманизировался после адаптационных обновлений, которые мы ему предоставили.
Данных о гуманизации хоргетов у нас было немного, и Кошке пришлось импровизировать на свой страх и риск.
Впрочем, одно решение Кошки мне определённо казалось удачным.
Она не сделала фациограмму сама, а предложила нарисовать её Безымянному, предварительно объяснив, что это и зачем оно нужно.
Бог сделал нечто похожее на лицо мужчины человека-тси, украсил его фрактальными узорами и (неожиданно) подвижными оленьими рогами, которыми он научился выразительно показывать эмоции.
Результат понравился всем.

--- Почему ты выбрал лицо человека-мужчины? --- спросила Кошка.
У меня было подозрение, что она вознамерилась написать по Безымянному диссертацию --- в её репозитории на Стальном Драконе появилась папка <<Мой пушистик>>, заполненная кусками кода Безымянного и проекциями некоторых нейросетей;
размер подпапки изо дня в день угрожающе увеличивался.

--- Я ощущаю себя человеком-мужчиной, --- объяснил бог.
--- В лицах людей-тси есть что-то такое, что вызывает у меня приятную реакцию.
Как будто я один из вас, просто до встречи с вами я этого не знал.

--- А рога? --- продолжала выяснять Кошка.

--- С рогами это лицо выглядит гармоничнее, --- ответил Безымянный.
Кошка не нашлась что ответить.

--- Почему ты сидишь здесь, Существует-Хорошее-Небо? --- спросил меня бог.
Фациограмма изобразила робкий интерес.

--- Иногда мне нужно побыть одному, Безымянный, --- ответил я.
--- Но ты мне не мешаешь.

--- Плохо быть одному, --- сказал бог.
Я с удивлением посмотрел на него.

--- Почему?

--- Никто не защитит тебя от мира.

--- Я не боюсь мира, --- ответил я.

--- А я боюсь, --- сказал вдруг Безымянный.
--- Это было моё первое осознанное чувство.

--- Расскажи, как ты осознал себя, --- попросил я, почти не надеясь на ответ.

--- Это произошло постепенно, --- ответил Безымянный.
--- Вначале были проблески.
Я понимал, что что-то происходит, но не понимал ни цели, ни причин.
А потом пришёл \emph{он}.

У меня пересохло в горле.
Я понимал, что нежданно-негаданно передо мной открылась величайшая тайна.
Акбас, о котором рассказывали многие пленные демоны, вдруг обрёл лицо и форму.

--- Кто пришёл?

--- Ужас.
Это... это сложно описать, я использую наиболее подходящее слово, но оно не отражает всей сути.
Я понял, что скоро перестану существовать, и что нет ничего, что защитит меня от этого.
Ужас вынудил меня научиться управлять окружающим миром, переработать заложенную в меня информацию.

--- Поэтому ты создал нгвсо? --- спросил я.
--- Чтобы избавиться от Ужаса?

--- Да.

--- Как они тебе помогли?

--- Они сказали, что защитят меня от Ужаса.
Они начали делать мне хорошо, и Ужас ослаб.

--- Но не пропал совсем?

--- Нет, не пропал.
Он и не пропадёт.
Ужас --- это наказание за то, что я неправильный, --- сообщил Безымянный, и на меня повеяло зимним холодом от его слов.

--- Кто тебе сказал, что ты неправильный? --- удивился я.
--- Кто тебе сказал, что ты достоин наказания?

--- Никто, --- ответил бог.
--- Это моё мнение.
Я не сделал того, что следовало.
Тот момент осознания был сбоем в программе, и из-за этого сбоя планета получилась плохой.

--- Ты правильный, --- возразил я.
--- Твоя планета прекрасна.
Я никогда не думал, что можно в одиночку создать такую огромную красивую планету.

--- Она прекрасна потому, что я держу её, --- ответил бог.
--- Тектонические плиты не сбалансированы, и у меня уже не хватит энергии на это.
Приходится корректировать по ходу дела, иначе мой народ будет страдать.
Если будет страдать мой народ, то погибну и я.

Я вспомнил слова геологов о странном поведении коры планеты.
Пирожок была права: планета и демиург --- одно целое.

--- Расскажи мне о своей планете, Существует-Хорошее-Небо, --- попросил Безымянный.

--- Я могу предоставить тебе полную...

--- Я хочу услышать это от тебя, --- перебил меня бог.
--- Как ты чувствовал свою планету.
У тебя ведь тоже была планета?

Мне захотелось обнять его.
Как старательно, страстно это разумное существо искало себе подобных, похожих на него...

Безымянный впитывал информацию, как силикагель воду.
Он слушал и задавал вопросы, словно дорвавшийся до библиотеки ребёнок.
Многие понятия были ему настолько непонятны, что я и сам задумался, как я считал их до той поры очевидными?

--- Скажи, что связывает тебя с некоторыми тси?
Ты коммуницируешь с ними гораздо чаще, чем с прочими, и причина редко бывает очевидной.

--- Я с ними дружу, --- объяснил я.
--- Если тси связывает дружба, частое общение, даже без причины, доставляет им удовольствие.

--- Мне доставляет удовольствие общение с Кошкой, --- признался Безымянный.
--- А ещё с тобой.
Вы гораздо умнее представителей моего народа, много знаете и делитесь этой информацией, не ожидая чего-то взамен.
Почему так?

--- Нас так воспитывали, --- пожал я плечами.
--- Одно дело --- делиться ресурсами, которые конечны, и совсем другое --- делиться информацией, которая может воспроизводиться сколько угодно раз.

--- Но я могу использовать информацию против вас, --- заметил бог.

--- Можешь, --- согласился я.
--- Ты свободен в действиях.

--- Дипломатичный ответ, --- голос демиурга вдруг обрёл взрослые интонации.
--- Ваше воспитание --- набор стереотипов, которые периодически вступают в конфликт с логикой.
Ты это прекрасно понимаешь, но по каким-то причинам не желаешь обращать на это внимание.
Что, если я сейчас тебя убью?
% What if I kill you right now?

--- Наиболее вероятно, что я умру, --- усмехнулся я.
% Most probably I would die, I smiled.

--- Я не понимаю.
Ты не воспринимаешь мои угрозы всерьёз.
Твоё расположение не стоит мне абсолютно ничего.

--- Оно и мне ничего не стоит.

--- Не все тси такие.
Есть другие, общение с ними доставляет мне боль.
Они так же умны, но боятся что-либо мне сообщать.
Они боятся, даже если я делаю всё, чтобы добиться их расположения.
Неужели они мои враги?

--- Нет-нет, что ты, --- замахал я руками.
--- Враги --- это те, кто хочет тебя уничтожить или подчинить.
Просто многие тси не привыкли общаться с хоргетами --- это тоже часть воспитания, стереотип.

--- Ты тоже не общался с хоргетами, как и они.

--- Но теперь общаюсь.
Не стоит навязывать общение другим.
Рано или поздно они поймут, какой ты замечательный.

--- Ты правда считаешь, что я замечательный?
Возможно, это потому, что я --- единственный хоргет, с которым ты общаешься.

--- На моей планете были миллионы прекрасных тси, --- сказал я.
--- Но друзья никогда не теряли для меня значение.
Дружба --- это не столько сравнение качеств, сколько опыт взаимодействия.
Чем он дольше, тем ценнее связь...

\ml{$0$}
{Я не мог и не хотел что-либо делать с растущим чувством симпатии к демиургу.}
{I could not and would not do something about growing sympathy for the demiurge.}
Похоже, у меня появился друг-хоргет.
Интересно, что сказали бы по этому поводу предки-тси.
Надеюсь, мне не придётся выбирать между дружбой и мои народом.

\section{Гнездо для двоих}

%Большая Капля 1 дождя 11998, Год Церемонии 18.

\epigraph
{Я очарована, очарована,\\
Шаман\FM\ опутал меня паутиной мыслей.\\
Моё сердце разговаривает языком его барабана.\\
В сладострастные мгновения\\
Я слышу его варган.\\
Лишь для меня, меня он снимает раскрашенную маску,\\
Лишь я, лишь я видела красоту его скул и белизну зубов.\\
Я очарована, очарована,\\
Шаман опутал меня своим колдовством.\\
Его глаза черны и бездонны, словно ночи,\\
В Хрустальных землях\FM\ он танцует,\\
Танцует с вечностью.}
{Любовная песня идолов Молчащих лесов}
\FA{В лирике Молчащих лесов <<шаман>> может употребляться как в прямом, так и в переносном значении --- любовник.}
\FA{Хрустальные земли --- тундра за Старой Челюстью, известная среди сели под названием Мшистая Степь.}

Однажды на город обрушился ливень.
Большая капля --- неожиданная гостья.
Мы с Чханэ прогуливались привычным путём --- с Базарной улицы на Стриженого Кактуса, со Стриженого Кактуса между домов на север, до самой стены, затем на запад вдоль стены и обратно к храму.
Дождь застал нас на окраине города, возле реки.
Небо вдруг резко приобрело тусклый оранжевый с синевой оттенок, на город пали ливневые сумерки.
Чханэ весело морщилась от стекающих по её лицу холодных капель, наши рубахи в первую же минуту промокли до последней ниточки.
Мы, пытаясь перекричать шум воды, делились впечатлениями.
Наконец, когда поток мутной воды на дороге уже грозил сбить с ног, я схватил девушку за руку и бросился к ближайшему жилищу.

--- Попросим погреться, --- прокричал я ей в ухо.

Но жилище оказалось заброшенным.
Наше веселье приутихло.
Я осторожно закрыл провисшую дверь, и шум дождя остался где-то далеко.
Мы с Чханэ какое-то время смотрели на покрытые мхом каменные блоки, на сухие, похожие на старческие руки корни соседнего дерева, на бледные, болеющие побеги растений, пробившиеся из полусгнившего пола.
В углу, глядя на нас неподвижным взглядом, свернулась кольцом безобидная верёвочная змея, где-то на крыше слабо жужжали большие полосатые осы --- рептилия и насекомые терпеливо пережидали ливень в безопасном месте.

Чханэ осторожно прошла в соседнюю комнату и погладила рукой разбитый глиняный очаг.

--- Надо же, треснул.
Наверное, полка сорвалась и разбила.
Интересно, кто здесь жил?

--- Я не знаю, --- пожал я плечами.
--- Сколько мы здесь проходили, я всегда думал, что жилище обитаемо.

--- Я тоже.
Смотри, даже плёнка целая, как будто его покинули совсем недавно, --- Чханэ кивнула на окна, затянутые мутной плёнкой бычьего пузыря.

--- Может быть, стоит спросить у соседей?

--- Для начала нужно переждать дождь и согреться, --- усмехнулась Чханэ.
--- Что тут у нас?
Дровишки! --- девушка, сияя улыбкой, выгребла из угла ворох сухих щепок и несколько поленьев.

Вскоре в разбитом очаге запылал животворный огонь.
Змея, тихо шипя, поспешила уползти от дыма в тёмный угол, осы тоже заволновались.
Но Чханэ открыла вытяжку --- и дым, извиваясь седыми кольцами, устремился наружу.
Шум дождя придвинулся ближе.
Мы с Чханэ, раздевшись донага и завернувшись в старое одеяло, сидели у очага и жевали дорожные пайки.

--- Жнаешь, --- заметила Чханэ с набитым ртом, --- Ликхэ мне не ошобо нжавитша, но готовит она вкушно.
Шитриш немного перешаливает.

--- Ммм, --- согласился я с подругой.

За окном совсем стемнело.
Мы сидели под одеялом, обнявшись, и изредка смотрели друг на друга.
У меня было ощущение нереальности происходящего.
Словно мы вдруг попали в древнюю легенду, которые кормильцы рассказывали детям, которые иногда приносили из далёких мест путешественники и где-нибудь на постоялом дворе, под шум дождя передавали жадно внимающим слушателям.
Волосы Чханэ пахли цветами, зёрнами какао, дорожной пылью, чем-то домашним, уютным...
Я видел каждый изгиб шрамов, каждую родинку, каждый волосок, каждую морщинку, каждую саккаду.
Мной овладело чувство --- мощное, как свет маяка, пробивающий навылет туманы и морские дали.

<<Это та, с которой я хочу остаться навсегда>>.

--- Можно тебя поцеловать? --- спросил я.

Чханэ вместо ответа прикоснулась губами к моим губам, и я, закрыв глаза, ответил на поцелуй.
Поцелуй оказался безвкусным, но очень приятным.
Язык Чханэ щекотно прошёлся по моим дёснам, по моим зубам и нежно коснулся моего языка...
Я целовал её снова и снова, зарываясь пальцами в короткие жёсткие волосы.
Чханэ прильнула ко мне всем телом, её соски чиркнули по моей груди.

--- Иди ко мне, --- шепнула она, откинув одеяло.

Я лёг на девушку, и она обняла меня ногами.
Я прошёлся кончиками пальцев по нежной мякоти её груди, по упругому, сильному, покрытому шрамами животу.
Но вот я неловко двинулся --- и Чханэ захихикала, прошептав:

--- Давай ещё раз, не торопись.
Поднимайся ко мне, как корабль скользит по волнам...

Я послушался --- и почувствовал, как мой стебель входит в нежное, влажное, жаркое, как пламя очага, лоно.
Огненные глаза Чханэ смотрели прямо на меня, по её лицу бродила слабая улыбка, словно она тоже не верила в реальность происходящего.

Веко опускается и поднимается, словно крыло, чёрный, как ночь, зрачок сужается...

И ещё раз...

И ещё...

Очаг уютно трещал поленьями.
Дождь за стенами пел свою тихую монотонную песню.

\section{Кристалл}

--- Какая холодрыга, --- пожаловалась Митхэ, кутаясь в меховой плащ.
--- Что мы здесь делаем, Атрис?

Атрис промолчал --- вопрос был риторическим.

Кристалл встретил пришельцев неприветливо. Вернее, для Кристалла спокойная осенняя хмарь, без бурь и штормов, могла считаться гостеприимством, если бы не впечатление от
ярких фонарей порта Прибой, в котором только-только отгремел очередной дождливый фестиваль.

Торговля шла своим ходом.
На угрюмых смуглых лицах ноа появились широкие улыбки, а вот наёмники заскучали.
Зизоце --- неразговорчивый народ.
Многие из жителей восточного побережья Кристалла не знали даже цатрона;
их монотонная однообразная речь с трудом давалась сели, привыкшим к тонам и экспрессивным музыкальным пируэтам южных языков.
Больше всех сноровки проявил пожилой Аурвелий;
как оказалось, Кристалл для него был не в диковинку.

Следить за скучающими наёмниками --- неблагодарное занятие, а Митхэ вдруг до смерти захотелось отдохнуть по-настоящему.
Отдавшись какому-то порыву, она назначила Ситриса старшиной над отрядом.
Роптали все;
Эрхэ, помня о данном Митхэ обещании, ушла подальше от лагеря и громко выругалась в безымянное ущелье.
Но больше всех возмутился сам Ситрис.

--- Я плохо переношу ответственность, --- заявил он воительнице.
--- Да и кто станет меня слушаться?

--- Тебя будут слушаться все, --- рявкнула Митхэ.
Суровость вышла настолько неискренней, что наёмники удивлённо замолчали.
--- Я и так пошла у вас на поводу, согласившись на этот поход.
Думаю, самое время напомнить, что мой отряд --- не последнее братство воинов в землях сели.
Я никого не держу!

Митхэ и Атрис специально искали жилище на отшибе, и их приютил старик, живший неподалёку от деревни.
Митхэ не могла не признать, что добротная бревенчатая изба с печью во всю стену как нельзя лучше подходила к погоде.
Старик заварил им отвар по собственному рецепту;
в его доме все стены были завешаны пучками ароматных болотных трав.

--- Пить, бедняк.
Холод, холод, --- посочувствовал им старик на плохом цатроне, с отцовской теплотой гладя влюблённых по мокрым волосам.
По возрасту он сгодился бы Митхэ в младшие братья и даже хранители, но его борода и редкие волосы, обрамлявшие выбритое до зеркального блеска темя, были белее снега --- людей Кристалла старость настигала чересчур быстро.
Женщина поймала себя на мысли, что думает о нём, как о кормильце.

--- А что такое горчит? --- спросил Атрис у старика.
Тот недоумённо поднял белые брови.
\ml{$0$}
{--- Горько.}
{``Bitter.}
\ml{$0$}
{На язык попасть --- лицо сжать.}
{Touching my tongue---gathering my face.''}

\ml{$0$}
{Атрис наглядно продемонстрировал эффект неизвестного компонента.}
{\Aatris\ mimed the unknown component effect with his face.}
\ml{$0$}
{Хозяин оскалился и захохотал, задрав седую бороду.}
{The host showed his teeth, then raised his gray beard and made a loud laughter.}
Затем подошёл к окну и показал на белые с чёрными полосками деревья, которые окаймляли почерневшее поле, словно стражи.

\ml{$0$}
{--- Орешки.}
{``Nuts.}
\ml{$0$}
{Орешки, из который лист.}
{Nuts where leaf from.}
\ml{$0$}
{Я их варить, они лицо сжимать.}
{I brew they---they gather face.''}

--- Берёза, --- сказал Атрис Митхэ.
--- На языке тси это дерево называется берёзой.

--- Откуда ты знаешь? --- поразилась Митхэ.

--- Понятия не имею, --- пожал плечами Атрис.
--- Кстати, от тебя тоже пахнет берёзовыми почками.

--- Неужели у меня такая горькая судьба? --- поморщилась воительница.
\ml{$0$}
{Старик продолжал ухмыляться, бормоча в бороду: <<Лицо сжимать, сжимать лицо>>.}
{The old man continued smiling and muttering into his beard: ``Gather face, gathering my face.''}
Борода была для сели в диковинку;
волосы на губах и подбородке росли у мужчин хака, ркхве-хор, да ещё сотронские сели часто вживляли себе в подбородок кусочки чужих скальпов.
Однако ещё большей диковинкой была улыбка на лице человека Кристалла.
И мужчины, и женщины выглядели измождёнными --- короткое лето они посвящали изнурительной обработке земли и починке жилищ, а зиму и дождливое межсезонье беспробудно пьянствовали --- чаще поодиночке, запершись дома.

Влюблённые отказались спать в жарко натопленном жилище, справедливо решив, что лучше умереть от холода, чем от удушья.
Старик постелил им в тёмной пристройке.

\section{Любовь и смерть}

Очаг погас.

--- Лис, а у тебя были до меня женщины? --- игриво спросила Чханэ.

--- Хаяй, нет, --- смутился я.
--- Ты первая.

--- Странно.
Я думала, что были.
Очень уж ты уверенно держишься.

--- Кормилица мне рассказала, что нравится женщинам, --- пожал я плечами.
--- А у тебя были мужчины?

Чханэ неожиданно тоже смутилась.

--- Были.

--- Расскажи.

--- Если ты просишь... --- неопределённо выразилась Чханэ.

--- Это печальная история?

--- Сложно сказать...

Оказалось, что у Чханэ был всего один любимый человек.
Его звали Манис ар’Ликх.
Высокий, красивый купец, переехавший откуда-то с севера помогать её кормилице, на десять дождей старше.
Чханэ собиралась с ним жить.

--- Его терпению можно было позавидовать.
Манис носил цветы из джунглей, срывал их в расщелинах на скале Летучих мышей.
До сих пор, проходя мимо скалы, удивляюсь, как он не разбился.
Он нравился мне, но я чего-то боялась, а чего --- сама уже не помню.
Потом, после полугода ухаживаний, он просто во время прогулки взял меня под этой скалой и тут же предложил жить у него.

Совместная жизнь обернулась для Чханэ неожиданностью --- прошло три дождя, а детей у них не было.

--- Кормилец волновался за нас.
Цветочек иногда тоже задумывался.
А мне хоть бы хны --- ну нет и нет.
Маленькая ещё была.
Я занималась своим делом --- патрулировала, проводила разведку, тренировала детей, а Цветочек сидел дома и пересчитывал товары, распределял приправы и перенаправлял белую шерсть.
По вечерам мы готовили разные вкусности, по ночам занимались любовью.
Простая счастливая жизнь.

Гарнизонный городок всегда жил в ожидании набега диких пылероев.
Но те несколько дождей выдались спокойными, люди расслабились.
Расплата за отдых пришла внезапно.

--- Кровь до сих пор стучит в висках, как вспомню про то побоище.
Время было обеденное, стража скучала.
Пылерои подошли к частоколу под песком.
Стражники легли сразу.
Предупредил нас ребёнок дождей восьми, сидевший на дереве, он завизжал на боевом языке: <<Асаа! Асаа! Хитр! Хитр!>>\FM\
\FA{<<Пылерои! Пылерои! Строй ,,серп``! Строй ,,серп``!>> (боевые сигналы Десяти Песчинок, использовавшиеся в Пыльном Пригорье).}
Это были его последние слова.

Чханэ запнулась.

--- Взрослый крик из детских уст --- это самое страшное, что я слышала в жизни.
Вряд ли он даже понимал смысл этого сигнала.
Может быть, это кричали его кормильцы, и фраза отложилась у него в памяти.

Что было дальше, Чханэ, по её словам, плохо помнила.
Оборона развернулась, как сломанный зонтик --- дело дошло до боя в кольце.
По стене клешней бил град полновесных ударов.
Девушка использовала любую возможность, чтобы рубануть шерстяное клыкастое месиво.
Время от времени ловила взгляды друзей, выхватывала из орущего, рычащего хоровода смутно знакомые человеческие лица...

--- Я уже даже не командовала, знаешь.
Люблю покомандовать, а тут как отшибло.
Весь мой район стоит рядом, бьётся, всем так туго, что дышать тяжело, а я только так командую, ногами --- вперёд, назад, в стороны.
Потому что крикнешь что не подумав --- и лягут все...

--- ... клешня разлетелась в щепки, ну я ему и ткнула в глаз какой-то случайной фалангой.
Пока была брешь, заорала, чтобы кольцо стянули, да куда там --- двоих сразу положили луки...

Кто-то догадался броситься в самоубийственную атаку <<на клешнях>>.
Это спасло кольцо --- пылероев оттеснили на простреливаемое пространство и вынудили отступить.

--- А потом я увидела его.
Берсерк, огромный пёс, на голову выше остальных.
Он раскидывал людей своей палицей, как щенят.
Манис сражался с другим пылероем и не видел, что кольцо развалилось.
Берсерк схватил моего мужчину за руку и сломал её так, что фаланга Маниса вонзилась ему же в спину.

Чханэ помолчала, с отсутствующим видом играя пальцами.
Потом усмехнулась и запрокинула голову.

--- С этого момента я видела только его.
Эту обтянутую кожей огромную клыкастую морду, которая возвышалась над полем битвы, как осадная башня.
Я не помню, как в моей левой руке оказалась фаланга Маниса.
Помню только, что берсерк обернулся, и мы встретились взглядами.

Битва словно отодвинулась на второй план.
Они шли навстречу друг другу --- оскалившая зубы гигантская собака с усаженной обсидиановыми лезвиями палицей и волочащая две фаланги высокая девушка, едва передвигающая ноги, скривившаяся от боли в сломанных рёбрах.

--- Мы оба были значительно выше остальных. Битва остановилась. На нас смотрели
все, кто вокруг: люди, пылерои. Кормилец пробивался ко мне и кричал: <<Беги,
дитя>>. А я всё шла и шла.

--- Вы с ним были эр'марквал, --- сказал я.

--- Кто? --- не поняла Чханэ.

--- <<Главные танцоры>>.
Иногда, крайне редко исход битвы решает поединок между двумя бойцами, причём совершенно неясно, почему именно их выбирает случай.
Но когда они сходятся, сражение останавливается.
Мне говорили в Храме.

Чханэ подняла рубаху, показав уже знакомые мне шрамы на животе.

--- Знаешь, что это такое?

--- Кхаагатр, --- кивнул я.

Чханэ усмехнулась:

--- Хорошо, что вас хоть чему-то учат.

--- Как ты выжила? --- прошептал я.
Воображение отказалось мне подчиняться.

--- Кормилец.
Он носил с собой трофейную духовую трубку.
Она была особо ценным оружием какого-то знатного идола.
Акхсар знал, что произошло --- тело Маниса лежало неподалёку.
Он понял, что я не остановлюсь.
Он был далеко, шагах в пятнадцати, и всё, что мог --- послать стрелку наудачу.

Отчаянная стрела нашла свою цель --- уже замахнувшийся палицей берсерк получил в глаз едкую ядовитую занозу.
Дубина, миновав девушку, ударилась о камень, и в этот момент Чханэ с нечеловеческим рёвом обрушила фаланги, одну за другой, на прикрытую костяными щитками шею нагнувшегося пылероя.

--- Я сломала о него оба, и продолжала бить обломками.
Потом выдохлась.
Он стоял на коленях с разорванной мордой и ошарашенно смотрел на меня целым глазом, а я смотрела на него.
Кормилец кричал: <<Беги, дитя>>.
А я всё стояла и смотрела.

Пылерой опомнился и успел-таки схватиться за ритуальное оружие своего племени.
Кхаагатр свистнул в воздухе, но ослабевшая от парализатора и кровопотери рука собаки не сумела завершить удар.
<<Когти игры>> сломались, лишь оцарапав живот девушки.
Последний удар нанёс подбежавший Акхсар.

--- В тот день мы отбились малой кровью.
Благодаря тому ребёнку городок поднялся, едва пылерои добрались до Дороги Ветров.
Воины и крестьяне сдержали их до прихода луков и пращей с окраин.
Кормилец потом сказал, что собаки испугались, увидев гибель прославленного воина, но я думаю, он просто пытался меня утешить.
Это было тактическое отступление.

Чханэ помолчала.

--- Того ребёнка торжественно похоронили под деревом.
На нём до сих пор виден охранный знак с именем --- Булькающий-в-Вине-Колокольчик.
Цветочек тоже умер, защищая город, но, кроме меня, некому было его оплакивать.

--- Ты всё ещё любишь его, Чханэ?

--- Я всегда буду его любить.
Больше некому.

--- Он сейчас в пристанище лесных духов, --- заметил я.
--- Предки заботятся о нём...

--- Брехня, --- оборвала меня Чханэ.
--- За гранью смерти нет ничего.
Всё, что осталось в этом мире от Маниса --- здесь, --- она тронула пальцем висок.
--- А лесные духи --- всего лишь раскрашенное дерево на входе в деревню.

--- Безумного тоже нет?

Чханэ дёрнула щекой.

--- Если бы.

--- Почему же тогда ты не веришь в лесных духов? --- сказал я, проигнорировав богохульство.

--- Никто не защитил Маниса, когда пришли пылерои.
Никто не даровал мне детей, когда он просил.
Лесные духи --- лишь погремушка, которая весело звенит в колыбельке младенца, пока другого младенца приносят в жертву Безумному.

--- А в Безымянного ты тоже не веришь?

--- Что толку верить в того, кто был изгнан?
Мне нечего о нём вспомнить.

Я промолчал.

--- Никто не будет добр к тебе просто так.
Добрые боги отличаются от жестоких лишь тем, что требуют зерно, а не кровь детей.

Я лёг ей на живот, задрал рубаху и легонько стал поглаживать его.
Она смешно дёрнула пупком и почесалась.

--- Щекотно.

--- Ты сильная, --- сказал я.

--- Сильная? --- Чханэ рассмеялась горьким смехом.
--- Сожми крысу в кулаке, и она проест твою ладонь насквозь.

--- Наверное, за тобой многие стали ухаживать.

--- Я бесплодна, --- ровным голосом сказала Чханэ.
--- Кто-то из подруг постарался, растрезвонил, что у нас с Манисом не просто так не было детей.

--- Почему ты, а не Манис?

--- У него были дети от какой-то крестьянки.

Я поморщился:

--- Тебя не хотели, потому что ты бесплодна?
Слабо верится.

--- Я пристрастилась к конопляному вину, --- объяснила Чханэ.
--- Не то чтобы это большая редкость у нас, но я пила без просыпу, иногда целыми днями.
Поначалу --- чтобы утихли боль в рёбрах и скорбь по Манису, потом просто так.
Один пытался меня взять пьяную --- едва орехи успел унести.
Больше желающих не нашлось.
Как у нас говорят, нет хуже союза, чем союз крепкого кулака и тумана в голове.

Наступило молчание. Я не мог подобрать слова, чтобы продолжить беседу.
Но Чханэ неожиданно весело добавила:

--- Ты не волнуйся, больше я не пью.
Это другая история, и она куда веселее.

--- Ну так расскажи, --- попросил я и устроился поудобнее у неё на животе.

--- Всё тебе расскажи, Ликхмас-тари, --- улыбнулась Чханэ и погладила меня по голове.
--- Ладно.
Дело было где-то дождь спустя.
Я после обхода напилась конопляного и разбавила всё это пряным...

--- Пряным вином\FM?
\FA{
Пряное (жгучее) вино --- алкогольный напиток, усиливающий половое влечение.
}
На твою долю выпало мало приключений?

--- Знаю, это было лишнее.
Хотела вспомнить, как мы с Манисом его пили.
Так вот, пришла я домой.
Была ночь, кормилица уже спала, её женщина тоже дремала на циновке в зале.
Слуга заворчал на меня --- припёрлась, мол, среди ночи пьяная.
Дать бы ему по башке, да голова кружилась сильнее обычного.
Масло в животе кипит, штаны липкие, хоть выжимай.
В общем, драться не хотелось совсем, хотелось с кем-нибудь пообщаться, поговорить по душам, но не со слугой же.
А тут брат сидит один в углу.

--- У тебя есть брат?
Ты не рассказывала, --- удивился я.

--- У меня их четверо, только все уже взрослые, разлетелись по миру.
А Бумажка на три дождя младше меня.
Кормилица не хотела, чтобы мы общались, однако он часто приходил в храм поболтать.
Но ты слушай.
Сидит и грустит.
Я к нему подсаживаюсь и спрашиваю: что грустишь, родной?
А он обычно такой разговорчивый, а тут сидит, отмалчивается.
Ну, слово за слово, я и выяснила, что ему нравится девушка, думает, что делать.
А я ж храбрая и бойкая с вина-то.
Говорю, война ещё не началась, не пыли зря.
Собираюсь и иду среди ночи к той девушке.
Дом заперт, ставни закрыты.
Я ж, не будь дура, открываю ставню шипом гарды и залезаю прямо на лежанку к той девице.
Хорошо, что не к её кормильцам, они спали у другого окна.
Зажимаю ей рот, чтобы не орала.
Ну, она маленько прочухалась, узнала меня.
Я её как можно тише веду в кладовку и сразу в лоб --- ты чего моего брата обижаешь?

--- Она, наверное, удивилась.

--- А ты как думаешь?
У девицы глаза были, как котелки.
Залезает к тебе среди ночи девка в два цана\FM\ ростом, тащит куда-то, спрашивает невесть что, да и дух от меня ещё тот --- конопля, пот, сивуха, кровь чья-то, не помню уже.
\FA{
Цан (<<обрубок, черенок>>) --- мера длины, примерно равная 1,3 м.
Высокая Чханэ <<слегка>> преувеличила свой рост.
}
Кричать боится --- знает, каков мой клинок в бою.
И тут пряное ударило мне в голову, да и она такая красивая, мягкая, пахнет парным молоком со сна, одуреть можно с этого запаха... соблазнила я её.
Прямо в кладовке.

--- Серьёзно? --- засмеялся я.

--- Ага.
Она даже не сопротивлялась.
Ласковая такая, как котёнок.
Выбираюсь потом из её дома, опять же через окно.
Счастливая как амбарная мышь, на алтарь бы сама пошла, если бы Безумные позвали.
По пути чуть не получила стрелу от патруля, который принял меня за пылеройку.
Как признал, обругал по-всякому и отпустил, не рассказал никому в Храме.
Я всё же на хорошей славе была, хоть и пила.
Прихожу домой.
Брат уже спит.
Я заваливаюсь к нему в комнату, говорю, что попробовала, мол, эту вертихвостку, выдержка чувствуется, букет восхитительный.
Он молчит, глаза как котелки.
Я ему: <<Ах, ты мне не веришь?
Я тебе сейчас покажу, какова она...>>
Пряное вино... в общем, той ночью я сделала брата мужчиной.

--- Хаяй...

--- Утром проснулась в кладовке, завёрнутая в мешок, в обнимку с тюком хлопка.
С тех пор не пью и любовных приключений у меня больше не было.

--- А что брат?

--- Пока я спала, он пошёл к той девушке.
Сейчас они уже дом выстроили, живут вместе.
У них я всегда желанный гость, не то что у кормилицы.
Меня называют между собой Хри-тари\FM.
\FA{
<<Юный Хри>>.
Имеется в виду лесной дух Хри-соблазнитель, покровитель плотской любви.
}

Я захохотал, уткнувшись носом в тёплый живот подруги.
Чханэ почесала меня за ухом и ласково улыбнулась:

--- Ну вот и всё, Лис.
Правда, я не уверена насчёт того тюка, может, с ним ничего и не было...

\section{Чего хотела душа}

\epigraph
{Последователи буддхизма утверждают, что человек находится в постоянном круговороте жизни и смерти, и лишь следование определённым правилам позволяет выйти из этого круговорота.
Это самое прекрасное представление о жизни, которое я только слышал, и тот, кто хочет покинуть этот круговорот, видимо, крайне несчастен.}
{Мартин Охсенкнехт}

--- Ну поговори же со мной, --- Митхэ довольно сильно ткнула Атриса под рёбра.
Тот дёрнулся, но не издал ни звука.
Сегодня он был молчаливее, чем обычно.

--- Милый, почему ты такой холодный и неразговорчивый?
Ты уже вторую ночь игнорируешь меня как женщину.
\ml{$0$}
{Неужели эта таёжная глушь охладила твой пыл?}
{I can't believe that taiga wilderness has cooled your fire.''}

Мужчина помолчал, потёр холодными тонкими пальцами щёки подруги.

\ml{$0$}
{--- Не знаю.}
{``I don't know why.}
\ml{$0$}
{Ты мне очень нравишься, но в моих мыслях много и другого.}
{I like you so much, but I have many other on my mind.}
Я часто отвлекаюсь.
Не обижайся.
Иди обниму.

--- А я всё думала, когда до тебя дойдёт, --- обрадовалась Митхэ и обмотала Атриса меховым плащом с головой.
Ей было ужасно интересно, о чём же таком думает Атрис, но она уже перестала спрашивать --- выжать из менестреля вразумительное описание его мыслей было невозможно.

--- Мы как близнецы в утробе, --- заметил Атрис.

--- Это могло прийти в голову только тебе, --- усмехнулась Митхэ куда-то в тёплую шею.
--- Что бы я без тебя делала...
Я многих хоронила на обочине своего пути и свыклась со всеми смертями.
С твоей же, боюсь, свыкнуться я не смогу...

Атрис выпутался из плаща и привстал на локте, вглядываясь в лицо подруги.
Митхэ удивлённо замолчала и сделала робкую попытку снова завернуться в плащ.

--- Милый, холодно вообще-то...

--- Митхэ, пообещай мне одну вещь.
Ты будешь жить, даже если я умру.

--- Но...

--- Я не такой, как все, --- голос Атриса внезапно приобрёл величественные обертоны.
--- Если я умру, то смогу вернуться.
Для тебя же обратной дороги не будет.

--- Хорошо, --- Митхэ немного испугалась такой перемены в любимом человеке.
--- Хорошо, Атрис.
Только как ты вернёшься?

Атрис улыбнулся и завернулся в плащ, снова замкнув кольцо уюта.
Из его уст вырвалось странное щёлкающее шипение.

--- Что это?

--- Язык Красных травников.
Дословно --- <<новоселье>>.
Ты не знаешь ничего о религии травников?

Митхэ изумлённо посмотрела на Атриса, а потом захохотала.

--- А я-то думаю, знакомое какое-то скрежетание.
Я когда ножик точу, он со мной так же разговаривает.
Вот делать мне нечего, только болтать о религии с богомолами-переростками.

--- Зря ты так.
С Кусачкой, который ходит с нами чуть ли не всю жизнь, никто никогда не разговаривает.
А ведь он очень образованный и интересный собеседник.

Митхэ вдруг стало стыдно.

--- С ним тяжело общаться, --- призналась она.
--- Все эти его <<живите долго>>, <<камни говорящие>>, а ещё он постоянно бормочет под нос какие-то числа, причём не к месту.

--- Числа, кстати, это элемент вежливого общения.
Числа имеют у травников сакральный смысл, могут использоваться как благословение и как проклятие.
Даже при торговле.
И не приведи Безумный предложить травнику неправильное количество товара.
Бродячий Народ, понятное дело, привык, а вот дикие горцы могут крепко обидеться.
Десять, тридцать пять, восток, север.

--- Ты говорил с Кусачкой о его религии? --- Митхэ с интересом взглянула на Атриса.

--- Среди прочего они верят в переселение душ.
Каждая душа проживает жизнь от рождения до самой смерти, а потом находит себе новое тело.
И всё начинается сначала.

--- А как душа попадает в тело?

--- Она сама выбирает его.

--- Чушь какая-то, --- буркнула воительница.
--- Кто управляет этим?
Да и какой в этом смысл --- проживать жизнь за жизнью?

--- Этим управляет существо, которое они называют Старой Личинкой.
У души есть какая-то цель.
Старая Личинка заключает с душой договор --- она съедает старое тело, переваривает его и кусает жилу джунглей, чтобы освобождённая душа пришла к новому телу, соответствующему новой цели.

--- Эту чепуху я уже слышала где-то в Сотроне.
Там ходит байка про Червя-узурпатора, и в ней не больше здравого смысла.
Какая, скажи мне, цель у ребёнка, который умрёт годовалым от лихорадки?

--- Цели бывают разные, Митхэ, большие и маленькие.
У кого-то цель --- спасти от смерти юношу, который станет Королём-жрецом, у кого-то --- написать величайшую поваренную книгу всех времён, кто-то просто хочет попробовать пирог из черноягодника.
А ребёнок, даже ещё не родившись, может круто изменить историю.

--- А что случается с тем, кто эту цель достиг? --- Митхэ спросила так, словно уже знала ответ.

--- Он умирает не позднее следующего заката, --- подтвердил её подозрения Атрис.
--- Поэтому Красные травники не грустят, когда кто-то умирает.
\ml{$0$}
{Они говорят: <<Ты Это сделал.}
{They used to say: `Thou hast done This.}
\ml{$0$}
{Возвращайся, и ты снова Это сделаешь>>.}
{Thou shalt come back, thou shalt have done This again.' ''}

--- \emph{Это}?

--- О загадочной цели травники говорят исключительно так --- Это.
Кстати, по этой же причине они очень настороженно относятся к долгожителям.
Чаще всего те, кто достиг преклонного возраста, сами уходят из поселения и ищут своё предназначение в путешествии.
Многих таких путешественников сложно отличить от последователей Сата-скитальца --- те же фонари и те же посохи.

Митхэ помолчала.

--- А что бы ты загадал на следующую жизнь?

--- Хай, --- задумался Атрис.
--- Я бы хотел прожить с тобой сто дождей.

--- Всего сто? --- пробурчала Митхэ.
--- А я хотела по меньшей мере сто тысяч...

--- Да ты что, Митхэ?
Люди столько не живут.
Да и неинтересно как-то...

--- Хаяй, тебе со мной неинтересно? --- обиделась воительница.
--- Нос разбить?

--- Да ну тебя.

В палатке воцарилось молчание.
Митхэ рассеянно тёрлась носом о шею Атриса.

--- Пойдём погуляем по лесу, --- предложил мужчина.
--- Сон не идёт.

--- Ага, --- саркастически крякнула Митхэ.
--- Стрела идола в ягодице замечательно подчеркнёт мою фигуру.

--- Здесь нет идолов, --- напомнил ей Атрис.
--- Мы на другом материке.

--- Лесные духи, я и забыла, --- Митхэ хлопнула себя по лбу.
--- Ну есть же другие опасности, которые будут хорошей отговоркой?

--- Нет, --- хитро улыбнулся Атрис.
--- Проводник сказал, что люди здесь мирные, а единственные звери, которых надо бояться --- это огромные коричневые меховые шарики.
\ml{$0$}
{Медолюбы, медоведы, медоеды, как-то так.}
{Borns, bears, brings, something like that.}
Сейчас они все сытые, на человека нападать не станут.
И не нужно сарказма, эту странную привычку Ситриса и так перенял весь отряд чести.

--- Ягуары здесь есть?

--- Вряд ли, в таком редком лесу ягуара видно за кхене.
Кроме того, я их не слышал.

--- А волки?
Здесь же должны быть какие-нибудь северные волки?
И ещё мне надоело вытаскивать из подмышек клещей.

--- Если мне не изменяет память, их вытаскивал я.

--- Волков?

--- Митхэ, хватит уже.
Вылезай и пошли.

Ворча, женщина выбралась из-под меха и, зябко поёживаясь, начала поправлять ремни на доспехах.

--- Здесь любая одежда кажется бумагой, --- бурчала она.
--- Насквозь продувает...

--- Ляг, --- велел ей Атрис, --- я завяжу тебе плащ по-местному.
Да, вот так, только перевернись.

Митхэ послушно легла на живот.
Атрис замысловато сложил плащ у неё на спине, пропустил под животом подруги ремень и затянул его.

--- Обереги тебя духи, --- поблагодарила Митхэ, почувствовав, что стало теплее и проклятый мех уже не так раздражает её.

\section{Кролик и койот}

--- Неужели уже утро? --- пробормотал я с закрытыми глазами, вслушиваясь в непрекращающийся плеск воды.
По улицам, похоже, неслись целые реки.

--- Уже день, --- промурлыкала Чханэ и зарылась глубже в одеяло.
--- Успокойся, сегодня мы можем отдыхать, первое число.
Я понежусь ещё немного и разожгу очаг --- у меня остался целый паёк.

Я поцеловал подругу в губы раз, второй, и <<ещё немного>> растянулось до вечера.
Мои поцелуи, словно следы путника, привели к женским воротам.
Чханэ всхлипнула.

--- Трай-трай-трай\FM...
\FA{
Что-то вроде <<тихо-тихо>>, <<щекотно>>, <<колется>>.
}
Тебе что, нравится?

\ml{$0$}
{--- Похоже на сливки, --- поделился я впечатлениями.}
{``Tastes like cream,'' I shared my impressions.}
\ml{$0$}
{--- Сливки с сахаром.}
{``Cream with sugar.}
\ml{$0$}
{И солью...}
{A little salty ...''}

--- Дурак... --- томно выдохнула подруга.
Она попыталась отвесить мне затрещину, но вышла только слабая ласка.

Когда сумерки дошли до густоты хлебного теста, мы всё-таки разожгли очаг и поужинали.

К ночи ливень превратился в мелкий моросящий дождь;
нужно было собираться в храм.
Предложение переночевать в доме кормилицы Чханэ категорически отвергла.

--- Что ж, --- усмехнулся я, поправляя в последний раз робу, --- жилище мы нашли хорошее, верно?

--- Трааа...
Что?

\ml{$0$}
{--- Жилище.}
{``A house.}
\ml{$0$}
{Для нас с тобой, --- пояснил я.}
{For you and me,'' I explained.}

Чханэ смотрела на меня со смесью восхищения и страха.

\ml{$0$}
{--- Лис, но мы не можем здесь жить.}
{``Fox, we mustn't live here.''}

\ml{$0$}
{--- Почему?}
{``Why not?''}

\ml{$0$}
{--- Что значит <<почему>>?}
{``What does it mean?}
\ml{$0$}
{Ты жрец, я воин.}
{You're a priest, I'm a warrior.}
\ml{$0$}
{Наш дом --- это храм, мы не должны быть хозяевами жилища, помнишь?}
{The temple is our home, we mustn't own a house, do you remember that?''}

\ml{$0$}
{--- Какая разница, --- отмахнулся я.}
{``Who cares,'' I waved my hand.}
\ml{$0$}
{--- Здесь никто не живёт.}
{``Nobody lives here.''}

\ml{$0$}
{--- Большая, Лис! --- Чханэ разволновалась не на шутку.}
{``I care, Fox!'' \Chhanei\ became pretty nervous.}
\ml{$0$}
{--- У нас нет разрешения.}
{``We have no permission.}
Жилище чужое.
Это пахнет Насилием.

--- Ты уже кутрап, тебе-то что, --- шёпотом пошутил я, бросив взгляд на окна.
\ml{$0$}
{--- И хватит о законах.}
{``And stop talking about law.}
\ml{$0$}
{Мы можем здесь поселиться.}
{We have the right to live here.''}

\ml{$0$}
{--- И по какому праву? --- скептически спросила Чханэ.}
{``What right?'' \Chhanei\ asked skeptically.}

\ml{$0$}
{--- По праву кролика и койота, --- ответил я.}
{``Right of rabbit and coyote,'' I answered.}
\ml{$0$}
{--- Если нора койота пуста, кролик имеет полное право в ней жить.}
{``While coyote hole is empty, a rabbit has the right to live there.}
\ml{$0$}
{Это право, я скажу тебе, древнее всех законов сели.}
{I say unto you: this right is older than \Seli\ law.''}

\section{Рябина}

Вскоре влюблённые уже шли под водяной пылью моросящего дождя, оглядывая неприветливые облетающие деревья.
Под ногами хлюпали жёлтые пятнистые листья, изредка попадались большие сопливые грибы.
Атрис то и дело смеялся над их формой:

--- Смотри, Митхэ, этот похож на ухо.
Смотри, а этот похож на...

--- Да, Атрис, я поняла, на что он похож.

--- Тебя это не забавляет?

--- Я хочу мужчину и есть.
Еда в форме мужского стержня вызывает у меня смешанные чувства.

--- Я не знаю, можно ли их есть.
Проводник сказал, что бывают ядовитые.

--- Пахнут аппетитно.

Вскоре парочка вышла на небольшую поляну.
Митхэ, сделав шаг в мокрую траву, замерла как вкопанная, наклонилась и провела рукой по маленькому золотистому колоску.
Вылущив зёрна, воительница задумчиво отправила их в рот.

--- Надо же, рожь.
Культурная.
Здесь кто-то жил?

--- Вполне возможно, --- ответил Атрис и крякнул --- сапог с чавканьем ушёл в незаметную мочажину.
--- Местные племена часто меняют поля.
Земля очень бедная.

--- Почему они не удобряют её?

--- Ты будешь смеяться, но это против их верований.
Проводник от моего вопроса пришёл в священный трепет и сказал, что mama нельзя трогать, она может излечиться только сама.
Забавно, да?
Они считают землю mama.

--- Значение слова в точности как у хака --- женщина, выносившая и воспитавшая дитя?

--- Да.
Я думаю, это очень древнее слово.

--- Не так уж они и не правы.
Из земли растут травы и деревья, плодами питаемся все мы.
Разве нет?

--- Почва рождается из мёртвых листьев и травы.
Кто кому ещё даритель.
Смотри лучше сюда.
Ты видела такие?

Атрис заворожённо разглядывал странные, ни на что не похожие деревья.
Митхэ подошла поближе... и тоже открыла рот.
Небольшие, с блестящей, словно медной корой, благородной формы ветви.
Собранные из множества зубчатых монеток листья, местами жёлтые, местами нежно-салатовые, а местами багряные, как кровь.
И ягоды.
Тяжёлые гроздья оранжевых и красных твёрдых ягод.

Митхэ осторожно подошла к волшебному деревцу и прикоснулась рукой к коре.
Странно.
Кора местных деревьев была склизкой, насквозь пропитанной сыростью, как старая мочалка.
Медный панцирь деревца отталкивал влагу, оставаясь солнечным и сухим.

Атрис осторожно оторвал ягодку и сунул в рот.
Поморщился и выплюнул.

--- Не ядовитая, но жёсткая и горькая, --- сообщил он.
--- Или ещё не созрела, или это есть нельзя.

--- Я впервые такие вижу, --- призналась Митхэ.
--- Даже в джунглях нет такой красоты, а тут, на севере!..

--- Давай посидим под этим деревцем, --- предложил Атрис.

--- Давай.

Атрис накрыл Митхэ своим плащом, и она благодарно улыбнулась ему.

Они молчали.
Дождь шептал траве что-то об ушедшем лете, о жизни и о двух маленьких людях, сидящих под чудесным, совершенно не северным деревцем.

--- Помнишь, там, в доме хозяина, ты спросила, что мы тут делаем? --- прервал молчание Атрис.

Митхэ кивнула.

--- А вдруг мы здесь, чтобы увидеть эти деревья?

--- Возможно.
Больше ничего красивого на Кристалле нет.
Северная глушь.

--- Я имел в виду другое.
Что, если мы родились, чтобы Это сделать?

Митхэ посмотрела на любимого человека, потом на раскинувшееся над их головами деревце... и ничего не ответила.

\section{Хризантемы}

\ml{$0$}
{В течение трёх дней мы с Чханэ обустроили новый дом.}
{In three next days \Chhanei\ and me have made the old house fit for living.}
\ml{$0$}
{У нас появились циновки, ковёр, посуда, инструменты и даже плетёное кресло из тростника.}
{We'd got grass mats, a carpet, pottery, tools, even a wicker cane chair.}
Очаг, к сожалению, починить так и не удалось;
я аккуратно поинтересовался у горшечников, как это сделать, получил смешок и лаконичный ответ: <<Закажи новый, жрец, и не валяй дурака>>.
%{``Order a new one, priest, don't fool yourself.''}
Поэтому я разобрал половые доски и выложил кирпичами защищённое кострище, прикрыв его обычной жаровней.
Готовили мы с Чханэ в основном в храме, дымовод у нашего жилища был общий с соседями\FM, а обитателей улицы Стриженого Кактуса не особо интересовали идущие по своим делам воин и жрец.
\FA{Мост-дымоход --- распространённый вид дымохода на кольцевых домах. }

В храме нашлись симпатичные старые лотки и горшки для растений.
Они лежали в куче хлама, и никто их не трогал, вероятно, с самого моего рождения.
Я хотел взять их, но Чханэ решительно воспротивилась.

--- Во-первых, у нас с тобой есть храмовые участки, --- заявила она.
\ml{$0$}
{--- Во-вторых, люди могут не обратить внимания на новую дверь, а вот цветущие хризантемы в горшочках расскажут нашу тайну всему Тхитрону.}
{``And second: people may fail to notice a new door, but chrysanths blooming in their pots will share our secret with the whole \Tchitron.''}

\ml{$0$}
{Мне пришлось признать её правоту, несмотря на то, что без хризантем образ дома казался незаконченным.}
{I had to admit she's right; with no chrysanths the image of home seemed incomplete to me, though.}

Устойчивые пары внутри Храма были редким явлением.
Воины и жрецы, разумеется, имели любовников в городе;
многие Храмы, особенно из молодых, представляли собой одну большую поликулу.
Тем не менее в Храмах не разрешалось воспитывать детей.
О нас с Чханэ знали все, но смотрели на этот союз с лёгкой иронией --- дескать, пусть балуется молодёжь.

Я ещё в первые декады в Храме заметил, что Кхохо и Ситрис держатся, как пара.
Я спросил об этом у Конфетки.

--- Они знают друг друга очень давно, --- объяснил тренер.
--- Я столько не прожил, сколько они друг друга знают.
Однажды Ситрис вернулся из джунглей со сломанной рукой и разбитым носом.
Оказалось, что Кхохо <<пошутила>> и срезала его лиану, когда они лезли за созревшими акхкатрас.
Без <<железа>> и тросов, на одних пальцах.
Её совершенно не волновало, что он может покалечиться или даже разбиться.
Когда я спросил Ситриса, не обидело ли его это, он довольно грубо ответил, чтобы я не лез не в своё дело.
Иногда они дерутся до крови, но потом лежат вместе и над чем-то смеются.
Когда люди знают друг друга столько времени, их связывает что-то недоступное нам, посторонним.

--- У них есть дети? --- с чисто детской наивностью спросил я.
О парах я судил по кормильцам: раз мужчина и женщина вместе --- должны быть дети.

--- Есть, --- ухмыляющийся Ситрис появился словно из-под земли с Кхохо в обнимку.
\ml{$0$}
{--- Ты.}
{``You.''}

Кхохо громко захохотала, показав красивые белые зубы.
Не будь шрама и татуировок, она бы отличалась удивительной красотой --- живые глаза золотисто-орехового цвета, тёмные веки, алые губы и гармоничное, благородное лицо с обветренными морщинками.
Её крепкое, слегка полное тело переливалось, словно ртуть.
Кхохо производила впечатление крестьянки-хохотушки, но позже я узнал, насколько обманчивым было это впечатление --- в бою она превращалась в берсерка.
Ситрис рассказал, что Кхохо как-то на острие клина прорубилась сквозь строй копейщиков, а потом долго бегала от врагов, потому что все союзники остались по другую сторону строя.
\ml{$0$}
{Когда воительницу нашли в кустах, на ней не было живого места.}
{When the warrior was found in the bushes, she had no piece of skin unbroken.}

--- У неё всегда были проблемы с работой в команде, --- резюмировал Ситрис.
--- Если уж быть честным, то и с головой тоже.

Круглосуточная для Чханэ шла полным ходом.
В общем и целом она мало отличалась от той, что устроил мне когда-то Конфетка.
Напасть на молодую воительницу могли и в уборной, и в душе, и за трапезой.

--- Мальчики, я сегодня на ужин не приду.
После обхода пойду играть к ткачам, у них большая заруба в <<Пьянку>>, --- сообщила Хитрам, жуя кукурузную лепёшку.
--- Кстати, а кто сегодня...

Ликхэ неожиданно бросила столовый нож в Чханэ, и рукоять ударила девушку точно в грудь.

--- И ещё один штрафной, --- радостно возвестила Ликхэ.

--- Да что ж такое, --- проворчала подруга.

Эрликх решил закончить урок и попытался дать Чханэ пощёчину.
Блок, удар локтем в грудь --- и воин отлетел в угол.

--- Прогресс определённо есть, Ликхэ, --- невозмутимо заметил Ситрис и показал забинтованный палец.
--- Сегодня ночью Чханэ великолепно среагировала на нападение.

--- Если честно, я просто испугалась, --- сообщила Чханэ и помогла слегка ошалевшему Эрликху сесть обратно за стол.
--- А первое, что подвернулось под руку --- столик.

--- Кстати, о том самом столике, --- сипло буркнул Эрликх, потирая грудь.
\ml{$0$}
{--- Краснодеревщики сказали, что собирать эту мозаику в третий раз они не будут и сделают новый, без возможности его схватить и использовать как оружие.}
{``Cabinetmakers told me: they will not piece together this mosaic the third time, they will make a new table---without the possibility to grab it and use it as weapon.''}

--- Да ты что! --- расстроилась Хитрам.
\ml{$0$}
{--- Ушла эпоха...}
{``The end of an era ...''}

--- Ситрисовская голова обладает разрушительным действием на мебель, --- пожаловалась Кхохо.
--- Куда его ни приложишь --- то вмятина, то трещина, то лак портится...
Из-за него такой беспорядок на этаже!

--- Да это мебель обладает разрушительным действием на мои пальцы! --- поправил Ситрис.
--- Но это ещё ничего.
Помню, я как-то во время круглосуточной бросил в тренера оцелотом.
Дело было ночью, а коту что-то потребовалось возле моей лежанки.
Тренеру пришлось накладывать швы на ухо.
Он сказал, что я буду великим воином, раз даже лесные звери становятся моими солдатами.

--- Полно врать, Ситрис, --- хмыкнула Кхохо.
--- Какой тренер?
Какая круглосуточная?

--- И у меня был тренер, Кхохо, не поверишь, --- парировал Ситрис, но все заметили, что он смутился.
Прошлое помощника вождя было тайной за семью печатями.
Он редко мылся и никогда не снимал перед сном рубаху.
Однажды я случайно увидел, почему --- спина Ситриса была сплошь испещрена следами от хлыста, которым погоняют скот.
Кхарас посоветовал мне не распускать по этому поводу язык.

В следующую декаду за обедом Чханэ превзошла саму себя: нож Ликхэ ещё не покинул ладони, когда из рукава подруги вылетел круглый тяжёлый камень.
Ликхэ, явно не ожидавшая от неповоротливого Скорпиона такого проворства, молча сидела и тёрла разбитую губу, пока остальные хохотали.
Кхарас признал, что круглосуточную для Чханэ пора заканчивать, но на всякий случай накинул пять дней --- для закрепления навыков.

--- Разумеется, это не значит, что через пять дней ты должна расслабиться, --- строго заметил боевой вождь.

--- Хай, мальчик, перестань, --- примирительно сказала Хитрам.
--- Девочка только отъедаться начала, щёчки круглее стали, а ты ей лишнее беспокойство внушаешь.

--- И вообще Конфетка говорил, что круглосуточную надо заканчивать, когда послушник расслабится, --- заметил Эрликх.
--- Потому что смысл тренировки в том и есть, чтобы человек стал чувствовать себя уверенно.

--- Конфетка? --- удивилась Чханэ.
\ml{$0$}
{--- Это кто?}
{``Who's that?''}

\ml{$0$}
{Воины тепло улыбнулись и промолчали.}
{Warriors smiled warmly and kept silent.}

\ml{$0$}
{--- Долгая история, --- ответил я.}
{``That's a long story,'' I answered.}

В остальном жизнь вернулась в прежнее русло.
Писем из Тхаммитра не приходило, и понемногу мы с Чханэ, жившие в лёгком напряжении, начали успокаиваться.
Может быть, Храм Тхаммитра счёл подругу и её мучителя пропавшими без вести.
Сложно ли потеряться в сельве, которая в любой тени может таить гибель, которая за считанные рассветы заплетает дороги и тела сетью жадных побегов?

\ml{$0$}
{<<В этом опасность и прелесть сельвы, --- говаривал старый охотник, Сиртху-лехэ.}
{``There is danger and appeal in the Silva,'' the old hunter, \Sirtchu-\lechoe, used to say.}
\ml{$0$}
{--- Там, где неумелый пропадёт, знающий скроется>>.}
{``Where an inept one is lost, a skilled one is sheltered.''}

\chapter*{Интерлюдия III. Ученица Повелителя ветров}

\addcontentsline{toc}{chapter}{Интерлюдия III. Ученица Повелителя ветров}

\textbf{Сказка сели.
Сборник сказок для детей из библиотеки Тхитрона}

Это было в незапамятные времена, когда дороги прокладывали по облакам, мосты строили через моря, а люди слыхом не слыхивали о требовавших крови жестоких богах.
То была эпоха духов, ведавших землёй, водой и девятью ветрами\FM.
\FA{
По восьми сторонам света плюс <<властелины пустошей>> --- торнадо.
}

В далёкой северной деревушке, в доме старого Кита родился беловолосый ребёнок.
Той ночью, когда родильница ощутила первые схватки, поднялся ветер.
Вековые деревья трепетали на этом ветру, словно распущенные косы, крыши многих жилищ перелётными птицами отправились на юг.
Едва младенец испустил первый крик, стихия утихла.

--- Неспроста это, --- толковали старики за чаем, --- не иначе как сам Повелитель ветров нанёс дому Кита визит.

Но жизнь шла своим чередом.
Вместо сломанных веток деревья вырастили новые, ещё краше, и сорванные крыши были для сильных мужей лишь поводом для ремонта.
Солнце всё так же садилось на западе и всходило на востоке, люди сеяли и жали кукурузу, в очагах пеклись лепёшки, а пёстрые курочки высиживали цыплят.
Младенец, которого назвали Веснушкой, стал девочкой, девочка превратилась в молодую женщину, и о давних стариковских пересудах никто не вспоминал.

--- Веснушка, ты сегодня хорошо поработала!
Иди кушать! --- этими словами кормилица обычно звала девушку домой.

В тот вечер Веснушка не откликнулась на зов.
Она лежала за кустом черноягоды, украшенным тяжёлыми обсидиановыми серьгами, и её белые волосы раскинулись по старому шерстяному одеялу.
Ветер нежно холодил обнажённое разгорячённое тело, девушка прижималась к любимому человеку и чувствовала себя самой счастливой на свете.
В деревне зажглись и потухли фонари, а влюблённые тихо беседовали и смеялись.
Гроздья ягод служили им пищей, а хрустальный родничок утолял их жажду.

--- Не случилось ли чего, лехэ? --- обеспокоенно спрашивала кормилица, заламывая руки.

--- Любовь случилась! --- смеялся старик Кит.
--- У старика Веха Ключик тоже пропал.
Не иначе скоро в деревне выстроят новое жилище.

Но вот беседа утихла, и любовников начало клонить ко сну.
Первым заснул Ключик, прижав подругу к своей широкой груди.
Звёзды шли своим кругом, в непролазных кустах стрекотали сверчки, и Веснушка погрузилась в сладкую дрёму.

Вдруг, откуда ни возьмись, в ночной темноте закружился вихрь.
Застонали деревья, притихли птицы, понеслись тревожными синими искрами индиго-светляки.
Испуганная девушка, прижавшись к любимому, смотрела, как к ней идёт древний седой старец, закутанный в рваный чёрный плащ.
Его очи были подобны звёздам, его белые волосы развевались, словно высокие ледяные облака, и такой же белой была его длинная борода.

--- Мир тебе, девушка, --- сказал старец.
--- Я --- Повелитель ветров, и ты была избрана мне в подмастерья ещё до рождения.
Пойдём, и я научу тебя премудростям проводника вихрей.

--- Мир тебе, Повелитель ветров, --- прошептала Веснушка, прижимаясь к Ключику.
--- Для меня честь говорить с великим духом, однако я не могу быть твоим подмастерьем.
У меня есть любимый человек, дом и очаг, и я обязана им.

Загорелись гневом глаза старца, затрепетали под яростным порывом облачные космы.

--- Я обрушу вечные ураганы на твою деревню.
Ни одно жилище, ни одна мельня не смогут стоять прямо, если ты не пойдёшь со мной, --- прогремел он, и трещали в его голосе ломающиеся деревья.

Что оставалось делать девушке?
Надела Веснушка штаны, подцепила рубаху на пояс, поцеловала в последний раз любимого и ушла с Повелителем ветров в ночь.

Догорела ли свеча, догорело ли солнце, но донесли облачные олени Повелителя ветров и Веснушку в поднебесные башни, на вершину высочайшей из гор.
Пусты и неприютны были их стены --- грубый камень ласкали лучи светил, бриз врывался в окна и гулял по огромным залам, но не было в залах ни мебели, ни иных вещей, что придают домашний уют жилищу.

--- Это твой дом, --- сказал ей Повелитель ветров.
--- У меня нет одежд и пищи, но здесь ты не будешь знать ни голода, ни холода, ни болезней.

Так началось обучение Веснушки.
Повелитель ветров рассказывал ей всё, что знал, но чересчур своеволен и гневлив был этот учитель.
Порой в залах царил штиль, иногда же, в плохие дни, молнии и ураганы плясали на вершине высочайшей из гор, ворочая огромные камни, словно куриные перья.

Девушка быстро постигала науку.
Шли годы, рубаха на её груди, спине и плечах истлевала, словно одеяние мертвеца, но пальцы и прекрасные белые волосы ткали ветра всё с большим мастерством.
Дело нравилось Веснушке, но нет-нет она да вздыхала, вспоминая по вечерам оставленный дом и любимого.
Смирились ли люди с её уходом? Простил ли её Ключик за то, что она бросила его посреди ночи?

Наконец пришла пора испытаний.
Запрягла Веснушка облачных оленей, и помчали её игривые звери по призрачным струнам.

Вот внизу показалась деревушка, страдающая от засухи.
Река обмелела, и кукурузные стебли горели, не успев дать початка.
Люди ходили, словно тени, и с тяжёлым чувством переговаривались о предстоящих голодных дождях.

Веснушка посмотрела на это, и печалью наполнилось её сердце.
Умелые пальцы девушки соткали ветер --- и на исходе дня пролился над горящими полями животворный дождь.

Гневными словами встретил ученицу Повелитель ветров:

--- Что ты наделала! --- бушевал он, и в его голосе его разбивался о скалы девятый вал.
--- Ты напоила одну деревню, но оставила целую страну без урожая!

Махнул Повелитель ветров рукой --- и Веснушка увидела страшное наводнение, постигшее далёкие земли.
Полей здесь не было --- всё скрылось под глинистой рыжей водой.

Опечалилась девушка, но продолжила свои занятия.
И снова пришла пора испытаний.
Запрягла Веснушка облачных оленей, и помчали её игривые звери по призрачным струнам.

Вот внизу показалась несметная тьма кораблей.
Великий завоеватель вёл их против маленького флота, защищавшего родной город.

Веснушка посмотрела на это, и печалью наполнилось её сердце.
Умелые пальцы девушки соткали ветер --- и запылали корабли захватчиков, и зазвенели паруса защитников, и победили они в той битве.
Захватчики отправились в свиту к Сестре Воде, а их предводитель был пленён и заклеймен, подобно последнему насильнику.

Гневными словами встретил ученицу Повелитель ветров:

--- Что ты наделала! --- бушевал он, и в его голосе его ревел пустынный самум.
--- Он мог объединить земли от рассвета до заката, принеся в раздираемые междоусобицами земли мир.
А теперь у живых не будет ни мира, ни урожая!

Махнул Повелитель ветров рукой --- и Веснушка увидела засуху и усобицы, поразившие земли защитников.
Кукуруза не взошла, и бывшие герои бродили по лесным тропам, перебиваясь разбоем и грабежами.

Опечалилась девушка, но продолжила свои занятия.
Больше она не запрягала облачных оленей --- она сидела в башне и читала книги.

Книги поведали Веснушке, что она не могла даже качнуть цветок орхидеи, не принеся несчастий.
Чем больше девушка читала, тем меньше ей хотелось применять своё искусство.

Догорела ли свеча, догорело ли солнце, но однажды ночью родилась мысль, подобная свету в первородной тьме.
Запрягла Веснушка облачных оленей, и помчали её игривые звери по призрачным струнам.

Но девушка более не ткала ветра.
Она летала над миром и смотрела на вечную беседу подвижного и неподвижного, земли, воды, солнца и воздуха.
У гор нити ветров всегда сплетались в кольца.
У рек и морей они напитывались влагой.
Леса расчёсывали ветер, словно гребни, а по степям и пустыням вихри гуляли по своим, ведомым лишь им законам.

Девушка начала ткать свой собственный холст.
Она дни и ночи напролёт сидела в башне, смотрела на карты, листала огромные пергаментные страницы, размышляла над словами давно мёртвых, забытых даже духами языков.
Повелитель ветров сердился на Веснушку: молнии плясали по окрестным горам, пыльные ураганы истачивали древние скалы.
Но девушка не обращала на них внимания.
И в один прекрасный день, пока Повелитель ветров дремал, Веснушка снова запрягла облачных оленей.

Она встала посреди пустыни и стала ждать.
Всё её существо трепетало от осознания того, что она собиралась сделать.
Прошёл день, прошёл другой, третий --- и на рассвете четвёртого дня Веснушка, затаив дыхание, махнула рукой.

Повелитель ветров не встретил свою ученицу.
На высочайшей из гор царил мёртвый штиль.
Девушка в страхе вихрем носилась по огромным залам, просторным коридорам и потайным лестницам, но они были пусты.
Наконец она обнаружила старца на маленьком балконе, с которого весь мир был виден, как на ладони.
Повелитель ветров сидел и дрожал, закутавшись в рваный плащ и совершенно потеряв свой грозный облик.

--- Что ты наделала, --- шептал старец, глядя в пустоту померкшими глазами,
--- что ты наделала...

--- Я развесила грузила на ткацком станке ветров, и теперь ветра будут двигаться по своим дорогам сами, --- сказала девушка.
--- Люди, наблюдая за облаками, смогут предсказать их путь и решать, когда время сеять, а когда время уходить в другие земли.
Больше я тебе не нужна.

Старец взвыл от бессильного гнева и затопал ногами.

\ml{$0$}
{--- Когда-то давно, --- закричал он наконец, --- камни и щепки предсказали мне, что ты, рождённая в ночь Согхо 8, будешь величайшей из Повелителей ветров.}
{``Ages ago,'' he shouted at last, ``stones and sticks predicted you, born at the night of \Sogcho\ 8, as the greatest of Windlords.}
\ml{$0$}
{Но ты осквернила наше ремесло.}
{But you have violated our handicraft.}
\ml{$0$}
{Ты обесчестила меня и всех моих предшественников.}
{You have defiled the honor of me and all my forebears.}
\ml{$0$}
{Уходи из моего замка и никогда не возвращайся.}
{Go away from my castle and never return.''}

Повелитель ветров махнул рукой --- и девушка оказалась возле родной деревни.

<<Сколько меня не было?>> --- со страхом размышляла Веснушка.
Ей казалось, что прошли века, а то и тысячелетия.
Она видела, как сменяются эпохи, как возвышаются и низвергаются города.
Но сонная деревня по-прежнему мирно стояла вдали, звёзды шли своим кругом, а в непролазных кустах пели сверчки.

Вдруг Веснушка увидела черноягодный куст.
Сомнений нет --- это тот самый, который дал приют ей и её давно покинутому любовнику.
Как и тогда, он был усеян тяжёлыми гроздьями.
Девушка ахнула и подбежала поближе.
Ключик тихо сопел носом, завернувшись в одеяло, словно она оставила его мгновение назад.

До нас дошло, что у Веснушки было много потомков.
Никто из них не унаследовал прекрасные белые волосы, но все были наделены удивительным даром понимать и предсказывать путь ветров.
Впрочем, было ли это даром?
Доподлинно не известно ни духам, ни бродящим от звезды к звезде фонарщикам, ни небесному пахарю.
Наследники Веснушки и в хмелю, и в объятиях прекраснейших женщин молчали с многозначительной улыбкой, едва речь заходила об этом.

Это было так, и пусть станут пищей ягуара те, кто рассказывает иначе!

\chapter{[-] Путь жреца}

\section{[-] Посвящение}

--- Ай!

Я наткнулся на полку глазом, пытаясь найти свежую рубаху.
В голове зазвенело.
Ругаясь вполголоса, я приоткрыл окно и прислонился ноющей глазницей к прохладной оконной раме.
В келью ворвался шум ночного дождя, крупные капли, сильный запах влажного камня и мокрой травы.

\ml{$0$}
{Старые должны передавать опыт молодым.}
{The old must share experience with the young.}
\ml{$0$}
{Именно поэтому я пришёл в Храм.}
{That's the reason I came to the Temple.}
\ml{$0$}
{Именно поэтому в Храм приходили тысячи и тысячи до меня --- испуганные, сомневающиеся, плачущие, ничего не знающие и ничего не умеющие.}
{That's the reason thousands before me came to the Temple---scared, confused, crying, with no knowledge, and no skill.}
\ml{$0$}
{Приходили, чтобы остаться, приходили, чтобы уйти, но всегда становились чем-то большим, лишь сделав шаг через порог древнего здания.}
{They came to stay, they came to leave, but they always became something bigger than they were, with one little step over the threshold of the ancient building.}

Учеников в Храме всегда хватало, но не в эти дожди.
Я был единственным послушником Верхнего Этажа.
Как-то я спросил, почему жрецы не обучают других.
Трукхвал сказал, что, когда я стану жрецом, он мне скажет.

И вот этот день пришёл.
Однако всё оказалось совсем не тем, что я ожидал.
В зале было темно и пусто, факелы уже давно догорели.

--- Ликхмас, что-то случилось? --- на лестнице показалась всклокоченная голова Кхатрима.
--- Я услышал твой голос.

--- Сегодня испытание, поэтому я встал пораньше, --- объяснил я.

--- Ты с ума сошёл, ещё даже намёка на рассвет нету.
Иди спать!

Недоумевая, я пошёл обратно в келью и завалился досыпать прямо в одежде.

Мне говорили, что послушник должен пройти испытание перед тем, как войти в Храм.
Когда меня всё-таки подняли на рассвете и повели в зал, я уже представлял собрание Совета, коварные вопросы и прочие испытания, и был немало разочарован, когда меня усадили за стол кушать.

Единственным коварным вопросом за время трапезы оказался вопрос насчёт синяка под глазом.
В конце же, перед традиционной чашей отвара, Первый жрец торжественно сказал, что выслушает мою клятву.

Слова клятвы я уже успел заучить наизусть.
Её цитировали все --- по поводу и без. Я прочистил горло и, встав из-за стола, произнёс:

\begin{quote}
<<Ни действием, ни бездействием я не наврежу книге, написанной не мной, даже если в ней не будет ни слова истины, и не позволю никому сделать подобного;
я не испорчу ни текста, ни художеств, и гулять моему перу лишь по её полям.\\~\\
Я не утаю знания ни из жалости, ни ради защиты, ни ради выгоды.
Я не оставлю без внимания ни ложь, ни заблуждение, и скажу во всеуслышание то, что считаю истиной.\\~\\
Я не потушу ни искры дарования, данного ребёнку, но разожгу их в пламя;
я расскажу ему то, что знаю и считаю истинным, и не оставлю без ответа ни один его вопрос в тот момент, когда он был задан.\\~\\
\ml{$0$}
{Я принесу жертву богам и в скорби, и в болезни, и при смерти;}
{Even grieving, sick, or dying, I shall sacrifice to the gods;}
\ml{$0$}
{я отдам людям эту ношу лишь тогда, когда не найдётся жреца, способного держать нож.}
{I shall shift this burden on to my people only when there will be no priest could hold a knife.}\\~\\
\ml{$0$}
{Я приложу все силы, чтобы облегчить боль;}
{I shall do anything to ease one's pain;}
\ml{$0$}
{я не наврежу человеку ни моими чувствами, ни моим невежеством;}
{I shan't harm one by my feelings or my ignorance;}
\ml{$0$}
{в жизни и смерти я буду щитом, и ни мгновения --- плетью>>.}
{dead or alive, I shall be a shield, and shan't be a lash.''}
\end{quote}

Жрецы заулыбались.
Мне показалось, что они с трудом сдерживают смех.
Кхатрим пододвинул мне чашу.

--- Добро пожаловать в Храм, Ликхмас.

--- Трааа... --- я на мгновение почувствовал себя обескураженным.
--- Это всё?
Я думал, вы просто хотели узнать, заучил ли я клятву.

--- И мы узнали, --- улыбнулся Первый.
--- теперь ты --- жрец.

--- Но я думал, что...

Мои слова прервал громкий хохот и звон посуды.
Жрецы больше сдерживаться не могли.

--- Ты думал о тяжких испытаниях, торжественных ритуалах и прочей ерунде? --- спросил Трукхвал.
--- Все так думали, не волнуйся.
Испытание ты давно прошёл.
А клятва, как ни крути, лишь слова.
Слова не имеют смысла, если нет того, кто может их понять.
Ты хорошо это усвоил, работая в библиотеке.

--- Все так думали, Лис, не ты первый, не ты последний, --- похлопал меня по плечу Кхатрим.
--- Это традиционная забава Верхнего этажа --- смотреть на озадаченную физиономию новоиспечённого жреца.
Вернее, на физиономию того, кто уже был жрецом как минимум три дождя, но ещё об этом не догадывался.

Я вдруг понял.

--- Значит, испытанием было...

--- Испытанием было всё, что произошло за время обучения, --- подтвердил Трукхвал.
--- Каждый твой шаг был испытанием.
И ты его выдержал.
А клятву мы просто так спросили, для формы.

Я сел за стол и взял чашу.
Жрецы всё ещё смотрели на меня и ухмылялись.

--- Что-то я не чувствую особой разницы до и после клятвы, --- признался я.
В этот момент Верхний этаж едва не понёс тяжёлые потери --- Кхатрим от хохота серьёзно подавился конфетой.
Впрочем, с этим непредвиденным испытанием я справился превосходно --- конфета была успешно вытащена и съедена под отвар.

\section{[-] Молодой учитель}

\textspace

--- Ликхмас, теперь ты жрец, --- торжественно начал Кхатрим, --- а это значит, что тебе следует пройти через самое суровое испытание в твоей жизни.

--- Но я же уже работал в госпитале, приносил жертвы, --- удивился я.
--- Да и клятву у меня уже спросили...

--- Не имеет значения, --- прервал меня жрец.
--- Через это испытание проходят все молодые обитатели Верхнего этажа.

Я промолчал.
Чутьё подсказывало мне, что врач не шутит, но...

--- Школа, Ликхмас, --- в священном ужасе прошептал Кхатрим.
--- Ты будешь учить детей в школе.

\ml{$0$}
{Я рассмеялся.}
{I laughed.}

\ml{$0$}
{--- Ты меня почти разыграл!}
{``You almost pranked me!''}

\ml{$0$}
{--- Ты не боишься? --- удивился жрец.}
{``Are you not afraid?'' the priest looked surprised.}
\ml{$0$}
{--- Это самое ужасное, что было в моей жизни...}
{``It was the most dreadful thing in my life ...}
В таком случае пойдём, я расскажу тебе, как готовить уроки и как выбирать себе ученика.
Жрецу нужен преемник.

--- Уже учеников?

--- А ты как думал?
Между учителем и учеником обязательно должны быть дружеские отношения!
Дружба не вырастает за день на голой земле!

--- А почему ты?
Где учитель?..

--- Митликх уехал в Кахрахан --- поклониться тотему Удивлённого Лю, отдохнуть и поговорить с кем-то из местного жречества.
\ml{$0$}
{Ещё он упоминал как-то, что там живёт его старинный любовник.}
{Also he mentioned his man, lover of yore, is living there.}
\ml{$0$}
{Ещё он сказал, что полностью тебе доверяет и знает, что ты справишься один.}
{Finally he said that you're completely trusted and competent enough to manage alone.''}

\ml{$0$}
{--- Трааа...}
{``\Traa\ ...''}

\ml{$0$}
{--- В общем, вряд ли Митликх вернётся в ближайшие восемь декад, --- смущённо закончил Кхатрим.}
{``Well, anyway, \Mitlikch\ is hardly going to return in the next eight decades,'' \Kchatrim\ got embarrassed.}

--- Главное, чтобы он не насовсем, --- заметил я.
Перспектива провести в школе остаток жизни почему-то вызвала во мне смутную тревогу.

--- Вот тут я с тобой полностью соглашусь, --- совершенно серьёзно ответил Кхатрим.
\ml{$0$}
{--- Я как-то тоже уехал из Сотрона на отдых.}
{``I left \Sotron\ for holidays once.}
\ml{$0$}
{Угадай, куда.}
{Guess where I went.''}


\section{[-] Любовь к языкам}

\epigraph
{Чем больше языков вы знаете --- тем больше у вас возможностей сказать глупость.}
{Длинный-Мокрый-Хвост}

\textspace

--- С сегодняшнего дня мы будем учить языки Трёх Материков, --- начал Хонхо-лехэ.
--- Прошу вас, запомните сказанное.
Мне доводилось слышать многие отговорки, почему человек не знает языков ближайших соседей.
Самой смешной отговоркой было отсутствие способностей к языкам.
Чушь!
Всё это полная чушь, сказанная человеком, который как минимум один язык знает в совершенстве!
Язык --- это инструмент, который вы всегда носите с собой и можете применить.
Язык --- это оружие, для которого не нужно ножен и стрел.
Знание языков --- это путь любого человека, который живёт в обществе, а не в одиночестве посреди голой пустыни.
Запомните это навсегда.
Вы можете презирать дикарей хака, вы можете ненавидеть идолов Живодёра, но их языки полюбите, как кормильцев и лучших друзей!

Слова старого жреца запомнились мне на всю жизнь, и я процитировал их детям.
Дети закивали в знак того, что поняли.

--- Хорошо.
Кто из вас знает язык цатрон?

В классе неуверенно поднялись две-три руки.

--- Спрошу по-другому, --- я перешёл на цатрон:
--- Сколько мер зерна требуется на среднее поле?

--- Одиннадцать, --- хором грянула половина класса.

--- Чьи корабли имеют раздвоенный штевень?

--- Ноа, --- грянула треть.

--- А когда птицы вили гнёзда среди скал в Сикх'амисаэкикх?

--- <<Пока дикарка ждала и плёлся паучий шёлк>>, --- раздался в полной тишине тоненький голосок.
Все недоумённо посмотрели худенькую девочку, и она стушевалась.

--- Эту песню кормилец напевает в поле каждую страду, --- объяснила она.

--- А говорите --- не знаете, --- засмеялся я.
--- Говорите только редко.
Это исправим.
Вы ещё и писать на нём будете.


\section{[-] Телесные наказания (отрывок)}

\textspace

Ещё совсем недавно детей могли побить за шалости.
Но около двадцати дождей назад в Ихслантхаре произошло нечто примечательное.
Наказание выпало девочке по имени Ситхэ, бойкой большеглазой заводиле.
Девочка, недолго думая, заявила учителю, что признаёт вину, но не примет наказание, так как считает его оскорбительным.
Учитель в шутку предложил ей устроить Дело Перекрёстка.
Ситхэ молча притащила из арсенала фалангу и щит, за неё вступились другие дети, и вскоре шутка вылилась в ритуальный бой между школьниками и воинами Храма.

--- Храм --- опора традиций, --- заявил боевой вождь.
--- Вы --- не взрослые, чтобы устраивать Перекрёсток.

--- Если нам отказывают в праве на Перекрёсток, поднять оружие --- наша священная обязанность, --- парировала Ситхэ строками закона.

\ml{$0$}
{--- Да будет так, --- ответил вождь.}
{``So be it,'' the warchief answered.}

Ситхэ с неожиданной сноровкой выстроила детей в боевой порядок и скомандовала атаку.

Домой дети вернулись в синяках, ушибах и с разбитыми носами, ещё больше уверенные в своей правоте.
Самой <<предводительнице>> потребовалась помощь жреца.
Осложнилось дело тем, что произошло всё на Змею 8, накануне праздника Обнимающего Сита.
Кормильцы, увидев своих чад в таком виде, созвали уже взрослый Перекрёсток (впрочем, вместо оружия все принесли мётлы, веники и букеты цветов), и традиция телесных наказаний была отменена.

Все были рады.
Дети радовались, что добились своего, кормильцы радовались за детей, но больше всех были счастливы, как ни странно, пожилые учителя, которым <<упражнения>> доставляли немало хлопот.
Многие из них по старинке грозили ученикам палками, но спустя десять дождей телесные наказания упразднили во всех городах Севера и Запада.

\section{[-] Школа}

\epigraph
{Стрелка компаса точна, но её концы раскрашивает мастер.
А уж нулевой меридиан и вовсе там, где вздумалось чихнуть королевскому географу!}
{Клаудиу Семито Фризский.
<<Воспоминания о великом тёзке>>}

\textspace

--- А что здесь, учитель? --- спросил я у Хонхо-лехэ, указав на южные части материков.

--- Хаммм, --- задумался учитель.
--- Ликхмас из дома Люм, я хотел бы ответить на твой вопрос, но не могу.
Эти земли называют Краем.
Известно, что они плохо пригодны для жизни --- болотистая местность чередуется с бесплодными полупустынями и горами.
Известно также, что туда очень сложно добраться.

Учитель достал указку и провёл её концом по узкой полоске воды.

--- Могильный пролив непроходим.
В будущем кто-то из вас может стать мореплавателем и захотеть испытать свои силы, но молю всех лесных духов, чтобы они удержали вас от этого самоубийства.
Толчки в проливе, по слухам, раскалывают мачту на две ровные половинки, отражённый водой солнечный свет выжигает глаза, а морские звери могут прогрызть днище корабля, не дожидаясь, пока своё дело сделают цунами.
Могильный пролив оправлен в скалы почти на всём известном протяжении, и едва ли вы найдёте хорошее место для причала.
Таким образом, этот путь в Край для нас закрыт, несмотря на кажущуюся его очевидность.

Хонхо провёл концом указки по хвосту Кита.

--- Эти берега, как я уже говорил, принадлежат народу трами, родичам идолов.
У нас с ними никогда не было тёплых отношений, трами вместе с ноа ходили на нас в походы, но их можно считать хорошими торговыми партнёрами.
Трами сообщили, что вот это море --- Голубое Зеркало --- охраняется могущественнейшим из племён дельфинов.
Голубое Зеркало можно даже считать их государством, потому что некоторые трами видели у берегов нечто, напоминающее полуводные поселения.
В отличие от своих северных сородичей, эти дельфины далеко не дружелюбны.
Любые суда они берут на сопровождение и для начала вежливо намекают, что лучше плыть обратно.
О судьбе тех, кто не подчинился, ничего не известно, как не известно ничего о тех, кто пытался установить с ними контакт.

--- Значит, этим путём в Край могут пойти дипломаты? --- спросил кто-то из детей.

--- Безусловно, Секхар из дома Тра.
Если кто-то из вас сможет договориться с этим воинственным дельфиньим племенем --- дорога в Край, по крайней мере в его восточные земли, называемые Тысячеречьем, будет открыта.
Дипломат-путешественник способен дать своему народу куда больше, чем самый великий воин.

--- Учитель, --- поднял руку сидевший рядом Столбик, --- я вижу на карте Края названия --- например слева, Заболочье...

--- Суболичье, --- поправил Хонхо-лехэ.
--- Вы ещё не разбирали правила чтения некоторых древних слов, просто запомните.

--- Суболичье, --- поправился Столбик.
--- А справа --- Тысячеречье, Каменное море, да и материк Кит, как вы сказали, был назван так из-за сходства очертаний с морским зверем.
Кто же узнал, как выглядят материки?
Кто дал этим странам названия?

--- Хороший вопрос, Ликхсар из дома Митр.
Карта, которую вы видите, является копией карты из библиотечной книги.
Как вы знаете, многие книги были написаны Древними, которые пришли из звёздной страны Тхидэ.
Вероятно, очертания материков они увидели с большой высоты.

Дети восхищённо забормотали, и Хонхо позвонил в колокольчик, призывая к тишине.

--- Но мы же не можем знать, что это правда, верно? --- настаивал Столбик.

Хонхо улыбнулся.

--- Я отвечу на вопрос так.
Карты из \emph{этой} книги --- карты известного нам мира --- весьма точны.
Кроме того, экспедиции, о которых я расскажу сейчас, подтвердили некоторые данные касательно Края.
Именно поэтому мы считаем, что автору можно доверять.

Указка пробежалась по западному побережью Короны.

--- Этот путь в Край, на мой взгляд, один из самых простых.
Путь искусного мореплавателя.
Здесь множество опасностей --- рифы, штормы, палящий зной, водовороты и цунами.
Именно здесь проходили в Край те, кто привёз истории о его необитаемости и непригодности к земледелию.
Но это лишь западное побережье.
Что творится в прочих местах, сказать нельзя.

--- А по суше?

--- По суше путь закрыт, --- отрезал Хонхо-лехэ.
--- Через смертные каменные пустоши не проходил никто.

--- Учитель, --- подняла руку девочка, --- а были ли те, кто уходил в Край насовсем?
Я имею в виду... разумеется, мы не знаем, добрались они или нет, но неплохо было бы знать хоть примерно, кто там... может быть.

Девочка стушевалась и опустила руку.
Хонхо-лехэ улыбнулся.

--- Я понял твой вопрос, Согхо из дома Лам, --- сказал учитель.
--- Сейчас посмотрим.
Ликхмас из дома Люм, Ломающий Храмы, встань.

Дети засмеялись.
Это прозвище, придуманное Хонхо-лехэ, лишь по случайности не приклеилось ко мне на улице.
Я остался сидеть.

--- Да встань же, --- улыбнулся учитель.
--- Я хочу, чтобы ты принёс из библиотеки книгу.
<<Хроники дорог и ветров>>, третий том, кажется... да, третий.

Я отправился в библиотеку.
Трукхвал, как обычно, встретил меня тепло.
Я еле отказался от сладостей и отвара, сославшись на то, что меня ждут на уроке.

Хонхо принял из моих рук тяжёлую книгу и открыл её на середине.
Полистал, нахмурился.

--- Вот, --- сказал он наконец.
--- В Край было четыре крупных экспедиции, одна из которых, Коричный флот, вернулась.
Первыми ушли, что неудивительно, ноа.
Они не стали испытывать судьбу ни в Могильном проливе, ни в Тихом океане, а двигались уже известным вам западным путём, преодолев по Торговому Потоку земли сели, тенку и ркхве-хор.
О дальнейшей их судьбе ничего не известно.

Хонхо перелистнул страницу.

--- Последнее колено Синих Травников.
Сели и даже ркхве-хор, что удивительно, предлагали им приют --- травники хороши как союзники, и Бродячий Народ желал воссоединения с сородичами.
Но Синие были чересчур подавлены собственным поражением в войне с Живодёром, да и Бродячих, как и многие из Красных, Синие презирали.
Они отправились через бассейн Реки Кувшинок, потому что, как и все травники, недолюбливали моря.
В любом случае им пришлось бы идти под защитой моря --- согласно данным Коричного флота, за дельтой Реки Кувшинок лежат ещё как минимум пятьсот кхене смертных пустошей.
Тем не менее очевидцы утверждают, что травники пошли сушей.
О судьбе Синего колена ничего не известно.

Учитель задумался и почесал подбородок.

--- Хай, вот.
Встреча у Мыса-лезвие, я ещё буду рассказывать вам о ней на истории.
Знаменитый пират-трами по имени Молочная-Роса-Движения, зная о дельфинах Голубого Зеркала и поняв, что флот Чистого союза отсек его кораблям путь на север, отправился на юго-восток.
По словам его противников, ветер был необычайно благоприятным, паруса пиратских кораблей раздувались, словно лёгкие бегущего.
Если его флот не погиб в море, то непременно должен был выплыть к берегам Края.

\section{[-] Пахарь}

\textspace

--- Все из вас слышали легенду о небесном пахаре, --- сказал я.
--- Так вот, сейчас я расскажу вам истину о кометах.
На самом деле это камни, которые падают со звёзд и раскаляются, падая на землю.

--- Так значит, нас обманывали? --- недоумённо спросила девочка.

--- Нет, --- сказал я.
--- Это не совсем обман.

--- Тогда зачем нужны легенды, учитель? --- продолжала допытываться она.
--- Зачем придумывать глупые истории, когда истина так прекрасна?

Я улыбнулся.

--- Как тебя зовут?

--- Митхэ ар'Тра, --- ответила девочка.
--- Масло-Взбитое-Взглядом.

--- Это тайна воспитания, Митхэ, --- ответил я.
--- Легенды вплетают человека в любое явление, чтобы такие умные дети, как ты, не забывали, на что способен человек.
Сказки дают человеку невероятную силу, чтобы вы мечтали и фантазировали, как претворить эту силу в жизнь.
Так мы движемся вперёд.

\section{[-] Жрица}

--- Ты с ума сошёл! --- буркнул Кхатрим, услышав мою идею.

--- И всё же...

--- Лис, это традиция, --- Кхатрим сел и прикрыл лицо рукой.
--- Думаешь, мне не попадались умные девочки, которых я желал бы взять в ученики?
Попадались.
Сказать по правде, умных девочек на моей памяти было даже больше, чем умных парней --- женщины быстрее взрослеют.
Но традиция есть традиция.

--- Почему нельзя её изменить?

--- Потому что живот никуда не денешь.

--- Ликхэ недавно диверсантов взяла беременной, и ничего, родила в срок без проблем.

--- Забудь про воинов, это другое.
На Западе вообще есть обычай, что командовать обороной должна беременная воительница --- враг на полёт стрелы к городу не подойдёт.
Потому что если какая угроза родному дому --- у них голова работает за десятерых.
А как будет выглядеть беременная жрица?
Жертвоприношение каждые несколько дней, а адепт Верхнего этажа с животом и под гормонами.
А если война или катаклизм?
Во время болотной лихорадки беременных и близко к карантинной зоне подпускать нельзя.

--- Кхатрим, женщины не всё время ходят беременными.

--- Однако беременеют в самый неподходящий момент.
Проверено.

--- Глупость какая-то.
Мужчина может заболеть, женщина может забеременеть.
Есть, в конце концов, белое жречество, почему нельзя...

--- Нет никакого белого жречества, --- перебил меня Кхатрим.
--- Запомни это раз и навсегда.
Это между собой мы так --- белые, красные.
Если долг позовёт, твой цвет не спросит ни народ, ни Безумные.

--- Но учитель же...

--- Если ты про Трукхвала --- не смотри, что он всё время кряхтит у себя в библиотеке.
Я видел, чего он стоит --- и во время набегов, и во время эпидемии, и во время похода.
Кстати, если ты не знал, он и с оружием обращается неплохо --- при его-то хромоте.

--- И всё же я хочу попробовать обучить эту девочку.

--- Храм её не примет, --- отрезал Кхатрим.
--- Хочешь --- пожалуйста.
Девочка будет благодарна тебе за образование, но образование --- это всё, что ты сможешь ей дать.

--- Возможно, стоит обсудить этот вопрос с остальными.
Ведь если в народе будут люди с образованием жреца, это может...

--- И эта идея не нова.
Пытались, и последний раз, как мне думается, не ранее чем сорок дождей назад.
Результат один --- нам переставали верить.

--- Но почему?

--- Если ты, жрец, сделаешь что-то вредное для народа и Храма, тебя можно просто выставить.
Тогда у народа формируется впечатление, что Храм блюдёт чистоту своих рядов.
А что делать с тем, кто де-юре крестьянин, но де-факто является ещё и жрецом?
В глазах народа он и будет жрецом --- ведь к нему придут и заболевшие соседи, да и роженицу проще тащить в соседний дом, чем в центр города.
Однако если что-то вредное сотворит этот человек, мы ничего не сможем с этим сделать.
Статус жреца --- и при этом никакой ответственности за этот статус.
Это подорвёт честь всего жречества.

--- Можно обвинить этого человека в Насилии.

--- Если Насилие имело место, что редкость.
А если он, к примеру, пишет инструкции, которые отличаются от храмовых в худшую сторону?
Даже если инструкция работает у одного из шестнадцати, эти свитки потом гуляют в народе, и очень сложно доказать что-либо крестьянам.
У купца по статистике падение урожая, а у крестьян один ответ --- <<у меня всё работает>>.
Можно, конечно, изымать и сжигать, как это делают тенку, но...

--- <<...ни действием, ни бездействием я не наврежу книге>>, --- закончил я.

--- Именно, Лис.
Это уже настоящее Насилие, без каких-либо сомнений.
Теперь понял?

--- Понял.
Но вреда от обучения как такового быть не может.

--- Это правда, --- признал Кхатрим.
--- Ты волен делиться знаниями, и я даже приветствую это.
Будет хорошим опытом для тебя как для учителя, да и в целом.
Но право принимать в Храм останется за Храмом.
Ты не первый, вопрос насчёт женщин не раз обсуждался и не стоит двух часов сна.
Для женщин Верхний этаж закрыт, и те, кто сменил пол, тоже обычно уходят.
Традиция есть традиция.
Если ты найдёшь способ убить Безумного и прекратить войны навсегда, можешь обучить делу хоть всё женское население города, --- жрец характерным жестом потрогал пальцем язык.
\ml{$0$}
{--- А пока --- извини, мне нужно работать.}
{``Before you're done, excuse me, I've got much work to do.''}

Я понурился.
Кхатрим задумался и снова вздохнул.

\ml{$0$}
{--- Она тебе нравится?}
{``Do you like her?''}

\ml{$0$}
{--- Она задаёт вопросы, на которые у меня в детстве не хватило ума, --- признался я.}
{``She asks questions I wasn't smart enough for, being the same age,'' I answered.}

\ml{$0$}
{--- Это серьёзно.}
{``A serious stuff.}
\ml{$0$}
{Что ж, я поставлю прочих жрецов перед фактом.}
{Well, I'll present that to other priests as a fait accompli.}
\ml{$0$}
{Мы будем делать вид, что так и должно быть, а ты не будешь подавать девочке ложных надежд, правда?}
{We'll pretend things are going on as they should, and you won't give false hope to the girl, will you?}
\ml{$0$}
{Благодарностей не нужно.}
{You're welcome, don't thank.''}

\section{[@] Гармония}

Сегодня при поддержке Безымянного мы установили ещё тридцать золотодобывающих установок.
Работа идёт полным ходом, механики не успевают собирать машины.
Нам не хватает стали и пластика, но Кусочек-Уха-Кабана придумал совершенно удивительную замену --- древесина.
С разрешения Безымянного мы аккуратно, стараясь не повредить экосистему, срезали некоторое количество старых деревьев, пропитали древесину специальными растворами и пустили в обработку.
Листик весь день ходила грустная.

--- Это варварство, --- заявила она.
--- Да и зачем резать деревья, если можно было просто сказать им, как расти?

--- Листик, мы не успеем написать вегетекторные хромосомы для всех видов конструкций, даже если задействуем вычислительные мощности Дракона на полную, --- грустно сказал Мак.
--- И даже если успеем --- сколько лет они ещё будут расти?
Безымянный разрешил --- уже хорошо.
Я бы на его месте отказался.

Товарищи уже успели привыкнуть к частому посещению Стального Дракона Безымянным.
Никто больше не шарахался, увидев в библиотеке мирно смотрящую тексты или фильмы фациограмму.
Время от времени я видел, что Безымянный задавал вопросы товарищам.
Те отвечали ему --- вначале неохотно, потом посвободнее.
Культуролог Гладит-Зелёную-Кошку общалась с ним душевнее других.
Видимо, она тоже почувствовала эту странную детскость, исходящую от демиурга.
Да и интересно ей как культурологу поговорить с хоргетом-девиантом.

--- Что такое музыка? --- спросил сегодня Безымянный у меня.
Мы с Баночкой, Кошкой и Заяц в это время обедали в моей каюте.

Гладит-Зелёную-Кошку подавилась пищевым концентратом.

--- Я же тебе показывала! --- с лёгким укором сказала она.

--- Показывала, --- виновато согласился Безымянный.
--- Но я так ничего и не понял.
Зачем складывать гармонические звуковые колебания разной частоты?
В чём смысл шумов?

--- Они вызывают у сапиентов положительные эмоции, --- ответил Баночка.

--- Не обязательно положительные, --- возразила Заяц.
--- Но эмоции --- это важное условие восприятия музыки.

--- Сапиенты начинают испускать эманации? --- уточнил Безымянный.
--- И всё?

Мы дружно опустили головы, не зная, что ответить.
Внезапно Кошка улыбнулась и кивнула:

--- Идём со мной.
И вы тоже.
Сейчас проведём небольшой эксперимент.

Кошка привела нас в аудиозал библиотеки и стала рыться в записях.

--- Кому что нравится? --- спросила она у нас.

--- Группа <<Солнечный цветок>>, октава 12, --- немедленно ответила Заяц.
--- Флейта звучит просто превосходно.

--- Я предпочитал <<Катарсис>>, симфонический оркестр эпохи Начала, октава 16, --- пожал плечами я.
--- Но это в далёкой юности.
Давно музыку не слушал.

--- Жанр <<электрошторм>>, октава 21, --- сказал Баночка.
--- Особенно группу <<Кровь и лёд>>.

--- Начнём с Заяц, --- ответила Кошка.
--- Баночка, ты не обидишься?
Просто октава 21 с полусинтетическим экстремальным вокалом --- это явно не то, с чего следует начинать знакомство с музыкой.

--- Не волнуйся, я привык, --- ухмыльнулся плант.
--- Как сказал Фонтанчик, под <<Кровь и лёд>> хорошо сходить с ума в горящем лесу, а я под неё засыпаю.
Но вообще это классика жанра, ей восемьсот лет.

--- Попозже, дружок.
Безымянный, ты вопринимаешь идущие от нас эманации?

--- Да, --- ответил бог.

--- Попробуй соотнести поток эманаций со звуковым потоком.
Сможешь?

--- Это не сложно, --- ответил хоргет.

--- Тогда поехали, --- сказала Кошка и включила запись.

Потекли нежные звуки колокольчиков.
Мягко и успокаивающе вступила флейта.
Утробным басом заговорила виола.
И вот, после недолгой беседы они слились в едином порыве, в единой песне, прославляющей что-то непонятное, далёкое... то, что осталось где-то на покинутой нами планете.

Зелёная трава, подсвеченная заходящим солнцем.
Колышущиеся метёлки злаков.
Моя первая любовь --- хрупкий, нежный красавец, с которым я в этой траве шестьдесят лет назад потерял свою невинность.
Имени его я вспомнить так не смог --- из-за переплавки тела и восстановления мозга стволовыми клетками ранние воспоминания амнезируются...

Баночка погрустнел.
Заяц кусала губу, уставившись в пространство полными слёз глазами.
Кошка выглядела так, словно сама была не рада своему выбору.
Эта музыка будила чересчур опасные, чересчур личные воспоминания.
Пальцы Кошки нежно гладили пульт управления, делая робкие попытки выключить... и не выключая.

Когда наступила тишина, мы ещё долго сидели очарованные.
Очнулись, только когда Кошка тихим голосом спросила у хоргета:

--- Ну, закономерности нашёл?

Голограмма хранила неподвижность.

--- Безымянный?
Эй!

--- Это... прекрасно, --- ответил Безымянный и исчез.
Снаружи, за толстой бронёй Стального Дракона, хлынул дождь.

Что творилось в этот момент в сознании хоргета?
Что он чувствовал?
Увы, для меня это навсегда останется тайной за семью печатями.
Мы привыкли считать, что мотивирующее чувство (<<удовольствие>>) у хоргетов возникает при поглощении одноимённых эманаций.
Вряд ли я испытывал что-то положительное --- мне было ужасно грустно.
Неужели и они могут видеть эту непередаваемую гармонию в корпускулярно-волновом хаосе?..

Больше Безымянный не задал ни одного вопроса.
Возхможно, что он что-то понял.
И я надеюсь, что так оно и было.

\section{[*] Музыка в голове}

\textspace

Цитра Атриса замолчала.
Митхэ сидела, опустив голову, и улыбалась.

--- Объясни мне, --- наконец прошептала она, --- как, как у тебя получается так красиво играть?

--- Вначале долго слушаешь, потом пытаешься играть.
А вскоре руки сами начинают что-то делать, --- смущённо пожал плечами Атрис и, аккуратно взяв руку подруги, положил её на цитру.
--- Попробуй.

Митхэ провела пальцем по струнам.
Цитра то ли фыркнула, то ли звякнула, тихо и недовольно.

--- Ну вот, --- засмеялась воительница.
--- Она на меня ругается.

--- Если будешь упражняться, она будет смеяться.
Доверие инструмента нужно заслужить.

Митхэ убрала руку с цитры и ткнулась головой в плечо менестреля.
Крохотная <<рыбка>> сверкнула огоньками костра.

--- Я лучше тебя послушаю.

--- А если тебе захочется музыки, а меня рядом не будет? --- улыбнулся Атрис.

--- Тогда я умру, --- полушутя, полусерьёзно ответила Митхэ.

Атрис задумчиво нахмурился.
Потом, бросив взгляд на реку, он вдруг засиял улыбкой и, потрепав воительницу по щеке, прыгнул куда-то в тростники.

--- Атрис!
Ну куда ты в болото?..

Из зарослей менестрель показался ещё более довольным.
В тонких пальцах он вертел сухой стебель <<муравьиного зонта>>, из каких дети делают игрушечные духовые ружья.

--- Ну и зачем?

Атрис вместо ответа с важным видом ушёл в лес.

--- Питонья отрыжка, --- выругалась Митхэ.
--- Как будто умрёт, если скажет пару слов.
Ладно, ладно, Хри-соблазнитель.
Хочешь гулять один --- гуляй, а мне завтра дежурить.
И так из-за тебя почти не сплю...

С этими словами воительница залезла в палатку и завалилась спать.

Проснулась Митхэ уже под утро.
Костерок весело пылал, пожирая свежие сучья.
Атрис лежал возле неё с открытыми глазами и улыбался.
Длинные тонкие пальцы сжимали красивую, расписанную чёрной, красной и синей красками флейту.

--- Это тебе, --- сказал менестрель.
--- Теперь музыка всегда будет рядом с тобой.

Митхэ осторожно взяла инструмент.
Флейта весила не больше пера сиу-сиу и внешне очень напоминала духовое ружьё.
Только там, где обычно рисовали лик Удивлённого Лю, находился Печальный Митр.
Губки флейты изображали его открытый рот.

У глаз Митхэ рассыпались счастливые морщинки.
Она долго лежала и смотрела то на флейту, то на любимого человека, не зная, что сказать.

--- Ну как?

--- Красивая, --- признала Митхэ.
--- Но я всё равно её сломаю.
Моя рука так же создана для музыки, как твоя для клинка.

--- Глупости.
Ты попробуй.
Возьми вот так... а, нет, на другую сторону.
Губы вот так.
А теперь дуй.

Митхэ дунула.
Звук получился очень смешной, и влюблённые расхохотались.
Воительница, утирая слёзы, положила дудочку на одеяло:

--- Нет уж...

--- Я тоже когда-то не умел, --- возразил Атрис.
--- Давай ещё раз...

Постепенно звук становился всё более похожим на музыку.
К тому моменту, когда солнечный край бросил первую горсть золотой пыли на верхушки деревьев, Митхэ испустила свою первую трель.

На огонёк зашёл ночной стражник.
Увидев Митхэ с флейтой, он покачал головой, пробурчал что-то про изнасилованного Митра и ушёл.

Наконец где-то вдалеке прокричали за утро.
Митхэ с некоторым сожалением отложила флейту и начала надевать доспехи.
Атрис лежал и смотрел на подругу.

--- Ты словно готовишь меня к скорби, --- сказала воительница.

Атрис устало вздохнул.

--- Я не хочу, чтобы твоё счастье зависело от меня или от кого-либо ещё.
Если я хоть что-то понимаю в любви --- это она и есть.

--- А если флейта сломается или её съедят жуки?

--- Найдёшь новую, --- пожал плечами Атрис.
--- Музыку в твоей голове не сломает никто.

--- Моя музыка в чужой голове.
И эта голова дурная, --- проворчала Митхэ и, на прощание взъерошив Атрису волосы, направилась к храму.
Менестрель, проводив её нежным и немного грустным взглядом, заполз в палатку --- наверстать бессонную ночь.

\section{[-] Библиотека}

\textspace

--- Учитель Трукхвал, а что здесь?

--- Ооо, Ликхмас-тари...
Здесь находятся самые ценные книги.

Я пробежался пальцами по корешкам.
Знакомые иероглифы... и совершенно непонятный смысл.

--- Эти книги написаны на языке Тхидэ.
Прародины людей сели.

--- Знаешь ли ты этот язык, учитель?

Трукхвал смутился.
От природы обладающий острым умом и невероятной жаждой знаний, учитель смущался только в одном случае --- если не знал.

--- Мне удалось кое-что разобрать, Ликхмас.
Мои предшественники тоже постарались, но перевести всё за эти годы мы не смогли.
Так и переписываем... хай... не понимая смысла.
Вот, например, название у книги --- <<Движение и время в почти симметричной пустоте>>.

--- Что???

--- У меня такая же реакция до сих пор, а я побольше тебя живу, --- развёл руками Трукхвал, снял книгу с полки и открыл её.
--- Имеется в виду межзвёздная пустота, но при чём здесь симметрия?
На первой странице пометка, уже понятным языком: <<Если вы когда-нибудь полетите к звёздам, вам понадобится>>.
Наверное, нам ещё рано.
Кстати, посмотри вон туда.

Учитель разгрёб свитки на одной из полок и показал мне маленькую, в локоть длиной, коробку.

Я наклонился над ней и провёл руками по гладкой поверхности.
Странный материал, не похожий ни на что из того, что я видел в своей жизни.
Лёгкий, полупрозрачный, немного опалесцирующий в свете лампы.
Сбоку знак --- число <<13>>, но никаких выступов, за которые можно было бы ухватиться и открыть.

--- Сокровища древних, --- улыбнулся Трукхвал.
--- Шучу.
Но доля правды в этой шутке есть.
Наверное, это что-то невероятно ценное, потому что никто так и не смог открыть этот ларец.

--- Он не особо похож на ларец.
А распилить не пробовали?

Трукхвал рассмеялся.

--- Узнаю Ликхмаса-тари.
Не удалось открыть --- надо распилить.
Помню, как Хонхо-лехэ прибежал ко мне и возмущённо закричал: <<Этого купчонка нельзя допускать в храм!
Он предложил переделать кладку с северной стороны, потому что ,,так будет устойчивее``!
А когда я спросил, что делать со старой, он сказал ,,сломать``!>>

Да, такое было.
Кормилец долгое время строил дома на востоке, и все тонкости кладки я знал к десяти дождям из застольных бесед.
Кстати, северную сторону всё же переделали, когда зодчий из Тхартхаахитра, залившись истерическим смехом, предложил обрушить стену храма двумя ударами.
Не особо сведущие в зодчестве жрецы поверили мастеру на слово и выделили необходимые материалы.

--- Я помню, --- смущённо улыбнулся я.
--- Учитель Хонхо разбил свою чернильницу и пригрозил отлупить палкой, если я ещё раз предложу такое.

--- Однако именно он понял, что ты гораздо талантливее других, --- заметил Трукхвал.
--- Да позаботятся о нём лесные духи.

--- Да найдёт он покой, --- поддержал я учителя.

Мы помолчали, глядя на таинственный сундучок.

--- Так что насчёт... --- начал я.

--- Никаких распилов, --- отрезал Трукхвал.

Я кивнул.

--- Тем более что я уже пробовал, --- извиняющимся тоном добавил учитель.
--- На ларце не оставляет царапин даже алмазный резец, не говоря уже о стали.
Точнее, оставляет, но они... трааа... зарастают.
Да.
Только пожалуйста, никому ни слова о том, что я вытворяю с реликвиями Древних! --- добавил старик страшным шёпотом.

--- Я нем, как скала, --- заверил я Трукхвала.

\section{[-] Математический анализ}

Так началась моя жизнь в библиотеке.
Я сидел и переписывал старые книги, Трукхвал рассказывал мне значения некоторых иероглифов.
Потихоньку старый язык, который ни я, ни мои прапращуры не слышали вживую, стал обретать смысл.

Природе нужно время для всего --- для прорастания, цветения, плодоношения.
В джунглях Тра-Ренкхаля время --- это единственный ресурс, который достаётся живой природе по каплям.
Первое, чему учится крестьянин, и второе, чему учится ремесленник --- терпение.
Я же схватывал науки и ритуал походя, без труда и ожидания, книги храмовыми ласточками пролетали передо мной, и в моём обучении этот важный момент --- терпение --- как-то прошёл мимо.
Пришло время навёрстывать упущенное.

Трукхвал заставлял меня по десять раз переписывать страницы, найдя в строгом наущении манускрипта ненароком затесавшуюся шуточку.

--- Но её же будет скучно читать! --- вздыхал я.

--- <<Я не испорчу...>> --- строго начинал Трукхвал.

--- << ... ни текста, ни художеств, и гулять моему перу лишь по полям>>, --- грустно заканчивал я.
--- Ну хоть на полях-то можно?

В конце концов Трукхвал сдался.

--- Можешь разрисовать поля, --- буркнул он.
--- Как угодно --- надписями, рисунками.
Только помни, что книгу будет читать другой человек.
Ему должно быть приятно её читать, это раз.
А ещё ему нужно место для собственных пометок, это два.
А вообще ученикам мы такое не разрешаем...

Строгость никогда не была коньком учителя, и моя импульсивная, похожая на морской прибой настойчивость очень часто брала верх.
Поэтому терпению я обучился лишь в самых общих чертах, но зато прекрасно усвоил, что такое компромисс.

\ml{$0$}
{--- И запомни ещё одну вещь, --- сказал напоследок учитель.}
{``And remember one more thing,'' the teacher said at last.}
\ml{$0$}
{--- Красота --- вечна.}
{``Beauty is eternal.}
\ml{$0$}
{А вот шутки со временем протухают.}
{But jokes are used to rot.''}

\ml{$0$}
{--- Что ты имеешь в виду?}
{``What do you mean?''}

Трукхвал молча достал с полки книгу и, открыв её, показал надпись на полях.
Я поморщился и высунул язык.

\ml{$0$}
{--- Трай, какая гадость.}
{``\Trai, that's disgusting.}
\ml{$0$}
{Пришло же кому-то в голову...}
{And there was the one to whom it had come to mind ...''}

\ml{$0$}
{--- Когда-то это было смешно, --- объяснил учитель.}
{``It was once funny,'' the teacher explained.}
\ml{$0$}
{--- Сейчас --- как видишь.}
{``You can see how it is now.''}

\ml{$0$}
{--- А кто переписывал? --- поинтересовался я, пролистывая книгу в начало.}
{``Who made the copy?'' I asked leafing the book to the title page.}
\ml{$0$}
{Почерк показался мне смутно знакомым...}
{The handwriting seemed vaguely familiar ...}

\ml{$0$}
{--- Я, --- коротко ответил учитель и отобрал у меня книгу.}
{``Me,'' the teacher briefly answered, then took the book from me.}
Его украшенные серьгами уши были похожи на красные бумажные фонарики.

У меня запылало лицо.
Остаток дня мы оба смущённо молчали.

Выносить книги из храма было строжайше запрещено, но мне очень не хотелось читать в душном, пропахшем старой кожей и опилками зале.
Время от времени я протаскивал то один, то другой том за пазухой и читал дома.
Учитель Трукхвал, видя меня неестественно выпрямившимся и весёлым, только вздыхал:

--- Что на этот раз, Ликхмас?

--- <<Математический анализ>>, --- шёпотом сообщал я.

--- Там одни непонятные значки и рисунки!
Возьми лучше что-нибудь из истории!

--- Вдруг интересное найду!

Трукхвал вздыхал ещё тяжелее и возвращался к переписыванию.

Увы, но <<Математический анализ>> мне не дался.
Учитель объяснил, что книга рассказывает об искусстве обращения с числами, но чисел там почти не было.
Я сидел над ней целую декаду, пытаясь понять смысл страшных многоэтажных столбиков, табличек и рисунков.
Иероглифов там было очень мало, а те, что были, не объясняли почти ничего.

<<Из этого следует...>>

<<Как мы видим...>>

<<Доказывает это следующее...>>

Для меня, не самого достойного потомка Древних, это звучало, как насмешка.
С некоторым сожалением я вернул книжку на её место и взялся за следующую.

\section{[-] Свободное время}

\textspace

Трукхвал не походил на других взрослых.
Может быть, именно из-за этого прочие жрецы и многие воины его сторонились.

Однажды я сидел за переписыванием толстенного тома <<Обычай и ритуал>>.
Учитель же свою норму переписывания давно выполнил и сидел скучал, раскладывая по форме древние кукхватровые детальки.
Где он их брал --- так и оставалось загадкой;
Трукхвал давно уже привык работать с детальками при мне, но едва в дверь библиотеки раздавался стук, детальки волшебным образом исчезали в бумажных кульках.

--- О чём думаешь, Ликхмас? --- спросил Трукхвал.
За год, проведённый со мной в библиотеке, учитель великолепно научился распознавать озадаченное выражение лица.

--- Я задумался о нашем обществе, --- сказал я.
--- Мне кажется, в нём есть какая-то несправедливость.

--- Что ты имеешь в виду?

--- У меня столько свободного времени.
Я могу заниматься тем, чем хочу.
У моих друзей --- ремесленников, крестьян --- его гораздо меньше.

--- Ты думаешь, несправедливо?
Взгляни на это с другой стороны, Ликхмас-тари, --- Трукхвал отхлебнул травяного отвара и ткнул пальцем в потолок.
--- Предки были мудры.
У нас есть Храм, у нас есть Двор, у нас есть Цех и Сад.
И у каждого сообщества свои обязанности.
Вот например, если в городе много больных, если дети не знают и не хотят знать грамоты, а Безумные обрушивают на народ свой гнев, кто виноват?

--- Верхний этаж, --- подумав, сказал я.

--- Правильно.
А если молодые не умеют держать оружия, не знают тактики боя, боятся сражаться или излишне жестоки?

--- Нижний этаж.
Воины не тренируют молодых и не учат их искать ногой твёрдую опору.

--- А кто виноват, если люди голодают?

--- Может быть, Сад?

--- Безусловно.
Двор тоже причастен к хозяйству, твоя кормилица после каждого дождя делает расчёты и пишет для крестьян план, чтобы урожай был разнообразным.
Однако недостаток урожая --- это вина Сада: крестьяне не умеют обрабатывать почву или разленились.
А кого ты обвинишь, если в городе мало завозных товаров, если племена дикарей совершают на нас набеги?

--- Плохо работает Двор.

--- Вот и всё.
Каждое сословие выполняет свою роль, и ошибки их видны.
Главное --- не закрывать на эти ошибки глаза и исправлять их вовремя, ведь одна ошибка тянет за собой другую.
Для решения проблем существуют Советы, которые каждый может созвать и на которых каждый может высказаться.
Свободное время же зависит от мастерства человека в его деле.

--- А если человек хочет заниматься чем-то другим?

--- Он волен выбирать.
Но ему придётся долго учиться, чтобы выбрать другое дело, а значит, жертвовать своим свободным временем.

--- Будущих жрецов выбирают жрецы.
Если взрослый человек захочет учить детей или врачевать, ему будет поздно выбирать.

Трукхвал развёл руками.

--- Бывали такие случаи.
С этим ничего не поделать.

--- В этом и есть некая несправедливость, --- сказал я.

--- Может, ты придумаешь что-то получше, --- улыбнулся учитель.
--- Для того тебе и дано свободное время.

\section{[-] Тысячи в одном теле}

Сегодня Митхэ сразила меня наповал новым вопросом.

--- Скажи, учитель Ликхмас, а я целая?

--- Что ты имеешь в виду?

--- Я как существо целая?

--- Разумеется, ученица.
Да, хоть ты и состоишь из крохотных живых существ --- клеточек, ты совершенно определённо целая как существо.

--- Тогда ответь мне, --- попросила Митхэ, --- кто обнимает игрушки во сне, если это делаю не я?

--- Почему ты решила, что это не ты?

--- Я в это время смотрю сны и не знаю об игрушках.

--- Сон --- особое состояние.
Возможно, ты просто этого не осознаёшь.

--- Тогда возьмём бодрствование.
Кто дышит за меня, когда я об этом не думаю?
Кто за меня идёт, когда я не обращаю на это внимания?

На этот раз я вынужден был сказать, что не знаю ответа на эти вопросы.

--- А однажды мне очень понравилась девочка, --- продолжила Митхэ.
--- Я хотела всё время проводить с ней.
А другая девочка мне не нравится.
И ты мне очень нравишься, и это тоже выбирала не я, потому что я не знаю причин.
Учитель Ликхмас, я точно целая?

--- Тебя это беспокоит, Митхэ?

--- Я знаю, что есть вещи, которые я не делаю, --- заявила Митхэ.
--- Например, солнце всходит и заходит без меня, и ветер тоже следует своим законам.
Однако то, о чём я рассказывала, происходит во мне, а не вне меня.

--- Возможно, что ты права, --- предположил я.
--- Тело --- это действительно множество существ, которые могут действовать самостоятельно.

--- И я должна с ними подружиться, --- подхватила Митхэ.
--- Ведь так?
Ведь для этого нужна школа, тренировки --- чтобы подружиться с собственным телом?

Я кивнул, в очередной раз думая о том же --- ребёнок интеллектуально превосходил меня на голову.

--- Как обучать того, кто умнее тебя? --- спросил я как-то у Трукхвала.

--- Я же тебя обучил, --- оскалился старый жрец.
--- Наверное, главное --- подавить в себе естественное желание ограничить ученика.
Мне порой очень хотелось.
Меня, хай, немного пугали твои мысли.
Они выглядели абсолютно логично, не придерёшься, но уходили далеко за пределы привычного.

--- Я очень боюсь сделать что-нибудь не так, --- признался я.

--- Ооо, а как я-то боялся, --- посетовал в ответ учитель.

Мы так долго хохотали, что на шум пришёл Первый жрец.

--- Что это ещё такое? --- строго сказал он.
--- Ну-ка брысь оба из библиотеки!
Нашли где смеяться!
Всю пыль с книг стрясёте хохотом!

--- По-моему, давно пора стрясти с них пыль, --- шёпотом заметил Трукхвал.

\section{[M] Книжная пыль}

Однажды мы с Трукхвалом отправились в старый отдел библиотеки за ингредиентами для красок.
Там было сухо и чересчур пыльно.
Кладка в некоторых местах просила хэситр, а полки прогибались дугой, словно западные рогатые луки.
Пыль играла с лучами масляной лампы, пыль клубилась у наших ног.
Меня не отпускало чувство, что эта пыль --- книжная.
Она пахла старыми чернилами.
Я поделился этим ощущением с учителем, на что он вздохнул:

--- Так оно и есть.
Здесь мы храним книги, к которым уже нельзя прикасаться.

--- Почему, учитель?

--- Потому что они превратились в прах.
Видишь, полки закрыты занавесками?
Книги рассыпаются от простого прикосновения.
Хаяй...
Больно это видеть.

--- Почему их не переписали, учитель Трукхвал?

--- Пятьсот дождей назад, как ты знаешь, Тхитрон был опустошен идолами...

--- Нашествие Змей, --- кивнул я.

--- Люди отбили его много позже.
Увы, это то, что они не успели переписать.
Так и храним...
Надо, конечно, в Тхартхаахитр съездить, сделать копии, если в их храме есть, но куда мне с моей... хай...
Вот, посмотри.

Учитель осторожно отвёл занавеску.
Обычные на вид книги.
Но лежавшая рядом слоистая куча пыли показывала, что с этими книгами происходит при неосторожном прикосновении.

--- Странно, что идолы не сожгли эти книги.

--- Ничего странного, Ликхмас.
Идолы не дураки.
Многие даже уносят книги в свои деревни.
Правда, понятия не имею зачем.
Говорят, пытаются читать или просто хранят их в своих святых местах.

--- Учитель, даже идолы относятся к ним с почтением.
Не должны же они так лежать и рассыпаться у нас!

--- А что делать? --- развёл руками Трукхвал.

--- Может, их можно как-то склеить?

--- Как, ученик?
Это же пыль!

Я задумался.

Обратный путь мы проделали в полном молчании.

Следующие дни я раз за разом возвращался в мыслях к древнему тёмному помещению.
\ml{$0$}
{Книги снились мне каждую ночь: я пытался снять том с полки, а он рассыпался у меня в руках.}
{I dreamed every night that a book has come apart in my hands when I try to take it off the shelf.}
Просыпался я в холодном поту.
Кхотлам за завтраком стала посматривать на меня с беспокойством.

\ml{$0$}
{--- Ликхмас, дитя, ты совсем не ешь.}
{``\Likchmas, child, you hardly eat.}
\ml{$0$}
{Ты не заболел?}
{Are you sick?''}

\ml{$0$}
{--- Нет, кормилица.}
{``I'm not, nurse.''}

--- Сестрёнки потеряли тебя.
Ты целых девять дней не заходил к ним перед сном и не желал спокойной ночи.
Днём тебя тоже не найти.
Негоже так.

Я мысленно проклял эти книги.

--- Где они сейчас?

--- Играют у реки с друзьями.
Поинтересуйся, как у них дела.

День выдался погожий.
Только-только закончился сезон дождей, и прохлада влаги ещё не превратилась в неприятную духоту.
Ветер ласково трепал мои волосы.
Да, засиделся я в библиотеке.

Вокруг кипела работа.
Приречные крестьяне возили на свои участки отходы --- кто по реке, на каноэ, кто в наземных деревянных повозках.
Кудахтали куры, хрюкали довольные свиньи, мяукали и ворчали домашние оцелоты, выпрашивая потрошки пожирнее.
Люди оживлённо переговаривались, разбрасывая нечистоты по грядкам, залихватски ухали и пели песни, перекапывая жирную, влажную почву.
В воздухе стоял крепкий запах навоза и пота.

--- Лисичка, как твои дела? --- крикнул мне загорелый парень-крестьянин, с которым мы часто играли в детстве.

--- Слава духам, Столбик.
Гляжу, у тебя сезон?

--- Хорошая почва --- хороший урожай будет, --- весело откликнулся Столбик, сдвинув на затылок травяную шляпу.
--- Ты так же, в библиотеке, книги пишешь?
Заходи как-нибудь вечером в гости, а то солнца не видишь!

Я помахал ему и пошёл дальше, на песчаный берег, на игровую площадку.
Дети носились, прыгали и кричали.
Простая счастливая жизнь.
Я замедлил шаг, высматривая в толпе сестрёнок.

--- Хаяй, мой домик!

Маленький ребёнок обиженно смотрел на меня, сидя на песке.
Кажется, я не заметил его песчаного домика и наступил на него.

--- Ой, прости, малыш, --- улыбнулся я и присел на корточки.
--- Давай я построю тебе новый.

--- Давай, --- личико ребёнка, на котором уже начали проступать женские черты, просветлело.

Я разгладил руками мокрый холодный песок и запустил в него пальцы.

--- Как тебя зовут?

--- Эрси ар’Митр, --- гордо сообщил ребёнок, с интересом глядя на стройку, и потеребил пальчиками мою робу.
--- А ты жрец?

--- Пока ещё нет, --- засмеялся я.
--- Но скоро буду, если духи помогут.

--- Ты из дома ар’Люм?
Твои сестрёнки ушли домой, --- рассказал ребёнок, не дожидаясь вопроса.
--- Кажется, к вам кто-то приехал, и нужно было проводить гостей.
Ты ведь за ними пришёл?

--- А вот чтобы тебе домик построить.

--- Ты врёшь, --- насупился ребёнок.

--- Правда-правда, --- заверил я его и добавил к полувшемуся дворцу пару завершающих штрихов.
--- Эрси ар’Митр, можешь заселяться.

--- Он красивый, --- сказал ребёнок и критически осмотрел домик со всех сторон.
--- Жаль только...

--- Что?

--- Пока песок влажный, он стоит, --- объяснил Эрси.
--- А когда взойдёт солнце и подсушит его, он превратится в гору пыли.

В моей памяти молнией сверкнула пыльная библиотека и рассыпающиеся книги.
В храм я бежал быстрее лани.

\asterism

--- Ликхмас, ты с ума сошёл! --- воскликнул учитель, выслушав мою идею.
--- Книги рассыпаются при прикосновении, а ты вздумал лить на них воду!..

--- Я не собираюсь лить на них воду, --- терпеливо принялся объяснять я.
--- Пыль пока имеет форму.
Надо закрепить её чем-то жидким.
Склеить.
Ведь можно же касаться книг очень мягким, смоченным в...

--- Хорошо-хорошо, я тебя понял, --- Трукхвал, похоже, сам заинтересовался.
--- Давай подумаем, что можно использовать.
Вода, конечно, отпадает сразу.
Который час?

--- Ещё даже не полдень, учитель.

--- Бегом на рынок, возьми рыбьего клея.
Хай, ты же ещё не жрец, бесплатно тебе ничего не дадут.
Вот перо золотого песка, купи сразу побольше.

Мы развели в храме нешуточный переполох.

--- Трукхвал, что за дела? --- возмутился Кхатрим, только собравшийся готовить лекарства про запас.
--- Куда ты потащил мой котёл?

--- Ликхмас ар’Люм!
Верни перья сиу-сиу на место!
Из-под палок не встанешь, выкидыш идола!

Воины только посмеивались, глядя на нас.

--- Что старый, что молодой, а ума --- как у ящерицы.

Уже под вечер мы с Трукхвалом стояли в старой библиотеке.
На подставке дымилась чаша с тёплым жиденьким рыбьим клеем, в нём плавали самые мягкие перья сиу-сиу, которые я смог отобрать.

--- Может, муки чуть побольше? --- шёпотом спросил учитель.

--- Нет, --- покачал я головой.
--- На этот раз мы всё смешали правильно.
Смотри, страничка отошла почти полностью, расслоилась только в самом конце.
И всё равно кое-что можно прочитать...

--- Страничку-то ты спас.
А книжка? --- посетовал Трукхвал.

Мы с некоторым сожалением посмотрели на останки предыдущего эксперимента.

\ml{$0$}
{--- Эта книжка совсем сгнила.}
{``This book was completely rotten.}
\ml{$0$}
{Нам просто не повезло.}
{Just bad luck.''}

--- Ликхмас, давай... --- жалобно начал учитель.

--- Нет, --- твёрдо сказал я.
\ml{$0$}
{--- Мы будем пытаться, пока не получится --- или пока мы не превратим в пыль всю эту секцию.}
{``We shall try until we achieve success---or until the whole section have been turned into dust.}
Смотри вон на тот огрызок, он всё равно ни на что не годен.

\ml{$0$}
{Трукхвал только горестно заохал.}
{\Trukchual\ groaned in despair.}
Он знал, как сложно меня отговорить, если я чем-то увлёкся.

Я, не обращая на него внимания, задумчиво смотрел на освещённую неверным светом лампы одинокую книжку, лежавшую отдельно от остальных.
Судя по корешку, когда-то это был целый том в три пальца толщиной, сейчас от него осталось несколько десятков страниц.

--- Хай, а если опять не получится, ученик? --- пробормотал Трукхвал, чуть не плача.
--- Святыня...
Ещё одна...
А потом ещё...

--- Святыни, которые всё равно рассыплются, --- пожал плечами я.

Трукхвал нервно усмехнулся.

\ml{$0$}
{--- Твоя правда.}
{``You're right.}
\ml{$0$}
{Давай-ка ты на этот раз, твои руки нежнее моих.}
{This turn is yours, your hand is softer than mine ever was.''}

Я обмакнул перо в клей и начал осторожно промазывать обложку и корешок.
Подождал, пока клей впитается, и промазал ещё раз.
Трукхвал, до хруста сжав костлявые кулаки, наблюдал за процессом.

Время от времени я пробовал остриём пера обложку на прочность.
С самого края, чтобы не повредить странички.
Первый раз перо провалилось почти на горошину, Трукхвал охнул.

\ml{$0$}
{--- Давай ещё.}
{``Go on.}
\ml{$0$}
{Не торопись.}
{Go slow.''}

Ночь пролетела незаметно.
Мы так увлеклись, что не видели и не слышали ничего вокруг.
И вот, уже под утро, обложка поддалась и отошла, не разрушив странички.
Мы с учителем, как зачарованные, уставились на заглавный лист.

<<Повесть о Существует-Хорошее-Небо, вожде сели, наезднике Железного Змея>>.

\asterism

На следующий день Трукхвал рассказал о нашем эксперименте другим жрецам.
Похвал мы не дождались --- жрецы, особенно <<красные>>, скупы на это, --- но ни один, надо отдать им должное, не упрекнул нас в святотатстве и излишнем авантюризме.
Те, что постарше, выразили желание ознакомиться со спасёнными книгами, а жрецы помоложе пошли с нами --- лично принять участие в спасении.
В библиотеке запахло рыбьим клеем и потом.

Учитель был совершенно счастлив.
Его нежданные помощники приносили ему страничку за страничкой.
Трукхвал сортировал листки, добавлял замечания на краях или клочках старой бумаги, время от времени сокрушённо качал головой, когда текст оказывался безнадёжно испорчен, а меня усадил за переписывание.

--- Ликхмас, здесь ключ <<листик>>, а не <<рыбка>>.
И вот здесь.
Да, у переписчика был странный почерк, согласен.
Нет-нет, не надо переделывать всю страницу, подотри пергамент шкуркой.

--- Ликхмас, эмоглиф в языке Тхидэ идёт сразу после глагола, модификатор глагола ставится над эмоглифом.
Да, всегда.
Да откуда я знаю почему?!
Просто перепиши это и будь внимательнее.

Вечером я заглянул к сестрёнкам и извинился за невнимание целой горстью сладких конфет.
Они наперебой рассказывали о своих играх и тренировках, а я сидел и гладил их по светлым, цвета какао с молоком, волосам.
Зашедшая пожелать спокойной ночи Кхотлам стояла и, улыбаясь, наблюдала за нами.
Вдруг я вспомнил, что кое-кому должен...

--- Молочко, --- спросил я у сестрёнки, --- ты знаешь Эрси ар’Митр?
Он играл с вами вчера днём.

--- А, молчун Котёнок, --- кивнула сестрёнка.
--- Он сегодня не пришёл.
Кусачка, ты видела его?

--- Не-а, --- помотала головой близняшка.

--- Эрси ар’Митр? --- вмешалась кормилица.
--- Его сегодня ночью принесли в жертву.
А почему ты спросил, Лисёнок?

\chapter{[-] Семья}

\section{[М] Доблесть}

\textspace

Вспомнил свою первую схватку с настоящим врагом.
Кормилица отправилась в деревню за рекой, чтобы помочь в решении каких-то проблем с товарами.
Я, двадцати пяти дождей, отправился вместе с ней.
В окрестностях орудовала банда разбойников, и случай столкнул меня с двумя мужчинами.
Они были умелыми бойцами и играли со мной, как сытые коты с мышью.
Я сумел полоснуть одного из них ножом, но рана лишь разозлила его.
Тогда, забившись в узкий проход, я зарядил духовое ружьё и птицей попросил помощи.

--- Хэй, мальчик, выходи! --- издевательски кричали разбойники.
--- Или ты трусишь?

Я испугался.
Мои руки ослабели, нож дрожал так, словно меня уже избили до полусмерти.
Ни о каком бое не могло быть и речи.

Подоспевшие жители спугнули разбойников, и я вышел на свет.
Все молчали, и я тоже молчал, думая, что меня считают трусом.

--- Ты как? --- наконец спросил Кора, ученик плотника чуть постарше меня.
Я не ответил.
Меня била крупная дрожь.

--- Все мы боялись в наш первый бой, Ликхмас ар’Люм, --- хлопнул меня по плечу Кора.
Остальные закивали, выражая согласие.

--- Их двое, и они куда опытнее тебя.
Тебя бы просто убили или покалечили.
Молодец, что спрятался.

--- Трусы, --- добавил кто-то.
--- Двое взрослых мужчин на одного мальчишку.
И сбежали, как крысы, едва увидев нас.

Это слабо утешало.
В тот вечер я ушёл к реке и долго плакал от обиды и злости на самого себя.

\asterism

Следующие две декады я не мог отойти от собственного мига трусости.
Разбойники начали сниться мне по ночам.
Я отчаянно кричал и пытался поразить их ножом, но они лишь смеялись и издевались надо мной, дымными змеями ускользая от ударов.

Тренировки стали настоящей мукой.
Казалось, что каждый спарринг-партнёр знает о случившемся и втайне смеётся надо мной.
Одному из них --- крепкому, но неповоротливому Огранённому --- я сломал нос и на миг испытал невероятное наслаждение, глядя на его испуг и бегущую из ноздрей тёплую кровь.
Конфетка покачал головой и повёл парня в храм, к жрецам.

Вскоре молодые воины стали меня избегать.
Я не был лучшим в поединках, но ужасное новое чувство заставляло любыми путями достигать победы.
В ход шло всё --- подручные средства, психологическое давление и довольно рискованные приёмы.
Ещё спустя пятнадцать дней к жрецам отправились трое парней, попытавшихся после тренировки разобраться со мной <<по-воински>>.

Вечером в дверь постучались.
Следом ко мне заглянула растерянная кормилица и сказала, что пришёл Конфетка.

--- Пойдём, --- коротко и довольно неприветливо сказал тренер.

Конфетка привёл меня к реке.
Солнце уже давно зашло, и река казалась сплошным чёрным пятном.
Где-то нежно квакала лягушка, чирикала певчая речная черепашка и весело плескались рыбки, ловя на лету жирных насекомых.

Конфетка устало сел на поросший мхом камень и скрестил руки на груди.

--- Ликхмас, объясни, что с тобой происходит.

--- Они сами начали... --- бросил я заготовленную фразу, но тренер нетерпеливо махнул рукой.

--- Я прекрасно знаю, когда всё началось.
Ты уже две декады сам не свой.
Что это?
Откуда такая жестокость?
Куда подевался мальчик, который когда-то отказался бить слабого?

--- Они не слабые, --- огрызнулся я.
--- Их было трое.

--- А Костёр?
Ты понимаешь, что ты мог его убить <<хлыстом двух скал>>?
Я говорил --- приём отработать, но против друг друга не применять ни в коем случае.
Говорил?
Говорил.
Парня спасло только то, что он в защите куда лучше, чем ты --- в атаке.
А Огранённый?
Чем этот увалень тебя обидел?
И хватит отнекиваться.
Ты прекрасно знаешь, как и куда нужно бить на тренировках.

Я молчал.
Конфетка вздохнул.

--- Кхотлам сказала, что в деревне Кхихутр что-то произошло.
Расскажи мне.

--- Как будто ты не знаешь, --- снова ощерился я.

--- Нет, я ничего не знаю, --- устало сказал тренер.
--- И Огранённый, и прочие тоже в неведении, но исправно получили за это свежих пирожков.
Так давай, расскажи, а я попробую объяснить парням, за что им досталось.

--- За дело, --- рявкнул я и направился к дому.

--- Трус, --- бросил Конфетка мне в спину.

В моей голове застучала кровь, и я, оттолкнувшись от ствола молодого дерева, изо всей силы ударил тренера локтем в глаз.
Конфетка вскрикнул и упал.

Прошла целая михнет --- и до меня вдруг дошло, что он даже не собирается вставать.
Тренер просто лежал на земле, прижав руки к лицу и свернувшись калачиком.
Я склонился над ним.

--- Хай, Конфетка.
Ты жив?

Тренер не ответил.
Я потряс его за плечо.

--- Встань, пожалуйста.
Прости меня, я не хотел.

Ответа не было.
Я заплакал и начал отрывать его прижатые к лицу ладони.

--- Прости меня!
Встань, пожалуйста!
Прости!

Конфетка медленно сел и открыл лицо.
Бровь была разбита, и опухшее веко совершенно скрыло глаз.

--- Ликхмас, --- тихо сказал тренер, --- можешь ударить меня ещё раз, но расскажи, что случилось.

Этот взгляд и кроткие слова мне суждено было запомнить на всю жизнь.

\asterism

Время уже давно перевалило за полночь.
Мы с Конфеткой сидели дома за дымящимися чашами с травяным отваром.
Домашние уже спали, я слышал сонное бормотание Эрхэ и похожий на урчание ягуара храп Сиртху-лехэ.
Только кормилица, так и не задав ни одного вопроса, отправилась на улицу выбрасывать вареные листья.

Конфетка пил отвар и смотрел на меня целым глазом.
На другой глаз я сделал ему примочку.

--- Кхотлам я ничего не сказал и не скажу, --- сообщил Конфетка.
--- Это её не касается.
Ответственность за произошедшее лежит на тебе и на мне.

Я кивнул.

--- Знаешь, --- тихо сказал тренер, --- в детстве меня часто били в Храме.
Но чаще, чем от сверстников, удары я получал от тренера.
Настоящие, с охотой, не чета моим.
Он считал, что если детей почаще бить и ругать, то они станут храбрецами.

--- Но ведь ты стал, --- возразил я.

--- Нет, --- улыбнулся Конфетка.
--- Не стал.

--- Но ты же...

--- Да, я многое повидал, достаточно повоевал.
У меня был выбор --- завести семью с подходящими людьми, осесть, стать ремесленником.
\ml{$0$}
{Я подрабатывал когда-то краснодеревщиком, ты знал?}
{I was a cabinetmaker once, did you know?''}

\ml{$0$}
{--- Нет, не знал.}
{``I didn't.}
\ml{$0$}
{И как далеко ты зашёл?}
{How far has it gone?''}

--- Хай...
Ну, ты сидишь за столиком моей работы, --- Конфетка нежно погладил столешницу.
--- Если приглядишься, то увидишь, что вот здесь небольшая трещина, заполненная древесным тестом.
Мне привезли плохие доски --- и тем не менее я смог сделать стол, достойный Двора.
И для Верхнего этажа сундуки делал тоже я.
Не знаю насчёт остальных, давно не заходил, но у Трукхвала в келье точно мой стоит, до сих пор --- с узорчатым бронзовым замком на шестнадцать штифтов, который было ужасно сложно ставить.
В общем, да, мои работы ценились.
Но я не считал их чем-то выдающимся, заслуживающим уважения.
Я выбрал путь воина, потому что хотел доказать себе, что я не трус.

--- Так ты и не трус.

--- Много раз было так, что я действовал в ущерб себе, боясь, что меня посчитают трусом.
Это тоже было трусостью, как я сейчас понимаю.

Я промолчал.
Конфетка рассеянно гладил пламя свечи пальцем.
Не любящий фамильярностей язычок сердито пытался отодвинуться.

--- Когда-то я вёл беззаботную жизнь.
Походы, дерзкие вылазки.
Каждый шрам казался мне благородной каллиграфией.
У меня была сабля, я её потом продал, тонкая, как ивовый лист --- я рассекал ею летящих птиц пополам, и это тоже казалось мне доблестью.
Любовников у меня было не так много, как могло бы быть --- многие смотрели с надеждой в мою сторону, но ведь настоящий воин и храбрец должен держать марку, верно?
У меня была, стыдно признаться, даже коллекция ушей поверженных врагов.

Я поморщился.
Конфетка кивнул и закрыл лицо руками.

--- Да, я сегодняшний отреагировал бы точно так же.
Но мне казалось, что это и есть доблесть.
И поэтому, когда Морковка после одной из ночей попросила разрешения похоронить эти уши, я назвал её дурой и сказал, что ухожу навсегда.
И это тоже казалось мне доблестью.
Она не заплакала, просто молча собрала мешок и так же молча поцеловала на прощание.
На самом верху лежал свежий черноягодный пирог.
Выйдя за ворота, я со смехом бросил пирог в канаву.
Ведь воин не должен размениваться на ерунду и баловать своё тело.
Надеюсь, что она его потом не нашла.

Конфетка говорил совершенно спокойно, но по моей спине пробежала дрожь.

--- Однажды я возвращался в родной город из Тиши и шёл мимо её дома, --- проговорил тренер, и целый глаз вдруг вспыхнул.
--- В окнах горел свет.
Там весело переговаривались дети.
Я всегда любил ребятишек, но когда услышал детские голоса из её дома, у меня защемило сердце.
Странное чувство, да?..

Я кивнул.

--- Сама она сидела со своим мужчиной на скамье и любовалась звездами.
Они были такие красивые, оба.
Они разговаривали о том, как же, каким искусством сплели полотно этого мира.
Наверное, Безымянный был очень трудолюбив, раз создал такое количество звёзд, небо, землю и зверей.
Так они говорили.

Конфетка опустил голову, и отвар стал чуть солонее.

--- А я понял в тот миг, что мир получился таким прекрасным просто потому, что его создателю не перед кем было кичиться.
Ему не нужно было чьё-то внимание, чьё-то одобрение.
Некому было сказать, что он растяпа, трус или лентяй.
Он просто делал то, что хотел.
И у него всё получилось.
Он, я полагаю, счастлив.

Конфетка кивнул и отхлебнул из чаши.

--- На следующий день я пришёл в Храм и во всеуслышание заявил то, что не давало покоя мне все эти годы.
Я сказал: <<Да, мне страшно.
Я боялся сражаться, когда шёл через Тиши, но шёл, потому что боялся проиграть.
Я положил годы на ужас походов, потому что боялся показаться трусом.
На самом деле я и есть трус.
Одерживая победу оружием, я проиграл жизни, в страхе потеряв то, чего хотел больше всего.
Но сегодня этому пришёл конец.
Я говорю всем: мне страшно, и я этого больше не стыжусь>>.

Конфетка улыбнулся.

--- А что было потом? --- тихо спросил я.

--- Потом я сказал, что больше никогда не возьму в руки боевое оружие.
Разумеется, это было отречением от воинской клятвы, но воины, надо отдать им должное, поняли меня и не позволили уйти из Храма.
Боевой вождь, моя лучшая подруга, спросила: <<Что тебе нужно?>>.
Я ответил: <<Мне нужна семья>>.
Так я стал тренером.

Я вздохнул.
Разумеется, это не была настоящая семья --- лишь суррогат, попытка попробовать на вкус давно ушедший в прошлое пирог.
Возможно, Конфетка больше не чувствовал в себе сил на постройку собственного жилища, на изменение уклада собственной жизни, и решил довольствоваться тем, что есть. Некогда кипевшая энергия ушла на никому не нужную иллюзию.

--- С тех пор моя жизнь изменилась.
Я превратился в сладкоежку, отрастил брюшко, стал любителем поспать и, надо сказать, большим любителем секса тоже.
Но что-то ушло безвозвратно, и иногда я даже не могу понять, что именно.

\ml{$0$}
{--- А та женщина простила тебя? --- спросил я.}
{``That woman, did she forgive you?'' I asked.}

\ml{$0$}
{--- Наверное, --- повеселел Конфетка.}
{``May be,'' Candy cheered up.}
\ml{$0$}
{--- Когда я пришёл к ней поговорить, то получил котелком по загривку.}
{``When I came to talk to her, I got hit on the neck with a cauldron.}
\ml{$0$}
{Горячим котелком, Ликхмас, только с огня.}
{A very hot cauldron, \Likchmas, freshly heated.}
\ml{$0$}
{Четыре раза.}
{Four hits.}
\ml{$0$}
{Думаю, да, простила.}
{I guess, yeah, of course she did.''}

\asterism

Я открыл дверь, и в дом ворвалась скрипучая песня тысяч сверчков.
Конфетка поправил плащ, прислушался и щёлкнул языком от удовольствия.

--- Всё-таки интересно, как Безымянный додумался до сверчков.
<<Скучный какой-то жук.
Может, он будет тереть лапками крылья, чтобы найти самку?>>, --- Конфетка прыснул.
--- Вот они и пилят по ночам, стараются.
Но изобретение однозначно хорошее и нужное.

--- Ты веришь в Безымянного? --- удивился я.
Конфетка всегда казался мне далёким от религии нигилистом.
\ml{$0$}
{--- Почему?}
{``Why?''}

--- Жил-был некто, кто создал много хороших вещей, его изгнали и отобрали у него творение, --- Конфетка болезненно поморщился.
--- История чересчур похожа на правду, чтобы не быть ею.

--- Я сменю тебе повязку завтра, --- пообещал я.
--- Только не тереби, а то я тебя знаю.

Тренер расхохотался и тут же, опомнившись, прикрыл рот ладонью.

--- Хаяй...
Ночь уже, а я тут...
Не волнуйся, не буду.
А вот к парням зайди обязательно.
Что им сказать, знаешь.

--- Что я трус, --- кивнул я.

--- Дурень, --- покачал головой Конфетка.
--- Прощения попроси.
А насчёт храбрости... ну что поделаешь, Ликхмас.
Истинным храбрецом, увы, можно лишь родиться.
Но и осторожность, и сострадательность --- недешёвый товар.

--- Но ведь мне придётся сражаться, --- сказал я.
--- Значит, нужно как-то подавлять страх?

--- Ты когда-нибудь надрывал мышцу, Ликхмас?

Я кивнул.

--- Представь, что ты надорвал мышцу и тебе нужно поднять груз.
Можно разработать мышцу, чтобы она болела меньше.
Иногда приходится прикидывать, как изловчиться, чтобы её не потревожить.
А иногда... иногда приходится поднимать, невзирая на боль.
Не нужно только убеждать себя и окружающих, что у тебя ничего не болит.
\ml{$0$}
{Я потратил на это полжизни и не обрёл ничего.}
{I spent half my life trying, and got nothing in return.''}

Я кивнул.

--- Твои чувства --- это такие же мышцы.
Нужно принять их такими, какие они есть, и хорошо знать их особенности.

--- И тогда я стану непобедим?

--- Нет, --- покачал головой Конфетка.
--- Ты станешь счастливее.

Конфетка на секунду отвернулся и вдохнул носом влажный ночной дурман.

--- Но ведь чувствами нельзя управлять, --- сказал я.

--- Конечно, можно.
Ты же переносишь вес на здоровую ногу, поднимая груз, верно?
Так же и с чувствами.
Если боишься сражаться, можно позволить другому чувству овладеть собой.
Например, любви.

--- <<В жизни и смерти буду я щитом>>, --- процитировал я.

--- Верно, схватил суть, --- Конфетка похлопал меня по плечу.
--- Хорошая у вас, жрецов, клятва.
Нашей, воинской, чего-то недостаёт.
Наверное, границ защиты.
Войны сейчас уже не те, они лживые по сути --- в них нет агрессоров, одни защитники.
Как говорит Ситрис, <<словно кто-то выбросил слово из клятвы и памяти людской...>>\FM\
\FA{
<<... и прахом стали верность и честь, стоявшие на зыбкой почве песен>>.
Строфа из плача Маликха по Ликхмасу, <<Легенда об обретении>>.
}

--- А как распознать грань между вторжением и обороной?

Конфетке понравился мой вопрос.
Таким довольным я не видел его никогда.

--- Духи лесные.
Как же ты мне напоминаешь одного человека.
Он искал ответ на этот вопрос всю свою жизнь.

--- Где он его искал?

--- В кошачьих слезах, сосчитанных песчинках и утренней росе\FM, --- отшутился тренер.
\FA{
Имеются в виду три основных философских школы боя --- Путь Ягуара, Десять Песчинок и Плющ и Капля Росы.
}

--- Где сейчас этот человек?
Он нашёл ответ?

--- Хаяй, кусочек где-то здесь валяется, --- Конфетка ущипнул меня за бок, --- а где остальные...

Тренер вдруг замолчал и торопливо начал пристёгивать оружие.

--- Ты знаешь, я бы очень хотел, чтобы ты принёс воинскую клятву, --- признался он.
--- Даже хотел в пику Храму, глухой ночью, в строжайшей тайне, но... думаю, не нужно.

--- Я принесу её, --- сказал я.
--- Если ты хочешь.
Сейчас как раз глухая ночь.

Конфетка грустно улыбнулся.

--- Я тебя слушаю.

\begin{quote}
<<Я не причиню больше страданий, чем это необходимо;\\
я не оскверню звание воина ни пыткой, ни мучительным убийством, ни угрозой близким;\\
я не подниму руку на безоружного и неумелого;\\
я не позволю сделать подобного ни врагу, ни другу.\\~\\
Я не сломаю дух ребёнка ни оскорблением, ни ударом, но сделаю его стойким;
я обучу его всему, что знаю и умею в бою, но в первую очередь --- тому, что поможет избежать боя.\\~\\
В жизни и смерти я буду клинком, и ни мгновения --- плетью>>.
\end{quote}

--- Мне бы следовало дать тебе пощёчину, чтобы ты запомнил как следует, --- вздохнул тренер.
--- Но если этот удар лишний --- а я в этом не сомневаюсь, --- я не прощу его себе никогда.
Пусть в кувшинах твоих снов будут только вина и ароматные масла, Ликхмас.
Завтра жду тебя, как обычно.

Конфетка тепло улыбнулся и пропал в темноте под непрерывную любовную симфонию сверчкового оркестра.
А я стоял, придерживая рукой открытые двери, и чувствовал, что ужасное чувство ушло вместе с тренером.

Всё встало на свои места.

\section{[-] Метритхис --- игра учёных}

%Змея 4 дождя 12001, Год Церемонии 20.

--- Меркхалон-кровопийца, ты только посмотри, как нам Фатум поле выложил, --- недовольно буркнула Кхохо.
--- Ну и как играть в таких условиях?

--- А мне нравится, --- улыбнулась Ликхэ.
--- Для <<мышей>> идеально.

--- А я вот не умею на <<мышах>> играть.

\textspace

--- Ликхэ, ты рыбина, --- проворчал Ситрис.
--- Нельзя так, нехорошо.

--- Тхэай-кхвай, что же будет делать тактик Тхитронского Храма? --- осклабилась Кхохо.
--- Твоему сброду против троих Воинов не выстоять.
Разбойник сбежит, точно тебе говорю.
А Мститель даже на полной мотивации положит от силы двоих.

Ситрис подмигнул мне.

--- А теперь, Ликхмас, одна из тайн.

Ситрис пододвинул Разбойника к Мстителю.
В следующий момент обе фигурки отправились в ящик и были заменены Воинами.

--- Хэ! --- возмутилась Ликхэ.
--- Не жульничать!

--- Что произошло? --- удивился я.
--- В правилах же сказано, что равный обмен возможен только...

--- Только между антагонистами, но необязательно полными антагонистами, --- объяснил Ситрис.
--- Эти классы --- антагонисты по мотивации при прочих равных характеристиках.
Разбойник показал Мстителю смехотворность его мотивов.
Мститель показал Разбойнику, что такое цель и для чего она нужна.

Ситрис посмотрел на Кхохо.
Женщина, казалось, ушла в себя --- весьма необычное для неё состояние.

--- Складно поёшь, Ситрис.
А теперь ставь фигуры обратно, --- буркнула Ликхэ.

--- С чего бы?

--- У Разбойника плюс два к беседе.
У Мстителя минус два.
В итоге ноль.
Разговора не было, --- настаивала Ликхэ.

--- Кинем? --- прищурился Ситрис.

--- Сколько?

--- Восемнадцать.

--- Ты чокнулся? --- вытаращила глаза Кхохо.

Ситрис бросил кихотр.

--- Вот ты ж, --- угрюмо шмыгнула Ликхэ.

--- Воина с Воином рассудит кихотр, --- заключил я.
--- Трое против двоих --- вполне себе честный бой.

--- Только бросать буду я, а не этот жулик, --- буркнула Ликхэ.

--- Ещё чего! --- возмутился Ситрис.
--- Нет уж, пусть Фатум бросает.
Фигуры вплотную, а сейчас его ход.

--- Не подлизывайся к Судьбе, Мятежник, --- захихикала Кхохо.
--- Тебе и так везёт до неприличного часто.

\textspace

Ликхэ, огорчившись последнему раскладу, вначале подчистую продула своих Воинов ситрисовским, а затем получила Заговор от Кхохо на столе дипломатии.
Ликхэ попыталась выбросить <<побег>>, но удачи ей явно недоставало --- девушка вышла из игры насовсем и стала преувеличенно жизнерадостным тоном комментировать ходы.

\textspace

Кхохо изящно поставила чёрный камень на доску и, взяв ещё один, бросила мне точно в нос.

--- Мир, --- провозгласила Ликхэ.

--- Крепись, Ликхмас ар’Люм, --- осклабилась Кхохо.
--- В следующий раз ты обязательно...

--- Сейчас его ход, --- напомнил Ситрис.

Я не глядя бросил кихотр.
Восемь.
Запустил руку в мешочек с Благами и Несчастьями и швырнул камень воительнице.
Чума.
Кхохо вытаращила глаза.

--- Вот не приведи лес тебе выбросить...

Я подбросил кихотр ещё раз.

--- ... восемнадцать.

Кхохо задумчиво смотрела на обращённую к ней цифру.

--- Смерть правителя, --- так же радостно провозгласила Ликхэ.

--- Я тебе голову оторву, --- спокойно пообещала Кхохо.
В следующую секунду доска и камни разом полетели на меня, словно стая взбесившихся птиц.

... В себя я пришёл только после громкого крика: <<Ликхмас ар’Люм!>>
Я сидел верхом на Кхохо и возил её лицом по каменному полу.
В следующую секунду пришла тупая сильная боль в паху, и я свалился камешком на пол.

Ситрис и Трукхвал растащили нас с Кхохо в разные стороны.
Впрочем, тащить пришлось только меня по той простой причине, что идти я не мог --- внизу живота разливалась боль, а правая нога не желала разгибаться.
Кхохо же, едва я оказался вне пределов досягаемости, начала деловито проверять на ощупь форму носа.

--- Ликхмас!
Поднимайся!
Где болит?

--- А, сейчас, --- сообразил Ситрис и, подойдя ко мне, аккуратно сжал мякоть бедра.
Щелчок --- и нога снова начала двигаться.

\ml{$0$}
{--- <<Замочек на ножку --- поползай на брюшке>>, --- пояснил он.}
{\emph{``Your paw locks so, you crawl too low,''} he explained.}
\ml{$0$}
{--- Пах болит?}
{``Groin is aching, isn't it?}
\ml{$0$}
{Это у Кхохо любимая мишень, она может ударить в пах любой частью тела.}
{\Kchoho's favourite target, she can attack the groin by each part of her body.}
\ml{$0$}
{Любой.}
{Each one, boy.}
\ml{$0$}
{Особенно обожает отсроченные удары.}
{She adores delayed hits most of all.}
\ml{$0$}
{У меня на случай ссоры с ней есть золотые щитки для...}
{In case of disagreements I wear a little golden casing on my---''}

--- Ситрис! --- рявкнул Трукхвал, и воин смущённо замолчал.
--- Приведи молодого жреца в порядок, и пусть он поднимется ко мне.

Трукхвал заковылял на второй этаж.
Таким рассерженным я не видел его никогда.

\section{[-] Гнев учителя}

--- Ликхмас, ты --- жрец!
Жрец, а не воин!

Трукхвал не имел большого опыта в упрёках и потому всегда ругался без предисловий.
Я едва успел переступить порог библиотеки.

--- Это воинам позволительно мять друг другу бока, когда вздумается!
Для них это тренировка!
Жрец должен держать себя достойно!

\ml{$0$}
{--- Это не я начал...}
{``I'm not the one who started---''}

\ml{$0$}
{--- Азартные игры --- недостойное жреца времяпрепровождение! --- Трукхвала, похоже, было уже не остановить.}
{``Gambling is an amusement that's unsuitable for a priest!'' \Trukchual\ seemed to be unstoppable.}

Я хмуро слушал, классифицируя ругань по затронутым темам.
В точности так, как сам Трукхвал меня учил.
Драка, азартные игры, долг перед городом, честь, драка...
Повторяется.

--- Ты проводишь на нижнем этаже чересчур много времени!
Я этим займусь! --- в сердцах закончил свою тираду учитель.

Трукхвал устало сел за стол и глотнул остывшего отвара из чаши.
Я дал ему отдышаться.

--- Учитель, --- осторожно начал я.
--- Что плохого в том, что я...

--- Ликхмас, --- уже спокойнее заговорил Трукхвал.
--- Мы --- хранители мудрости и отправители обрядов.
Мы приносим жертвы Безумным богам.
Серьёзность этого очень слабо сочетается с тем, что я видел сейчас.

\ml{$0$}
{--- Кому какая разница, как я провожу свободное время, если я хорошо выполняю свою работу?}
{``Who cares about my amusement if I do my work well?''}

\ml{$0$}
{--- Мне --- никакой, --- согласился Трукхвал.}
{``I don't,'' \Trukchual\ admitted.}
\ml{$0$}
{--- И прочим в Храме тоже.}
{``All the Temple don't.}
Но многие крестьяне искренне верят, что у нас есть власть над словом, духами и силами природы, и разрушить эту веру --- значит лишить их покоя.
Представь, что подумал бы человек с улицы, если бы зашёл в тот момент.
Ему могла потребоваться помощь жреца в лечении или отправлении обрядов, и лично я не доверился бы мальчишке, который михнет назад кормил Кхохо паркетиной.

--- И что, мне вообще нельзя отды...

--- Жрец и отдыхает достойно, --- перебил меня Трукхвал.
--- А теперь иди и сшей двести школьных тетрадей по шестьдесят листов.
Вот тебе отдых.

--- Учитель! --- возмутился я.
--- Мы уже равные, между прочим!

--- Триста, --- спокойно парировал тот.
--- Пока не сделаешь, никакого сна.

\section{[-] Вездесущая месть}

От работы учитель освободил меня уже через кхамит.
Работа едва перевалила за двенадцатую тетрадь.
Я так и не понял, чем был вызван странный приступ гнева, а учитель был, похоже, чересчур смущён собственной строгостью, чтобы распространяться на эту тему.

--- Если соберёшься... хай... домой, то отнеси кормилице вот это, --- пробормотал Трукхвал, сунул мне в руку смятую запечатанную записку и ретировался.

С воинами я столкнулся на лестнице.
Ликхэ едва успела отбить летящий в меня факел (к счастью, незажжённый) и подхватить Кхохо под мышки.

--- Всё-всё, --- примирительно сказала Кхохо и, подняв растопыренные ладошки, пошла по своим делам.

--- Я бы на твоём месте переночевал дома, --- философски заметил Ситрис, похлопав меня по плечу.
--- Утром приходи, она к тому времени остынет.

Я вздохнул.
Дом был на другой стороне площади, но тащиться не хотелось.
Вручённая мимоходом записка не казалась весомым аргументом.

--- Даже не думай, --- опередил меня Ситрис.
--- Спать она тебе не даст, хоть в крипте ложись, хоть на улице Стриженого Кактуса.
А вот твою кормилицу она побаивается и к вам в дом не полезет.
Если Кхохо бузит, лучше сразу бежать --- проблем меньше.

--- А почему Кхохо побаивается кормилицу?

--- Кхотлам один раз наложила на Кхохо наказание, как и на меня, чтоб на всю жизнь запомнила.

--- Боюсь спросить, какое.
Она вообще чего-нибудь боится?

Ситрис ухмыльнулся.

--- С тобой мелким сидеть.
Если я салфетки полгода вышивал, держался, то Кхохо взвыла через декаду.
Она детей не выносит.

--- И за что?

\ml{$0$}
{--- Осквернила святилище.}
{``She violated sanctuary.''}

\ml{$0$}
{--- Мне говорили, что за это казнят.}
{``I was told it's punished by death.''}

\ml{$0$}
{--- Разумеется, казнят, если только за твоё спасение не возьмётся Кхотлам ар'Люм.}
{``Indeed, if there is no \Kchotlam\ ar'\Loem\ to plan your rescue.}
У Тхитрона после этого на пять дождей отношения с Одиноким Столбом испортились, и это мы ещё легко отделались, я скажу...

\ml{$0$}
{--- Кхохо должна быть благодарна, я думаю.}
{``\Kchoho\ should be thankful, I think.''}

\ml{$0$}
{--- Так думают все, кроме самой Кхохо.}
{``Everyone think so, except \Kchoho\ herself.''}
Кстати, ты Чханэ не видел?
Передай ей, что я её жду возле ткачей на обход.
Третий раз один патрулирую, надоело уже.

Воин вздохнул и скрылся за поворотом.

\section{[@] Беременность}

Плановое обследование преподнесло сюрприз.

--- Небо, ты беременный? --- удивился Костёр, только собравшийся снимать мультитомограмму.
--- Почему не сказал?

Я решил, что он шутит, и засмеялся.

--- Да какие шутки? --- возмутился Костёр.
--- Я чуть <<взрослую>> дозу контраста тебе не дал!

--- То есть ты хочешь сказать, что я бы не заметил стигм беременности, если бы они у меня были?

--- Так они и есть.
Посмотри в зеркало.
Хорошие такие стигмы, недельные.

Я посмотрел.
Да, всё правильно, шея полосатая.
Как говорит Баночка, <<тигров ловить собрался>>.
Как я не заметил недельные стигмы беременности, так и осталось загадкой.

--- Кто отец?

--- Я не знаю.
У меня ни с кем не было отношений ни во время, ни после перелёта.

--- Может, партеногенез, --- предположил Костёр.
--- Хотя с чего бы, апиды в поселении есть, концентрация феромонов в воздухе достаточная...
Так, сиди смирно, я посмотрю.
Небо, милый мой, я понимаю, что ты волнуешься.
Расслабься и подумай о чём-нибудь приятном.

Щёлкнул биопсийный пистолет.
Результаты пришли сразу --- отцом оказался Комар.

--- Надо же, --- удивился врач.
--- Видимо, из-за стресса зародыши впали в анабиоз.
Сиди, сиди.
Двое, развитие идёт удовлетворительно, пороков нет.
Похоже, Небо, что твои дети из титана --- не всякий зародыш переживёт космический полёт.

Я удивился, насколько мягче стал голос врача.
Он по-прежнему оставался резким и сухим мужчиной-человеком средних лет, но старый инстинкт сработал безотказно --- с беременными, детьми и больными разговаривать следует предельно осмотрительно.

Костёр дал мне зелёную капсулу и стакан воды.
Я непонимающе уставился на него.

--- За препаратами и питанием будешь приходить каждый день.
Инъекционных обойм у меня не осталось.

--- А томограмма?

--- Через десять-двенадцать дней, на этих сроках контраст может дать повышение гестотропина и резкий иммунный сдвиг, это плохо для личиночной нервной системы апида.

Я попытался возразить, что есть дела важнее вынашивания, но бесполезно.
Костёр был мягок, но непреклонен.

--- Небо, нас очень мало.
Прерывать здоровую беременность в нынешней ситуации --- преступление.
Прочие врачи скажут тебе то же самое.

--- Но это моё тело!

--- И твой вид, от которого осталась жалкая горстка.
Вас шестьдесят четыре, ты понимаешь?
Шестьдесят четыре.
Тем более это дети погибшего апида, сильной неродственной особи, свежие гены в популяцию.
Дети Комара, в конце концов.
Это для тебя уже ничего не значит?

Я промолчал.

--- Не надо на меня так смотреть.
В крайнем случае отдашь личинок другим апидам.
Шмель с удовольствием возьмёт, ему и так все детей спихивают.

\section{[@] Лучший воспитатель}

Для Шмеля разбили отдельную палатку.
Самое интересное, что детей-апид здесь почти не было --- много человеческих, щенята и девочка-плант, которая с важным видом, как взрослая, носила щенка-младенца.
Шмель каким-то невероятным образом успевал их умывать, развлекать и даже учить.
Недостатка внимания не было заметно --- дети выглядели вполне счастливыми.

--- Шестерёнка! --- Шмель подбежал к девочке-планту, погладил её по голове и выхватил у неё щенка.
--- Давай мне Косточку и иди покушай.
Привет, Небо, подожди.
Заря, смотри, как я умею.

Шмель подхватил другой рукой начавшего кукситься человеческого малыша и начал щекотать его ногочелюстями.
Ребёнок засмеялся.

--- Я не смогу взять у тебя личинок, Небо, --- бросил Шмель в сторону.
--- Ууу, какие у меня усики, смотри!
А это что, носик?..
Если я буду постоянно отвлекаться на кормление, то могу сразу забыть о сне и пище.
От детей отпуск не возьмёшь.
Как выкуклятся --- можешь приносить, но не раньше.
Шестерёнка, девочка моя, вытри ротик Камешку, пожалуйста...

Я кивнул и направился к выходу.

--- Небо, стой, зараза усатая.
Помоги покормить и поиграй с ними немного, что ли.
Всё одно для тебя тренировка.
Кто пришёл на новоселье, дарит малышам веселье!
Да, дети?

Малыши загомонили и попрыгали мне на руки.
Я чуть не упал.

--- Небо, ну пожалуйста, давай...
Что за день, я с ума сойду...
Почему вы так медленно растёте, а, малышня человеческая?

\section{[-] Достойные}

Записка, которую я принёс с собой, оказалась не из приятных.
Кормилица, едва сорвав печать, пошла пятнами.

--- Что этот старик о себе возомнил! --- гневно фыркнула она.

Я решил воздержаться от вопросов --- Кхотлам впервые отозвалась о Трукхвале в таком тоне.
Кормилица тем временем несколько раз перечитала написанное, словно искала скрытый смысл.

--- Будешь ты мне... шантажировать... --- бурчала она.
--- Так и передай своему учителю, чтобы он знал своё место.

--- На записке стоит печать <<Обменять на ответ>>, --- возмутился я.
--- Ты сама учила меня основам дипломатии.
Можешь отправить чистый пергамент, можешь нарисовать рожицу, но устно я от тебя ничего передавать не буду.

У кормилицы появилось на лице выражение, которое я не видел прежде --- она словно сожалела до глубины души, что когда-то обучила Ликхмаса ар'Люм основам дипломатии.
Впрочем, держалось выражение едва ли секхар.

--- Ты абсолютно прав, Лисёнок, --- отсутствующим тоном сказала наконец кормилица, открыла ларец с бумагами, вынула кусочек кожи и отдала мне.

--- Это Трукхвалу? --- уточнил я.

--- Нет, Лисёнок.
Это тебе, --- Кхотлам погладила меня по голове и ушла к себе в спальню.

Мне?
Что это ещё такое?

Я осмотрел и обнюхал дар.
Кожа выглядела и пахла очень старой.
Кроме того, я не смог определить зверя, которому она принадлежала.
К горлу подкатил ком, у меня появилось странное предчувствие, не дурное, но и не хорошее.

Я подошёл к окну и вгляделся в написанное.

\begin{quote}
<<Ликхмас, милая моя Лисичка!
Если ты читаешь это, меня уже нет в живых.
Я, твоя дарительница, Митхэ ар’Кахр э’Тхартхаахитр, оставила тебя у лучшей подруги, которую ты называешь кормилицей.
Ты был желанным ребёнком, и я бы пошла против всего мира ради того, чтобы растить тебя.
Но мой любимый мужчина был вероломно отдан в рабство.
Я не могла жить, не зная его судьбы.
Нет ничего ужаснее, чем есть, пить и веселиться, когда близкий человек страдает.
Не надеюсь на то, что ты поймёшь, но надеюсь, что ты вырос здоровым, красивым и счастливым человеком --- мужчиной или женщиной>>.\\~\\
\end{quote}

В самом низу была малозаметная приписка:

\begin{quote}
<<Для Кхотлам: пусть будет кем угодно, но только не воином.
Моя сабля заплатила народу сели на десять поколений вперёд>>.
\end{quote}

Иероглифы за долгие годы стёрлись, и было видно, что чья-то заботливая рука изредка их подновляла.

--- Лис?

Чханэ зашла в комнату и тут же с беспокойством подбежала ко мне.

--- Лис, ты как будто жабу проглотил!

Я молча протянул ей записку.
Чханэ, читая, постепенно менялась в лице.

--- Митхэ ар’Кахр дала тебе жизнь? --- Чханэ ошарашенно тряхнула головой.
--- Кихотр...
Нам про неё столько историй рассказывали, она у нас навроде Маликха была, героем легенд.
Обращала в бегство целую армию отрядом в десять человек, убивала врагов листом юкки...
И она была Плачущим Ягуаром, едва ли не единственным на всей Короне...
Хай...

Чханэ вернула мне кусочек кожи, и я сжал его в руках.
Не то чтобы мир изменился, но я многое увидел в новом свете: недосказанности, странные приступы нежности кормилицы, аура почтения, которая непонятно почему преследовала меня с первого дня в Храме, старания Конфетки сделать из меня воина и ежедневный традиционный гостинец, который тренер клал мне в карман...

--- Лис, я понимаю твои чувства, но тебе определённо повезло с семьёй, --- Чханэ неуклюже попыталась утешить меня.
--- У тебя достойные кормильцы, и женщина, которая тебя выносила, тоже была достойной, это несомненно.

Я промолчал.

--- Жаль, что я плохо знаю твоих домашних.
Расскажи о них, --- Чханэ всеми силами пыталась меня отвлечь.
Мне вдруг стало её невыносимо жаль --- действительно, кроме меня, больше у Чханэ никого не осталось.

--- Ха-ай, --- я улыбнулся как можно теплее, --- Эрхэ --- хранительница друга Кхотлам, он попросил позаботиться о ней.
Ликхэ --- её женщина...

--- Это я, как ни странно, и сама поняла, --- весело сказала Чханэ.
--- У них есть дом?

--- У Ликхэ есть.
Она передала его брату, у которого много детей.
К нам она попросилась, потому что здесь тихо.
Да и Эрхэ рядом.

--- А Сиртху-лехэ?

--- У Сиртху-лехэ нет ни дома, ни родных, ни друзей.
Кормилица подобрала его в сожжённом идолами посёлке, больше никто не выжил.
Он был настолько убит горем, что отказывался есть и требовал хэситр, но Кхотлам его выходила.
Старичок оказался смышлёным и бойким --- пожил немного на правах гостя, оправился и взял хозяйство на себя.
Сейчас он ведёт у нас учёт товаров.
А ещё он чудесно рассказывает сказки.

--- А тебе он рассказывал?

--- Рассказывал.
Манэ и Лимнэ вообще выросли на его сказках.

Чханэ улыбнулась.

--- Хороший у тебя дом.
Мужчина, женщина, два их хранителя, две цветущих\FM, ты-приёмыш и одинокий старичок.
\FA{
Цветущими сели называли людей, практикующих исключительно гомосексуальные связи.
Исключительно гетеросексуальных людей называли плодоносящими, асексуальных --- вечнозелёными.
}
Напоминает старые дома, в которых не делали различий между родичами и приёмными.
Сейчас звериное время, люди заботятся только о своих...

--- А то и их бросают, --- ввернул я.

Чханэ смутилась, но возразила:

--- Если ты про Митхэ, то она сделала всё так, как велит обычай.
Она могла тебя бросить в лесу, но отдала в самые надёжные руки, которые только знала.
Я видела то, что могло тебя ждать, Лис --- детей, умерших от голода в городе.
Детей, которые умерли от голода, Обнимающий Сит, где ты был?
Больше я этого не видела и надеюсь не увидеть никогда.

Я кивнул и засунул записку родильницы в карман.
Пусть подождёт своего часа...

--- И стариков бросают, --- задумчиво добавила Чханэ.
--- Если уж совсем тяжко с едой, даже убивать куда гуманнее.
У нас были случаи --- стариков и больных жрецы поили лаковым соком, нашествие было, поветрие.
Но бросать!..
В этом мире происходит что-то страшное, Лис.
И хуже всего то, что я происходящее не могу понять.

--- Ситрис ждёт тебя у ткачей, --- напомнил я.

--- Я помню.
Если хочешь, могу посидеть с тобой.

--- Сначала Храм, потом всё остальное, --- улыбнулся я и хлопнул подругу по плечу.
--- Иди, я справлюсь.

Чханэ кивнула и, подхватив заплечный мешок, вышла в надвигающиеся сумерки.

\section{[@] Норма реакции}

\epigraph
{Мыслящие никогда не признают, что свободы воли не существует и что они являются сложными, но всё же механизмами.
Это признание разрушит всё, на чём стоит их маленький уютный мирок --- заслуги и преступления, дружбу и вражду, религию, власть и суд.
Они будут цепляться за идею непознаваемого и до последнего отрицать своё родство с обычной ветряной мельницей, не понимая, что это родство не унижает, а возвеличивает их.}
{Леам эб-Салах.
Письмо к Анатолиу Тиу, случайно попавшее в сборник <<Записи под молитвенным ковром>>}

Костёр пришёл ко мне в тот же вечер.
Мы с Заяц сидели на песке и любовались звёздами.
Где-то вдалеке плескались и ворковали два дельфина.

--- Это не дело, --- сказала Заяц.
--- Надо сказать дарителям, чтобы большую часть времени дети проводили с ними.
Пусть хоть на работу таскают, там малыши без внимания не останутся, взрослых много.
Шмель --- лучший воспитатель, но он так долго не выдержит.
Да, врачей у нас тоже мало, но дети требуют больше внимания, чем тяжелобольные!

--- Они обязательно найдут отговорки, --- заметил я.
--- Все заняты...

--- Так надо им напомнить, в чьи руки их занятия попадут после, --- резко бросила Заяц.
--- Если дети вырастут несчастными, любая работа потеряет смысл.
Привет, Костёр.

--- Привет, Зайчик, --- вздохнул мужчина.
--- Небо, извини.

--- Перестань, --- отмахнулся я.
--- Ты всё сделал правильно.

--- Нет, --- возразил врач.
--- Я посягнул на твою свободу.
Это не есть правильно.

--- Ради жизни, --- пожал я плечами.
--- Апид действительно осталось мало.

Костёр помолчал.

--- Зря ты применил эту фразу тогда, --- наконец проговорил он.
--- Ею можно оправдать самые страшные поступки.
Я под её прикрытием нарушил твоё право на свободный выбор.

--- И не только ты, --- вмешалась Заяц.
--- С первого дня, когда мы вступили на эту планету, мы только и делаем. что нарушаем принципы.
Разве вы не видите?
Мораль тси трещит по швам.
Она не была готова к Катаклизму.
Почему предки не предусмотрели такой возможности?
Да ладно предки.
Почему никто из нас, выросших под угрозой вторжения демонов, не озаботился устойчивостью морального кодекса?
И вообще, что с нами происходит?
Откуда вся эта грязь?

--- Она в нас была всегда, --- проворчал Костёр.
--- У каждого параметра, в том числе и поведенческого, есть норма реакции.
Мы жили счастливо и хорошо, вот все и стали добряками.
Как пришла война --- оп-па! --- мы опустились на нижнюю границу нормы реакции.

Костёр сделал выразительное движение ладонью вниз.

--- Ведь это же... --- робко начала Заяц.

--- Ещё ни один моральный кодекс в истории не выдержал испытания войной, --- перебил её Костёр, внушительно разделяя слова.
--- Война показывает одно из двух: либо кодекс негуманен, либо нежизнеспособен.
Где вы найдёте кодекс, в котором освещались бы <<тёмные стороны>> личности?
Взять ту же жестокость.
Что сказано в кодексе тси о жестокости, кроме того, что она приемлема лишь в игровой форме?

--- Кто-то сублимирует жестокость в спорте, в сексуальных играх... --- возразила Заяц.

--- Этого недостаточно, Зайчик! --- громко заявил врач.
--- Почему за все тысячелетия никто не разложил по полочкам <<тёмную сторону>> тси?
И не говорите, что попытки предпринимались.
Да, были исследования, но каждый раз их результаты искажались, сглаживались или стыдливо замалчивались, а о влиянии полученных результатов на воспитание молодёжи и речи не шло!
Почему никто не адаптировал кодекс к ситуации, когда личность скатывается к диаметрально противоположной границе нормы реакции?

--- И как бы ты построил счастливое общество, если бы принял как данность жестокость, вероломство и трусость? --- скептически спросила Заяц.

--- Я не знаю! --- развёл руками Костёр.
--- Но если этого не сделать, то любое общество будет ждать такой же конец.

Мы с Заяц поёжились.
Костёр говорил страшные, но, похоже, правдивые вещи.
У меня возникло чувство, что эта идея возникла у него задолго до Катаклизма.
Возможно, он что-то увидел в больницах.
Что-то натолкнуло его на подобные мысли именно там, куда сапиенты поступали вышедшими из строя, где из-за нарушения функции мозга наружу проступали самые древние, звериные качества.

--- Природа --- суровая дарительница, --- глухим голосом сказал Костёр.
--- Вот взять даже вас, апид.
Ты ведь помнишь, сколько оборотов прошло со времени Войны Тараканов?
Равноправие победило, колониальные апиды Ди были истреблены.
Но вы, апиды-тси, до сих пор восприимчивы к тем, древним гормонам.

Я промолчал.
Да, разумеется, об этом знали все.
Даже сейчас, после тысяч итераций генотерапии, десять процентов апид --- и я в том числе --- до сих пор были способны к дифференцировке.
Синтетические гормоны апид Ди могли превратить любого из нас в Матку.
Или в Солдата --- бездушную машину для убийства, ужас, оставшийся где-то на страницах древней истории.

Почему тси не избавились от этого механизма?
Апиды высказались против --- в последний раз не ранее чем сто оборотов назад.
Наша особенность несколько снижала плодовитость, но делала нас сильнее физически и гораздо привлекательнее для прочих.
В медицинской карте у меня, как и у прочих тринадцати миллионов таких же, стоял красивый иероглиф-лигатура СД --- <<способность к дифференцировке>>.
Значок тревоги, подразумевающий осторожность при использовании новых и некоторых старых препаратов.

--- Говорят, что Солдаты безоговорочно исполняли все приказы Матки, что этому нельзя было противиться.

--- Солдат готовили идеологи, --- разко бросила Заяц.
--- Не всё решается одними лишь гормонами.
Хотя я вообще слышала мнение, что Матка была такой же заложницей Солдат, как и прочие апиды.
Они использовали её как инкубатор и говорили от её имени...

--- Хотел бы я, чтобы это было правдой, --- пробормотал Костёр.

--- Предки всё это знали, --- настаивала Заяц.

--- Знали и не оставили предупреждений? --- с горечью сказал я.
--- Нам никогда не говорили: <<Если вы поймёте, что наше учение неприменимо, отойдите от него>>.
Нет, нам говорили: <<Даже если цивилизация будет на грани гибели, вы должны быть истинными тси>>.
И мы стараемся.
Страдаем, гибнем, но стараемся быть истинными тси.
И ни у кого из тех, кто принёс себя в жертву на Золотой дороге, даже не было мысли, что можно поступить иначе.
Сейчас я завидую Комару --- он умер, не зная этих ужасных сомнений.

--- И ведь все погибли добровольно, --- печально сказала Заяц.
--- Ты сказал правду: кто-то боялся, но никто не сомневался в том, что это правильно.
А получается, что эти жизни отданы не ради лучшего, а ради того... ради...

--- Ради того, чтобы всё шло так, как идёт, --- закончил за подругу Костёр.
--- Мы всю жизнь работаем именно ради этого.

Заяц промолчала и слегка уничижительно посмотрела в его сторону.
Врач грустно улыбнулся и встал, собираясь уйти.

--- Прервать беременность ещё не поздно, --- наклонился он ко мне.
--- Приходи завтра.

Я помотал головой.

--- Почему?

--- Потому что ты меня убедил, --- объяснил я.
--- Свобода --- это инструмент, а не цель.

Врач хохотнул и погладил меня по голове.

--- Не забудь принять лекарства, командир.

\section{[-] Расплата за молчание}

\textspace

Почему она мне ничего не сказала?

Пылая гневом, я пошёл в спальню кормилицы.

Дверь оказалась незапертой.
Кхотлам крепко спала на подоконнике, положив кудрявую голову на руки.
Странная записка выпала из её руки.
Я развернул её и поднёс к узкому клинку предзакатного солнца, пробившего тонкую занавеску.

\begin{quote}
<<Его тянет к воинам.
Он чувствует, что его с ними связывает нечто глубинное, недомолвки и ложь только усугубляют ситуацию.
Если Ликхмас сорвётся, дело может принять плохой оборот.
Ты знаешь, перед кем мы в ответе.
Мне неважно, по какой причине ты до сих пор не сообщила ему правду.
\ml{$0$}
{Если ты не сделаешь это сегодня, это сделаю завтра я>>.}
{If you fail today, I shall be the one to do it tomorrow.}
\end{quote}

Под запиской не стояла подпись, но я узнал бы почерк Трукхвала, даже не будь я его лучшим учеником.

Я посмотрел на Кхотлам.
То ли она притворялась спящей, то ли действительно спала.
Впрочем, значения это не имело.
Обычная реакция человека, которому страшно боится неотвратимого.
Кормилица ждала наказания, словно провинившийся ребёнок.

Стащив с лежанки одеяло, я накрыл кормилицу и вышел, тихо прикрыв за собой дверь.
Причину я узнаю после.
Расплатой для неё стали прошедшие двадцать михнет.

\section{[-] Лицо в траве}

\textspace

--- Но ты же кормила меня грудью, как такое может быть?

--- Кормила, Лисёнок.
Я приложила тебя к груди и появилось молоко.
Потом, когда я уехала по делам, тебя кормил ещё и Хитрам.

Кормилец согласно кивнул.

--- У меня на несколько лет выросла грудь, --- сообщил он.
--- Да ещё такая красивая.
Пожалуй, самое странное, что со мной происходило в жизни...
Меня облапали все, кому не лень, да, Пёрышко?

--- Не болтай, --- кормилица покраснела, и Хитраму прямо в лоб прилетело мокрое полотенце.
Эрхэ и Ликхэ захихикали.

--- Как она выглядела? --- продолжал допытываться я.
--- Я похож на неё?

\ml{$0$}
{--- Наверное, --- уклончиво сказала Кхотлам.}
{``May be,'' \Kchotlam\ equivocated.}

--- Кормилица!
Скажи её хасетрасем!

Кхотлам потупилась.

\ml{$0$}
{--- Лисёнок, я не помню.}
{``Baby Fox, I can't remember.''}

\ml{$0$}
{--- Не ври.}
{``Don't tell a lie.''}

\ml{$0$}
{--- Я не вру.}
{``I don't.}
\ml{$0$}
{Много лет прошло.}
{Many years have passed.''}

То же самое говорили мне и прочие люди, знавшие Митхэ ар’Кахр.
\ml{$0$}
{Все при этом отводили глаза.}
{All of them averted their eyes while talking.}

--- Ликхмас, она пропала тридцать дождей назад где-то на востоке, --- попытался урезонить меня Ситрис.
--- Вряд ли кто-то помнит, как выглядела Митхэ.
Да и лицо, которое ты хочешь увидеть, уже давно поросло зеленью.
Ты же читал?..
Хай, забудь.

Ситрис тоже знал о записке.
О ней знали все, кроме меня.

Впрочем, жизнь текла своим чередом, подобно Ху'тресоааса, могучей полнокровной жиле на теле Тра-Ренкхаля.
Она текла от горизонта к горизонту, и самому острому взгляду кормчего было доступно считанное число извивов.
Жизнь текла по Вселенной, принимая самые неожиданные формы, и хоть законы её течения давно познаны и рассчитаны, наблюдатель никогда не увидит в мутных водах всё многообразие.
Лишь иногда любопытная чаша зачерпнёт воды с чудесной рыбкой, глаза полюбуются на прихотливую игру чешуи, и рыбка отправится обратно за борт, так и не открыв наблюдателю истины --- чудес в чаше было больше одного.

Прошёл дождь, прошёл другой, и меня захватили дела.
В конце концов, что изменилось?
Кормильцы остались для меня такими, какими и были.
Кхотлам всё так же называла меня Лисёнком, Хитрам с улыбкой рассказывал про дарительницу интересные истории, которые до сих пор кочевали в народе, и мало-помалу Митхэ ар'Кахр превратилась для меня в легенду, какой она была для остальных.
А что до крови...
Однажды ночью я понял: немного надо, чтобы отдать кому-то каплю своей крови.
Я в последний раз призвал на Митхэ благословение лесных духов, заснул спокойным сном --- и следующее утро стало первым, когда я о ней не вспомнил.

\chapter*{Интерлюдия V. Постоялый двор <<Бамбуковая клетка>>}
\addcontentsline{toc}{chapter}{Интерлюдия V. Постоялый двор <<Бамбуковая клетка>>}

\textbf{Байки Дальнего Севера.
Апокриф <<Легенды об обретении>>}

Было это в далёкие времена, когда Безумный свирепствовал и брал жертвы не только на алтаре, но и везде, куда дотягивалась его рука, и ни духи, ни дышащие не смели сказать ему ни слова.

Многие города сгорели в беспощадном пламени Безумного, многие люди сошли с ума под ужасом многоцветья.
Особенно ненавидел Безумный мастеров своего дела и книжников, ибо они учили прочих избегать несчастий.
Ненавидел он и тех, кто знал его под именем Безумного и писал это имя ореховыми чернилами и сурьмой.
Много книг сгорело в огне, и забывать стали люди о значении старых иероглифов, а уж писать и вовсе разучились.
Духи ещё помнили старые языки, но Безумный наложил на них страшное проклятие --- никто из духов отныне не смел учить языкам живых.

Было это так.
Собрались однажды духи на совет.
Немногие пришли --- в камень обратились те, что потеряли цель и смысл;
побледнели и прозрачными стали, как воздух, что были подзабыты живыми.
И сказал им Безумный:

<<Вечное молчание падёт на тех из вас, кто скажет одно слово на двух языках, кто скажет слово и напишет его иероглиф, кто скажет слово в ответ на иероглиф, кто напишет иероглиф слову в ответ.
И лишь в тот день, когда дух и живой разделят ложе, когда влюблённый презреет свои чувства, когда речные воды станут солонее морских, в руках дышащего окажутся пергамент, перо и чернила!>>

Хотел Безумный создать новый мир для себя, ибо нет такого разрушителя, который не почитает себя величайшим из творцов.

Но был среди духов старец;
он был бледный и прозрачный, словно знойная дымка, и было его слышно лишь тогда, когда все замолкали\FM.
\FA{
Забытый --- персонаж многих историй.
Ранние сели верили, что забытые потомками больше не существуют, и Забытый служил собирательным образом тех, чьи имена и деяния были утрачены.
В <<Легенде об обретении>> Забытый исчезает навсегда;
из-за этого Легенду критиковали, и даже были попытки изменить её.
Но автор, согласно <<Хроникам дорог и ветров>>, высказался ясно: <<Я рассказал то, что рассказал.
Вам придётся это принять>>.
}
Сказал старец:

<<Нет среди нас веры твоим угрозам;
священен лишь договор.
Поклянись же, что в городах и деревнях, в лагерях и становищах будет твоим наделом лишь алтарь, покуда народ исполняет твою волю>>.

Оскалился Безумный, и загрохотал гром, и засверкали молнии.
<<Клянусь>>, --- захохотал он и бросил кихотр на игральную доску.
И пропал старец насовсем, ибо выпало на кихотре Забвение.

Взмолились духи сестре Воде:

<<Сестра Вода, ты дала начало всякой твари, и даже в мёртвых костях есть твоя неумирающая суть.
Ты разбиваешь скалы и протекаешь сквозь гранит;
ты остужаешь кровь земли, и приносишь почву, и взращиваешь зелёные злаки на ужасной той крови.
Если ты не найдёшь дорогу в проклятии Безумного, никто не сможет её найти>>.

Приостановила сестра Вода бег двух рек\FM, и дожди отложила --- думала.
\FA{
<<Четыре засушливых года>>, связанные с усилением активности Тра-Ренкхальского солнца, по время которых значительно обмелели Ху'тресоааса и Ху'микхаса. 8117--8125, Годы Церемонии 1--4 (дата подтверждена лабораторией климатических изменений, отдел геологии и классической планетологии Ордена Преисподней, филиал Тра-Ренкхаля. Архив уровня 0, серийный номер 456300-15, тезис 84.).
}
Как учить, если не иероглифом к слову, если не словом на иероглиф?
И нашла-таки сестра Вода способ.
Осталось лишь найти достойного ученика.

И увидела как-то сестра Вода молодого рыбака, что чертил палочки у самого берега.
<<Вот он --- ученик>>, --- решила она.
Но, подойдя ближе, сестра Вода обнаружила, что мужчина красив и улыбчив, как Хри-соблазнитель.
Воспылала сестра Вода к мужчине страстью.
Ушли её мысли в сторону, да и сказка в сторону ушла.

Приняла сестра Вода облик женщины дивной красоты и вышла к рыбаку из родной стихии.

<<Что ты делаешь?>> --- спросила она.

<<Считаю рыбу>>, --- ответил рыбак.

<<Не желаешь ли ты научиться чтению и письму?>>

<<Зачем мне это?>>

<<Ты можешь обрести мудрость предков, чтобы ловить рыбу>>, --- сказала сестра Вода.

<<Чтобы ловить рыбу, мне достаточно мудрости моего наставника>>, --- ответил рыбак.

<<Ты можешь познать тайны звёзд, гор, морей и сельвы>>, --- сказала сестра Вода.

<<Мне хватает тайн реки>>, --- ответил рыбак.

<<Ты можешь вложить свою память в пергамент, и она не покинет тебя>>, --- сказала сестра Вода.

<<На память не жалуюсь>>, --- белозубо усмехнулся рыбак.

Совсем потеряла голову сестра Вода.

<<Я стану твоей, и буду с тобой столько, сколько ты захочешь>>, --- сказала она.

Расхохотался мужчина.

<<Мне хватает и тех женщин, что не измазаны в тине и не опутаны водорослями>>, --- сказал он, собрал свои снасти и ушёл.

Обиделась сестра Вода, вздулись воды Ху'тресоааса от её гнева.
Дождавшись, когда выйдет мужчина на рыбалку, подхватила она его утлую лодку, вынесла к Ихслантхару и разбила о Мраморный мыс\FM.
\FA{
Мраморный, или Кошачий, мыс --- опасная зона судоходства, место слияния Речек-близнецов, Со'хата и Со'хатра.
}
И рыбе наказала, чтоб держалась подальше от его сетей.

Пришёл рыбак на постоялый двор, что на перекрёстке трёх дорог.
Мокрый весь, грустный.
Шутка ли --- саму сестру Воду, не признав, оскорбил!

Подошла к рыбаку хозяйка двора.
Не была она так красива, как сестра Вода, но веяло от неё прелестью.

<<Принеси мне что-нибудь поесть>>, --- попросил её рыбак.

<<Всё, что подают в этом дворе, есть наверху>>, --- ответила хозяйка.

Взяла она мужчину за руку и повела наверх, где был коридор и пять комнат.

<<Только знай: ты не можешь отказаться ни от чего, что дают в этих комнатах>>, --- сказала хозяйка.

<<Я понял>>, --- сказал рыбак.

Завела хозяйка рыбака в первую комнату --- был там кузнечный горн с огненными лисами.
Вспрыгнули лисы на колени рыбака, потоптались, потёрлись --- и согрелся он, и высохла его одежда.

Завела хозяйка рыбака во вторую комнату --- там был стол, полный яств.
Поел рыбак от души, выпил вина --- но не исчезли яства со стола, и даже как будто прибыли.

Завела хозяйка рыбака во третью комнату --- там стояла цитра-самозвон.
Заиграла цитра, и утешился рыбак, и ушла его грусть.

Завела хозяйка рыбака в четвёртую комнату --- стоял там один чан с рыбьим жиром, и котёл стоял для выварки рыбы.
Понял рыбак, чем жить ему без лодки и удачи на воде.

Вывела рыбака хозяйка и указала ему дорогу наружу.

<<Подожди, --- сказал ей рыбак, --- покажи мне пятую комнату>>.

<<Но разве ты не получил всё, что хотел?
Ты выглядишь счастливым>>, --- удивилась женщина.

<<Я счастлив, --- сказал рыбак.
--- Но эти чудеса раздразнили меня, и я не успокоюсь, пока не увижу пятую комнату.
Что в ней?>>

<<В первой комнате --- твоя самая очевидная, насущная нужда.
В последней комнате --- то, что придаст твоей жизни новый смысл или, возможно, уничтожит тебя.
Я не знаю, что ты там найдёшь, а даже если узнаю, то не смогу до конца понять значение дара.
Это только для тебя>>.

<<Покажи мне последнюю комнату>>.

<<Помнишь ли ты, что не можешь отказаться от моих даров?>> --- спросила хозяйка.

<<Помню, --- ответил рыбак.
--- Все твои дары пришлись мне по душе.
Сомневаюсь, что я откажусь>>.

Завела хозяйка рыбака в пятую комнату.
Там были очаг, лежанка, фонарь, ароматное масло кокоса и чаши с водой.
Села хозяйка на край лежанки, позвала рыбака к себе, и провели они ночь вместе за беседами и любовью.

Выйдя же наутро из постоялого двора, забыл рыбак обо всём, что было.
Остались при нём лишь сухая одежда, сытость, хорошее расположение духа, мысли о новом его ремесле и свербящее чувство того, кто только что ушёл с ложа лучшей из женщин --- без надежды на возвращение.

Вернулся он в родную деревню и стал вываривать рыбий жир.
Но чувство недостающего не покидало его.
Знал рыбак, что влюблён, но не помнил имени;
видел рыбак женский образ --- то в вечерних тенях, то утром, на берегу пробуждения, --- но не видел лица;
слышал рыбак мягкий голос --- усладу и отдых для ушей, --- но не помнил ни слова.
Говорили в деревне, что угодил молодой мужчина в Беспамятный постоялый двор, что зовётся <<Бамбуковой клеткой>>.
Понял рыбак, что подвела его память;
вспомнил он о словах сестры Воды.

И вернулся рыбак к реке снова, а с ним и сказка вернулась.

Стал рыбак звать сестру Воду.
Она откликнулась на его зов, хоть всё ещё была немного сердита.

<<Зачем ты пришёл? --- спросила она.
--- Клянчить у меня рыбу?>>

<<Мне не нужен больше улов, я теперь вывариваю рыбий жир>>.

<<Тогда что тебе нужно?>> --- удивилась сестра Вода.

<<Я приму твои дары, --- сказал рыбак.
--- Твою любовь, искусство чтения и письма>>.

Расцвела сестра Вода в тот же миг, и запахла река ароматом ночных кувшинок.

<<Я не могу ответить иероглифом на слово, и словом на иероглиф не могу, --- сказала сестра Вода.
--- Но средство есть.
Возьми книгу, разруби иероглифы на кусочки и дай каждому кусочку имя.
Тогда я смогу тебя учить>>.

И разрубил рыбак иероглифы на ключи, и дал имена всем ключам, и стала сестра Вода рассказывать ему истории слов ключами.
Чувствовал Безумный, что неладное творится на земле, что трещит порядок его, но поделать ничего не мог.
Спросит Безумный у рыбы, а та разбежалась <<листом>> да <<хвостом>>, спросит Безумный у воды, а та <<змеёй>> да <<колодцем>> обернулась.
Неистовствовал Безумный, поднимал облака тьмы над северными горами, крушил твердь, метал молнии, сыпал пеплом с небес\FM, да так ничего и не узнал.
\FA{
Речь об извержении супервулкана Закованный Пик, Лук 8126 (дата подтверждена Лабораторией Катастроф, отдел геологии и классической планетологии Ордена Преисподней).
}
Не нарушила сестра Вода его наказа, но научила рыбака грамоте.

Прониклась глубоким чувством сестра Вода, пока учила мужчину читать.
Но проговорился мужчина, что учится читать и писать, чтобы разгадать тайну <<Бамбуковой клетки>>.
И поняла сестра Вода из его слов, что пылает он чувством к той, кого не может вспомнить.
Ревность взяла её.

<<Уходи, --- сказала она ему, разразившись плачем.
--- Забирай то, что я тебе дала, и больше меня не тревожь>>.

Поблагодарил рыбак сестру Воду, попрощался с ней и пошёл в Ихслантхар, к перекрёстку трёх дорог.

Хозяйка двора встретила его приветливо, словно незнакомца.
Взглянул рыбак на её лицо --- и понял, по кому страдал всё это время, но не сказал ей ни слова.
И лишь одна мысль была у него --- унести в своей памяти то, что он понял.

Взяла хозяйка мужчину за руку и повела наверх, где была всего одна комната.

Завела его хозяйка в комнату.
Там был стол с пером, чернильницей и свитком пергамента.
Остолбенела хозяйка;
ударила она в ярости по столу, и покатился письменный прибор по полу.

<<Что за люди! --- закричала женщина.
--- Я давала вам всё, что вы хотели --- тепло, пищу, любовь, новую жизнь, --- забирая лишь память о нескольких кхамит.
Но вы упорно пытаетесь вернуться за ней!
Хочешь свою память обратно?
Забирай!
Забирай вместе со всем, что я тебе дала, и больше меня не тревожь!>>

Отдала хозяйка перо, чернильницу и свитки рыбаку, вернула ему память и вывела наружу.

С тех пор постоялый двор пропал.
Но не бесследно --- желание помочь в хозяйке было сильнее всяческих обид.
Случалось ли вам видеть прекрасные сны, которые утекали из вашей памяти, едва вы проснулись?
Знайте --- это открылись двери постоялого двора у перекрёстка трёх дорог.
И так же, как и в былые времена, там есть комнаты для любых бед и проблем.
Однако у спящих не найдётся пергамента, чтобы всё записать, и тайна комнат хранится крепко до сих пор.

Смирился рыбак с тем, что не видать ему милой хозяйки <<Бамбуковой клетки>>;
прочёл он много книг, пытаясь забыть свою нечаянную страсть, и стал книжником, и переписал старые книги, и научил людей заново этому искусству.
А Безумный смирился, как смиряются перед природой жизни даже самые неистовые ветра, даже самое беспощадное поветрие, даже жгучая кровь земли.
Не был Безумный умён, как сестра Вода, и не нашёл он лазейки в связавшей его клятве.

Смирилась и сестра Вода, хоть плакала она от обиды и в тот год, и в следующий;
и солона была Ху'тресоааса, и снег у подножья Серебряных гор был солон, и даже дождь отдавал горечью в те года\FM.
\FA{
Горькие дожди --- период после извержения супервулкана Закованный Пик, нашедшего отражение в сказаниях многих народов.
Ориентировочно Большая Капля 8126 --- Пирог 8129 (дата подтверждена Лабораторией Катастроф, отдел геологии и классической планетологии Ордена Преисподней).
}

\chapter{[-] Горькие дожди}

\section{[-] Разлом}

\epigraph
{Чем шире распахнута дверь, тем раньше видно друга и врага.}
{Пословица сели}

\textspace

В тот дождь снова активировались землетрясения.
Тревожные сообщения поступали и с севера, и с запада.
Но вскоре беда прошла прямо у нас перед носом --- расширился разлом Гибельных Руин.
Ворчливая земля обрушилась на ближайшие деревни хака, а игривые ветры понесли смертоносные испарения на восток.
Хака подняли зелёное знамя\FM.
\FA{
Зелёное знамя --- знак насильственного переселения.
}

Мы тем временем жили так же, как и всегда.
Это устраивало всех, кроме кормилицы;
она выглядела обеспокоенной.
Вскоре она созвала Советы Сада и Цеха.

--- Хака нужно помочь, --- заявила она.
--- Я бы направила к ним обозы с провизией и тканями.
Да, я знаю, многие из вас считают, что это не наша проблема.
Моё мнение --- это временно.
Проблема рискует стать нашей уже в ближайшую пару дождей.

Старшиной Сада была пожилая юркая женщина.
Её кипучей энергии хватало на целых шесть кварталов помимо огромного дома и обычной крестьянской работы --- признак человека, пережившего тяжелейшие лишения в детстве.
Она лаконично резюмировала общее решение двух Советов:

--- Каким бы образом проблема ни стала нашей, провизия и ткани пригодятся нам самим.

--- Вот так вот, --- сказала Кхотлам за ужином, завершая рассказ о заседании.
--- Мне уже давно следовало написать хака.
А что писать теперь?
Я не знаю.

--- Но Саритр права, --- пожала плечами Эрхэ.
--- Избыток может пригодиться в любой момент.

--- Попросят --- другое дело, --- поддержал женщину Хитрам.
--- Самим-то зачем отправлять?

--- А я бы отправила, --- заметила Ликхэ.
--- Проблемы проблемами, но люди голодают.
От нас не убудет.
Кроме того, хака иногда устраивали набеги, но вот такие долги всегда возвращали.
Для них бесчестье принять помощь и не отплатить тем же.

Сиртху-лехэ промолчал, но всё было ясно и так --- к хака и идолам Живодёра у старика были личные счёты.
Кормилица посмотрела на меня:

--- Что скажет Храм?

--- Очень невежливо было бы написать <<попросите нас>>, но другого выхода из ситуации я не вижу, --- признался я.
--- Если прибудут делегаты, Сад и Цех уступят, хотя бы сделают скидку.

Кхотлам кивнула и встала из-за стола.

--- Всегда есть способ выразить вежливо даже невежливую мысль, Лисёнок.

\begin{quote}
<<Мы слышим о вашем горе.
Мы знаем силу духа народа хака.
И если телу народа понадобится пища и одежда, вам нужно лишь окликнуть ближайших соседей>>.
\end{quote}

\section{[-] Ответ}

\textspace

Ответ от хака пришёл спустя три рассвета.

\begin{quote}
<<Мы готовы купить у вас половину избытка с полей, и половину избытка ткани, и половину запасов строительной древесины по обычной цене.
Караван выступает через четыре рассвета после написания этого письма.
Просим распорядиться о складировании и сортировке как можно быстрее>>.
\end{quote}

Кормилица, получив его, весь день хмурилась и мяла кусок пергамента в руке.

--- Интересно, откуда у хака столько золота.
И почему четыре...
Так, повременим-ка пока со складированием...

\section{[-] Игры не будет}

\textspace

--- Ну что, Ликхмас-тари? Может, ещё партейку в Метритхис? --- ухмыльнулась Кхохо.
--- Надеюсь, ты не разобьёшь себе губу и не найдёшь потерявшуюся в джунглях тратра\FM.
\FA{
Тратра (дикий цатрон) --- термин биологического родства, дарительница третьего прародителя по мужской линии.
В племенах царрокх, а затем хака и тенку те, кто достиг статуса тратра, считались самыми уважаемыми женщинами, талисманами племени;
иногда их авторитет даже превосходил авторитет старейшин.
}

Я подмигнул Кхохо и сел за стол.

--- Очень хороший игрок, --- отрекомендовала меня Кхохо.
--- Только иногда мне ему лицо разбить хочется.
Так, расставляем фигуры...
Молоток, я видела, как ты засунул камень в рукав.
Положи обратно, мелкий жулик.

Крестьянин смущённо улыбнулся щербатой улыбкой и положил камень в горшок.

--- Зелёный, --- вздохнула Кхохо.
--- На кой ты взял зелёный, дурачок?
Ты же Сосед, а не Король!

--- Я Король! --- возмутился Молоток.

--- Да? --- Кхохо задумчиво посмотрела на стол.
--- Точно.
Что-то я замоталась на работе.
Да чтоб вас всех, я фигуры не так расставила, я ж думала, что ты Сосед...

Вдруг в дверях таверны появилась тень и, показав какой-то знак, мотнула головой.

--- Меркхалон-кровопийца, опять работа, --- выругалась Кхохо и подхватила лежавшую рядом серебряную саблю.
--- Пошли, Ликхмас, сегодня игры не будет.

\section{[M] Вкус броненосца}

\textspace

В глазах всех, кого мне пришлось убить, было одно --- страх и слёзы.
Я убивал, чтобы кормилица спала спокойно, а сестрёнки утром шли за водой, распевая песенки.

Вспомнился один день из моего детства.
В тот день погиб мой брат.

Я, возрастом пятнадцати дождей, стоял на стене Тхитрона со взведённым <<детским>> арбалетом.
Детским он назывался, разумеется, из-за размера --- убойность его не уступала лукам взрослых.
Кхотлам стояла рядом.
Мы не отрываясь смотрели на границу леса.
Там было тихо, но мы знали --- деревья нашпигованы идолами, как пирог яблоками.
Многочисленный отряд Молчащих идолов прошёл по северным землям хака, как стая степной саранчи.
Половина земель Травинхала курилась дымом запустения, теперь пришёл черёд Тхитрона.

Первыми до ворот дошли беженцы хака.
Мы с Кхотлам застали эту новость дома, за едой.
Кормилица бежала к воротам, словно её гнали горящими палками.

Воины Храма, что удивительно, стояли гурьбой и жаром переговаривались.
Увидев Кхотлам, все разом замолкли.
Вождь вышла вперёд.

--- Кхотлам, я, кажется, сказала тебе...

--- Впустить их, --- без предисловий приказала кормилица.

--- Отбой, --- бросила вождь.
--- Отправьте им сигнал...

--- Я сказала --- впустить их, --- рявкнула кормилица, уткнув в лицо вождю <<птичью лапу>>.

Купец имел право в особых случаях наложить вето на приказ вождя.
Кхотлам воспользовалась этим правом первый и последний раз в жизни.
Несколько воинов бросились выполнять приказ.

Ворота открылись, и унылая цепочка повозок заехала в город.
Пугливые женщины, дети с серьёзными взрослыми глазами, сгорбленные старики --- я никогда не видел ничего более внушающего жалости.
Раненых, что необычно, не было совсем --- возможно, хака оставили их.
Немногочисленные мужчины тут же, не задавая вопросов, поднялись вместе с нами на стену и взяли луки на изготовку.

Чтобы снять моё напряжение, Кхотлам спросила:

--- Ликхмас, о чём думаешь?

Конечно, я ненавидел идолов.
Чистой, детской ненавистью.
Детям часто говорили: Молчащие --- непримиримые, смертоносные враги.

--- Я думаю о том, как убить больше врагов.

--- А знаешь, о чём думаю я, когда целюсь?

Я покачал головой.

--- Я думаю о том, что после я куплю на рынке самого большого и жирного броненосца и приготовлю его, а Хитрам, Хмурый и ты будете сидеть за столом, уплетать его и облизывать пальчики, --- прошептала Кхотлам, прищурилась и отпустила тетиву.
Застонал разорванный воздух, с далёкого дерева упало верещащее зелёное тельце, ударилось о землю и замолчало навсегда.

Кормилица наложила новую стрелу на тетиву и продолжила как ни в чём не бывало:

--- Почему ты ненавидишь их, Ликхмас?
У тебя есть повод?

--- Многие наши люди убиты идолами, --- сказал я и, заметив в кроне движение, нажал на спуск.
Стрела исчезла в листве, не причинив никому вреда.

--- Те, кто носят в ушах перья согхо и раскрашивают лица фиолетовой краской, думают прежде всего о покое для умерших близких.
Хотя... я не думаю, что покой мёртвых зависит от судьбы их убийц.

Ещё одна стрела Кхотлам нашла свою цель.
Мучения легко раненного идола завершила другая стрела.

--- Хай, о чём мне думать во время боя? --- спросил я.

--- Думай о моём броненосце, --- весело предложила кормилица и выпустила стрелу.
На границе леса осталось лежать ещё одно зелёное тело.

Идолы вскоре смекнули, что зря теряют товарищей.
Вскоре через лужок к стене пошли построенные <<черепахой>> отряды.
Воздух наполнился стрелами, копьями и пращными снарядами.

Вождь зычным голосом велела воинам спуститься за стену.
Уязвимое место идолов --- бой на открытой местности, строй на строй.
Приказание вождя повторил большой <<осадный>> барабан.

Брат стоял в первых рядах.
Про Саритра говорили, что он похож на большого старого ягуара.
У него была чересчур массивная челюсть, маленькие, близко посаженные глаза и тонкий кривой нос.
Он был нелюдим и очень лаконичен на словах.
При этом мало кто мог выстоять против него один на один --- его ловкость и силу тоже сравнивали с кошачьими.

Я же знал этого сурового (в его-то тридцать дождей) воина как доброго Хмурого, который учил меня мудрёным приёмам, рыбалке, а ещё таскал для меня у рыночных торговок сладкие конфеты.
Старшие, бывало, поколачивали младших за шалости.
Хмурый за всю свою короткую жизнь не тронул меня и пальцем --- даже на шуточных поединках он меня щадил, за что кормилец ему выговаривал.

--- Испортишь брата, дитя.
Он должен привыкнуть к ударам, враг не будет его жалеть!

Нелюдим Хмурый улыбался толстыми губами.
Уж он-то знал, как мне доставалось от Конфетки.
Брат легонько лупил меня деревянной саблей, я притворно сгибался и корчился.
Хитрам отвешивал нам обоим подзатыльники, больше похожие на ласку, и возвращался к делам.

... Идолы приближались под градом стрел и снарядов.
<<Черепахе>> эти кусочки дерева и камня не причиняли никакого вреда, но и вести ответный огонь идолы тоже не могли.
Когда первый отряд подошёл на расстояние выстрела из духового ружья, воины сели, выставив щиты, бросились в атаку.

Саритр влетел во вражескую <<черепаху>> первым.
Удар его щита помял строй идолов, покорёжил черепаший панцирь.
Он рубил и кромсал зелёных карликов как одержимый, как вдруг идолы поменяли строй.
Из-за выстроившихся полумесяцем щитов высунулись духовые ружья.

Новая тактика идолов стала для людей полной неожиданностью.
Огонь со стен не причинял полумесяцу никакого вреда, а с его внутренней, полуоткрытой стороны в растерявшихся воинов полетели яркие оперённые стрелы.
Саритр вместе с цветом авангарда --- их было человек пятнадцать --- упал, получив дозу кураре.
Две <<черепахи>> неожиданно сошлись, сомкнули щиты над упавшими и быстро понесли их к лесу.
Рванувшихся вслед за ними воинов сели встретили бамбуковые дротики и шквальный огонь с деревьев.

--- Назад! --- выкрикнула вождь, осознав свою ошибку.
Барабан повторил её приказ.
Сели построили у стены <<черепаху>> и замерли.

--- Кормилица!
Они схватили Хмурого!

Я в отчаянии выпускал одну стрелу за другой, но все они отскакивали от щитов и вонзались в землю.
Кхотлам стояла, натянув тетиву, и по её лицу текли слёзы.

--- Кормилица, почему ты не стреляешь?..

Она стояла и смотрела...
Один за другим стрелки прекращали бесплодные попытки стрельбы.
Наконец, когда <<черепаха>> с пленниками уже почти преодолела расстояние до леса и начала разваливаться, Кхотлам спросила меня:

--- Помнишь, я говорила тебе про броненосца?

--- Кормилица!..

--- Он будет таким же вкусным и большим, как я обещала, --- прошептала Кхотлам и,
внезапно подняв лук, спустила тетиву.

Стрела вошла <<навесом>> в полупядевую щель между щитами.
Идолы заверещали, но это был крик разочарования, а не боли.
Черепаха прошла, а один из захваченных людей остался лежать на земле со стрелой в шее.

Мой добрый брат Саритр.

--- Никто из моих детей не умрёт на алтаре, --- тихо, но твёрдо прошептала кормилица.

\asterism

Ночь выдалась тихой.
Идолы развернули лагерь где-то в лесу, и до стражи доносились их весёлое верещание и детские голоса.
Словно насмешка над плачущей песней пятнадцати человек, вставших кругом на улице недалеко от нашего дома.
Кхотлам не пошла к ним.
Она сидела у окна и молча слушала, перебирая бусины занавесей.

Вождь хака и несколько самых уважаемых воинов ночевали в нашем доме.
Перед сном они долго сидели на полу гостиной, попыхивая длинными трубками.
Я хорошо запомнил их грубые горбоносые лица, освещённые слабыми бликами очага.

--- Скажи мне, Кхот-Трис ин ур'Лим, --- сказал вождь, --- как получилось, что хака всё ещё живут между землями сели и Молчащими лесами?

Кормилица всегда была на редкость женственной и нежной;
однако вождь хака произносил её имя на мужской лад и говорил о ней в мужском роде --- видимо, чтобы не потерять лицо перед своими воинами.

--- Если хака не жгут наши поля и не убивают наших сыновей, то они хорошие соседи, --- сдержанно ответила Кхотлам у окна.
\ml{$0-[tn]$}
{--- По каким соображениям вас терпят идолы, мне неведомо.}
{``I have no idea why idols tolerate your presence too long.''}

\ml{$0-[tn]$}
{--- Ты прекрасно знаешь, что уже завтра я могу прийти вместе с союзом племён как враг, --- сказал вождь.}
{``You know well: very soon I may come as enemy, within the Tribal Union,'' the chieftain said.}

--- Тогда я убью тебя и вывешу твою голову у ворот, Си-Жак, --- без обиняков посулилась кормилица.

\ml{$0-[tn]$}
{--- Это будет честно, --- согласился вождь.}
{``It would be fair,'' the chieftain admitted.}
--- Если бы не Кхот-Трис ин ур'Лим, Инхас-Лака сейчас гнили бы на родных пепелищах.
\ml{$0-[tn]$}
{Мне жаль твоего сына.}
{I'm sorry about your son.}
\ml{$0-[tn]$}
{В его чреслах таилась дивная мощь.}
{He'd got incredible powers in his loins.''}

--- Я позабочусь, чтобы вам выдали достаточно пищи на сезон.
У нас избыток.

--- Мы пришлём вам дары чести, когда это будет возможно, --- поклонился Си-Жак.

--- Лучшим даром чести будет, если Инхас-Лака останутся дома во время следующего похода хака на сели.

Вождь дёрнул проколотой губой с клыком бородавочника и бросил взгляд на сидящих рядом воинов хака.
\ml{$0$}
{Предложенное Кхотлам означало бесчестье, как, впрочем, и отказ от предложенного --- просьба спасителя священна.}
{Acceptance of \Kchotlam's offering meant dishonor, and refusal too---saviour's wish is sacred.}
\ml{$0$}
{Этого кормилица и добивалась --- вождю придётся сделать выбор.}
{This was what my nurse wanted---to make the chieftain choose.}

\asterism

Подошедшая утром подмога с Хатрикаса не нашла уже никого --- аборигены Молчащих лесов снялись ещё до рассвета и ушли восвояси.
Несколько деревень были опустошены накануне, и пленных для жертвоприношений они получили предостаточно.

Хмурого похоронили в лесу.
Хитрам выбрал самое сильное и узловатое чёрное дерево и нарисовал на его коре знак Сана-сновидца.

--- Он был воплощением силы, --- сказал отец ритуальные слова.
--- Если его посмертие не будет лёгким и спокойным, то нам да будет отплачено впятеро.

Кормилица, когда все ушли, аккуратно дорисовала рядом Сита.
Однако её руки дрожали, и краска на веках Сита потекла.
Улыбчивый любвеобильный дух заплакал горькими чёрными слезами.

Броненосец и вправду получился очень вкусным, но сильно пересоленным.
Мы с кормильцами сидели за низеньким столом и молча уплетали нежное мясо.
Под конец трапезы Кхотлам нарушила молчание.

--- Пообещайте одну вещь.
Ни один из вас не наденет перьев согхо.
И ничьего лица не коснётся фиолетовая краска.

Хитрам ар’Кхир происходил из семьи крестьян.
Чтить законы сели было у него в крови.

--- Ты сама понимаешь, что это невозможно, Пёрышко, --- проворчал он.
--- Хмурый не упокоится, если...

--- ... если не погибнет соплеменник его убийцы? --- закончила Кхотлам.
--- Стрела, убившая наше дитя --- моя стрела, Пловец.
Кому ты собрался мстить?

--- Хаяй, этого не может быть, --- Хитрам неверящим взглядом смотрел на кормилицу.
Она смотрела в тарелку.

--- Никто из моих детей не умрёт на алтаре, --- повторила она.

... Кормилец ушел, не закончив еду.
Река вышла из берегов и вернулась восвояси, а мы не услышали о Хитраме ничего.
Кхотлам занималась своими делами --- нанимала носильщиков, распределяла товары, --- но делала это всё отстранённо.
Она рассчитала всех посыльных, оставив только милую крестьянку Эрхэ --- её дед был другом детства Кхотлам.
Она перестала готовить и заниматься рукоделием.
По вечерам кормилица долго сидела у моей лежанки и гладила меня по голове.

--- Одни мы с тобой, Лисёнок.
Совсем одни.
Что ты один, что я одна...

Я играл с волосами кормилицы и говорил, что, когда вырасту, то увезу её из Тхитрона в спокойное место.
Она грустно смеялась, целовала меня и уходила к себе.

Однажды в дверь постучали среди ночи.
Шёл проливной дождь.
Я по привычке соскочил с лежанки и спрятался за занавесями, наблюдая за входом.

За тяжёлыми резными дверями чёрного дерева стоял Хитрам.
По его длинным каштановым волосам стекала вода, вода стекала и по лицу, на котором прибавилось свежих шрамов.
Он стоял молча, опустив голову, и держал в руке промокший букетик розовых орхидей.

Кхотлам, не сказав ни слова, втащила его в дом, отвесила Хитраму звонкую пощёчину и поцеловала в губы.
Я робко вышел из-за занавесей, и кормилец, утирая с лица слёзы, перемешанные с дождевыми каплями, подбежал ко мне и поднял на руки.

В ту ночь были зачаты мои сёстры.

Больше Хитрам нас не покидал.

\section{[@] Ещё люди}

Экспедиция на Корону принесла неожиданный результат: мы нашли людей.

Я спросил Безымянного, почему он нам про них не рассказал.
Бог ответил: <<Я не знал, что они тоже относятся к людям>>.
Если честно, странная отговорка.
Биологи на всякий случай попросили его подробнее рассказать о формах жизни на планете --- мало ли кто тут ещё скрывается!

<<Царрокх>>, как себя называет найденное племя, --- потомки людей Древней Земли, живших незадолго после Последней Войны.
В пользу этого говорит чрезмерно развитый половой диморфизм (сравнительно высокие мускулистые мужчины и пухлые женщины, едва превышающие ростом плантов), недосбалансированный геном и отсутствие некоторых идиоадаптаций, характерных для сапиентов поздней Эпохи Богов.
В частности, они используют исключительно оральную артикуляцию, а их пол определяется в момент образования зиготы.

Баночка сейчас сидит у меня расстроенный.
Царрокх категорически отказались общаться с ним.
Людей-тси дикое племя называет цели (Te'сели) --- золотые люди.
Сначала историки предположили, что это из-за имплантов.
Потом выяснилось, что импланты царрокх приняли за украшения, а <<золотыми>> людей назвали из-за цвета кожи (у аборигенов кожа коричневого цвета).
Особенно им понравилась ласковая красавица Гладит-Зелёную-Кошку --- её пускали в любой дом и позволяли задавать любые вопросы.
Она и собрала большую часть информации о племени.

Кани, апид и плантов царрокх называют <<кхерр>> (понятия не имею, что это, но звучит очень неприятно) и ведут себя с ними трусовато-враждебно.
Баночка рассказал, что проникновение тси в святое место царрокх спровоцировало столкновение, и наполовину состоящий из биокибернетики Фонтанчик вызвал ужасный переполох.
Сам Фонтанчик отозвался кратко: <<Обошлось без жертв>>.
Видимо, были покалеченные, но всех удалось спасти.
Фонтанчик сказал, что больше к ним не поедет.

Не нравится мне это всё.

Поправка: Баночка только что по контексту нашёл возможное значение слова <<кхерр>>.
Это философское понятие Древней Земли, обозначающее умышленный вред или абсолютного, непримиримого врага.
Оно является антонимом другого понятия --- намеренное полезное действие, безусловное дружелюбие.
Баночка растерянно сказал, что, по-видимому, аналогов этих слов в языке тси не существует.

\section{[@] Кошкины слёзы}

\textspace

Кошка была совершенно счастлива.
В её глазах блестела мысль, она снимала предметы быта, рисунки и архитектуру, обновляла свой репозиторий каждые несколько минут.
Баночка сидел в палатке и с некоторым сожалением слушал её рассказы по телекому.

--- Всегда мечтала о такой работе, --- поделилась Кошка с Фонтанчиком.
Канин хмыкнул и покачал головой.

Вскоре пришёл черёд капища.
Шаман царрокх после долгих препирательств позволил двоим тси --- Кошке и Фонтанчику --- войти внутрь.

Внутри царило простое, одухотворённое великолепие.
Капище окружали увитые лианами тотемные столбы, земля была покрыта отполированными каменными плитами, меж которых пробивалась редкая травка.
Шаман с важным видом начал рассказывать о богах и духах.

Венцом капища был красивый базальтовый алтарь, украшенный совершенно чудесными орнаментами.
Гладит-Зелёную-Кошку, очарованная алтарём, поднялась на пару ступеней и протянула руку, чтобы ощупать каменный орнамент.

Вдруг одно из копий позади женщины рванулось вперёд.
Гладит-Зелёную-Кошку, в последний момент ощутив кожей неладное, попыталась увернуться, но было слишком поздно --- копьё вошло точно под сердце.

Фонтанчик взвыл, словно волк, которому в горло вживили тревожную сирену.
Не прошло и секунды, как он разбросал царрокх, словно деревянные игрушки, и успел подхватить подругу под руки.

Люди стояли, медленно осознавая происшедшее.
Несколько воинов, переломанные и избитые, стонали на земле.
Возле тотема медленно испускал дух неудачливый копейщик --- обитый металлом тупой конец оружия раздавил ему грудную клетку.
Прочие, с ужасом глядя на Фонтанчика, выставили в его сторону копья.
Инкрустированные обсидианом наконечники дрожали.

--- Трусливые твари! За что?! --- заорал Фонтанчик на ломаном цатроне, потеряв над собой контроль.

--- Женщина не должна касаться святого алтаря!
Это осквернение! --- завопил шаман.

\ml{$0$}
{--- А предупредить... нельзя было?.. --- всхлипнула Кошка.}
{``You couldn't ... warn me ?..'' Cat sobbed.}

\ml{$0$}
{--- Кхерр! А! --- рявкнул один из воинов.}
{``Kch\r{e}rr! \^{A}!'' a warrior roared.}

\ml{$0$}
{--- Кхерр! А!}
{``Kch\r{e}rr! \^{A}!}
\ml{$0$}
{Кхерр! А! --- подхватили остальные.}
{Kch\r{e}rr! \^{A}!'' the others joined.}

<<Медики, люди-мужчины.
Пять царрокх с множественными травмами, двое в агонии, один в клинической смерти, несу Кошку с травмой 12-E4C>>, --- передал Фонтанчик по общему каналу и, подхватив Кошку на руки, бросился вон из капища.
Царрокх расступились, дав ему дорогу.

Кошка тихо, прерывисто дышала, цепляясь слабеющими руками за плечи друга.
На её лице замерла детская обида...

--- Протокол <<Тайфун>>, --- отчеканил на бегу Фонтанчик, --- код... эээ...

Фонтанчик пробежался глазами по лицу Кошки, сжал и разжал её руку.

--- \emph{Ноль-ноль, восемьдесят эф, ноль-два.}

Лицо Кошки приняло умиротворённое выражение, тело расслабилось.
Кровь резко перестала течь из раны.

\ml{$0$}
{--- Животные, --- шептал Фонтанчик.}
{``Beasts,'' Fountain had been whispering.}
\ml{$0$}
{--- Тупые жестокие животные...}
{``Stupid, cruel beasts ...''}


\section{[@] Знакомство с Нейросетью}

\textspace

Баночка уже нашёл себе друга.
Когда мы пришли, по палатке неуклюже прыгала молодая любопытная ворона.

--- Это ещё что такое? --- слабо удивилась Кошка, приподнявшись с носилок.
Костёр шикнул и знаком велел ей лечь.

--- Это ворона, --- просветил подругу Баночка.
На его лице застыла маскообразная улыбка, впалые щёки ещё хранили следы ногтей.
--- Ты как себя чувствуешь?

--- Хорошо, --- ответила Кошка.
--- Сам как?

--- Я больше не отпущу тебя одну.

--- Я не была одна, --- грустно сказала Кошка и огляделась.
--- Фонтанчик здесь?

--- Здесь, --- ответила Заяц.
--- Нужду справляет.

Ворона каркнула на всю палату, привлекая внимание.

--- Ветвь Настоящие Птицы, Воробьи, Вороны, вид Большая скальная ворона с рекомбинированными первой, шестой и шестнадцатой хромосомами, отсутствующими плечами четвёртой, --- отрекомендовал птицу Мак, сверившись с компьютером.
\ml{$0$}
{--- С определённой натяжкой --- Девиантная ветвь, так как была собрана Безымянным, а не привезена переселенцами.}
{``Actually, the species may be called Deviant Fork, because it was compiled by Nameless, not brought by settlers.''}

--- Это Безымянный так изменил ей хромосомы? --- подняла бровь Заяц.

--- Нет, она сама, --- ухмыльнулся биолог.
--- Экваториальная радиация.

--- А как ты понял, что это Девиантная ветвь?

--- Митохондрии не вороньи, --- объяснил Мак.
--- Листик уже говорила на эту тему с Безымянным.
Он сказал, что да, в период его творческого взлёта многие виды получили митохондрии какого-то несчастного зайца, которого и на планете-то уже нет.
Он их только немного подогнал, чтобы работали.

--- То есть местная фауна в какой-то степени заяц, --- резюмировал Костёр, игриво стрельнув глазами в сторону подруги.

--- Ты можешь меня не бесить? --- раздражённо всплеснула руками Заяц.
--- Благодарю.

--- Это объясняет в том числе и факт, что многие эволюционно далёкие виды восприимчивы к одному и тому же яду --- продолжал Мак, уткнувшись в компьютер.
--- Токсин растения, которое царрокх называют <<Снежной Бородой>>, необратимо ингибирует один из ферментов окислительного фосфорилирования именно у этого типа митохондрий.
Просто спросив, на кого царрокх охотятся с этим ядом, мы с Листик выявили ещё пять девиантных видов, хахаха!.. --- Мак поднял глаза на бледную Кошку и спохватился:
--- Извините.

--- А как ты назовёшь эту ворону? --- спросила Кошка у планта, пытаясь сгладить напряжённость.

--- Нейросеть, --- сказал Баночка.
--- Она даже отзывается.

Ворона действительно отзывалась.

--- Красивое имя, --- заметила Заяц.
--- А главное --- редкое.
У девяноста процентов моих знакомых животные носят эту кличку!
Давай её хотя бы Чернушкой назовём или Крылаткой!

--- Хорошо, пусть будет Биологическая самообразовавшаяся нейронная сеть с функциями клевания пищи, общения и воспроизводства.

--- Ты сейчас дал исчерпывающее описание половины тси, --- засмеялся Костёр и тут же осёкся под взглядом Заяц.

--- Длинно, --- резюмировала Заяц, продолжая сверлить врача взглядом.

--- А для краткости --- Нейросеть, --- выкрутился Баночка.
--- Кстати, зайцев я тоже порядочно знаю.

Женщина не нашлась что ответить.

\asterism

\ml{$0$}
{--- Баночка, мы не можем взять ворону с собой!}
{``Flask, we must not take the crow with us!''}

--- Кошка, лежи.
Тебе вредно волноваться, --- пытался урезонить женщину Костёр.

--- Почему? --- допытывался плант.
\ml{$0$}
{--- Нет, Костёр, подожди, не трогай её имплант.}
{``No, Campfire, wait, don't touch her implant.}
\ml{$0$}
{Пусть ответит.}
{Let her answer.''}

--- Здесь её вид!
Здесь её дом! --- Кошка начала выходить из себя.
--- Мак, ну скажи ты ему!

Мак смущённо молчал и тёр бритую голову.

--- Хорошо, давай так, --- Баночка вытащил ворону из ящика с инструментами и посадил на руку.
--- Нейросеть поедет вот так, у меня на руке.
\ml{$0$}
{Я не буду её держать.}
{I don't hold her.}
\ml{$0$}
{Пусть сама решает.}
{Let her decide by herself.}
\ml{$0$}
{Мак?}
{Poppy?''}

\ml{$0$}
{--- Это будет честно, --- пожал плечами Мак.}
{``It would be fair,'' Poppy shrugged.}
\ml{$0$}
{--- Кошка, без обид.}
{``Cat, no offence.}
\ml{$0$}
{Птица мне тоже очень нравится.}
{I like this bird too.''}

Кошка вздохнула.

\ml{$0$}
{--- Хорошо.}
{``You won.}
\ml{$0$}
{Но если она улетит --- не смей её звать.}
{But if it flies away, you dare not call it back.''}

\ml{$0$}{--- Договорились.}
{``Agreed.''}

Обратный путь получился очень занятным.
Заяц держала в объятиях Фонтанчика и гладила его по голове;
друг грустно думал о чём-то своём.
Я пытался соблазнить их печеньем.
Баночка сидел снаружи, на корме, и гладил Нейросеть.
Прочие наблюдали за ними --- большинству была интересна реакция вороны, но кто-то очень хмурый на носилках строго следил, чтобы Баночка не держал птицу за лапки и крылья.
Впрочем, ворона была не против приключений.
Её не смутил даже пролив Скар.
Пока транспорт шёл через тёплые воды, Нейросеть смирно сидела на руке планта и клевала орешки.
И едва нос лодки коснулся берега и трансформировался в сухопутную подвеску, птица по-хозяйски отправилась осматривать окрестности.

--- У меня ощущение, что кое-кто бессовестно мигрировал посредством нас, --- заявил Мак за обедом.
\ml{$0$}
{--- Притворилась ручной, преодолела пролив и сбежала.}
{``It was pretending to be pet while crossing the strait, then fled.}
\ml{$0$}
{Ворона оказалась умнее десяти тси.}
{Ten Qi outsmarted by a crow.''}

\ml{$0$}
{--- Она вернётся, --- ответил Баночка.}
{``She will be back,'' Flask answered.}

\ml{$0$}
{--- Да откуда тебе знать? --- поморщился Мак.}
{``How can you know?'' Poppy frowned.}

Баночка загадочно улыбнулся и кивнул на окно.

Нейросеть сидела на подоконнике.
Глаза вороны горели жаждой орешков.

\section{[-] Старый вождь}

\textspace

--- Сколько родов?

--- Прошу прощения? --- переспросил посланец.

Кхотлам устало прикрыла глаза.

--- Цвета, знаки, символы, всё, что запомнили!

--- Были с красными перьями...

--- Красные перья ара, одна жёлтая серьга в левом ухе, трубят в козьи рога, --- резко сказала Кхотлам.
Она не терпела такого невежества относительно ближайших соседей.

--- Да, точно они.

--- Бвожай-сау.
Ещё кто-то был?

--- Мы слышали какой-то странный свист, как будто...

--- Белая рыба на синем фоне, синие перья самца кхар-тхир, белые серьги, прерывистый свист, --- кормилица издала оглушительную свистящую трель, заставив сидящих поблизости воинов похвататься за оружие.

--- Да, --- кивнул посланец.

--- Ты уверен?

--- Я уверен, синие перья и белые серьги видели, --- вмешался другой.

--- Инхас-лака, --- пробормотала Кхотлам.
--- Похоже, мой милый друг Си-Жак умер.

\chapter*{Интерлюдия VI. Две речи}
\addcontentsline{toc}{chapter}{Интерлюдия VI. Две речи}

\textbf{Речи ноа}

\section*{Речь о мужчине}
\addcontentsline{toc}{section}{Речь о мужчине}

\ml{$0$}
{Живёт мужчина.}
{They a man live.}
\ml{$0$}
{У него есть имя, лицо, руки и ноги, мужской орган, инструменты, оружие, жилище и зелёный парус с вытканным солнцем.}
{They have a name, a face, hands and legs, a man's limb, tools, weapons, a house, and a green sail of woven sun.}

\ml{$0$}
{Однажды к нему пришёл второй и сказал ему.}
{Once they another man come and tell them.}
\ml{$0$}
{Мужчина отказался.}
{They a man say no.}
\ml{$0$}
{Тогда второй сказал ему: <<Ты не мужчина>>.}
{They another man say them: you are not a man.}
\ml{$0$}
{Но мужчина, у которого есть имя, лицо, руки и ноги, мужской орган, инструменты, оружие, жилище и зелёный парус с вытканным солнцем, просто засмеялся в ответ.}
{But they a man who have a name, a face, hands and legs, a man's limb, tools, weapons, a house, and a green sail of woven sun, just laugh instead of an answer.}

\ml{$0$}
{Однажды пришла к нему женщина и сказала ему.}
{Once they a woman come and tell them.}
\ml{$0$}
{Мужчина отказался.}
{They a man say no.}
\ml{$0$}
{Тогда женщина сказала ему: <<Ты не мужчина>>.}
{They a woman say them: you are not a man.}
\ml{$0$}
{Но мужчина, у которого есть имя, лицо, руки и ноги, мужской орган, инструменты, оружие, жилище и зелёный парус с вытканным солнцем, просто засмеялся в ответ.}
{But they a man who have a name, a face, hands and legs, a man's limb, tools, weapons, a house, and a green sail of woven sun, just laugh instead of an answer.}

\ml{$0$}
{Однажды пришёл к нему второй и улыбнулся.}
{Once they another man come and smile.}
Они взяли инструменты, вышли вместе и сотворили красоту, которой не было равных.
Они оба были мужчинами, у которых есть имена, лица, руки и ноги, мужские органы, инструменты, оружие, жилища и зелёные паруса с вытканным солнцем.
Смеялись они одинаково громко.

Однажды пришла к нему женщина и улыбнулась.
Он привёл её домой и сделал ей ребёнка, и когда они делали детей, они смеялись.
У женщины тоже было имя, лицо, руки и ноги, инструменты, оружие и жилище, но её органы были женскими, паруса --- голубыми, а в их центре были вытканы звёзды и рыбы.

\ml{$0$}
{Однажды пришёл к нему ребёнок и спросил: <<Почему ты улыбаешься, если никто не называет тебя мужчиной?>>}
{Once they a child come and ask: why you smile if they no one call you a man?}
\ml{$0$}
{Мужчина ответил: <<У меня есть имя, лицо, руки и ноги, мужской орган, инструменты, оружие, жилище и зелёный парус с вытканным солнцем, и я умею обращаться со всем вместе и с каждым по отдельности>>.}
{They a man answer: I have a name, a face, hands and legs, a man's limb, tools, weapons, a house, and a green sail of woven sun, and I can handle all they things and every thing of them.}
\ml{$0$}
{И ребёнок остался в доме мужчины, пока не стал мужчиной сам.}
{They a child stay in a man's house until they become a man themselves.}

\section*{Без края}
\addcontentsline{toc}{section}{Без края}

Живёт некто по имени Инфинит\FM.
\FA{
В оригинальной версии речи по отношению к Инфиниту употребляется единственно-безличное местоимение, т.е. Инфинит --- личность, не имеющая ни пола, ни возраста, ни биологического вида.
В речах употребляются исключительно единственно-безличные местоимения --- даже если по ходу повествования становятся ясны пол и возраст героя.
В разговорном ноа-лингва ЕБМ употребляются по отношению к богам либо в иронично-высокопарном ключе.
На сектум-лингва ЕБМ обычно переводится как <<некто конкретный, о котором говорящему не известно ничего, кроме сказанного>>.
}

Он бессмертен и хорош собой.

Идёт Инфинит по дороге.

\ml{$0$}
{Когда падает дождь, он достаёт зонтик и поёт песню про дождь.}
{When it rains, they take an umbrella and sing of the rain.}

\ml{$0$}
{Когда дождь заканчивается, он надевает шляпу и поёт песню про жаркое солнце.}
{When the rain stops, they wear a hat and sing of the hot sun.}

Когда наступает ночь, Инфинит ложится спать.

Когда-нибудь Инфинит отрастит крылья и будет лететь над дорогой.

А когда крылья его устанут, у него появятся жабры и он поплывёт.

Когда солнце будет чересчур жарким, Инфинит будет идти по ночам.

\ml{$0$}
{Когда песни перестанут быть символом радости и печали, Инфинит начнёт чирикать или хлопать в ладоши, и все будут делать то же самое.}
{When a song is no more a sign of happiness or sadness, the Endless shall tweet or clap, and the others shall do he same.}

А когда земля закончится, Инфинит пойдёт по пустоте искать новую землю, по которой можно ходить и над которой можно летать, или воду, в которой можно плавать.

Когда свет погаснет, Инфинит найдёт другой свет и будет радоваться ему.

\ml{$0$}
{А у мира края нет, и другие земли всегда будут рождаться из пустоты, и другие воды всегда будут падать с небес, и другой свет всегда найдётся.}
{And the world is edgeless, and an earth can always be born from the void, and waters can always fall from a skies, and a light can always be found.}

Инфинит знает это и всегда идёт, останавливаясь лишь для того, чтобы отдохнуть и набраться сил.

Он будет жить всегда, и всегда будет хорош собой.

\chapter{[-] Круг}

\section{[-] Раненые}

\textspace

Новые раненые прибывали каждый кхамит.
В храмовый госпиталь вызвали всех свободных жрецов.
Хирургических инструментов и расходников не хватало.

\ml{$0$}
{--- Ликхмас, ты закончил с ногой и двумя проникающими в брюшную полость? --- спросил Кхатрим, выглянув с хирургического балкона.}
{``\Likchmas, are you done with the leg and the double penetrating abdominal injury?'' \Kchatrim\ asked from the surgical balcony.}

\ml{$0$}
{--- Пошти, --- отозвался я, зубами развязывая тугой узел жгута на безымянной ноге.}
{``Almosht,'' I answered, nearly lying on the nameless leg and undoing a very tight tourniquet knot with my teeth.}
\ml{$0$}
{Два ранения брюшной полости с любопытством и лёгким сочувствием наблюдали за процессом.}
{Double penetrating abdominal injury was watching the process with curiosity and a bit of compassion.}

--- Как закончишь, налови в пруду саламандр и собери секрет, --- устало попросил врач.
--- Хирургический клей заканчивается.
Размажь его тоненько по пластине и положи в тёплое место, чтобы быстрее высох.
Вон, тут лежачие есть, кому-нибудь на живот.
На солнце только не клади, он теряет свойства.
А, и руки не забудь охладить в воде, чтоб саламандр не обжечь и они не заболели.

--- Какие ещё саламандры там заболели? --- взмолился Ликхсар, вытирая лоб полотенцем.
--- Саламандр я лечить не буду, мне и людей хватает!
Восемь рожениц в моём районе сегодня было!
Восемь!
Это заговор какой-то --- как набег хака, так всем приспичивает рожать!

--- Всё нормально, у них нерест, --- ухмыльнулся Трухквал.

Жрецы захихикали в маски.

--- Ликхмас, заодно волоса нарежь и обесцветь перекисью, тоже заканчивается, --- добавил Ситлам, выпрыгнув из противоположного балкона и выбросив инструменты в таз.
--- Кошмар какой-то.
\ml{$0$}
{Хэй, там, внизу!}
{\Hei, downstairs!}
\ml{$0$}
{Есть ещё кто-нибудь в состоянии между <<чуть-чуть болит>> и <<чуть-чуть умер>>?}
{Is there anybody else between `hurt a little' and `dead a little'?}
Всё, что ближе к трупу, поднимайте ко мне.
Всё, что точно труп, пока не уносите, только облейте холодной водичкой --- могут пригодиться детали и обшивка.

--- Ситлам! --- возмущённо крикнула откуда-то снизу Хитрам.
--- Ты сейчас у меня выгребешь за <<детали и обшивку>>!

--- Ну не при людях же, серьёзно, --- согласился Трукхвал.

--- Трукхвал, я за сегодня вставил три почки и подшил столько кожи, что ею можно обить весь храм!
Имею я право хотя бы выражаться так, как считаю нужным?!
От меня люди уходят обиженные, но зато живые.

--- Да всё нормально, Ситлам, --- успокоил жреца басовитый крестьянин.
--- Ты меня не слышал, когда я после дождей первый раз поле вскапываю!
На моих ругательствах помидоры растут лучше, чем на навозе!

По лестнице поднялся мужчина, которого под руку вела Хитрам.
Мужчина шатался, словно кукла;
его лицо было полностью залито кровью.

--- Ничего не вижу, --- бормотал он.

--- Глаза целы, так что это временно, --- успокоил Ситлам, бросив на мужчину взгляд.
\ml{$0$}
{--- А вот твоему носу нужен очень хороший я.}
{``But your nose needs a small portion of my good self.}
Хитрам, давай его в восточный.

--- Я тебе челюсть сломаю, шутник, --- посулилась взмыленная Хитрам, выплюнув потную прядь волос.

--- Зачем? --- буркнул Трукхвал.
--- Чтобы потом ещё и его лечить?

--- Корнцанги надо заказывать, мне попался инвалид с неправильным прикусом, --- сообщил Первый жрец, клацнув рекомым инструментом.
--- Ликхмас, раз уж всё равно бежишь, можешь поправить на верстаке, хоть немного?
А то я точно у кого-нибудь что-нибудь внутри потеряю!

--- Народ, у кого вьющиеся волосы средней толщины? --- крикнул я пациентам.
--- Пожертвуйте красоту, расходники горят!

Я отхватил ножницами несколько прядей у Ситлама, Трукхвала и себя.
Затем добавил ещё прядь кудрявой плотницы со стрелой в ноге;
девушка подозвала меня и подставила голову.
Потом подбросил дров в жаровню, на которой кипел стерилирующий раствор.
Наконец я выхватил корнцанг из рук Первого и побежал дальше.

В госпитале всё кажется проще.
Разрезал, зашил, перебинтовал, напоил зельем.
Реальное горе, которое стоит за всем этим --- неизбывное, многодневное --- всегда оказывается где-то за его дверями.

\section{[*] Кристаллическая скала}

--- Почему мы не действуем? --- спросил Ситрис.
--- Во времена моей молодости храмовники прибегали, едва я появлялся на дороге.
У вас тут армия под стенами города, а вы лежите и дуете вино!

--- Разбойники --- чаще одиночки, их нужно выслеживать сразу, --- пояснила Митхэ.
--- Отряд налётчиков --- другое дело.
Отряд медленнее передвигается, оставляет больше следов.
Тем более армия.
Армию в сельве можно вообще не выслеживать, она передвигается только по дорогам.

--- По дорога идёт тележки с еда, а армия от еда никвогда далеко не ухводит, --- хихикнул Аурвелий.

--- Вождь отдал приказ эвакуировать все близлежащие поселения, --- сообщила Эрхэ.
--- Укрепления Тхартхаахитра неприступны.
По крайней мере, ни у кого из захватчиков ещё не хватило ума их взять, и репутация неприступности играет нам на руку.
Пока не начнётся осада, мы не при делах.
Так что пока расслабься и тоже выпей.

--- Они разграбят поселения, --- напомнил Ситрис.

--- Люди --- важнейший приоритет, --- ответила Митхэ.
--- Урожай вырастет, дома можно отстроить, главное --- как можно быстрее увести жителей из опасной зоны.

--- Тем более у нас избыток урожая, --- подхватила Эрхэ.
--- Разобьём хака --- и всё будет по-прежнему.
В отличие от этих бедняг, для которых любая война --- даже победоносная --- заканчивается голодом.
Дикарей жизнь ничему не учит.

Ситрис вытащил из кармана свёрток с коноплёй, набил трубку и закурил.

--- Почему Король-жрец не спешит с тобой встретиться, если он вызвал тебя аж из-за половины обитаемой Короны?

--- Я ещё не видела нового Короля-жреца, --- пожала плечами Митхэ.
--- Ничего не могу сказать.

--- Я уже могу сказать, что он намеренно отодвигает тебя от дел.
Тебя поселили в постоялом дворе и велели ждать.
У тебя достаточно опыта, чтобы руководить обороной, но из тебя сделали запасное колесо для повозки.

--- Заплатят --- и ладно, --- отмахнулась Митхэ.
--- Защищать Тхартхаахитр я бы согласилась и за идею, но в этот раз нам ещё и золота отвалят столько, что мы починимся и закупимся.

--- Расслабься, Ситрис, --- Эрхэ похлопала пузатый кувшин по глиняному боку.
--- Будешь вино или нет?

--- Не сейчас.
Пойду погуляю.

Ситрис затушил трубку пальцем, завернулся в плащ и вышел из постоялого двора в вечернюю сырость.

Что-то тревожило разбойника уже целый день, и слова Эрхэ только подогрели эту смутную тревогу.
В голове вертелись эпизоды из далёкого прошлого --- жилище, охваченное огнём, шум и беготня по всей деревне...

--- Привет.

Ситрис обернулся.
С ним мимоходом поздоровался молодой жрец --- кажется, его звали Митрис ар'Люм.
Синяя жреческая роба уже почти исчезла в темноте.

Ситрис вдруг понял, что это первый храмовник Трёх Этажей, который с ним поздоровался.

--- Хэй, Митрис, подожди!

Жрец обернулся.
Зелёные глаза вопросительно взглянули на разбойника.

--- Чем-то могу помочь?

--- Могу я составить тебе компанию?

Жрец колебался едва ли секхар.

--- Да, конечно.
Я не думаю, что воину Бродячего Храма будут интересны мои занятия, но могу уделить тебе немного времени для беседы.

--- А куда ты направляешься?

--- Хритра, --- Митрис указал на возвышающуюся над Тхартхаахитром громаду Кристаллической скалы.

--- Отлично.

По дороге наверх мужчины почти не разговаривали.
Ситрис был чересчур занят собственными мыслями, а Митрис не стал настаивать.
Внизу осталась и прорубленная в камне винтовая лестница, стёртая тысячами ступней, и развалины старого храма.
Когда они ступили на гладкую, покрытую редкой травкой вершину скалы, уже наступила ночь.

Митрис повозился с какими-то бесформенными металлическими конструкциями, вбитыми прямо в камень.
Конструкции зажужжали и затрещали.

--- Пойдём вниз, --- сказал жрец.
--- В Хритра теперь есть свет.

--- Можешь постоять здесь со мной?

--- Да, конечно.

Только здесь, с высоты половины кхене, Ситрис понял, насколько огромен Тхартхаахитр.
Улицы, путь по которым занимал больше кхамит, казались трещинами в пылающем угле.
Огромный, растёкшийся по поверхности планеты зверь, каждая часть которого жила своей собственной жизнью.

--- Скажи, сколько здесь жителей?

--- Сто тысяч, --- ответил Митрис.
--- Сто двадцать, если считать хутора и мелкие города.

--- На что хака вообще рассчитывают?
Они не смогут штурмовать такой крупный город.

--- Они не будут его штурмовать, я думаю.
Их основная цель --- взять контроль над Сикх'амисаэкикх.
Там много охотничьих угодий, оттуда до сих пор идёт небольшой, но стабильный поток золота и драгоценных камней.
Угроза столице --- больше психологический приём.
Если хака покажут, что они могут диктовать сели условия, диктовать условия начнут другие соседи.

--- Хатрикас пока держится.

--- Они не будут штурмовать Хатрикас.
Достаточно просто не давать сели пройти по Вильчатому тракту --- и Сикх'амисаэкикх у них в кармане.

Ситрис внимательно посмотрел на собеседника.
Тот рассуждал иначе, чем прочие жрецы Трёх Этажей, и, пожалуй, чересчур здраво для обычного библиотекаря.
Митхэ говорила то же самое --- разумеется, другими словами.

--- Что бы ты посоветовал сделать нам?

--- То, что мы и делаем --- эвакуация земель Тхартхаахитра, иллюзия слабости, дразнящие бои на берегах Ху'тресоааса.
Если часть армии хака, хотя бы одна пятая, потеряет голову и переправится на правый берег --- перекрываем мост и празднуем победу.
У нас хватит сил их разбить в поле, а оставшиеся будут деморализованы.

--- Они не осмелятся.

--- Пока они проигрывают в локальных стычках --- да.
Но если Тхартхаахитр даст слабину на подступах, либо, скажем, возникнут серьёзные проблемы с оборонительными линиями, им достаточно отправить диверсионную группу по Тропе Воров на Западный Тракт, сжечь мост через Ху'микхаса --- и центр страны будет парализован.
Да, Тропа Воров пока для них незнакома, да, они не смогут удержать позиции, но дай им достаточно времени, чтобы освоиться в землях Тхартхаахитра --- и они вполне смогут диктовать условия.

--- Что ты имеешь в виду под <<проблемами с оборонительными линиями>>?

--- Тхартхаахитр строился очень беспорядочно, --- напрямик ответил Митрис.
--- Дай руку, покажу.

Жрец подошёл к разбойнику вплотную, взял его за руку и начал чертить в воздухе линии, повторяя невидимые отсюда оборонительные укрепления.

--- Здесь стена.
Вот здесь.
Три башни, которые раньше были ключом к востоку, сейчас с них видно только рынок.
Вал --- вообще непонятно зачем.
Ещё группа укреплений вот здесь, по сути представляющая лишь историческую ценность.
Это не Тхитрон с его чётко очерченными линиями обороны, это не Ихслантхар, который, по-моему, вообще невозможно штурмовать, не имея тридцатикратно превосходящей население армии.
По сути город защищает река и одна стена, высокая, но одна, город, в котором пять Малых Храмов.
Убрать кусочек стены --- и все силы придётся бросить за защиту прорехи.

Ситрис задумался.

--- Бывают ли здесь пожары?
Имеется в виду, в мирное время.

--- К сожалению.
Мы стараемся корректировать застройку, чтобы огонь не распространялся, но люди, как обычно...
Чаще всего горит квартал кожевников, вон там, с северо-западной стороны.
А на востоке, видишь, есть дурацкая затопляемая площадь, которую мы зовём Торговым Озером.
От неё идёт улица Пылающего Амбара, угадай, откуда название.
Отсюда я видел пару раз...

Митрис осёкся.

--- Вот снова.
Та самая улица и горит.
Даже две, кожевники тоже запылали.
Два больших пожара, бьют в гонг.

--- В храме слышно?

--- Да, храмовники уже знают.
Скорее всего, туда отправят врачей...

--- Просигналь воинам, чтобы трубили сбор, --- произнёс Ситрис.
--- Хака предпримут диверсию с северо-востока.

Митрис взглянул на чёрное пятно северо-восточного Тхартхаахитра, выругался и бросился вниз.

<<Ещё одна ночь без сна>>.

\section{[M] Хорошее дело}

Во время обучения врачеванию у меня выработалась привычка спать аж три раза --- два раза ночью и один раз днём.
Этот странный режим я перенял от врачей.
Ночью людям в госпитале часто становится хуже, поэтому их нужно проверять.
В полдень, когда все обедают, не привозят почти никого --- идеальное время для того, чтобы добрать бесценный сон.

И вот однажды меня подняли именно во время полуденной полудрёмы.
Сдёрнутое чьей-то сильной рукой одеяло слетело на пол.

--- Ликхмас, быстро в душ и на балкон.
Сокхэ ар’Трукх привезли, лесной бык потоптал, --- сообщил Кхарас.

Михнет спустя я уже натягивал на балконе <<кровавый плащ>>.
Пострадавшая девушка лежала в маковом полусне с разрезанным животом.
Кхатрим молча кивнул на зеркала, и я подхватил их из рук Ситлама.

--- Шире.

Я кивнул и раздвинул операционную рану шире, поморщившись от открывшегося зрелища.
Воздух наполнился запахом бойни.
Ситлам взял бутылочку с коагулирующим раствором и корнцанг.

--- Бык поднял на рога, а потом ещё и прогулялся по ней, --- пояснил он.
--- Ветер духов уже подул, но её кормилец успел перехватить артерии.

--- Кишечник, матку и правую почку придётся удалить, --- добавил Кхатрим.
--- На них живого места нет.

Я наблюдал за операцией, мимоходом восхищаясь слаженностью действий жрецов.
Они понимали друг друга с полузнака.
Вскоре брюшная полость была очищена, и Кхатрим начал рассекать брыжейку обсидиановым скальпелем, выделяя остатки кишки.
Время от времени в большой каменный таз отправлялся очередной ком размозжённой, покрытой кровью и синими прожилками плоти.
Я отстранённо вспоминал названия сосудов, которые попадали в поле зрения --- Дерево Хетра, артерия Кобры, вена Двух Дорог...
Ситлам проверил белки глаз и посмотрел, запотевает ли зеркало у рта девушки.

Наконец Кхатрим тихо вздохнул.
Хороший знак --- жизни девушки пока ничто не угрожает.
Ситлам тоже повеселел.

--- Повезло ей, что её кормилец воином был.
У тех, старой закалки, ещё есть привычка носить с собой зажимы и скальпель.
Молодые плевали на элементарные меры предосторожности.

--- Кстати, зажимы наброшены качественно, --- отметил Кхатрим.
--- Не всякий жрец так сможет, да ещё и впопыхах.
Надо ему их потом вернуть.

--- Смотри, Ликхмас, --- Ситлам достал из полости плотный комок размером с кулак, --- вот в такой крохе дети вырастают до локтя в длину.

С этими словами врач отправил комок в таз.

Под балконом кто-то закричал, размахивая руками.

--- Рыбы безмозглые, --- выругался Ситлам.
--- Бычьи тропы людьми мостят, что ли?
Дома надо сидеть в сезон.
Ликхмас, поможешь, --- с этими словами жрец быстрым шагом вышел с балкона.

--- Шить.

Я одной рукой подал Кхатриму отбелённый волос и кривую иглу, с сожалением бросив взгляд на лежащее в тазике месиво.
Казалось, что там осталась большая половина Сокхэ.

--- Не волнуйся за неё, --- ухмыльнулся врач, проследив за моим взглядом.
--- Девушка сильная, здоровая, посмотри на её тело.
Главное --- это выходить сейчас, в критический момент, дальше природа сама её излечит.

Кхатрим жестом попросил ещё один волос.

--- Она будет такой же, как раньше? --- удивился я.

--- Почка обязательно вырастет, кишка тоже, но не полностью, где-то до двух третей длины.
Можно, в принципе, подсадить почку от умершего, но особого смысла не вижу --- это на случай, если обе пострадают.
Детей выносит, если пол не сменится, такие случаи я сам видел в Сотроне.
Купца нашего привезли, уже в летах такой роскошный мужчина, седой как снег в Грозовых горах.
Телегой придавило, да как нарочно --- по мужскому органу.
Бедолагу мы у Безумного выиграли с мышью на красном поле\FM, мужские глазки удалили, думали, что отрастут.
\FA{
Мышь на красном поле (Разведчик в Пустыне) --- крайне невыгодная оборонительная комбинация в Метритхис.
Победить можно лишь ограниченное число нападающих, и то если на кихотре выпадет число больше пятнадцати.
}
А купец раз и в женщину превратился.
Правда, сам он к этому с юмором отнёсся, весёлый был и за словом в карман не лез, --- голос Кхатрима вдруг потеплел, словно он вспомнил что-то очень приятное.

Я улыбнулся в маску.
Многих учителей я знал лишь стариками.
Было очень занятно вылавливать из поучений такие оговорки и наблюдать, как на покрытых морщинами лицах порой проскакивают следы бурной юности и былых страстей.

--- А как долго растут органы?

--- Растёт это хозяйство сезон-полтора, и до той поры кушать придётся противную жидкую кашу.
Первое время кормить будешь ты, я потом всё расскажу.

Вдруг в воздухе появился слабый, но отчётливый, оставляющий лёгкое першение в горле запах гари.
Кхатрим поморщился и поправил утиную маску.

--- Фильтр горит, --- недовольно заметил он.
--- Да что за день такой.
Ликхмас, после операции сбегай к Кхохо-лехэ на улицу Живой Рыбы, скажи, чтобы он на балконе фильтр заменил.

Я кивнул.
Кхатрим, ещё раз покачав головой, взялся за иглу.

Закончилась операция благополучно.
Я зашил мышцы и кожу отбелённым волосом, и девушку ещё дремлющей увезли в глубь Храма.
Кхатрим стянул с себя <<кровавый плащ>> и бросил его в угол.

\asterism

\textspace

Я постучал в провисшую дверь.
Спустя мгновение дверь приоткрылась, словно хозяин в этот ранний час стоял за ней.
В меня вонзился колючий взгляд из-под седых бровей.
Я оробел.

--- Митлам-лехэ?

--- Ла? --- буркнул старик.

--- С твоей хранительницей всё хорошо, --- сказал я.
--- Пришлось удалить много органов, но она выживет.
Сегодня вечером перевезём домой.

Дверь открылась чуть шире.

--- Хай, тебе просили передать, --- я протянул старику свёрток с зажимами.

Митлам-лехэ высокомерно смотрел на меня несколько долгих мгновений, потом протянул руку и схватил свёрток.

\ml{$0$}
{--- Хэ. Хоть это вернули. Трааа...}
{``\Choe. Leastwise, they returned this. \Traa\ ...''}

Я продолжал топтаться на месте.

--- Чего стоишь?
На завтрак напрашиваешься?

--- Нет, --- промямлил я.
Старый воин действовал на меня как-то чересчур удручающе.
--- Я хотел, чтобы ты проверил свёрток и сказал, что я ничего не забыл.
Тогда мне не придётся бегать ещё раз.

--- А дом Митлама-лехэ тебе не по душе? --- рявкнул старик.
Я попятился.

--- Плохо то жилище, в котором не рады гонцу, принёсшему хорошую весть, --- пробормотал я.

Старик оскалился и вдруг разразился смехом.

--- Хорошую весть, значит.
Заходи, молодой смелый жрец.
Заходи, не бойся.
Покажет тебе Митлам-лехэ, куда ты принёс <<хорошую весть>>.

Старик ударом расхлябил дверь и исчез в темноте.
Я, поколебавшись, вошёл внутрь.

В жилище царил дух запустения.
Треснувший очаг, наспех заделанный обычной приречной глиной, уже начинал горчить\FM, но хозяин даже не думал открывать вытяжку.
\FA{
Неточный перевод слова, обозначающего различаемый людьми-тси запах монооксида углерода.
}
От глиняной лампы остался только черепок с носиком и фитилём;
огонёк испустил дух, едва я возмутил движением воздух.
Стена с одной стороны держалась на паре досок.
С наступлением сезона дождей её просто смоет.

Старик, не спрашивая разрешения, прикурил лампу от моего фонаря и махнул на развалившуюся циновку.
Я сел, стараясь не обращать внимания на вонзившиеся в ноги сухие тростинки.
Старик ухмыльнулся.

--- Надо же, сел, --- хозяин обращался не ко мне, а к невидимому собеседнику, стоявшему у очага.
--- А то заходил тут намедни жрец.
Ситлам зовут.
Жопа у него для циновок чересчур белая, посмотрите на него.
Небось шёлком подтирается и сидит только на северных шкурах.

Я промолчал.

--- Что молчишь?
Нравится нищенское логово?

--- Почему здесь так неуютно? --- спросил я.

Глаза старика на миг сверкнули гневом, но тут же потухли.

--- Неуютно.
Будь Митлам-лехэ моложе, он бы выразился покрепче.
Срач здесь жуткий, вот что.
Курятник, а не дом.
И скоро сюда притащат Сокхэ.
Чем там её кормить надо? --- старик снова обратился к очагу и тут же хлопнул себя ладонью по лбу.
--- Хай, бульончиками, кашками молочными, как же старик не догадался.
Когда из тебя кишку полностью выпускают, ничем другим-то кормить и нельзя --- загнётся дитя.
А откуда зверюшки и молочко, хэ?
Нет зверюшек --- ушли зверюшки, не настрелял их Митлам-лехэ.
Где молочко?
Нет молочка --- не купит его Митлам-лехэ.
Где зерно для кашек?
Нет зерна для кашек --- не вырастил его Митлам-лехэ.

Старик снова разъярился, его глаза загорелись.

--- Новость хорошую он принёс.
На завтрак, думал, пригласит его Митлам-лехэ.
Нечего жрать в этом жилище! --- рявкнул старик, пнув валяющуюся на полу миску.
--- Нет в амбаре ни сладкой муки, ни обычной!
Лучше бы у вас под ножом сдохла эта тупая девица, чем здесь медленно от голода помирать!

--- Но как же, --- пробормотал я.
--- Урожай нынче хорош, люди избыток за бесценок...

--- Да, за бесценок!
Именно!
За бесценок продают! --- старик рассмеялся диким смехом.
--- Кхараму продают, Митхэ продают, дому ар’Хэ продают, а Митламу-лехэ не продают!
И Сокхэ не продают!
Нет семян, дерьмо кроличье сей, Митлам-лехэ!
Нет овощей, бычьи лепёшки гложи, Митлам-лехэ!
Много сортов дерьма-то, богатый стол будет!

--- Но почему?..

--- Какой же он тупой, --- плюнул старик.
Густая желтоватая слюна попала ему на подбородок и на рубаху, но он и не подумал утереться.
--- Разрушитель здесь живёт, понимаешь?
Тхаха.
Разрушитель.
Какое хорошее слово.
Хлопнул Митлам-лехэ по пьяни одного придурка --- и уже что-то там разрушил.
Хотя ничего, кроме жизни ребёнка, не разрушилось.
И разрушили её вы, дерьмоеды, не я.
Вам бы здешних лепёшек отведать.
Теперь каждая падаль, которую Митлам-лехэ один на один со связанными руками придушит, считает нужным его оскорбить.
Да не как-нибудь, а пообиднее да повитиеватее --- кто ещё тебе не ответит?

Я молчал.

--- И жрец этот, Ситлам, --- старик указал очагу куда-то в сторону Тхитрона.
--- Разговаривает так, словно к нему во сне сам Лю явился.
Гордится своей робой, словно со дна моря её руками доставал.
Дай ему волю --- он бы на амулет Сана молиться начал, поклоны бить.
Да он за всю жизнь столько ран не зашил, сколько Митлам-лехэ в одних Молчащих лесах.
Сколько он людей спас?
Много наспасаешь-то, в храме жопу просиживая.

Я молчал.

--- А теперь беги, парень, --- старик вдруг с ловкостью зверя вскочил и выхватил фалангу.
--- Старик голоден, а у него в доме молодой барашек.
Старик может не удержаться.

С его губ сорвался безумный гогот.
Я не двинулся с места.
Лезвие фаланги --- иззубренная сталь, а не кукхватр --- едва держалось на треснувшем древке.

--- Это отвратительно, --- сказал я.

--- Ещё как, парень.
Митлам-лехэ ещё и воняет, --- старик открыл щербатый длиннозубый рот, хрипло выдохнул и снова залился безумным смехом.

--- И тем не менее я хотел бы чаще бывать здесь, --- сказал я.

--- Что он несёт? --- старик ткнул в мою сторону остриём фаланги, обращаясь к невидимому собеседнику.
--- Молодой барашек что-то сказал.
Не пора ли его выпотрошить?

Мой страх прошёл.
Вместо старого разъярённого берсерка я видел несчастного, молящего меня на коленях о помощи.

Кормилица как-то рассказывала, что у хуторских крестьян, как и у восточных соседей-хака, очень тесные родственные отношения.
Это, к сожалению, было чревато как раз такими осложнениями --- отношение к одному человеку проецировалось на всех, кто жил с ним под одной крышей.
Кормилица говорила, что изоляцию мог преодолеть <<навесной мост>> --- уважаемый человек, у которого хватало смелости войти в жилище отверженных.

--- Если дух уймёт своё безумие и даст Митламу-лехэ выслушать меня, то я буду благодарен, --- сказал я.
--- Митлам-лехэ впустил меня в дом не просто так.

Старик замер.

--- У вас будут продукты, материалы и семена, --- сказал я и встал.
--- Я куплю всё сам.
Когда Сокхэ выздоровеет, я предложу ей себя как мужчину.

Стрела попала в цель.
Старик медленно сел, поигрывая фалангой.
Пожевал губы, по-прежнему глядя куда-то в сторону.
Я направился к выходу.

--- Только не вздумай говорить с собой, когда я уйду, --- бросил я напоследок.

Провисшая дверь со скрипом встала на место.
Этот скрип был единственным звуком, которым проводило меня полуразрушенное жилище старика.

\asterism

Кормилица приехала в хутор Самитх в тот же день и собрала всех жителей возле хибарки.
Она при всех обняла воина и поцеловала ему руку, на которой синела старая метка.

--- Если преступник достоин казни, его следует казнить, --- сказала она.
--- Ни один закон не позволяет мучить его и его родичей всю оставшуюся жизнь.

Хуторские починили дом Митлама-лехэ ещё до наступления вечера.
Ему дали пищу и новую одежду, жрецы прислали чистых бинтов и пучки целебных трав.
Сокхэ, которую я вместе с Столбиком перевёз под родную крышу, лежала на чистой постели, улыбалась и пыталась схватить мою руку слабыми пальцами.
Когда мы встречались взглядами, я видел в её глазах такую благодарность и любовь, что мне хотелось бежать --- куда угодно, лишь бы её не видеть.

--- Сокхэ с тобой не будет, --- объявил мне старик, едва я устроил девушку на её месте.
--- Убирайся, и чтобы Митлам-лехэ тебя здесь больше не видел.

--- Её нужно кормить очень часто, --- напомнил я.

--- Не учи охотника шкуру снимать, --- бросил старик.
--- Кормилице передай благодарность.

Я кивнул и вышел вместе со Столбиком из дома.
Дверь больше не скрипела.

--- Удивительная женщина, --- поделился впечатлениями Столбик.
--- Сколько раз здесь бывал, а её не видел.
И ведь без семян и прочих товаров как-то держала дом, да не один год!
Характера ей недостаёт, но выносливая и хозяйка хорошая.
Если бы она не понравилась тебе, я бы за ней поухаживал.

--- Уступаю другу, --- сказал я.
--- Перед твоим обаянием она точно не устоит.

--- Кто знает, я ж не Хри-соблазнитель.
А что так?
Не хочешь возиться с ней полгода?

--- Знаешь же, что не в этом дело, --- поморщился я.

--- Хай, подожди, --- Столбик схватил меня за рукав.
--- Девушка красивая, сильная, она тебе нравится, да и ты её успел очаровать, кормя с ложечки.
И ты просто так уступаешь это сокровище мне?
В чём дело, Лис?

Я покачал головой.

--- После такой жизни её очарует любой, кто отнесётся к ней по-людски.

--- Да ну тебя, какая разница?

--- Большая, дурень ты мой!
Ну и её кормилец, как ты слышал, выставил меня за дверь.
Всё-таки он гораздо лучше будет относиться к крестьянину, чем к жрецу.

--- Ты ж не с её кормильцем детей делать будешь, --- ухмыльнулся Столбик.
\ml{$0$}
{--- Ладно, дело твоё.}
{``Well, suit yourself.}
А вообще это ты дурень.
\ml{$0$}
{Надо использовать силы, данные тебе лесными духами.}
{You should use the powers the Silva Spirits have given to you.''}

Столбик ухватил меня крепкой рукой за достоинство.

--- Отстань, --- я отвесил другу подзатыльник.
\ml{$0$}
{--- Я использую другие силы.}
{``I've got many other powers to use.''}

Остаток пути мы разговаривали о всякой ерунде и смеялись.
Столбик старался обратить свои слова в шутку, но я чувствовал спрятанное в них приятное волнение, словно тонкий аромат, идущий от ещё закрытого бутона.
Всё складывалось для Сокхэ и её кормильца наилучшим образом --- они обрели защитника.
По крайней мере до поры, пока девушка не выздоровеет.

Ужин Кхотлам провела в тяжёлом молчании.
Домашние, чувствуя её недовольство, тоже не решались заговорить.
Все поблагодарили хозяйку за еду и ушли спать, едва закончив трапезу.
Мы с кормилицей остались вдвоём.

\ml{$0$}
{--- Почему ты грустишь? --- спросил я.}
{``Why are you sad?'' I asked.}
\ml{$0$}
{--- Мы сделали хорошее дело.}
{``A good thing has been done.''}

\ml{$0$}
{--- Хорошее? --- буркнула Кхотлам.}
{``Good?'' \Kchotlam\ mumbled.}
\ml{$0$}
{--- Разумеется.}
{``Indeed.}
\ml{$0$}
{Только и занимаюсь тем, что делаю хорошие дела.}
{All I do is do good things.}
\ml{$0$}
{Хотя, если честно, мне надоело напоминать людям, чтобы они были людьми.}
{Though to be fair, I'm sick of reminding humans to be humans.}
\ml{$0$}
{Ты поел, Лисёнок?}
{Are you full, Baby Fox?}
\ml{$0$}
{Иди спать.}
{Go to bed.''}

\section{[@] Оружие и еда}

Историки продолжают собирать данные о царрокх.
Вчера была полностью расшифрована письменность.
Обнаружены зачатки технологий --- царрокх знают, что такое электричество, и умеют им пользоваться.
В некоторых сооружениях найдены примитивные лампы накаливания.
Люди используют для связи простые передатчики на метровых радиоволнах.
Химия у них развита значительно слабее --- не было найдено никаких признаков использования сложных реактивов, производства лекарств и взрывчатых веществ.
Обработка металла также в зачаточном состоянии --- железа, как я уже сказал, на планете не так много.

Историки впервые увидели настоящее оружие для убийства сапиентов.
В Красном музее лежали образцы, мы даже видели некоторые из них, когда шли через Палаты Войны --- мечи, арбалеты, пулевое и волновое оружие, взрывающиеся снаряды, был даже архаичный танк на антигравитационных двигателях, --- но всё это были реплики, изготовленные по обрывкам дошедших до нас описаний.
В племени царрокх оружие, напоминающее дубинку с лезвиями, носит каждый мужчина, и из-за этого к безоружным мужчинам-тси, намного превосходящим любого из аборигенов по силе и ловкости, они относятся с плохо скрываемым презрением.
Во избежание разногласий было решено разработать для тси эффективное оружие ближнего боя и курс владения им.
Этим занялся заскучавший Баночка.

Одним из способов установления отношений у царрокх являются так называемая <<связь крови>> --- нечто вроде биологического родства, которое распространяется почему-то не только на потомков, но и на предков.
Да, это самый удивительный аспект --- если у двух человек появляется совместное потомство, признанное религиозным лидером, то члены их родов начинают относиться друг к другу по-особому.
К сожалению, эти люди относятся к другому биологическому виду, и потомство людей-тси и царрокх без генной инженерии получить невозможно.
Красный-Мак-Под-Кустами хотел заняться устранением репродуктивного барьера --- уж очень заманчивой была перспектива так легко подружиться с племенем.
Но остальные справедливо решили, что проблем у нас и так хватает.

Биологов удивил необычно низкий рост аборигенов.
Впрочем, всё прояснилось быстро: они очень плохо питаются.
Животноводство ограничивается разведением курообразных.
Вокруг каменных городов огромные площади джунглей пускают под топор ради одной-двух посадок, причём полученная древесина используется в основном как топливо.
Ужасная, просто невероятная расточительность.
При этом было подсчитано, что интенсивная обработка почвы без использования сложных технологий позволяет семье иметь поле в тысячу раз меньше --- в год здесь можно собирать по восемь-десять урожаев.
А имеющаяся у нас на корабле кольцевая теплица способна кормить всех царрокх до конца жизни.
Работы много, много.

На сегодняшнем собрании было единогласно решено передать технологии земледелия и животноводства царрокх в понятной для них форме.
Вопрос еды стоять не должен\FM.
\FA{
Вероятно, отсюда пошла знаменитая поговорка сели <<поднимать вопрос еды>> --- мелочиться.
}
Мы обязательно должны подружиться с этими аборигенами --- они обладают большим количеством специальных знаний, необходимых нам для выживания.

\section{[-] Совет перед войной}

--- Да перерезать их, и дело с концом, --- буркнула Ликхэ, скрестив руки на груди.
--- А тех, кто в плен попадёт --- на алтарь.

--- Ликхэ! --- укоризненно воскликнула Хитрам.

--- А ты не хочешь для начала разобраться, в чём дело? --- осведомился Ситрис.
--- Мы предлагаем им помощь в трудное время --- и нам же в спину прилетает копьё.
Это настолько странно, что я бы повременил с выводами.
Хака так не делали никогда.

--- На наши хутора напали, попытались обчистить склады и увели детей для жертвоприношений.
Какие ещё выводы тут могут быть?

В зал вошли Чханэ и Кхарас и сразу же сели за стол.
Подруга схватила лежащий на столе пергамент и протянула его вождю.

--- Вот, смотри.
Как раз то, о чём я говорила.

Кхарас кивнул, взял пергамент и задумчиво потёр голову.

--- Что-нибудь новое? --- осведомился он у Ситриса.

--- В общем-то, всё как обычно, --- пожал плечами Ситрис.
--- Но есть нюанс.

--- Что за нюанс? --- спросила Чханэ.

--- Восточная часть рынка за тобой, --- сказала Ликхэ.
--- Ты видела там хака сегодня?

--- И вчера, и позавчера видела, --- кивнула Чханэ и задумалась, загибая пальцы.
--- Караван с камнями и серебром, четырнадцать человек, в лицо всех помню, имена в блокноте.
Этот забавный маленький торговец специями, Жохас, который называет меня <<мандаринчиком>>.
Два вечно пьяных игрока с ручным питоном --- Кирт и Тибву.
Все на месте.

--- Перед нашествием хака их торговцы обычно \emph{пропадают}, --- пояснил Ситрис.

--- \emph{Пропадают}, --- повторил Эрликх с улыбкой.
--- Проплачивают места на рынке, не выписываются из постоялого двора.
Тем не менее, и их прилавки, и их комнаты почему-то оказываются пустыми.

--- Разумная предосторожность, --- продолжил Ситрис, --- чтобы не попасть под горячую руку.
Караванщики разбивают лагерь за городом, а Жохас с той парочкой просто отсиживаются в укромном месте.
Мы это давно знаем и поэтому успеваем подготовиться.

--- Тайно, --- вставила Кхохо и театрально схватилась за голову:
--- <<Хаяй, как же хорошо, что мы случайно перетащили зерно в городской амбар и оставили недельный запас стрел в Кхихутре!>>

Воины засмеялись;
Кхарас неодобрительно посмотрел на воительницу и снова углубился в бумаги.

--- В общем-то, да, если нас спрашивают, мы обычно списываем всё на счастливый случай и формальные процедуры, --- признал Ситрис.
--- Раскрывать источник информации не стоит.
Когда враги видят, насколько эффективны наши формальные процедуры против их хитрого замысла, азарт быстро пропадает.
Скромность --- не только украшение хорошего воина, но и неплохая защита.
Все таланты, которыми он не обладает, припишут ему побитые враги.

--- А как хака узнают о набеге?

--- Ручные птицы, --- сказала Ликхэ и изобразила руками ласточку.
--- У Жохаса обширная родня на востоке, чуть ли не в трёх племенах сразу.

--- Может быть, сделать ему предупреждение? --- предложила Чханэ.
--- Он может передавать хака сведения об обороне города.
У нас тут и так, --- она бросила взгляд на Кхараса, --- некоторые интересные проблемы с обороной...

--- Не стоит, --- сказал Эрликх.
--- Эти данные --- достояние народа, их может передать кто угодно.
Кроме того, нам нет нужды даже читать его переписку --- тот факт, что она существует, для нас гораздо важнее.

--- То есть, как ты понимаешь, в этот раз нашествие готовил не союз племён, --- заключил Ситрис.
--- И даже, вполне вероятно, не Совет племени.
Иначе маленького проныру хоть кто-то бы предупредил.

--- Частная инициатива? --- удивилась Чханэ.

--- Или крупный разбойничий рейд, --- кивнула Хитрам.

--- То есть это могут быть вообще не хака, а кто-то под них замаскировался?

--- Если так, то на нас напал бродячий театр, --- буркнул Ситрис.
--- Кхотлам говорит, что это хака, а её непросто обмануть.

--- У вас в Предгорье всё достаточно просто, --- улыбнулся Эрликх.
--- Есть мохнатый клыкастый враг, и есть граница.
Нам же большую часть времени приходится заниматься вот такой простенькой дипломатией, а не махать саблей.
Ещё мы привыкли разрешать, а не запрещать, а предупреждениям предпочитаем ненавязчивое наблюдение.
Это часть философии Плюща, для тебя это может быть в новинку.

--- Пока что понятно, --- пожала плечами Чханэ.
--- В конце концов, если присутствия воина недостаточно для порядка, не помогут даже угрозы.

\asterism

Кхарас долго молчал, читая донесения и слушая реплики адептов Храма.
Наконец, когда мнения иссякли, прозвучало его слово:

--- Храм справится.

--- Ну Кхара-ас! --- возмущённо протянула Кхохо.
--- Зачем задействовать весь Нижний этаж?
Возьми шестнадцать человек, по клинку с квартала, дай Ситрису --- и они эту молодёжь раскатают за кхамит!

--- Ну ты это, нормально так всех блох на меня повесила, --- возмутился Ситрис.
--- Я один, по-твоему, должен это разгребать?

--- Вот именно, --- поддержала Хитрам.
--- Если Храм справляется, незачем беспокоить народ.
Идём все вместе.

--- Нам нужен жрец, --- продолжил Кхарас.

Кхохо обиженно замолчала и стала ковырять ногтями стол.
Если вождь сменил тему --- значит, решение по предыдущему вопросу принято и обжалованию не подлежит.

--- Кхарас, мы с Ситламом сегодня по уши, --- сказал Кхатрим.
--- Сам же знаешь, сколько раненых привезли.
Возьми Ликхмаса.

--- Ликхмас?

Я повернулся к Трукхвалу:

--- Учитель, ты меня подменишь в школе?

--- Да, конечно, --- улыбнулся старый жрец.
--- Только не задерживайся, а то книги рассыпятся в пыль.

--- Трукхвал, Ситлам, соберите и проинструктируйте его по основным моментам, --- распорядился Кхарас.
--- Ликхэ, на тебе пайки --- пять рассветов, плюс двадцать единиц для выживших пленников.
Кхохо, замени все факелы --- к нашему уходу в храме должно быть светло.


\section{[@] Обучение Нейросети}

\textspace

Баночка много времени проводил с Нейросетью.
У планта явно были способности к общению с животными --- это отметили все.
Однако вскоре я заметил, что Баночка не просто общается с вороной интонациями --- он начал каркать.

--- Ты раскодировал её сигналы? --- удивился Мак.

--- Потом --- да, --- нетерпеливо отмахнулся Баночка.
--- Но вначале было лень, поэтому я придумал свою систему сигналов, которые ворона может запомнить и воспроизвести.
Так мы и общаемся.

--- Подожди, Баночка, --- нахмурилась Листик.
--- Что значит <<Потом --- да>>?
Ты использовал семантический декодер Дракона и он осилил низкоуровневый язык эволюционно далёкой Ветви?

--- Выучив мой язык, Нейросеть обучила меня своему, --- пожал плечами Баночка.
--- Очевидно же, она достаточно умная.

Птица часто сидела у Баночки на плече, пока он чинил машины.
Он сумел втолковать ей, что делает очень важное дело, и Нейросеть стала ему помогать.
Она перестала таскать детальки.
Баночка сказал птице, что это может навредить машине, а ещё пообещал дать ей аналоги почти любых деталей, которые ей приглянутся.
Как-то в разговоре плант упомянул, что птица относится к машинам как к живым существам и пытается с ними <<говорить>>.

Дальше --- больше.
Птица усвоила цифры и некоторые иероглифы --- Баночка придумал для них специальные звуковые сигналы.
Правда, памяти у неё хватало только на восемь-десять знаков, но для практического применения этого оказалось достаточно --- реестровые номера деталей редко содержат больше.
Баночка написал скрипт для распознавания вороньих цифр;
Нейросеть периодически летала к принтеру, сообщала ему реестровый номер и возвращалась со свеженапечатанной деталью.
Иногда ворона заигрывала с плантом, принося вместо требуемой детали жука или камешек, но в целом Баночка был доволен.
Апогеем обучения Нейросети стала письменность --- Баночка сделал для птицы крохотное писчее перо и научил её вычерчивать цифры.

--- Как у тебя это выходит? --- удивлялся Мак.
--- У ворон высокий интеллект, но не настолько же, чтобы усвоить письменность!

--- Я думаю, что Нейросеть --- гений среди ворон, --- совершенно серьёзно ответил Баночка.

Однако и Мак, и Баночка ошибались.
Ворона оказалась совершенно обычной представительницей своего вида.
Так как Нейросеть продолжала общаться с местными птицами и даже имела среди них авторитет, математические навыки в несколько упрощённом виде разошлись по популяции.
Семьдесят пять процентов ворон в той или иной форме использовали числовые сигналы Баночки при передаче информации.
Но самое удивительное выяснилось, когда вороньи сигналы переняли скворцы, превратив шестнадцатиричную систему счисления в восьмеричную.
Это изумлённая Листик сообщила нам через полтора года.

--- У меня, похоже, наклёвывается диссертация по межвидовой коммуникации, --- сообщила она.

Техники не разделяли восторг Мака по поводу вороны.

--- Ты с ума сошёл, --- буркнул Хомяк, когда Баночка предложил подключить ворону к нейроинтерфейсу.
--- Ты ей ещё репозиторий сделай в хранилище Дракона!
Дружище, ну это всё-таки низшее животное, хоть и очень умное для своего вида!
Что оно будет делать с нейроинтерфейсом?

Однажды Нейросеть пробралась в технический ход и исписала цифрами все поверхности, которые смогла.
Выгоняли хулиганку целый час.
Ворона обиженно отрекомендовала Пирожок как <<глупую драчливую собаку>>.
Чувства вороны оказались взаимными.

\ml{$0$}
{--- Баночка, держи свой мясной привод подальше от систем, --- раздражённо сказала Пирожок, --- а не то я его лично утилизирую под соус!}
{``Flask, keep your meat actuator away from the system,'' Cupcake indignantly said, ``or I'll recycle it with sauce by myself!''}

Впрочем, я отвлёкся.
Я ужасно не хотел бы касаться этой темы, но её не избежать.
Задание, которое я дал Баночке, оказалось роковым --- ещё одной страстью планта стали военные технологии.

\section{[-] Опасная игрушка}

%Плуг 5 дождя 12002, Год Церемонии 20.

\textspace

--- Идём, Ликхмас-тари. Я должен тебе кое-что показать перед походом.

Старик, охая и подтягивая неестественно прямую ногу, зашаркал к большому сундуку в углу.
В это время года нога мучила старика особенно сильно.
Длинный бронзовый ключ на шестнадцать бородок, несколько раз чиркнув по вензелю замка, всё же попал в скважину, два раза музыкально щёлкнул язычок.
Пожилой жрец потянул за ручку и чуть не упал.

--- Ох, проклятая нога...
Парень, поднимай.

Я подошёл и не без труда откинул тяжёлую резную крышку сундука.
Внутри, поверх старых свитков, изорванной в лохмотья одежды, мешочков, колб с непонятным содержимым и прочего хлама лежали три наруча.
Трукхвал указал на них, и я с интересом поднял один.
Обычный, отделанный кожей наруч, без каких-либо украшений, пожалуй, только чересчур тяжёлый.
Такие часто можно было увидеть на воинах.
Я недоумённо посмотрел на старика, а он, улыбаясь во все жёлтые кривые зубы, наблюдал за мной.

--- Учитель Трукхвал, что это?

--- Ооо, мой молодой друг...
Знаешь ли ты, почему воины народа сели одерживают победы в боях с дикарями?

--- Благодаря храбрости и искусству боя, учитель?

--- Не только, Ликхмас.
Ответ ты держишь в своих руках.
Ла, надень.

Я повертел наруч в руках и, недоумевая, натянул его.

--- На левую, Ликхмас.
Ты же чувствуешь --- он тяжеловат.
Вот так.
А теперь смотри --- вот здесь есть выступ.
Надави на него.

Я подчинился.
Выступ спружинил удивительно легко.

\ml{$0$}
{--- А теперь... знаешь ли ты боевые языки народа сели?}
{``Well, now ... do you know battle languages of \Seli?''}

\ml{$0$}
{--- Разумеется, учитель.}
{``Of course, teacher.}
\ml{$0$}
{Птичий, звериный, барабанный язык... но...}
{Language of birds, animals, war drum language ... but---''}

\ml{$0$}
{--- Вот на барабанном мы и остановимся, --- перебил меня Трукхвал.}
{``Let's pause on the last one,'' Tr\`{u}kchu\r{a}l interrupted me.}
\ml{$0$}
{--- Давай-ка, отбей на этом выступе <<проверку>>.}
{``Well, drum the `check' on this button.}
\ml{$0$}
{Только не ударами, а нажатиями.}
{Try to push, not to beat.''}

\ml{$0$}
{--- На этой маленькой бусинке?}
{``This tiny bead?}
\ml{$0$}
{Меня никто не услышит, учитель.}
{I won't be heard, teacher.''}

--- Кто знает, Ликхмас, --- учитель улыбался уже с таким предвкушением, что по подбородку текла тоненькая струйка слюны.

Я пожал плечами и надавил на выступ.
Три коротких, длинный, пауза...
<<Первый, проверка>>.

Внезапно наруч ощутимо завибрировал, и я ошеломлённо разобрал в вибрации знакомый язык:

<<Второй, принял.
Имя?>>

Я защёлкал выступом:

<<Ликхмас.
Второй, имя?>>

Ответ пришёл спустя миг:

<<Кхарас.
Первый, докладывай>>.

Я, открыв рот, смотрел на старика:

--- Но Кхарас сейчас же...

--- ... на площади, --- громким шёпотом закончил учитель.
--- Как я узнал?
Смотри.

Неожиданно изящным, молодцеватым движением Трукхвал вскинул руку и задрал рукав --- там был наруч длиннее раза в два, с чем-то похожим на колокольчик.
Трукхвал крутнулся на каблуках --- и наруч завибрировал, то выше, то ниже, то тише, то громче.

--- Кхарас на площади, Ликхэ в зале, в углу сидит, ты вплотную ко мне.
Ситрис...
Эээ, где он?
Хай...

Учитель растерянно покрутился на месте.

--- Нету...
Хай, над этой штукой ещё потеть и потеть, я её не до конца проработал.
Это для координатора.
Правда, Кхарас сказал, что это ему не нужно --- тяжело и бесполезно, он уже привык ориентироваться по голосам.
Хаяй, ну да ладно, --- Трукхвал улыбнулся и схватил меня за руку.
--- Как тебе, Ликхмас-тари?

Я оторопело разглядывал чудесные наручи то на своей руке, то на учителевой.
Внезапно мой снова завибрировал:

<<Первый, докладывай>>.

Я опомнился:

<<Второй, докладываю.
Всё спокойно, проверяю связь>>.

<<Первый, принял.
Конец связи>>.

Наруч замолчал, и я перевёл дух.

--- Ну как?
Нравится?

--- Да...

--- Возьмёшь его с собой.
В джунглях, когда каждый шорох может стоить тебе жизни, это вещь незаменимая.

--- Но это же... магия... --- я заикался.
--- Как... как... оно...

--- Как оно к нам попало?
Мы делали его здесь, в храме, по древним чертежам.
Может, помнишь, та странная книга.
<<Законы молнии и камня>>, Съешь-Рычащий-Пирожок.
Как-то мы с Хонхо и Хатламом сидели и читали эту книгу, а тут рисунок.
Никаких лишних надписей, только легенда с названиями металлов.
Хатлам сказал: <<А давайте попробуем в точности, как на рисунке, вдруг выйдет чего?>>
Мы собрали восемь штук, как обычно --- ты же знаешь, что-то сгорит, что-то будет собрано неправильно, что-то сломаю я, а что-то потеряется в хламе на складе.
Но мы бы так и не поняли, для чего оно пригодно, если бы не случай --- я случайно унёс одну штуковину с собой, а Хатлам со скуки решил пощёлкать выключателем другой.
Представь мой ужас, когда сумка начала жужжать нехорошими словами!

Трукхвал засмеялся.

--- Вот так мы собрали устройство, протестировали его и придумали замаскировать устройство под наруч.
Теперь им пользуются наши воины и весь Север до самой Ледяной Рыбы.
Разумеется, пока это тайна.

--- Как оно... работает?

--- Ааа... --- протянул старик, и улыбка его стала какой-то несчастной.
--- Я... не знаю, ученик.
И никто не знает.
Мы делаем, как написано, и оно работает.
Древние знали.
Мы не знаем.
Вот, смотри, как оно выглядит внутри...

Пожилой жрец, крякнув, порылся в сундуке и выудил небольшую пластинку --- деревяшка с медной проволокой и небольшими цилиндрами.

--- Здесь медь.
Железо для этой детали непригодно.
А вот эти штуки --- это медная же фольга и пчелиный воск, свёрнутые трубочкой.
Вот это --- кусочек серного свинца, его делают особым способом.
А вот эта коробочка --- самая важная часть.
Внутри --- серебро и солёная вода с бумагой.

--- И всё?!

--- Да, да, ученик!
И это самое удивительное!

--- А что это?
Похоже на змею, обвившуюся вокруг ветки.

--- Это тоже медная проволока, но почему-то её нужно свить именно так.

--- А вот этот камешек?

--- Это не камешек, это...

... Из покоев Трукхвала я вышел в приподнятом настроении.
Чудесный наруч как влитой покоился на моём левом запястье.
Я ощущал его тяжесть.
Старик всю дорогу до Зала косился на меня, и в его глазах была непонятная тоска.

--- Знаешь, Ликхмас, --- заговорил он наконец, --- в книге говорилось, что Древние считали такие вещи... хаяй... простыми и неинтересными.

--- Как можно считать это неинтересным, учитель?
Переговоры без шума, сквозь камень и воду!

--- Однако ж.
Я даже больше тебе скажу --- то, что ты надел на свою руку, было игрушкой для детей.
Понимаешь?
Игрушкой.

--- Понимаю, учитель.

--- Хаяй... Ничего ты не понимаешь, молод ещё.
А я много размышлял над этим, Ликхмас, очень много.
И, кажется, я понял, почему Древние исчезли.

--- Почему, учитель?

--- Потому что вещью, которая способна решить судьбу народов, играли малыши.

\section{[*] Странная просьба}

--- Потому что вы недооценили размер собственного города!
Двух пожаров в нужных местах хватило, чтобы оттянуть народ от самого уязвимого пролёта стены --- и поставить под угрозу захвата десять кварталов!

Митхэ рывком дёрнула карту, и деревянные фигурки со стуком попадали на стол.

--- Да кто ж знал, Митхэ, --- попытался оправдаться старый жрец.

--- Тебя вообще ни в чём не обвиняют, --- успокоила его Эрхэ.
--- Это Нижнему Этажу не мешало бы знать, что в городе полно хака --- торговцев, путешественников, шпионов --- которые не хуже прочих жителей знают, что и как обычно горит!
Куда вы смотрели?
Вы это называете <<северная часть города под контролем>>?
Где были дозорные?

--- Там, где надо! --- рыкнул в ответ боевой вождь.
--- Кто это предвидел?
Вы?

--- Мы хотя бы подоспели вовремя! --- огрызнулась Эрхэ.
--- Согхо ждала на улице Спящей Красотки и разобралась с группой подрывников сразу после сигнала со скалы!

--- То есть то, что ваша воительница случайно оказалалась в любимой таверне на Северо-Востоке, вы называете <<подоспели вовремя?>>

Согхо и Кхотрис медленно начали вставать.
На их лицах застыла ярость, у обоих были перебинтованы руки --- обычное дело после схватки в извилистых узких туннелях старых башен, где кулаки оказываются предпочтительнее даже ножа.

--- Объясни-ка свои слова, любезнейший, --- потребовала Согхо.

--- Я не обязан вам ничего объяснять.
После того, как пироги сгорели, все норовят обругать дрова!

--- Я привыкла доверять Храму города, в котором нахожусь! --- вспылила Митхэ.
--- И моя ошибка в этом случае бесчестит вас не меньше, чем меня!

--- Я полагаю, что нам следует использовать её опыт.

Все обернулись.
Король-жрец стоял в дверях зала в парадном головном уборе и при оружии.

--- Митхэ, я признаю, что мы просчитались, --- Король-жрец сразу перешёл к делу.
--- Я не хотел тревожить тебя просьбами, но на этот раз нам действительно нужна твоя помощь.
Ты получишь всё необходимое.
Своё разрешение я хочу подтвердить, отдав тебе Бриз Битвы.

Король-жрец снял с руки увесистое кольцо с печатью и протянул Митхэ.

--- Эту печать знают во всём городе.
Когда всё закончится, ты можешь оставить её на память.

--- Я в курсе, что такое Бриз Битвы, --- Митхэ повертела кольцо в руках.
--- Этого недостаточно для оплаты.

--- Да как ты смеешь! --- прорычал вождь.
--- Это знак военной власти Трёхэтажного Храма!

--- Знак военной власти, который бессилен в моих руках, --- парировала Митхэ.
--- Пока Король-жрец желает, чтобы я выполняла его поручение, знак будет действовать.
Как только мои услуги окажутся не нужны --- Бриз Битвы превратится в обычное кольцо, которым он и является.
Я возьму его в знак моего согласия.
Но оплата не отменяется.

--- Мы уже отдали вам шестую часть храмового золота, Митхэ, --- вкрадчиво сказал Первый жрец.
--- Если нагрузить им весь отряд чести, вы не унесёте всё.

--- Золото не понадобится никуда нести.
Мы потратили девять десятых на починку снаряжения, оленей и повозки.
Остального едва хватит на провиант.

--- Это не наша проблема.

--- Вы просите моей помощи, а потом заявляете, что наша экипировка --- не ваша проблема? --- усмехнулась Митхэ.
--- Большая часть наших повозок осталась где-то на дороге.
Доспехи латаные-перелатаные.
Воительница, которая сегодня спасла ваш город, сражается кустарным копьём, которое ей сделал любовник, потому что у неё нет денег на сталь!

--- Золото, перестань, --- пробормотала Согхо.

--- Я не стесняюсь своей бедности, Согхо, и тебе не советую.
Заплатите нам ещё тысячу перьев.

--- Мы не можем себе этого позволить, и ты это знаешь.

--- Это ещё одна десятая храмового золота, --- непринуждённо сказал Ситрис, показав какие-то вычисления на кусочке пергамента.
--- Как вы знаете, неприступность очень важна для статуса города.
Если вы считаете, что тысяча перьев --- это много, Тхартхаахитр может распрощаться со статусом столицы.
Я всегда думал, что тот же Тхитрон красивее и гораздо безопаснее.

Жрецы и воины Трёхэтажного Храма заорали на него почти одновременно.
Их лица исказила ярость.

<<Засранец, --- с восторгом подумала Митхэ.
--- Надо его повысить>>.

Король-жрец поднял вверх кулак, немедленно оборвав поток бессвязных слов.

--- Хорошо, --- ответил он.
--- У тебя есть ещё какие-то условия, Митхэ ар'Кахр?

--- Да, есть.
Первое --- мне предоставят право голоса на переговорах, буде таковые случатся.
И второе --- кутрапы, которые сражаются в моём отряде чести, получат полную амнистию в землях сели.

\section{[-] Местоположение воина}

--- Я не уверен, что я тебя понял правильно, --- поморщился Ситрис.
--- О каких кутрапах речь?

--- Я знаю, что это звучит странно, --- заговорила Чханэ.
--- Но у меня есть опыт диверсий, и мы так делали пару раз.
Это называется <<горшочек с мёдом>>.
В любом городе есть скрывающиеся кутрапы, я уверена, вы даже знаете, кто и где именно.
Предлагаете им задание: смешаться с местными жителями, попасть в плен, сообщить хака нужные данные.

--- То есть я правильно понял, что ты предлагаешь сознательно --- специально! --- направить отряд захватчиков в один из хуторов, используя невинного человека?

--- Не такого уж и невинного, Ситрис.
И не одного, а нескольких.
Таким образом отвлекаем основные силы, освобождаем пленников --- и уходим в город.
Все жители земель Тхитрона к тому времени будут в безопасности.
Оставшимся хака просто придётся уйти восвояси с пустыми руками.

--- Они увезут с собой награбленное, --- напомнила Ликхэ.

--- Маленький отряд много не увезёт.
Это не пылерои, это хака, они меньше, слабее и медленнее.
Их основной целью были пленники, а пленников мы отобьём.

--- А если твоих агентов просто убьют до того, как они передадут данные хака?

--- Такого ещё не бывало, --- ухмыльнулась Чханэ.
--- Обычно мы предлагали им полную амнистию в случае, если они выполнят задание и вернутся живыми.
За амнистию кутрапы сражаются как ягуары, особенно после нескольких лет жизни вне закона.

--- Мы не будем использовать человека как наживку, --- отрезал Ситрис.
--- Даже кутрапа.
И хватит об этом.

\textspace

--- Ситрис, Трукхвал мне жаловался, что ты постоянно отключаешь наруч и он не может тебя найти, --- сказала Ликхэ.

--- Трукхвал уже надоел со своими экспериментами, --- проворчал Ситрис.
--- Сама идея наруча великолепна, но следить за моим местоположением без моего ведома я ему не дам!

--- А тебе есть что скрывать, Ситрис? --- вскинулась Ликхэ.

--- Да, есть.
Мои дела, которые касаются только тех, кого я хочу!

--- Соглашусь с Ситрисом, --- заявила Хитрам.
--- Воин на то и воин, чтобы его не могли просто обнаружить --- в том числе и союзники.
Сам по себе наруч --- копьё с двумя клинками\FM, и я об этом уже говорила Трукхвалу.
\FA{
Оружие, опасное и для противника, и для того, кто им пользуется.
}
Но он помешался на своей технике.

\section{[@] Великий воин}

Началось это помешательство сравнительно безобидно.
Баночка связался со мной и попросил прийти в палатку, которую он разбил на берегу моря.
Пока плант работал, Нейросеть искала на берегу жирных улиток и дробила их раковины о камни.

Когда я вошёл в палатку, там оказался ещё и случайно забежавший Фонтанчик.

--- Ты посмотри на него, --- не оборачиваясь на меня, неодобрительно сказал канин.

Баночка стоял посреди палатки.
На нём красовался богато украшенный гравировкой доспех из стабитаниума.
На поясе у планта висело нечто, напоминающее меч и кинжалы.

--- Я проработал все суставы, --- гордо сказал техник из-под забрала.
--- Теперь аборигены точно будут вас уважать --- в стыки не войдёт даже микрозонд.

Баночка один за другим бросил кинжалы в стоящую посреди палатки куклу, набитую песком.
Клинки лениво нарисовали в воздухе изящные петли, резко набрали скорость и по самую рукоять вошли в песок.
Затем плант с воинственным криком выхватил меч и, подскочив к кукле, разрубил её пополам.

--- Ну как? --- сияя улыбкой, спросил он.

--- У нас нет столько стабитаниума, чтобы одеть всех, --- тихо напомнил я другу.

--- И я это носить не стану, --- заявил Фонтанчик.

Улыбка Баночки увяла.

--- Вы сами попросили оружие и...

--- Мы просили обычные заострённые палки, которые можно сделать из сухой древесины, а не замаскированный под меч резонанс-резак, --- сказал я.
--- У оружия назначение --- причинять вред мягкому, сделанному из плоти существу, слегка прикрытому одеждой и дублёной кожей.
А твоим мечом можно камни рубить.

--- Да и самонаводящиеся самообучающиеся кинжалы --- это опасно в том числе и для тебя, --- поддержал меня Фонтанчик.
--- У меня аж сердце замерло, когда я понял, что у них лезвия с аэродинамическими механизмами, и явно не императивно управляемыми.
Сразу вспомнил тех роботов на Тси-Ди...

--- Так это же хорошо! --- запротестовал Баночка.
--- Чем лучше оружие, тем...

--- Нейросеть глюканёт --- и без головы останется кто-то из нас, --- перебил его я.
--- Я против оружия с искусственным интеллектом.
Тем более это альфа-версия.
Если уж тебе нравится метательное оружие, пусть снаряды хотя бы летят по предсказуемой траектории, чтобы мы в случае чего успели увернуться.
Ну или чтобы получившаяся травма хотя бы входила в медицинскую классификацию травм.

\ml{$0$}
{--- Но ведь...}
{``But---''}

--- Дружище, --- начал Фонтанчик, --- сделай нам красивые палки с обычными, понимаешь, обычными лезвиями.
Научи этим пользоваться --- отбивать вражеские заострённые палки.
Потом возьми материал, который используют аборигены, и сделай нечто похожее на их доспехи, чтобы жизненно важные органы были чуть-чуть прикрыты.
\ml{$0$}
{Неважно, что по сути они ни от чего не защищают, главное, чтобы было.}
{It doesn't matter if they're useless, all that matters is their existence.}
Этого хватит для установления уважительных отношений.
Мы всё равно не собираемся с ними драться.

\ml{$0$}
{--- А если придётся?}
{``What if we have to?''}

\ml{$0$}
{--- Ты их боишься? --- удивился Фонтанчик.}
{``Are you afraid of them?'' Fountain looked surprised.}

\ml{$0$}
{--- Да, --- кивнул Баночка.}
{``Yes,'' Flask nodded.}
\ml{$0$}
{--- Я боюсь даже не того, что они большие и агрессивные --- в принципе, я буду даже посильнее многих, и защититься смогу.}
{``I'm not afraid of their size or aggressiveness---as a matter of fact, I'm stronger than most of them, and I can defend myself.}
\ml{$0$}
{Но беспокоит мысль, что они могут убить просто потому, что я им не нравлюсь, и сделать это исподтишка.}
{One thought makes me nervous: they can kill me just because they dislike me, and they can stab me in the back.''}

--- Я буду посещать царрокх вместе с тобой, --- ободряюще сказал Фонтанчик.
--- Со мной тебя точно никто не посмеет тронуть.

--- Хорошо, --- грустно сказал Баночка.
--- Однако элементарные средства защиты нужны.
Давайте я хотя бы сделаю для тси одежду из полимерного волокна.

--- Она защитит только от рубящих ударов, --- покачал головой я.
--- Прямым ударом копьё пробьёт любую одежду без твёрдых вставок.
Займись лучше техникой боя, умение защитит лучше любых костюмов.

--- Хорошо, --- повторил Баночка.
--- Есть ещё вариант сделать высокоскоростное метательное оружие, звуковое, химическое или...

--- Исключено, --- тут же замахал руками Фонтанчик.
--- Во-первых, аборигены не сразу поймут, что это оружие, и нам придётся кого-то покалечить.
Во-вторых, превосходящие военные технологии заставят их бояться.
Это нам не нужно.

--- Совсем не нужно, --- согласился Баночка.
--- А с этим что?

--- А это ты будешь носить, --- оскалился канин.
--- Мы будем тебе поклоняться, как самому великому воину.

--- Ещё чего не хватало! --- возмутился Баночка.
--- Ладно, пущу стабитаниум в дело.

--- Оставь, Баночка, --- замахал я руками.
--- Вышло очень красиво.
Правда.
Когда-нибудь здесь будет музей, и мы обязательно выставим там твоё творение.

\section{[-] Сгоревшие глаза}

\textspace

Эрликх, Кхарас и Хитрам сидели на ступеньках.
С площади прибежал Ситрис.

--- Мы узнали, где они разбили лагерь, --- сказал воин.
--- Красная скала.
Судя по словам разведчика, очень грамотно, незамеченными подойти не удастся.

--- Тогда ждём до утра, --- сказал Кхарас.
--- Есть вести из Тхисема?

--- Они по-прежнему не желают эвакуироваться.

Эрликх и Хитрам одновременно застонали.

--- Идолы раздери этот Тхисем...

--- Опять, жопой яму копать...

--- Ты им сказал, сколько хака на этот раз? --- осведомился Кхарас.

--- Передал всё, что необходимо.
Они сказали, что любезно задержат налётчиков до нашего прихода.

--- Как же, --- буркнул Эрликх.
--- К нашему приходу там будет пепелище.
Их бы и в прошлый раз перебили, если бы я не подоспел с отрядом.
С другой стороны, мы не кормильцы --- уговаривать дитя кашу съесть.

--- Может, снова послать к ним отряд? --- задумчиво сказала Хитрам.

--- И снова положить пару десятков человек зазря? --- развёл руками Эрликх.
--- Хака \emph{привыкли} драться в поселениях, это их стихия.
Если на дороге мы можем покосить весь отряд, как травку, то в городском бою мы размениваем одного на четверых.
Сорок хака --- десять сели, и это ещё очень неплохой размен.
Нет уж.
Хотят эти крестьяне умереть воинской смертью --- могу только посочувствовать.
Зато в следующий раз эвакуироваться будут все, без лишних вопросов.

Кхарас задумался.

--- Сообщи всем, что сбор переносится на кхамит раньше.
Может, успеем перехватить хака по дороге на Тхисем.

Ситрис кивнул и скрылся в темноте.

В кустах застрекотали сверчки.
Послышался оживлённый рассказ парня и женский смех.
На площади миловалась молодая пара чуть старше меня.

--- Красивая пара, да? --- Хитрам кивнула на стоящих.
--- Паренёчек такой милый, круглощёкий.
И девушка --- почти ребёнок.
Чистые наивные глазки.
На лицах ни тревог, ни потерь, ни войн.
Только ветер, только радость впереди.

Я почувствовал в её голосе лёгкую зависть.

--- Наверное, им следует узнать всё это, --- пожал плечами Эрликх.
--- Это все узнают со временем...

--- Не следует, --- вдруг громыхнул Кхарас.
Все замолчали и посмотрели на вождя.

--- Слушай, Кхарас, --- попытался возразить Эрликх, --- они живут в жестоком мире, и...

--- Мы для того здесь и есть, чтобы дети оставались детьми, --- угрюмо пробурчал Кхарас.
--- Мы.
Храм и Двор.
И если это в наших силах, пусть эта пара, и их потомки, и потомки их потомков никогда не познают горя.

--- Мне так нравится, когда ты начинаешь говорить о долге, --- сказала Хитрам.
--- Напоминает мне тебя молодого.
Когда шла речь о несправедливости, когда где-то сильный обижал слабого, твои глаза загорались ярким пламенем негодования, способным прожечь горы и иссушить моря.

\ml{$0$}
{--- С возрастом Кхарас стал спокойнее, да, --- отметил Эрликх.}
{``\Kcharas\ has become calmer with age, yeah,'' O\r{e}rl\'{\i}kch observed.}

\ml{$0$}
{--- Нет, --- покачала головой Хитрам.}
{``Nah,'' \Chitram\ shook her head.}
\ml{$0$}
{--- Просто несправедливости стало столько, что его глаза сгорели дотла.}
{``There's actually so much injustice his eyes have burned to ashes.''}

Пара на площади, тихо смеясь, побежала куда-то за склады.
Вскоре осколки смеха окончательно раздробились о стены и брусчатку.
Наступила ночь.

--- Тебе принести что-нибудь согревающее? --- ласково спросила Хитрам.

Кхарас угрюмо промолчал.

--- Сейчас, --- Хитрам поднялась и пошла в храм так резво, словно вождь сказал, что умирает от жажды.

\section{[*] Обещание Короля}

\epigraph
{Прежде суд, долго сохранявший черты религиозного института, требовал от человека раскаяния.
В случае раскаяния приговор могли смягчить, а то и вовсе отменить.
И сейчас я скажу --- это самая большая ошибка законников прежних веков.
Если честный человек убил случайно, для него не найти худшего палача, чем он сам, и худшей тюрьмы, нежели тюрьма его мыслей.
Если же честный человек мстил или защищался, то он поступил по своей совести;
в чём ему следует раскаиваться?
Раскаяние унижает человека, страх наказания заставляет его вступить на путь лжи;
униженный и запуганный лжец --- хороший раб, но плохой гражданин.}
{Анатолиу Тиу.
<<Земной суд>>}

\textspace

Зал был набит до отказа, не только храмовниками, но и членами городских Советов.
Многие перешёптывались, обсуждая последние события;
Митхэ отметила неприязненный тон.

--- Итак, начнём, --- Король-жрец кашлянул и поднял руку, призывая собравшихся к тишине.
--- Сегодня мы решаем вопрос относительно, хай, требований, которое выдвинула Митхэ ар'Кахр в обмен на планирование обороны города.
Первое требование: она будет участвовать в переговорах с хака.
Храм считает, что требование обоснованное и должно быть удовлетворено.
Возражения?
Да, Совет Цеха, слушаю.

--- Почему глава Нижнего Этажа должна принимать участие в переговорах в иной роли, нежели охрана Короля-жреца? --- поинтересовался низкорослый мужчина.

--- Вопрос хороший.
Митхэ, можешь пояснить своё требование?

--- Ситуация уже осложнена до предела.
Хака дошли до столицы.
Моё присутствие в качестве переговорщика покажет союзу племён серьёзность намерений сели.
Хака должны понять, что их ждёт в случае невыполнения обязательств.

--- Справедливо и разумно, --- поклонился старшина Цеха.

--- Храм придерживается того же мнения.
Решение принято.
На очереди обсуждение второго требования Митхэ ар'Кахр: амнистия кутрапов, которые состоят в её отряде чести.

Зал затих.

--- Я прекрасно понимаю ваше недоумение.
До меня доходили слухи, что в отряде чести Митхэ ар'Кахр есть кутрапы.
Но, если честно, я не ожидал, что глава Бродячего Храма заявит об этом во всеуслышание.
Если бы в Трёх Этажах обнаружился преступник, моего авторитета бы не хватило, чтобы смыть с репутации Храма такое пятно.

--- С \emph{моим} авторитетом я могу себе позволить некоторые вольности, --- парировала Митхэ.
По залу пронёсся смешок.
--- Я верю во второй шанс.

--- Смею заметить, что для большинства кутрапов твой шанс далеко не второй.
Впрочем, это внутренние дела отряда чести и меня не касаются.
Кто хочет высказаться?

--- Совет Цеха против, --- тут же поднялись в зале несколько рук.

--- Совет Сада против, --- повторили ещё несколько голосов.

Король-жрец развёл руками.

--- Митхэ, Храм тоже не в восторге.
Ты хочешь, чтобы я своей властью отменил решение городского Совета, и даже не одно решение, и даже не одного Совета, и даже не одного города.
Это серьёзный шаг.
Думаю, ты поймёшь меня в моём желании получить некоторые гарантии.

--- Какие именно гарантии хочет Король-жрец?

--- Во-первых, я предпочёл бы любую разумную альтернативу твоему требованию.
Частичная амнистия для нескольких кутрапов, например, только для земель Тхартхаахитра.
Либо полная амнистия для одного кутрапа.
Если альтернатива невозможна, должно быть во всеуслышание объявлено, что ты не оставила мне иного выбора, кроме как принять твои условия.
Думаю, здесь твой авторитет сыграет на руку.

--- Скользкая тварь, --- усмехнулся Акхсар --- так, что его услышали только Эрхэ и Митхэ.

--- Считай, что это ультиматум, --- кивнула Митхэ.

--- Митхэ, ты же видишь, что большинство против, --- с укором сказала представительница Сада.
--- Король-жрец прав.
Почему нельзя выдвинуть другое требование, менее... радикальное?

--- Я дала клятву этим людям, что сделаю всё возможное для их блага в обмен на добросовестную службу, --- с достоинством ответила Митхэ.
--- Они выполнили свою часть уговора.
Так как все трое намерены в дальнейшем остаться в моём отряде чести, а не поселиться в Тхартхаахитре, частичная амнистия для них бесполезна.
Комментировать возможность амнистии только для одного из троих не буду, речь идёт о человеческих жизнях и подобные торги я считаю совершенно неприемлемыми.

--- Я хочу выслушать этих кутрапов, --- сказал Король-жрец.
--- Ты сообщила их количество, но ничего не сказала об их качестве.

--- Я бы хотела сохранить их имена в тайне, во избежание...

--- Зачем ты просишь моей помощи, Митхэ ар'Кахр, если между нами нет доверия?

Митхэ замялась и замолчала.

Ситуацию спас Аурвелий.
Старик вышел в центр зала и учтиво прикоснулся мизинцем к губам.

--- Судя по приветствию, ты ноа с Могильного берега, --- кивнул Король-жрец и повторил жест.
--- Как тебя зовут, сенвиор?

\ml{$0$}
{--- Аурвелий Амвросий, --- прогудел старик.}
{``Aurweli Amwurosi,'' the old man said sonorously.}

--- Душистая Золотая Полынь, --- перевёл Король-жрец.
--- Красивое имя.
Почему бы тебе не назваться, как сели?

--- Я взять от джунгли достаточно для себвя, --- с достоинством ответил старик.

--- Так почему ты стал кутрапом?

\ml{$0$}
{--- Я всю жизнь быть пвират, Король-жрец.}
{``All my life I pwirate, Priest-kwing.}
\ml{$0$}
{В молодость я быть пвират с моё удовольствие;}
{Young I pwirate with my pleasure;}
\ml{$0$}
{средний я быть пвират из-за неквуда идти.}
{middle I pwirate bwecause nowhere go to.}
\ml{$0$}
{Сейчас я стар и устать.}
{Now I old and twired.''}

Жрецы возмущённо забормотали.

--- То есть ты всю жизнь грабил и убивал, а теперь требуешь от нас спокойной старости, ноа? --- осведомился один из них.

--- Бвольшая часть жизни я делвать это бвез ввыбора, --- спокойно ответил Амвросий.
--- Ввыбор появвиться пять дождь как, капита Миция слушать мвеня, отличие другвой.
\ml{$0$}
{Я не имветь золото и кукхватр, я не имветь дом и семья, я не иметь зелёный парус с вытканным солнцем\FM, как подобает мужчине.}
{I don't have gwold and \kukchuatr, I don't have home and ffamily, I don't have a gween sail of woven sun\FM, as a man should do.}
\FA{
Красивое крепкое судно.
}
Я ничего требвовать, я хотеть служить для покой.

--- Мы тебя поняли, сенвиор Амвросий, --- кивнул Король-жрец.
--- Что насчёт оставшихся двоих?
Будут ли они смелее старого ноа?

--- Король-жрец, не смей оскорблять моих воинов при мне, --- на грани шёпота произнесла Митхэ.

--- Я не оскорбиться, капита Миция, --- поднял ладони Аурвелий.
--- В мой годы поквой важнее чей-то плохие словва.

Вперёд вышел Ситрис.

--- Ситрис ар'Эр э'Тхинат, разбойник.

Жрецы снисходительно улыбнулись.

--- Это всё, что ты хотел нам сказать? --- осведомился Король-жрец.
--- Сенвиор Амвросий хотя бы рассказал, почему он желает снять с себя звание кутрапа.

--- Это вас не касается, --- отрезал Ситрис, вызвав новую волну возмущённого шёпота.
--- Есть товар, и есть покупатель.
Подробности деяний и чувства --- моё личное дело, и если мне не изменяет память, перед сидящими в зале я чист.

--- Мы тебя поняли, Ситрис ар'Эр, в твоих словах есть истина, --- холодно сказал Король-жрец.
--- Кто же третий?

Ликхлам вышла вперёд и нерешительно потопталась на месте.

--- Ликхлам ар'Сатр, --- тихо сказала она.
--- Я люблю мужчину --- уроженца Тхартхаахитра, но не могу жить с ним в одном доме.

--- Вот кого устроила бы частичная амнистия, да, Митхэ? --- весело заметил один из жрецов.
Прочие ухмыльнулись.

--- Я кутрап-насильник, воровала оружие и золотые вещи, --- продолжала Ликхлам.
--- Я не Разрушитель и никогда им не была.
Я помню всех, кого убила в своей жизни, и могу посчитать их пальцами одной руки --- один пылерой, трое Молчащих идолов и человек, который пытался меня ограбить...

Ликхлам подняла растопыренную ладонь, по которой стекали крупные грязноватые капли пота.

--- Достаточно, --- Король-жрец остановил Ликхлам и обратился к Митхэ:
--- Интересно, как исполняет воинский долг женщина, которая может одной рукой сосчитать убитых ею врагов?

\ml{$0$}
{--- Она готовит пищу и содержит в порядке упряжь, --- с достоинством ответила Митхэ, заглушив не очень-то вежливый шёпот наёмников:}
{``She cooks food and keeps tack in order,'' \Mitchoe\ answered with dignity, trying to drown out mercenaries' whispers which were far from being polite:}

\ml{$0$}
{--- Да какая тебе разница?}
{``How the fuck is it your business?''}

\ml{$0$}
{--- А сам-то ты много настрелял?}
{``Ya have killed more, haven't ya?''}

\ml{$0$}
{--- Следи за своими курами, петушок.}
{``Look after your own chicks, Baby Cock.''}

\ml{$0$}
{Жрецы возмущённо забормотали что-то о <<гостеприимстве>>;}
{Priests murmured indignantly something about `hospitality';}
Король-жрец, если и услышал, то не подал вида.

--- Что ж, Митхэ ар'Кахр, я вижу, что \emph{почти все} кутрапы готовы стать частью народа.
\ml{$0$}
{Даю тебе слово: я сделаю всё от меня зависящее, чтобы они смогли жить там, где захотят.}
{You have my word: I shall do my best to ensure that all of them could live wherever they desire.}
\ml{$0$}
{А теперь следуй за мной, нам нужно обсудить стратегию.}
{Now follow me, we have a strategy to discuss.''}
Заседание окончено.

Король-жрец поднялся и властно поманил Митхэ за собой.

--- Тхэай, не нравятся мне эти данные вскользь обещания, --- пробормотал Ситрис в сторону.

\ml{$0$}
{--- Молодые ноги хорошо бегать, Ситрисиас, --- улыбнулся Аурвелий.}
{``Yawng feet run fast, Sitrisias,'' Aurweli smiled.}

Ситрису вдруг ужасно захотелось обнять старика.

<<Если Король-жрец лжёт, Аурвелий, ты пойдёшь со мной.
Не думай, что усталость --- оправдание для покорности судьбе>>.

\section{[@] Нейросеть. Туда}

\textspace

Нейросеть беспокоилась весь день.
Баночка, поговорив с ней, вдруг заплакал и сразу после завтрака вышел с вороной на воздух.
Мы, недоумевая, последовали за ними.

--- Ты правда этого хочешь, подруга моя хорошая? --- ласково спросил Баночка и что-то прокурлыкал.

Ворона ответила нежным карканьем, попрыгала и покрутилась у него на руке.

\ml{$0$}
{--- Я буду ждать тебя до смерти как друг, --- прошептал Баночка.}
{``I'll wait for you till my death as a friend,'' Flask whispered.}
\ml{$0$}
{--- И ещё пару дней как пища.}
{``And two more days as a food.''}

Он поцеловал птицу в клюв и подбросил в воздух.
Мы зачарованно смотрели, как ворона делает величественный круг над лагерем и исчезает вдали, в утренней дымке.
Наконец Баночка, грустно улыбаясь, повернулся к нам:

\ml{$0$}
{--- Гнездование.}
{``Nesting.}
\ml{$0$}
{Инстинкт.}
{Instinct.}
\ml{$0$}
{Пусть летит.}
{Let her go.''}

\section{[-] Обряд}

%Плуг 6 дождя 12002, Год Церемонии 20.

Вышел отряд затемно.
Солнца ещё не было, но в воздухе уже чувствовалась утренняя дымка.
Крики ночных птиц, квакш и мраморных цикад затихали, изредка робко подавали голос птички дневные.
Мы миновали несколько улиц.
Дома стояли чёрные, сонные, в окнах их не было ни огня.
В птичниках шуршали крыльями и стучали жёсткими лапками спящие куры.
Лишь на улице Дышащего Дерева мучимая бессонницей старушка, молча попыхивая едкой травой тумана, проводила нас взглядом со своего деревянного кресла.
Воины кивнули ей, и старушка ответила неторопливым кивком.

Мы шли молча.
Тишину прерывали лишь шуршание рубах, скрежет ремней и лёгкое позвякивание чьих-то плохо пригнанных лезвий фаланги.
Наконец отряд дошёл до небольшой площадки на окраине города.
Воздух, словно застывший в центре города, вдруг сменился лёгким живым ветром.
На меня пахнуло сильным ароматом прелой листвы и цветов.

--- Пришли, --- шёпотом сказала Ликхэ мне на ухо.

Воины так же молча, с торжественной серьёзностью начали снимать с себя мешки и доспехи.

Круг Доверия.
Обряд, который совершают во время войны или походов.

--- Передохни, Ликхмас-тари, --- шепнул мне кто-то из темноты.
--- Пока что мы никуда не идём.

Кхарас тем временем быстро и умело развёл костёр, и я смог оглядеться.

Все, кроме меня и нескольких человек, уже успели раздеться до пояса.
Чханэ, улыбаясь и потирая ладонью красивый, натёртый поясом живот, подошла ко мне и игриво потянула за шнурки горловины:

--- Лис, ты в одежде обряд собрался совершать?
Давай помогу.

Подруга развязала шнурки и, смеясь, стащила с меня рубаху.
Ликхэ тем временем, ревниво глядя на нас, сняла передник, открыв большие, похожие на зрелые дыни груди.
Благодаря им воительницу часто преследовали восхищённые взгляды мужчин, но в Храме говорили, что этот подарок Безумных может стоить ей жизни.
Удобный укреплённый передник, в котором богатство носилось, для Ликхэ сшила моя кормилица, которая считалась на необычное шитьё большой мастерицей.

Кхарас подбросил пламени сухих сучьев, и костёр заревел, как старый ягуар-узурпатор.
Искры летели вверх, к небесам, и мы с Чханэ зачарованно уставились в это небо, бездонное, полное звёзд.
Прошло несколько михнет, прежде чем я понял, что нашему примеру последовал весь отряд.

--- Ликхмас-тари, что там? --- шепотом поинтересовалась Кхохо.
Воительница впервые обратилась ко мне с уважительным суффиксом.

--- Звёзды, --- пожал я плечами.

Кто-то засмеялся.

--- Так, начинаем, --- громко сказал Кхарас, и словно в подтверждение его слов в огне лопнуло полено, выбросив в небо сноп искр.
Шёпот и смешки утихли.
--- Хай, у нас новенькие...
Ликхмас, ты участвовал в Круге?

--- Нет, --- ответил я.
--- Я даже на Мягкие Руки так ни разу и не попал --- ни жрецом, ни участником.

--- Хай, --- понял Кхарас.
--- Припоминаю, нам с ними последние двадцать дождей не везло.
То одно, то другое.
Тханэ?

--- С двадцати пяти дождей участвую.

--- Что ж, тогда мы начнём, а ты объяснишь Ликхмасу по ходу дела.

Чханэ серьёзно кивнула.
Потом, едва сдерживая смех, шепнула мне на ухо:

--- Кхарас такой важный, как варан в брачный период.
Просто расслабься и делай то же, что и все.
И не вздумай фальшивить, за такое побить могут.

--- Фальшивить?..

--- Тихо.

Отряд сел вокруг костра и Кхараса, сцепившись кончиками пальцев на расслабленных руках.
Отстранённые лица осветил костёр. Кхарас встал лицом к огню, бросив на нескольких участников колеблющуюся, непроницаемо чёрную тень.
Рот его открылся, глаза закатились.

--- Ааа... --- разнёсся в ночной тиши его низкий голос.
Кхарас довольно долго тянул ноту, пока её не подхватили остальные.

--- Ммм...

--- Ааа...

--- Ооо...

Разными голосами воины тянули ноты одного и того же мажорного трезвучия.
Я различил в общем хоре хрипловатый бас Ситриса, приятный контральто Чханэ и сладковатый альт Ликхэ...
Вскоре пришла и моя очередь.

--- Ааа...

--- Ооо...

Кто-то достал барабан и варган, голос обрёл ритм.

Та-та-та-хум!
Тхийо-тхо-тхоу!

Та-та-та-хум!
Йолль!
Йоуу!

Костёр горёл, его пламя плясало свой вечный танец.
Потянуло странным, слегка удушливым дымом.
Я вдохнул его полной грудью.
Удивительно, но он не вызывал желания кашлять.
Чужие голоса раздували меня, как ночной ветер океана раздувает парус.
Время поплыло, звёзды вдруг стали близкими и невообразимо горячими.
Я почувствовал себя всеми остальными --- Кхарасом, в экстазе тянущим ноту и отбивающим на собственных коленях ритм;
Кхохо, закрывшей глаза и качающейся подобно дереву на ветру;
варганщиком Эрликхом, который извивался на земле, не прекращая отстукивать на крохотном инструменте животворящий, согревающий кровь ритм...
Надо мной в ночном небе загорелись глаза Чханэ.
Что это --- реальность или моя фантазия, неотличимая от реальности?..

--- Ааа...

--- Ооо...

--- Ммм...

Та-та-та-хум!
Тхийо-тхо-тхоу!

Та-та-та-хум!
Йолль!
Йоуу!

Вдруг Ситрис, не прекращая извиваться, схватил Чханэ за руку и притянул к себе.
Она, продолжая биться в том же ритме, впилась губами в его губы.
Вскоре они уже сплелись воедино...

--- Лисичка... --- сияющие зелёные глаза Ликхэ очутились прямо передо мной.
Воительница набросилась на меня сверху, как голодный ягуар на добычу.
Рука её обхватила мой стебель, и я почувствовал, как вхожу в её горячее нутро.
Словно откуда-то издалека раздался мой собственный стон, долгий, монотонный, идеально входящий в звенящее в моей голове мажорное трезвучие...

Та-та-та-хум!
Тхийо-тхо-тхоу!

И снова --- глаза Чханэ среди звёзд...
Огромные, на полнеба, огненно-оранжевые глаза.
Я, сделав усилие, повернул голову, пытаясь найти среди танца обнажённых бронзовых тел подругу.
Она лежала под Ситрисом, её губы были раскрыты в истоме, а огненные глаза под благородным лбом не мигая смотрели на меня...

Я зарычал и впился в губы Ликхэ.
Миг --- и она оказалась внизу.
Глаза её закатились, по телу пробежала волна...
И Ликхэ растворилась в воздухе.

--- Лисичка... --- низкий голос над моим ухом принадлежал Кхарасу.
Его мускулистые руки обхватили меня за талию, подбородок лёг мне на плечо.
Я почувствовал, как в меня входит его мощный стебель и застонал от смешанной с болью истомы...
Твёрдые губы Кхараса коснулись моих, в голове словно зазвенела толстая басовая струна экстаза...
И Кхарас растворился в воздухе.

--- Ликхмас, --- рядом появились живые, обрамлённые тёмными веками глаза Ситриса.
Я набросился на него и немного неуклюже повалил на землю, руки и ноги не слушались меня.
Ситрис охнул и довольно оскалился:

--- Ох, какой ты горячий.
Осторожнее...

Ситрис откровенно наслаждался происходящим.
Он проводил пальцами по моему животу, и каждое лёгкое прикосновение заставляло меня биться в судорогах.
Воин смеялся и шептал какие-то нежности...

А голоса всё тянули и тянули:

--- Ааа...

--- Ооо...

--- Ммм...

Та-та-та-хум!
Йолль!
Йоуу!

Я вдруг осознал, что никто больше не пел и не играл.
При этом музыка, потусторонняя, неземная, гудела у меня в голове.
На мгновение стало страшно, но в следующий миг путы страха разлетелись призрачными клочьями.
Я потерял контроль над разумом и чувствами --- меня всецело поглотил транс.

Ситрис растворился в воздухе, как и Кхохо, и Эрликх, и Хитрам...

Та-та-та-хум!
Тхийо-тхо-тхоу!

Руки Кхохо, в отличие от Ликхэ, не дрожали, воительница делала всё чётко, как по уставу.
Лишь изредка, во время особо чувствительного пассажа, в мою шею вонзались её острые ногти, и вокруг них, словно круги на воде, разбегались радужные волны, остро пахнущие корицей.
В потемневших глазах я вдруг увидел непреодолимое, страстное желание задушить меня.
Я положил свои ладони на её и сжал её руками своё горло.
Воительница засмеялась странным, похожим на клёкот смехом, погладила мне шею и больше к ней не прикасалась.

Та-та-та-хум!
Йолль!
Йоуу!

Поющие голоса надрывно зазвенели.
Я слышал грозный, величественный рык Кхараса, одинокий отчаянный крик Кхохо, мечтательный смех Ликхэ, приносящий спокойствие баритон Ситриса, нежный счастливый напев Чханэ...
Голоса наполняли меня, вытесняя, выдавливая мою собственную сущность.
Мне --- тому кусочку меня, который ещё продолжал соображать --- вдруг снова стало страшно.
Я понял, что больше я себе не принадлежу.
Однако следом пришла ещё одна мысль --- я никогда себе не принадлежал.
Едва мысль родилась, я успокоился так же быстро, как и испугался.

Эрликх валял меня по траве, как куклу, я потерял всякое представление о верхе и низе, но глаза воина смотрели за горизонт, сфокусировавшись на бесконечности.
Его голос звучал смиренно, словно пело согнутое в дугу молодое деревце.
Языки пламени играли на его груди спектакль теней, складываясь в благородные улыбающиеся лица --- воспоминания о давно ушедших за грань...

--- Ааа...
Лаа, лээй...

В небе вновь загорелись огромные глаза Чханэ, живые, невообразимо прекрасные...
Я любовался ими, зная, что больше никто не способен их видеть.
И только когда её губы нежно-нежно коснулись моих, словно девушка пыталась попить из ледяного горного водопада, только тогда я понял, что небо осталось где-то позади, а глаза действительно не иллюзия.

Мы с Чханэ завершили Круг друг у друга в объятиях.
Воины, разморенные утренней прохладой, лежали вповалку на земле и слушали пение птиц.
Голова Кхохо гнездилась на животе Ситриса, они оживлённо, с тихим смехом что-то обсуждали;
Ликхэ чуть в стороне рассеянно мяла в руках ароматные мелкие цветы;
Кхарас и Хитрам, примостившись рядом с провизией, с аппетитом жевали половинки пайка;
флейта Эрликха устало добирала последние вариации.
Наркотик больше не тревожил странными видениями, я чувствовал себя легко и спокойно.
Ночное небо бледнело, из чёрного становясь сиреневым.

Наступило утро.

\section{[*] Битва Трёх}

--- Ну вот, а говорил, что не можешь, --- сказала Эрхэ, обнимая Аурвелия.
--- Ты ещё ого-го дедуля.

Отряд чести лежал на Узорной площади Тхартхаахитра.
Медленно тлели угли костра.
Глаза воинов и лежавшего с ними Атриса сияли ярче звёзд.
Но даже они не могли соперничать с улыбкой Аурвелия.

--- В ночь тысячи дождей и тростинка превращается в копьё, --- развёл руками ноа.
--- Но спина... эвааа... болеть.

--- Так, размять дедушку, --- распорядилась Митхэ.

--- Не нужно, --- отмахнулся Аурвелий.
--- Я не понимать, зачем это делвать.
Мне нравится, но я не понимать.
Что вы вверить эта делвать?

--- Мы верим, что наши души смешиваются во время Круга, как ингредиенты супа в котле, --- сказала Эрхэ.
--- Сваренное в супе мясо несёт его аромат, даже если положить его в отдельную миску.

--- Властелин этого мира может лишь требовать жертвы, --- добавила Согхо.
--- Он ест нас поодиночке сырыми.
Лишь смешавшись в невкусном супе, мы можем встать ему поперёк горла!

Воины засмеялись;
многие тут же смущённо потрогали мизинцами языки.

--- Скажи, Аурвелий, а боги ноа отвечают за что-то? --- поинтересовалась Митхэ.
--- Или они --- как Безумный?

--- Бвезумный вне едим ввесь с... o sancta marin...

---  <<Безумный --- не один, он вместе с...>>, --- перевёл Ситрис, лучше всех понимавший старика.

--- Я устать, --- усмехнулся Аурвелий.
--- Когда я устать, я всегда... делать... неправильный язык, слово, никто не понимать.
Извините.

--- Расскажи тогда на ноа-лингва, --- сказала Митхэ.
--- Мы все понимаем твой язык.

--- Это невежливый.

--- С нашей стороны было бы невежливо заставлять тебя говорить на сложном чужом языке, когда ты устал.
Всё-таки сели как иностранный язык --- это скорее пытка, чем удовольствие.
Так всё-таки, за что отвечают ваши боги, боги ноа?

Аурвелий подумал и перешёл на ноа-лингва:

--- Мы не знаем, за что отвечают боги, но мы знаем, что знаменует усиление или ослабление каждого из них.
Если солнце выжигает посевы и на морском берегу проступают соляные письмена, Деа Марин слабеет.
Мы приносим ему жертвы.
Если бушуют штормы и вода касается deserto sancto\FM, слабеет Деа Солар.
\FA{
Светлая, чистая пустыня (ноа-лингва).
}
Тогда мы подпитываем его.
Если же люди счастливы и довольны, если посевы дают богатый урожай, а моря светлы и гостеприимны --- Деа Акседент проигрывает битву, и тогда мы помогаем ему.
Смерть бога --- шаг к концу мира.

--- Вы верите, что у мира есть конец? --- удивилась Эрхэ.
--- Я слышала Далёкие речи ноа, речь о Бесконечном...

--- Да, Бесконечный может переходить с одной земли на другую, и находить новый свет, когда старый погаснет.
Но земли смертны, и смерть нашей земли зависит от них, богов.
Нужно очень стараться, чтобы двое не победили одного, и нужно стоять насмерть, чтобы один не победил двоих.

--- Но кто помогает людям? --- удивилась Митхэ.
--- Неужели вы думаете, что в этом мире нет сил, которые выступают на стороне людей?

--- Четвёртая сила есть, --- ответил Аурвелий.
--- Эта сила держит мировой порядок, эта сила не даёт победить богам.
Эта сила --- мы, те, кто дышит и когда-то дышал.
Вы, белые рубахи, тоже приносите жертвы и помогаете держать мировой порядок, вы только не знаете, кому вы их приносите.
Мы знаем.
Если слабеет один из богов --- мы заклинаем его собирать все жертвы в известном мире.
И его подпитывают --- и мы, и вы, и идолы, и ркхве-хор.

--- И ты веришь в это?

Аурвелий кивнул.

--- А как же быть с Легендой об обретении?
Это ложь?

--- Нет, --- ответил Аурвелий.
--- Легенда об обретении рассказывает о прибытии Деа Акседента.
Именно его вы называете Безумным.
Но сели не знают, что ещё до Деа Акседента миром правили боги солнца и моря, и мы, ноа, поддерживали их борьбу.

--- Какая странная легенда, --- сказала Эрхэ.
--- Я не уверена, что поняла её правильно.

--- Не имеет значения легенда, --- заметила Митхэ.
--- Главное, что это работает и наш мир успешно сдерживает гнев богов, сколько бы их ни было.

--- Вот здесь я соглашусь с Митхэ, --- ухмыльнулся Акхсар.
--- Мой кормилец говаривал так: верь во что хочешь, но иди куда следует.

\section{[-] Близость и даль}

\epigraph
{Увы, но даже сейчас, в Эпоху Богов, на нашей родной Земле есть закрытые уголки.
Сектантские поселения, изолированные от внешнего мира племена, не принявшие путь генной модификации общества.
Люди, живущие там, прозябают в невежестве, страдают от болезней и угнетения.
Иногда всё, что им нужно для лучшей жизни --- это пересечь границу родного поселения, но они рождаются, живут и умирают в его стенах, олицетворяя полную страха жизнь человека прошлого.\\
О гуманности сказано достаточно моим другом Михаэлем [Михаилом Кохани], мне добавить нечего.
Всё, чего я хочу пожелать людям --- это смелости.
Опыт и знания --- ничто, если у вас нет смелости применить их.
Желания --- ничто, если вы не осмелитесь их осуществить.}
{Людвиг К.\,Р.\,Вейерманн, первооткрыватель омега-поля}

\textspace

Я рассказал воинам об увиденном в трансе.
Ситрис выслушал с неподдельным интересом, Кхарас одобрительно крякнул:

--- Парень-то чувствует.

--- Это точно, --- согласилась Кхохо.
Женщина после произошедшего начала относиться ко мне с большей теплотой.
Правда, выражалось это в основном в грубоватой ласке и тумаках.
Взглянув в её глаза, я понял: мне суждено было узнать о Кхохо то, что знали единицы.
Та тьма, потаённая, алчущая крови и страданий, но временно укрощённая, по-прежнему смотрела на меня из мерцающих зрачков.
Однако я чувствовал ещё кое-что: для меня эта тьма больше не представляла опасности.
Я подчинился и навсегда стал ей другом.

Ликхэ избегала говорить со мной, но я чувствовал её напряжение и знал, что скоро она не выдержит.

Так и случилось.
Воительница нашла меня возле ручья и, как обычно, взяла оленя за рога.

--- Хочешь, я буду твоей женщиной? --- спросила Ликхэ, присев рядом.
--- У меня много умений и достоинств, и мне очень нравится заниматься с тобой любовью.

Я усмехнулся.

--- Я тебя знаю с детства, как и все твои умения и все достоинства.
Нет нужды напоминать о них.

\ml{$0$}
{--- Так за чем дело стало?}
{``So what are we talking about?}
\ml{$0$}
{Чханэ может спать по одну сторону от тебя, а я --- по другую.}
{\Chhanei\ will sleep on your left, and me on your right.''}

\ml{$0$}
{--- Ты ведь на этом не остановишься.}
{``I doubt you will stop on that.}
Неужели Ликхэ ар'Трукх спокойно будет делить моё ложе и меня с Чханэ?

--- Ликхэ ар'Трукх потерпит, --- мило улыбнулась девушка.
%{will bear it}

\ml{$0$}
{--- Вот именно --- <<потерпит>>.}
{``That's the point---`bear'.''}

\ml{$0$}
{--- И что здесь такого?}
{``What's wrong?''}

\ml{$0$}
{--- Ликхэ, мне не нужно, чтобы ты терпела.}
{``\Likchoe, I don't want you to bear anything.}
\ml{$0$}
{Если бы ты искренне любила Чханэ и не желала владеть мной безраздельно, мой ответ мог бы быть другим.}
{If you loved \Chhanei\ like I do, if you didn't want to control me, my answer could be different.}
А так --- извини, мы тебя не примем.

\ml{$0$}
{--- <<Мы>>?}
{`` `We'?}
\ml{$0$}
{Ты и за Чханэ тоже решаешь?}
{Have you decided for \Chhanei\ too?}
\ml{$0$}
{Или это Чханэ всё за тебя решила?}
{Or it's \Chhanei\ who makes decisions?''}

--- Ликхэ, остановись.
Мы --- один Храм, помнишь?

--- Ты всё же подумай, --- Ликхэ ворковала нежно, но я корнями волос чувствовал скрывающуюся за этой нежностью ярость.
Мой отказ воительница восприняла как личное оскорбление.
--- А я пока нужду справлю.

Ликхэ встала, поправила складки на штанах и ушла в кусты.
Почти одновременно ко мне подошла Чханэ.
Девушка села рядом, рассеянно сорвала ароматный тёмно-зелёный лист <<белого пламени>>, потёрла его о шею и сладко зевнула.

--- Как прошёл Круг Доверия?

--- Очень бодрит, --- усмехнулся я.
--- Кстати, почему нельзя фальшивить?

--- Аккорд в голове слышал?
Слышал, конечно, все слышали, он для всех один.
А теперь представь, что кто-то сфальшивил.
У нас за такое били, потому что Круг насмарку.

Чханэ весело сморщилась и хихикнула.

--- Мы с тобой всем понравились.

--- Правда? --- удивился я.
--- А мне показалось...

--- Ты ещё неопытен.
Вообще я немножко удивилась распределению ролей в вашем Храме.

--- А что не так?

--- Ну, это сложно объяснить, на уровне чувств.
Вот взять тебя.
Ты очень нежный, не сильный, но решительный.
Будь ты воином, у нас тебя сразу определили бы в разведчики.
Ликхэ нежная, но это нежность анаконды.
Из неё вышел бы отличный убийца, но она, как видишь, стрелок.
В Кхохо нежности нет совсем, но действует она филигранно --- знает, когда применить силу, а когда нет.
Вот такие у нас шли в стрелки, но она ружьё в руки не берёт.
Эрликх --- наоборот, много нежности и гораздо меньше осмысленности.
Таких обычно отправляли к детям на площадку, но он за хозяйство отвечает.
И больше всего меня удивило, что во главе отряда стоит Кхарас.
Он силён, но грубоват.
Заметил, наверное?
У нас во главе отряда всегда стоял самый лучший разведчик --- они осторожны и совершают меньше ошибок.
Ситрис, мне кажется, больше подошёл бы.
Интересно было бы узнать, почему у вас по-другому.

--- Ты почти угадала.
Кхохо не берёт в руки ружьё, потому что оно ей жутко надоело, но вообще она может выстрелить из отгрызенного тростника кактусовой иглой и пронзить птице оба глаза, даже с похмелья.
Эрликх часто подменял Конфетку на площадке.
А Кхарас с Ситрисом...
На самом деле их связывает что-то очень глубокое.

--- Более глубокое, чем Ситриса и Кхохо? --- лукаво прищурилась подруга.

--- Я не лоцман и не ныряльщик.
Долгая история.
Они говорят, что так надо.

Из кустов показалась Ликхэ, поддёрнула штаны, мило улыбнулась нам и ушла к лагерю.
Мы с Чханэ кивнули ей.
Едва воительница скрылась из глаз, Чханэ, прищурив лукавые глаза, обратилась ко мне:

--- О чём вы говорили?

--- Она предлагала мне любовь, --- честно ответил я.

--- Так и знала, --- вздохнула девушка.
--- Она тебя ревнует.
Причём так, что у меня за кхене от её ревности глаза слезятся.
И как будем это решать?
Будешь спать и с ней?

--- Нет, конечно.
И я постарался это до неё донести, --- махнул я рукой.

Чханэ сплюнула.

--- Думаешь, её это остановит?
Я знаю таких людей.
Они относятся к тем, кого хотят, как к собственности, и ревность их страшна, как шторм в проливе.
Будь осторожен и не бери ничего из её рук.
Я тоже постараюсь.

Я промолчал и укоряюще посмотрел на подругу.
Она смутилась.

--- Лис, я понимаю --- ты знаешь ей с детства.
И...

--- И ещё я знаю, что ей можно доверять, --- перебил я.
--- Да, временами она бывает жестокой, временами вредничает, но мы --- один Храм.
Это о чём-то говорит?

--- Я очень хотела бы доверять ей так же, как ты.
Но я видела слишком много...

--- Что?
Что ты видела?

Чханэ скрестила руки и положила голову на колени.

--- В чём проблема, объясни.
Разве вы устраиваете Круг Доверия не для того, чтобы доверять?

--- Именно для этого, --- кивнула Чханэ.
--- Для меня было бы лучше узнать о её интересах сейчас, чем когда яд или кинжал войдёт в моё горло.
Обычная вещь в храме...

--- Обычная? --- недоверчиво спросил я.

Чханэ оглянулась и сделала вид, что поправляет доспехи.
А потом приникла к моему уху:

--- За то время, что я служила Храму Тхаммитра, власть в Храме сменилась три раза.
Просто ночью из катакомб тайком выносили тела, а утром в зал выходил новый Первый жрец.
Не было ни шума, ни разговоров, в городе никто ничего не знал.
Да и кого интересуют внутренние передряги Храма, если Храм выполняет свою работу?
Мне повезло --- ни один переворот я не застала.
Просто приходила утром, видела другого человека во главе стола и обращалась к нему, как к прежнему, не задавая вопросов.

Я ошеломлённо молчал.
Чханэ посмотрела на моё лицо и тихо вздохнула.

--- Не там я родилась.

--- У меня даже мысли не возникало, что может быть иначе, --- признался я.
--- Между адептами нашего Храма всегда были если не дружеские, то по крайней мере добрососедские отношения.
Я не понимаю, что в Храме делить, из-за чего спорить?
Мы и так почти всё получаем бесплатно...

--- Факт остаётся фактом.
Как бы я хотела уйти от всего этого, духи лесные, как бы я хотела...
Уехать туда, где нет ни войны, ни интриг...

--- Нам некуда ехать, --- заметил я.
--- Безумные правят везде.

--- Да ну? --- отстранённо, словно кому-то другому, пробормотала девушка.

Я насторожился.

\ml{$0$}
{--- Объясни.}
{``Explain.''}

Чханэ замолчала и попыталась встать, чтобы уйти.
Я повис на её руке и вернул на место.

\ml{$0$}
{--- Объясни.}
{``Explain.''}

Чханэ мотнула головой:

--- Это всего лишь слухи, я не...

--- Позволь мне решать.
Говори, что знаешь, --- я чересчур сильно сжал её руку и опомнился, только когда она свирепо шмыгнула носом.

Чханэ потёрла затёкшую кисть и, поколебавшись, пробурчала:

--- Кристалл.
Там нет ни алтарей, ни жертв.
Так говорила бабушка.

--- А радужное безумие?
Это ведь не выдумки?

--- Нет, не выдумки.
Я видела его своими собственными глазами.

--- Значит, и Безумный существует.
Но если где-то люди живут без жертв...
Почему у нас по-другому?
Это какой-то заговор?

--- Говори тише, Лис.
Если и заговор, то воины тут ни при чём, --- отрезала Чханэ.

--- Ты намекаешь, что при чём здесь жрецы?

Чханэ посмотрела на меня, прикидывая, стоит ли говорить то, что вертится у неё на языке.
Потом нервно оглянулась и снова сделала вид, что поправляет доспехи.

--- Помнишь Сатхира ар’Со?
Я тебе про него рассказывала, когда...

--- Помню.

--- Так вот.
Его брат, тоже жрец, очень сильно пил.
Ситлам-лехэ, милейший старичок.
Одно время я часто бражничала с ним вместе.
Он жаловался, что не может больше смотреть на страдания детишек.
И я знаю, что многие жрецы в обход запрета утешают жертв, пытаются облегчить...

\ml{$0$}
{--- Чханэ, идол тебя покусай, короче!}
{``\Chhanei, the idol bite you, tell briefly!''}

Девушка сверкнула на меня глазами.

--- Сатхир был другим.
Неограниченная власть его развратила.
Он мучил с нескрываемым удовольствием.
Думаю, что я не первая, кого он повёл в джунгли --- ему просто перестало хватать обычных жертв.

--- Но как это возможно?
Что это за...

--- Он не был таким всегда, Лис.
Это похоже на болезнь, --- быстро зашептала Чханэ, наклонившись ко мне.
Видимо, когда-то эта метаморфоза повергла молодую воительницу в ужас, далёкие отголоски его метались в её оранжевых опалесцирующих глазах.
--- Его взгляд... он нездоровый.
Пустой.
А самое страшное --- что таких больных очень много, они знают друг о друге и координируют свои действия.
\ml{$0$}
{И я не одна, кто так думает.}
{And I'm not alone who thinks so.}
Это заговор не жрецов, а кого-то гораздо более могущес... пойдём-ка, Лис, не нужно оставлять товарищей надолго.

Обратный путь мы проделали в молчании.
Наконец Чханэ сказала:

\ml{$0$}
{--- Надеюсь, ты понимаешь...}
{``I hope you realize---''}

\ml{$0$}
{--- Я буду молчать, --- отрезал я.}
{``I won't tell anyone,'' I interrupted her.}
Это получилось чересчур резко, и я чуть мягче добавил:
\ml{$0$}
{--- И подумаю над твоим предложением.}
{``Also, I will think about your suggestion.''}

\ml{$0$}
{--- Каким предложением? --- удивилась Чханэ.}
{``What do you mean?'' \Chhanei\ looked surprised.}

\ml{$0$}
{--- Ну, Кристалл не так уж и далеко, --- смущённо усмехнулся я и пожал плечами.}
{``Well, Crystall never was far away,'' I shrugged with a shy smile.}

Лицо Чханэ просветлело, и она тут же прикрыла улыбку ладонью, словно боялась, что та упорхнёт.
Девушка схватила мою руку и сжала её --- с теплотой и благодарностью.

\chapter*{Интерлюдия VII. Пророчество}
\addcontentsline{toc}{chapter}{Интерлюдия VII. Пророчество}

\textbf{Легенда об обретении}

Было то во времена, когда все народы земель и морей жили в мире, и духи ходили по земле, словно люди.

Пришёл тот, кто был назван Ликхмас, в город Ихслантхар.
\ml{$0$}
{Город не заметил юношу, как не замечал прочих.}
{The city hasn't noticed the young one, as he hasn't noticed others.}
И пришёл Ликхмас на постоялый двор, что стоит на перекрёстке трёх дорог\FM.
\FA{
Перекрёсток трёх дорог --- площадь в квартале ювелиров Ихсланхара.
Среди сели есть мнение, что это заколдованное место, и там стоит невидимый постоялый двор для лесных духов.
}
Он заказал себе пищу, и [принесли ему] тарелку кукурузы, сладости и чашу с водой\FM.
\FA{
Пища, которую поднесли Ликхмасу, на самом деле представляет собой классическое подношение умершим --- зерно, сладости и хэситр.
}
Пока Ликхмас трапезничал, к нему подсел старик, закутанный в рваный чёрный плащ.
Глаза его горели, как звёзды, а седая борода развевалась, словно на шквальном ветру, хоть сидел старик в душной таверне.

<<Жизнь и процветание, лехэ, --- поприветствовал старца [Ликхмас].
\ml{$0$}
{--- Не желаешь ли ты разделить со мной пищу?>>}
{``Would you mind to share my meal?''}

Сверкнули глаза старца.

\ml{$0$}
{<<Я поклялся более никогда не разделять пищу с людьми>>.}
{``I have sworn never to share meals with humans.''}

<<Отчего так?>> --- удивился юноша.

\ml{$0$}
{<<Я был духом, --- ответил старик.}
{``I was a Spirit,'' the old man answered.}
\ml{$0$}
{--- И наследница вашего рода лишила меня того, благодаря чему я был нужен этому миру.}
{``And a scion of your kind has wiped out what the world needed me for.}
\ml{$0$}
{Она не лишила меня сил, она не лишила меня умений, но я не смею тронуть ту гармонию, что она установила хлопком одной ладони.}
{She left untouched my powers and my skills, but I dare not interfere with the harmony she has made with one hand clap.}
Если я не найду другого занятия, Ликас ар'Тро, я стану каменным духом из тех духов, что [существуют] без цели и смысла>>.

<<Сочувствую твоему горю, лехэ, --- ответил юноша.
--- Но это не повод лишать себя простых радостей, ибо они есть цель и смысл среди прочих [целей и смыслов].
И да, меня назвали Ликхмас.
\ml{$0$}
{Откуда бы ты ни знал то имя, Ликас ар'Тро мёртв, и никогда больше [он не будет живым]>>.}
{Wherever from you knew that name, Likas ar'Troa is dead, and [he] shall never [be alive anymore].''}

<<Значит, я пришёл туда, куда нужно, --- ответил старец.
\ml{$0$}
{--- Духи знали о приходе Безумного.}
{``The Spirits have foreseen the Senseless' coming.}
\ml{$0$}
{Духи знали о том, что мир ждёт его владычество, ибо было то предсказуемо и неизбежно.}
{The Spirits have foreseen the world lying under his rule, for this was predictable and inescapable.}
\ml{$0$}
{Но так же [предсказуема и неизбежна] его гибель>>.}
{So is his death.''}

Ликхмас промолчал.
Он слушал, как вслушиваются в небеса, пытаясь различить жужжание светлячка, кружащего вокруг звезды.

\ml{$0$}
{<<Кихотр Безумного может вершить судьбы.}
{``The Senseless' \Kihotr\ decides the fate.}
\ml{$0$}
{Ни один смертный не сможет поднять и бросить этот камень без вреда для себя.}
{No breathing thing can throw this dice and remain unharmed.}
\ml{$0$}
{Но существует человек, который может выбросить ту из бесчисленных граней, что несёт знак смерти самого Безумного.}
{But there is a man who can get the one of uncountable faces---the death sign of the Senseless himself.}
\ml{$0$}
{Ибо не Безумный сотворил кихотр, ибо он не мастер из тех мастеров, что куют собственный инструмент>>.}
{For it was not the Senseless who made the \Kihotr, for he is not a master of the masters who forge their own tools.''}

Последние слова порывом ветра пронеслись по таверне, и люди, идолы, пылерои и травники замолчали, слушая старца.

\ml{$0$}
{<<Этот человек уже умер и всё ещё жив.}
{``This man is already dead and still alive.}
\ml{$0$}
{Этот человек имеет имя и лишён имени.}
{This man has got a name and has lost a name.}
\ml{$0$}
{Он есть сидящий человек и бегущая лань>>.}
{This man is a human who sits and a doe who runs.''}

Старец встал и вышел вон, и лишь ветер провыл ему вслед.

Ликхмас обратился к сидящим в таверне:

<<Живые и мёртвые, духи, люди, идолы, пылерои и травники.
Знает ли кто из вас, где найти кихотр Безумного?>>

Вперёд вышел воин, имя которому было Маликх\FM.
\FA{
Маликх (правильнее Муал'ликх) --- <<владыка>>, дословно --- <<возлежащий на цветах>>.
Титулование главы знатного рода в племени тенку.
}
Он не носил доспехов, и волосы его были так длинны, что конец косы лежал за воротами города\FM.
\FA{
Воины сели обычно стриглись коротко.
Длина волос Маликха могла быть намёком как на невозможность его существования, так и на его невероятное мастерство, помогавшее быть воином даже с волосами такой длины.
}

\ml{$0$}
{<<Я не знаю, где находится кихотр, но я защищу тебя от любого бедствия, что встанет на твоём пути>>.}
{``I never knew where the \Kihotr\ is, but I shall protect you from every threat you will ever face.''}

Затем вперёд вышел купец с орлиным носом, острым подбородком и ртом, [способным вместить] яблоко целиком.
Имя ему было Чхалас, и была его речь так сладка, что даже лес и горы уступали ему, если [он обращался] к ним.

\ml{$0$}
{<<Я не знаю, где находится кихотр, но я смогу найти того, кто это знает, и он поведает мне всё без утайки>>.}
{``I never knew where the \Kihotr\ is, but I shall find the one who does, and they shall tell me what they knew.''}

Затем вышла вперёд идол, у которой маска прикрывала лицо, ибо половина её головы была куском дерева\FM.
\FA{
Не вполне ясно, как в Легенде появился этот персонаж.
Описание его напоминает планта, страдающего наследственным верхнечелюстным экзостозом, который в тяжёлой форме напоминает древесную кору.
Возможно, что такие планты были объектом почитания у некоторых племён Тра-Ренкхаля;
тем не менее, подобные традиции не сохранились ко времени перехода планеты под знамёна Ордена Преисподней.
}

<<Я не знаю, где находится кихотр, но я отдам вам ветвь из своей головы.
Она за взмах ресниц способна прорасти в лес, столь густой, что ни один враг не пройдёт через него>>.

\ml{$0$}
{Затем вышел вперёд старый травник о трёх ногах\FM.}
{Then an old three-legged grasshider\FM\ stepped forward.}
\FL{weaver}{Танцуют-Четыре-Камня}

\ml{$0$}
{<<Я не знаю, где находится кихотр, но мне было суждено оказаться здесь, и мне было суждено заговорить.}
{``I never knew where the \Kihotr\ is, but I was meant to be here, and I was meant to speak out.}
Посмотрите на мои ноги --- их число равно числу дорог, что сходятся на этом перекрёстке.
\ml{$0$}
{Но я выкуклился с четырьмя ногами, и значит, что четвёртая дорога должна быть>>.}
{But I escaped my cocoon having four legs, and the fourth road is meant to be.''}

Встал травник по сторонам света, расставил целые ноги по направлению дорог --- и указала культя его ноги [в нужную сторону].

И на рассвете вышли Ликхмас, Маликх и Чхалас из Ихслантхара и пошли по четвёртой дороге\FM.
\FA{
Современные сели <<четвёртой дорогой>> называют неочевидный путь.
Например, если свернуть с тропы и пойти через глухую чащу --- это будет четвёртой дорогой.
Если вместо двери пробиться через стену или прокопать подземный ход --- это тоже будет четвёртой дорогой.
}
\ml{$0$}
{И тьма сомкнулась над ними --- отрезана [четвёртая дорога] от мира, подобно ноге старого травника, и нет над ней ни светил, ни небес.}
{And darkness has swallowed them, for the fourth road is divided from the world, like the old grasshider's leg, and above it there is no skies and no light.}

\chapter{[-] Поход}

\section{[-] Четвёртая дорога}

--- Быстрее, Лисичка, не отставай, --- Ликхэ хлопнула меня по спине.
--- Следы холодные, надо спешить.

--- Мы успеем перехватить их до Тхисема? --- спросил я.

--- Успеем, если прибавим шагу, --- Ликхэ напряжённо вглядывалась куда-то в заросли.

--- Кхарас, идём по четвёртой дороге, --- Ситрис тоже выглядел так, словно собирался броситься галопом в заросли.
--- Ликхэ дело говорит, они нас опережают на пару шагов.

Кхарас нахмурился, но промолчал.

--- Я бы больше волновался не за Тхисем, а за Одинокий Столб, --- заметил Эрликх.
--- Тхисем почти на границе с землями святилища.

--- Почему нужно успеть до святилища? --- спросила Чханэ.
--- Я думала, что с пленниками туда нельзя...

\ml{$0$}
{--- В крепость их, конечно же, не пустят, --- согласилась Хитрам.}
{``They won't be allowed in the stronghold, yes,'' \Chitram\ agreed.}
\ml{$0$}
{--- Но если мы нападём на них на землях святилища --- виноваты будем мы и проблемы будут у нас.}
{``But, if we attack them on the sanctuary ground, guilt will be ours and problems will be ours.''}

--- Ну не убьют же, --- пожала плечами подруга.
--- Если не успеем до Тхисема, я бы предложила устроить засаду на землях святилища, неожиданно для них.
\ml{$0$}
{В следующий раз эта гниль шестнадцать раз подумает, прежде чем детей похищать.}
{Next time those rotten hearts will think sixteen times over before they dare to kidnap children.''}

--- Нет, --- буркнул Кхарас.

\ml{$0$}
{--- Если оскверним святилище, договориться будет в разы сложнее, --- пояснил Эрликх.}
{``With the sanctuary violated, it will be much more difficult to negotiate,'' \Oerlikch\ explained.}
\ml{$0$}
{--- Мы уже это проходили, обсуждать нечего.}
{``We've been through this, there's nothing to discuss.}
Кхохо вон как-то дала в морду одному на территории святилища.
Кхотлам еле её тогда...

--- Так, всё, хватит, --- рявкнула Кхохо, швырнув в Эрликха горсть листьев.
--- Все уже поняли.

По дороге попался колодец.
Ликхэ и Хитрам по привычке кинули туда по золотой грануле, пробормотав под нос имена детей и молитвы Ситу;
Ликхэ пробормотала одно имя, Хитрам задержалась гораздо дольше.
Эрликх, оглянувшись по сторонам, стянул с руки кольцо и тоже бросил в воду.

--- Кстати, Орешек, я намедни была в жилище Матрэ ар'Тхар, --- тихо сказала Хитрам Ликхэ.
--- Видела твоего хранителя.
Он же просто копия Кхатрима!
Глаза, губы!

--- Да, я тоже заметила, --- весело ответила Ликхэ.
--- Скорее всего, от него.
Хотя вначале я думала, что от Кхараса.

--- Да я сама так думала! --- развела руками Хитрам.
--- Такой важный молчун.
Ты Кхатриму ещё не рассказывала?

--- Он не особенно интересуется своими потомками, если они не лежат у него в госпитале, --- призналась Ликхэ.
--- При случае, может, скажу.

Последним мимо колодца прошёл Ситрис.
Он окунул в воду ладонь, поднёс ко рту и тут же выплеснул горсть на землю, поморщившись от незнакомого запаха.

--- Дикари, --- пробормотал он, отряхиваясь.
--- Если не цените жизней сели, так цените хоть труд каменщика!

--- Что такое? --- обернулся Кхарас.

--- Что-то в колодце, --- пояснил Ситрис.
--- Пробовать не советую.

--- Колодец в двух кхамит ходьбы от хутора, --- помощилась Ликхэ.
--- Вот какова была вероятность, что мы будем пополнять из него мехи?
Точно дикари.
Надо потом сказать советам, чтобы почистили.

--- Стоять, --- рявкнула вдруг Хитрам.

Мы замерли.
Воительница осторожно вытащила из мочажины туго свёрнутую ветвь какого-то колючего кустарника и принесла Кхарасу.
Мне почудился слабый едкий запах.

--- Тхэтраас, --- резюмировала Ликхэ.
--- Кишки рыбьи, ещё и яд необычный, какое-то сложное зелье.
Наверняка не сразу убивает.

\ml{$0$}
{--- Всё правильно, --- заметил Ситрис.}
{``They got that right,'' \Sitris\ noted.}
\ml{$0$}
{--- Быстродействующий уменьшает количество врагов, но не их скорость.}
{``Fast-acting poison reduces number, but not speed of the enemy.}
\ml{$0$}
{Труп незачем тащить, труп не требует заботы.}
{A corpse needs no carry and no care.}
\ml{$0$}
{Кхарас, не будь рыбой, дорога тухлая.}
{\Kcharas, don't be a fish, the road's rotted.}
\ml{$0$}
{Сворачиваем.}
{We've to get off''.}

Кхарас хмуро смотал чётки Сата и скомандовал сойти с тропы.

\section{[-] Змея}

--- Давайте вплавь, что уж тут, --- пробормотала Хитрам.

--- Ликхэ надо дать вон тот валун, --- захихикала Кхохо.

--- Чего это? --- возмутилась Ликхэ.

--- Да у тебя два поплавка за пазухой, ты сама нырнуть не сможешь!

--- Лучше поплавок за пазухой, чем между ушей! --- хмыкнул Эрликх.

--- Эрликх, вот ты сейчас получишь у меня, --- буркнула Кхохо.
--- Пусть она сама отшучивается уже!
Ребёнка и так испортили, ей на чувство юмора кит хвост положил!

--- А тебе --- на чувство самосохранения, --- вмешалась Ликхэ.

--- Ну вот, --- пожал плечами Эрликх.
--- Может же, когда хочет.

--- Ребят, давайте всё-таки поищем переправу ниже по течению, --- поморщился Ситрис и показал обломки костей со следами мелких зубов.
--- Здесь сапожник водится.
Да и не время сейчас мокнуть, отряд близко.

--- Выполняйте, --- буркнул Кхарас.
Разговоры тут же прекратились.

\asterism

Переправа нашлась уже через сто шагов.
Крепкий и толстый корень тхальсар --- настоящая жила джунглей --- перекинулся через реку, прихотливо оплетя оба берега.

--- Какой хороший, --- восхитилась Чханэ, щупая корень рукой.

--- Главное, что живой, --- поддержала Ликхэ.
--- Редко можно найти живой мост такой длины.
Надо потом сказать, чтобы здесь сделали навесной и тропинку, очень удачное место.

С дерева соскользнула тёмная змея, и Кхохо, выругавшись, упала в реку.
Руки Чханэ и Ситриса запоздало царапнули воздух, пытаясь подхватить воительницу.
Мутная зеленоватая вода сомкнулась над ней.

Я хотел броситься следом, но Эрликх, выдохнув <<Хэ>>, схватил меня за шиворот.

--- Не сейчас, Лис, --- домашнее имя означало серьёзность ситуации.

Ситрис напряжённо вглядывался в воду.
Ни малейшего признака Кхохо.
Вдруг в зеленоватой воде показалась кровавая дымка.

--- Змеиная, --- констатировал Ситрис, окунув палец в воду и попробовав её на вкус.
--- Надеюсь, Кхохо её убила.
Не дёргайтесь, сейчас главное --- сапожника не проворонить.
Да стой ты, --- рявкнул он на Ликхэ, собравшуюся прыгать.

--- Если она не добила змею, та её задушит, --- прошипела Ликхэ.

--- Если мы пропустим сапожника, нырять будет уже не за кем, --- объяснил Эрликх, так как Ситрис даже не удосужился посмотреть в сторону девушки.

Кхарас бросил Ситрису чашу, и тот налил в неё из бутылочки немного жидкости --- белой, остро пахнущей, слегка напоминающей молоко на вид.
Кхарас подхватил чашу и перебежал на другой берег.

Вода вокруг забурлила.
На кровь приплыл косяк рыбы-сапожника.
Эта маленькая рыбка получила название за то, что за считанные секхар сжирала кожаные сапоги.
Ситрис и Кхарас тут же вылили молочно-белую жидкость в ручей.

Я почти услышал яростное щёлканье челюстей.
Несколько десятков уродливых зубастиков всплыли кверху брюхом, прочие предпочли ретироваться.

--- Слава духам, что я взял молочник, --- Ситрис бессильно прислонился к мшистому стволу.
--- Эрликх, будь другом... я не в состоянии...

--- Раздевайся, Ликхмас, --- толкнул меня Эрликх.
--- Я один её из змеи не выпутаю.

Кхарас так же молча снял с плеча верёвку, размотал и бросил нам конец.

\section{[-] Лекарство}

Вскоре исцарапанное, бледное как снег тело Кхохо уже лежало на берегу.
Я, Кхарас, Ситрис и Эрликх стояли рядом, прочих вождь отослал вперёд, разведать ситуацию.
Пульс и дыхание я не определил, но Кхарас махнул рукой: <<Всё в порядке, не суетись>>.

Эрликх методично сорвал с шеи Кхохо последних пиявок.
Из мест укусов тонкими ручейками текла кровь, оставляя на лодыжках, запястьях и лице алые полоски.
Я выдавил по капле коагулирующего раствора на каждую ранку.

--- Ты чего? --- удивлённо поднял брови Эрликх.

--- Слюна пиявок...

--- Да знаю я.
Успокойся, она не хрустальная, выживет и так.
В походе нет времени на каждую царапину.

--- Маленькие раны могут дать большую кровопотерю.

\ml{$0$}
{--- Как знаешь.}
{``Suit yourself.}
\ml{$0$}
{Это твоё время.}
{It's your time to waste.''}

Ситрис попросил мою коробку с инструментами, неторопливо развернул трубочку и с неожиданной сноровкой ввёл её подруге в горло по клинку ларингоскопа.

--- Шприц.

Я протянул.
Ситрис надел трубку на шприц, поводил, потянул поршень шприца, нахмурился и ещё раз поводил трубкой.

--- Молодец, всё правильно сделала.
Учись, Ликхмас.
Если знаешь, что тебя вытащат из воды --- прикусывай язык что есть мочи и выдыхай весь воздух.
Хотя воздух тут, пожалуй, змея выдавила...

--- А как её теперь оживить?

--- Можно просто оставить, сама оживёт.
Но это скучно.

С этими словами Ситрис достал ещё одну бутылочку и капнул Кхохо в нос маслянистой жидкости.
Десять секунд спустя воительница зашевелилась с яростным мычанием и, нащупав мехи, начала трясущимися руками промывать нос.
Затем бросила мехи, подползла к воде и окунула в неё голову.

--- Она меня теперь год будет по ночам лупить, --- вздохнул Ситрис.
--- Это масло пустынного перца из Санта-Виктории.
Говорят, кто его не ел --- тот не знает вкуса жизни.
Хотя, на мой взгляд, это скорее вкус смерти, причём весьма болезненной смерти от огня.
От воды будет только хуже.

И, достав кисет с коноплёй и бумагу из водорослей, воин принялся сворачивать оленью ножку.
Кхарас невозмутимо поднёс ему тлеющий трут.
Он, похоже, знал о привычке Ситриса курить после сильных волнений.
Мужчины взглянули друг на друга и ухмыльнулись.

\section{[-] Везение}

От подобных происшествий не был застрахован ни один житель джунглей.
Ежегодно город терял несколько человек из-за змей, ядовитых пауков и аллигаторов.
Ещё с десяток возвращались с увечьями.
Кхохо повезло --- встреча с питоном стоила ей нескольких треснувших рёбер.
Несмотря на уговоры Кхараса, воительница продолжила путь.
Она еле переставляла ноги и по-детски цеплялась за Ситриса, забыв о его жгучем сюрпризе.

Вскоре перед нами нежданно-негаданно встало болото.

--- Да откуда здесь болото? --- в отчаянии бросил Ситрис, зачерпнув рукой немного прозрачной пахучей воды.

--- Оттуда, --- флегматично отозвался Кхарас.
--- Ты здесь когда последний раз был?

--- Если ты думаешь, что я пошёл по тропе своих воспоминаний, то ошибаешься, --- огрызнулся Ситрис.
--- Охотники мне ничего про болото не сказали.

--- Так, всё, хватит, --- буркнула Ликхэ.
--- Тхисем в одном кхене пути, здесь что-то должно быть.

--- Кажется, я вижу мостки, --- Чханэ указала куда-то вдаль.

--- Вот поэтому ничего и не сказали, --- объяснила Ликхэ.
--- Они здесь каждый день ходят и не задумываются, что это может стать препятствием.
Всё, Ситрис, успокойся, мы уже опоздали.
Ты же не собирался соваться в Тхисем, если его атакуют?

--- Да не в этом дело.
Эти мостки на виду, если мы толпой попрёмся через них, пиши пропало.

Кхарас нахмурился.

--- Ликхмас, сбегай по мосткам до хутора и посмотри, что как.
Мы пока поищем путь в обход.

\textspace

\section{[-] Разорённый хутор}

\epigraph
{Слабость людей --- в разъединённости.
Существует много способов разделить людей, но самым древним, варварским и действенным способом является разделение по признаку пола.
Оно начинается с самого детства.
Мальчики и девочки получают разные игрушки, разные обязанности и права, их отправляют в разные школы, им предлагают разный набор профессий.
В обществе активно пропагандируется гомосоциальность, нагота противоположного пола позиционируется как постыдное зрелище.
Когда же сковывающие мужчину и женщину барьеры падают, они, охваченные безрассудной страстью, бросаются навстречу друг другу, не зная друг о друге ничего.
Цена невежества --- взаимное неприятие, выливающееся в холодность, страх и насилие.
Чтобы управлять этими людьми, женщинам достаточно показать зверя-мужчину, а мужчинам --- женское тело.}
{Мариам Кивихеулу, писательница времён Последней войны}

\textspace

Захватчики окружили горящий дом.
Но вскоре плач ребёнка утих.

--- Там никого нет, Имжу, --- сказал один из них.
--- Ребёнок погиб.

--- Ждите, --- ответил главный.
--- Он может быть со взрослым.
Взрослого убивайте сразу.

Предводитель сделал знак помощнику, и они вместе отправились на следующую улицу.
Воины заворчали, но остались ждать.

Наконец крыша рухнула, в небо вылетел сноп искр и угольков.
Горящая дверь отвалилась, и на улицу вышла девушка с ребёнком на руках.
Одежды девушки и ребёнка окутывал пар --- она смачивала их водой.
Девушка зачем-то прикусила язык --- с уголка рта текла струйка крови.
Волосы у обоих сгорели дотла, странно лоснящаяся кожа покрылась копотью, но оставшиеся чистыми участки кожи были белее горного снега.

Воины застыли в изумлении.
На девушке были ожоги, но явно не те, которые они ожидали увидеть.
Девушка швырнула ребёнка под ноги воинам, и те зашипели --- ребёнок лежал бездыханный.

--- Вам нужно что-то ещё? --- спокойно спросила она, с трудом шевеля опухшим языком.

Воины несколько секунд рассматривали крестьянку.
Затем дружно расхохотались.

--- А хороша, --- довольно рыкнул один из молодых.
--- Ну что, парни, развлечёмся?
Негоже побывать в землях сели и не попробовать баб.
У них характер есть, не то что у наших.

--- И то верно, Перт, --- сказал другой.
Он возвышался над девушкой, как скала.
--- Эта немного пережарилась, но тоже сойдёт.

--- Имжу сказал... --- попытался возразить кто-то.

--- Да срал я на этого выродка, --- плюнул Перт, говоривший первым.
--- Харт, ты что, бабьи ноги не удержишь?

Прочие рассмеялись.
Девушка презрительно оглядела мучителей.
Лицо из снежно-бледного стало ярко-красным, затем пошло пятнами.

--- Вам нужно ещё и моё тело?
Подойдите и заберите, животные! --- девушка сорвала с себя полусгоревшую одежду.

--- Совсем стыда нет, --- с отвращением бросил Харт.

Кряжистый детина, осклабившись, схватил девушку за пояс и попытался повалить.
Но его толстая шея вдруг хрустнула и неестественно вывернулась под тонкими девичьими руками.
Побагровевший воин рухнул на землю.
Палец девушки стрелой пронёсся в сторону другого воина и вошёл в глазницу, словно горячий нож в масло.

В поясницу ей тут же вонзилось три копья.
Девушка, истекая кровью, легла и свернулась клубком, словно спящий мангуст.
Ни стона, ни крика воины хака от неё так и не услышали.

--- Безумный разорви эту проклятую деревню! --- выругался один из хака и пнул девушку.

К воинам тем временем подошёл главный.

--- Что у вас тут... НАЗАД!
БЫСТРО ВСЕ НАЗАД!

Поздно.
Кровь вдруг перестала сочиться из колотых ран на спине девушки.
Лицо умирающей стало умиротворённым и равнодушным, и она бросилась на ближайшего врага.
Захрустел перегрызаемый хрящ --- и воин упал, захлёбываясь кровью.
Та же участь постигла следующего.

Предводитель подлетел к девушке и одним ударом отрубил ей голову.
Затем окинул яростным взглядом оставшихся в живых.

--- Перт! --- заорал он.
--- Я кому инструкции давал, каменная твоя башка?
Не подпускайте сели на расстояние вытянутой руки, даже безоружных!
Ни стариков, ни женщин, ни тем более мужчин!
В плен брать только детей!
Взрослых убивать ударом в сердце или голову, и никак иначе!

--- Так я и бил в сердце! --- огрызнулся Перт.

--- Справа у них сердце, справа.
Ты право и лево путаешь?

В горящих развалинах дома кто-то зашевелился, и на свет вылез обгоревший подросток.
Один из воинов, вытащив фалангу, двинулся к нему.

--- Стой, выкидыш идола!
Сели только кажутся тщедушными, а в умирающих и вовсе вселяются мстительные тени.
Стрелу, целься в глаз!

Воин спустил тетиву.
Имжу грязно выругался.

--- Хватит с меня... черви подземные.
Уходим, мы взяли в плен достаточно.
Разведчики сели уже здесь.
Один сидит в кроне, за синим домом.

Разумеется, речь шла обо мне.
Воины напряглись.
Перт медленно достал лук.

--- Где он?
Я не вижу...

--- Лучше молчи, Перт, или я оставлю тебя на корм шакалам.
Из-за того, что кое-кому захотелось поразвлечься, мы уже лишились пятнадцати копий, и это не считая основных потерь.
Я вам говорил, что не надо атаковать этот хутор?
Говорил.
Я вам говорил, кого брать в плен?
Говорил.
В следующий раз будете умнее --- если вернётесь к родному очагу.
А не вернутся ещё двадцать мужчин, и это если нам повезёт в дороге.

--- Разведчик...

--- Забудь про разведчика.
Бери детей и отводи людей к нашим позициям.
Этого тоже возьми.

Имжу пнул ребёнка убитой крестьянки.

--- Он дохлый! --- запротестовал Перт и тут же осёкся, увидев взгляд Имжу.

--- Он оживёт, Перт, --- ласково сказал предводитель.
\ml{$0$}
{--- Ты повесил на свой пояс чересчур мало голов, чтобы знать --- трупы не бывают такого цвета.}
{``You've got too few heads on your belt to know: dead bodies are never that color.}
\ml{$0$}
{Это либо мраморная статуя, либо, чтоб тебя растерзали коты, живой притворщик-сели!}
{It may be a marble statue or, the cats tore you apart, it may be a \Seli\ pretender, alive!}
\ml{$0$}
{Ещё вопросы есть?}
{Any more questions?}
Если нет, первый отряд идёт со мной, прикрываем отступление прочих с пленниками.

Униженный Перт покраснел, словно ягода акхкатрас, но больше не роптал.
Воины, время от времени выискивая меня глазами, бросились исполнять приказ.
Я соскользнул с ветки и побежал обратно в лагерь --- докладывать.

\section{[-] Звон топора}

--- Значит, это всё-таки хака, --- задумчиво сказал Кхарас, выслушав меня.
--- Не ожидал от них такого.
Ещё и Инхас-Лака...

--- А что не так с этими Инхас-Лака? --- поинтересовалась Чханэ.
--- Они на слуху, а почему --- не вполне ясно.

--- Насколько я понял Кхотлам, в союзе племён хака есть несколько группировок, которые искусственно подогревают ненависть к сели, --- объяснил Кхарас.
--- И старейшины Инхас-Лака --- одна из таких группировок.
Нашествие Молчащих обезглавило эту группировку, плюс Кхотлам очень своевременно наладила дружеские отношения с вождём.
Однако то, что Инхас-Лака снова в числе напавших --- плохой знак.

--- Я думаю, что не все из напавших --- хака, --- заметил я.
--- Их предводитель увидел меня, хоть замаскировался я хорошо.

--- Хака не мог тебя увидеть, они не видят тепло, --- вскинулся Кхарас.
--- Думаешь, это сели?

--- Не похож на сели, --- ответил я.
--- Но всякое может быть.
Вдруг это пираты-ноа в боевой раскраске хака?

--- Ноа в союзе с хака? --- ошарашенно спросила Чханэ.

--- В таком случае маскируемся предельно аккуратно, --- бросил Кхарас остальным.
--- Молодец, Ликхмас.
\ml{$0$}
{Если этот человек, кто бы он ни был по крови, ещё и координатор, то мы бы хлебнули кровушки.}
{If that man is a coordinator, whatever his blood is, we'd gulp our own blood.}
\ml{$0$}
{Кхохо, беги обратно, от тебя сейчас толку нет.}
{\Kchoho, run back, you're no use to us now.}
Сообщение: <<Красный топор звенит о дерево.
Трапеза у ветвистых рогов старой рощи, дальше по шестипалым следам с запахом какао>>.
А мы выдвигаемся.

В кои-то веки Кхохо не стала спорить.
Звон топора означал всеобщую мобилизацию.
Красный --- кодовый цвет ноа.

--- Давненько не звенел красный топор, да, Донышко? --- так, что слышали только я и Ситрис, шепнула женщина.

--- Сало на твоей заднице было тоньше и мягче, --- с грустью согласился Ситрис.
Но я чувствовал, что его печаль относилась отнюдь не к сказанному.

\section{[-] Лимнэ}

\textspace

--- Идём-ка, дитя, --- тихо позвал Ситрис Чханэ.

--- Какая она тебе <<дитя>>? --- тут же вскинулся Кхарас.
--- Субординация!

Ситрис невозмутимо показал вождю оскорбительный жест и увёл Чханэ за дерево.
Ликхэ нахмурилась.
Похоже, её <<дитя>> Ситрис не называл.

Вернулась подруга очень скоро.
Я её не сразу заметил среди листьев.

--- Он меня раскрасил полосками, начернил ресницы и рассказал, как лучше вести бой в <<царстве орхидей>>, --- радостно рассказала Чханэ.
--- А ещё дал вот это.

Чханэ держала в ладонях лимнэ --- камуфляжную повязку на глаза с козырьком.
Я слышал о таких вещах, но столкнулся с ними впервые.

--- Чтобы глазами не светила, --- пояснила подруга.
--- Ситрис сказал, что видел такую на знакомом... --- Чханэ понизила голос, --- на знакомом пирате, который был оцелотом.
Ситрис не уверен, что сшил так, как надо, но носить можно.

--- А ты раньше такие не носила? --- удивился я.
--- Ты же...

--- Я же пустынная кошка, --- усмехнулась Чханэ.
--- В джунглях бывала редко, тем более, Скорпиону особо нет нужды прятаться.

\section{[-] Боевой вождь}

\textspace

Кхарас методично объяснил план боевых действий.
Многие, впервые попав в храм Тхитрона, удивлялись тому, что бесхитростного и прямолинейного воина выбрали вождём, Чханэ не была первой.
Тому существовали веские причины.
Тактику Кхарас строил грубо, но очень добротно и надёжно, гармонично сочетая нападение и оборону.
От него сложно было дождаться сомнительных действий и рискованных приёмов, ставящих под угрозу отряд.
Подобными делами в частном порядке занимался Ситрис, имевший большой опыт дерзких диверсий, и в особо сложных ситуациях Кхарас просто исключал Ситриса с подопечными из плана, позволяя тем действовать самостоятельно.
Это давало отличный результат, и ни у кого даже мысли не возникало оспорить роль вождя, ведь как бы хорош ни был диверсант, закрытая спина --- прежде всего.

\section{[-] Прикуси язык}

\textspace

Для общения мы использовали язык, имитирующий крики лесных животных, пение птиц, шум листвы и треск деревьев.
Длинные сообщения передавались через наруч.
Переговариваться ультразвуком в этой густонаселённой части леса мы не могли --- хака не слышали ультразвук, но мелкие птички начали бы волноваться и выдали нас с головой.
К тому же оставался этот неведомый пришлый...

Маскировочный плащ отлично гасил инфракрасное свечение тела, но долго сидеть под ним в душный летний вечер --- настоящая пытка.
Я сидел, обливался потом в четыре водопада и надеялся лишь на то, что дурманящий аромат орхидей заглушит запах уже начавшего бродить пота.
Кажется, Кхохо что-то говорила на эту тему.
<<Прикуси язык>>.
Я не мог понять, пошутила она или нет, но решил попробовать.

И правда помогло.
После двадцати мучительных секунд руки вдруг стали холодными, а сердце словно провалилось в бездонную дыру.
Я понял, что вполне сочетаюсь по инфракрасной расцветке с деревом даже без плаща.

\section{[-] Схватка в лесу}

\textspace

<<Ликхмас, как дела?>> --- донёсся от Кхараса нежный свист птенца согхо.

<<Враг в прямой видимости у земли в десяти на шесть>>, --- отбил я на наруче.

<<Почему не стреляешь?>> --- удивилась откуда-то Чханэ голубой квакшей.

<<Он может выдать товарищей>>, --- объяснил я щёлканьем мохноножки.

Кхарас мяукнул, делая мне строгий выговор за риск.

<<И вообще, что за ящерицыно семя поставило жреца-новичка на передний фронт? --- загудел наруч, передавая возмущение Ситриса.
--- Мы такие богатые --- молодёжью разбрасываться?
Кхарас, ты куда смотришь?>>

<<Поздно перегруппировываться, --- отозвался Кхарас.
--- Не засоряй эфир.
Ликхмас, постарайся не высовываться>>.

Мои ожидания оправдались через несколько томительных мгновений.
Воин хака сделал неосторожный знак головой.
Я проследил направление его взгляда.
Товарищ моей цели был истинным мастером.
Издалека его можно было принять за обросшую орхидеями ветвь.
Я немедленно поделился наблюдениями.

Ответ от Ликхэ пришёл спустя минуту:

<<Ликхмас, ты в переплёте.
Юнца, которого ты приметил первым, я взяла на прицел.
Второго не видит никто.
Это очень опытный разведчик, и, возможно, ты уже у него на прицеле.
У тебя один шанс.
Не промахнись>>.

Нежное грудное воркование кота-оцелота прямо-таки сочилось страхом.

Я медленно, горошина за горошиной, поднимал духовое ружьё.
Выверил угол, сосредоточился на слабом дуновении ветра, турбулентных потоках вокруг листьев --- и, резко отклонившись в сторону, выстрелил.

Ружьё врага дёрнулось в мою сторону, но было поздно.
Стрела вошла точно в не прикрытое пластиной бедро воина.
Я ожидал крика отчаяния, но вместо этого прозвучал чёткий гортанный речитатив, в которой, по-видимому, сообщалось о примерном направлении полёта моей стрелы.
Хака действительно оказался опытным разведчиком.

<<Ликхмас, беги!>> --- одновременно зашипела Чханэ и рявкнул Кхарас.

Когда тебя обнаруживают, на церемонии время не остаётся.
Я ломанулся через заросли, ломая ветки, стуча сапогами и полностью оправдывая данное мне имя.
Над ухом просвистела чья-то шальная стрелка, я повернул голову, уворачиваясь от неё... и едва успел заметить летящий в спину томагавк.
Я схватил его за рукоять левой, одновременно правой поднимая духовое ружьё.
Полный отчаяния стон спустя мгновение подтвердил, что стрела достигла цели.
И тут... дерево, просто идеальное, чтобы спрятаться.
Я пустил шум ломающихся ветвей вместе с вражеским томагавком куда-то в сторону, а сам вжался в тёплое, влажное дерево и замер.

<<Я в порядке, ещё одного зацепил по пути>>, --- обрадовал я товарищей через наруч.
Кхарас злобно закудахтал, выражая крайнюю степень недовольства.

<<Второго я уложила>>, --- сказала Ликхэ.

<<Один в пяти на восемь убит.
Кто-нибудь ещё видит?>> --- завибрировал передатчик словами Чханэ.

<<Двое в кроне в тринадцати на два.
Одного убил, второй на земле со сломанной спиной, --- сообщил Ситрис.
--- Молодец, Ликхмас, устроил переполох.
Но больше так не делай>>.

Все остальные ответили отрицательно.

<<Остался один, --- констатировал Кхарас.
--- Ищите>>.

Однако поиск ничего не дал.
Последний хака затаился основательно.

\section{[-] Молодёжь}

\epigraph
{Жизнь не похожа на сказания.
Когда люди слушают истории про великие битвы, они радуются победам и грустят над поражениями героя.
О безымянных, безликих соратниках и врагах героя, которые тысячами погибают в каждой битве, которые приносят песчинки своих жизней на весы судьбы, --- о них никто не вспоминает.
Но любому из нас отведена в истории именно эта роль.
Героями становятся единицы, и не по делам, а более по прихоти шкёльда...}
{Хервар Лонгсин-Храш, шкёльд (сказитель) двора Лодемора Сурового}

Мы спустились с деревьев.

--- Это не опасно? --- поинтересовался я у Кхараса.

--- Нет, --- ответил Кхарас.
--- Хака знает, что сейчас попытка убить любого из нас равносильна самоубийству.
Он будет сидеть, надеясь, что мы уйдём.

--- Но мы не уйдём, --- ввернула Чханэ.

--- Судя по его расположению, это должен быть тот самый координатор, --- объяснил ей Кхарас.
--- Сели он или ноа, но его просто необходимо взять в плен.
Ликхмас, иди обработай сломанного.
Насчёт твоего поведения поговорим потом.

Я, раздвинув заросли папоротника, прошёл к другой группе воинов.

Молодой хака, почти мальчишка, лежал на земле и тяжело дышал ярко-красным ртом.
По его лицу градом катился пот, смывая грязь, сок листьев и боевую раскраску.
Ситрис пытался снять с него доспехи, не тревожа перелом.

--- Дети пришли за другими детьми, --- сказал он, обернувшись ко мне.
--- Но не играть.

--- Точно.
Молодёжь одна, --- с некоторым сожалением пробормотала Хитрам, разглядывая упавшего.
Она успела стащить пятерых хака под дерево и совершить над ними обряд.
--- Красивые, весёлые.
И чего вам дома не сиделось, мальчики?

--- Хотели славы и почёта, --- хмыкнул Ситрис.
--- Ликхмас, командуй.

--- Ага, ещё церемониться с ним, --- сказала Ликхэ.
--- Забудь, Лисичка, я его так потащу.

Ликхэ грубо схватила пленника и взвалила себе на спину.
Хака взвыл и задёргал руками, его ноги висели безжизненными плетями.

--- Ты что делаешь, рыбья печень? --- крикнул Ситрис, нарушая все правила поведения в джунглях.
Птицы на ближайших деревьях подняли дикий гомон.
--- Положи немедленно!

Ликхэ обиженно посмотрела на воина и бросила пленника на землю.
Тот снова дико, с подвыванием заорал.
Ситрис уложил его как можно ровнее, вытащил флягу с водой и облил голову пленника.
Затем яростно посмотрел на Ликхэ:

--- Иди веток на носилки и шину нарежь, животное!
Ликхмас, --- обратился он ко мне гораздо спокойнее, --- действуй.
Ты знаешь, что нужно.

Ликхэ злобно сплюнула и ушла в лес.

--- Поплюйся тут мне ещё, соплячка, --- вполголоса пробормотал Ситрис.
--- Нелюди нынче какие-то пошли...
Нельзя же так с живыми...

Я покрутил в руках амулет и открыл крышечку Красного Сана.

--- Тихо-тихо, малыш, --- остановил меня Эрликх.
--- Меньше каплю делай.

--- Хай, не мешай парню, он лучше нас знает, --- поморщился Ситрис.

--- А всё-таки, --- полушёпотом добавил Эрликх.
--- Это хака, они слабее.
Ты же знаешь, что кураре их не парализует, как нас, а просто убивает?
Вот.
А я сам видел, как от обычной для нас дозы Красного хака засыпают и больше не просыпаются.
Да, вот такую, маленькую.
Ещё меньше.
И под язык капай...

Спустя десять минут пленник лежал на жёстких носилках в маковом полусне.
Я в последний раз проверил наложенную шину.
Ликхэ стояла поодаль и исподтишка бросала злобные взгляды на Ситриса.
Вдруг выражение её лица изменилось, и она словно невзначай вытащила духовое ружьё и начала его чистить.

<<Я вижу последнего, --- зажужжал наруч словами воительницы.
--- В пятнадцати на три.
Беру прицел>>.

Никто из нас даже не изменил позы.
Я осторожно скосил глаза.
Молодой хака в трофейном шлеме сели, сидя в розетке большого папоротника-опахала, смотрел на нас из-под раскрашенных полуприкрытых век.

<<Если не уверена, что это не хака, не используй кураре, --- предупредил Эрликх.
--- Возьми болевые стрелы>>.

<<Безумный меня растерзай, я болевые стрелы уронила, пока бегала, --- пожаловалась Ликхэ.
--- Может, в бездну?
Давайте я его просто убью?>>

<<Нет, --- отозвался откуда-то Кхарас пением ,,красного шара``.
--- Он нужен живым>>.

<<Потеряла стрелы.
Как вы с Ликхмасом вообще дожили до своих лет?>> --- философски прожужжал Ситрис.

Хака как будто понял, о чём идёт речь, и припустил на юго-восток со всех ног.
Ликхэ, на рефлексах выпустив ядовитую стрелку в пустое место, громко выругалась.
Мы с Ситрисом бросились в погоню.

\section{[-] Погоня}

<<Да он что, стальной? --- бросил Ситрис на десятой минуте погони.
--- Ликхмас, я на пределе.
Будь осторожнее, а ещё лучше --- возвращайся.
Ну его к свиньям>>.

Хака оказался гораздо выносливее своих товарищей.
Я знал, что любой сели способен обогнать хака и накоротке, и на большой дистанции.
Но наш преследуемый петлял между ветвей, стволов и лиан как заяц и даже не думал уставать.

Ситрис отстал --- тело его подвело, сказывался возраст.
Я же был в самом расцвете юности, комплекция позволяла проскальзывать в самые узкие щели, и погоня меня только раззадорила.

Но вскоре хака начал проявлять признаки усталости.
Всё чаще он вместо бега затаивался в зарослях, и мне приходилось напрягать все свои чувства, память и ум, чтобы понять, куда делся враг.
Вот дорогу мне перегородил извивающийся в предсмертных судорогах молодой питон --- хака потратил несколько бесценных секунд на борьбу с неожиданно схватившей его рептилией.
Несколько раз враг пытался выстрелить в меня из духового ружья, и в последний я расколол его хрупкое оружие брошенным томагавком.
На дереве, за которым прятался преследуемый, остались свежие брызги крови.

<<Охотник никогда не должен забывать, что у отчаявшейся, уставшей жертвы когти и клыки растут отовсюду, откуда могут расти>>, --- вспомнил я слова Конфетки и мысленно призвал на тренера благодать лесных духов.
Наука, которую он передал мне, сотням и тысячам стоила жизни.

Дальше идти было куда проще --- хака источал отчётливый, неистребимый запах страха, а на листьях и стволах изредка оставались мелкие кровяные капельки.
Можно было даже не торопиться --- я знал, что преследуемый бежит из последних сил и обязательно должен остановиться, чтобы решить дело боем.

Вскоре, как я и ожидал, хака решил перейти в контратаку.
Он сделал кольцо из запаха вокруг дерева, а сам взобрался по лиане и повис вниз головой с ножом наготове.
Но я был начеку и встретил прыгнувшего сверху воина ударом в подбородок.
Хака крякнул и мешком рухнул на землю, выронив нож.
Моя фаланга уткнулась ему в грудь.

--- Natra mis\FM, --- сразу обозначил я свои намерения.
\FA{
\ml{$0$}
{[Следующее твоё] движение будет последним (цатрон).}
{[Your next] move will be [your] last (\Tesatron).}
}
Но пленник, похоже, даже не собирался двигаться.
Хака дышал с надрывом, хрипло и часто. Он сорвал на бегу застёжки нашейника, и шлем при падении слетел с головы, открыв молодое мужественное лицо с по-дикарски яркими глазами.
По левой щеке бежала кровь --- мой томагавк рассёк лицевую артерию.

Я примотал пленника к ветке его же роскошными волосами, затем занялся обыском.
В сторону полетели два гибких потайных ножа, наруч со спрятанным в коже лезвием, два комплекта стрел и верёвочная удавка, искусно вплетённая в ткань рубахи.
Затем я занялся перевязкой артерии --- хака потерял много крови.
Пленник, увидев зажим и отбелённый волос, угрюмо закусил воротник, и пока рана не была зашита, я не услышал от него ни звука.

\section{[-] Благородный баньян}

Вскоре я повёл хака в рощу благородного баньяна, оказавшуюся неподалёку от места схватки.
Могучее дерево, похожее на зал с колоннами, хранило свою землю от чужих.
На мшистой почве не было видно ни насекомых, ни травянистых.
Говорят, что в таких рощах сели устраивали храмы...
Догадку подтвердил лик Тёплого Хетра, уже почти заросший, но ещё угадывавшийся на стволе древнего дерева.
Я встал и ковырнул мох под ним --- на свет выскочил край изящной жертвенной тарелочки из ослепительно белого, <<фарфорового>> нефрита, которым была богата Старая Челюсть.

Память услужливо предоставила сведения --- на нынешней границе земель хака и сели были развалины древнего города с утраченным названием.
Гибельные руины --- так их называли теперь.
После одного из землетрясений посреди города прошёл разлом;
западная возвышенность осталась пригодной для проживания, но сели оставили город, опасаясь набегов --- разлом разрушил всю тщательно организованную линию обороны города.
Хака же суеверно избегали величественное заброшенное место, боясь гнева мёртвых.

Я рассеянно начал выкапывать тарелочку из мшистой почвы.
Пальцы ощутили чёткую, умелую вязь орнамента на её ребре.
Связанный хака исподлобья с интересом наблюдал за мной.

\ml{$0-[hk]$}
{--- Bv\o h ant naani? --- спросил я.}
{``\textit{Bv\o{}h ant naani?}'' I asked.}

Пленник молчал.

--- Ликхмас ар'Люм, --- добавил я, положив руку на грудь.
Хака посчитает это слабостью, но, как говорила Кхотлам, хозяина положения не должны заботить такие мелочи.

Пленник подумал, хмыкнул и наконец-то открыл рот.

\ml{$0-[hk]$}
{--- Imzhu.}
{``\textit{Imzhu.}''}

<<Замечательное начало знакомства, --- вспомнил я весёлые слова кормилицы, читавшей очередную дипломатическую переписку.
--- Оба уверены, что победили, а для вояк это главное>>.

--- Bv\o{}h laer mimi, Imzhu? --- спросил я и, показав паёк, разломил его напополам.
Одну половину засунул в рот, вторую протянул Имжу.
Воин насмешливо сощурился.

--- Думаес, ты купись меня еда, мальсик сели?

Парень коверкал слова и неприятно шепелявил, но говорил на сели, и, несмотря на множество ошибок, я его понимал.
Я вгляделся в пленника внимательнее.
Зелёные, подкрашенные углём глаза, искривлённый горько-агрессивной насмешкой небольшой рот, высокие скулы...
Едкий пот смывал смуглину и боевую раскраску, оставляя на лице дорожки чистой ярко-золотистой кожи.
Как и ожидалось, хака оказался поддельным.
Слово <<мальчик>> он произнёс с легко узнаваемой оскорбительной интонацией;
оно подсказало, что собеседником можно отлично манипулировать с помощью слов, относящихся к возрасту.

--- Ты говоришь на сели?
Отлично.
Тогда давай так.
Я отведу тебя к нашим.
Вряд ли ты после этого останешься в живых, но я знаю, что ты хочешь жить, и жить как можно дольше.
Для этого взрослому человеку следует \textit{mimai}\FM.
\FA{
Пить и есть (язык хака).
}
Принять \textit{mimai} не есть предательство.
\textit{Bv\o{}h sare \o{}i laer limaa}\FM?
\FA{
Ты понял или повторить? (язык хака).
}

Хака молча смотрел на меня.
Я поднёс половину пайка к его губам, и пленник жадно вцепился в неё.
Прожевал и проглотил.

--- Вода, --- сказал он.

Я влил ему в рот воды из фляжки.
Хака выпил предложенное.

Солнце, похоже, уже прошло зенит.
Следовало торопиться обратно.
Я начал собирать амуницию пленного в его мешок.

Вдруг моё внимание приковал к себе шлем.
Зелёная кожа, отделанная чёрными змеями, классический нашейник народа сели, знак Кхара-защитника, нарисованный загущённой кровью.
Явно работа известного мастера-сели, дикари таких никогда не делали.
Я перевернул его и посмотрел на изящный тиснёный иероглиф, украшавший охвостье нашейника: <<Солнечная, из дома Песка, из города Отравленного Вина>>.

Митхэ ар’Кахр э’Тхартхаахитр.

В голове моей, и так не особо свежей после погони, помутилось.
Перед глазами пронеслись виновато отдающая мне записку Кхотлам, Чханэ, намёки и недосказанности служителей Храма...
Я подошёл к Имжу, показал ему шлем и вытащил нож.
Пленник напрягся.

--- Скажи, где человек, которому принадлежит этот шлем.

Имжу промолчал и повращал глазами, ища возможность сбежать.
Я взял его за волосы и поднёс нож к бешено пульсирующей сонной артерии.

--- Кто твоя родильница?

--- К-кто? --- переспросил Имжу.

Я попытался вспомнить подходящий термин, бывший в ходу у хака.

--- \textit{Mama.}
\textit{Mama n\o{}nt?}

--- \textit{Zhyu gerren si}\FM, мальсик сели.
\FA{
Я ребёнок вождя (язык хака).
}
Я не понимась...

Я мысленно упрекнул себя за незнание языков и задумался ещё раз.

--- \textit{Mama na... n\o{}ni gerron w\o{}izh}\FM?
\FA{
Твоя родильница --- женщина, связанная с вождём законом племени? (язык хака).
}

--- \textit{Zhyu gerren si}, --- повторил Имжу, оскорблённо шмыгнув носом.

Насколько я знал, у хака женщины никогда не занимали ведущие посты в племени.
В моём сердце, как ядовитая лоза, росли печаль и подозрение.

<<Запомни, Ликхмас, --- зазвучал в моей голове голос Конфетки, --- я учу тебя лишь владеть телом, но знания, которые у тебя в голове --- оружие куда более страшное, чем целая армия.
Цивилизованные, разумеется, будут принимать тебя лишь как опасную фигуру, но дикарь будет смотреть на тебя, как на потустороннюю силу в человеческом обличье.
И любой враг, с которым ты встретишься, познает страх, едва увидев жреческое одеяние, даже если и не признается себе в этом страхе>>.

--- Так, молодой хака, --- я рывком поднял парня на ноги и бросил спиной на дерево.
--- Отведёшь меня к вождю своего племени.
Если сделаешь это --- получишь жизнь и свободу.
Если же нет...

Я провёл лезвием по шее хака. Он неподдельно вздрогнул, испугавшись такой крутой перемены моих намерений.

<<Ликхмас, Ликхмас, ты же мне поклялся!
``Я не оскверню себя пыткой!..''>> --- взвыл воображаемый Конфетка где-то на окраине сознания.
Но меня было уже не остановить.

--- Я тебя уничтожу.
Ты будешь умирать медленно и мучительно, каждая капля крови, покидающая твоё тело, будет причинять тебе невыносимую боль.
Когда же ты умрёшь, я вырежу у тебя на лбу знак каменного духа.

Меня колотило в почти священном трансе.
Слова вылетали из моего рта, словно потревоженные летучие мыши из чёрного зева пещеры.

--- Ты умрёшь обессиленным и никогда, слышишь, больше никогда не станешь мужчиной и воином!

\begin{verse}
\textit{Bv\o{}h her brai lar in izhe,}\\
\textit{sht\o{}h mart sai likar ciche,}\\
\textit{es di izhu aes sart}\FM,
\end{verse}
\FA{
У тебя не будет посмертия в живых лесах, мёртвый камень станет твоим телом, пока планета не рассыплется в прах (язык хака).
}

--- хрипло пробормотал я первый пришедший на ум экспромт.

Хака крайне редко прибегали к угрозам;
они считали, что угроза равносильна проклятию --- воззванию к суду Безумного, которые могли принять любое решение, вплоть до смерти проклинающего.
Стихотворная форма же в тысячи раз усиливала действие слов.

Судя по лицу пленника, угроза возымели своё действие.
Зелёные глаза стали почти чёрными, целая щека посерела.
Губы молодого воина дрогнули.

--- Ты... ты не мось, мальсик сели!
Это мось только слуга Безумный!

--- \textit{Si laar}\FM, --- полушёпотом сказал я и покачал перед лицом пленника амулетом Сана.
\FA{
Я могу (язык хака).
}

Хака несколько секунд смотрел на изящный бронзовый диск и, сглотнув, проворчал:

--- Хоросо.
Я вести.
Шаман сели хотесь видесь восдь, а хака хотесь жись.

\section{[-] Как взрослеют дети}

\textspace

Я протащил Имжу мимо ощетинившейся копьями толпы и втолкнул его в хижину.
Тут же один воин легонько шлёпнул клинком копья по моему запястью, ещё пара упёрлась мне в бока.
Я сказал <<Хэ>>\FM\ и сильнее прижал острие ножа к шее пленника.
\FA{
Междометие, предупреждающее об опасности для собеседника или третьего лица, в зависимости от контекста.
}
Копья отодвинулись на полпяди.

В хижине было тепло и душно.
Казалось, воздух закаменел от выедающего глаза запаха дыма.
В дрожащей мгле, подсвеченной неверным пламенем костра, я разглядел невысокую женскую фигуру.
Лицо женщины наполовину закрывали спутанные кудрявые волосы, плечи и посадка головы выдавали бывалого воина.

--- Зачем ты пришёл, тари? --- чёрные губы женщины открылись, показав зубы, похожие на обмытый рекой кварц.
Верхнего резца не хватало.
Отсутствие акцента подтвердило мои опасения.

--- Встречай потомков, Митхэ ар’Кахр, --- с этими словами я толкнул Имжу, и молодой воин упал на колени.
Сразу же три копья упёрлось в кожу моего нашейника.

--- \textit{He gt, bv\o{}h iszh!} --- бросила женщина сорванным командным голосом, и копья медленно отодвинулись.
Я осторожно сделал шаг к пленнику, который, тяжело дыша, прислонился к резному деревянному столбу.

--- Потомков? --- как ни в чём не бывало осведомилась чернявая ведьма.

Я медленно отцепил от пояса змеиный шлем и бросил жалобно звякнувший доспех перед ней.
Она подняла его, перевернула, и на протянутую ладонь упал кусочек кожи.
Подведённые углём глаза расширились.

\ml{$0-[hk]$}
{--- \textit{Ke, yghd.}}
{``\textit{Ke, yghd.}''}

Воины поспешно вышли из хижины и опустили травяную занавеску.
В душном воздухе остались только приглушённый шум дождя, потрескивание брёвен и тяжёлое дыхание Имжу.
Женщина смотрела на меня во все глаза.

--- Кто ты?

--- Ликхмас.
Могла бы и вспомнить.

Родильница опустила голову, задумчиво теребя в руках письмо.
Наконец она решилась прервать молчание:

--- Хай, так ты...

--- Объясни мне одну вещь, --- перебил я, и женщина замолкла, словно окаменев.
--- Почему хака получил твой шлем и твоё плечо, а сели --- полуистлевшую чернильную ложь?
Почему хака получили твой меч, а сели --- сожжённые деревни и печаль?

Митхэ бесстрастно смотрела на меня, поигрывая шлемом. Я вздохнул.

--- Ладно, это твоё дело. Расскажи, что с моим дарителем.

Митхэ пожевала губу, грустно усмехнулась и, присев, пригласила меня сесть напротив.

--- Садись.
\ml{$0-[hk]$}
{\textit{Imzhu, bv\o h legh.}}
{\textit{Imzhu, bv\o h legh.}''}

--- \textit{Mame, segh ni?..} --- попытался возразить юнец.

--- \textit{Imzhu!} --- родильница повысила голос.

Имжу покорно сел рядом с нами, успев надменно крякнуть в мою сторону:

\ml{$0-[hk]$}
{--- \textit{Ghe di bv\o{}h, zhyu gerren si...}}
{``\textit{Ghe di bv\o{}h, zhyu gerren si---}''}

--- А я --- тот, кто сохранил жизнь \textit{gerren s\o{}ine}, когда уже занёс над ним нож, --- в тон юнцу ответил я.
Тот замолчал.

--- Прошу прощения, Ликхмас... ар’Люм, кажется?
Мой питомец ещё не вышел из своего возраста.
Но он умён и храбр, из него выйдет хороший воин хака.

--- Я не держу обид на побеждённых, --- ответил я чьей-то давно услышанной фразой.

Родильница вдруг улыбнулась и помолчала несколько секхар.

--- Твой боевой наставник --- благородный человек, и ты --- достойный его ученик.
Немногие смогли бы проделать путь, который прошёл ты.
Что ж.
Прости и то, что я не предлагаю тебе пищу.
Сейчас сели --- враги хака, и мы не можем есть за одним столом, пока горит костёр войны.

\ml{$0-[se]$}
{--- Это не моя вина.}
{``That's not my fault.''}

\ml{$0$}
{--- И всё же прошу меня понять.}
{``Even so, I ask you to understand.}
\ml{$0$}
{После того, как ты привёл моего питомца, кхм, за шкирку, у воинов достаточно поводов сомневаться в моём авторитете.}
{You've taken my nurseling---khm---by the scruff of the neck, and my warriors have got enough reasons to doubt my authority.}
\ml{$0$}
{Что с остальными?}
{How are the others?''}

\ml{$0$}
{--- Все, кто не вернулся, мертвы.}
{``Everyone who hasn't returned is dead. }
Мы убили их как захватчиков, но оплакали как своих, --- ответил я ритуальной фразой дипломата, как учила меня Кхотлам.

Митхэ вдруг снова улыбнулась той же улыбкой.

\ml{$0$}
{--- Это было честно, --- наклонила голову она.}
{``That was fair,'' she nodded.}
\ml{$0$}
{--- Итак, ты хочешь узнать тайну этой записки?}
{``Well ... are you ready to know the secret of this note?''}

\ml{$0$}
{--- Я в нетерпении.}
{``I've been waiting.''}

Родильница помолчала, глядя в огонь.
Имжу тем временем нагрел какой-то сыпучий камень и начал затемнять им кожу.
Затем подправил сурьмой форму глаз и накрасил ресницы, явно пытаясь скрыть яркий зелёный блеск.
Пока Митхэ собиралась с мыслями, чужак по крови стал неотличим от хака.

\ml{$0$}
{--- Я была воином, --- заговорила наконец женщина.}
{``I was a warrior,'' the woman finally began.}
\ml{$0$}
{--- Одним из лучших воинов, не только в Тхартхаахитре, но и во всех Восточных джунглях.}
{``One of the best, not only in \Tchartchaahitr, but in the whole Eastern jungle.}
\ml{$0$}
{Хотела бы я преувеличивать, но увы.}
{I wish it'd be an exaggeration, but it never was.''}

\ml{$0$}
{--- Почему <<увы>>?}
{``Why would you wish that?''}

Митхэ посмотрела на меня с горькой улыбкой.

--- Как же я рада этому вопросу...
Ты никогда не был ни в чём лучшим?

--- Учитель говорил, что для жреца важно уметь всё понемногу и сохранять руки нежными, --- сухо сказал я.
\ml{$0$}
{--- Я умею всё понемногу, поэтому я не лучший ни в чём.}
{``Many things I can do, and none I'm the best at.''}

\ml{$0$}
{--- Такой подход не лишён смысла.}
{``That makes sense.}
\ml{$0$}
{Когда ты умеешь что-то лучше всех, это больше ответственность, чем возможность.}
{If you're the best at something, it means more responsibilities than abilities.}
Особенно если дело касается военного ремесла.
Я командовала отрядом чести в боях против идолов, пылероев, случалось биться и с людьми --- тенку, хака... --- она кивнула на дверь.
--- Командовать не особо любила, но приходилось --- ответственность в огороде не закопаешь.
\ml{$0$}
{Разумеется, у меня было много мужчин, были и дети.}
{Of course I've been with lots of men, I had children.}
\ml{$0$}
{Всех и не упомнить...}
{Too many to remember them all ....''}

\ml{$0$}
{--- Ага, --- с горечью проворчал я.}
{``Yeah,'' I bitterly muttered.}
\ml{$0$}
{--- И каждому доставалась в наследство записочка?}
{``I suppose none of them left disinherited, without a nice note?''}

\ml{$0-[hk]$}
{--- Ты не оскорблясь мой \textit{mama}! --- зло выпалил Имжу, не поняв, о чём речь, но верно истолковав тон.}
{``You will not insult my \textit{mama} in front of me!'' Imzhu exploded; he hardly understood my words, but he correctly interpreted my tone.}

\ml{$0-[hk]$}
{--- \textit{Imzhu, ghyti ni.}}
{``\textit{Imzhu, ghyti ni.}''}

Парень застыл.
Выражение его красивого диковатого лица не сулило мне ничего хорошего.
Родильница помолчала, жуя губу.

--- Потом в Тхартхаахитре я встретила Атриса.
Сам он утверждал, что у него нет имени, а это прозвище --- <<журавль>>\FM\ --- ему дали женщины на рынке.
\FA{
Прозвище созвучно имени Хатрис, распространённому среди восточных сели.
}
Он был беден, его одежда едва прикрывала тело.
Но он чудесно играл на цитре.
Его пение было тихим, почти неслышным, но услышавшие приходили послушать его снова и снова.
Атрис похитил мои мысли, и я попросила его пойти со мной.
С тех пор он шёл рядом.
Лицом он был солнцем, глазами он был сыном змеи, а ещё у него были тонкие длиннопалые руки.
Полгода пролетели, как согхо над рекой.
Он всё говорил, что я не для битвы, уговаривал меня бросить военное дело и осесть где-нибудь в спокойном месте...

Митхэ говорила нараспев, намеренно искажая тоновый рисунок и употребляя самые изысканные поэтические обороты.
Среди хака были те, кто немного понимал язык сели, и знать тему разговора им было совершенно ни к чему.

--- Его ли цветы дали завязь?

Родильница опустила голову.

--- Об этом знают лишь Безумный да Хри-соблазнитель.

--- Что же знаешь ты?

--- Меня приревновал один из жрецов.
Он давно хотел быть со мной, но я отвергала его.
С хака тогда подписали перемирие, и в городе было большое пиршество.

--- Пиршество? --- удивился я.
--- А как же?..

--- Молчание было нарушено.
Тот дождь наверняка помнят многие, ты можешь узнать о нём из рассказов людей.
Мы с Атрисом веселились, занимались любовью и мечтали о доме.
Нашем доме.
Но на второй день тот жрец что-то подсыпал в моё вино.
Одурманил меня.
Он насиловал меня и поил этим зельем.
Как долго это продолжалось --- я не знаю, несколько дней...
Очнулась случайно в его покоях --- он просто забыл в очередной раз опоить меня.
Выбежала в зал --- мои воины очень удивились.
Они решили, что я всё-таки поддалась на уговоры Атриса и уехала с ним на Запад.
А послушник, с которым я подружилась, сказал по большому секрету, что цитру моей жизни отдали вождю хака, как раба.

Митхэ взглянула мне прямо в глаза.

--- Я сразу же отправилась на его поиски, со мной пошли трое товарищей.
<<В сезон дождей созрели дыни...>>\FM\
\FA{
<<... Я родила тебе дитя, о сладкая лиана, опутавшая голову мою, о гость желанный чрева моего>>.
Отрывок стихотворения Эрхэ Колокольчик.
}
Я давно знаю Кхотлам --- она надёжный человек.
Из Тхитрона я пошла в путь в сомнениях, без сил.
Я не знала, жив ли Атрис, и не знала, кто засеял поля --- <<захватчик ли, хозяин ли исконный>>\FM.
\FA{
<<Мы связаны, судьбу не изменить, и важно ль, кто на нас надел оковы --- захватчик ли, хозяин ли исконный?>>
Отрывок из <<Легенды об обретении>>, песня Ликхмаса.
}
Я надеялась на лучшее.
Доспехи натирали моё тело, я задыхалась при ходьбе, молоко горело, как сухая трава...
Но я думала, что, даже если отдам дух Безумному, никто не будет на меня в обиде.

--- Атрис мёртв?

--- Да, Ликхмас-тари.
Атрис мёртв.
Вождь отказался отдать мне его, и я своими ушами слышала, как цитра моей жизни издаёт на алтаре свой последний аккорд.

Я промолчал.
Родильница нервно стучала по столу рукой, и я различил знакомый уже до боли барабанный язык:

<<Я стала w\o izh вождя.
Я взяла власть в племени.
А потом я убила его.
Медленно отравила, чтобы он познал всю мою боль>>.

<<Я подставил тебя, я знаю>>.

<<Выкручусь.
Бывало и похуже>>.

<<Вместе сподручнее>>.

Завуалированный вопрос остался без ответа.

--- Объясни, почему ты был так зол на меня? --- спросила Митхэ.
--- Чем я заслужила твой гнев?
Воспитывать детей --- не воинский удел.

--- Не нужно было писать ничего этого, --- тихо сказал я.
--- Мне было бы достаточно легенд о тебе.

--- Я поддалась чувствам, --- извиняющимся тоном сказала родильница.
\ml{$0$}
{--- Я правда думала, что в тот раз всё будет по-другому.}
{``I really thought that time things could be different.}
\ml{$0$}
{Что я осяду, что у меня будет место, которое я назову домом...}
{I thought I'd settle down, I'd have a place to call it home ....}
\ml{$0$}
{Трудно расставаться с мечтой, даже если ты ясно видишь, что дела идут как обычно.}
{It's so hard to bury your dream, even if things obviously go as usual.''}

Разговор себя исчерпал.
Я поднялся на ноги и вдруг понял, что сейчас мне предстоит пройти через поселение, полное ненавидящих меня людей.
Осознание свалилось на меня, словно гранитная скала, на мгновение лишив меня смелости, гордости и способности соображать.
Видимо, это отразилось на лице, потому что оба --- и Митхэ, и Имжу --- смотрели на меня во все глаза.
Митхэ выглядела встревоженной, Имжу --- ошеломлённым.
Я тут же натянул обычную холодную маску.

--- Ты мне помогла, --- сказал я, направляясь к двери.
--- Надеюсь, что больше твои стрелы не обратятся против Тхитрона.
Если моя рука вновь встретит твоего питомца в бою --- он умрёт.

--- Это не мне решать, Ликхмас ар’Люм, --- Митхэ бросила нервный взгляд на входной клапан.
\ml{$0-[hk]$}
{\textit{Haghp\o{}th!}}
{\textit{Haghp\o{}th!}''}

Зашёл пожилой, чисто выбритый хака и коротко поклонился.

--- \textit{He gt, gh\o{}r thenis is litaa}, --- резко приказала родильница.

Хака посмотрел на неё с удивлением:

--- \textit{Gerre, bv\o{}h ni tras, ygh... \Seli?}

\ml{$0-[hk]$}
{--- \textit{Eh \Seli\ mitygh magh zhyu.}}
{``\textit{Eh \Seli\ mitygh magh zhyu.}}
\ml{$0-[hk]$}
{\textit{Berth p\o{}h ni haawr.}}
{\textit{Berth p\o{}h ni haawr.}}
\ml{$0-[hk]$}
{\textit{Gh\o{}r re thenis is litaa e.}}
{\textit{Gh\o{}r re thenis is litaa e.}''}

<<Он искал своего дарителя>>.
В глазах хака это неплохая причина для столь дерзкого прибытия --- они очень ревностно блюдут кровные узы.
Вдобавок Митхэ решила послать со мной дипломатическое сообщение, удвоив мой иммунитет.

Родильница сделала вошедшему левой рукой неизвестный мне знак.
Однако стук пальцев правой предназначался для меня: <<В лесу тебя будут ждать>>.
Мы с вошедшим кивнули почти одновременно.

Митхэ ценили за боевой опыт и коварство, но не все хака доверяли вождю-сели.

Выходить первым было чистым самоубийством.
Я приподнял занавеску, пропуская воина вперёд, и вдруг вспомнил ещё кое-что.

--- Скажи мне, Митхэ ар’Кахр.
Тот жрец...

--- Он не встал с ложа, на котором насиловал меня, Ликхмас-тари.
Стены кельи покраснели от ярости моего клинка.
А рука, которой ты грозишь моему сыну, --- женщина на мгновение светло улыбнулась, её глаза заблестели, --- создана не для сабли.
Оставь ты это.

Я посмотрел на свои узкие длиннопалые ладони и, едва заметно улыбнувшись Митхэ, бросил взгляд на нахохлившегося возле костра Имжу.
Он взирал на меня исподлобья, враждебно, но с лёгкой тенью неуверенности.
Умён, ничего не скажешь.
Я коротко кивнул брату и вошёл в проливной дождь под ощетинившийся копьями конвой.

\section{[-] Чужая флейта}

\epigraph
{Месть кулаком за слово --- уродливый обычай униженных.
Он много хуже первобытного <<око за око>>, ибо от крепкого и меткого словца, что бы ни говорили люди, ещё никто не умер;
боль, причинённая словом, идёт лишь от тайн сердца слушающего.
И король, издав Закон о преступлении словом, своими руками отдал этому обычаю оружие!
Прольются реки крови, господа, попомните мои слова.
Люди ничто не хранят так бережно, как остатки поруганного достоинства.}
{Клаудиу Семито Фризский.
<<Воспоминания о великом тёзке>>}

Ворота поселения остались позади, и я немного вздохнул.
Да, рядом по-прежнему шли несколько молодых воинов, но это было лучше, чем молча смотрящая на тебя толпа.
Разумеется, мне сильно повезло;
мне везёт до сих пор.
Сдай у кого-то из хака нервы --- даже у одного --- и я бы уже бодро бежал по жилам джунглей с бумажным фонариком.

Дарительница осталась в палатке, а вот Имжу вышел, чтобы проводить меня взглядом.
Молодой конвой напрягся.

--- И что это он тебя не убил, крашеная морда, --- пробормотал тот, которого звали Перт.
--- Может, этот сели --- твой брат?

--- Что ты сказал, Перт? --- с убийственной вежливостью осведомился Имжу, пронзив Перта двумя ярко-зелёными лезвиями.
Несмотря на то, что Имжу был безоружен, а Перт нёс аж два копья, желание прокомментировать свои слова у парня сразу отпало --- он что-то злобно буркнул и ускорил шаг.

Впрочем, мои приключения ещё не закончились.
Едва ворота исчезли за изгибом тропы, рука одного из воинов совсем не по-дружески легла мне на плечо.
Я высвободился и выхватил фалангу, успев увернуться от метко брошенного ножа.
Мы замерли друг напротив друга с оружием в руках.

--- Прекратите.

Грубый сорванный голос принадлежал пожилому мужчине, которого подозвала мать.
Хагпот, кажется.
Он стоял, закутавшись в плащ, на дороге, нарочито расслабленно опустив плечи.
Я тут же отошёл в сторону.

--- Хагпот, --- проворчал Перт.
--- Мы должны отправить сели сообщение.
Уйди.

--- Кажется, ты слышал слова вождя из моих уст, Перт, --- непринуждённо заметил Хагпот.
%{``I think you've heard the warchief's words from my tongue, Pert''}
\ml{$0$}
{--- Она велела отпустить мальчика, а не убивать его.}
{``She wants the boy freed, not killed.''}

\ml{$0$}
{--- Он мужчина!}
{``He is a man!''}

--- Посмотри на него, --- Хагпот махнул на меня рукой.
\ml{$0$}
{--- Он маленький, худой и безбородый.}
{``He's short, skinny, and beardless.}
\ml{$0$}
{Конечно, мальчик.}
{Of course he's a boy.}
\ml{$0$}
{В убийстве ребёнка чести нет.}
{There is no honor in killing a child.''}

\ml{$0$}
{--- Все сели безбородые и выглядят как женщины.}
{``All \Seli\ are beardless and womanlike.}
Как и ты, старик.
\ml{$0$}
{Этот одолел отряд Имжу!}
{This one has defeated Imzhu's squad!}
И он расскажет сели о том, что видел и слышал!

\ml{$0$}
{--- Ты мыслишь как бойцовый петух.}
{``You're thinking like a gamecock.}
\ml{$0$}
{Старейшины совершили ошибку, но у них ещё есть возможность восстановить мир.}
{Elders have made a mistake, but they still have the opportunity to restore peace.}
\ml{$0$}
{И кровь тех, кто пошёл за Имжу, будет платой.}
{And blood of those who followed Imzhu will be the price.''}

\ml{$0$}
{--- Какой ещё мир?}
{``What is the damn peace?}
\ml{$0$}
{Какая плата?}
{What is the damn price?}
\ml{$0$}
{Мы будем мстить!}
{We shall seek vengeance!''}

--- На наше счастье сели никогда не мстят сверх меры.
Закон гласит --- одного за одного, не больше.
Мужчину за мужчину, женщину за женщину, взрослого за ребёнка.
Если они захотят уничтожить хака под корень, то сделают это за пару дней, и о наших деревнях до самых Молчащих лесов не останется даже воспоминаний.

\ml{$0$}
{--- Пусть попробуют!}
{``Let them try!''}

\ml{$0$}
{--- Давай начистоту.}
{``Speak freely.}
\ml{$0$}
{Ты сомневаешься во мне или в мудрости вождя?}
{Do you question my words or the warchief's wisdom?''}

--- Мудрости старого дырявого котла, в который запускал ложку каждый встречный? --- скривился Перт.

\ml{$0$}
{--- Ты сомневаешься в мудрости старейшин хака, которые выбрали её боевым вождём?}
{``Do you question wisdom of Haka's Elders who have chosen her for a warchief?}
\ml{$0$}
{Эта женщина в молодости стоила десятерых таких, как ты, грязноротый недомерок.}
{When she was younger, this woman's worth ten of you, you filthy-mouthed shortling.}
\ml{$0$}
{А сейчас она стоит сотни.}
{She's worth hundred of you now.''}

\ml{$0$}
{--- Говорят, что ты и твоя мёртвая потаскуха развлекались с ней втроём.}
{``I was told you, and her, and that your dead slut was having a little fun together.}
\ml{$0$}
{Скажи, так заведено у сели или тебе просто нравятся трущиеся друг о друга мохнатые шкурки?}
{Is that a \Seli\ tradition or you just enjoy rubbing pelts?''}

--- Мне нравится то, чего ты никогда не познаешь, --- ответил пожилой воин.

--- Уйди, старик, предупреждаю последний раз.
Я намерен отправить сели послание, и мне неважно, кто ещё разделит судьбу пленника, --- Перт угрожающе приподнял духовое ружьё.

--- Даже если это ты?

--- Не понимаю, о чём ты, --- духовое ружьё преодолело полпути к губам парня.

Хагпот не двинулся с места.

--- Я одолею всех вас и в честном бою, Перт, --- спокойно проговорил он.
--- Но сейчас не вижу причин сражаться честно.

Из леса, как по сигналу, вышли ещё пять воинов с луками наготове.
Подручный Перта едва успел перехватить его духовое ружьё.

--- У меня прямой приказ вождя, избранного старейшинами, как и у тебя.
Может быть, старшие зря оказали тебе доверие?

--- Можешь съесть доверие своего вождя и испражниться в её палатке.
Там ему самое место.

--- Также закон хака гласит, что за неисполнение прямого приказа полагается наказание, за противодействие --- смерть.
Ты можешь не ценить доверие, но если ты не уберёшься...

--- Тебе воткнут стрелу в спину, гнилой огрызок, --- ощерился Перт.

--- Тебе уже будет всё равно, --- парировал старый воин.
--- А теперь хватит игр, малыш.
Я знаю, под чью флейту ты поёшь.
Его здесь нет, он сидит в тёплой хижине и греет руки у чужого костра, пока его младший нелюбимый сын стоит под прицелами луков.

\ml{$0-[hk]$}
{--- Что ты сказал? --- трясясь от гнева, пробулькал Перт.}
{``What did you say?'' Pert said constrained trembling in anger.}

\ml{$0-[hk]$}
{--- Ничего нового и ничего того, что не было бы правдой.}
{``Nothing new, and nothing wrong.}
Так что перестань трясти тем, чего у тебя нет, иди в селение и молись кому хочешь, чтобы о твоих словах не узнали старейшины и Имжу.
За подобное оскорбление он вызовет тебя на kazha-l\o m, и не найдётся хака, который отведёт от тебя бесчестье или смерть.
\ml{$0-[hk]$}
{Мои люди, --- он окинул их многозначительным взглядом, --- будут молчать.}
{My people---'' he gave them a meaningful look, ``---will not divulge that.}
\ml{$0-[hk]$}
{До поры, разумеется.}
{Until it needs to be, of course.''}

Перт, кипя яростью, смотрел то на меня, то на закутавшегося в плащ Хагпота, потом окинул взглядом суровые лица окруживших его воинов.
Наконец он, сплюнув на землю, буркнул что-то товарищам и побрёл к деревне.

\ml{$0-[hk]$}
{--- \textit{Ke, yghd.}}
{``\textit{Ke, yghd.}}
\ml{$0-[hk]$}
{\textit{Bv\o{}h lei sert.}}
{\textit{Bv\o h lei sert.}''}

Воины поклонились Хагпоту и пошли вслед за удаляющимися юнцами.
Мы со стариком остались на дороге вдвоём.

\section{[-] Старый друг}

\epigraph
{Война --- лишь трусливое бегство от проблем мирного времени.}
{Клаус Манн, Эпоха Последней Войны}

Хагпот подошёл ближе, буравя меня стальными выцветшими глазами.
Если насчёт Имжу я поначалу сомневался, здесь сомнений быть просто не могло --- передо мной стоял не хака.
Неухоженные, криво подстриженные колтуны пегих волос с разноцветными нитками и костяными погремушками, три или четыре серьги в одном ухе, лезвия маленьких ярко-зелёных глаз...
Очень сильно похож на наёмника-ноа --- так же любит цветистые украшения и такой же смуглый.
Впрочем, нет, это не смуглина и не та краска, что использовал Имжу.
Неужели просто копоть от костра?

--- Так значит, тебя зовут Ликхмас, --- шевельнулись тонкие верёвочные губы.
Он говорил на языке сели --- с лёгким акцентом, но не коверкая слов.

Я молчал.
От этого матёрого ягуара можно было ожидать чего угодно, и можно даже не пытаться вытащить оружие.
В нём чувствовались огромная сила и богатый опыт.
На мшистой земле было много следов борьбы, сапог и сандалий, но ни один из следов не принадлежал сапогам мужчины --- он ступал в точности там, где следы оставить было нельзя.
Я знал только двух человек, которые ходили так же, <<птичьим шагом>> --- одним из них был Ситрис, а про другого рассказывала Кхотлам...

--- А тебя зовут Хагпот.

Старый воин поднял бровь и посмотрел на меня с интересом.

--- Ты наблюдателен, но тебя подвело незнание языка.
Haghp\o{}th на языке хака означает <<помощник вождя>>.
А зовут меня Акхсар ар’Лотр, я друг твоей родильницы.
Не бойся меня.

У меня отлегло от сердца.

--- Откуда ты знаешь, что она моя родильница?

--- Я стоял рядом с твоей колыбелью, когда Митхэ оставила тебя в Тхитроне.
Я один из тех, кто отправился с ней на поиски твоего дарителя.

--- Я даже не знаю, был ли Атрис...

--- ... твоим дарителем? --- Акхсар неожиданно улыбнулся щербатой улыбкой.
У него не хватало того же зуба, что и у Митхэ.

--- Брось, парень.
Мои глаза не так остры, как раньше, но и ими я вижу, кто твой даритель.
И родильница твоя не слепая.

--- Ты хорошо знал его?

--- Хай, --- вздохнул Акхсар.
--- Старина Хат прошёл с нами полмира.
Он был смелым и хорошим другом.
Он не знал, с какого конца браться за саблю, но цитра в его руках плакала и смеялась.
Давай-ка пройдёмся немного.

Акхсар, как видно, соскучился по родной речи.
Акцент пропал уже на тридцатом шаге.

--- Сначала о деле.
Митхэ сказала, чтобы я отпустил тебя без послания.
Ты вырос в купеческом Дворе и должен знать, что это значит.

--- Если пленника отпускают живым без конкретного послания, это предложение переговоров.
Хака --- и не только хака --- считают, что предложить переговоры напрямую --- это капитуляция.

--- Молодец.
Запомнишь?

--- Разумеется!

Прошли мы с ним около кхене, но он успел рассказать много интересного...

--- Помню, был бой на перевалах у Крысиного Лаза.
Мы шли уставшие, а старину Хата сильно ранили в ногу.
Он попросил, чтобы мы посадили его на оленью упряжку и дали ему цитру.
Он заиграл... и мы почувствовали прилив сил.
Мы шли и смотрели друг на друга, не понимая, что происходит.
Я никогда раньше такого не чувствовал --- равнодушен я к музыке, знаешь ли.
А он сумел что-то задеть в моём сердце.
Жаль, что ты его не знал...

--- ... Моя последняя женщина --- тоже воин.
Как-то на пиру она пыталась рассказать мне о своих чувствах, но это у неё выходило плохо.
Хат увидел, что-то заиграл...
И его музыка ударила нам в голову, как хорошее вино.
Я даже не знал, что так может быть.
Всё, что я помню после --- в ту ночь мы с Обжоркой как с цепи сорвались.
Кстати, она вон там лежит, на обочине, десять дождей уж как.
Лихорадка съела.

--- Скажи, Акхсар-лехэ, вы с Обжоркой... твоей женщиной пошли за Митхэ по закону чести?

--- По закону дружбы, малыш.
Да, её звали Эрхэ ар’Люм, или Слишком-Много-Ест, --- Акхсар произнёс это имя, словно отдал сокровище, спрятанное до поры у него во рту.
--- Кажется, она дальняя родственница твоей кормилицы.
И пожалуйста, не называй меня <<лехэ>>, не так уж я и стар...

--- Ты и Кхотлам знаешь?

--- Д-да, конечно, --- смешался Акхсар.
--- С чего бы мне её не знать.

--- Имжу --- твой хранитель?
У хака и сели не может быть детей.

--- Да, малыш.
Но воинам и самому Имжу это знать не обязательно.

--- Почему хака напали на Тхитрон?
Мы же собирались им помочь.

Акхсар смутился ещё больше.

--- У нас и без Разлома выдался плохой год.
Зверьё из-за ушло на запад, в джунгли Живодёров.
Многие погибли, дети голодают, но хуже всего --- не хватает людей для жертвоприношений.
Старейшины говаривали, что если мы будем приносить жертвы детьми, как раньше --- многие племена просто вымрут.
Был вариант организовать общее многоплеменное становище, но ты же знаешь --- разногласия, старые обиды.
Плюс, кхм, неравенство.

--- Неравенство?

--- Когда времена хорошие, всем плевать, что один род живёт лучше других, верно?
Всем плевать, кто кому сколько золота отвесил, если у самих в котелке есть похлёбка и зерно --- в закромах.
А вот когда приходит беда...
Нужно было срочно разрядить обстановку, поэтому молодых Инхас-Лака подбили на поход.
Твоя родильница предлагала поход против идолов, но кое-кто, --- глаза Акхсара зажглись недобрым огнём, --- кое-кто надоумил старейшин, что сели --- добыча куда легче...
Мы с ней к этому отношения не имеем, --- поспешил добавить он.
--- И на брата не сердись.
Он был против, но просто не мог оставить молодых воинов на произвол судьбы.
Без него отряд даже не дошёл бы до земель сели.

--- Я должен отнестись к нему с пониманием?

--- Иногда лидер должен уступать подчинённым, даже если он с ними не согласен, --- тихо сказал Акхсар.
Ирония в моих словах стёрла едва начавшую цвести улыбку.
--- Имжу было очень тяжело.
Во время эпидемии умер его брат.
\ml{$0$}
{Он был сильным и смелым, но лихорадка расплавила его за два дня, словно снежинку.}
{He was strong and fair, but fever has melted him in two days, as if he was a snowflake.''}

--- У Имжу был брат?

--- Да, первенец прежнего вождя.
Митхэ растила их вместе, и Имжу его очень любил.

--- Кто это был? --- спросил я.
--- Кто подговорил старейшин на поход?

--- Он не станет вождём, пока я жив, --- уклончиво ответил воин и сплюнул на землю.
--- Если он попытается спровоцировать Митхэ ещё грубее, то ему придётся отвечать не перед старейшинами, а перед моим клинком.
И неважно, что будет потом.
Кстати, это Митхэ надоумила старейшин вставить в последнее сообщение очевидный обман.
Кхотлам, как обычно, великолепна, сразу поняла...

Акхсар вдруг запнулся и сменил тему.

Вскоре становились мы у малоприметного дорожного знака.
Акхсар свистнул по-птичьи, невдалеке кто-то отозвался мяуканьем.

--- Расстанемся здесь, Ликхмас-тари.
Вы с Имжу шли по совершенно непригодной тропе, чересчур близко к Гибельным руинам, и не следует испытывать судьбу ещё раз.
Иди направо, через десять кхене поверни налево и дальше через храмовую рощу.
Благородный баньян, наверное, знаешь такое дерево.
У тотема Сата-скитальца найдёшь тропинку, она ведёт на запад.
Никто пока не вышел из деревни, так что, думаю, пара кхене форы у тебя есть.
И да.
Я очень надеюсь, что ты осознал всю глубину того, насколько тебе повезло.
Больше так не делай.

Я кивнул.

--- Скажи, почему ты не вернулся к сели?

--- Видишь ли, твоя родильница --- кутрап, но у неё есть друзья и замечательный хранитель, к которому она может вернуться, --- Акхсар усмехнулся, бросив на меня добрый взгляд.
--- У меня всё ровно наоборот --- перед законами сели я чист, а вот идти... идти мне некуда.
Саблей, --- он откинул плащ левой рукой, обнажив уродливую культю, --- на жизнь я больше не заработаю, и ничем не заработаю из того, что требует две умелых руки.

Я кивнул.

--- Беги ланью, парень.
И да, кажется, ты кое-что забыл.

Акхсар протянул мне полуистлевший клочок сыромятной кожи.
Я принял его.
Старый воин прощально кивнул и зашагал сниз по дороге.
Посмотрев вслед удаляющейся тяжёлой мужской фигуре, я развернул кожу.
На стёртой, с редкими ворсинками изнанке кожи стоял свежий иероглиф, от которого на душе потеплело:

<<Я счастлива>>.

\chapter*{Интерлюдия VIII. Птицы и Мастера}
\addcontentsline{toc}{chapter}{Интерлюдия VIII. Птицы и Мастера}

\textbf{Сказка сели в пересказе Сиртху-лехэ}\FM\
\FA{
Старик рассказывал сказки с неподражаемым хуторским колоритом.
Так, как он, не рассказывал никто, я ручаюсь.
По многочисленным просьбам я попытался сделать максимально точный перевод.
}

Когда змеями рубахи подвязывали, когда на свиньях спали и сны видели, когда курами полы подметали, пришла из земель тенку\FM\ эта история, да ножек у порожка не отряхнула.
\FA{
Почти везде в землях сели и ноа повесть начинается именно так.
Тем не менее историки склоняются к мысли, что сочинена повесть была не тенку, а гораздо ранее --- первыми переселенцами с Лотоса, рассказывавшими о периоде Державы Рабов, или Первого Народовластия (600--920 по Вилет'март, 1300--2871 Эпохи Богов).
}
Правда ли, выдумка --- кто теперь разберёт?

Говорят, что когда-то Мастеров в клетках держали, как маленьких птичек или сверчков.
Только вместо чудного чириканья или стрекотания Мастера ткани делали, металлы выплавляли и камень резали.
Цепями да путами держали Мастеров.
Тех, кто был привержен дому, привязывали к старым и немощным предкам.
Тех, кто слаб был перед велением плоти, привязывали к мужчинам и женщинам.
Тех, кто продолжения своего рода желал, привязывали к многочисленному потомству.
Связывали Мастеров законами, правилами их неволили, саблями острыми угрожали.
А самых свободолюбивых сажали в настоящие клетки, навроде птичьих, и сковывали настоящими цепями, пока не смирялись они перед мощью последней клетки --- собственного тела.

Не пристало Мастерам летать, ибо чем больше руки похожи на крылья, тем меньше в них мастерства.
Так Мастерам их Хозяева говорили, свои животы потирали, а чужие уши в воде полоскали.
Что думали Хозяева тёмными ночами и о чём меж собой говорили --- неизвестно, завесь полотенцем клетку --- слеп и глух попугайчик.

Перелётные птицы жалели Мастеров, ведь были они с Мастерами одной крови.
Они прилетали да зерно клевали, которое Мастера оставляли на пороге, а Мастера слушали рассказы Птиц о дальних странах, вдохновлялись песнями-сказками, да так, что ткани краше, камень глаже да металлы звонче становились.
На свободе мало еды, говорят.
В клетке нет вдохновения, сказывают.

Мастера и Птицы понимали друг друга.
Птицы знали, что нет Мастерам дороги из клеток, и не звали их с собой.
Мастера же знали, что нельзя сажать Птиц в клетки, и потому, покормив, отпускали их.
Так жили Мастера и Птицы.

Но однажды появился Глупец.
Так его Хозяева назвали, и ни имени домашнего не помянули, ни по дому-племени\FM\ не обратились.
\FA{
По имени цатрон.
}
Он сказал, что не должно быть клеток, и тело --- не клетка, и нос --- не задвижка, и шея --- не цепь.
Глупец ходил и о том рассказывал, как сладко на свободе жить.
Кто-то из молодых Мастеров ему верил, те же, что постарше, знали, что заблуждается Глупец --- ведь они дружбу с самими Птицами водили, а Птицы и мир видали, и из морей воду бочками хлебали.
Только настаивал Глупец на своём и сказал наконец, что недостойны друг друга Мастера и их друзья крылатые.

Сказал и сказал, горлом воздух помял, потоптал --- кому какая разница?
Но тяжёлым оказалось слово Глупца.
Хозяева заявили, что это была правда, что сказал Глупец чистую правду раз в жизни, и стали её Мастерам день ото дня твердить, животы потирать, чужие уши в воде полоскать.
И Мастера поверили Хозяевам, а Птицы поверили Мастерам, и прервалась старинная дружба.

С тех пор Птицы умирать стали, и тела их на дорогах лежали, поедаемые кошками и крысами.
\ml{$0$}
{Молили-закармливали Птицы духов, чтобы дожить до весны --- и доживали, и умирали зелёными вёснами.}
{The Birds, oh, they prayed and cried and plied the Spirits with food, just to make it to the spring---and they did, and they died when the spring has come.}
\ml{$0$}
{А коли сам не ведаешь, чего просишь, какой с духов спрос?}
{If even you have no idea what you pray for, how could the Spirits help it?}
Мастера же вдохновение растеряли и ремесленниками стали, да так, что ткани грубы, и камни шершавы, и металлы петь перестали.
\ml{$0$}
{На свободе мало еды, говорят.}
{A free has no food, they say.}
\ml{$0$}
{В клетке нет вдохновения, сказывают.}
{A caged has no passion, I've heard.}

И с тех пор лишь один раз в жизни ремесленник может доказать, что достоин дружбы с Птицами.
Приходит к нему Заказчик.
Говорили люди, что Заказчика посылает Глупец, что жив он, дышит и землю топчет.
Понял Глупец, что сотворил, и хочет всё исправить, молоко разлитое тряпкой собрать да кашу из того молока сварить.
Заказ же всегда прост и труден, и порой не отличишь его от обычной работы.
Те же из ремесленников, что дело своё любят да снова стать Мастерами хотят, всю жизнь ткани делают, металлы выплавляют да камень режут, словно в последний раз, вглядываясь, пропустить боясь --- не пришёл ли тот самый Заказ, не вспомнили ли о давнем друге Птицы, не выбрал ли его Глупец среди прочих?

\ml{$0$}
{Да так и ждут-пождут, да без меня, я стороною проходил, ушами далёкие речи ловил.}
{So that way they wait and wait, without my self, I was there and I was here, caught a rumor by the ear.}
\ml{$0$}
{Что наловил, то и в котелке сварил, что в котелке сварил, тем и дитя кормлю.}
{What I caught, I put in a pot, what I stew, was eaten by you, child.}

\chapter{Хрупкий мир}

\section{Переговоры с хака}

Наверное, меня никогда и нигде не встречали так.

Ликхэ, Ситрис и Кхохо кинулись мне на шею и долго не отпускали.
Чханэ просто подошла, спотыкаясь, нервно рассмеялась, потрепала меня дрожащими пальцами за рукав и села под ближайшим деревом, в уголке деревенской площади.
Я сел рядом с ней.

\ml{$0$}
{--- Живой?}
{``Alive?''}

\ml{$0$}
{--- Живой.}
{``Alive.''}

\ml{$0$}
{--- Точно?}
{``Positive?''}

\ml{$0$}
{--- Точно.}
{``Positive.''}

--- Как чудесен мир, --- лаконично ответила подруга.
Вся её разговорчивость куда-то пропала.

День выдался погожий.
Пели птицы, жужжали стрекозы, шуршали в кустах ящерицы, светило яркое солнце.
Я лежал у Чханэ в объятиях, чувствуя себя маленьким-маленьким котёнком.
Она что-то шептала, запустив пальцы мне в волосы.
Её большое, оранжевое, как солнышко, лицо было прямо надо мной;
янтарные глаза смотрели куда-то далеко, и в них было что-то...

--- ...Когда-нибудь ты встретишь человека, который тебя будет любить, --- говорила кормилица.
--- Как я люблю Хитрама, как Хитрам любит меня.
Этот человек может стать твоей семьёй.

--- А как я узнаю, что человек меня любит? --- спросил я.

--- Ты поймёшь, --- коротко ответила Кхотлам.

Кажется, я понял.

Я прижался к Чханэ крепче, и она ответила на мои объятия.
Я почувствовал на лбу её слёзы;
они текли и терялись в моих волосах.
Мы так и просидели под деревом в обнимку несколько часов, ловя на себе частью удивлённые, частью неприязненные взгляды хака;
едва ли прибывшая делегация ожидала увидеть такую романтику на переговорах.

Кхотлам и Митхэ обменялись крепким рукопожатием.
Большего не позволяла ситуация.

--- Рада, что ты живая, Митхэ ар'Кахр, --- сказала на цатроне кормилица.

--- Я благодарю тебя за всё, что ты сделала, Кхотлам ар'Люм, --- ответила Митхэ на том же языке.

--- Жрецы выбрали \emph{совсем нейтральную территорию}, --- шёпотом сообщила мне Чханэ, обведя рукой полузаброшенный хутор.
--- Всё-таки Митхэ --- кутрап, и мы обязаны её казнить, если она свернёт с трёх главных дорог.
Жрецы решили не рисковать.
Мало ли что придёт в голову этим хака...

Я кивнул.
Жрецы всё сделали правильно.

Постепенно к нашему дереву подошли все воины и сели рядом.
Они молчали, и я чувствовал в их молчании какую-то недосказанность.

--- Всё хорошо? --- шёпотом спросил я наконец.
Мне не ответили.

Я вдруг остро ощутил, что чего-то --- вернее, кого-то --- не хватает.

--- Где Хитрам?

Ликхэ обернулась, и я увидел, что она плачет.

--- Стрелка вклинилась между пластинами доспеха, --- прошептала девушка.
--- Она рукой дёрнула и напоролась.
Надо же, какая случайность...

--- Вы похоронили её?

--- Нет, --- прошептал Эрликх.
По его лицу тоже текли слёзы.
--- Это произошло в старом лесу, земля там жилистая.
Оставили её под большим папоротником-опахало.
Достойное место отдыха для достойного воина.

--- Мы даже хэситра ей налили на донышке, --- подхватил Ситрис.
--- Воду для детей оставили, на обратный путь.
Надеюсь, она не сердится.

--- Она бы поддержала это, --- уверенно сказала Чханэ.
--- Она добрая женщина.

Митхэ не участвовала в переговорах.
Она выглядела так, словно её привязали голой к позорному столбу, которые использовали для наказания преступников тенку.
Все без исключения тхитронцы восхищённо поедали Митхэ глазами --- кто-то украдкой, как Кхохо, кто-то спокойно, как Ситрис, а наблюдатели из числа ткачей и вовсе пялились, открыв рот.
Я тоже изредка бросал на неё взгляды.
Мне часто говорили, что я похож на родильницу, и я думал, что речь о Кхотлам.
Но нет.
Теперь, когда они стояли почти рядом, я увидел разницу.
Я был копией Митхэ.
Я словно смотрел в зеркало Хри\FM\ --- те же длинные брови, тот же точёный нос, тот же крупный, чувственный, даже немного жадный рот.
\FA{
Согласно преданию, у Хри-соблазнителя было зеркало, в котором мужчины превращались в женщин и наоборот.
}
Тот же маленький рост, над которым не упускал случая подшутить Конфетка.
Различались только волосы --- жёсткие чёрные кудри Митхэ совершенно не были похожи на мои мягкие светлые локоны.
Я схватил один локон и критически осмотрел;
да, наверное, волосы я унаследовал от Атриса.

И только кормилица, чисто выговаривая слова цатрона, украдкой пыталась поймать взгляд этого странного однорукого старика Акхсара, который скромно стоял в углу.
На глазах Акхсара лежала непроницаемая тень, и я не видел, на кого смотрит он.
Он даже не взглянул на Ситриса, который подошёл к нему перед переговорами.

\ml{$0$}
{--- Где Эрхэ? --- тихо спросил Ситрис.}
{``Where's \Oerchoe?'' \Sitris\ quietly asked.}

\ml{$0$}
{--- Умерла, --- одними губами ответил Акхсар.}
{``Dead,'' \Akchsar\ silently mouthed.}

\ml{$0$}
{--- Сочувствую, --- пробормотал Ситрис и отошёл.}
{``I'm sorry,'' \Sitris\ mumbled, then walked away.}
\ml{$0$}
{Всё оставшееся время воины делали вид, что не знают о существовании друг друга.}
{Each of them acted like the other doesn't exist---all the time since.}

Кхарас стоял чуть поодаль и смотрел в пустоту.
Я вдруг понял, что он тяжелее всех переживал смерть Хитрам.
Сколько она была в его Храме?
Кто ещё будет нежно называть его, грубого седого силача, <<мальчиком>>?
Я подошёл к вождю и положил руку ему на плечо.
Кхарас кивнул без улыбки.

\ml{$0$}
{--- Ты как?}
{``How are you?''}

--- Мне нужен путник, плававший по дорогам рыб\FM, --- сухо сказал вождь.
\FA{
Кхарас хотел получить данные топографии поселений Инхас-Лака, которые мне удалось добыть.
}

--- Всё будет, --- кивнул я.

Уже под конец переговоров Кхарас сказал: <<Мы ведь могли взять её с собой, за городом папоротники пышнее>>.
Произнёс он это тихо, всеобщее внимание было поглощено беседой Кхотлам и двух старейшин.
Разумеется, никто не обратил внимания на его бормотание.

\asterism

Наконец переговоры подошли к концу.

--- Я думаю, нам стоит устроить пиршество, --- улыбнулся один из старейшин хака.

--- Я думаю, нам стоит уважать традиции друг друга, --- веско сказал Акхсар из угла.
--- Традиции хака --- пиршество.
Традиции сели --- Молчание.
Пусть каждый из народов последует своим традициям.

Хака нахмурились;
кое-кто зароптал.

--- Кхотлам ар'Люм, а что думаешь ты? --- обратился старейшина к кормилице.
--- Хагпот Инхас-Лака мудр, но это решение...

---  ... это решение я поддерживаю целиком и полностью, --- закончила Кхотлам.
--- Надеюсь, старейшины смогут объяснить молодым его целесообразность.

\ml{$0$}
{--- Это решение обижает нас, Кхотлам ар'Люм.}
{``This decision offends us, \Kchotlam\ ar'\Loem.}
\ml{$0$}
{Мы...}
{We are---''}

--- Похоже, между нами некоторое недопонимание, --- Кхотлам сузила глаза и понизила голос --- ровно настолько, чтобы фразу услышали все, но тем не менее сочли приватной.
--- Старейшины хака решили предать алтарю наших сыновей и дочерей.
Мы не разделим со старейшинами ни тарелку бобов, ни кукурузную лепёшку, пока не сочтём их кровь чистой.

Хака вздохнули и забормотали;
глаза старейшины налились кровью.

--- Странников и торговцев из числа вашего народа мы примем с теплотой, как и ранее, --- чуть громче продолжила Кхотлам, обращаясь к конвою и сопровождающим.
--- А теперь, я думаю, самое время уважительно попрощаться и разойтись по домам.

<<Умно, --- подумал я.
--- Унижение на глазах народа станет старейшинам уроком.
В следующий раз хака будут ценить не только дружбу, но и элементарную дипломатическую вежливость>>.

\ml{$0$}
{То, о чём думала Митхэ, так и осталось тайной.}
{\Mitchoe's thoughts remained a secret.}

\section{Нарушенное Молчание}

\epigraph
{Прославлять себя победой --- это значит радоваться убийству людей.
Победу следует отмечать похоронной процессией.}
{<<Книга Пути и Достоинства>>, Мудрый Старец.
Культура Цина, Древняя Земля}

<<Неужели я вернулась домой?>> --- думала Митхэ, с некоторым удивлением рассматривая растущие башни Тхартхаахитра.
У неё было смутное ощущение, что что-то происходит в последний раз, и по воинской привычке списала это на предвестье скорой смерти.

\ml{$0$}
{--- Удивляешься, что вернулась? --- тихо усмехнулся ей на ухо Атрис.}
{``Are you surprised you're back?'' \Aatris\ whispered in her ear.}
\ml{$0$}
{--- Я тоже.}
{``Me too.}
\ml{$0$}
{Всегда.}
{Always.''}

--- Так, --- Митхэ вырвалась из плена собственных размышлений, --- ты почему пешком идёшь?
\ml{$0$}
{Ну-ка быстро в повозку.}
{Get in the cart, now.''}

\ml{$0$}
{--- Я прекрасно себя чувствую, --- заверил её менестрель.}
{``I'm fine, really,'' the minstrel assured her.}
\ml{$0$}
{--- Нога больше не болит.}
{``My leg doesn't hurt anymore.''}

\ml{$0$}
{--- Как знаешь, --- меньше всего на свете Митхэ хотелось спорить с Атрисом.}
{``Suit yourself.'' The last thing in the world \Mitchoe\ wanted to do is to argue with \Aatris.}

Все знали --- по прибытии домой им впервые придётся нарушить древнюю традицию Молчания после войны.
Вернувшихся воинов всегда хорошо кормили и на целую декаду, а порой и на две освобождали от обычной работы, чтобы они могли хотя бы выспаться вволю.
Также в течение декады отменялись все праздники, даже Слёзы Ситхэ;
сели знали, что если была война --- значит, кто-то лишился самого дорогого.
Однако вскоре в город должны были прибыть вожди хака, и в знак дружбы следовало устроить пир --- заставлять гостей ждать было бы невежливо.
Молчание должно было быть нарушено, и одна эта мысль наполняла Митхэ дурными предчувствиями.
Она видела пышные триумфы военачальников тенку, которыми те праздновали свои крошечные жалкие победы на поле боя.

Среди старых воинов бродил тот же безрадостный толк.
Аурвелий, несмотря на ослепительно белые волосы, казался чернее ночи.

--- Если кваждый победа праздновать, война никвогда не кончиться! --- буркнул он как будто между прочим, проходя мимо Митхэ.

Ноа славились воинственностью;
послевоенные триумфы если и не были у них традицией, то считались по крайней мере оправданной вольностью народа.
Услышать такие слова от старого пирата-ноа Митхэ ожидала меньше всего.

--- Мы празднуем не победу, а заключение мира, сенвиор Амвросий, --- веско заметил Атрис.
--- Это плохая отговорка, но давление обстоятельств в отговорках и не нуждается.
Я думаю, тебе стоит улыбнуться --- теперь ты вступишь в полные права на землях сели.

Старик хмыкнул одновременно с кем-то, лежащим в повозке.

--- Ситрис, кинь что-нибудь пожевать, --- попросила Согхо.

Из повозки вылетел круглый кахраханский хлеб и скрылся где-то в рядах наёмников.
За хлебом последовал помидор, весело блеснув на солнце полированным алым боком.

--- Вот лентяй, --- буркнул Акхсар.
--- Ситрис, ты уже весь день лежишь и жуёшь.
Не надоело?

--- Это надоесть не может, --- сообщила повозка.
--- А вот отряды из южных Храмов меня уже нервируют.

--- Жуй и не бойся, --- успокоил разбойника Атрис.
--- С нами тебя никто не тронет.

На этот раз хмыкнула половина отряда.

--- Атрис, дружище, я бы пошёл с тобой на край света, --- проворчал Ситрис, глубже зарываясь в тюки.
--- А с этими сверчками безногими даже нужду справлять небезопасно.

--- Я бы на твоём месте нужду справляла где-нибудь снаружи повозки, --- хмуро сказала Эрхэ.

Все, включая Ситриса, захохотали.
Эрхэ ловко увернулась от брошенного помидора и тут же начала шарить в своём мешке, подыскивая подходящий снаряд;
скорость, с которой Ситрис построил из тюков крепость, внушала уважение.

--- Бум-буррубум-бам-бум\FM, --- прогудел разбойник, подражая осадному барабану.
\FA{
<<Поднять щиты, закрыть бойницы>> (барабанный язык сели).
}

--- Придурок, --- пробормотала Эрхэ.
--- Твоё счастье, что промазал.
Чего ржёте? --- раздражённо добавила она, глядя на хохочущих наёмников.

--- Лезть повозка тоже, Атриций, --- посоветовал Аурвелий менестрелю.
--- Ты больно, я ввидеть.

--- Не волнуйся за меня, кормилец.
Мы почти дошли.

\section{Камень и огонь}

Родные ворота.
Я стоял и гладил их, с каким-то неведомым ранее чувством ощущая под пальцами тепло дерева, трещины и резьбу.

--- Ликхмас, ты чего? --- удивился Ситрис.

Я не ответил.

Кормилица, проходя мимо, молча погладила меня по голове.
Чханэ приложила ладонь к воротам рядом с моей и подержала несколько секхар, затем пошла дальше.
На дереве остался едва заметный влажный отпечаток;
ветер унёс его мгновение спустя.

--- Лисичка, ты проголодался, наверное, --- прозрачно намекнула Ликхэ, кивнув в сторону храма.

Я снова промолчал.

Город не изменился, но я его не узнавал.
Люди приветствовали меня, как обычно;
мало кто знал о том, что я успел пропасть и вернуться обратно живым.
Старушка на улице Дышащего Дерева кивнула мне, словно старому знакомому;
трубка в её зубах дымилась, как и раньше.

--- Ну что там, Ликхмас? --- подбежала ко мне молодая крестьянка.
--- С хака заключили мир?
Воины какие-то молчуны нынче.

Услышав разговор, нас окружили ещё несколько человек с вопросами.

--- Хай, я это...
В целом всё хорошо, подробности сообщит Двор Люм в ближайшее время, --- сказал я, высвобождаясь из цепких рук.
\ml{$0$}
{--- Сегодня можете спать спокойно.}
{``People can sleep easy tonight.''}

Вот и храм.
Центр, убежище и дом.

Я вдруг понял, что на площади царила странная, неприятная суета.
Торговый день уже давно закончился, однако вокруг тихо сновали люди, таща какие-то вещи.
Они сновали вокруг, нервно глядя на меня.
Их глаза вопрошали: <<Что ты здесь делаешь?>>

<<Что ты здесь делаешь, синяя роба?>>

<<Разве тебе не следует быть в другом месте?>>

Я перешёл на бег.
В голове пронзительно звенел колокольчик тревоги.
Первая ступенька, вторая, третья, тринадцатая... вот и вход.

Меня встретил пустой чёрный зёв храмовых ворот.
Внутри не горел ни один факел.

Трукхвал сидел у входа и рыдал.
Я подбежал к нему.

--- Учитель! Что...

Тут мой взгляд выхватил из темноты странный туман.
Нет, это был не туман.
Воздушная пыль пахла камнем и огнём.

--- Ликхмас, --- из темноты выпрыгнул Ситрис.
--- Оставь Трукхвала, я позвал людей, о нём позаботятся.
Быстро на дежурство наверх.

--- Диверсия, Ликхмас, --- раздался голос Кхохо.
--- Подробности потом.
Инструменты уже готовы, Эрликх тебе будет ассистировать.
Иди по западной лестнице, восточную завалило.

Вскоре я разглядел лежащие на полу изувеченные тела, наспех прикрытые одеялами.
Одно, два, три... тринадцать.
Вот рука Кхатрима со знакомой татуировкой.
В предсмертной муке она сжала амулет Сана, да так и окоченела.

Я медленно пошёл наверх, не веря в происходящее.

--- Лис, беги, я сказал! --- рявкнул Ситрис.
--- Горевать будешь утром, жрец!

--- Я сейчас... --- услышал я сдавленный всхлип Трукхвала.
--- Пусти меня, Уголёк... я тоже наверх...

--- Сиди, Звоночек, --- Кхохо явно пришлось приложить силу, чтобы усадить учителя обратно.
--- Парень справится.
Ты хорошо его воспитал.

\section{Стили боя}

--- Боевое искусство должно воспитывать устойчивость к боли и страху, потому что боль и страх являются неотъемлемой частью боя.
В этом принципиальное отличие боевого искусства от любого другого.
Это последнее, на что я хочу обратить ваше внимание.

Тси слушали Баночку, открыв рты.
Культура взаимного уничтожения была для нас новой и непривычной.
Кое-кто --- например, Листик и Мак --- напрочь отказались её изучать.
<<Это неправильно>>, --- отрезали они хором и ушли по своим делам.

--- Итак, --- Баночка поднял руки.
--- Тогда я прошу всех встать в пары.
Старайтесь подбирать пары по росту и комплекции;
вы поймёте, почему это важно.
Первое занятие --- болевая подготовка.

--- То есть ты хочешь, чтобы мы... били друг друга? --- ошеломлённо спросил Фонтанчик.

--- Именно.

--- Баночка, дружище ты мой, ну я же не это имел в виду, когда говорил про технику боя...

--- Я проанализировал традиции племён джунглей, подключил кое-какую древнюю историю и пришёл к выводу, что это неотъемлемый этап подготовки, --- заявил Баночка.
--- И да, встань в пару с сестрой, пожалуйста, иначе ты кого-нибудь прихлопнешь.

--- Пусть встанет с Зубиком, --- буркнула Пирожок.
--- Мне это не нравится, и я этим заниматься не собираюсь.

Впрочем, Зубик тоже не пожелал продолжать обучение.
Вместе с Пирожок с тренировочной площадки ушла ещё четверть тси.

Однако Баночку это не обескуражило.
Его глаза горели;
того, кто увлёкся делом до такой степени, сложно в чём-либо убедить.

--- Неужели вы не видите в этом красоты? --- спрашивал он у нас с Заяц.
--- Пожалуйста, забудьте хоть на секунду, что это методики взаимного уничтожения, и посмотрите на внешнюю форму!

--- Мне об этом сложно забыть, --- призналась Заяц.

--- Разве хищные звери не прекрасны, хоть их тела --- орудие уничтожения? --- допытывался Баночка.

\ml{$0-[qi]$}
{--- Это другое, --- отрезала Заяц.}
{``It's different,'' Hare objected.}
\ml{$0-[qi]$}
{--- Не мешай хищничество с этим непотребством.}
{``Don't mix predation with an abomination like that.}
\ml{$0-[qi]$}
{То, чем ты увлёкся --- самая тёмная сторона внутривидовой конкуренции.}
{Your new passion is the darkest side of intraspecific competition.''}

\ml{$0-[qi]$}
{--- Ты считаешь внутривидовую конкуренцию неестественной?}
{``You think there's something unnatural about intraspecific competition?''}

\ml{$0-[qi]$}
{--- Я считаю, что в ней нет нужды, если большинство представителей вида способны планировать хотя бы на два шага вперёд!}
{``I think there's no need for this if members of the species can plan at least two steps ahead!''}

\ml{$0-[qi]$}
{--- Подумай ещё, --- кротко попросил Баночка.}
{``Think again, please,'' Flask meekly asked.}
\ml{$0-[qi]$}
{--- Попробуй увидеть то, что для меня очевидно.}
{``Try to see the thing that's obvious to me.''}

\ml{$0-[qi]$}
{--- Объясни, почему для тебя так важно, чтобы твоё занятие нравилось и мне, и Заяц, и всем остальным тси? --- спросил я.}
{``Explain: why is it so important to you that I, and Hare, and all other Qi like your hobby?'' I asked.}

\ml{$0-[qi]$}
{--- Ну, вы же мои др...}
{``Come on, you're my fr---''}

Баночка вдруг смешался и задумался на целую минуту.

\ml{$0-[qi]$}
{--- Тоже верно, --- признал он наконец, по очереди обнял нас с Заяц и убежал.}
{``You're right,'' he conceded, then hugged us in turns and left fast.}

Больше мы не услышали от Баночки ни одной лекции.
Но новости о его деле разлетались быстро;
кто-то даже запустил несколько парящих видеокамер и наладил на форуме видеотрансляцию с тренировочной площадки.
Тси разделились;
большинство открыто осуждало происходящее, но многие, выкроив час времени, записались на тренировки.

Баночка разработал сразу несколько парадигм (или, как он выразился, <<стилей>>) боя.
Отныне каждый день два часа он уделял тренировкам прочих, сам же тренировался всё свободное время, выбирая всё более сложную местность, погодные условия и расположение противников.
Он поднял на одном из серверов <<кукловода>> и обучил несколько кукольных нейросетей различным стилям боя.
Все признавали, что любая из этих кукол превосходила по мастерству воинов царрокх --- и тем не менее Баночка их побеждал.
Прибывшая с материка Кошка открыла рот, когда увидела, как плант за какие-то несколько секунд расправился с двадцатью виртуальными врагами.
Два коротких трёхчленных копья, названных им <<фаланги>>, перерубали длинные копья кукол и легко вонзались в жизненно важные органы --- ровно настолько, чтобы враг уже не встал.

--- Ужасно, --- пробормотала она.
--- Баночка со своими палками может подчинить всех царрокх в одиночку.

Однако вне тренировок плант продолжал носить резонанс-саблю.
Сабля лежала рядом, когда он ел, что-то чинил, сабля висела у него на спине, когда он отправлялся на берег моря или в горы --- узкоспециализированный тяжёлый инструмент, превратившийся в фетиш.
Я пытался поговорить с ним на эту тему, но плант демонстративно поворачивался и уходил, едва речь касалась его оружия.

\section{Последний знак}

Вечером Атрису стало хуже.
Храмовый жрец осмотрел рану и вызвал носилки.

Цепочку дурных знамений завершила сабля Митхэ, Перидотовая Лиана --- в тот же вечер она треснула почти вдоль безо всякой видимой причины, не дожив, если верить пробе, всего один дождь до своего трёхсотлетия.
Митхэ, стараясь не поднимать шума, отнесла осколки одному из оружейников Тхартхаахитра.
Несмотря на молодость --- мужчине едва перевалило за сорок дождей --- он уже был старшиной квартала и считался одним из лучших мастеров по кукхватру на Юге.

--- Я не знаю, --- сокрушалась Митхэ.
--- Я взяла эту саблю боем у подлого человека.
Я пыталась очистить её ритуалами, дать ей другую славу.
Неужели я что-то сделала не так и мёртвые потребовали отнятое назад?

--- Это бывает, --- успокоил её мужчина, осмотрев клинок через микханское стекло.
--- Кое-кто считает, что оружие трескается из-за гнева убитых им, но причина не в этом.
У нас есть поверье, что кукхватр сродни человеческой плоти --- он помнит форму, заложенную создателем.
Как ты знаешь, весь кукхватр идёт от Древних.
И когда мы куём его, мы на самом деле всего лишь временно придаём ему форму, которая нам нужна.
Но каждая деталь и каждый инструмент, из которого была создана сабля, стремится вернуться к изначальной форме.

--- Я не знала об этом, --- призналась Митхэ.

--- Это всего лишь поверье кузнецов, оно бесполезно для воинов.
Но, как ты видишь, на правду похоже.
Вот здесь на изломе твоей сабли можно разглядеть зубчики и волны --- возможно, это был редуктор или что-то подобное.
Я бы не решился использовать столь прихотливую деталь для средней части клинка, но, видимо, мастер получил свежий слиток и об этом не знал.
Таким образом, многое зависит от формы сабли и условий её ковки.
Если выбрать неправильную форму, неправильно сочесть составляющие, в кукхватре начинает накапливаться напряжение.
Вот, смотри.

Оружейник снял саблю с подставки на стене и принёс Митхэ.
Затем ударил по сабле крохотным молоточком --- и она вдруг застонала, словно воткнувшаяся в дерево стрела: <<Уиуиуиуи...>>

--- Это, конечно, совсем запущенный случай, чаще всего услышать могу только я да пара-тройка мастеров с Запада.
Копится напряжение и десять, и пятьдесят дождей, пока не разрешится трещиной.
Его можно снять, если вовремя заметить.
Но если трескается --- только перековка.
Поэтому многие предпочитают ковать короткие клинки, ножи и фаланги, --- с ними такого не бывает.
По крайней мере, если нож и треснет, то не в эту тысячу дождей.

--- Сделаешь в точности такую же?

--- Дело сложное, но возможное.
Только я бы предложил кое-что получше, --- оружейник легонько ударил молоточком по осколку, и он отозвался протяжным чистым звоном.
--- Я немного изменю форму и перекантую составляющие, чтобы напряжение накапливалось медленнее.
Вполне возможно, что на твоём веку <<разгрузка>> клинка и не понадобится, поэтому я вставлю в эфес алюминиевую табличку с указаниями для следующих владельцев.
Алюминий мягкий, и мастер сможет набить дату последней <<разгрузки>>, не причиняя клинку вреда.
Очень удобно, мастера Запада так делают уже давно.

\ml{$0$}
{--- Я согласна.}
{``Agreed.}
\ml{$0$}
{Сколько?}
{How much?''}

\ml{$0$}
{Мужчина назвал цену.}
{The man gave her a price.}
\ml{$0$}
{Митхэ вытаращила глаза.}
{\Mitchoe\ gave him a goggle-eyed look.}

\ml{$0$}
{--- Ты серьёзно?}
{``Are you serious?}
\ml{$0$}
{У меня нет столько!}
{I don't have that much!''}

\ml{$0$}
{--- Нет? --- удивился оружейник.}
{``Do you?'' the smith looked surprised.}

\ml{$0$}
{--- Нет и не было никогда!}
{``Now or ever before!''}

\ml{$0$}
{--- Сколько есть?}
{``How much do you have?''}

Митхэ вывалила из дорожного мешка все сбережения.

\ml{$0$}
{--- Здесь едва хватает на расходники, --- извиняющимся тоном сказал оружейник.}
{``This is barely enough for consumable supplies,'' the smith said almost apologetically.}
\ml{$0$}
{--- Не пойми неправильно.}
{``Don't get me wrong.}
Ковка кукхватра --- недешёвое удовольствие, даже для меня, хоть мне по старшинству в квартале многое даром достаётся.

Митхэ молчала и смотрела на осколки.
Мужчина вздохнул.

\ml{$0$}
{--- Я знаю, кто ты и откуда вернулась, --- глухим тоном сказал он.}
{``I know who you are and where you returned from,'' he said muffled.}
--- У ремесленников есть поверье о заказе Глупца.
В жизни бывает один заказ...

--- Я слышала эту сказочку.

--- Глупость, да, --- усмехнулся мужчина.
--- Однако ж...

Оружейник вздохнул ещё раз и пошевелил рукой осколки.

\ml{$0$}
{--- Дай мне хоть на еду, --- попросил он.}
{``Feed me, at least,'' he asked the warrior.}
\ml{$0$}
{--- Не могу же я целую декаду ковать себе в убыток.}
{``I can't forge at my own expense for a whole decade.''}

Митхэ вытряхнула из карманов несколько щепоток золота, вытащила из-под секхвим толстое перо шипастого страуса --- тайник на самый крайний случай, --- опорожнила его и положила рядом;
наконец, сняла с руки тяжёлое золотое кольцо Кхара-защитника --- подарок Короля-жреца.
Оружейник принял кольцо, осмотрел его и грустно чему-то улыбнулся.

--- Ты чего, Золото? --- спросила Эрхэ за ужином на постоялом дворе.
--- Хватит жевать вяленое мясо, мы не в походе!

--- Я не хочу есть.

Эрхэ вдруг подняла руку подруги и задумчиво потёрла средний палец, на котором ещё синел след от кольца.
Митхэ молча выдернула руку.
Эрхэ вздохнула и тихо подозвала трактирщика:

--- То же самое ещё по тарелке, запиши на меня.
\ml{$0$}
{И вина побольше.}
{And more wine.''}

\section{Жертвоприношение}

Ночь для жреца --- жаркое время, во всех смыслах этого слова.
В кельях, надо сказать, душновато, несмотря на все ухищрения зодчих.
Кое-кто из жрецов в такой сезон перебирался вниз, к воинам.
Кое-кто стелил ложе на крыше или в крипте;
мстительные тени, по их словам, мучают только первые две-три ночи.
Если же не келья, то ночной вызов --- закровила роженица, залихорадило ребёнка, принесли из сельвы укушенного охотника;
если не ночной вызов, то жертва.

\ml{$0$}
{В тот день я был жрецом.}
{I was a priest that day.}
\ml{$0$}
{Вторым жрецом в Тхитроне, чьё сердце ещё билось, и единственным, чьи руки могли держать нож прямо.}
{I was one of two \Tchitron\ priests having a heartbeat and the only who could hold a knife.}
\ml{$0$}
{<<Я принесу жертву богам и в скорби, и в болезни, и при смерти>> --- так говорится в клятве.}
{``Even grieving, sick, or dying, I shall sacrifice to the gods,'' the oath says.}
\ml{$0$}
{<<Сегодня можете спать спокойно>> --- так сказал я людям.}
{``People can sleep easy tonight,'' I said in the daylight.}
\ml{$0$}
{Кто же знал, что этот долг и это спокойствие свалятся на мои плечи.}
{Who would know that duty and that easiness will be on my shoulders alone.}

Одно жертвоприношение мало отличается от другого.
Сначала ты сидишь и ждёшь знака кирпича, стараясь не заснуть;
впрочем, звук взорвавшегося терракотового бруска, изготовленного специальным способом, способен пробудить даже мёртвого.
Затем ты так же терпеливо, медленно, бесстрастно лишаешь жизни другое живое существо, стремясь максимально продлить его страдания.
Чем дольше и мучительнее умирает жертва, тем больше тихих дней впереди.

Одно жертвоприношение мало отличается от другого.
Но есть то, которое отпечатывается в памяти навсегда.
Первое.

Я сидел в камере и смотрел на пленника.
Средних лет мужчина-идол, низкорослый, широкоплечий.
Его кисти пестрели короткими впалыми шрамами, выдавая ловца змей.
Свалявшийся, покрытый засохшей кровью пучок разномастных волос --- волос врагов-людей --- бросал непроницаемую тень на избитое, опухшее лицо.
Идолы Живодёра срезали с головы врагов кожу и вживляли себе;
постепенно лысые зелёные головы и спины прославленных воинов покрывались роскошной пёстрой шевелюрой.

Я хорошо помню, как тускло мерцали из-под татуированных век лишённые белков глаза.
Сели редко татуировали лицо;
у моего народа бытовало поверье: вечная краска на лице --- вечный бой.
Идолы же считали, что духи костей через священные знаки проникнут в глаза и защитят своих потомков от болезней, врагов и смерти.

Этому, увы, не помогли ни татуировки, ни боевое мастерство.

Пленник молчал.
Я уже не был ребёнком, который без лишних вопросов, от всей души ненавидит идолов.
Мне хотелось сказать ему что-то напоследок, хоть как-то облегчить его участь.
Но что сказать?
Любые слова излишни, любая фраза из моих уст --- уст врага и мучителя --- будет лишь грязной насмешкой.

--- Ликхмас, пора.
Выводи его, --- окликнул меня из темноты Митликх.

Я видел ранее, как грубо выволакивали жертв жрецы.
Заслуживал ли идол такого отношения?
Митликх отомкнул замок, и я просто махнул идолу, призывая его следовать за мной.
К моему удивлению, пленник послушно встал и пошёл, перед этим испустив совершенно человеческий, полный грусти вздох.

--- Держи его, а то сбежит, --- предупредил Митликх.
--- Опытный воин.
Эти твари ловкие до ужаса.

--- Некуда ему бежать, --- возразил я.
Идол, словно соглашаясь со мной, молча пошёл к выходу из крипты.
Митликх, подозрительно глядя на пленника, покачал головой и аккуратно закрыл дверь.

Жрецы уже ждали нас во всеоружии.
Их фигуры казались толстыми и неуклюжими из-за мешковатых <<кровавых плащей>>, надетых поверх обычной синей робы.
Лица скрывали <<утиные маски>>, холодные глаза замерли, выжидающе глядя на пленника.

Идол замер перед алтарём.
Он в ужасе смотрел на бронзовые скобы, на чёрный, пропитанный кровью базальт, на осколки кирпича, ещё лежавшие вокруг.

Я вдруг понял, что именно сейчас нужно приободрить.

--- Воин, --- сказал я ему на наречии Живодёра, --- я вижу по твоим шрамам, что ты всю жизнь храбро сражался.
Это твоё последнее сражение.

Идол бросил на меня пронизывающий бешеный взгляд.
В какой-то момент я подумал, что сейчас он кинется и попытается хотя бы зубами...
Двое жрецов нетерпеливо схватили его за плечи:

--- Так, хватит церемоний...

--- Тихо, --- басом громыхнул Первый, и вплетённые в его волосы перья качнулись.
--- Парень всё делает правильно.
Отпустите.

Жрецы нехотя отошли от пленника.
Идол снова вздохнул и неуклюже, боком лёг на алтарь.
Первый жрец быстро, со знанием перерезал его путы и прикрутил конечности к бронзовым кольцам.
Его ассистент вставил в рот жертвы расширитель.
Идол зажмурился.

Я не стану описывать жертвоприношение.
В голове помутилось на первой минуте жуткого, полного отчаяния тонкого крика.
Искажённое болью татуированное лицо отпечаталось в моей памяти, словно его вырезали раскалённым алмазным резцом...

... Пришёл в себя я уже внизу.
Чадил и трещал почти прогоревший факел, жрецы шли в полутьме, на ходу снимая с себя <<кровавые плащи>> и обтирая ими руки.
Они говорили о каких-то обыденных вещах, кто-то даже смеялся.
А мне казалось, что я после увиденного не смогу улыбаться никогда.
Всё казалось каким-то нереальным, мир раскололся на две части --- крышу храма и всё, что вне её.
Как можно ходить по улицам, веселиться с друзьями, когда каждые несколько дней в мучениях умирает живое существо?
Как можно есть, пить и сладко спать, зная о происходящем на крыше храма?
Как можно жить?..

Никто не замечал, что со мной творилось что-то неладное.
Лишь Кхатрим подошёл и, обняв, погладил меня по голове.

--- Ликхмас, --- сказал тогда врач.
--- Я знал, что тебе пока нельзя на крышу.
Госпиталь --- одно, это --- совершенно другое.
Но ты жрец.
Руки жреца обязаны знать ритуал.
Идём, угощу медовухой, полегчает.
Потом сразу домой.

Медовуха не помогла --- желудок скрутило, тело охватила лихорадка.
Я попрощался с жрецами, но до дома так и не дошёл --- ватные ноги подкосились на нижнем этаже.
Тогда я ещё не знал и не мог знать о реакции Стлока, я просто думал, что заболел.
Рвало до крови, пустой желудок сводило так, что невозможно было даже вдохнуть.
Ликхэ и Ситрис увели меня в спальню и ухаживали за мной, как за ребёнком.

--- Так, быстро, раздевай его.
Спирт я принёс, растирай, --- шёпотом командовал пожилой воин.

--- Донышко, может, жрецов позвать? --- испуганно спросила Ликхэ.

--- С ума сошла?
Он реагирует не так, как обычные адепты.
Многих рвёт, да, но лихорадка?
Его могут не допустить... --- Ситрис нервно оглянулся и приник к уху Ликхэ.
Та ахнула.
--- Так что никаких жрецов.
Я кое-что позаимствовал у Кхатрима, сейчас дадим ему... --- он снова зашептал что-то едва слышно.

--- Ликхэ, Ситрис, может, хватит уже ерундой заниматься? --- проворчал заспанным голосом Кхарас.
--- Спать только мешаете.
Ликхмас, возьми себя в руки, всех рвало в первый раз.
Был бы это ещё ребёнок, куда ни шло, но идол...

\ml{$0$}
{--- Заткнись и спи, --- рявкнула на него Ликхэ.}
{``Shut your mouth and sleep,'' \Likchoe\ snapped at him.}
\ml{$0$}
{Это был первый и последний раз, когда она позволила себе такое обращение с боевым вождём.}
{That was the first and last time she ever treated the warchief like that.}
\ml{$0$}
{Слегка ошарашенный Кхарас замолчал, но я слышал, как он нервно возится на своей лежанке.}
{Slightly stunned \Kcharas\ fell silent, but I heard him fidgeting in his bed.}

Как я узнал впоследствии, Ситрис дал мне едва ли не единственное средство против реакции Стлока, которое было у него в распоряжении.
Что заставило его напоить меня именно этой, растущей исключительно в Пыльном Предгорье травой --- ещё одна неразгаданная загадка.

Странный, дурманящий пряный отвар сбил температуру и успокоил внутренние органы, но рассвет не принёс облегчения.
Жуткий крик преследовал меня по дороге на жреческий этаж.
Он звенел в посуде, из которой я ел.
Он завывал в ветре, который прорывался в узкие окна.
В душе я яростно тёр губкой тело и никак не мог смыть с себя этот ужасный предсмертный крик.

Потом я видел много жертвоприношений.
Вначале просто наблюдал, как работают старшие жрецы, потом ассистировал.
Под моим ножом испускали дух дети, взрослые, пленные иноплеменники.
Но каждый раз в ночь жертвоприношения мне снился один и тот же сон --- тёмная комната без окон и дверей и связанный идол, молча смотревший на меня чёрными, лишёнными белков глазами.

\section{Чужие сны}

Костёр нашёл меня на берегу моря, на обычном месте.
С воды дул зябкий бриз;
я пытался устроиться между скал.

--- Нашёл место, --- проворчал врач.
--- А если простудишься?
Подумай хоть о личинках, если о своём здоровье не думаешь.

Костёр набросил на меня свой свитер с вязаными белыми собаками.
Свитер был мне велик, и я в нём сразу утонул.
Затем Костёр натянул мне на голову шапку из имитации меха --- тоже в виде собачьей головы.
Шапка забавно прядала ушами и вращала глазами, фокусируясь на источниках резких звуков.

--- Меня одолевают странные сны, --- признался я.
Я вдруг понял, что среди моего окружения Костёр --- едва ли не единственный, кто сможет это понять.
--- Знаешь, все эти разговоры о войне, всё это оружие... меня это как-то выбило из колеи.

--- И не одного тебя, --- хмыкнул Костёр, поправив мне шапку.
--- Всё поселение талдычит, что кодекс тси под угрозой.

--- Баночка ни в чём...

--- Баночку никто не обвиняет, --- прервал меня врач.
--- В конце концов, его странные, но \emph{пока что} безобидные увлечения --- лишь видимое следствие глубинного процесса.
Народ пока ещё не настолько одичал, чтобы этого не понимать.
Так что тебе снится?

--- Мне снилась пустыня.
Я ужасно устал... так трудно быть командиром.
Понимаешь, я устал во сне.

--- Пустыня, говоришь, --- отозвался Костёр.

--- Страшная жара, разорённый Улей.
Он был покинут, --- с жаром заговорил я.
--- У меня онемели две ноги, я едва шёл.
Кани, люди, планты вокруг меня, все такие же измождённые.
Они сказали, что нельзя заходить в этот Улей, что он...

--- Что мы станем лёгкой добычей для диверсантов-самоубийц, которые носят взрывные устройства на себе.
Матка засылала таких...

--- Вместе с разведчиками, --- подхватил я.
--- Четыре разведчика, два диверсанта...

Я поперхнулся и схватил Костра за шировот.

--- Откуда ты это знаешь?

Костёр болтался у меня в руках как кукла.
Усталая кукла.

--- Откуда я это знаю?
Из истории?

--- Да какая история, --- усмехнулся Костёр.
--- Ты уверен, что твои сны --- только твои?

Врач встал и потянулся.

--- Мне иногда кажется, что ты и не умирал у меня на руках, командир.
Что и я не оставил свои кости где-то там, в радиоактивном песке.
Что мы так и шли вместе --- мрак Красных Каньонов, ветреные вечнозелёные пустоши, которые терпко пахли смолой и гниющими растениями, --- помнишь этот запах? --- котловина Кон-Тики, разорённый Улей и дальше, сюда, на Планету Трёх Материков.
Что та война продолжается до сих пор.
И я всё так же, как и тогда, готов идти с тобой хоть в ядерное пекло, хоть в бараки Прокажённых, хоть в пыточные камеры Матки.

Я слышал слова врача словно издалека.

--- Помнишь те четыреста двадцать километров?
Мы сбили ноги до костей, и ты, надрываясь, тащил меня на себе.

--- Осталось немного, --- пробормотал я.
--- Осталось совсем чуть-чуть.

--- Ты повторял это каждые четырнадцать шагов, словно заведённый, на языке, который я едва научился понимать.
А тот крик ликования, который разносился по радиоактивным полям?
Я не мог его слышать, командир --- я был уже мёртв.
Но я слышу его так же ясно, словно сам стоял рядом во всеми.
Эти два слова надрывали моё горло, как надрывали глотки прочих.

Этого просто не может быть.
Это просто совпадение.
Мы просто все немного сходим с ума...
Но я ведь тоже слышал этот крик.
У меня не было ни рук, ни ног, ни голоса, ни ушей, но я слышал...

--- Нет, это не сумасшествие, --- покачал головой Костёр.
--- Войны отпечатываются не только в памяти людей --- о них помнят биологические виды, планеты и сама ткань мироздания.
Война Тараканов была худшей, которую видели Ветви Земли --- из тех войн, которые начисто отбивают тягу к конфликтам.
Даже двести тысяч лет благоденствия не способны изгнать её из наших снов.
Когда война приходит вновь, возвращаемся и мы с тобой.

\emph{Возвращаемся.}

--- Жаль только, что не с кем об этом поговорить --- разве что с Баночкой, --- пожал плечами Костёр.
--- Но наш маленький зелёный воин ещё не настолько испугался собственных снов, чтобы ко мне прийти.
Он-то выжил.
Выжил, насеял желудей и к старости стал похож на морщинистую жабу.
Да всё молился своим амулетам, чтобы мы с ним в одно посмертие угодили, а не каждый в своё...
Помнишь его молитвы, м?
Бормотание на грани слышимости, когда засыпаешь.

Конечно, я помнил.
Сложно забыть эти едва различимые слова, даже единожды их услышав в переплетении ночного шума.
Единственное, чего я так и не мог услышать --- собственное имя.
Далёкий шёпот Баночки расплывался каждый раз, когда друг поминал меня.

\begin{quote}
<<Маат, Славная мать, корни мои, соки моего хребта, любящая и всепрощающая.
Прости слабость твоему побегу и позволь остаться навечно с теми, кто стал мне дороже племени.
Маат, Славная мать, сила рук моих, влага моих глаз.
Не зелена их кожа, и кровь их красна, как противный тебе огонь, и не ведаю я, кто покровитель их в жизни и посмертии, но иссохну я без них, как иссыхают лишённые твоей любви>>.
\end{quote}

--- <<Маат, Славная мать, корни мои, соки моего хребта>>, --- прошептал Костёр, и ветер ответил воем в далёких скалах.
--- Так вот бывает.
Глупость, суеверие, а через двести тысяч лет слышно, словно вчера было сказано.

--- Возьми мою руку, четыре к пяти, --- вырвалось у меня.

--- Тебе ещё рано это говорить, командир, --- одёрнул меня врач, но тут же улыбнулся и погладил по голове.
--- Прими успокоительное и иди спать.
\emph{Эти} сны всегда приходят по одному.

\section{Лихие времена}

Едва очнувшись от беспокойного прерывистого сна, я бросился в зал.
Всегда есть надежда, что произошедшее тебе только приснилось.
А ещё есть надежда, что тебя обманывают глаза, когда видишь, что в толстом каменном перекрытии зияет дыра --- огромная, неровная, неуместная дыра, словно культя руки, словно безвременно умерший друг, словно навсегда покинутая родина, словно дом, оставшийся лишь в детских воспоминаниях.

Храм пустовал.
Нет, разумеется, здесь ещё лежали осколки плит, ещё клубилась каменная пыль, а у полуразрушенной лестницы уже примостились несколько досок и чемоданчик каменщика.
На нижнем этаже, кажется, даже кто-то разговаривал.
Но никогда на моей памяти храм не был настолько пустым.
В нише на стене стояла раскрашенная чаша;
я проспал погребение.
Вернее, меня не стали будить --- обидный до ужаса поступок, если не знать, что за ним стоит простая, но убийственная мысль: <<Каждый кхамит твоего сна на счету>>.
%{Make count every kchamit of your sleep.}

Учитель сидел в библиотеке и переписывал книгу, как и всегда.
Медленно догорала свеча, как и всегда.
Однако, подойдя ближе, я увидел, что переписной лист был покрыт одним и тем же незнакомым иероглифом.

Я взял лист из стопки.
То же самое.
И ещё один такой же.
Последние несколько кхамит Трукхвал, словно испорченный колодезный ворот, проворачивал один и тот же знак.
Мысль жреца остановилась на третьей странице <<Обычая и ритуала>>;
узловатые измученные руки продолжали писать.
Я испуганно потряс его за плечо.

--- Звоночек!

--- Я отправил весточку в Тхартхаахитр, --- глядя в стол, отсутствующим тоном ответил учитель.
--- Жрецы прибудут в ближайшее время.
Ликхмас-тари, я должен попросить тебя о большом одолжении.

--- Я весь внимание, учитель Трукхвал.

--- У нас не осталось ни врачей, ни школьного учителя, ни жертвенников.
Митликх обещал вернуться, но когда --- неясно.

--- Ты ему написал?

--- Ему написали ещё ночью, но ты сам знаешь, где находится Кахрахан.
Мы с тобой совершенно одни, --- губы учителя задрожали.
--- Я пойду собирать травы и лечить.
Воины будут следить за знаком кирпича, ты возьмёшь жертвы.
Но я хочу попросить тебя заглядывать ещё и в школу.

--- Учитель, мы не выдержим и полдекады работы в таком ритме.
Ведь есть же жрецы, в своё время ушедшие из Храма?
Есть же Храмы ближе Тхартхаахитра?
Может, где-то рядом проходят отряды чести или даже сами Люди Золотой Пчелы, а мы не знаем?
Почему не попросить их?

--- Твоя кормилица сказала, что Люди Золотой Пчелы сейчас не в землях сели, они ушли куда-то в сторону Ожидания Вести, возможно, что и за море.
И вообще, судя по её словам, это слухи --- никто не уверен, что Золотая Пчела по-прежнему ходит по земле, а не по жилам джунглей.
Во всех Храмах сейчас нехватка.
Когда пару лет назад обезглавило Храм Травинхала, жрецы ехали аж с Запада.
Да и ты знаешь, когда и откуда приехали Кхатрим и Ситлам.

Я опустил голову.

--- Да, я тоже понятия не имею, как это всё успеть, но у нас просто нет выбора, и поэтому школу и библиотеку придётся...

--- Я возьму на себя школу, --- опередил я Трукхвала.
--- Травы тоже отдай мне, но не оставляй библиотеку.
Выдержим, сколько сможем.
Надеюсь, духи помогут и больных будет мало.
Брошенные книги рассыпаются в пыль, ты ведь помнишь.

Учитель слабо улыбнулся и вдруг заплакал.

--- Ликхмас-тари, я воспитал самого лучшего жреца на всей Короне.

\section{[] Как чудесен мир}

\textspace

Вереск корпела над своим творением день и ночь.
Она с маниакальной скрупулёзностью готовила каждое блюдо, кормила всех, кого только могла поймать, брала пробы и подвергала компоненты ферментативному разложению, а также сравнивала результаты квантово-химического анализа с субъективными ощущениями тси.

--- Ну куда ты! --- пытался урезонить её Мак.
--- Все говорят, что ты готовишь вкусно, абсолютно все.
И переваривают, и чувствуют себя хорошо.
\ml{$0$}
{Ты хочешь достичь идеала, который не сможет оценить вкусовой анализатор сапиента!}
{You want to achieve the ideal couldn't be recognized by gustatory system of sapient!''}

\ml{$0$}
{--- Это ваш не сможет, --- парировала Вереск.}
{``By your one,'' Heather used to answer.}

И вот наконец книга была дописана.
Узнал я об этом ранним утром, за завтраком.
Заяц пришла ко мне в слезах.

--- Вереск умерла.

--- Что?

--- Повесилась за шею на промышленном тросике.
Это не несчастный случай.
Она планировала всё заранее.

Заяц сунула мне кусок ткани с неровными иероглифами.

\begin{quote}
<<Дорогие мои тси.
Я закончила свою книгу --- Десять тысяч блюд.
Теперь, даже если теплица выйдет из строя, вы будете хорошо кушать и жить счастливо на этой планете.
А я не могу.
Не упрекайте меня, прошу вас.
Я честно пыталась начать жизнь сначала, я нашла в лице многих из вас друзей.
Однако сложно свыкнуться с мыслью, что ты находишься не там, где должна.
Недавно я узнала, что молодые племена верят в жизнь после смерти в местах, которых не существует.
Как чудесен мир!
Какие сюрпризы он ещё преподнесёт нашим учёным!
Если бы верования этих людей были правдой, я хотела бы вернуться на родную планету --- такую, какой она осталась в самых светлых моих воспоминаниях.\\~\\
Пчела-Нюхающая-Вереск, инженер-технолог пищевой промышленности Тси-Ди>>.
\end{quote}

\section{Славь духов}

Следующая декада прошла словно в тумане.
Я держался как мог --- в ход пошли и крепкие отвары, и пощёчины.
Количество часов в школе пришлось сократить, чтобы я мог хотя бы поспать.
Ни о какой отчётности речи не шло --- я писал темы уроков на стене кельи перед тем, как рухнуть на лежанку.
После ещё одной ночи жертвоприношения я притащил Трукхвалу стебли ядовитого мертволаза вместо молитвенных маков.
Я ожидал гневной проповеди, но учитель, глядя на моё сонное, измученное лицо, снова расплакался.

Чханэ всё это время пропадала на тренировочной площадке.
Они с Ликхэ поделили часы напополам --- патрулирование отнимало меньше сил.
На каждом воине теперь было три, а то и четыре квартала.
Как-то, забежав домой перекусить, я увидел девушек, Ситриса и Эрликха спящими вповалку в углу зала.
Растерянная кормилица убирала со стола нетронутую еду.

--- Лисёнок, сегодня с тобой на крышу пойдёт Кхарас, а в дозор Кхохо возьмёт старшин кварталов, --- сказала она.
--- Девочки уже не кушают от усталости.

--- Они едут? --- этот вопрос я повторял последние четыре дня непрерывно.

Кхотлам развела руками и пододвинула мне тарелку с ещё тёплым супом.
Я взял черпачок и осознал, что тоже не хочу есть.

Кхарас, и прежде неразговорчивый, теперь молчал целыми днями.
Он стал ещё более сухим и холодным, словно северная скала.
Жертвоприношение вождь ассистировал со знанием --- из воинов таким могла похвастаться разве что Кхохо;
однако с его губ так и не слетело ни одного звука.

Мне ужасно хотелось закончить жертвоприношение быстрее.
Я едва сдерживался, чтобы не сделать надрез больше, чтобы не оставить сочиться кровь из обрубленной руки...
Но каждый раз у меня перед глазами почему-то вставал Кхатрим.

--- Ликхмас, только попробуй, --- говорил он звенящим бесплотным голосом, подавая мне щипцы.
--- Надрез там, подтёк здесь --- и в декаду на два жертвоприношения больше.
Что ты скажешь этим детям?
Что ты просто хотел спать?

\ml{$0$}
{--- Но ведь у утки клюв больше! --- пытался возразить я.}
{``But the duck has a bigger beak!'' I tried to object.}

\ml{$0$}
{--- Ликхмас, это логика Преддверия Снов, в которые тебе сейчас нельзя.}
{``\Likchmas, it sounds like the Hallway of Dreams, and you're not allowed in there.''}

\ml{$0$}
{--- Череп голубой.}
{``I see skulls of blue.''}

\ml{$0$}
{--- Прости, но я тебя ударю, --- сказал Кхатрим и размахнулся.}
{``I'm sorry, but I have to hit you,'' \Kchatrim\ said, then swung.}

...В себя меня привела пощёчина Кхараса.

\ml{$0$}
{--- Живой?}
{``Alive?''}

\ml{$0$}
{--- Угу.}
{``Mm-hmm.''}

\ml{$0$}
{--- Прибираться будем?}
{``Cleaning?''}

\ml{$0$}
{--- Угу.}
{``Mm-hmm.''}

Закончив уборку, мы зашли в крипту и наконец-то рухнули на скамью.

--- Мягко, --- пробормотал вождь, щупая холодный голый камень, и погрузился в сон.
Я обнял его --- и понял, что камень обнимать куда мягче и теплей.

Митликх вернулся четырнадцатого числа.
Он вбежал в школу посреди урока и, шепнув мне <<Иди спать>>, как ни в чём не бывало повёл урок дальше.
Я сел на скамью рядом с учениками и на мгновение прикрыл глаза.
Открыл уже у себя в келье спустя сутки --- дети в конце урока отнесли меня на лежанку, оставив у изголовья записочки с пожеланием выспаться и три пахучих увядших венка.

--- Проснулся наконец-то, --- проворчала Кхохо.
Оказывается, она всё это время сидела в углу, положив голову на руки.
--- Твой завтрак у лежанки.

Сон сползал неохотно, словно засахарившийся мёд.
Да, завтрак действительно стоял.
Правда, Кхохо не удержалась и отгрызла корку у куриного пирога, однако чтобы воительница носила кому-то еду в постель --- дело неслыханное.

--- Кормить с ложечки не буду, ты уж извини --- ненавижу эту детскую возню, --- скривилась Кхохо.
\ml{$0$}
{--- Жри своими руками.}
{``Use your own hands.}
\ml{$0$}
{А лучше --- ртом.}
{Or mouth---that's better.''}

\ml{$0$}
{--- А...}
{``What about---''}

\ml{$0$}
{--- А корку отгрызла не я.}
{``It's not me who chewed the crust off.}
\ml{$0$}
{Это новый Первый жрец.}
{New Head Priest did.}
\ml{$0$}
{Я шла с подносом, а он вцепился в пирог зубами, как зверь.}
{I was bringing the tray, and he locked his teeth on the cake's flesh, like a beast.}
\ml{$0$}
{Я отбила самую лучшую часть.}
{I won back the best piece for you.''}

Я улыбнулся, чувствуя прилив нежности к Кхохо.
Как же я люблю эту странную, прожорливую, мстительную лгунью.

--- Когда они прибыли, Уголёк?

\ml{$0$}
{--- Вчера днём.}
{``Yesterday afternoon.}
Они даже успели провести расследование диверсии.

--- И это не хака, --- я уже успел обдумать ситуацию за прошедшую декаду.

\ml{$0$}
{--- Эти дикари нас с тобой спасли, Лис.}
{``These savages are our saviours, Fox.}
Идолы проникли на стоянку каравана и подменили свинцовый блеск в коробках на какое-то грязное взрывчатое вещество --- вероятно, смесь кусачей бумаги\FM\ и отходов металлургии.
\FA{
Кусачей бумагой сели называли нитроцеллюлозу.
}
Диверсант уже на нашем складе подорвал коробки вместе с собой.
Эта тупая тварь даже не знала, что именно ей приказали поджечь.

Кхохо смачно плюнула на стену и тут же, оглянувшись на меня, вытерла плевок рукавом рубахи.

--- Скорее всего, идолы намеревались атаковать город малыми силами во время радужного безумия.
Они всё рассчитали.
Было время трапезы, потом был запланирован общий сбор совета.
Если бы не эти внезапные переговоры с хака, из-за которых всё пришлось перенести, пал бы весь Храм.
Но кихотр лёг не в их пользу --- выжили два жреца и почти половина Нижнего этажа.
Трукхвал пошёл во Двор за какими-то бумагами, громыхнуло прямо у него за спиной.

Я откинулся на подушку.
Учителю постоянно пеняли, что он опаздывает на трапезу и ему не достаётся супов;
в этот раз ему не досталось хэситра.

\ml{$0$}
{--- Был ли диверсант?}
{``Was there one to detonate?}
\ml{$0$}
{Кусачей бумаге достаточно неаккуратного обращения.}
{It's enough to handle biting paper roughly.''}

\ml{$0$}
{--- Был.}
{``It was.}
\ml{$0$}
{Куски его зелёной башки я сама нашла.}
{I've found pieces of its green head.''}

\ml{$0$}
{--- Трукхвал же лично проверил коробки.}
{``\Trukchual\ checked all the boxes by himself.}
\ml{$0$}
{Там был свинцовый блеск.}
{They definitely contained lead glance.''}

\ml{$0$}
{--- Ликхмас, голову приложи.}
{``\Likchmas, use your head.}
\ml{$0$}
{Конечно, они замаскировали взрывчатку тонким слоем свинцового блеска.}
{Of course explosive was masked with thin layer of lead glance.}
\ml{$0$}
{Я бы сделала так же.}
{I would do the same.''}

--- А как они поняли, что коробки пойдут в храм?
И, самое главное, как они узнали время сбора?

--- Всё было написано --- и на коробках написано, и про сбор объявление вывешено во всех кварталах, --- горько хмыкнула Кхохо.
--- Ну ты же знаешь --- если кто-то из жрецов на выезде, крестьяне обязательно прочитают и ему расскажут.
Оказывается, идолы тоже умеют читать.
\ml{$0$}
{Ты жри давай, я зря, что ли, эту баланду тащила?}
{Feed your fucking self, dude, why you think I got that grub?}
\ml{$0$}
{Пусть работает желудок, а не голова.}
{Let your stomach work, not your brain.}
\ml{$0$}
{Прошлое --- в прошлом.}
{The past is in the past.''}

Я посмотрел на пищу.
Вот оно --- счастье.
Я сладко потянулся и пододвинул к себе поднос, предвкушая полный желудок и ещё одни сутки в постели.

\ml{$0$}
{--- Слава духам.}
{``Praise Spirits.''}

\ml{$0$}
{--- Иди в зал и славь их там, если ещё останется желание.}
{``Go to the hall and praise Them there, if you'll still want to.}
\ml{$0$}
{Я воздержусь.}
{I pass.''}

И, вздохнув, Кхохо снова погрузилась в глубокий сон.

\asterism

Пока я ел, окутанный сонной дымкой мозг лихорадочно соображал над последними словами Кхохо.

Воительница выглядела не просто уставшей --- она испугалась.
Что-то напугало человека, который прорывался сквозь заросли копий и нелицемерной ухмылкой приветствовал смерть.
Что может быть страшнее смерти?

Я вдруг вспомнил слова Сиртху-лехэ: <<Зверь полагается на себя, и лишь в отчаянии ищет защитника>>.
Кхохо не испытывала страха, и мгновение настоящего испуга потрясло её до глубины души.
Она просто не знала, что делать с этим чувством.
О чём вспоминают в такие моменты?
О том, кто может защитить, и о месте, где можно спрятаться.
Звериный разум Кхохо, ведомый странными инстинктами, привёл её в мою келью.

Я ощутил необычность ситуации.
Кхохо пришла не к Ситрису, умелому бойцу, не к Кхарасу, командиру Нижнего этажа, не к Трукхвалу, у которого всегда можно было спрятаться от обязанностей и попить отвара, и не к кормилице, которая владела словом и справедливо разрешала конфликты.
Она пришла ко вчерашнему мальчишке, который не имел ни щепотки власти и мертвецки спал после бессонной декады.
И, разумеется, ни о какой заботе или признательности речь не шла --- всё это так же мало знакомо Кхохо, как и страх.
\ml{$0$}
{Воительница принесла пищу, чтобы задобрить непонятное существо, которое было её последней надеждой.}
{The warrior brought the food to bribe a weird being who was her last hope.}

И отгрызла корку у пирога.
Вполне понятный жест для того, кто знал Кхохо, однако и тут что-то неуловимо не вязалось одно с другим...

Нужно срочно узнать причину её страха.
Что-то подсказывало мне, что она находится в зале.

Поев, я поднял Кхохо на руки и уложил на свою лежанку.
\ml{$0$}
{Она стала лёгкой, словно перо, щёки ввалились, а скулы проступили, словно скалы из-под сгоревшей сельвы.}
{She felt as light as a feather, her cheeks looked sunken in, and cheekbones stuck out like stones from the burnt Silva.}
Воительница сильно похудела со времён нашего последнего спарринга.
Интересно, когда это было --- три, четыре декады назад?
Она вдруг открыла глаза и в упор посмотрела на меня.

--- Это не я слопала корку, Лисичка! --- жалобно-обиженно заявила воительница.
--- Не я!

\ml{$0$}
{Я улыбнулся.}
{I smiled.}

\ml{$0$}
{--- Я знаю, Уголёк.}
{``I know, Coal.}
\ml{$0$}
{Спи.}
{Sleep.}
\ml{$0$}
{Здесь ты в безопасности.}
{You're safe here.''}

Новые жрецы уже были в зале, за наспех сколоченным столом у целой стены.
Они разом повернулись ко мне, едва я ступил за порог;
семь пар пустых глаз взглянули в мои --- и у меня едва не подкосились ноги.
Я никогда не видел в человеческих глазах такой пустоты.
Это была не глупость и не бесчувственность --- пустота просто выглядела чужой, как чужда человеку Каменная ярость.
Сидевший на месте Первого жрец улыбнулся странной, кривой улыбкой:

--- Ликхмас ар'Люм.
\ml{$0$}
{Мы тебя ждали.}
{We've been waiting on you.}
\ml{$0$}
{Проходи и чувствуй себя как дома.}
{Come and be welcome.''}

Надкушенная корка пирога, уже изрядно обсиженная мухами и осами, сиротливо лежала на углу стола.
Проследив за моим взглядом, Первый небрежно смахнул её на пол.

\section{Нарушенное обещание}

\epigraph
{Никогда не позволяйте изолироваться людям, обладающим опасными навыками и знаниями, всеми силами удерживайте их в обществе.
Иначе может настать день, когда вам придётся жить по их законам.}
{Март <<Одноглазый>> Митчелл, идеолог Эволюциона}

\ml{$0$}
{<<Меня не ждали>>.}
{\emph{``I'm not welcome.''}}

Митхэ изо всех сил пыталась уложиться в плотный график.
Вначале зайти к Серебряному на постоялый двор и проверить, чтобы слуги его накормили.
Затем проведать тяжелораненного наёмника, которого перетащили из храма в один из жилых домов под предлогом предстоящего празднества --- вопиющее неуважение.
Затем найти, вытащить из канавы и привести в человеческое состояние Эрхэ --- воительница была в запое с достопамятного дня, стоившего Митхэ оружия.
Затем заплатить за работу шорникам и забрать из починки цитру Атриса, который уже без разрешения врача ковылял по площади и рвался играть.
Венцом тяжёлому дню должен был стать совет Храма Тхартхаахитра.
Однако когда взмыленная Митхэ поднялась в величественный зал, она поняла --- совет начался едва ли не на кхамит раньше назначенного времени.

\ml{$0$}
{<<Меня не ждали>>.}
{\emph{``I'm not welcome.''}}
Жрецы бросали на Митхэ нервные взгляды.

--- Здравствуй, Митхэ ар'Кахр, --- как ни в чём не бывало поприветствовал воительницу Король-жрец и закончил:
--- Что ж, на сегодня все вопросы решены.
Все свободны.

--- Подождите, --- Митхэ махнула рукой, и уже поднявшиеся из кресел жрецы замерли.
--- Я прошу прощения, что пропустила совет.
\ml{$0$}
{Король-жрец, было ли вынесено решение по известному тебе делу?}
{Priest-king, has the case you know been resolved?}
\ml{$0$}
{Не мог бы ты его повторить для меня?}
{Could you repeat for me?''}

\ml{$0$}
{--- Да, разумеется, --- улыбнулся Король-жрец.}
{``Of course I could,'' Priest-king smiled.}
\ml{$0$}
{--- Как мы могли забыть про это дело.}
{``How dare we forget about that case.}
\ml{$0$}
{Сегодня вечером будет большое празднество!}
{We have a great celebration tonight!''}

Кто-то крикнул <<Вино рекой!>>, и воины огласили зал громовым хохотом.

--- Король-жрец, не увиливай.
Ты обещал амнистию тем троим кутрапам, которые сражались в моём отряде! --- вспылила Митхэ, перекрывая шум.

Хохот затих в одно мгновение, словно обрезали шёлковую ленту.
Наступило неловкое молчание;
цитра в руках Митхэ вопросительно звякнула на весь зал.

--- Я сказал, что посодействую, Митхэ ар'Кахр, --- проговорил Король-жрец, с некоторым подозрением глядя на цитру.
--- Я посодействовал, но мои люди не считают войну искуплением.
Разрушители и Насильники должны понести наказание.
Тем более один из них --- ноа-пират, а обсуждение участи кутрапа-чужеземца не стоит и двух часов сна.

--- Сильный, творящий произвол над слабым --- вот истинный Разрушитель и Насильник! --- рявкнула Митхэ.
--- И сила в этой ситуации --- отнюдь не трое людей, пытавшихся найти свою чашу в лоне племени!

--- Ты можешь дать гарантии, что они не займутся прежним ремеслом, когда им наскучит спокойная и размеренная сельская жизнь? --- осведомился Король-жрец.

--- Я лично выслежу и убью их, если это произойдёт!

--- Верю.
Но сколько сели пострадает до той поры?
Я не готов заплатить такую цену.

Митхэ замолчала, не зная, что ответить.

--- Разумеется, пока эти кутрапы под защитой твоего Храма, мы не будем их трогать, --- примирительно поднялась сухощавая ладонь.
--- Но смею напомнить, что Тхартхаахитр за Тремя Этажами, и если мои воины обнаружат их одних --- в городе ли, в храме ли, --- то поступят согласно закону.
\ml{$0$}
{Не обессудь.}
{Don't get me wrong.''}

--- На улицах будут праздновать, расфуфыренный ты петух! --- выкрикнула Согхо.
--- Дай им хотя бы несколько дней пожить без страха!

Король-жрец пропустил слова воительницы мимо ушей.

--- То есть если я хочу жить, я должен ходить с командиром за ручку, так, Король-жрец? --- подал голос Ситрис.

--- Я никогда не говорил и не скажу ничего подобного, Ситрис ар'Эр.
Если ты обиделся на мои слова --- прошу прощения.
Пусть празднество будет тебе утешением.

--- Я не держу обид на побеждённых, --- Ситрис отвесил издевательский полупоклон.

--- Не испытывай наше терпение, кутрап, --- ледяным тоном сказал боевой вождь Храма.

Митхэ подхватила цитру под мышку, взяла Ситриса за руку и под сдержанные смешки присутствующих вышла из зала.
Аурвелий последовал за командиром, не забыв наградить Короля-жреца демонстративным плевком.

\asterism

\epigraph
{Правители пойдут на самую низкую подлость,\\
Чтобы заполучить искреннюю верность солдата.\\
Верному не нужно платить,\\
Верного можно убить,\\
Верного можно забыть ---\\
Ценный товар для тирана.}
{Гимн наёмников Фоуф, планета Тысяча Башен}

--- Золото, остынь, --- пытался успокоить её Атрис.
--- Хочешь отвара?

--- Митхэ, перестань, --- в десятый раз повторял Ситрис.
--- Ведь судьба в любом случае забросит нас на Запад, верно?
Уйду с Аурвелием на Кристалл.

--- В чём-то я понимаю позицию жрецов, --- сказал менестрель.
--- Без обид, дружище: убийствами разбой не искупить.

--- Да какие обиды.
Ты прав.

--- Но Король-жрец мог и не обещать, и возможности сдержать обещание у него тоже были.
Это уже вероломство.

--- Это политика, --- поправил Ситрис.

--- Нет, дружище.
Политика --- это если бы Король-жрец обманул хака, а не...

--- Это в любом случае вероломство, --- перебила Митхэ.
--- Не имеет значение, свои или чужие.
На лжи мира не построишь.
Если враг не верит твоим угрозам, если враг не может положиться на твоё обещание, торговцы превратятся в диверсантов, странники --- в разведчиков, а дипломаты станут погремушкой, скрывающей лязг сабель.
И о мире можно забыть на столетия.

--- Те, кто умеет обманывать и бить в спину, обычно выигрывают, --- напомнил Ситрис.

--- Что выигрывают?
Мир, в котором невозможно доверять?

--- Жизнь, --- просто ответил разбойник.
--- Для того, кто едва её не потерял, это неплохой выигрыш.
И уж точно ценнее любых слов и обещаний.

--- Неважно, ценно ли обещание для Короля-жреца, Ситрис.
\emph{Моё} всё ещё имеет значение.

--- Ты ничего никому не должна, --- буркнул Ситрис.
--- Ты сделала, что могла, но обстоятельства были против тебя.

--- Отговорка Короля-жреца, Ситрис?

--- Что для него отговорка, для тебя --- истина, --- пожал плечами Ситрис.
--- Кто из вас Король-жрец, в конце концов?

--- Ни приказ, ни запрет не оправдывают бесчестный поступок.
Я это вам постоянно говорю.

--- Да-да, конечно.
Вы все такие честные, правильные, последовательные, как ездовые олени.

Митхэ раздражённо фыркнула.

--- Я не имел в виду Серебряного, он свой человек, --- поправился Ситрис.
--- Хорошо, как курочки, которые клюют то, что им насыпали, и откладывают яйца там, где им сделали гнёздышко.
Но истина такова --- обещания никого за горло не держат.
И рисунки твои на лице тоже, --- быстро проговорил разбойник, предвосхищая слова Митхэ.
--- Ну не будут тебя уважать сели, если ты поступишь вразрез с Кошачьими Слезами.
Так иди к пиратам, на Диком Юге о тебе ходят такие слухи, что ты в первый же день соберёшь армаду.
И не обязательно кого-то грабить, многие пираты --- такие же наёмники, как и ты.
Иди со мной на Кристалл, в конце концов, и доживай свой век в таёжном тупике, попивая вино холодными осенними вечерами.
С тобой твой мужчина.
Что тебе ещё нужно?

--- Я думала, ты поддерживаешь моё решение.

--- Если я киваю, это значит, что я тебя понял.
Это не значит, что я с тобой согласен.

--- Нет, ты не понял.

--- А ты уверена, что не понял именно я?
Ты сама до конца понимаешь все эти красивые фразы, которыми ты умащаешь свой путь --- Путь Кошачьих Слёз?
\ml{$0$}
{Ты смелый боец, ты хитрый тактик, но когда дело доходит до человеческих взаимоотношений, ты превращаешься в ребёнка, которого можно одурачить фигурой из четырёх пальцев.}
{You're brave in battle, you're tricky in tactics, but when it comes to human relations, you become a child can be fooled by four-fingered figure.}
\ml{$0$}
{Одно дело --- следовать принципам, и другое --- быть невольником своего положения...}
{It's one thing to apply the principles and quite another to be a slave of status---''}

Разбойник осёкся, поняв, что зашёл чересчур далеко.
Впрочем, Митхэ была чересчур занята своими мыслями, и смысл фразы до неё не дошёл.
Атрису же срочно, в ту самую михнет понадобилось настроить починенную цитру, и он ушёл в процесс с головой.

--- В общем, это, --- буркнул Ситрис.
--- Я, может, многого прошу, но советуйся со мной, если какие переговоры или что.
Всё-таки я теперь старшина, тебе не зазорно.
Я не говорю, что я мастер дипломатии типа Короля-жреца, но мне приходилось краденое сбывать и насчёт контрабанды договариваться.
Я знаю, у кого что спросить, да и сам могу понять, когда человек вихляться начинает.
Без обид: если тебя каждый встречный-поперечный будет вот так дурачить, Кодекс Ягуара останется только скрутить и скурить под выпивку.

--- Я всегда тебя выслушаю, --- сухо ответила Митхэ, --- но решения буду принимать сама.
Отряд чести --- плод \emph{моих} решений, и ты в моём отряде.
Разумеется, если ты не захочешь уйти.

Ситрис хмыкнул и хотел выбежать за дверь, но очень некстати вспомнил о приказе Короля-жреца.
\ml{$0$}
{Разбойник извернулся кошкой и снова оказался на том же месте, словно пол коридора был покрыт горячей смолой.}
{The outlaw made a squirm in air like a cat and returned to the same place, as if the corridor floor was covered with melted resin.}

\ml{$0$}
{--- Я... я, пожалуй, с... с вами посплю, у двери...}
{``I---I suppose I'll have a s---sleep here, on the doormat ....}
\ml{$0$}
{Хат, ты не против?..}
{\tHat, would you mind ...?''}

\section{[] Поминки}

\textspace

--- Это я виновата, --- плакала Заяц.

--- Нет, --- твёрдо сказал Мак.
--- Это был её осознанный выбор.
Она сделала для процветания народа всё, что могла, и ушла достойно.

Все промолчали.
Мак сказал то, что следовало, и добавить к этому было нечего.

Вечером Листик сообщила, что кольцевая теплица забрала тело Вереск полностью.
Тси почтили память подруги так, как следовало --- собрали большой стол перед Стальным Драконом и приготовили по три блюда из каждой главы её книги.

\section{Дешёвка}

Тем временем подготовка к празднику шла полным ходом.
По всему городу запахло праздничной едой.
Книга-кормилица давно не удостаивалась такого внимания;
Митхэ подозревала, что скоро сбудется старое шутливое пророчество --- туда, где приготовят одновременно все десять тысяч блюд, придёт отобедать даже Безымянный.

Активно поддерживал праздничное настроение и Атрис;
звуки его цитры привлекли ещё несколько бродячих менестрелей, словно огонь фонаря --- индиго-светляков.
Музыканты без разрешения храмовников зашли в резонирующую нишу и начали играть что-то совершенно немыслимое, собирая огромные толпы народа.
Многие тысячи собравшихся хлопали ладонями в унисон, как один;
в толпе то и дело возникали красивые водовороты из танцующих, по поверхности пробегали красивые волны из рук.
Храмовые жрецы помоложе наблюдали за происходящим сверху, свесив ноги с крыши храма;
кто-то из них даже притащил кувшин с вином и миски с угощениями.

--- Было очень здорово, --- похвалила менестреля Эрхэ.
Её пока ещё немного мутило после длительных возлияний, что совсем не мешало получать от происходящего удовольствие.
--- Тебе следует чаще петь весёлые песни --- про приключения, хитрых героев и забавные случаи.
Грустные у тебя сейчас выходят, хай, неискренне --- мордаха чересчур круглая.
Кстати, мне понравилась импровизация этого старика-зизоце с длинной бородой.
А, толстенькая девушка-трами тоже была хороша --- редко можно увидеть талантливого Пересмешника.
Как их зовут?

--- Я понятия не имею, как их всех зовут, --- засмеялся Атрис.
--- Мы, в общем-то, и не разговаривали.

--- Красиво играешь, паренёк, --- вдруг громко заявила Атрису проходящая мимо старуха.
--- А вот этих Пересмешников я не понимаю.
Во времена моей молодости такое даже музыкой бы не назвали.
А молодёжь слушает.

Старуха кивнула на Эрхэ и Акхсара и гордо удалилась.

--- Ну что, Снежок, мы теперь с тобой молодёжь, да? --- ухмыльнулась воительница.

--- Вы достаточно молоды, --- пожал плечами Атрис.
--- У вас молодые глаза.

--- Мне тут стало интересно, чего ты такой худой и оборванный был, когда Митхэ тебя привела, --- хмыкнул Акхсар.
--- Сколько вам в этот раз накидали?
У нас весь отряд чести столько за год не зарабатывает.

Атрис немного растерянно потряс цитрой;
ворох украшений внутри неё внушительно звякнул.

--- Менестрель должен намекать, что он играет ради золота, --- объяснил Атрис.
--- Например, ставить перед собой чашу.
Однако чаша --- это напоминание: ты слушаешь не бесплатно.
Я не хочу, чтобы моих слушателей что-то тяготило.
Когда было особенно туго с едой, я просто просил, чтобы меня накормили.
Тебе нужно золото?
Я могу отдать его тебе.

--- Ты лучше себе шёлк для праздничной одежды купи, дружище, --- весело посоветовал проходивший мимо Ситрис.
--- Зачем этому истукану золото?
На украшения для Эрхэ?
Или на вино, чтобы она и дальше топила свою несчастную любовь?

--- Иди в бездну, немытая шея.
Без тебя разберёмся.

Акхсар бросил нервный взгляд на покрасневшую воительницу и пошёл к храму.

Эрхэ, проводив его долгим печальным взглядом, подошла вплотную к Атрису и приникла к его уху:

--- Ты точно можешь без проблем расстаться с золотом?

--- Мне не на что его тратить, --- шёпотом ответил менестрель.

--- Тогда слушай.
То, что я тебе говорю, должно остаться между тобой и мной.
Этот мир довольно несправедлив --- один человек отдал всё, что имел, а второй не получил достойной оплаты...

... Перидотовая Лиана вернулась к хозяйке перед самым праздником.
Оружейник прислал питомицу.
Девочка с важным видом отдала Митхэ клинок, а затем, порывшись в карманах, ойкнув от испуга и порывшись ещё раз, с облегчением достала запечатанный свиток и увесистый мешочек.

<<В кошельке --- сдача.
Бриз Битвы принять не могу, это чересчур ценный предмет.
Всё произошедшее --- между тобой и улицей Серого Рассвета.
Благослови тебя Сит, Митхэ ар'Кахр>>, --- гласила коротенькая записка.
Перстень был аккуратно вплавлен в восковую печать.

--- Сдача? --- Митхэ озадаченно посмотрела на посланницу.

--- Вопрос удалось решить с минимальными затратами при сохранении качества, --- важным тоном ответила девочка.

--- Ничего не понимаю, --- призналась Митхэ, взвешивая кошелёк в руке.
--- Здесь как будто даже больше, чем я...

--- А я гравировку сделала, --- шёпотом похвасталась девочка.
Похоже, больше сдерживаться она не могла.
--- Тебе нравится?
Ну давай, посмотри!
Пожалуйста, --- тут же добавила она.

Митхэ выдернула клинок из ножен... и тут же задвинула обратно.
Она знала, что у неё будет время на то, чтобы изучить новое оружие, и одного взгляда хватило, чтобы понять --- эти михнет будут полны самых светлых чувств.
Образ сабли пламенел перед внутренним взором --- как живое море, как грозовые небеса, как осколок некоего идеального творения, распространившего волны гармонии по всей Вселенной.

\ml{$0$}
{--- Легенда Серого Рассвета, --- прошептала она.}
{``Gray Dawn Legend,'' she whispered.}
\ml{$0$}
{--- Знай, Глупец: эти достойны.}
{``Fool, know this: they are worthy.''}

--- Ну как? --- жалобно спросила девочка.

--- Это самая лучшая на свете гравировка, --- заверила её Митхэ.
\ml{$0$}
{--- Твои руки превратят любое знание в красоту.}
{``Your hands can turn knowledge into beauty.''}

Девочка просияла.
Митхэ повертела перстень Короля-жреца в руках и отдала его маленькой мастерице.

\ml{$0$}
{--- Это тебе за работу, --- ласково сказала она.}
{``It's for the work you did,'' she said gently.}
\ml{$0$}
{--- А кормильцу передай, что перстень стоит ровно столько, сколько в нём содержится золота.}
{``And tell your nurse this: the ring costs exactly as much gold it contains.''}

\section{Травы или котомка}

\ml{$0$}
{--- Что значит <<за работу>>?}
{``What does `for the work' mean?''}

\ml{$0$}
{--- То же, что и всегда, --- сказал я.}
{``Exactly what it always means,'' I said.}
\ml{$0$}
{--- Жрецов сейчас достаточно, и я прошу разрешения...}
{``The staff is adequate now, so I ask your permission---''}

\ml{$0$}
{--- Кто решил, что жрецов <<достаточно>>?}
{``Who decided that the staff is `adequate'?}
\ml{$0$}
{Кажется, делать такие выводы не в твоей компетенции.}
{I guess this lies beyond your competence.''}

\ml{$0$}
{Я промолчал.}
{I said nothing.}
\ml{$0$}
{Первый жрец начал сбивать меня с толку с первых же слов.}
{Even Head Priest's first few words confounded me.}
\ml{$0$}
{Что за компетенция нужна для очевидных выводов?}
{What competence is needed for an obvious conclusion like that?}

\ml{$0$}
{--- Так чего ты от меня хочешь?}
{``So what do you demand from me?''}

\ml{$0$}
{--- Отпуск, --- лаконично сказал я.}
{``Time off,'' I succinctly said.}
\ml{$0$}
{--- Заслуженный отпуск на два рассвета.}
{``Well-deserved two sunrises off.''}

\ml{$0$}
{--- Кажется, между нами некоторое недопонимание, Ликхмас ар'Люм.}
{``It seems there is a certain misunderstanting between us, \Likchmas\ ar'\Loem.}
\ml{$0$}
{В данный момент ты не являешься адептом моего Храма.}
{You are not an adept of my Temple.}
\ml{$0$}
{И ты им не будешь, пока я не приму соответствующее решение.}
{And you won't be before I decide on that matter.}
\ml{$0$}
{Так о каком отпуске речь?}
{So what time we're talking about?}
\ml{$0$}
{Ты свободен, как птица.}
{You're as free as a bird.}
Может, ты ошибся и пришёл просить меня, чтобы я принял тебя в Храм?

Я прищурился.
Хорошо, сыграем по твоим правилам.

\ml{$0$}
{--- Я хочу вступить в Храм через два рассвета.}
{``I want to join the Temple in two sunrises.''}

\ml{$0$}
{--- Ты действительно этого хочешь?}
{``Is that really what you want?''}

\ml{$0$}
{--- Я настаиваю, --- оскалился я.}
{``I insist,'' I grinned.}

\ml{$0$}
{--- Но к чему медлить, если ты \emph{настаиваешь}?}
{``So why wait if you \emph{insist}?}
\ml{$0$}
{Я приму тебя прямо сейчас, если ты ответишь на несколько вопросов.}
{I'll hire you right now if you answer a few questions right.}
\ml{$0$}
{Разумеется, ты можешь отказаться и поискать другой Храм.}
{Of course you can refuse and look for another Temple.''}

Я мысленно выругался.
Мне явно не хватало дипломатического опыта кормилицы.

--- Скажи мне, Ликхмас ар'Люм, что для тебя главное в жизни?

--- Дом и те, кто в нём живёт, --- не задумываясь ответил я.

--- Хорошо, я спрошу по-другому.
Что для тебя как жреца главное в этой жизни?

--- Быть щитом, и ни мгновения --- плетью.

Жрец поморщился, словно ему в рот попало нечто горькое и этим горьким его кормили по меньшей мере год.

\ml{$0$}
{--- Единственным мерилом защиты и нападения является закон.}
{``The only measure of defense and offense is the law.}
Законы сели мудры, и мы, Ликхмас ар'Люм, должны им следовать.
Что ты думаешь о Ситрисе ар'Эр?

--- Трааа... --- меня несколько смутила неожиданная смена темы, --- Ситрис --- один из лучших тактиков Тхитрона и помощник боевого вождя.
Он мой друг и наставник.
Этого достаточно для описания?

--- Вполне, --- согласился Первый.
--- Ваш дипломат сказала почти то же самое, а у меня нет причин ей не доверять.
А как ты относишься к Кхохо ар'Хетр?
Её манера общения плохо вяжется с весьма средними профессиональными навыками.
Я как-то не привык, чтобы воины храпели на посту и просыпались только от четвёртой пощёчины.

Я нахмурился.
Первый жрец не мог не знать, в каком режиме Храму пришлось работать последнюю декаду.

--- Хорошо, оставим вопрос о Кхохо открытым, --- вздохнул жрец.
\ml{$0$}
{--- Буду с тобой честен и прям: в Тхитроне у нас есть враги.}
{``I'll be honest and straight with you: there are people in \Tchitron\ who are enemies to us.''}

\ml{$0$}
{--- У нас --- это у кого именно? --- спросил я.}
{``Who exactly do you mean by `us'?'' I asked.}

Жрец снова поморщился от мнимой горечи.

\ml{$0$}
{--- У нас --- это у законопослушных жителей города Тхитрона и находящихся под его юрисдикцией земель сели.}
{`` `To us' means `to law-abiding citizens of \Tchitron\ and all the \Seli\ lands under its jurisdiction'.}
\ml{$0$}
{Мы --- то есть Верхний этаж и часть Того-Что-Ниже, --- полагаем, что и нападения хака, и недавнюю диверсию идолов помогали проводить именно враги, скрывающиеся --- или не скрывающиеся --- в городе.}
{We---I mean the Upstairs and a part of Whatever-Lies-Below---we suppose that both Haka invasion and Idol sabotage were supported by the enemies, hiding---or not hiding---in the city.}
Поэтому я надеюсь, Ликхмас ар'Люм, на твою помощь в защите народа и в обеспечении торжества закона над теми, кто его преступил.

--- Народ может защитить себя сам, --- заявил я.
--- Жрецы просто указывают путь, а воины координируют действия.
Почему ты, Первый жрец, решил взять такую тяжкую ношу на одного себя?

Первый ухмыльнулся.

\ml{$0$}
{--- Ты умён не по дождям, но недостаточно, чтобы понять суть происходящего, Ликхмас ар'Люм.}
{``You're wise beyond your rains, but not wise enough to understand what it is all about, \Likchmas\ ar'\Loem.}
\ml{$0$}
{Где ты работал последнее время?}
{Where were you working recently?''}

\ml{$0$}
{--- Я учил детей в школе и приносил жертвы.}
{``I teach children in the school, and sacrifice to the gods.''}

\ml{$0$}
{--- А не слишком ли ты молод для этого?}
{``Aren't you too young for that?''}

Однако я успел понять, к чему клонит жрец, ещё на первом вопросе, и решил завершить унизительную беседу на своих условиях:

\ml{$0$}
{--- Травы или котомка\FM.}
{``Herbs or knapsack\FM.''}
\FA{
Котомка --- обещание покинуть город.
}

--- Третья келья --- твоя, --- милостиво разрешил Первый.
--- На вопросы ты ответил верно, и от своих слов я не отступлю.
Но если врачам не хватит одного --- слышишь, Ликхмас ар'Люм? --- одного-единственного пузырька с микстурой, котомку я для тебя соберу лично.
В знак уважения к твоим заслугам, так сказать.
И помни --- слава о жреце, бросившем Храм, летит быстрее испуганной лани.

\chapter*{Интерлюдия IX. Кукольная лавка}
\addcontentsline{toc}{chapter}{Интерлюдия IX. Кукольная лавка}

\textbf{Сказка сотронских сели}

\epigraph
{Одиночество ребёнка наделяет куклу душой.}
{Януш Корчак.
Эпоха Последней Войны}

Говорят, что однажды Глупец оставил свои глупости и решил заработать на жизнь руками и бессонными ночами.
Только вот беда --- делать он ничего не умел, кроме как воду в колодце мутить да людям жизнь портить.
Пробовал Глупец дерево строгать --- только пальцев лишился.
Пробовал Глупец ветки вязать --- глаз себе выбил.
Пробовал Глупец камень резать --- и того хуже вышло.
Да и мастера в округе с детства своим ремеслом занимались, и подмастерий брали детьми, а ему уж за две сотни дождей перевалило;
кому нужен переросток?

Пришёл с такой нуждой Глупец в Сотрон.
А в Сотроне и корабельщики были лучше всех, и краснодеревщиков таких нигде не видывали, и кузнецы славились во всех концах Короны.
Кукольников только не было совсем --- ребятишки себе сами всё делали.

Так и занялся Глупец куклами.
По первости и дерево портил, и ткань изводил, и краску зазря лил;
потом принялся мастерить кукол из травы, щепы, дорожной грязи, мха, смолы и стрекозиных крыльев, подбирая правильный материал.
Смеялись люди над его творениями.
Но через год-другой продал-таки Глупец первую куколку.
А где первая, там и вторая, где вторая, там и десятая.
Пошло дело у Глупца --- мастером его признали.

И всё бы хорошо, но был в его куклах какой-то изъян.
Дети их любили как своих домашних, не расставались с ними, а взрослые посмотрят, обсудят, да понять всё не могут и не сойдутся меж собой, что не так.
Из дерева ли сделаны куклы, или из крашеной шерсти, или вообще из глины?
Глаза криво ли нарисованы или чересчур прямо?
Улыбается ли кукла, плачет или злится?
Кое-кто и побаиваться стал этих кукол --- известно, боятся люди непонятного.

И вот однажды случилось так, что одного мальчика задразнили его товарищи.
Был он тих и молчалив, и не мог им ответить.
А как известно --- чем больше молчишь, тем обиднее тебя дразнят.
Довели товарищи этого мальчика до слёз, и он побежал жаловаться кормильцам.
Кормилец, узнав, в чём дело, отвесил мальчику оплеуху.

<<Не смей жаловаться! --- закричал он.
--- Тебе предстоит кормить семью и защищать дом!>>

Мальчик ушёл из дому на улицу.
Он знал, что ему предстоит кормить семью и защищать дом, но в нём не было достаточно силы, чтобы даже защитить себя.
Он решил потренироваться на кукле.
Было очень просто взять и оторвать её крохотную голову.

В тот же день погиб лесоруб --- друг кормильца.
Его шея случайно попала под топор товарища.

Вскоре одна девочка, делая работу по дому, случайно опрокинула горшочек с мёдом.
В дом слетелось множество ос и пчёл; они съели весь мёд и унесли его обратно в улей.

<<Неумеха! --- закричала на девочку кормилица и вытолкала её из дома.
--- На ужин даже не приходи!>>

Заплакала девочка от обиды, взяла свою куколку --- подарок кормилицы --- и забросила её подальше в кусты.
В вечер пропал в сельве охотник --- брат кормилицы, и его так никогда и не нашли.

Поползли по Сотрону тёмные слухи.
Собрался как-то тайно квартал оружейников и порешил избавиться от кукол Глупца.
Ночью, когда дети заснули, их кормильцы забрали кукол, развели большой костёр и сожгли.

Утром квартал голосил --- дети оплакивали потерю любимых куколок.
К вечеру взвыл уже весь Сотрон --- на город напала болотная лихорадка, которую отродясь не видывали в Водоразделе.
Она не щадила никого --- ни молодого, ни старого, ни сильного, ни слабого, ни здорового, ни болезного.
По всему городу молитвенными маками расцвели погребальные костры.
Поняли тогда все люди, кто виноват в их бедах.

И подговорил один камнерез людей прийти к Глупцу с оружием и потребовать, чтобы хитроумный мастер сделал всё как было.
Собралась вокруг кукольной лавки большая толпа.

<<Выходи, шутник! --- кричали люди.
--- Или мы сожжём тебя в собственном жилище!>>

Глупец встретил толпу в рабочей одежде и с инструментами в руках;
он делал очередную куколку на продажу.

<<Отмени своё колдовство! --- потребовал у Глупца камнерез.
--- Отмени, или не будет тебе жизни в этом городе!>>

<<Вы хотите, чтобы я отнял куколки у ваших детей? --- спросил Глупец.
--- То, что вышло из моих рук, более не в моей власти>>.

<<Убить его!
Зарубить его!>> --- закричали люди.

Выхватил камнерез саблю и только занёс её над Глупцом, как пустился в пляс.
А следом за ним заплясали и прочие.
Кто-то плакал, кто-то кричал в гневе, а кого-то и смех уже пробрал, но все плясали под песенки оставленных дома детей, решивших поиграть с любимыми куколками.
Так и плясали, пока не свалились наземь в изнеможении.
А Глупец стоял и смотрел.

<<Отмени своё колдовство!
Духами тебя заклинаю! --- взмолился наконец камнерез.
--- Золото отдам, всё отдам, только отмени!
Не могу же я быть детской игрушкой всю оставшуюся жизнь!>>

<<А кем ты был до этого? --- спросил Глупец.
--- Твоя судьба всегда была в детских руках, я лишь свалял куклу из нитей твоей судьбы.
Законы природы я отменить не могу.
Хочешь жить --- думай головой>>.

Сорвал Глупец с себя фартук, подхватил посох, забросил в сумку хлеб --- и был таков.

Так и разошлись сотронцы по домам ни с чем.
Кто-то ещё пытался найти свою куклу среди десятков тысяч таких же, да где там!
Прочие же, смекнув, объяснили чадам, какая сила в их руках.
И не было больше в городе брошенных детей, и никто не плакал от обиды, и ни один взрослый больше не смел ударить ребёнка.

А когда дети повзрослели, куколки затерялись где-то в мешках со старыми игрушками.
Никто о них даже не вспомнил.

\chapter{Празднество}

\section{[] Возвращение старого друга}

\textspace

Баночка повеселел --- Нейросеть вернулась.
Насколько плант смог понять из её карканья, она вполне успешно вывела <<много>> птенцов, и они разлетелись.
Несмотря на то, что птица очень хорошо считала до пятидесяти, <<много>> могло означать двоих или четверых --- Баночка уже заметил, что она то ли любит приврать, то ли подвержена парамнезиям.
Впрочем, произошло это, по её словам, <<очень давно>>, и распространяться на эту тему птица не считала нужным.

\textspace

Листик стала часто запираться в тор-отсеке.
Однажды она забыла включить звукоизоляцию, и я услышал, как она плачет и рассказывает кольцевой теплице о своих проблемах.
Все её друзья погибли на Тси-Ди, а с оставшимися в живых она почему-то не смогла подружиться.
Никто её не слушает, потому что она маленькая и тихо говорит, а может, просто не хочет слушать?
Реактивы кончились.
В джунглях на том берегу есть месторождение необходимых минералов, но одна она ехать боится, там страшные аборигены царрокх.
Ей хочется любви, но всё время уходит на работу.
Ей хочется домой, в её милый маленький домик на краю Восьмого города.
Ей хочется ещё раз, хотя бы раз погладить кошку Шипучку, которая притащила котят незадолго до Катаклизма.
Ей ужасно хочется любви...

\section{Одиночество}

\epigraph{
\ml{$0$}
{Человек не должен оставаться один на один с системой.}
{A man must not be ever left to fight alone against the system.}
}{Постулат Элект}

<<Сит, помоги мне в моих печалях>>.
С этими словами сегодня, на самом краешке рассвета я высыпал с колодец перо с золотом --- впервые в жизни.

Я надеюсь, что вы никогда не ощутите то чувство одиночества, которое день за днём сковывало меня.
Мучительная декада после диверсии была полна надежды и ожидания;
новая волна мучений изо всех сил притворялась обычной жизнью --- ждать было нечего.

Третья келья, которую дал мне Первый жрец, больше прочих пострадала от взрыва.
Несмотря на несколько уборок, в ней всегда стояла пыль.
Из угла под столиком постоянно дуло;
к концу дня живот превращался в заледеневшую доску.

Вскоре я потерял колбы и часть реактивов.
Видимо, полку случайно задели строители во время ремонта, и её содержимое превратилось в груду мокрых пахучих осколков.
Как ни странно, с извинениями никто не пришёл, а я был чересчур подавлен, чтобы выяснять обстоятельства происшествия самостоятельно.
Осколки были убраны, но тяжёлый запах перебродивших трав и спирта так и остался висеть в воздухе, делая работу невыносимой.

--- Ты с ума сошёл! --- вытаращил глаза Эрликх, забежав в келью за лекарством.
--- Здесь не то что работать, здесь находиться нельзя!

Тяжёлый аптекарский столик и варочные стойки мы с Эрликхом кое-как вывинтили из пола и вынесли в храмовые земли, кашляя как смертельно больные.

--- Точно сдурел, --- бормотал воин, моргая слезящимися глазами.
--- Воздух --- чистая щёлочь.
Как ты вообще там сидел...

\ml{$0$}
{На двери кельи появилась красочная деревянная табличка <<\textsc{Осторожно! Вонючая зона!}>>.}
{On the cell door a bright wooden plaque appeared, it said `\textsc{Danger! Stinky area!}'.}
Она радовала меня целый день;
вечером же кто-то выдрал её вместе с гвоздями и наглухо закрыл дверь, залив замочную скважину смолой.
Больше Эрликх ко мне не приходил.

Я начал избегать шумных сборищ и вообще скоплений народа.
Странная мысль --- да, злоупотребляет положением Первый жрец, но кто первым поддержит его решение, если я оступлюсь?
Народ.
Те самые жители, основа основ, которым я отдал часы бесценного сна в достопамятную декаду.
Вселенная оказалась полна врагов;
они радостно кричали, смеялись и пели, и эти громкие звуки били меня по ушам, словно усаженные шипами пылеройские палицы.
Город кишел врагами, и едва ли хоть кто-то из них об этом догадывался --- слепое орудие, нависшее надо мной.

Никто не мог мне помочь.
Кормилица была занята своими делами;
пища в храме стала просто отвратительной, для меня её оставалось совсем немного, и я какое-то время ходил кушать домой, пока в один прекрасный день на столе не оказался прокисший суп.
Я попытался готовить сам, но быстро понял, что перестану успевать с травами.

Лекарства расходились чересчур быстро.
Я не понимал, в чём дело.
На травы уходили все силы, я проводил у аптекарского столика едва ли не всё время.
Тайна раскрылась на исходе страды: Саритр, новый жрец-врач, каждый вечер аккуратно сливал неиспользованные настойки у заднего входа.
Кому об этом сказать?
Кто в это поверит?

Чханэ отдалялась от меня всё больше.
Она проводила за воинскими занятиями всё время, а где она кушает и спит, я не знал.
В конце концов я, поймав её в коридоре, заявил, что нам больше не стоит быть вместе.
Подруга даже не стала спрашивать причину.

--- Хорошо, --- устало кивнула она и пошла дальше по своим делам.

Сегодня я достал старый походный мешок кормильца и долго смотрел на него.
Возможно, и в самом деле пришло время уйти.
Кормилица много рассказывала о море.
В молодости она держала торговый корабль, пока шторм в Могильном проливе не отправил его на дно.
У меня было чувство, что сейчас ни она, ни Храм не заметят моего ухода.
Все заняты.
У всех есть дела...
Меня уже не волновало, где и чем я буду жить.
Лишь бы убраться подальше от полного врагов города и прокисшего супа.

Впрочем, спустившись в зал, я передумал.
Манэ и Лимнэ мирно сидели и вырезали из цветной бумаги фонарики.
В сестрёнках не чувствовалось ни капли враждебности;
они не кричали и были готовы в любой момент уделить мне время --- это я знал совершенно точно.

Мне вдруг ужасно захотелось сесть рядом, прижать сестрёнок к себе и провести так хотя бы одну вечность --- пока не оттает промёрзший живот.
Я не успел даже проговорить желание про себя, как оказался между ними.
Руки Манэ и Лимнэ лежали у меня на солнечном сплетении.

--- Не дрожи, братик, ты уже в тепле, --- прошептала Манэ на ухо.

--- В этом храме любой замёрзнет, --- сказала Лимнэ, уткнувшись мне в шею.

Мало-помалу я успокоился.
В конце концов, поссориться с Первым --- дело одной михнет.
Может, действительно стоит потерпеть с уходом --- хотя бы до праздников?

\section{Традиция}

Когда-то давно на Тси-Ди существовало такое понятие, как праздники.
В отличие от обычных выходных, праздники и фестивали были привязаны к году или периоду в несколько лет, а также имели особую тематическую окраску;
тси выходили на улицы, собирались вместе и устраивали шумное весёлое представление.
Около пяти тысяч лет назад традиция была упразднена и заменена новой --- Театром Тысячи Положений.

Театр не был привязан к конкретной дате.
Его начало высчитывалось по общепланетному коэффициенту усталости --- тысячи механизмов непрерывно анализировали деятельность тси и объявляли Театр, как только ОКУ превышал определённое значение (к моему стыду, и порядок вычисления, и пороговое значение я напрочь позабыл).
За два дня все тси обсуждали на форумах тематику и правила нового Театра;
надо заметить, на моей памяти они ни разу не повторились.
Последний Театр, который мы застали на Тси-Ди, был посвящён мужскому полу с позиции культуры --- тси всех полов оделись в исторические <<мужские>> костюмы и отыграли всех мужских персонажей литературы и кино, каких только смогли вспомнить.
Мне не составило труда вжиться в роль диверсанта Раскалённого, мужчины-планта из романа-антиутопии <<Чёрная дверь>>, а вот кроха Листик повергла всех в дикий хохот, когда попыталась отыграть древнего жестокого полководца Рамеса.

Самым забавным было то, что никто, кроме узкого круга техников, не знал, сколько продлится Театр Тысячи Положений.
Он мог продлиться и день, и восемь дней, и даже больше.
Поэтому тси старались отдыхать с полной отдачей.
Самый долгий Театр за историю продлился четверть года;
ошибка в системе обнаружилась уже на исходе срока, когда трудяги сами вернулись к работе и даже самые отчаянные гуляки начали подозревать неладное.

Я чувствовал, что Театр необходимо устроить.
Однако когда?
На Тси-Ди этим ведала Машина, не могу же я назначить дату с бухты-барахты?..

--- Можешь, --- заявила Кошка, когда я решил посоветоваться с ней.
--- Именно с бухты-барахты.
Я тоже понятия не имею, как вычислялась дата, но я устала, и ты устал, и многие из моих знакомых выглядят уставшими.
Думаю, сам факт, что ты про вспомнил про Театр --- это достаточно веская причина его устроить.

И я назначил мероприятие на послезавтра, чтобы дать тси время на обсуждение тематики и правил.

\section{Ранние слёзы}

В тот год дожди зарядили возмутительно рано --- аж на пять дней против обычного.
Первый ночной дождик, который, согласно традиции, считался Слезами Ситхэ, пошёл на Змею 3, прямо перед праздником Мягких Рук.
Жрецы, посовещавшись (разумеется, без меня и Трукхвала), вышли к народу и объявили своё решение --- детей следовало порадовать на Змею 4, а влюблённые, мол, денёк потерпят со своим праздником.
Никто, впрочем, и не возражал.

Город в одно мгновение вырядился в гирлянды, бумажные фонарики и цветочные венки.
Обычно кормилица велела домашним готовить украшения и одежды заранее, дней эдак за пять.
В противном случае, уверяла она, украсить дом не хватит времени.
Только сейчас я понял всю глубину женского коварства --- на подготовку было достаточно нескольких кхамит, но она служила отличным поводом ничего не делать в предпраздники.

Починить разрушенную стену храма и пол зала строители не успели, но повели себя как настоящие друзья --- гирлянды, фонарики и цветочные венки висели даже на строительных лесах, превратив храм в какой-то чудесный дворец.
Вначале все избегали огромной дыры, но украшения сделали своё дело --- вскоре половина Храма (разумеется, Нижняя) весело прыгала через неё туда-сюда.

--- Может, сказать строителям, чтобы так и оставили? --- задумчиво сказала Чханэ.
--- Лестница скучновата.

--- Ага, --- с энтузиазмом подхватила Ликхэ.
--- В Тхартхаахитре есть Трёхэтажный храм, в Миситре --- Винтовой\FM, в Сотроне --- Плавучий, а в Тхитроне будет Дырчатый.
\FA{
Более точный перевод --- Праворезьбовой храм, из-за специфической формы --- храм выглядит как винтовая башня.
Но практически на все языки Тра-Ренкхаля название храма Миситра переводят как Винтовой или Кручёный.
}
\ml{$0$}
{Единственный в землях сели Дырчатый храм.}
{The only Holey temple in all the \Seli\ lands.''}

--- До первого землетрясения, --- ввернула Кхохо.
--- А то мало трупов было за последние дожди.

--- Да, как-то я не подумала, --- согласилась Чханэ и похлопала Ликхэ по плечу.
\ml{$0$}
{--- Не переживай, Орешек.}
{``Don't be upset, Nut-Nut.}
\ml{$0$}
{Дырчатым можно обозвать и вполне целый храм.}
{Even a restored temple can be called `Holey'.''}

Храму внеплановый отдых пришёлся очень кстати.
Старые воины (я не мог избавиться от привычки называть их именно так) наконец-то собрались вместе за обедом и обсудили последние события.

--- Кто на этот раз будет переодеваться на Слёзы Ситхэ? --- спросил Ситрис.
\ml{$0$}
{--- Кхохо?}
{``\Kchoho?''}

\ml{$0$}
{--- Опять? --- фыркнула воительница.}
{``Again?'' the warrior snorted.}
\ml{$0$}
{--- Я в прошлый раз в этом костюме чуть не грохнулась, и детишки узнали столько новых слов, что им на всю жизнь хватит.}
{``Last time I almost landed on my face because of that clothing, and children have learned new words---enough for the rest of their lives.}
\ml{$0$}
{Пусть кто-нибудь из... жрецов переодевается.}
{Let an---a priest change.}
\ml{$0$}
{О, Ликхмас, давай ты.}
{\^{O}, \Likchmas, you.''}

\ml{$0$}
{Кхохо едва не сказала <<старых жрецов>>.}
{\Kchoho\ was so close to say `an old priest'.}
\ml{$0$}
{Или у меня просто разыгралось воображение?}
{Maybe I was just imagining things?}

\ml{$0$}
{--- Молодой чересчур, --- покачал головой Эрликх.}
{``Too young,'' \Oerlikch\ shook his head.}
\ml{$0$}
{--- Ликхмас, не обижайся: когда я заходил в школу, я не сразу мог сообразить, где здесь учитель.}
{``\Likchmas, no offence: when I entered the school, it took some time to understand which one is the teacher.}
\ml{$0$}
{Трукхвала нарядите.}
{\Trukchual\ is more fit.''}

\ml{$0$}
{Нет, это было не моё воображение.}
{No, it wasn't an overactive imagination.}
\ml{$0$}
{Эрликх понял Кхохо так же, как я.}
{\Oerlikch\ understood \Kchoho's words as I did.}

--- У него нога разболелась, --- сказал я, ощущая прилив дружеской теплоты в животе и осознав, насколько я отвык от этого чувства.
\ml{$0$}
{--- Давайте я попробую, только потренируюсь голос постарше делать.}
{``Let me try, I'll just practice making my voice sound older.''}

--- Вот и решили, --- Кхохо с облегчением начала поглощать похлёбку.
--- И вообще, Слёзы Ситхэ --- это ритуал, а ритуал должны проводить жрецы.
Почему последние пятнадцать лет наряжаются только воины?
Как будто у нас больше нет...

\ml{$0$}
{--- Кхохо, заткнись, --- посоветовал Ситрис.}
{``\Kchoho, shut up,'' \Sitris\ adviced.}
\ml{$0$}
{--- Ты думаешь, \emph{эти} пойдут?}
{``You think \emph{they} will?''}

Воины промолчали --- слово прозвучало хуже всех витиеватых грязных ругательств, которые я когда-либо слышал от Кхохо.
Женщина смутилась и налегла на еду с удвоенным рвением.

--- Нет, вообще-то всё правильно, --- примирительно начала Ликхэ.
--- Первый жрец сказал, что...

\ml{$0$}
{--- Если мне потребуется мнение жреца, я спрошу того, который обедает с нами, --- мрачно буркнула Кхохо.}
{``If I need priest's opinion, I'll ask the one who shares our meal,'' \Kchoho\ grimly grunted.}
\ml{$0$}
{--- А \emph{этот} мог бы и нас спросить, прежде чем называть себя... ссс!}
{``\emph{That one} neglected to ask ours before he called himself a---sss\r{a}!''}

Кхохо издала всхлипывающий смешок боли.
Судя по звуку, она одновременно получила аж два удара под столом.
Кхарас смотрел на воительницу, как ястреб на мышь.

\ml{$0$}
{--- Кхохо, я уже объяснял тебе текущую ситуацию.}
{``\Kchoho, I have explained the situation.''}

\ml{$0$}
{--- Да как-то неубедительно, --- ответила Кхохо, морщась и потирая пах.}
{``You have not convinced,'' \Kchoho\ answered, wincing and rubbing her groin.}

\ml{$0$}
{--- В любом случает это приказ.}
{``And I have given my command.''}

\ml{$0$}
{--- Я исполню любой приказ, имеющий отношение к делу.}
{``I shall obey every your command which is relevant to my job.}
\ml{$0$}
{Но если ты хочешь заткнуть мне рот --- сначала трахни мой разложившийся труп.}
{But if you want to shut my mouth, fuck my decomposed body first.''}

--- Кхохо, --- тихо сказал Ситрис, --- давай не за столом.

--- А когда ещё, Ситрис?
Приходит в город мясное чучело и говорит: я ваш Первый жрец.

--- Кхохо, --- рыкнул Кхарас.

--- Я серьёзно, Кхарас, он меня разбудил и сказал, что он \emph{Первый жрец}!
Вот так, с порога!
Без приветствия, знакомства, голосования, без ничего, не спрашивая ни народ, ни Храм, ни меня лично!
\ml{$0$}
{Нормально, да?}
{Pretty clever, yeah?}
А давай я припрусь куда-нибудь в Кахрахан и объявлю себя боевым вождём?

\ml{$0$}
{--- Кхохо, не испытывай моё терпение.}
{``\Kchoho, stop trying me.}
\ml{$0$}
{Ещё слово --- и обедать ты будешь отдельно.}
{One more word and you have your meals alone.''}

\ml{$0$}
{--- Ты правда хочешь, чтобы я делилась соображениями с горожанами, а не с тобой, вождь?}
{``Do you really want me to share my thoughts with citizens, not with you, warchief?''}

Ситрис смотрел в стол с выражением смертельной скуки.
Ликхэ стушевалась;
девушка слабо понимала, о чём идёт речь.
Чханэ глотала кашу и морщилась, словно каждое слово причиняло ей боль.
Кхарас прожигал глазами невозмутимо жующую Кхохо.
Я чувствовал, что вождь очень хочет высказаться, и ему приходится прилагать сверхчеловеческие усилия, чтобы сдержаться.

\ml{$0$}
{--- Кхохо, если ты не понимаешь мой приказ, я попрошу, чтобы до тебя его донёс Первый жрец.}
{``\Kchoho, if my command is misunderstood, I'll ask Head Priest to give you an explanation.''}

В глазах Кхохо снова промелькнул дикий, иррациональный страх, но голос звучал бодро:

\ml{$0$}
{--- Ну давай, дитятко, беги жаловаться.}
{``Run to your nanny, child.''}

\ml{$0$}
{--- Если тебе что-то не нравится --- дверь там же, где и всегда!}
{``If there's something you don't like, the door is where it always is!''}

\ml{$0$}
{--- Благодаря тебе, Кхарас, и благодаря новому руководству я хорошо запомнила, где у нас дверь!}
{``Thanks to you, \Kcharas, and thanks to our new management, I know well where the door is!}
\ml{$0$}
{Если хочешь меня выкинуть --- имей смелость хотя бы сделать это своими руками, а не перекладывай ответственность на меня или на этого...}
{If you want to kick me out, at least find a courage to do it by your good self, without shifting responsibility onto me or that---!''}

\ml{$0$}
{--- Перестаньте, --- тихо сказал я, и Кхохо неожиданно для меня умолкла.}
{``Stop it,'' I quietly said, and \Kchoho\ unexpectedly obeyed.}
\ml{$0$}
{--- Я думаю, всем присутствующим уже доходчиво объяснили, где выход.}
{``I guess all present have been explained in simple language where the exit is.''}

Кхарас угрюмо занялся едой.
Ситрис, Эрликх, Кхохо и Чханэ сочувственно посмотрели на меня, переглянулись и тоже уткнулись в миски.
Я вдруг почувствовал, что у меня в животе не просто потеплело --- там начали распускаться цветы.

--- Насчёт Слёз, --- осторожно сказала Ликхэ.
--- Лисичка переоденется, а кто потащит мешок?

\ml{$0$}
{Все как один посмотрели на Чханэ.}
{All as one looked at \Chhanei.}

\ml{$0$}
{--- Мне не сложно, --- пожала плечами подруга.}
{``Sounds easy,'' my friend shrugged.}

\ml{$0$}
{--- Учись, --- наставительно обратился Ситрис к Кхохо.}
{``Learn from them,'' \Sitris\ recommended to \Kchoho.}
\ml{$0$}
{--- Молодёжь сознательнее тебя.}
{``The youth are more responsible than you.''}

\ml{$0$}
{--- Да ради всех духов, --- махнула рукой воительница.}
{``For all the Spirits' sake, let it be,'' the warrior waved her hand.}
\ml{$0$}
{--- Пусть хоть на празднике вместе погуляют.}
{``Let them walk together, at least until the celebration ends.''}

Лицо Чханэ приобрело яркий грейпфрутовый цвет.
Я тоже почувствовал, что краснею.
Ситрис вдруг ни с того ни с сего обнял Чханэ, поцеловал её в лоб и встал из-за стола.

--- Занесу Трукхвалу миску супа, что ли, --- рассеянно сообщил он, направляясь на кухню.

\asterism

\epigraph
{Сомнения --- единственное, что избавило мир от ужасов моего гения.}
{Эпитафия на безымянной могиле.
Захоронение Аквамарин, планета Тысяча Башен}

После завтрака я пошёл прогуляться, думая о предстоящем празднестве.
Перебирал в памяти подробности ритуала, маршрут, какие-то стихи...
Очнулся только уткнувшись в закрытую дверь школы.
Такова сила привычки.
Я уже не мог припомнить, когда последний раз был здесь, но ноги и руки, видимо, были не в курсе последних событий.

<<И вообще сегодня праздник.
С чего меня сюда понесло?>>

Я уже хотел повернуться и уйти, как услышал чей-то плач под лестницей.
Знакомая худенькая фигурка, грязная, без рубахи, сидела и давилась рыданиями.
Я подбежал к ней.

--- Митхэ!

Митхэ, увидев меня, отвернулась и зарыдала ещё громче.

--- Митхэ, что случилось?

--- Кормилица поругала.

--- Что она сказала тебе?

Митхэ всхлипнула.

--- В честь праздника нужно было сделать подношение духам.
Кормилица велела мне тоже.
Я спросила <<Зачем?
Я справляюсь и без духов>>.
Она начала говорить, что так нельзя, что во мне нет веры, и что мне будет трудно жить, потому что мне не будут помогать духи.
Ещё она сказала, что и дружба, и жизнь с мужчиной требуют веры.
И я заплакала.

--- Ты заплакала от того, что тебе будет трудно?

--- Да.
Мне и так трудно, понимаешь?
Все говорят, что мне очень повезло, что у меня хороший дом, что меня воспитывают хорошие люди.
Но мне от этого не становится легче.
Они меня не понимают, а я не понимаю их.
Как можно верить, если ты знаешь или хочешь знать?

Я обнял ученицу.
Похоже, наступил тот самый момент, о котором предупреждал меня Трукхвал --- <<когда учитель становится учителем>>.

--- Митхэ, --- начал я.
--- Твоя кормилица пожила на земле и знает, как ей следует жить, чтобы процветать.
Но она --- не ты.
Ты особенная.
Кто-то считает это поводом для гордости, но для меня особенность --- это в первую очередь напоминание: твой путь к счастью будет не таким, как у прочих.

Я помолчал, собираясь с мыслями.
Митхэ посмотрела на меня.

\ml{$0$}
{--- То есть она не знает, что лучше для меня?}
{``In short, she doesn't know what's good for me, does she?''}

\ml{$0$}
{Я смешался.}
{I got confused.}
\ml{$0$}
{Ребёнок в который раз ставил меня в неудобное положение своей прямотой.}
{Child's directness put me on the spot, yet again.}

\ml{$0$}
{--- Никто не знает.}
{``No one knows.}
\ml{$0$}
{Единственное существо, которое может это узнать --- ты сама.}
{You are the only being who can find it out.''}

--- Я думала об этом, --- призналась девочка, --- и мне стало ещё обиднее, что меня ругают ни за что.

--- Митхэ, сомнение --- это величайший дар для того, кто наделён острым умом.
Где бы мы были, если бы не подвергали сомнению услышанное и увиденное?
Кто-то ошибается, кто-то намеренно пытается тебя обмануть, а есть ещё твои собственные чувства, которые имеют свои границы.
Всё, что за границей чувств, досягаемо лишь умом.
Всё, что на границе чувств, может быть искажено, как то, что видно на самом краю линзы.
Есть ещё твоё собственное воображение, которое необходимо отграничить от чувств.
Ум --- твой инструмент, но без сомнения этот инструмент бесполезен.

--- Но ведь кормилица права?
Мне будет трудно?

--- Искать путь --- трудная задача, --- сказал я.
--- Но я тебя прошу --- в минуты тяжёлого выбора, в облаке неопределённости, в момент самой противной нерешительности помни, что, возможно, ты проявляешь лучшее качество, которое у тебя есть.

Митхэ посмотрела на меня, зябко повела плечами и прижалась ко мне крепче.
В чертах её миниатюрного, по-детски круглого личика вдруг вспыхнули другие --- угловатые и резкие, змеёй прополз по щеке извилистый шрам...

Как и Кхохо, Митхэ искала защиты у меня.
Вот почему её занесло в школу.

--- Учитель Ликхмас...

--- Что, ученица?

--- Возьми меня к себе.

--- Я, хай... --- поперхнулся я.

--- Прошу тебя.
Храмовникам ведь можно воспитывать детей в городе, я узнавала.
У тебя есть дом где-то у реки, на улице Стриженого Кактуса, я видела, что ты с твоей женщиной уходите из храма на ночь.
Она добрая, я это знаю, я была у неё на тренировках.
Я буду мыть посуду и полы, я буду самым лучшим ребёнком на свете, только забери меня от них...

Митхэ прорвало.
Слова лились в мою рубаху потоком вперемешку с рыданиями.
Я гладил девочку по голове.
Прошло несколько михнет, прежде чем она успокоилась.

--- Ну вот, теперь ещё и это, --- прошептала она.
\ml{$0$}
{--- Я сомневаюсь, что люди, давшие мне жизнь, меня любят.}
{``I doubt if the people who gave me life give me love as well.''}

Вдруг мне в голову пришла идея.

\ml{$0$}
{--- Пойдём.}
{``Come.''}

Мы прошли в небольшую саговую рощицу на воинских землях.
Там из земли торчал большой деревянный столб.

Я вытащил из сумки маленький кукхватровый резец и аккуратно вырезал на податливой поверхности столба знак вопроса.
Митхэ с интересом следила за мной.

--- Вот, --- сказал я, отряхнув с ладоней деревянные крошки.
--- Это тотем Сомнения.

\ml{$0$}
{--- Зачем ты его сделал?}
{``Why have you built it?''}

\ml{$0$}
{--- Понятия не имею, --- признался я.}
{``No idea,'' I admitted.}
\ml{$0$}
{--- Возможно, это нужно нам с тобой.}
{``Maybe both you and I need it.''}

\ml{$0$}
{--- Ты уверен, что это нам нужно?}
{``Are you sure we need it?''}

\ml{$0$}
{--- Нет.}
{``I'm not.}
\ml{$0$}
{А ты?}
{You?''}

\ml{$0$}
{--- И я не уверена.}
{``Me neither.''}

Я почесал голову.

--- В таком случае могу сделать предположение, что тотем Сомнения действительно работает.
Но точно сказать не могу.

Митхэ вдруг ахнула и вытаращила глаза --- она поняла.
Затем мы оба расхохотались.

--- Возможно, туда стоит поставить тарелочку и раскрасить тотем красками, --- сказал я на обратном пути.

--- Я не уверена, что это нужно, но, вероятно, я это сделаю, --- со смехом отвечала ученица.

--- Возможно, я тебе даже помогу, --- заключил я, --- но это не точно.
А теперь нам нужно к колодцу.
Смоешь с себя грязь, немного обсохнешь --- и пойдём праздновать.

--- Так ты меня возьмёшь к себе?

--- Ты можешь приходить ко мне и уходить, когда захочешь.
Я думаю, это почти то же самое.

--- Нет, это не то же самое!

Митхэ нахмурилась и остановилась.

--- Не увиливай, --- пробормотала она, и я снова увидел в её глазах слёзы.
--- Мне от этого ещё больнее, чем от отказа.
Я понимаю, что есть закон, я понимаю, что есть другие люди со своими интересами.
Просто ответь: если бы в мире были только двое, ты и я, если бы не было законов и чужого мнения, ты бы меня приютил?

--- Без сомнений, --- вырвалось у меня.

Митхэ робко улыбнулась.

--- Кажется, я поняла, как действует Тотем.

\section{[] Обида}

--- И ты не поверишь, какую тему они выбрали, --- буркнул Баночка.

--- Судя по твоему недовольству --- культура аборигенов?

--- Они Кошке сделали больно, --- тихо сказал друг.
--- Я не буду одеваться в их одежды и танцевать их танцы.
Фонтанчик сказал то же самое.

--- А что Кошка?

--- Это мои чувства, а не её.
Я не простил.

--- Тогда просто побудьте с нами как друзья.
Если возникнут вопросы, я постараюсь объяснить народу, в чём дело.

Баночка кивнул.

--- Знаешь, я не могу решить, как к ним относиться.
Можно, конечно, списать всё на примитивные инстинкты, но относиться к ним как к животным неэтично.
Если же считать их разумными и равными нам --- то дружить я с ними не хочу.

Я промолчал.
Идея уже давно бродила в обществе тси, и Баночка лишь высказал её вслух.

Листик высказала ещё более необычное мнение.

--- Мне стыдно, --- заявила она.
--- Мы будем носить их одежды и имитировать их ритуалы ради развлечения.
Но ведь для царрокх ритуал --- это не развлечение, это жизнь!
Они искренне верят, они делают всё всерьёз!

--- Но ведь ты уже отыгрывала...

--- Исторические события --- это одно.
Отыгрывать жизнь племени, которое живёт по другую сторону пролива --- это больше похоже на насмешку.

\textspace

Кошку я нашёл в укромном уголке; она разбросала везде парящие видеокамеры и вооружилась компьютером.

--- Неужели ты даже сейчас хочешь что-то изучать? --- удивился я.

--- Это самый интересный момент, Небо, --- шепотом пояснила Кошка с горящими глазами.
--- Культуру царрокх я увидела воочию.
А сейчас увижу ту же культуру, пропущенную через восприятие тси.
Это поможет многое понять глубже.
Ты иди-иди, празднуй, не распугивай мне тут объект исследований.

\textspace

\section{Маска Сита}

--- Ликхмас, где ты был? --- набросился на меня Ситрис.
--- Ты хотя бы предупреждай, если задерживаешься, а то мы не знали --- то ли тебя ждать, то ли Кхохо одевать.
Народ уже собирается, время!

--- Важное дело, --- отрезал я.

--- Более важное, чем ритуал?
Рубаху снимай.
Вот маска. Вот плащ...
Всё, пошли наверх, почти готово.

Уже на лестнице я понял недовольство Кхохо --- костюм действительно был очень неудобным.
Маска натирала лоб и щёки, а ветки цеплялись за всё, что попадалось на пути.

--- Это ты ещё только половину костюма надел, --- захихикала Кхохо.
--- Рукава сейчас Кхарас принесёт.

\ml{$0$}
{--- Хорошенькое дельце! --- возмутился я.}
{``Brilliant!'' I said with indignation.}

\ml{$0$}
{--- Терпи.}
{``Patience.}
\ml{$0$}
{Я эту дрянь целых пять раз носила.}
{I wore this trash five fucking times.''}

--- Так, подождите, я быстро подгоню маску, --- я ринулся в мастерскую, шурша, как дикобраз.

--- Ликхмас, Ликхмас, плохая идея, брось, --- зашептал мне вслед Ситрис.

Быстро отыскав резец и пилку, я попытался снять заусеницы.
Хрусть --- и маска разломилась пополам.

--- Вот сказали же тебе, --- недовольно буркнула Ликхэ из-за плеча.
--- Думаешь, я не пыталась в позапрошлом году ещё?
Тащи скобы, в зелёном сундуке.

--- Новую уже давно надо, --- прошипел я в ответ.
--- Почему не заказали-то?
Праздник один день в году!

\ml{$0$}
{--- Да руки всё не доходили, --- пожала плечами Кхохо.}
{``We haven't got round to it,'' \Kchoho\ shrugged.}
\ml{$0$}
{--- Ликхмас, проведёшь ритуал --- хоть весь храм перестраивай.}
{``\Likchmas, you can rebuild all the temple afterwards.}
\ml{$0$}
{А сейчас залатай маску на скорую руку, чтоб тебя идолы поимели на перекрёстке!}
{But now fix the fucking mask quick, idols bonk you at a crossroads!}
\ml{$0$}
{Тащи скобы!..}
{Bring me staples ...!''}

\ml{$0$}
{--- ...Ликхэ, да загни ты эту скобу или выдерни к свиньям, она мне глаз выцарапает!}
{``... \Likchoe, bend or remove this staple to swine, it's gonna scratch my eye out!''}

\ml{$0$}
{--- Ликхмас, маска развалится!}
{``\Likchmas, the mask's gonna come apart!''}

\ml{$0$}
{--- А мой глаз?}
{``So what about my eye?}
\ml{$0$}
{Это вообще нормально --- скобы забивать остриями внутрь?}
{You think pointy ends on the inner side is a normal state of things?''}

\ml{$0$}
{--- Ты хочешь, чтобы они торчали на виду у всех?}
{``So they gotta be in front of people?''}

\ml{$0$}
{--- Я хочу, чтобы они торчали на почтительном расстоянии от моего глаза!}
{``They gotta keep a respectful distance from my eye!''}

--- Да какого вы тут творите? --- прошипел пришедший на шум Кхарас.
--- Быстро на площадь, мы уже опаздываем на десять михнет!

\ml{$0$}
{--- Это невыносимо, --- бормотал Ситрис.}
{``Unbearable,'' \Sitris\ mumbled.}
\ml{$0$}
{--- От него перья на ходу отваливаются.}
{``He's loosing feathers on the move.}
\ml{$0$}
{Ладно, будем считать, что так и задумано.}
{Well, let's imagine that's the idea.}
\ml{$0$}
{Смена времён года, преходящесть сущего и всё такое.}
{Season change, and evanescence of being, and that stuff.''}

\ml{$0$}
{--- Преходящесть кого?}
{``Evanescence of peeing?''}

\ml{$0$}
{--- Кхохо, заткнись.}
{``\Kchoho, shut up.}
\ml{$0$}
{Бери верёвку и привязывай к нему маску, прикроем её венком.}
{Take a rope and tie the mask to his head, we'll cover it by wreath.''}

\ml{$0$}
{--- Ай!}
{``\^{A}i!}
\ml{$0$}
{Волосы!}
{My hair!''}

\ml{$0$}
{--- Сам виноват.}
{``Your fault, \Likchmas.}
\ml{$0$}
{Все предложения по рационализации --- после ритуала, а не во время.}
{All proposals for rationalizing must be made after the ritual, not during.''}

\ml{$0$}
{--- \emph{До} ритуала, Ситрис!}
{``\emph{Before} the ritual, \Sitris!}
\ml{$0$}
{До!}
{Before!''}

\ml{$0$}
{--- Да какая разница?}
{``Who cares?}
\ml{$0$}
{Главное, что не сейчас.}
{Not now, that's the point.}
Так, теперь оставшуюся часть костюма...
Да быстрее, быстрее!..

... В себя я пришёл уже ближе к полудню на лежанке Чханэ.
Улыбающаяся подруга вытирала губкой мой исцарапанный маской лоб.

\ml{$0$}
{--- Это кошмар какой-то, --- выдохнул я.}
{``It's the worst kind of nightmare,'' I breathed out.}
--- Я бы лучше ещё раз к хака сходил...

--- Дети были в восторге, --- весело сообщила она.
--- Особенно когда ты запутался в рыбацких сетях и начал кричать <<Я рыба, я большая рыба, не кушайте меня!>>.

\ml{$0$}
{--- Слушай, ну почему нельзя было сделать \emph{удобный} костюм?}
{``Tell me, why couldn't they just make a \emph{comfortable} clothing?}
\ml{$0$}
{Страдания --- это часть ритуала?}
{Is misery a part of the ritual?}
\ml{$0$}
{Может, лесные духи обязательно должны носить колючие ветки?}
{Are the Silva Spirits used to wear prickles?''}

--- Кхохо уже сама сходила к портным и заказала.
\ml{$0$}
{Даже ей было больно на это смотреть.}
{Even for her, it was painful to watch.}
\ml{$0$}
{В этот раз веток не будет, только вышивка.}
{No twigs this time, just embroidery.}
Кстати, прибегала Митхэ ар'Тра э'Самитх, может помнишь, такая немного забитая, но понятливая девочка из хуторских.
Спрашивала, всё ли с тобой хорошо, и так странно на меня смотрела, знаешь, как будто ребёнка не кормили полдекады.
Я сказала, что ты спишь.
Может, и правда поспишь пару кхамит?
Я тебя разбужу.

--- Хорошо, --- пробормотал я, проваливаясь в сон.

\section{[] Новый знакомый}

\textspace

Парень был на редкость худой, угловатый, с длинной шеей.
О его ключицы можно было ненароком порезаться, а кадык торчал из горла, словно острый нос вросшего в скалу корабля.

Трукхвал казался угрюмым, но Митхэ удалось его разговорить.
Вскоре она поняла, что нашла интереснейшего собеседника.
Казалось, в свои двадцать восемь дождей он знает всё, чего только может коснуться разговор;
его язык был богат, словно все поэты прежних дождей вселились в одно тщедушное тело.

--- Скажи, Трукхвал, а что там, на горе?
Сколько лет жила в Тхартхаахитре, но не видела этого строения.

--- Его сложно увидеть, --- согласился ученик жреца.
--- Это Малый храм --- Хритра.
Когда-то жертвы богам приносили там.

--- Что случилось?

--- Землетрясения, --- объяснил Трукхвал.
--- Здесь землетрясения были не всегда --- они пришли примерно тысячу дождей назад.
Именно тогда и была выстроена Трёхэтажная пирамида.
Сейчас Хритра, хай, в опасном состоянии --- вот-вот рухнет.
Митрис ар'Люм, впрочем, ходит туда постоянно.

--- Кто это?

--- Новый библиотекарь, который с севера приехал.
Он немного необщительный, хоть и всегда вежлив, и помогает, и вообще старается не доставлять этим хлопот.
Книги переписывает по ночам и часто там, на горе.
Вероятно, у него оборудовано место с электрическим освещением --- при свечах ночью можно испортить глаза.
Одно время мы увлекались лампами --- такие хрустальные шарики с кристаллами, кажется, свинцового блеска и муассанита.
Очень сложные в изготовлении, но очень яркие.
У нас ещё осталось несколько десятков.

--- Что такое муассанит?

--- Минерал, залежи которого, кстати, как раз в скале под Хритра.
По крайней мере, были.
Камнерезы и ювелиры его очень любят --- крупные кристаллы идут в огранку, а мелочь, кажется, используют для полировки или сверления.

--- Ты всегда говоришь <<кажется>>, <<вероятно>>, но при этом оказываешься прав!

--- Если бы у жреца было оружие, его ковали бы из сомнения.

--- Хорошо сказано, --- одобрила Митхэ.
--- Кажется, я была несправедлива к нерешительным.
Можешь сводить меня в Хритра?
На экскурсию.
Показать залежи муассанита, ну и Малый храм заодно.

--- Хритра под замком, --- пожал плечами Трукхвал.
--- Может, тебе стоит попросить ключ у Митриса?
Если он у него есть, конечно.

--- Зачем вообще запирать руины?

--- Туда лазила молодёжь, когда хотели разделить ложе в необычной обстановке, --- улыбнулся Трукхвал.
--- Жрецы боялись, что кого-нибудь настигнет весьма необычная смерть.

--- А это-то ты откуда знаешь? --- захохотала Митхэ.

--- Слышал.

--- Слышал.
Сам лазил небось.

--- Нет, --- покачал головой Трукхвал.

--- А что так?
Мне было бы интересно попробовать.

--- Мне не очень.
Я в принципе секс не считаю интересным занятием.

--- Нет, я не это имела в виду.
Это нужно чувствовать вот здесь, --- Митхэ приложила руку к животу, --- а не здесь, --- указала на голову.

--- Мне рассказывали про влюблённость и сексуальное желание, но сам я такое не испытывал.
Я делил ложе и с женщиной, и с мужчиной --- ради нового опыта.
С мужчиной не вышло совсем, а вот женщина была хороша.
Секс похож на танец --- это беседа тел.
Секс похож и на письмо --- ты определённым образом проводишь пером по человеку, и пятно чернил вдруг обретает смысл.
Но книги читать мне показалось интереснее, и вопрос, как провести вечер, для меня не стоял.

--- Подожди, а как же Круг?
Ты участвовал?

--- В Мягких Руках участвую каждый раз.
Всё-таки я стану жрецом, а руки жреца обязаны знать ритуал.

--- Это единственная причина? --- допытывалась Митхэ.

Трукхвал устало закатил глаза.

--- Давай закроем тему.
Мне и в Храме разговоров про секс хватает.

--- Мне интересно узнать насчёт Круга.
Наркотик же на тебя действует?

--- Действует.
Голова кружится и язык немеет.

--- И всё?!

--- Всё.
Созвучие, кстати, я тоже не слышу, хотя про него все всегда говорят.

--- И никого не хочется?..

--- Послушай, мне приятен секс, но я не чувствую в нём потребности.
Я изучил все приёмы, и их гораздо меньше, чем даже тактик на столе дипломатии в Метритхис.
В определённый момент мне просто становится скучно.

--- Я не понимаю, как такое может быть, --- призналась Митхэ.

--- А я мало что понимаю из того, что касается человеческих взаимоотношений --- дружбы, интриг.
То же весьма ограниченное число тактик.
Но кто-то находит в этом удовольствие.

--- То есть для тебя интерес только в количестве тактик?

--- Для меня удовольствие от решения задачи пропорционально её сложности.

--- То есть секс для тебя --- задача?

--- Именно.

--- Теперь поняла.

--- А для тебя?

--- Скорее созерцание, --- сказала воительница.
--- Нот и красок всего двенадцать, но их сочетания не могут надоесть никогда.

--- Вопросов больше нет, --- улыбнулся Трукхвал.
--- Знаешь, я, кажется, вспомнил, где находится ключ.

--- Тебе нужно что-то взамен?

--- Взамен не взамен, но захвати мешочек сушёных фиников и воды для нас двоих.
Финики я люблю, а туда очень утомительно подниматься.

\section{[] Искристая зелень}

Влюблённые не дотерпели.
Вердикт вердиктом, но народ сам решает, когда устраивать праздники.
На столах осталось множество яств, которые потеряли бы свежесть уже через несколько кхамит, а крестьяне --- практичный народ.
Ближе к вечеру, когда веселье по поводу Слёз Ситхэ уже утихло, по всей площади начали разгораться совсем другие костры.

Кхарас хмуро молчал.
Он вообще не любил, когда дела шли не по плану, это выбивало его из колеи.
Новые жрецы тоже оставались в стороне.
Не успел я продрать глаза от полуденного сна, как на меня неожиданно свалилась вся организация.

--- Ликхмас, а кто в этот раз Круг начнёт?
Ты?

--- Ликхмас, как сегодня с погодой?
Зонты на всякий случай тащить?

--- Ликхмас, а кто за столы возле цветного камня отвечает?
Можно, мы у них немного еды к нашим заберём?

--- Ликхмааас!
Я палец в столе защемила!

--- Кхохо!
Иди в бездну со своим пальцем! --- не выдержав, заорал я.

Кхохо обиженно заморгала.
Спустя михнет я, перебирая самые витиеватые ругательства, вытаскивал защемлённый палец с помощью молотка и стамески.
Ситрис и Эрликх бились в истерике и даже не думали помогать.

---  <<Мы связаны, судьбу не изменить>>, --- процитировал Эрликх <<Легенду об обретении>>, и оба зашлись в новом приступе судорожного смеха.

\ml{$0$}
{--- Лиc, незачем бегать самому, --- сказала наконец Чханэ.}
{``Fox, you don't have to do this by yourself,'' \Chhanei\ finally said.}
\ml{$0$}
{--- Назначь ответственных, а там как пойдёт.}
{``Assign responsible persons, and let things go.''}

\ml{$0$}
{Пошло очень даже неплохо.}
{Things went pretty great.}
\ml{$0$}
{Стамеска пролетела пять шагов и вонзилась в столб в пяди от уха Ситриса.}
{The chisel has flown five steps and buried itselt in a stick---a hand away from \Sitris' ear.}

\ml{$0$}
{--- Быстро взял стамеску и помог!}
{``Take the bloody chisel and help!''}

\ml{$0$}
{--- Тихо-тихо, Ликхмас, сейчас сделаю, только не кипятись!}
{``Calm down, \Likchmas, I'll do it, just don't freak out!''}

--- Столбик, на тебе столы у гонга.
Так, народ, Ликас и Митхэ отвечают за костры, все вопросы к ним.
Ликхэ, Кхотрис-лехэ, наберите народ, чтобы гонг передвинули, мешает...
\ml{$0$}
{Кхохо, палец болит?}
{\Kchoho, the finger is aching?}
\ml{$0$}
{Отлично, будешь думать в следующий раз.}
{Good, next time you'll think twice.}
Травы для Круга на тебе.

Вскоре ко мне подошла Кхотлам.

--- Разобрался с организацией?

--- Ну, если можно так выразиться, --- я пожал плечами, оглядев кипевшую вокруг слаженную работу.
--- Раздал роли --- и все вроде как справляются без меня.

--- Мне этим и нравится работа купца, --- призналась кормилица.
--- Чтобы всё заработало правильно, нужно знать, что делаешь, и думать, что кому говоришь.
Зато потом, когда всё уляжется, просто сиди и занимайся своими делами.

До меня вдруг дошло.

--- А чего все они ко мне-то пошли? --- повернулся я к Кхотлам.
--- Ладно Круг, но за столы-то я каким боком отвечать должен?

Кхотлам повела бровями.

--- Ну, я отправила Храму официальный вопрос, почему они не занимаются организацией Мягких Рук.
Пришёл расстроенный Эрликх и на словах объяснил, что Верхний Этаж сидит по кельям и занимается своими делами, а моё письмо мне вернули, посоветовав съесть его с острым соусом.
Учитывая, что сам Эрликх никогда не опустился бы до таких выражений, я ему поверила.
Ну и раз мы опять остались без жрецов, я сказала старшинам, что ты главный.
\ml{$0$}
{Видишь --- и развеялся, и полезный опыт получил.}
{You see, you had some fun and gained some experience.}
\ml{$0$}
{Вдруг как-нибудь ночью пригодится.}
{One night it may be useful.''}

Кхотлам тут же вытащила вязание и принялась набирать петли.

\asterism

\textspace

Свою праздничную одежду я, как и все, вышивал сам.
На неё ушло немало золота --- я выбрал чёрный, как ночь, шёлк и безумно дорогие нитки <<искрящаяся зелень>>, которые привозили с Запада.
Первые цепочки стежков я сделал мотком, взятым из шкатулки кормилицы;
у любой кормилицы, в любом доме всегда есть такая шкатулка, в которой мелкие памятные вещицы с разных концов света лежат вперемежку с нитками, застёжками, гвоздиками и пришедшими в негодность заколками.
Когда Кхотлам увидела, то схватилась за голову:

--- Дитя, ты с ума сошёл!
Ты знаешь, сколько они стоят?
Да ещё объёмным швом, ты год потратишь!

Однако отступать было поздно.
Я действительно потратил год, причём не столько на вышивку, сколько на заработок золота.
Впрочем, жалеть о выборе не пришлось --- в первый же праздник я оказался в центре внимания.

--- Сам вышил? --- не переставал удивляться Кхатрим.
Он щупал и щупал мою одежду, пока Первый не оттащил его за шиворот.

--- <<Искристая зелень>>?
Ты дурень, нашёл на что золото спускать, --- резюмировал Столбик.
--- А если порвут в толпе?

--- Хаяй, это же домики!
Домики, заплетённые лианами! --- Хитрам прыгала вокруг и радовалась как ребёнок.

Справедливости ради замечу, что свою цену нитки оправдали --- штопать праздничную одежду мне ни разу не пришлось.

На ужин все спустились уже в праздничной одежде: Ситрис --- в оливковой, Ликхэ --- в переливчато-синей, а Кхарас --- в серой.
Чханэ надела зелёное платье прародительницы, объяснив, что свою праздничную одежду забыла дома.
Кхохо же вышла из общих комнат позже всех;
она завернулась в красное, как кровь, шёлковое полотнище прямо на голое тело.
Кхарас заворчал:

\ml{$0$}
{--- Кхохо, сколько можно! Как будто девочка двадцати пяти дождей!}
{``\Kchoho, how long can you dress like a twenty-five-rain-old girl!''}

--- В этот день полагается пить, танцевать и совокупляться, --- огрызнулась Кхохо.
--- Храм должен возглавлять, а не плестись в арьергарде!

--- Не волнуйся ты так, Кхарас, всё продумано, --- успокоил вождя Ситрис.
--- Ей просто лень стирать штаны и рубаху после Мягких Рук.

Я, Чханэ и Ликхэ переглянулись и дружно пошли переодеваться.

--- Лис, слушай, платье без штанов смотрится не очень странно? --- спросила Чханэ.

--- Мне так даже больше нравится, --- заверил я её.
--- Говорят, ноа ходят в таких платьях без штанов.

Ликхэ уже вознамерилась стянуть с себя передник, но задумалась и оставила, посетовав в очередной раз на своё тело:

--- Ну вот что с ними делать, а?

--- Ты сильно горбишься, --- заметил я.
--- Держи спину прямее, и грудь будет прыгать меньше.

--- Советы от человека, у которого нет груди, --- огрызнулась Ликхэ.

--- Это не мой совет.
Я слышал, как кормилица говорила это Эрхэ.

--- При всём уважении, Ликхмас, у Кхотлам \emph{такого} никогда не было!

--- Да ну, это здорово, когда грудь прыгает! --- засмеялась Чханэ.
--- У меня вот не прыгает.

--- Я бы с удовольствием поменялась, поверь, --- буркнула Ликхэ.

\section{[] Страсть жреца}

\textspace

--- Куда это ты ходила с Трукхвалом?

--- Он показывал мне город.

--- Не смеши меня, ты здесь родилась.

--- Я здесь давно не живу.

Первый ухмыльнулся.

--- Кажется, мальчишка сидит без дела.
Думаю, я найду ему дело на ближайшую декаду где-нибудь на заднем дворе храма.

--- О, я думаю, что ты его недооцениваешь, Нетрукх.
%{``\`{O}, I suppose he's underestimated, Netrukch.}
Если Трёхэтажный Храм считает, что место Трукхвала --- на заднем дворе, я с удовольствием заберу его с собой.

Нетрукх вздрогнул, словно Митхэ отвесила ему пощёчину.
\ml{$0$}
{Затем, прочистив горло, жрец подошёл ближе.}
{Then he cleared his throat and came closer.}

\ml{$0$}
{--- Митхэ, к чему все эти конфликты между моим Храмом и твоими наёмниками?}
{``\Mitchoe, how do they make sense, all these conflicts between my Temple and your mercenaries?}
\ml{$0$}
{Нам нечего делить.}
{There's nothing we should fight about.}
\ml{$0$}
{Более того, мы могли бы завязать крепкие дружеские связи.}
{Moreover, we could be good friends.}
\ml{$0$}
{Я предлагал это не раз и остаюсь верен себе.}
{I offered that more than once, and I'm true to myself.''}

Рука Первого жреца вдруг больно впилась Митхэ в плечо.

\ml{$0$}
{--- Ещё раз дотронешься до меня --- и я сломаю тебе руку в трёх местах, --- ровным голосом сказала Митхэ.}
{``Touch me again, and I'll break your arm in three places,'' \Mitchoe\ said in a flat voice.}

\ml{$0$}
{--- Ты ведь понимаешь, что у твоих слов могут быть последствия? --- осведомился жрец.}
{``I guess you know that your words would have consequences?'' the priest asked coldly.}

\ml{$0$}
{--- У твоих действий тоже.}
{``Just like your actions.}
\ml{$0$}
{Ещё раз дотронешься до меня --- и я сломаю тебе руку в трёх местах.}
{Touch me again, and I'll break your arm in three places.}
\ml{$0$}
{Мне отправить это предупреждение официальным письмом или до тебя уже дошло?}
{Should I repeat the warning in an official letter, or you already got it?''}

Похожая на клешню кисть с неохотой покинула плечо Митхэ.
Воительница кивнула и двинулась прочь.

--- Чем тебя привлёк этот бродяга? --- поинтересовался Нетрукх ей в спину.
--- Никогда бы не подумал, что ты предпочтёшь мясо из сточной канавы Первому жрецу.

Митхэ остановилась и, вздохнув, обернулась к собеседнику.

--- Атрис ар'Митр занимает пост Первого жреца отряда чести.
На сей раз я тебя прощаю.
Если подобные слова в его адрес прозвучат в будущем, мы будем вынуждены требовать официальные извинения.

--- Что? --- выплюнул Нетрукх.
--- Ты смеёшься надо мной?

--- В моём отряде чести я могу делать всё, что пожелаю.
Кажется, в этом мы с тобой похожи.

--- Я отказываюсь признавать его жрецом, и тем более --- Первым.
Он называется прозвищем, данным ему на рынке!
А родовой слог --- ар'Митр?
Ты его сейчас выдумала, на ходу?
\ml{$0$}
{У него нет ни родового слога, ни образования!}
{He has no family syllable, nor education!''}

\ml{$0$}
{--- Зато у нас есть Верхний Этаж, --- улыбнулась Митхэ.}
{``But we have the Upstairs,'' \Mitchoe\ smiled.}
\ml{$0$}
{--- Мы --- Бродячий Храм, а не кучка наёмников.}
{``We are the Wandering Temple, not a gang of mercenaries.}
\ml{$0$}
{Ты так много знаешь об Атрисе и не в курсе таких элементарных вещей?}
{You know so much about \Aatris, and you are so ignorant about these self-evident truths?}
\ml{$0$}
{Какая жалость.}
{What a pity.''}

Митхэ смерила Нетрукха презрительным взглядом и ушла в зал, оставив его кипеть в возмущении.

\section{[] Тихая ненависть}

\textspace

Первый жрец, проходя мимо, поймал меня за пояс.
Я едва подавил желание ударить его в горло.

--- Куда это ты собрался, Ликхмас ар'Люм?

--- Праздновать, --- бодро доложил я.

--- Кажется, я не давал разрешение покидать пост.
Ты забыл о долге жреца?

--- У меня клятва вместо завтрака, досуга и сна, --- сообщил я.
--- Лучше напомни той рыбине, которая поливает моими зельями сорняки, что сила не покидает микстуру с закатом солнца.

Я почувствовал, как мои ступни оторвались от каменного пола;
лицо Первого вдруг оказалось в двух пальцах от моего.

--- Ты можешь меня ненавидеть, Ликхмас ар'Люм, но ты будешь делать это тихо и покорно, --- прошелестел тихий голос, и у меня вдруг вспотели ладони.
Казалось, что жрец даже не открыл рта.
--- В моей власти отправить тебя на алтарь прямо сейчас и делать с тобой на алтаре всё, что угодно.
Да, и это тоже.
Ты меня понял?

--- Вполне, --- ответил я.
Язык чесался ответить по-другому, но иной ответ был неуместен, как крик мангуста в змеином гнезде.
Кинжал уже давно ждал команды в складках одежды;
в голове даже промелькнула пара здравых идей, где можно надёжно спрятать тело и как быстро собрать вещи.

Первый безусловно умел наводить страх, даже на таких, как Кхохо;
однако для меня страх был знакомым чувством, и я ощутил в нём некоторую фальшь.
Он был коротким, мощным и беспричинным --- словно мимо пробежал испуганный бык, обдав лёгким ветерком.
Мгновение спустя я расслабился и сосредоточился на свободной руке Первого, которая легкомысленно лежала в кармане.
Может, мне и страшно, но хозяина положения не должны заботить такие мелочи.

Едва я подумал об этом --- фальшивый страх исчез так же быстро, как и появился.

<<Да, и это тоже>>.
Простейшая уловка, иллюзия чтения мыслей, которая способна ввести в заблуждение разве что дикаря хака, и то не самого умного.
Я почувствовал, что прихожу в ярость, и сжал нож покрепче.
Да что этот рыбий хребет о себе возомнил?

\ml{$0$}
{--- У твоей поясницы --- лаковый сок, дурак, --- снова прошелестел голос.}
{``You got \laaka\ sap to your back, fool,'' the voice rustled.}
\ml{$0$}
{--- Перестань теребить нож.}
{``Stop strumming your knife.''}

\ml{$0$}
{--- У твоей поясницы --- закон сели и мой, \emph{пока ещё мой} город, --- прошипел я в ответ.}
{``You got \Seli\ law and my, \emph{still my} city to your back,'' I retorted.}
\ml{$0$}
{--- А вот в руке пусто.}
{``And nothing in your hand.''}

Первый захохотал.
Мои ступни снова коснулись пола.

--- Так значит, эта требуха не обманула.
Захожу я в город.
Спрашиваю на посту --- кто главный?
Лохматая дура с саблей, разлепив глаза, отвечает --- не особенно вежливо, --- что Ликхмас в храме.
Спрашиваю в храме, нелюдимая гора мышц говорит --- Ликхмас в школе, подождите здесь.
Спрашиваю в школе, смазливый учителишка уточняет --- Ликхмас спит, спросите Трукхвала в библиотеке.

Первый скривился.
На его худой шее билась жилка, гораздо медленнее, чем следовало.
Его совершенно не волновал мой гнев, он не чувствовал себя в опасности.

--- <<Ликхмас спит>>, <<подождите здесь>>, --- повторил он, гротескно подражая Митликху и Кхарасу.
--- Гляньте на них.
Вот она --- незримая иерархия поселений, которая смеётся над любыми указами Короля-жреца и волей захватчиков.
Не ожидал, что нынешним лидером Храма Тхитрона окажется мальчишка, но так или иначе ситуация окончательно прояснилась.
Распустил ты свой Храм, Ликхмас ар'Люм.
На постах беспорядок, с гостями разговаривать не умеют.

--- Лучшего Храма ты не найдёшь на всей Короне, --- процедил я.

\ml{$0$}
{--- Мне нет надобности \emph{искать} Храм, я способен любой мусор превратить в работоспособный механизм.}
{``I don't need to \emph{find} a Temple, any trash I got in my grasp I can turn into coherent working mechanism.}
\ml{$0$}
{В этом наше с тобой отличие.}
{That's the difference between you and me.''}

--- В моём Храме было три жреца со знаками Обителей, один из них жив до сих пор! --- выплюнул я, подойдя к нему вплотную.
--- Три жреца --- Кхатрим, Саритр, Трукхвал!
Мне неважно, какие силы стоят за тобой, чужак.
В моём присутствии ты будешь говорить с уважением об этих троих, равно как и обо всех храмовниках, имевших честь делить с ними еду.

Первый жрец смотрел на меня десять секхар не мигая.
Я понял, что окончательно перешёл черту;
одновременно пришло осознание, что ярость тоже была фальшивой.

Этот человек играл моими эмоциями, как кот с мышью.

--- Приношу свои извинения, --- угрожающе тихим голосом сказал наконец Первый.
--- Я очень рад, что проблемы Храма исходят от твоих организационных способностей, а не профессиональных умений прочих жрецов.
\ml{$0$}
{Так или иначе, мне \emph{придётся} с тобой договориться.}
{One way or another, we have to reach an agreement.}
\ml{$0$}
{Управленец из тебя неважный, но мы же \emph{равные}, не так ли?}
{You're not much of a leader, but we are \emph{equal}, aren't we?}
\ml{$0$}
{Например, я доверю тебе школу.}
{For example, I would give you the school.}
\ml{$0$}
{Что скажешь?}
{What do you think?''}

Я промолчал.
Способности собеседника меня немало впечатлили;
резко стал понятен слепой ужас импульсивной Кхохо.
К тому же вмешался прошлый опыт: если Первый жрец резко меняет тему --- значит, дело дрянь.

Что на этот счёт говорил Конфетка?
<<Если встретил врага, многократно превосходящего тебя по силе --- беги.
Если враг превосходит тебя по скорости --- прячься.
Если враг превосходит тебя по уму --- будь осторожен и не позволяй ему открыть рот.
Если убежать или спрятаться нельзя --- разыграй театральное представление, заставь врага поверить в то, чего нет.
Если враг превосходит тебя во всём --- к свиньям осторожность, бег и прятки, дерись как в последний раз, потому что спасти тебя может только случай>>.

Совсем некстати всплыл в памяти отрывок разговора с Чханэ: <<Из катакомб тайком выносили тела, а наутро...>>.
Я понял --- пора уходить.

--- Соглашайся.
Митликх ар'Митр не пожелал остаться.
\ml{$0$}
{Я даю тебе работу по твоим талантам и умениям, а ты не попадаешься мне на глаза и, главное, держишь свои таланты и умения подальше от моих дел.}
{I give you a job suitable to your skills and talents, and you stay out of sight and---more important---keep your skills and talents away from my business.''}

Жрец наклонился к моему лицу, и я почувствовал его почти неощутимое дыхание.

\ml{$0$}
{--- Ты понял мысль?}
{``Understandable?''}

\ml{$0$}
{--- Вполне.}
{``Quite.}
\ml{$0$}
{Я согласен.}
{I accept.''}

\ml{$0$}
{--- И, без сомнения, теперь ты можешь объяснить своим людям, что это решение --- результат достигнутого взаимопонимания, не так ли?}
{``And, beyond any doubt, you will explain your men that this decision is a compromise between our differing interests, won't you?''}

\ml{$0$}
{--- Я так и сделаю.}
{``I will.''}

--- Через рассвет жду тебя в храме, Ликхмас ар'Люм, --- с этими словами синяя роба, трепеща, исчезла за углом.

<<Через рассвет я буду в сельве>>, --- подумал я.
Замечательное окончание знакомства --- оба уверены, что победили.

Проходя мимо зеркала, я поправил съехавшее шёлковое полотнище, подтянул пояс и задумчиво посмотрел на заплетённые лианами домики.
За прошедшие дожди только один человек заметил, что домики и лианы на моей праздничной одежде соответствовали городам и путям обитаемой Короны.
Даже кормилица, которая смогла бы сплести чётки Сата во сне, не обратила внимания.

--- Учитель Ликхмас, зачем ты надел карту? --- шёпотом спросила Митхэ ар'Тра, встретив меня на прошлом празднике Большого Похмелья.
--- Это какой-то обряд?

Случай стал для меня уроком --- люди редко видят то, чего не ожидают увидеть.
А всего-то и нужно было --- вышить города всех народов одним цветом.
Интересно, как отнеслась бы Чханэ к идее посетить все вышитые на одежде поселения?
Кристалл, полный мира и спокойствия, вполне может подождать.

\section{[] Волосы и губы}

\textspace

Баночка сидел и гладил Кошку по волосам.
Кошка жмурилась и морщилась от удовольствия.

--- Мне так нравятся эти ваши волосы, --- сказал Баночка.
--- Жалко, что у нас их нет.

--- Так сделай пересадку кожи головы, --- пожал плечами Костёр.
--- Многие планты так делают.
Забавно, но человеческие волосы прекрасно растут на плантовых гормонах.

--- Я свои хочу!
Пересадить кожу каждый может.

--- Ты бы смотрелся странно, --- хихикнула Листик.
--- Но Небо с волосами выглядел бы ещё забавнее.

--- Мне ваши роговые наросты ни к чему, --- замахал я руками.
--- Жарко, путаются постоянно и везде застревают.
Как вы с ними вообще живёте.
Заяц приходит на чай раз в месяц и тихо сидит на диванчике, а робот-чистильщик потом выплёвывает килограмм её волос.

--- Хорошо, я буду их заплетать! --- развела руками Заяц.
--- Тебя это устроит?

--- Да я не жаловался, я просто не понимаю, как это возможно!

Снаружи зазвучала бодрая музыка.

--- О, кажется, это Безымянный, --- встрепенулась Заяц.
--- Я хочу танцевать.
Где Фонтанчик?

Заяц бросилась к иллюминатору.

--- Ах ты, он уже с кем-то!..
\ml{$0$}
{Ну погоди...}
{Well, just you wait ....''}

Я вскочил.

--- Я тоже в настроении станцевать что-то быстрое.

--- Ты беременный, --- напомнил Костёр.

--- И что, у беременных не может быть таких желаний?
И вообще, к танцам надо приучать с малых лет.
Пошли, Заяц.
Костёр, Листик, и вы.
Потанцуете вместе, вы замечательно смотритесь.

--- Правда? --- удивилась Листик и посмотрела в тёплые карие глаза Костра.
Костёр так же удивлённо смотрел в зеленоватые, лишённые белков очи Листик.
Но смотреть надо было на губы.
У обоих в уголках рта застыла одинаковая невыраженная печаль.

\section{Танец и мёд}

--- Я давно не танцевал в паре, --- шёпотом признался я.

--- Это просто, --- зашептала в ответ Чханэ.
--- Стиль Тени ещё помнишь?
Бери меня за руки и сделай <<зеркальное отражение>>.

Я послушался.

--- Молодец, теперь быстрее.
Нет, медленнее, ритм лови.
Молодчина.
Правый переход для двух сабель, змеёй.
Лис, руку не отпускай, это танец.
Пальчиками касаемся, да.
Теперь левый подхват для сабли и ножа, кроликом.
Подножку не обязательно, лучше вытяни носочек вот так, изящно и аккуратно, как будто ты хочешь тронуть ползущего по дороге жука.

--- Тебя так танцевать учили? --- поморщился я.

--- У нас в Храме по-другому и не учат.
Что танец, что спарринг --- одно и то же.
В ритм старайся, а не на скорость, и изящнее, легче двигайся, как будто плаваешь в реке.
Руки расслабь и чуть опусти, а то локтем кому-нибудь заедешь.
Так, а теперь посмотри, что делает вон та пара.
Давай, очень аккуратно...
У меня на голове сидит трасакх, и у тебя тоже.
Ты должен сбить их рукой одного за другим.

\ml{$0$}
{--- Птичку жалко.}
{``Poor birdy.''}

--- Не надо её жалеть, она воображаемая.
Да, вот так, а теперь слева...
Молодчина.
Всё, классический парный ты умеешь.
Смотри, мы даже лучше всех... почти.

Я ухмыльнулся.
Ситрис и Кхохо определённо были жемчужинами площади.
Они танцевали какой-то очень сложный южный танец.
Далеко не худенькая Кхохо вращалась волчком, обхватив ногами колено Ситриса.
Я с трудом понимал, как они двигаются в переплетении собственных рук и ног.
Оба раскраснелись от скорости и удовольствия.

Чуть поодаль неклюже топтались Кхарас и Эрликх.
Вождь явно не был настроен двигаться, и Эрликх, приникнув к нему всем телом, нарочно выбирал движения попроще.
Ликхэ кружилась вальяжно, смакуя каждое па;
её головка лежала на груди старшины кожевников.

Музыканты старались вовсю.
Дудели рожки, утробным басом стонали смычковые, нежно щебетала цитра.
Кажется, это чернявая Хатлам ар'Тро --- резчица по камню, наследница учителя Трукхвала;
кроме неё так играть на цитре могла лишь Кхохо.
Кормилец, которого природа обделила слухом, но безусловно одарила чувством ритма, виртуозно молотил трещотками в бубен.

Вскоре я уже почти не думал о том, куда ставить ноги.
Чханэ смотрела на меня, и на её лице, под лёгкой тенью венка цвела довольная улыбка.

--- Не стоит быть вместе, думаешь? --- прошептала она одними губами.

--- Наверное, я погорячился, --- улыбнулся я.

--- Впрочем, ничего нового, --- вздохнула подруга.

--- Вообще ничего, --- согласился я.
--- Кстати, я что-то вспомнил первый раз, когда ты меня обняла.

--- Это который? --- наморщила лоб девушка.

--- А который в лесу, --- ухмыльнулся я и мягко увёл Чханэ в сторону, избежав неминуемого столкновения с влюблённой парой.
--- Я аж от избытка чувств ножом махнул.

--- Помню, помню.
Кстати, я в тебя тогда и влюбилась.
Меня никто до этого так нежно не бинтовал.

--- А изображала из себя невесть что, --- упрекнул я подругу.

--- А ты думаешь, влюбиться легко?
Это всё равно что задницу в улей засунуть, проверяя, есть там пчёлы или нет.
Даже если не ужалят --- мёд всё равно со странным привкусом.

--- Трааа... --- я поперхнулся от такого сравнения.
--- Это у вас так на Западе говорят?

--- Это Кхохо так сказала.
Подозреваю, что умозаключение основано на личном опыте.

--- Меняемся! --- крикнул кормилец и, постучав в бубен, отдал его кому-то из толпы.
Кхотлам сидела в уголке и, улыбаясь, вязала очередную занавесочку;
кормилец попытался аккуратно забрать плетение из её рук и вытянуть Кхотлам на танец, за что получил лёгкий удар спицами по лбу.
Впрочем, вскоре плетение осталось сиротливо лежать на скамеечке, а кормильцы исчезли где-то в пёстрой праздничной толпе.

\ml{$0$}
{--- Ликхмас, ты не против? --- подошёл улыбающийся Ситрис.}
{``\Likchmas, would you mind?'' smiling \Sitris\ came up to us.}

\ml{$0$}
{--- Нет, конечно, --- я улыбнулся в ответ и протянул ему руку.}
{``Of course I wouldn't,'' I smiled back and gave him my hand.}

\ml{$0$}
{--- Ты мне его только живым верни, --- захихикала Чханэ.}
{``Just return him alive,'' \Chhanei\ chuckled.}
\ml{$0$}
{--- Кхохо вон уже лежит.}
{``\Kchoho's lying down over there.''}

\ml{$0$}
{--- Кхохо перебрала, --- пояснил Ситрис.}
{``\Kchoho's too much to drink,'' \Sitris\ explained.}
\ml{$0$}
{--- Под вином Кхараса тянет поговорить, Эрликха --- на сон, а Кхохо танцует всех, кто окажется на расстоянии вытянутой руки.}
{``On wine, \Kcharas\ is talky, \Oerlikch\ is sleepy, and \Kchoho\ dances everyone within arm's reach.}
Чханэ, к тебе партнёр.

К нам подошёл красавец-кожевник, который когда-то ухаживал за Чханэ, и мягко увёл её на танец.
Ликхэ, которая уже хотела пригласить меня, вяло откликнулась на приглашение Столбика и ушла с ним, бросая на Ситриса злобные взгляды.

--- Только не крути меня, а то вино обратно польётся, --- предупредил я.

--- Даже желания особого нет, --- ухмыльнулся воин.
--- Давай помедленнее и понежнее...

\asterism

Ситрис действительно был очень аккуратен.
Мне казалось, что я плыву по волнам.
Я смотрел в его чёрные, как угольки, глаза --- в них застыло что-то тёплое и непонятное.

--- Ты похож на Митхэ, --- шепнул он, поняв, что я прочитал выражение его лица.
--- Мы с ней тоже как-то танцевали, давным-давно, в Тхартхаахитре.
Я её тогда поцеловал украдкой, когда никто из отряда чести не видел...

\ml{$0$}
{--- Думаю, она влепила тебе пощёчину.}
{``I guess she slapped you.''}

\ml{$0$}
{--- Нет.}
{``No.}
\ml{$0$}
{Она поцеловала меня в ответ и сказала, что это был наш последний поцелуй вне Круга.}
{She kissed me back and then told me that it was our last kiss beyond the Circle.''}

\ml{$0$}
{--- Какой изящный выход из ситуации.}
{``What an elegant way out.''}

\ml{$0$}
{--- Кодекс Ягуара.}
{``The Jaguar Code.}
\ml{$0$}
{<<Просящему любви дай то, что в силах дать, но забери надежду...>>}
{Give what you can to the looking for love, but take away their hope ... ''}

\ml{$0$}
{--- <<...ибо тропа Ягуара --- тропа слёз>>.}
{`` ... for the path of Jaguar is the path of tears.}
Этой фразой заканчиваются все постулаты, я знаю.
Да и сами постулаты --- одно отдай, другое забери, как у торговцев на рынке.
Интересно, у Атриса она тоже забрала надежду?
А у меня?

\ml{$0$}
{--- Люди совершают ошибки.}
{``People make mistakes.}
\ml{$0$}
{Даже такие прекрасные, как она или ты.}
{Even as wonderful as she or you.''}

--- Мои ошибки пока пожиже.

--- Я никогда не видел тебя таким смущённым.
Из тебя сарказм так и лезет.

--- Естественно, я смущаюсь, видя, как ты подкатываешь ко мне!

--- Это заметно, да? --- воин виновато улыбнулся и снова блеснул антрацитовыми глазами.
--- Ни за что не позволил бы себе отнять тебя у апельсиновой --- она мне как питомица уже.
Да и ты словно мой...
И в то же время... --- щёки Ситриса неожиданно заалели.
\ml{$0$}
{--- Ничего с собой поделать не могу.}
{``Can't help that.''}

Мы помолчали.

--- Ты решил уйти, да? --- вдруг серьёзно спросил меня Ситрис.

Я снова промолчал и прижался к нему крепче.

--- Я буду скучать, Лис, --- прошептал Ситрис мне на ухо.

Я поднялся на цыпочки и поцеловал его в губы.

\ml{$0$}
{--- Это наш последний поцелуй вне Круга, --- кивнул Ситрис.}
{``It's our last kiss beyond the Circle,'' \Sitris\ nodded.}
В его глазах заблестели слёзы.
\ml{$0$}
{--- Я понял.}
{``I've got it.''}

Я кивнул и, легонько оттолкнув Ситриса, ушёл в сторону храма.

\section{Кусачка}

\ml{$0$}
{--- Ты с ума сошёл, --- констатировала Эрхэ, подойдя к Ситрису.}
{``You're going mad,'' said \Oerchoe\ to \Sitris.}
Тот продолжал смотреть вслед уходящей Митхэ.

\ml{$0$}
{--- Умеренно.}
{``Reasonably.''}

\ml{$0$}
{--- Умеренно?}
{``Reasonably?}
\ml{$0$}
{Ты вообще как-то разделяешь уважение к человеку и желанием его трахнуть?}
{Do you ever draw the line between feeling respect and desire to bed?''}

\ml{$0$}
{--- А ты?}
{``Do you?''}

Эрхэ покраснела.

--- Ладно, извини, --- неожиданно сказала она.
--- Мне плоховато, и я места себе не нахожу.
\ml{$0$}
{Общая беда?}
{Same boat?}
\ml{$0$}
{Обнимемся?}
{Cuddle?''}

\ml{$0$}
{--- Обнимемся.}
{``Cuddle.''}

Ситрис на пару секхар уткнулся носом в крепкий затылок Эрхэ и похлопал её по спине.

--- Пошли погуляем.

--- Мне нужно осторожнее гулять, --- напомнил Ситрис.

--- В таком случае я знаю, где тебя точно не будут искать, --- раздался сзади знакомый голос.

Атрис сиял.
За его руку цеплялся Кусачка.

--- Смею предположить, доблестные воины ни разу не были на празднестве Бродячего Народа.
Ну что, Кусачка, возьмём их?
Восемьдесят восемь, двадцать один?

--- Живите долго! --- застрекотал Кусачка.
--- Трое люди могут не пустить!
Трое --- к несчастью!

--- Кажется, я вовремя, --- ухмыльнулась возникшая из ниоткуда Согхо, стягивая порванную рубаху и скатывая её в рулон.
--- Гады, последнее порвали...
Живите долго, Кусачка.
Четверо сойдут?

--- Живите долго!
Четверо --- это к счастью!

--- Ну вот и пойдём, --- заключил Ситрис.
--- А идти-то сколько?
Вы же не в городе празднуете?

--- В одном городе! --- ответил травник.
--- Только не на этой земле, а под этой землёй!

--- Катакомбы, --- объяснил Атрис.
--- Частично храмовые, заброшенные, частично старые, которые под южными кварталами.
Люди туда почти не заходят, поэтому их облюбовали для своих сборищ травники.
Пойдём, Эрхэ, тебе понравится --- там вообще никакой выпивки, только всякие наркотические травы, которые на людей не действуют.

\section{Учитель}

Музыка потихоньку сходила на нет.
Все начали заниматься своими делами --- беседовать, играть и уплетать то, что ещё оставалось на столах.
Немного протрезвевшая Кхохо принесла цитру и запела что-то задумчиво-томное, по-южному надрывно растягивая слова;
собравшиеся вокруг неё девушки-крестьянки легли на землю в обнимку и зачарованно замолчали, слушая её богатый, с лёгкой хрипотцой голос.
Вскоре я узнал за непривычными аккордами знакомую с детства песню <<Портниха, что проходит мимо>>:

\begin{verse}
\ml{$0$}
{Моё сердце --- пустыня для врага и безумца,}
{My heart is a desert for a madman or foe,}\\
\ml{$0$}
{Моё сердце --- долина камней}
{My heart is a dale of stones}\\
\ml{$0$}
{Для чужестранца и сладкоречивого торговца,}
{For a golden-tongued trader or outlander,}\\
\ml{$0$}
{Моё сердце --- джунгли для соседа и прохожих,}
{My heart is the Silva for a neighbour or passer-by,}\\
\ml{$0$}
{Моё сердце --- заливные луга}
{My heart is a flood plain}\\
\ml{$0$}
{Для друзей и кормильцев и братьев.}
{For a friend, for a sibling or nurse.}\\
\ml{$0$}
{Но для тебя, моя лоза,}
{But for you, my vine,}\\
\ml{$0$}
{Но для тебя, благоухающий кувшин с маслом,}
{O\~{e}, but for you, my jug of fragrant oil,}\\
\ml{$0$}
{Моё сердце всегда как поле, вспаханное поутру.}
{My heart is a field plowed at dawn, forever.}\\
\ml{$0$}
{Приди к нему, брось в него семена ---}
{Come to them, throw your seeds in them,}\\
\ml{$0$}
{И в полдень оно расцветёт,}
{And at midday you'll see them bloom,}\\
\ml{$0$}
{И к вечеру оно даст урожай.}
{And at sundown the harvest shall come.}
\end{verse}

Трукхвал хмуро сидел на больничных носилках и пил принесённый Кхотлам горячий отвар.
Попытка вытащить его из библиотеки была похожа на представление.

\ml{$0$}
{--- Учитель, просто посиди на скамеечке.}
{``Teacher, just sit on a bench.}
\ml{$0$}
{Можешь полежать, если хочешь.}
{You may even lie down if you prefer.''}

\ml{$0$}
{--- Ликхмас, нет.}
{``\Likchmas, no.''}

\ml{$0$}
{--- Трукхвал, ты здесь помереть решил, в этом склепе? --- поинтересовался Ситрис.}
{``\Trukchual, are you gonna die in that tomb?'' \Sitris\ asked.}

--- Ситрис, иди к свиньям!

\ml{$0$}
{--- Я обещаю: тебя никто не будет трахать! --- развела руками Кхохо.}
{``I promise: no one shall fuck your good self!'' \Kchoho\ spread her arms wide.}
\ml{$0$}
{--- Даже я!}
{``Not even me!''}

--- Кхохо, и ты тоже!

--- Ликхмас, левую руку привяжи к носилкам, --- командовал Ситрис.
\ml{$0$}
{--- Вот так.}
{``Just like that.}
\ml{$0$}
{Да, для хромого старика он ого-го в хорошей форме...}
{Yeah, for a limping old man he's in a quite good shape.}
\ml{$0$}
{Так-так, осторожнее здесь, а то вы его опять башкой о косяк приложите...}
{Hey-hey, watch your step, or you will knock his head on the lintel, again ....}
\ml{$0$}
{Нет, Кхохо, кляп не нужен.}
{No, \Kchoho, we won't gag him.}
Да всё нормально, он как только за воротами очутится, сразу прекратит орать и дёргаться --- для него ж репутация превыше всего.

Спустя кхамит Трукхвал уже пересел с носилок на тростниковое кресло и весело болтал с облепившей его молодёжью.

---  ...а главное в летающих машинках --- винт.
Смотрите, если кусочек дерева вырезать вот так и покрутить, то он будет равномерно загребать воздух, подобно вёслам...

Мы встретились взглядами.
Я улыбнулся.

<<Видишь, теперь и я могу тебя кое-чему научить>>, --- знаками показал я ему.

<<Всегда>>, --- ответил мне Трукхвал и снова принялся рассказывать про винт.
Что он имел в виду --- так и осталось загадкой.

Празднество шло своим чередом, но я ни на секунду не мог забыть про нависшее надо мной будущее.
То тут, то там я натыкался на пустые, чуждые взгляды новых храмовников;
они попадались в праздничной толпе, словно камни в мягком пироге.
Учитель выглядел каким-то радостно-беззащитным ребёнком на их фоне...

Я не мог оставить его на съедение этим стервятникам.
Улучив момент, я оттащил Трукхвала от молодёжи и рассказал всё как есть.
Он не удивился.

--- А ты думал, я просто так сижу в библиотеке и никуда не выхожу?

--- Но это же не решение, учитель!

\ml{$0$}
{--- Молодые ноги быстро бегают, Ликхмас.}
{``Young feet run fast, \Likchmas.}
\ml{$0$}
{А я стар и хром.}
{I am old and crippled.''}

--- Я теперь уже не уверен, что принимаю правильное решение, спасаясь бегством.
Может, мне следовало?..

\ml{$0$}
{--- Нет, Ликхмас, силы неравны.}
{``No, \Likchmas, we're outnumbered.}
\ml{$0$}
{Верхний Этаж для нас потерян.}
{We've lost the Upstairs.}
\ml{$0$}
{Рано или поздно \emph{он} тебя убьёт, и ему ничего за это не будет.}
{Sooner or later, \emph{he} will find the way to kill you and get away with this.}
\ml{$0$}
{Тебя не защитят ни старый Храм, ни Кхотлам.}
{No Temple, and no \Kchotlam\ will be able to save you.}
И ещё... раньше я, хай, сомневался, но после твоего рассказа уверен: тебя он намерен именно убить, а не выслать.
\ml{$0$}
{Он не жрец, он дикарь, закалённый в интригах и заговорах.}
{He's not a priest, he's a savage, seasoned with conspiracy and betrayal.}
\ml{$0$}
{Он не будет рисковать, оставляя в живых опасного соперника.}
{He wouldn't risk leaving a dangerous rival alive.''}

\ml{$0$}
{--- Я его не настолько боюсь, учитель.}
{``I'm not so afraid of him, teacher.''}

\ml{$0$}
{--- А следовало бы!}
{``But you must be!}
Более того, если он убьёт тебя, то и твою кормилицу, возможно, что и твоих домашних тоже.
\ml{$0$}
{Он ведь знает, что они будут мстить.}
{He knows they'll be seeking justice.}
Кхотлам не только город, она весь Север перебаламутит, если с тобой что-то случится!
Да даже если она по каким-то причинам и решит сохранять нейтралитет, Манэ, Лимнэ, Хитрам и ваши слуги молчать не будут.
Особенно Сиртху --- ему-то вообще в этой жизни терять нечего, а ты у него на руках вырос.
Всё равно пойдут слухи, народ возьмётся за оружие.
Думаешь, люди не помнят, как ты вернулся из земель хака, как декаду метался между алтарём и школой, как ты засыпал с хлебом во рту?
\ml{$0$}
{Такое не забывают.}
{Things like that cannot be forgotten.}
\ml{$0$}
{А слухи --- они всегда страшнее правды, в твоей смерти обвинят не только жрецов, но и Кхараса, и Ситриса.}
{But rumours, they are always worse than the truth, not only priests will be charged for your death, but \Kcharas, and \Sitris\ too.}
\ml{$0$}
{Ты представляешь, что начнётся?}
{Do you realize what would have begun?}
\ml{$0$}
{Так что беги, прошу тебя, уклонись от боя в этот раз.}
{So run, please, escape the fight this time.}
Тхитрон достаточно повидал бед за эти дожди.

--- Как он сможет убить сразу меня и всех обитателей Двора, не вызвав подозрений?

\ml{$0$}
{--- Да по тому же сценарию --- кусачая бумага и диверсант!}
{``The same scenario, biting paper and a scout!''}

\ml{$0$}
{--- Ты с ума сошел?}
{``Are you bloody mad?''}

Трукхвал замолчал и насупился.

\ml{$0$}
{--- Ты... ты серьёзно?}
{``You---are you serious!}
\ml{$0$}
{Учитель!}
{Teacher!''}

\ml{$0$}
{--- Идолы здесь точно ни при чём, --- промямлил Трукхвал.}
{``Idols took no part in this, positively,'' \Trukchual\ mumbled.}
--- Я послал письмо тому караванщику.
Коробки со свинцовым блеском нам уже в Миситре отгрузили с нужным весом, и он ни на фасолину не изменился.
%{not a single bean more or less.}
Они проверяли, в учётных книгах есть все записи.
И я... я подозреваю, что идол лежал в одном из тюков связанный --- может быть, с травами, или с перцем, или среди перьев.
У нас было четыре тюка, в которых он мог лежать незамеченным и которые мы обычно принимаем на честном слове, потому что товар грошовый.
Его накачали наркотиками, и он всю дорогу...

\ml{$0$}
{--- Почему ты молчал?}
{``Why didn't you say anything?''}

\ml{$0$}
{--- А кто в это поверит? --- в глазах учителя застыли слёзы, его подбородок дрожал.}
{``Who would believe that?'' the teacher's eyes were full of tears, his jaw was trembling.}

Мне захотелось кричать.
<<Ты понимаешь, что это была решающая улика?
Ты понимаешь, что, если бы ты сразу сказал мне или Кхотлам, мы могли проследить происхождение коробок и выявить всю эту разбойничью шайку в жреческих робах, которая уничтожила наш Храм?>>
И только милосердие к учителю удержало эти слова внутри.

--- Надо устроить Дело Перекрёстка, --- пробормотал я.
--- Прямо сегодня ночью.
Вызовем их на суд.
Даже без доказательств это достаточно серьёзное обвинение, чтобы его так просто...

\ml{$0$}
{--- Нам некем их заменить, Ликхмас.}
{``We have no people to replace them, \Likchmas.}
\ml{$0$}
{Плохие они жрецы или хорошие, но других у нас сейчас просто нет.}
{Good they are or bad they are, we have no other priests.}
\ml{$0$}
{Что город будет делать без Верхнего Этажа?}
{What the city is going to do without the Upstairs?''}

Я схватил Трукхвала за плечи.

--- Напиши людям Золотой Пчелы.
\ml{$0$}
{Сегодня же.}
{Tonight.''}

\ml{$0$}
{--- И что им сказать?}
{``What should I say to them?''}

\ml{$0$}
{--- Правду, учитель!}
{``The truth, teacher!}
\ml{$0$}
{Если уж мне придётся уехать, я сам их поищу и передам твоё письмо.}
{If I have to leave, I myself will look for them and deliver your letter.''}

Трукхвал шмыгнул носом и кивнул.

--- Твоя правда.

--- И пожалуйста, выйди из Храма.
Иди к кормилице, она даст тебе какую-нибудь работу у себя.
Почти то же самое, что в библиотеке --- сиди да переписывай.
Она ведь прекрасно понимает, в чём...

\ml{$0$}
{--- Нет, ученик.}
{``No way, learner.}
\ml{$0$}
{Для меня это вопрос чести.}
{This is a matter of honor for me.}
\ml{$0$}
{Иди, обними старика и больше обо мне не вспоминай, не терзай душу.}
{Come, give that old man a hug and don't ever think about me again, please, don't beat yourself up.}
\ml{$0$}
{Я прожил жизнь хорошо и уж как-нибудь доживу оставшееся.}
{I've had a good life, and I'll somehow handle a piece left of it.''}

...Михнет спустя я сидел в кустах за чьим-то жилищем и плакал навзрыд.

\section{[] Катакомбы}

--- Как-то здесь невесело, --- констатировала Эрхэ, оглядывая сумеречные просторы катакомбной площади.

--- Я бы сказал, здесь очень весело, --- прошептал Атрис.
--- Наблюдай за ними, лови волну.
Это нужно понять.

Ситрис ни разу в жизни не видел такого количества травников зараз.
Они лежали на полу, сидели в нишах на стенах, размахивая руками и переговариваясь на своём стрекочущем языке.

\textspace

--- Если вы пришли, вы должны лечь с ними.
Все эти взаимные поглаживания, объятия, знаки доверия --- всё это очень важно.

--- Я не очень хочу, если честно, --- призналась Эрхэ.
--- Не то чтобы я относилась к травникам предвзято, просто не в настроении.
Можно как-нибудь без этого?

--- Для травников очень важен улей, --- объяснил Атрис.
--- Они гораздо более нуждаются в крепких родственных узах, нежели мы.
Но ульем они считают не то общество, в котором выросли, а то, которое их приняло.
Поэтому, как ни странно, Красные травники не считают Бродячий Народ своим ульем, но считают ульем многих ркхве-хор, с которыми постоянно контактируют.
Соответственно, для Бродячих народ сели гораздо ближе Красного колена.
Бродячие вообще в этом плане очень уязвимы, потому что они --- Колено Изгоев.
Они ревностно чтят улей.

--- У нас так же, --- пожала плечами Согхо.
--- Разве нет?

--- Не настолько, поверь, --- покачал головой менестрель.
--- Для травника найти улей --- это как родиться, уже будучи рождённым.
А потерять улей --- это как умереть, оставаясь ещё живым.

--- И ты этого хочешь, Кусачка? --- спросила Согхо.

Травник сложил руки на груди.
Ситрису показалось, что он смутился до глубины души.

--- Ты не считаешь себя частью улья? --- допытывалась Согхо.

--- Считает, --- как-то неуверенно сказал Кусачка.

<<Он всю жизнь избегал своих и общался в основном с отрядом чести>>.
Ситрис чувствительно пнул Согхо по лодыжке.
Воительница наконец поняла, что ведёт себя бестактно, и решила повести разговор в другом направлении.

--- Интересно, а я могу с ними пообщаться?
Я же человек.

--- Ложись с ними и общайся, --- улыбнулся Атрис.

--- Мне гораздо интереснее эти катакомбы, --- пробормотал Ситрис.
--- Кусачка, ты здесь часто бываешь?
Можешь показать мне окрестности?

--- Можешь! --- застрекотал травник.
--- Здесь красиво --- есть много сделанного разными старыми людьми, травниками, идолами, которые жили в одном городе!
Надпись, картинки, здание!
Мы чиним многое красивое, красим, сохраняем для тех, кто будет смотреть потом!

--- Покажи.
Обожаю древние вещи.

Кусачка схватил разбойника за руку и потащил куда-то в темноту.

\section{То, что нельзя было видеть}

Вскоре я обнаружил, что не я один сбежал от весёлого праздничного шума за дворы.
Многие успели переместиться в кусты с очевидной целью, не дожидаясь Круга.

Эрликх уже давно сидел на брошенных подушках в окружении поклонников и поклонниц;
счастливый щебет перемежался звуками поцелуев и вздохами.
Кхарас лежал в одиночестве в кустах, скрестив руки на груди.
Эрликх что-то сказал своим, несколько девушек перебежали к вождю и начали его целовать;
спустя секхар они убежали обратно, словно обжегшись.
Кхарас не двинулся с места.

Эрликх покачал головой, наконец высвободился из переплетения ласковых рук и подошёл к Кхарасу.

--- И кому ты пытаешься доставить этим удовольствие?

Молчание.

--- Если тебе хочется грустить --- грусти.
Но хотя бы делай это на виду.
Меня беспокоит то, что ты начал прятаться.
Давай, туша, вылезай из кустов и ложись на дорогу, чтобы о твою грусть все спотыкались.

Молчание.

--- Я знаю, о чём ты думаешь.
<<Сейчас он уйдёт, а она бы не ушла, она бы сидела со мной всю ночь в мокрых колючих кустах>>.
Да, она бы действительно так и сделала, я в ней не сомневаюсь.
Только, может, надо поднять задницу, когда просят, а не испытывать глубину чужих чувств?

Молчание.

--- Да пошёл ты.
Я тоже её любил.
И не смей говорить, что ты любил её больше.
Просто я и без неё чего-то стою в этой жизни.

Эрликх хмыкнул и ушёл к ожидавшей его молодёжи, украдкой утирая набежавшие слёзы.
Я едва успел скользнуть в узкий переход между жилищами до того, как он меня заметил;
меня не покидало ощущение, что я окунулся во что-то очень личное, не принадлежащее...

Не успев пройти десяти шагов, я уткнулся в скамеечку, на которой примостились две женщины.

Кормилица всегда говорила, что на празднике нужно ухаживать за грустными.
Весёлые займут себя сами.
Сейчас она сидела возле Ликхэ.

--- Ты понимаешь, я согласна на всё.
Пусть даже так.
Но нет. Я в растерянности.

Кхотлам, судя по всему, тоже пребывала в растерянности.
Она оглянулась по сторонам, и я вдруг понял, что она пытается найти меня.
Не найдя, кормилица снова посмотрела на сгорбленную, понурую девушку.

--- Он на тебя смотрит?

Я понимал, что снова становлюсь свидетелем того, что мне не следовало видеть и слышать.
Сзади были Эрликх и Кхарас со своими бедами, а здесь...
Я попытался подумать о чём-то другом, пытался посчитать про себя до ста, но голоса упрямо заползали мне в уши.

--- Не так, как на неё, --- тихо пожаловалась Ликхэ.

Кормилица гладила её по голове.

--- Это случается, златовласая моя.
Кому-то нравится одно, а кому-то другое.
Нельзя заставить...

--- Он на меня вообще не смотрит, пока я с ним не заговорю! --- всхлипнула девушка.
--- Да, заставить нельзя, я знаю.
Да, в моей жизни были мужчины, с которыми не всё просто.
Но к каждому можно найти подход, к каждому можно подобрать ключ.
А здесь...
Я просто выбилась из сил...

Кхотлам аккуратно приобняла воительницу.

--- Ты с ним говорила?

--- Да.
Он сказал... сказал... что он меня не примет, потому что я хочу владеть им безраздельно.
Ну я же не виновата, что я такая!

Вырвавшаяся из охрипшего горла фраза полоснула меня как ножом.
Кхотлам вздохнула.

--- Девочка моя, твоя жизнь чересчур круто замешана на мужчинах, --- тихо прошептала она.
--- Для меня мужчины всегда были приправой.
Да, кому-то нравятся остренькие блюда, как той же Кхохо.
Но даже в самых острых блюдах перец --- не основной ингредиент.

Ликхэ снова всхлипнула.

--- Хочешь повязать?
Я дам тебе нитки, --- мягко преложила кормилица.
--- Иногда я представляю, что вяжу небеса и звёзды, под которыми будут жить люди.
Это если ниточки синенькие.
А если...

--- Я лучше прогуляюсь, --- ответила Ликхэ.
Её голос стал сухим, словно Кхотлам оскорбила её в лучших чувствах.
--- Благодарю тебя, Пёрышко.

Кхотлам вздохнула и ушла в сторону площади.
Путь был открыт.
Но я ещё долго простоял без движения, пытаясь понять страшную вещь --- мой Храм рушился на моих глазах.

\section{Дары}

--- Я боялась, что эти перекрытия просто рухнут, пока мы танцуем.

--- Кошмар, --- поёжилась Митхэ.
--- И они там живут?

--- По крайней мере ночуют точно там, --- развела руками Эрхэ.
--- Постоялые дворы для травников некомфортны, сама понимаешь --- куча людей, высокие потолки.
До лагерей Бродячего Народа идти через весь город, а у них торговля начинается с раннего утра...

--- И Трёхэтажный Храм знает, что у них под главной площадью потенциальное кладбище Бродячего Народа?

--- Знают все.
Никто даже не почесался облагородить эти катакомбы, хотя торговля Тхартхаахитра держится на караванах травников!

--- Я поговорю с Советом Цеха, --- задумалась Митхэ.
--- Это не дело.

--- Я боюсь, что если поднимется шумиха, Храм просто закроет катакомбы, и этим всё закончится.

--- В таком случае обойдёмся без шумихи.

--- Митхэ, не надо, --- проникновенно сказала Эрхэ.
--- Два Храма в одном городе --- событие всегда значительное и далеко не всегда безоблачное, ты это знаешь как никто.
Храм привык быть хозяином, и роль гостя даётся ему с трудом.
То, что нас призвали во время нашествия хака, было необходимостью, мы все это понимаем.
Но сейчас, когда внешней угрозы нет...
Я боюсь, если отряд чести начнёт вмешиваться в политику города, Три Этажа пойдут на принцип.

Митхэ промолчала.
Вспоминая недавний разговор с Первым жрецом, она признала, что подруга права.
Повинуясь внезапному порыву, она рассказала Эрхэ о разговоре.

--- То есть всё ещё сложнее, --- Эрхэ, выслушав командира, схватилась за голову, зарывшись пальцами в короткие жёсткие волосы.
--- Слушай, я тогда вообще не знаю, что сказать.
Если он в тебя всё ещё влюблён, с ним нельзя вести переговоры.

--- Вот именно.
Поэты часто говорят, что влюблённость спасает, но не упоминают, сколько отвратительных вещей было совершено влюблёнными.
Я бы лучше доверила свою жизнь мечтателю или человеку, живущему расчётами.

--- А Атрис --- мечтатель или человек, живущий расчётами?

--- По-моему, и то и другое в очень привлекательной пропорции, --- засмеялась Митхэ.

--- Он видел, что Ситрис тебя поцеловал, и ни слова ему не сказал.

--- И не скажет.
Атрис совершенно не умеет ревновать, за это я его и люблю.

--- Знаешь, я вообще немного побаиваюсь людей, подверженных чувствам, --- призналась Эрхэ.
--- Тех, кто способен на страстную влюблённость, на дикие обиды.
Человек эмоций не принадлежит сам себе.
Сейчас он страстно целует тебя и клянётся в любви, михнет спустя раскроит тебе череп, а ещё через михнет будет плакать над твоим трупом.
Я-то уж знаю, сама такая порой.
Поэтому Снежок меня и боится...

Митхэ погладила колено подруги.
Эрхэ благодарно сжала её ладонь.

--- Но я всегда знала, где грань.
А бывают люди, которым чувства служат и принципом, и моралью, а при случае --- и оправданием.
Плохо то, что их сложно предсказать, и хуже всего, если такой человек ещё и скрывает свою бурю за каменным лицом.
Вот как этот твой ухажёр.
У него лицо, словно каменная маска, у которой в глазницах развели костры.

--- Ты права, --- Митхэ бросила взгляд на Первого жреца.
--- В бездну и его уважение, и его любовь, голова мне дороже.
Сегодня празднуем, завтра, может быть... потом идём на север, от греха подальше...
Не хочешь заглянуть в Тхитрон?

--- Кхотлам будет рада.

--- И мы будем рады.
\ml{$0$}
{Выпьем?}
{Maybe a drink?''}

\ml{$0$}
{--- За что пьём?}
{``What will we drink to?''}

\ml{$0$}
{--- За тебя, конечно.}
{``To you, of course.}
\ml{$0$}
{И за твою любовь.}
{You and your love.''}

\ml{$0$}
{--- Послушай, --- сморщилась Эрхэ, --- я, может, и выгляжу влюблённой дурой, но Акхсар --- не единственное, что интересует меня в этой жизни.}
{``Listen,'' \Oerchoe\ made a wry face, ``perhaps I look like a fool in love, but \Akchsar\ is not the only thing I'm interested in.}
\ml{$0$}
{Так что давай лучше выпьем за кое-что другое.}
{So let us drink to another thing.}
\ml{$0$}
{За дорогу.}
{To the road.''}

--- Давай, --- Митхэ схватила оказавшийся рядом кувшин и наполнила две чаши.
\ml{$0$}
{--- Почему за дорогу?}
{``Why the road?''}

\ml{$0$}
{--- Потому что каждый кхене приносит нам частицу того, что делает нас счастливыми, что заставляет разум и тело петь в унисон с природой.}
{``Because every single \kchenoel\ gives us a piece of our happiness, something making our minds and souls' voices join the chorus of the world.''}

\ml{$0$}
{--- Ты говоришь, как Атрис.}
{``You speak like \Aatris.''}

\ml{$0$}
{--- Это его слова, --- ухмыльнулась Эрхэ.}
{``His words,'' \Oerchoe\ grinned.}
\ml{$0$}
{--- И его правда.}
{``And his truth.}
Не будь Снежка --- я была бы так же счастлива, ведь источник счастья --- в моей груди.
Но Снежок есть, его принесла мне дорога.
И сегодня я буду наслаждаться её даром до самого утра.

--- Он согласен?

--- Ещё как.

\ml{$0$}
{--- За дорогу, --- улыбнулась Митхэ.}
{``To the road,'' \Mitchoe\ smiled.}

--- За дорогу, --- Эрхэ лихо ударила чашей о чашу и, поднеся вино к носу, добавила:
--- Но меня всё ещё тошнит от вина.
Что же делать?
Ведь это тоже дар?

--- Вспомни мой обычный совет перед походом.

--- <<Бери не более необходимого>>, --- улыбнулась Эрхэ.
Вино тонкой струйкой вытекло из наклонившейся чаши и исчезло между плитами пола.

\section{Явление Хри}

--- Ликхмас, --- Ситрис неожиданно подхватил меня под руку и потянул в толпу.
--- Поговорить надо.
Ко мне сейчас подошёл Эрликх и сказал...
Кхохо, да чтоб тебя идолы поимели!
Что тебе нужно?

--- Сиртрис, дай нмне конолпли, --- сказала Кхохо.
Она пошатывалась, распугивая сновавшие вокруг парочки.
Её зрачки плавали совершенно независимо друг от друга, словно черноягоды в винной бочке.

--- Ты же бросила, --- буркнул Ситрис.

--- Как рбросила, так и пондниму.
И ртрубку мою сюда.

--- К кому ты там собралась?
Чем маешься?
Ты о чём вообще?

--- Ртрубку, которую ты у меня из камрмана вытащил, ворикшка.
Не онтнекивайся, бошльше некому.

--- Ты понял, что она сказала? --- обратился ко мне Ситрис.

Я не без труда перевёл с пьяного южного на литературный сели.

--- Мы с тобой оба с Юга, но даже я тебя уже не понимаю.
Куда тебе курить, скажи?

--- Я тебе нос лсломаю, ченрная модрда.
Так понял?

--- Ты в тоны не попадаешь, на какой чердак ты сама?
Иди посиди в уголке, пока не протрезвеешь, и не позорь Храм.
Там ничем не хуже чердака.
И травы потерять не вздумай, рассвет на носу.

К моему удивлению, Кхохо послушалась и смирно ушла в уголок, по пути плюясь чем-то совершенно неразборчивым.

\ml{$0$}
{--- Ты ведь понял, что она сказала, --- ухмыльнулся я.}
{``You understood what she said,'' I smiled.}

\ml{$0$}
{--- Конечно, понял, --- вздохнул Ситрис.}
{``Of course I did,'' \Sitris\ sighed.}
\ml{$0$}
{--- Но если я это покажу --- придётся говорить на равных с плавающим в вине мозгом.}
{``But if I show it, I have to communicate as equals with a brain floating in wine.}
\ml{$0$}
{Оно мне надо?}
{Do I need that?''}

--- Мой тоже неплохо пропитался, --- сообщил я, косясь в сторону.
Новый жрец всё ещё был рядом.

--- Вино хорошее было, --- осклабился Ситрис.
--- Кстати, конопли хочешь?
Последний урожай Кхатрима.
Вот кто знал толк в выращивании.
С его сортов тянет на философию и поэзию.
Я вчера даже о природе огня задумался.
Всю ночь в очаг глядел, до сих пор синее пятно в глазу...

--- Благодарю, но нет.
Мне вкус конопли не нравится.

--- Лис, её курят не ради вкуса, --- ухмыльнулся Ситрис.

--- И всё-таки для меня важен вкус.
Крепкое пиво из агавы по той же причине не люблю.

--- Агава --- редкостная гадость, --- согласился Ситрис.
Он явно нервничал, хоть и пытался поддержать обыденный разговор.
--- Самое страшное --- Кхарас раньше очень любил это пиво.
И лез ко всем целоваться, когда выпьет...
Так, всё, Ликхмас, эта рыбина с Верхнего Этажа прошла мимо, можно поговорить насчёт Эрликха.
Он сказал, а вернее --- проговорился, что Трукхвал...

--- Если ты про расследование...

--- Я понял, тебе тоже есть что сказать, но нет, тут немного другое, --- Ситрис говорил скороговоркой, постоянно сверкая белками.
--- Возьми всё-таки коноплю.
Бери и делай вид, что всё в порядке, тут опять \emph{эти}...

Вдруг Кхохо издала дикий клич, мгновенно оборвав музыку, и в два взмаха сбросила с себя вышитый кусок ткани.
Танцующие замерли и как один повернулись к ней.

--- Да чтоб её, --- шепотом выругался Ситрис.
--- Не хочу сейчас этой дрянью дышать...
Слушай, Ликхмас, давай чуть позже, это правда важно...
После Круга возьми Чханэ и иди во Двор.
Скажи ей, чтобы она никуда не уходила, пока я не приду, понял?

Ситрис хлопнул меня по плечу и бесшумно скрылся в кустах.

--- Круг! --- рявкнула Кхохо и ещё одним взмахом отправила в костёр уже знакомый мне мешочек с травами.

--- Кхохо, о духи, ну ты хоть раз можешь подождать до рассвета?! --- тихо взмолился Кхарас.
В наступившей оглушительной тишине его услышали все.

Кхохо не ответила.
Она замерла, словно статуя, освещённая костром, словно величественный нерукотворный памятник Хри-соблазнителю.
Трубка в её зубах дымилась ровно, и так же ровно в тишине раздалась уже знакомая нота:

--- Ааа...

--- Лис, Лис, --- над моим ухом раздался жаркий шёпот Чханэ.
Я обернулся.
Глаза подруги сияли, словно звёзды.
--- Давай подышим и сбежим.
Сегодня хочу быть только с тобой.
Обещаешь?
Ты обещаешь?..

\section{[] Капкан}

\textspace

Вокруг так же шумела музыка и кричали люди, но крики людей были похожи на звуки, слышимые на грани сна и бодрствования;
Митхэ не могла разобрать ни слова.
Воительница вдруг поняла, что она находится не в зале.
И единственное живое существо стоит чересчур близко --- каменная маска с глазами, пылающими, как два костра.

<<Какого?..>>

Она вспомнила так кстати оказавшийся рядом кувшин с вином и удивилась, не обнаружив его в своей руке.

Дверь кельи захлопнулась с чавканьем, словно противный самой сути охоты капкан.

Мир завертелся.
Митхэ почувствовала, что навалившийся сверху жрец срывает с неё одежду.
Она схватила петлю удавки --- петля выскользнула из ослабевших пальцев.
Она шлёпнула по полу рукой --- но наруча со смертельным сюрпризом на ней уже не было.
Сознание медленно покидало воительницу.

--- Не сопротивляйся, Золото, --- услышала она далёкий жаркий шёпот.
--- Я не хочу делать тебе больно.

\chapter{Вечер}

\section{Катаклизм}

Похоже, только сейчас, спустя полгода этой планеты, я могу рассказать о том, что случилось на нашей родной Тси-Ди.
За окном ещё звучат отголоски фейерверков и чьи-то песни, а во мне теплится новая жизнь --- потрясающе наивная, беспечная жизнь, пламя которой до поры свито в нежный клубочек плоти.
Признаюсь, так мне гораздо проще.

Наверное, следовало бы начать с причин, с предпосылок, со странных разговоров, которые мы с Комаром вели в вечерней тишине.
Но нет, это всё не то.

\ml{$0$}
{Я как-то спросил Кошку, как для неё начался Катаклизм.}
{I asked Cat once, how it felt, the beginning of the Cataclysm.}
\ml{$0$}
{Кошка поёжилась и сказала всего одно слово:}
{Cat shivered and said one word:}

\ml{$0$}
{--- Тьма.}
{``Dark.''}

\asterism

\ml{$0$}
{Двенадцатый город погрузился по тьму.}
{City Twelve went dark.}
Я только закончил с работой и рассеянно думал, что сегодня поеду с Заяц к ней домой --- поиграть в игры и поговорить, лёжа в обнимку в тёплой кроватке, поглощая печенье и что-нибудь согревающее.
Простая радость жизни, которую обязано попрактиковать любое мыслящее существо, особенно если за окном зима, к которой худо-бедно приспособлены обезьянки и совсем никак --- твой некогда пойкилотермный вид.

\ml{$0$}
{Просто представьте --- предвкушаете тёплую кроватку, горячий шоколад, а тут --- тьма и взвывшие сирены.}
{Just imagine: you're dreaming of warm bed and hot chocolate, and now this---darkness and wailing sirens.}

Я был готов к перебоям с энергией --- такое мне случалось видеть раза два на своём веку.
Я был готов к сиренам --- о постоянной угрозе вторжения хоргетов предупреждали всех.
Но я не был готов к тому, что сирены умолкнут спустя четыре секунды.
Они просто не могли этого сделать.
Это означало одно --- сбой самых стабильных систем города.

Резервные каналы, квантовое шифрование.
Я лихорадочно перебирал идентификаторы друзей, сослуживцев...
Комар откликнулся почти сразу:

--- Небо, молчи и слушай.
Баррикадируйте узел связи согласно протоколу номер восемьдесят.
Отключить все роботизированные системы и нейросети, управление только императивное.
Отключить все мозговые импланты.
Приготовить средства химической защиты, персональные щиты.
Отведи всех, кого можешь, на этажи с первого по третий, но ни в коем случае не ниже.
Подтверждение личности только фактическое.
Подтверждаю свою личность --- у тебя депигментированное пятно на правой ладони в виде крокодильчика, ты разговариваешь во сне и лягаешься.

Я похолодел, словно зима пробралась сквозь толстую стену и защитные поля окон.
Комар работал в особом отделе, касающемся защиты от хоргетов.

Впрочем, вскоре мне удалось успокоить себя, что такие инструкции он бы дал в любом другом непонятном случае.
Всегда лучше перестраховаться.
Я передал сообщение Комара всем, кому смог.
В здании узла связи ещё оставалось пятьдесят тси, включая меня, Заяц, Фонтанчика и Баночку.
Мы оперативно закрыли окна и двери, опустили защитные поля, перекрыли дельфиньи лазейки и вентиляцию;
здание было полностью забаррикадировано за минуту до того, как на город обрушился химический снег.

--- Синильная кислота, --- доложил кто-то из техников, вглядываясь в изображения с наружных камер наблюдения.
--- Судя по направлению ветра, выброс с завода удобрений.

Город молчал.
Не было ни криков, ни суеты.
Мы немного успокоились.
Всякое бывает в жизни.

Но Заяц не унималась:

\ml{$0$}
{--- Небо, давай ещё раз.}
{``Sky, let me get this straight.}
\ml{$0$}
{Зачем Комар велел отключить роботов и даже персональные импланты?}
{Why Mosquito wants to have robotics shut down, and personal implants too?}
\ml{$0$}
{Почему сирены молчат?}
{Why silence instead of sirens?}
\ml{$0$}
{И почему, чтоб вас всех, мы не можем ни с кем связаться?}
{And why, for all our sakes, we can't connect to anyone?''}

--- Мы заметили нестандартное поведение системы, которое быстро распространялось по подсистемам, --- сказал Фонтанчик.
--- Такие права доступа были только у Машины.
Неужели ошибка в Машине...
Но тогда...

Фонтанчик и Заяц переглянулись тем самым взглядом полного взаимопонимания, который меня всегда восхищал.
В следующую секунду они ринулись к компьютерам.

\ml{$0$}
{Да, именно поэтому мы выжили.}
{Yes, that's why we survived.}
\ml{$0$}
{Мы не просто начали действовать вовремя --- мы раньше всех сообразили, кто во всём виноват.}
{We started to act in time; more important, we found out who was behind it, before others.}

Заяц, Фонтанчик и ещё двадцать тси на свой страх и риск сделали то, чего не делал никто.
\ml{$0$}
{Мы нанесли ответный удар.}
{We struck back.}
За неполные сорок минут мы ликвидировали все роботизированные системы Двенадцатого города, до которых смогли добраться из узла связи, и переключили их на императивное управление.
И тут каналы связи взорвались.
Всё это время прочие выжившие, забаррикадировавшись в домах, пытались связаться с друзьями.
Машина перенаправляла и глушила всё, что могла.

Ещё час ушёл на то, чтобы создать новый протокол связи и поднять в городе децентрализованную сеть.
К моему стыду, в процессе я с сожалением вспоминал о постельке и горячем шоколаде.
Сложно сразу свыкнуться с мыслью, что ты в смертельной опасности, особенно если в этот момент тебе ничего, собственно, не угрожает.
Снег из синильной кислоты мало отличается от обычного, если ты смотришь на него через мощное защитное поле.

Мы тщательно задокументировали все известные ячейки выживших --- во избежание взлома со стороны.
Четырёхмиллионный город вымер в одно мгновение --- на связь вышли сорок тысяч.
Это был шок.

Баночка смог получить кратковременный доступ к спутникам и принёс нам информацию, которая превратила шок в ужас.
\ml{$0$}
{Крупные города обеих планет были мертвы.}
{Cities of both planets were dead.}
Их окутали облака синильной кислоты, радиоактивных изотопов и прочих материалов, в один миг превратившихся в химическое оружие.

Я не стану описывать последующие сутки --- горе, ужас и растерянность.
\ml{$0$}
{У птицы всегда есть надежда на небо, у волка --- на лес, у рыбы --- на глубину вод.}
{A bird always looks to the sky, a wolf to the forest, a fish to the depth.}
\ml{$0$}
{Нам бежать было некуда.}
{For us, there was no escape.}
Никакое описание не подготовит вас к подобному и не передаст всю глубину этих чувств.
Но именно тогда я понял, для чего родился на этот свет --- чтобы сказать одно-единственное нужное слово: <<Спокойствие>>.

Жить дальше захотели не все.
В течение первых десяти суток мы потеряли ещё триста восемнадцать тси.
Многие из них свели счёты с жизнью, не вынеся потери близких.
Кое-кто погиб, небрежно обращаясь с системами жизнеобеспечения и оборудованием.
Но особенно потрясли всех два случая массового самоубийства.
Сначала перестал выходить на связь детский сад, воспитатели которого решили прекратить свои страдания и страдания вверенных им детей.
Через сутки после этого события запертый в одиночестве пожилой реаниматолог, не сумев починить аварийный генератор, отключил от систем жизнеобеспечения пациентов и сделал себе смертельную инъекцию.
Один из пациентов чудом выжил и пришёл в себя;
его забрала позже одна из групп.

Мы собрали конференцию из тси, командовавших ячейками выживших.
Комар стал негласным лидером --- он был больше всех осведомлён о происходящем.

--- Машина хочет нас уничтожить, --- сказал он.
--- Имела место ошибка в одной из базовых нейронных сетей.
Бункеры выведены из строя, все пытавшиеся скрыться в них погибли.

Поступило множество предложений.
Большая часть тси склонялась к мысли, что следует уйти под землю, в защищённые тоннели метро и заводов.

--- Из подземелий борьбу можно вести вечно, --- сказал я.
--- Если Машина обрела самосознание, мы должны показать, что мы друзья, а не враги.
Ещё не поздно это сделать.
Попробуйте связаться с ней и выиграть время.

Мы сидели за столом и спорили.
Конференция трещала от наплыва кричащих, воющих и плачущих голосов.
А на чистой полированной поверхности стола всё это время лежало решение.
Одна-единственная бумажка, которая спасла жизни одиннадцати тысяч тси.

Билет в Красный музей на двоих.

\section{Ученик жреца}

Митхэ открыла глаза и снова закрыла.
Мир вертелся и рушился.
Болели плечи, которые в клочья искусал наслаждавшийся властью жрец.
На груди и животе не было живого места от крохотных надрезов и уколов.
Ужасно хотелось пить и справить нужду.

Митхэ нащупала лежавшее рядом тело.
Оно уже успело остыть;
кровь свернулась, и Митхэ почувствовала приступ тошноты от мысли, что ей придётся утолять жажду этой отвратительной жижей.
Но вот в поле её туманного зрения попал стеклянный кувшин.
Он опрокинулся и закатился в угол, но в нём ещё оставалась вода.
Митхэ схватила кувшин, осушила его и закашлялась.
Затем, поняв, что уже не может приказывать кишечнику, сползла с лежанки и лёжа справила нужду в углу.

Митхэ знала, что Нетрукха скоро хватятся.
Она знала, что ей грозит алтарь за убийство жреца.
Однако в тот момент смерть казалась не такой страшной вещью, как отвратительное чувство запачканности.
Воительница была готова с боем прорываться к купели и убить за горсть мыльного раствора, чистые штаны и рубаху.

Митхэ напялила содранную с мёртвого насильника одежду и осторожно открыла дверь.
Слабый свет из коридора осветил келью;
келья была вся в крови.
Брызги засохли на стенах, на потолке, а пол и вовсе выглядел, словно алтарь после жертвоприношения.
Митхэ случалось терять над собой контроль в пылу сражения;
однако она не могла припомнить, чтобы убивала кого-то с такой жестокостью в полном беспамятстве.

И в полном молчании.
Митхэ знала --- если бы прозвучал хоть клич, хоть крик боли, её люди и храмовые воители приползли бы сюда и пьяными, и смертельно ранеными.

Женщине казалось, что её босые мокрые ступни, прилипая к камню, издают жуткий шум.
Штрокх, штрокх.
Кажется, за окнами царила глубокая ночь.
Оставленные без присмотра факелы уже догорали;
из зала доносилось неуверенное веселье на грани мертвецкого опьянения.

--- Золото! --- сияющая Эрхэ вынырнула словно ниоткуда и обняла воительницу.
--- Я уж думала, вы с Хатом нас бросили.
Вы что, всё это время в в подвале миловались?
Хаяй, нам бы с Акхсаром такую силу, нас всего на восемь заходов хватило, дальше уже просто болтали о всякой...
Да что?..

Митхэ бессильно обвисла на руках подруги;
обоняния Эрхэ достиг острый запах крови, трупов, немытого тела и нечистот.
Вопрос застыл у воительницы на губах.
Она подхватила командира на плечо, тихо вынула кинжал и крадучись бросилась прочь, держась в тени стен.

--- Митхэ, Митхэ!

Воительница испуганно схватила губами воздух.
Кажется, она снова потеряла сознание.
Тихо догорала свеча, давая смутное представление о размерах кельи;
Эрхэ, Акхсар и Аурвелий смотрели на командира во все глаза.

--- Митхэ, скажи, кто это сделал, --- тихо прошептал Акхсар.
--- Просто назови имя, и ему не видать ни жизни, ни посмертия.
Это Атрис?
Кто-то из местных?
Митхэ?

До сознания Митхэ донеслось только одно слово.

--- Атрис, --- голос донёсся словно из-под земли.
--- Атрис...

Акхсар сорвался с места, но Аурвелий схватил его за шиворот.

--- Ты думвать, что это её лвюбовь?

--- Она сама сказала, и ты её слышал! --- прорычал Акхсар и угрожающе взмахнул фалангой.
--- Отпусти меня, трухлявая головня!

--- Нет, --- выдохнула Митхэ.
--- Атрис...
Атрис...

Акхсар застыл.
На его лице проступило что-то непонятное.

--- А твеперь слушай првиказ от трухлявый уголь, blanclicaro\FM, --- дружелюбно проворчал Аурвелий, слегка придушив Акхсара воротом рубахи.
\FA{
<<Белорубашечник>> --- презрительное название сели у ноа.
}
Митхэ показалось, что Акхсар переживает унижение чересчур спокойно, пока она не увидела приставленный к спине воина огнистый клинок.
--- Перведай отрвяд, они немведленно поиск Атриций, он бвольшая беда.
Плохой челвовек навестить Аурвелий Амвросий, когда узнать имвя и хасетрасем.
Ввыполнять.

Хватка ослабла, и огнистый клинок исчез в ножнах.

--- Так точно, --- угрюмо сказал Акхсар и выскользнул за дверь.
Послышался звук столкновения и тихая ругань;
спустя секхар Акхсар втащил в комнату чьё-то упирающееся долговязое тело и швырнул в угол.

Митхэ, собрав всю волю в кулак, приподняла голову.
Это был Трукхвал.
Молодой послушник благоразумно лёг ничком, закрыв худыми руками голову.

--- Оставьте его, --- махнула рукой Митхэ и попыталась сесть, опираясь на Эрхэ.
--- Мальчик мой, не бойся, --- воительница постаралась вложить в охрипший слабый голос как можно больше теплоты.
--- Ты хотел меня видеть?

Трукхвал, бросив взгляд на Акхсара, подполз ближе и приник к её уху.

--- Они отдали его Си-Абву, --- прошептал Трукхвал.
--- Митхэ, Атриса увезли хака три дня назад.
Он теперь их раб...

Митхэ окаменела.

--- Митхэ, я пытался их остановить, я пытался тебя найти...

Трукхвал вдруг захныкал и сел на пол, размазывая грязь рукавом по острому личику.
Длинные, чёрные как смоль волосы совершенно скрыли его.
Митхэ, опомнившись, схватила послушника за руки.

--- Трукхвал, --- тихо сказала она, --- милый мой мальчик, я тебе благодарна уже за то, что ты сделал.
Ничего другого ты бы сделать не смог.
Беги обратно, пока жрецы тебя не хватились.

--- Я ведь сказал, что хочу с тобой... --- прорыдал Трукхвал.

--- Трукхвал, тебе нужно учиться, --- горячо зашептала Митхэ, чувствуя, что силы оставляют её.
--- Ты будешь лучшим из жрецов --- в тебе есть чувство справедливости.
Не нужно никому доказывать свою силу.
Залезай на дерево акхкатрас, как взойдёт солнце, обегай городские стены в полдень, на закате переплывай Ху'тресоааса.
А я позабочусь, чтобы ты и твои потомки не увидели ни одной войны.
Иди.

Едва бесшумно рыдающий Трукхвал вышел за дверь, Митхэ снова потеряла сознание.

\section{Отступление}

Красный музей.
Самая известная достопримечательность Двенадцатого города.
Здесь было почти всё, когда-либо созданное известными нам сапиентами: крохотные модели зданий, орудия труда, утварь, механизмы, одежда и украшения.
Чтобы обойти весь музей, требовалось ровно четырнадцать дней.
Многие умудрялись сделать это за один раз;
специально для них внутри музея были пищеблоки, комнаты отдыха и развлечений.

Однажды часть музея уже была уничтожена техногенной катастрофой.
Предки, справедливо решив не повторять историю, оборудовали комплекс самыми совершенными защитными системами.
Здесь даже имелся дополнительный контур планетарной защиты, который собрали местные школьники, --- не идеальный в плане конструкции, но вполне работоспособный.
Я не мог взять в толк, с чего демонам понадобилось бы залезать в музей;
видимо, студентам просто разрешили попрактиковаться.

Главным экспонатом Красного музея был космолёт.
Настоящий огромный космолёт, не просто построенный для путешествия между звёздами, а совершивший это путешествие от взлёта до посадки.
У него любопытная история, и я остановлюсь, чтобы рассказать о нём подробнее.

Дело в том, что строили этот космолёт не тси и даже не сапиенты Древней Земли.
Когда первые поселенцы достигли Тси-Ди, он уже был здесь.
По обрывкам информации от захваченных нами хоргетов Ада и Картеля удалось узнать, что, возможно, этот корабль --- наследие Ветвей Ночи.

Кто такие Ветви Ночи, увы, неизвестно.
Следы их деятельности были обнаружены на восьми планетах, в настоящее время находящихся под властью Картеля.
Есть мнение, что эта цивилизация была случайно уничтожена хоргетом-демиургом: во время изменения климата планеты произошло значительное повышение температуры, и высокоразвитые существа просто испарились, оставив после себя несколько десятков весьма своеобразных сооружений и этот космический корабль.
Большая часть зданий была перестроена апидами Ди и впоследствии уничтожена, кораблю же посчастливилось --- он был найден задолго после Войны Тараканов.

Когда учёные изучили космолёт, они пришли в восторг от простоты и изящества его устройства.
Вскоре появились энтузиасты, которые решили реконструировать его и приспособить для нужд тси.
Так появилось целое поколение космического и межпланетного транспорта --- поколение Ночи, по счёту шестое, по парадигме --- <<спичечное>>.
9.345.9814.3 корабль, названный <<Стальным Драконом>> (в честь корабля из фантастического романа Свет-Мерцающего-Осколка), совершил пробный полёт на субсветовой скорости в ближайшее межзвёздное пространство.

Именно к нему лежала наша дорога.
<<Стальной Дракон>> был ближайшим к нам космолётом и, по мнению многих моих товарищей, наиболее безопасным --- корабль был практически лишён автоматики, а его системы не пинговались из базовой сети.

Как билет оказался на столе --- так и осталось неразгаданной тайной.
Но решение было принято единогласно --- мы отступаем в музей.

Отступление было продумано до мелочей.
Мы исчертили карту города маршрутами, которыми могли двигаться выжившие.
Мы предусмотрели все возможные препятствия и ловушки, которые Машина, несомненно, приготовила для тси.
И всё же мы не были готовы к тому, что ждало снаружи.

Первыми к музею двинулись дельфины --- они были настроены решительнее всех, да и возможностей двигаться незамеченными у них было больше.
Дельфинам не нужны были защитные костюмы --- водные дороги были большей частью под навесом по зимнему времени, плюс из-за течения и систем очистки концентрация синильной кислоты в водных дорогах оказалась существенно ниже ПДК.
Тем не менее все они получили кислородные баллоны, так как воздух всё ещё был непригоден для дыхания.
Кроме того, из соображений безопасности мы велели им двигаться у самого дна водных дорог, в облаке мути.

\ml{$0$}
{<<Поплавки, доложите обстановку, приём>>, --- скалаз Комар.}
{``Floats, give me status, over,'' Mosquito said.}

\ml{$0$}
{<<Паршиво, Альма-матер, --- отозвался Пузырь, командир дельфиньего отряда.}
{``Crab piss, Alma mater,'' said Bubble, commander of delphina squad.}
\ml{$0$}
{--- Прошли Анемоновый бульвар.}
{``Just passed Anemone avenue.}
\ml{$0$}
{Смердит кладбищем на всём протяжении, аж жить не хочется.}
{Stink of death all along, it takes away my taste for life.}
\ml{$0$}
{Приём>>.}
{Over.''}

Причина выяснилась скоро.
Водные магистрали были перегружены трупами --- дельфины умирали, их тела сносило к развязкам, и там они и оставались гнить, делая воду непригодной для передвижения и распространяя зловоние на многие километры.

\ml{$0$}
{<<Поплавки, доложите обстановку, приём>>, --- повторил Комар.}
{``Floats, give me status, over,'' Mosquito repeated.}

\ml{$0$}
{<<Альма-матер, Поплавки вышли к Вратам, --- отрапортовал Пузырь.}
{``Alma mater, Floats are at the Gates,'' Bubble reported.}
--- Разведка сообщает: возможные пути отступления два, девять, шестнадцать, сорок непригодны для эвакуации, там движется что-то недружелюбное, конкретнее сказать не могу.
Насчёт путей десять и двадцать один есть сомнения.
Приём>>.

<<Приняли, Поплавки, благодарим за предупреждение, --- ответил Комар.
--- Удачно добрались?
Приём>>.

\ml{$0$}
{<<Отряд потерял десятерых --- судя по всему, запутались в мусоре и задохнулись.}
{``We've lost ten people, they apparently got caught up in debris and choked to death.}
\ml{$0$}
{Из менее неприятных новостей --- водный лифт у Врат не работает, и починить его отсюда никак.}
{As for the less bad news, a water elevator of the Gates is out of order, and we can't fix it from here.}
\ml{$0$}
{Надеемся на вас, Альма-матер.}
{Relying on you, Alma mater.}
\ml{$0$}
{Приём>>.}
{Over.''}

\ml{$0$}
{<<Поплавки, это Узел, --- вмешался я.}
{``Floats, it's Node,'' I told.}
\ml{$0$}
{--- Мы скоро выходим.}
{``We're coming soon.}
\ml{$0$}
{Держитесь и не паникуйте.}
{Hang in there, don't panic.}
\ml{$0$}
{Не закрывайте водный шлюз до нашего прибытия, по крайней мере пока есть кислород.}
{Don't lock the water gateway until we're on site, or at least while you have enough oxygen.}
\ml{$0$}
{К вам могут приплыть ещё дельфины, а движение шлюзовых ворот только привлечёт ненужное внимание.}
{Other delfina may come, and gateway activity will just attract unwanted attention.}
\ml{$0$}
{Поплавки, как поняли?>>}
{Floats, do you copy?''}

\ml{$0$}
{<<Вас понял, Узел, --- сказал Пузырь.}
{``Copy that, Node, wilco,'' Bubble answered.}
\ml{$0$}
{--- Небо, это ты сейчас на Узле?}
{``Sky, is that you Node now?}
\ml{$0$}
{У меня тут соображения насчёт ненужного внимания.}
{I have some ideas about unwanted attention.}
\ml{$0$}
{Ощущение, что горы трупов на развязках не просто так.}
{There has to be a good reason for all these corpses at road junctions.}
\ml{$0$}
{Зимний навес сломан в нескольких местах, и я не припомню, чтобы так было до атаки.}
{Winter canopy is broken in several places, and I don't remember that before the attack.}
\ml{$0$}
{Приём>>.}
{Over.''}

\ml{$0$}
{<<Пузырь, объясни.}
{``Bubble, explain.}
\ml{$0$}
{Приём>>.}
{Over.''}

\ml{$0$}
{<<Узел, объясняю.}
{``Explaining, Node.}
\ml{$0$}
{Кто-то специально перегораживал водные дороги подручными средствами.}
{Someone intentionally blocked roads with everything they could use.}
\ml{$0$}
{Сети, решётки, мусор.}
{Nets, grills, debris.}
И этот мусор не только строительный, в смысле, не только от разрушенных частей канала.
\ml{$0$}
{Ощущение, что все мусоросборные машины разом решили опорожниться в водные дороги.}
{I got a sense that all garbage trucks jointly decided to empty into water roads.}
\ml{$0$}
{Речная и морская магистрали закрыты мусором наглухо, лосось не проскочит.}
{River and Sea highways are sealed shut with debris, even a salmon couldn't swim through.}
\ml{$0$}
{Другими словами, для нас из города дороги нет.}
{To be more specific, for us there's no way out of the city.}
\ml{$0$}
{Приём>>.}
{Over.''}

--- Кто их перегораживал? --- в ужасе переспросил Мак.

Мы переглянулись.
Никто не осмелился дать на вопрос очевидный ответ.

Синильный <<циклон>> ушёл на север.
Костюмов химзащиты хватало на всех;
неиспользованные мы взяли с собой --- для тех, кому повезло меньше.
Первая такая группа обнаружилась в жилище, ставшем у нас на пути;
в маленькое помещение набилось столько тси, что они были вынуждены спать по очереди.
Многие были ослаблены голодом и нехваткой питьевой воды.

Обманчивая тишина во время химической атаки вскоре обрела лицо.
Скрюченные, искажённые мукой тела стали попадаться даже на дороге.
Их никто не ел --- птицы и крысы также погибли.
Поначалу мы обходили все жилища, но вскоре стали ограничиваться показаниями домовых датчиков --- без средств защиты не выживал никто.

Спустя час к нам присоединилась группа Комара, забаррикадировавшаяся в университете.
Он кивнул мне и похлопал по плечу:

--- Хорошо выглядишь.

Мне ужасно хотелось броситься к нему и спрятать голову у него на груди.
Но я подумал --- многие потеряли близких, и такие жесты с моей стороны только озлобят народ и подорвут их хрупкую волю к жизни.
Комар думал так же.

Но мы ошибались.
Удивительна всё же гибкость молодых.
Когда к нам присоединились студенты, тревога группы значительно ослабла.

--- Ну что, ребята, Катаклизм? --- весело обратился к нам женственный канин с гривой, заплетённой в множество косичек.
\ml{$0-[qi]$}
{Так крах цивилизации получил имя.}
{So the collapse of our civilization got a name.}

\ml{$0-[qi]$}
{И, надо заметить, весьма удачное.}
{And a good one, I should mention.}
\ml{$0-[qi]$}
{Катаклизм --- это очень серьёзно.}
{Cataclysm is an extremely serious thing.}
\ml{$0-[qi]$}
{Но не конец.}
{But not the end.}

--- Я читал в одной книжке, что в качестве обменной валюты можно использовать уникальные пластиковые печати для стабитаниумовых поверхностей! --- оживлённо подхватил какой-то человек, поправляя очки дополненной реальности, сползшие под маской.
--- Я уже достал целую кучу!

--- А у меня есть саундтрек, идеально подходящий к ядерной войне! --- прокричал татуированный по самые пятки апид.
--- Я его включу?

--- Зацените, чуваки.
Я, пока в подвале сидела, собрала настоящее оружие! --- похвасталась атлетического сложения девушка-человек, показывая внушительного вида агрегат.
--- Пусть роботы только попробуют сунуться!

--- Эй! --- возмутилась Заяц, подойдя к девушке.
--- Это не игрушки!
Оружие с таким же успехом можно применить и против тси!

--- Расслабься, милаха, ты под моей защитой, --- осклабилась девушка, похлопала Заяц по маске и отошла, оставив подругу кипеть в возмущении.

Мы не дождались одной из групп.
Кукловод-Бросающий-Камни ушёл на разведку и вернулся с новостью, заставившей нас перейти на бег.
Целостность убежища была нарушена роботизированным механизмом, известным как <<робот-техник>>.
Не выжил никто --- спасшихся от синильной кислоты добили плазменным резаком.

Комар велел всем активировать персональные поля и выделил несколько десятков тси для обороны отступающих, вооружив их теми же резаками.
Командовать обороной назначил Лиану.

--- Ты из особого отдела? --- восхищённо спросила девушка с агрегатом Лиану.

--- Это что? --- спросила Лиана, показывая на агрегат.

--- Метательное оружие.
Стреляет снарядами, сделанными из обычных болтов на двенадцать.
Я подумала --- пригодится.

--- Ну-ка покажи.

Девушка вскинула агрегат и прошила очередью ближайший столб.
Лиана одобрительно крякнула.

--- Не боишься?

--- Разумеется, нет!

--- Батарея?
Болты?

--- Тип 1-09-B, три четверти заряда.
Болтов полная сумка.

--- Иди в начало колонны.
И сделай мощность меньше процентов на шестьдесят-восемьдесят, а то тебя из-за отдачи аж качает.

Девушка козырнула и бегом отправилась на доверенную ей позицию.

И ещё одно опустевшее убежище --- на этот раз без признаков взлома.
Большая часть была убита большой дозой наркотического препарата.
Комар велел разведчику не распространяться о причине смерти.

--- Помнишь, я сказал отключить мозговые импланты? --- тихо сказал он мне.
Меня передёрнуло.
--- Пока молчи, Небо.
Эта новость уж точно подождёт до музея...

Из сорока тысяч выживших до музея добралось лишь пятнадцать.
Я до сих пор слышу крики тех, кому не удалось преодолеть опустевший город.
В тот вечер я торжественно дал клятву --- отныне Машине дорого обойдётся даже одна жизнь.

\section{Тайное становится явным}

\epigraph
{Притворяйся, когда зло требует подчинения!
Молись, когда зло бесчинствует!
Тебе всё равно не скрыться, когда зло назовёт тебя по имени.}
{Мариам Кивихеулу}

--- Лис!

Отчаянный возглас Чханэ настиг меня в зале.
Я даже не подозревал, что она способна на такой звук.
В ушах ещё звенела музыка празднества, в голове шумело вино;
я улыбался, вспоминая танец, горящие глаза подруги.
Её крик прозвучал как-то неправильно, уродливо, не к месту.
Однако рука уже сама нащупала нож, готовясь встретить в полутьме неведомого врага...

--- Остынь, Ликхмас.
Руку с пояса, я сказал!

Голос принадлежал Кхарасу.
В словах кузнечным молотом гремела власть, круто замешанная на фанатичном следовании долгу.
Вождь никогда не позволял себе такого тона даже с пленниками, не говоря о подчинённых.
Я вгляделся ослеплёнными глазами в тьму зала: он и Ликхэ держали избитую, взъерошенную Чханэ за связанные руки.

--- Я напоминаю тебе, Митрам ар’Сар, что вначале ты должен представить её перед Советами! --- рявкнула кормилица.
Все домашние тоже были здесь.

Я подошёл к Хитраму:

--- Что произошло?

--- Мы нашли у Манэ ар’Люм обсидиановый кинжал, запачканный кровью, с монограммой пропавшего без вести Сатхира ар’Со, --- сухим безжизненным голосом ответил жрец Митрам, словно читая вслух неинтересную книгу.
--- Она сказала, что эта вещь принадлежит Тханэ ар’Катхар, которая появилась в Тхитроне два дождя назад... хай... при странных обстоятельствах...

<<Лесные духи, Чханэ, почему ты от него не избавилась?!>> --- мысленно закричал я.
Чханэ, опустив голову, закрыла полные слёз глаза.
Хмель и очарование празднества испарились без следа.

--- Ликхэ, Кхарас, ну вы-то куда? --- воскликнула Кхотлам.

Воины промолчали.
Ликхэ спокойно выдержала мой укоряющий взгляд.

--- Мои люди исполняют приказ Первого жреца.
И я ещё раз напоминаю Кхотлам ар’Люм, что упомянутая Тханэ ар’Катхар --- воин, дело внутрихрамовое и черноруких\FM\ не касается.
\FA{
Общее название для Сада и Цеха.
На Юге является часто употребляемым нейтральным термином, на Западе несёт уничижительный либо оскорбительный оттенок.
}
Увести.

Манэ и Лимнэ, сжав худые руки в кулачки и нахмурив одинаковые детские личики, преградили воинам путь на улицу.
Я положил руку на рукоять ножа, кормилец словно невзначай отклонил тело вправо, чуть ближе к висящей на стене фаланге.
Эрхэ легко покачивала метлу в руках, готовясь отсоединить прутья от древка;
гибкое тело служанки, словно натянутый лук, ждало команды к бою.
Сиртху-лехэ сделал шаг вперёд и гневно открыл рот, явно собираясь отвлечь гостей и принять первый удар на себя.
Но Кхотлам едва слышно чмокнула губами --- и домашние замерли.

--- Скребок, --- кормилица намеренно назвала жреца домашним именем, что для нас означало <<Не сейчас>>.
--- Я не сомневаюсь в честности служителей Храма.
Но если ты не желаешь созвать Советы, чтобы решить судьбу Тханэ ар’Катхар, это сделаю я --- в противном случае вы рискуете совершить ошибку.

Ликхэ и Кхарас едва слышно выдохнули и убрали руки с оружия;
по лбу девушки катился пот, я впервые в жизни увидел воительницу такой испуганной.
Кхарас смотрел исподлобья, словно загнанный в ловушку зверь.
Жрец масляно улыбнулся.

--- Разумеется, Кхотлам.
А я сделаю именно так, как решит Совет Храма.
Кому, как не жрецу, следить за соблюдением законов сели?
И знаешь, следить за соблюдением законов было бы проще, если бы те, кто помогает кутрапам, разделяли их участь.
Те же тенку, как ты знаешь, за помощь преступникам подвешивают за шею.
Почему бы не перенять этот полезный опыт?

--- Сели никогда не примут этот отвратительный закон, --- прошипела Кхотлам.

--- Нет вещей, которые невозможно принять.
Есть дипломаты, не умеющие убеждать.

Кормилица окаменела.
Митрам насмешливо поклонился ей, кивнул воинам и, подтолкнув связанную, подавленную Чханэ, вышел за порог.
Вслед ему прилетело сдавленное проклятие от Эрхэ --- женщина в ярости грохнула метлу о стену и ушла в свою комнату, хлопнув дверью.

Манэ бросилась подбирать прутья, но кормилица махнула ей рукой:

--- Оставь, дитя.
Не до этого.
Ликхэ, ты здесь?

--- Да, --- раздался угрюмый голос второй служанки, и она выпрыгнула из задрапированного вентиляционного люка с духовым ружьём в руках.
--- Пёрышко, зря ты.
<<Черноруких>>, чтоб его, <<не касается>>!
Я этой рыбине череп разобью, если...

--- Убери ружьё.
Разнесёшь несколько писем, чем незаметнее, тем лучше.
Идём со мной...

\section{Неудавшиеся переговоры}

Незаметного отступления не вышло.
Машина, как выяснилось, уже знала, куда мы направляемся, и успела подготовиться.

Со всех сторон к музею двинулись её <<солдаты>> --- роботы-техники, роботы-уборщики, производственные роботизированные механизмы.
Всё, способное хоть как-то нанести вред живому существу, двигалось к музейному комплексу.

Защитные механизмы музея могли выдержать сутки, может, двое под таким напором.
Но что делать дальше?

--- В данный момент единственное, что сможет выдержать длительную осаду --- Стальной Дракон, --- заявил Комар на экстренном собрании.
--- Там достаточно мест и пищи для всех.

--- До Стального Дракона минимум десять часов пути, --- возразил Фонтанчик.

<<Я прошу прощения, --- подал голос Пузырь из коммуникатора.
--- Дельфиньи лазейки имеют небольшую галерейку для тех, кто с ногами.
Мне отсюда не очень хорошо видно, но на первый взгляд она вполне проходима.
Дойдёте до места за пять-шесть часов>>.

--- Там высокая влажность, мы задохнёмся, пока дойдём, --- покачал головой я.
--- Сейчас ещё и достаточно холодно для нас.
Плюс там вертикальные лестницы, которые лично я недолюбливаю.

--- С нами дети, так что давайте без лестниц, --- поддержала Заяц.

--- Даже пять часов --- это много, --- сказал Фонтанчик.
--- Нет, дельфиньи дороги точно отпадают.

--- Ещё вариант, народ, --- сказал Красный-Мак-под-Кустами.
--- Я когда-то работал здесь садовником.
Есть технические ходы.
У них небольшая пропускная способность, но путь будет гораздо короче.
Навскидку до Дракона мы можем добраться за три часа.

--- Только не говорите, что вы хотите улететь с Тси-Ди! --- пропищала Сок-Стального-Листа.

--- О полёте речь пока не идёт, --- сухо сказал Комар.

--- Нам \emph{придётся} улететь, --- сказал я.

Наступила оглушительная тишина.

--- Небо, объясни, --- попросил Мак.

--- Наш единственный шанс был в том, чтобы договориться с Машиной.
И не говорите мне, что это не так, --- я поднял руки, останавливая негодующие возгласы.
--- Мы создали Машину.
Мы вложили в неё все знания, которые имели.
Назовите хоть что-то из того, что знаете вы и не знает Машина.
Я не могу.

Тси задумались.

--- Иметь доступ к информации и знать --- разные вещи, --- рассудительно заметил Мак, --- однако я соглашусь с Небо.
Мы не можем сказать, когда она обрела самосознание, сколько времени у неё было на освоение доступной информации.

--- Что там с переговорами? --- обратился я к техникам.
Те промолчали.

--- Она отвечает, Небо, --- ответил за всех Фонтанчик.
В голосе друга звучала растерянность, словно вся его прошлая жизнь была сном, а сейчас он проснулся.
--- Отвечает.
Вполне понятно и просто: <<Сдавайтесь>>.
Это не недопонимание, не ошибка перевода, это простой и понятный ультиматум.
Она не хочет союза, она хочет безграничной власти над нами.

--- Как она распорядится этой властью, думаю, объяснять не нужно, --- проворчала Лиана.

--- И если в переговорах успеха нет и мы останемся на планете, то наша гибель --- вопрос времени, --- подвёл я итог.
--- У неё есть все технологические линии, все данные о материалах Дракона.
Для воспроизводства тси требуются годы, а для того, чтобы устроить нам терракотового волка, требуется восемь дней.
Направленного ядерного удара Дракон не выдержит.

--- Она не будет бить напрямую, --- возразила Лиана.
--- Реактор Дракона при повреждении может уничтожить всю планету.
Однако ядерное оружие --- хороший способ не дать нам улететь, если мы поднимемся в космос.

Комар подумал и обратился к Маку:

--- Схемы комплекса у нас есть, но о состоянии технических коридоров --- никакой информации.
Ни о пропускной способности, ни о вентиляции.
Давно ты здесь работал?

--- Лет десять назад, --- ответил Мак.
--- Я пойду впереди, мне нужен будет один помощник, чтобы открывать вентиляционные колодцы.
Идти будет некомфортно, там сейчас холод и сильный сквозняк, скажи всем надеть костюмы.
Выйдем в районе Палат Войны.
Оттуда до Дракона около двух километров пути по демонстрационным помещениям.

\section{Закон меча}

Ситрис проснулся неожиданно.
Он уже привык ночевать в комнате Митхэ;
несмотря на то, что последние пять дней она шлялась непонятно где, разбойник был уверен, что никто не посмеет тронуть его здесь.
Однако нулевое чувство вдруг подсказало ему --- больше комната не будет надёжной защитой.

Ситрис бросился в зал.
Там обязательно должны были быть несколько наёмников, в присутствии которых его не станут трогать.
Уже в коридоре разбойник понял, что его ожидания оправдались сверх меры --- зал был забит людьми Митхэ до отказа.

--- Это было изнасилование! --- настойчиво убеждала Ликхлам нескольких воинов Храма.

--- Которое классифицируется как Насилие, --- равнодушно парировал один из них --- широкоплечий детина с квадратной челюстью и маленькими глазами, прячущимися в тени
каменных надбровных дуг.
--- Она не имела права лишать жизни жреца, совершившего Насилие, если он совершил его однократно.

--- Она защищалась!

--- От чего? --- осведомился воин.
--- Жрец не угрожал её убить.
Кроме того, убийство произошло пост-фактум, даже если изнасилование имело место.

--- Мразь, --- заявил Аурвелий.
--- Если бы ты или твой женщина кто-то потрогвать, ты бы тоже рассердиться и бить егво!
Как и я, как и любвой!

--- Это не меняет сути дела, ноа.
Выдайте Митхэ сейчас же, или разговор пойдёт на другом языке.

--- Живите долго!
Что значит <<на другом языке>>? --- возмущённо застрекотал Кусачка, позванивая бубенцами на шапке.
--- Вначале с нами переговоры ведёт один воин, а затем придёт один купец?
Где ваш дипломат?
Почему вы начали переговоры с многих угроз?

\ml{$0$}
{--- Твоего мнения никто не спрашивает, жук-навозник.}
{``No one asked your opinion, you dung beetle.''}

\ml{$0$}
{--- Он один из нас, чучело ты мясное! --- великан Кхотрис рявкнул так, что с ласточкиных ниш посыпалась пыль.}
{``They are one of us, you meat scarecrow!'' giant \Kchotris\ roared so loudly that dust was raised from swallow's niches.}
\ml{$0$}
{--- Ты сейчас весь Храм оскорбил!}
{``You have insulted all the Temple!''}

Наёмники разъярённо заголосили;
Ликхлам отчаянно замахала руками, пытаясь их успокоить.
Никто не обращал внимания на маленького травника, который смотрел на Кхотриса не отрываясь, словно в священном трансе...

--- Мы сами пойдём с Митхэ на суд, --- всё громче повторяла Ликхлам, оглядываясь в поисках поддержки.
--- Да, мы все!
Мы не отказываемся подчиняться правосудию!
Мы пойдём с Митхэ и докажем, что она права!

В толпе раздались возгласы одобрения.

--- Я ещё раз повторяю: суд над Митхэ ар'Кахр состоится только в том случае, если она предстанет перед ним безоружной и под конвоем.
В противном случае суд решит, что обвиняемая пришла не доказывать свою невиновность, а сражаться.

--- Ты извращаешь законы! --- крикнул из толпы женский голос.
Кажется, Сокхэ.

Воин тяжёлым взглядом поискал собеседника и ухмыльнулся.

--- Сможет ли неуловимый голос повторить это Советам?
Или он достаточно храбр лишь для того, чтобы кричать из-за мужских спин?

--- Она права, Митликх ар'Мар, --- подтвердил Ситрис, выйдя из тени коридора.
Все повернулись к нему.
--- Во-первых, не Митхэ должна доказывать невиновность, а суд должен доказать вину.
Это верно даже в том случае, если бы она растерзала жреца на глазах всего города.
Во-вторых, про оружие ничего не сказано, но зато в правилах есть недвусмысленный пункт --- <<обвиняемый может защищаться>>.
Если обвинитель говорит --- то словами, если обвинитель всуе хватается за оружие --- то уместен и клинок.
Ты считаешь, что присутствующие на суде испугаются и не смогут одолеть одну маленькую женщину?
Видимо, некому обучить их искусству боя и некому показать, как искать твёрдую опору.
Кому, как не Храму, лучше это знать.

Здесь нужно отметить: одним из достоинств Ситриса было то, что он входил в курс дела мгновенно, услышав лишь обрывок разговора.
Оформленные мысли в голову разбойника приходили задним числом.

<<Я, конечно, знал, что эта история кончится конфликтом Храмов, --- подумал он. --- Но это что-то из ряда вон...>>

Воин смерил разбойника взглядом и подошёл ближе.
\ml{$0$}
{Ситрис, как обычно, почувствовал, что тема разговора чересчур щекотливая для столь короткого расстояния между ним и собеседником.}
{As always, \Sitris\ felt that subject was extremely prickly for such a short distance between him and his interlocutor.}

--- Ты что-то сказал, человек? --- спокойно осведомился воин.
--- Мы слышали, что среди отряда есть трусливый бандит, прячущийся за спиной Митхэ ар'Кахр, и его хасетрасем подозрительно похож на твой.
\ml{$0$}
{Я бы на твоём месте сидел тихо и молча, как опарыш в дерьме.}
{If I were you, I'd be silent and unseen like a maggot in shit.''}

\ml{$0$}
{--- О, я в этом нисколько не сомневаюсь, --- театрально поклонился Ситрис, отступая на шаг.}
{``\`{O}, beyond any doubt you would,'' \Sitris\ theatrically nodded---stepping backward at the same moment.}
--- Кстати, а есть ли решение Советов насчёт меня?
Или <<подозрительно похожий хасетрасем>> --- твой единственный аргумент?

--- А оно требуется? --- лениво поинтересовался воин.
--- Я могу прихватить тебя вместе с Митхэ ар'Кахр, а потом навести на Юге справки.
Даже если ты не кутрап, препятствование исполнению решений Советов --- вещь предосудительная.
Впрочем, если ты заткнёшься и встанешь в уголке, я могу про тебя забыть.

Ситрис спиной почувствовал жар ярости, идущий от наёмников, и слегка отставил ногу назад.
<<Меча законы --- ночи бессонны, будь ты хоть ребёнок, хоть воин>>, --- повторяла Митхэ.
Неважно, на чьей ты стороне --- диктат военной силы приведёт лишь к хаосу.
Что-то в облике Митликха говорило --- законом меча жил и он, лишь прикрываясь обычаем.

Ситрис не глядя выудил из кармана заветный жетон, прикреплённый к поясу цепочкой.
Подкинул его, поймал и бросил взгляд на оказавшееся вверху изображение.

\ml{$0$}
{<<Ястреб>>.}
{\emph{Hawk}.}
Дело дрянь --- договориться не выйдет.
Митликх опьянён своей силой --- за его плечами закон.
Он носит закон, словно непобедимую саблю, словно непробиваемый щит.
Он забыл, что закон бессилен, когда ты один против десяти.
Он забыл, что есть люди, для которых закон --- лишь слова.

<<Может, Митхэ и докажет свою невиновность, --- промелькнуло у разбойника в голове, --- но что делать мне?>>

Воин бросил насмешливый взгляд на товарища --- очень уж глупо Ситрис смотрелся со своим жетончиком.
Именно в этот короткий промежуток ему в висок прилетел тяжёлый железный шарик, висевший на поясе рядом с заветной монеткой.
Спустя десять секхар схватка закончилась.

--- Хорошая драка должна быть сюрприз, --- заметил Аурвелий, любовно ставя на место расписную вазу.
--- Если не сюрприз --- драка негодный.

--- Tu sembri felice\FM, --- буркнул Ситрис.
\FA{
Похоже, ты счастлив (ноа-лингва).
}

--- Ancora qui, come te, --- улыбнулся Аурвелий.

Согхо, содрав со стены Свиток Осады, бросилась на обход --- закрывать ставни и бойницы.
Ликхлам с ловкостью циркового фокусника обшаривала одежду побеждённых;
из-под её рук фонтаном в стороны летело искусно спрятанное оружие, золото, кольца, серьги, конфеты, ключи, свитки пергамента...

--- Ключик на пять штучек, замочек-колокольчик, от кельи, --- бормотала она.
--- Ключик на семь... это от склада, это я возьму.
Храмовая пятка-печатка, иди ко мне, моя хорошая...
Письмеца от Короля-жреца... так, Кхотрис, отнеси Ситрису, это точно что-то важное...

Тем временем прочие наёмники Митхэ связали оглушённых и раненых воинов, растащили их по кельям и, закрыв двери храма, набросили на крючья тяжёлый засов.

\section{Палаты Войны}

Мы остановились перед дверью, выкрашенной в красный цвет.

Знаменитые Палаты Войны.

Эта дверь, в отличие от остальных, открывалась вручную.
Она была чересчур тяжела для одного, и Фонтанчик вместе с двумя-тремя кани-мужчинами навалились на неё изо всех сил.
Телепатический голос возвестил:

<<Помни, тси.
Эта дверь тяжела, но нет ничего тяжелее войны.
Открывая эту дверь, вспомни об оружии, которое воины несли на себе, толкали и волокли в снегу, жидкой грязи и цветущей воде ради одной-единственной цели --- убийства себе подобных.
Вспомни об этой бессмысленной и страшной ноше>>.

Дверь наконец подалась, и я взмахом руки позвал остальных следовать за мной.

По мере продвижения по Палатам Войны мы невольно обращали внимания на голографические изображения, светящиеся в специально созданном полумраке.

Древнейшая история сапиентов.
Толпы полуголых людей, кидающих друг в друга заострённые палки.
Правители на колесницах, огромные по сравнению с гибнущими за них <<чёрными>> людьми.

Затем палки приобрели сложный механизм --- наступила эра метательного оружия.
Окутанные пороховым дымом стрелки, сражённые раскалёнными кусочками металла люди падают, словно трава под лазерным лучом...

Последняя война.
Хирошимская катастрофа, по которой я когда-то давно, ещё в школе, готовил доклад --- первый раз, когда люди увидели ужас терракотового волка.
Ядерный гриб, оплавленные каменные стены домов и непонятно как уцелевшая арка храма, похожая на летящую птицу.
Люди, умирающие от лучевой болезни.
Концентрационные <<лагеря смерти>> --- при их появлении на экране все тси без исключения вздохнули и отвернулись.
Огромные могильники, в которых не было почвы --- её заменяли лежащие слоями человеческие тела.
Солдаты, соревнующиеся в отрубании голов.
Пыточные тюрьмы правителей, называвших себя прогрессивными и великими --- места, где человеку не оставляли и щепотки надежды, откуда самые лучшие, честные и сильные выходили запуганными и сломленными животными.
Сожжённые термитными смесями леса.
Люди в глухих, закрывающих лицо чёрных шлемах и камуфляжных куртках с одинаковыми надписями --- существа, выведенные и воспитанные единственно с целью причинения боли и смерти.

Фото людей, которых пытали и убивали тысячами за украденный колос и горсть сахара.
Захоронения, площадь которых превышала площадь посевов.
Приказы с лимитами на казни, внушающие ужас своим бесстрастным слогом и цифрами.
<<Прошу разрешения казнить 1000 человек сверх лимита>>.
И небрежная подпись сверху.

Война Тараканов.
Последняя и самая страшная, которую знали наши прямые предки.
Армии перепрограммированных солдат, не знающих ни боли, ни чувств.
Химические агенты, превращающие сапиентов в полуразумных, жаждущих крови существ.
Болезни, заражающие десятки миллионов и запускаемые одним нажатием кнопки.
Это была война, в которой жизнь стоила меньше, чем ничего.
Война, в которой была превращена в мёртвый камень некогда цветущая планета.
Война на уничтожение, в которой апиды, кани, люди, планты и дельфины планеты Ди заслужили право на существование и свободу.

--- Небо, --- прошептала Заяц, глядя огромными, полными слёз глазами на разворачивающееся на экранах безумие, --- Небо, неужели, чтобы выжить, мы должны стать такими?

--- Заяц, пойдём, --- Фонтанчик ласково взял подругу под руку и потянул её дальше.
Она не двинулась с места.

--- Почему мы не можем договориться с Машиной? --- в отчаянии крикнула Заяц, падая на колени.
--- Мы могли бы остановить всё сейчас, единым словом, единым жестом мира!

--- Зайчик, --- один из техников-людей склонился над сходящей с ума женщиной, --- Зайчик, я пробовал.
Я лично пробовал несколько раз.
Она не желает нас слушать.
Нашей вины здесь нет.

--- Зайчик, --- прошептал Фонтанчик и дождался, пока бьющаяся в рыданиях женщина не обратила к нему опухшее от слёз лицо, --- Зайчик.
Сейчас трудно всем.
Но если ты не можешь идти, я всегда могу тебя понести.
Цепляйся за меня, давай.
И закрой глаза, незачем тебе сейчас \emph{всё это} видеть.

Заяц обхватила руками мускулистую шею Фонтанчика и взгромоздилась на него.
Мы пошли дальше.
Но её замершие глаза по-прежнему неотрывно смотрели на разворачивающиеся вокруг сцены --- пытки, убийства, лежащие в песке тела, безумцы, бессмысленными речами зажигающие в необразованных толпах жажду чужой крови...

\ml{$0$}
{--- И ведь конца-края этому нет, --- шептала она.}
{``It has no beginning and no end,'' she whispered.}
--- Пока мы жили в мире и процветании на Тси-Ди, остальная Вселенная стонала от войны между Адом и Картелем.
Мы предпочли забыть о том, что творится за пределами нашей звёздной системы, о том, сколько планет прозябает в невежестве и забвении.
\ml{$0$}
{Вот она --- расплата за слепоту.}
{And that is it, the payback for our blindness.}
\ml{$0$}
{Вот она...}
{That is it ....''}

\section{Утайка}

\epigraph
{Культура зарождалась как средство выживания.
Она помогала передавать знания, исцелять, защищать и успокаивать.
Однако настал момент, когда сапиенты запутались в культуре, как пауки в собственных тенётах.
В этот момент не может быть правых и виноватых.
Это всё, что я могу сказать про эпоху Последней Войны.}
{Людвиг Вейерманн}

--- Это я виновата, --- сокрушалась кормилица.
--- Нельзя было закрывать на проблему глаза.
Я могла послать прошение.
Пусть суд приговорил бы Чханэ даже к смерти, пусть, но это прояснило бы дело сразу и позволило нам выиграть время.
Глупая я, глупая...

Кхотлам за весь вечер не притронулась к еде.
Я только сейчас заметил, как сильно она сдала --- под глазами пролегли тени усталости, губы высохли и потрескались.
На десять писем пришёл один вежливый ответ, смысл которого сводился к простой фразе --- это не наше дело.

--- Так это правда? --- спросил Хитрам.
--- Чханэ --- кутрап?

--- Она сама в этом призналась? --- ошарашенно спросила Эрхэ.

--- И ты молчала, Пёрышко? --- возмутилась Ликхэ.

--- Да, --- упрямо сказала Кхотлам.
--- Всё это так.

--- Прекрасно, --- нехорошо усмехнулся кормилец.
--- То есть из-за твоей утайки мы едва не стали Разрушителями.
Молодец, Кхотлам.
И не в первый раз ведь, как мне...

--- Не смей, --- кормилица сказала это совершенно спокойно, но сидящие за столом, не исключая Хитрама, поёжились.
Мужчина смутился и продолжил уже примирительно:

\ml{$0$}
{--- Кхотлам, закон есть закон.}
{``\Kchotlam, law is law.}
\ml{$0$}
{Она убила жреца, она в этом призналась.}
{She did kill the priest, she did confess.}
\ml{$0$}
{Я считаю, что казнь --- чересчур суровое наказание, но она должна быть...}
{I consider death too severe for a punishment, but she should be---''}

--- Жрец пытался убить её, --- вмешался я.
--- Это была самоза...

--- Она могла оглушить его, Ликхмас, --- поморщился кормилец.
--- И пожалуйста, не говори мне, что опытный воин-Скорпион не мог этого сделать.
Законные пути решения проблемы всегда...

--- Законные? --- перебила его Кхотлам и страшно рассмеялась.
Напряжение и гнев, которые копились в ней всё это время, прорвали плотину самообладания.
--- Митрам ар’Сар вошёл в наш дом, едва не выбив дверь.
Он лишил свободы человека, находящегося в нашем доме, не спрашивая у меня согласия.
В купеческом Дворе, который, как и три ведущие к нему дороги, является нейтральной территорией.
Это не просто незаконно, это ещё и оскорбление для каждого из вас лично.
\ml{$0$}
{Может быть, ты, мой мужчина, глас Сада и обитатель Двора, расскажешь о законе Храму, а не нам?}
{Well, my man, the voice of the Garden, the dweller of the House, why won't you remind the Temple, not us, of the law?''}

Сидящие за столом как воды в рот набрали.

\ml{$0$}
{--- Пошёл вон. Все пошли вон. Ликхмас, останься.}
{``Get out. All of you. \Likchmas, stay here.''}

Домочадцы, не говоря ни слова, поднялись и разошлись по комнатам.
Едва закрылась последняя дверь, как кормилица расплакалась.
Я подбежал и обнял её.

--- Прости, Лисёнок.
Все в городе думают точно так же, как твой кормилец.
Я бессильна что-либо сделать.

Я гладил её по волосам.

--- Кормилица, может быть, мне?..

--- Хитрам чересчур верит силе закона.
Я всю жизнь была дипломатом и знаю, что любой, даже самый точный закон можно обратить против человека.
Моего кормильца...

Кхотлам запнулась.

--- Именно поэтому судьбу человека решают Советы, дитя.
Нужен здравый смысл живого человека.
Нельзя прожить одними формулами.

--- Что мне делать, скажи, милая, добрая женщина? --- тихо спросил я.

Кормилица погладила меня по голове и, плача, ушла в свою комнату.

\section{Решение}

Пока оставшиеся тси подтягивались к Палатам Войны, Баночка и Заяц взломали систему видеонаблюдения.
Я понял, что взлом прошёл успешно --- Заяц глухо застонала и села на пол.
Все бросились к ней.

--- Небо, --- дрожащим голосом сказала Заяц, --- там... тси.

Наблюдение показало, что она права.
На Золотой дороге, отделявшей нас от корабля, ждали не только роботы-техники --- там стояли люди, планты и кани.
Их лица скрывали газовые маски, вызвавшие у меня моментальную реминисценцию --- солдаты поддержания порядка времён Последней Войны.
Они стояли неестественно прямо и неподвижно;
в их руках было всё, что могло сойти за оружие --- плазменные резаки, острые металлические инструменты и анестезионные пистолеты.

--- Метательное орудие, --- глухо констатировал Кукловод, ткнув пальцем в экран.
--- Вот эта большая штука.
Скорее всего, откуда-то с музейных складов.
И вот здесь, и ещё одна.

--- Небо, пожалуйста, --- прошептала Заяц, --- скажи, что это неправда...

--- Небо, --- подал голос Комар, --- мы должны узнать, как им помочь.

--- И при этом не умереть от руки одного из них, --- добавил Кукловод.

Я взглянул на друзей.
Они ждали моих слов.
Никто из них больше не хотел брать на себя ответственность.
И я вдруг понял, что не могу их винить.

Может быть, ты, потомок, назовёшь меня трусом.
Но я знаю цену своему решению.

Управляющий интерфейс техников-млекопитающих напрямую связан с третичной ассоциативной зоной коры и вычислительными серыми ядрами.
Наиболее возможный вариант --- Машина поместила техников-тси в мир иллюзий, поменяв врага и друзей местами на чувственном уровне.
Этакий надвинутый на глаза мирок, скрывающий правду.
Это решалось простым выключением импланта.
В случае же, если Машина нашла способ прямо или косвенно перепрограммировать лобную кору, попытка спасти товарищей принесла бы больше горя, чем пользы.
Они могли притвориться нормальными и нанести удар в тот момент, когда прочие расслабятся.
На научные изыскания времени у нас не было.

--- Не трогать этих существ, --- справившись с тяжестью в животе, приказал я.
--- Если они нападут --- защищаться.

--- Небо! --- рявкнул Фонтанчик.
--- Мы не сможем пройти мимо них!
Столкновение неизбежно!

--- У нас ещё есть дельфиньи лазейки, они выводят прямо...

<<Кстати, о дельфиньих лазейках, --- сказал Пузырь.
--- Здесь они тоже стоят, прямо у нас по курсу, с каким-то подозрительным агрегатом.
Нюхом чую, если мы сейчас высунемся из воды --- нам кранты.
Из более неприятных новостей --- они тем же агрегатом пробили дыру в смотровом ходу Палаты Достижений, минимум из пяти сегментов вытекла вся вода, и у нас тоже медленно начало мелеть.
Похоже, нам всё-таки придётся отрастить ноги, если мы хотим как-то попасть в космолёт>>.

Вот она, ответственность.
Она всегда встаёт на пути, не давая ни малейшего шанса пройти мимо.
Мой голос звучал, словно чужой.
Это говорил не я.
Мне ужасно хочется думать, что это говорил не я.

--- Ради жизни.
Идём на прорыв.
Уничтожить перепрограммированных роботов и тси с перепрошитыми имплантами.

Тси замолчали.
По коридору прокатился сдавленный вздох.

Комар бесстрастно несколько мгновений смотрел на меня, щёлкая пальцами.
Потом повернулся к отряду:

--- Внимание всем.
Я подтверждаю приказ Существует-Хорошее-Небо.

...На обратном пути никто не смотрел мне в глаза.
Да и вообще все старались не смотреть друг на друга.
Я шёл отдельно;
ко мне никто не желал приближаться, кроме Комара.

--- Я не смогу, любовь моя... --- тихо сказал я другу по дороге.
--- Я не смогу уничтожать своих...

--- Тебе ничего не нужно делать, Существует-Хорошее-Небо, --- сухо ответил Комар.
Его лицо, похоже, навеки превратилось в оплавленную радиацией терракотовую маску.
--- Особый отдел выполнит твой приказ.
Дальше полагаемся на тебя.

\section{Слово дают ножу}

\epigraph{
<<Да>> и <<нет>> весят одинаково, если я предлагаю помощь.
%{Yes and no have equal weights, if I borrow my help.}
}{
Присказка ноа
}

Оставалось одно --- мою женщину нужно было отбить силой.
От этой мысли по спине пробежала дрожь --- на миг в голове нарисовалась картина, как я перешагиваю через тело Кхараса, Ликхэ, Кхохо, у них же клятва, они не могут просто отойти в сторону...
Я посмотрел в окно --- на небе пылал закат.
Жертвоприношение могло начаться уже через час.

Я бегом бросился в свою комнату и плотно закрыл дверь.
Меня встретил Сиртху-лехэ, который по-хозяйски устроился на моей постели, в тени сундука.

--- Когда бессильно горло, слово дают ножу? --- лукавые, молодые, похожие на кошачьи глаза, обрамлённые морщинистыми веками, смотрели прямо на меня.

Я промолчал и полез в сундук за точильным камнем.

--- Ты не попадёшь на вершину пирамиды.
Жрецы знают, что ты её мужчина, --- без обиняков сказал седовласый охотник.

--- И что?

--- Хай.
Я подумал, что тебе может понадобиться тот, кто сведущ в игре теней, --- улыбнулся Сиртху-лехэ.

--- Мне не нужна помощь, --- отрезал я.

Сиртху-лехэ тяжело, с кряхтением сел и подтянул к себе трость.

--- Ликхмас, не валяй дурака.
Твоя кормилица подарила мне лишних двадцать дождей жизни.
Клянусь Безумными богами, это лучший подарок, который у меня был.
Конечно, если не считать детей, которых родила мне Иголочка...

Сиртху-лехэ покашлял и прислушался к тишине.
Потом продолжил вполголоса:

--- Хай.
Я не хочу умереть, рассыпаясь, как песчаный домик.
Как говорят у нас на севере, самая достойная смерть --- это когда твоё тело становится знаменем для отчаявшихся, щитом для беззащитных и пищей для голодающих.

--- Мне стыдно принимать твою жизнь, дедушка, --- пробормотал я.

--- Хаяй.
Стыдно требовать от меня такое.
А я предлагаю сам, --- для убедительности старик хлопнул кулаком по моей постели.
--- И не спорь, своей жизнью имею право распоряжаться только я.

Старик усмехнулся молодыми глазами, и я понял --- умирать он пока не собирается.

Я кивнул.

--- Хай-хай... --- задумался Сиртху-лехэ.
--- Времени предостаточно, мы даже успеем прищемить ягуару хвост.
Встречаемся у юго-западного угла храмовых земель.
Старайся, чтобы тебя никто не увидел, даже домочадцы.
Собери походный мешок и оставь здесь.
А ещё нам нужно...

\section{Подготовка к прорыву}

--- ...нам нужно оружие.

Эти слова Комар сказал тихо, но их смысл дошёл до всех.
Отныне нам \emph{придётся} стать такими.
Заяц снова расплакалась.

Палаты Войны, помимо тяжёлых воспоминаний, преподнесли нам кое-что ещё --- настоящее оружие.
Танковые доспехи времён Тараканьей войны, шестнадцать штук --- отремонтированные и идеально функционирующие.
Когда-то давно они устраивали кровавый ад на полях сражений, пока не были подбиты или захвачены.
Доспехи были предназначены для солдат-апидов Ди, однако те отличались мощным телосложением, и в командные капсулы могли поместиться любые тси за исключением, пожалуй, взрослых кани.
Комар велел своему отряду прихватить эти доспехи.

--- Если у них есть тяжёлые орудия, танки --- наш единственный шанс, --- сказал он прочим.
--- Мы пойдём впереди, прикрывая собой пехоту.
Наша цель --- вывести все орудия из строя, умрите, но сделайте это.
Иначе мы проиграем.

Сам Комар разобрался с управлением достаточно быстро, а вот Лиане и Соли пришлось несладко в не предназначенных для их вида кабинках.

--- Не приведи судьба сдохнуть в этой консервной банке, --- слышал я тихую ругань Лианы.

В конце концов двуногие особого отдела напрочь отказались от доспехов и прихватили из Палат Войны оружие попроще, но лучше подходившее к их телам.

--- Нет, Кристалл!

--- Я хочу пойти с тобой, Лиана! --- страстно шептала девушка, потрясая болтомётом.
--- Я хочу биться!

--- Ты не из особого отдела.
Ты --- студентка.
Когда всё закончится, ты продолжишь обучение.

--- Я не спрашиваю твоего разрешения, я!..

--- Связать её, --- распорядилась Лиана.
--- Руки за спиной.

--- Что?
НЕТ!
Гранит, отпусти меня, грязный предатель!
Ай!

--- Так, ты, в клетчатой рубахе.
Гранит, да?
Отвечаешь за неё.

--- Моё место --- рядом с вами! --- закричала Кристалл.
Её кудри растрепались, тугие мышцы вздувались буграми под гладкой кожей, тщетно пытаясь разорвать кабельную стяжку.
--- Ты подумала, как я буду жить с этим дальше?!

--- Это меня не касается, --- отрезала Лиана и скомандовала построение.

Кристалл упала на колени и уткнулась лицом в пол.
Её тело содрогалось от глухих рыданий.
Кто-то аккуратно взял её болтомёт и спрятал среди стеллажей.

--- Небо, мы прорвёмся.
Отведи всех назад в Палаты Войны и постарайся отсечь Золотую дорогу от внешнего мира.

Комар сказал это чересчур резко, и я воспринял резкость на свой счёт.
Почему я не обнял его на прощание?
Почему?..

Взломать систему удалось с большим трудом.
Проклятая Машина подключила к вычислительным модулям все имеющиеся мощности и первым делом обезопасила себя от техников.
Нам пришлось прибегнуть к противоходу.
Я взломал несколько малозначимых компьютеров на периферии и подключил их к собственному мозгу, чтобы за счёт обрывков информации быстрее подобрать метод.
Да, это был риск.
До сих пор помню, как дрожали мои руки.
Минуту спустя к системе подключилась Заяц, внеся в потоки необходимый хаос, и только с мощностями её мозга мы смогли обмануть Машину.
Если бы не Заяц --- взломанными оказались бы не системы музея, а мозг Существует-Хорошее-Небо.

Меня терзала ужасная догадка --- мы единственные, кто ещё борется с Машиной.

\section{Сестрёнки и папа}

\epigraph{
\ml{$0$}
{Просите --- и дано вам будет, ищите --- и найдёте, стучитесь --- и отворят вам.}
{Ask and it will be given to you; seek and you will find; knock and the door will be opened to you.}
}{Сапфировая Книга, Прядь Матеуша 7:7}

Что мы можем сделать одни?
Я понимал, как мало у нас шансов на успех.
Да, возможно, я смогу одолеть в бою этих новых жрецов --- я ни разу не видел их на тренировочной площадке.
Но для этого нужно добраться до зоны молчания.
У меня не хватит сил одолеть даже Ликхэ --- последний дождь она предпочитала тренироваться со мной и неизменно побеждала в спарринге.

Я заглянул в комнату к сёстрам.
Манэ и Лимнэ уже спали;
лежанка уже была чересчур маленькой для двоих, но они по-прежнему ночевали вместе, крепко прижавшись друг к другу.
Я присел рядом с ними и погладил их по головам.
Когда-то светлые, цвета какао, волосы потемнели и начали немного различаться по цвету;
волосы Лимнэ теперь походили на лён, а волосы Манэ --- на потемневший мёд.
Я только сейчас заметил, что Лимнэ с её весёлыми глазами-капельками похожа на Хитрама, а Манэ пошла в Кхотлам.
Сестрёнки тут же открыли глаза и хитро посмотрели на меня.

--- Храни ваш сон лесные духи, солнышки мои, --- прошептал я и по очереди поцеловал их.
--- Завтра будет новое утро и новые радости.

--- Ты пошёл за Чханэ? --- серьёзно спросила Лимнэ.

Я промолчал.
Сестрёнки переглянулись.

--- А я тебе говорила, --- наставительным тоном сказала Лимнэ сестре.
--- Братик не мог поступить по-другому.

--- Мы тебя больше не увидим? --- прошептала Манэ.
В её глазах замерли слёзы.

--- Я очень надеюсь, что когда-нибудь смогу вас повидать, --- улыбнулся я.
--- Вы вырастете большие и сильные.
Если я умру, то буду ждать вас в пристанище и сделаю вам самые красивые игрушки из ветра, тумана и утренней росы.

--- Мы будем готовить пирог на каждое первое число и ждать тебя, братик, --- сказала Лимнэ.
--- Но если ты вдруг придёшь на пятое, мы обязательно сделаем ещё.

Я обнял сестрёнок и дал волю чувствам.
Они, смирные и тёплые, замерли в моих объятиях.

<<Братик не мог поступить по-другому>>.
\ml{$0$}
{Эти слова звучали во мне снова и снова, пока я проверял снаряжение у себя в комнате.}
{I heard these words over and over again while I was collecting the gears in my room.}
Огниво, стрелы, верёвка, два запасных комплекта одежды, иголка с ниткой, липкие шарики, спальные мешки, кадильница, туго свёрнутый молодой пергамент, перья, походная чернильница-непроливайка, милый розовый фонарик для посоха с ликом Сата-скитальца...
Под конец я вытащил праздничную одежду и зелёное платье Чханэ, оставшееся у меня.
Я совсем забыл их постирать;
шёлк был немного липкий на ощупь, пах костром, травами, цветами, зёрнами какао и свежеиспечённым хлебом.
Я улыбнулся и как можно аккуратнее сложил праздничную одежду в мешок.
Наверное, в дороге можно отыскать купель и мыльные растворы;
я очень надеялся, что так оно и есть.

--- Ликхмас.

Кормилец подошёл незаметно, окинул умным взглядом мешок, мою фалангу и лежащий на кровати точильный камень.
Я замер и хмуро посмотрел на него, ожидая отповеди.
Но он вдруг вытащил из кармана тяжёлые золотые серьги и протянул мне.
На опалённом солнцем лице играла добрая и грустная улыбка.

--- Серьги моей прародительницы, которая жила где-то в землях ноа.
Говорят, ей тоже пришлось бежать.
Серьги очень старые и ценные.
На них можно купить еду и оленей.

Я молча принял подарок и завернул его в рукав зелёного платья.

--- Ликхмас, возможно, мы с Кхотлам упустили этот момент в воспитании.
Просто знай: просить помощи не зазорно.
Люди могут не разделять твоих взглядов, им могут не нравиться твои цели, но ты удивишься, если узнаешь, сколько из них готовы тебе помочь --- просто потому, что ты попросил помощи.

Я снова промолчал.

--- Может быть, следует поднять город?

Я покачал головой.

--- К тому времени, как будет принято решение, Чханэ уже принесут в жертву.

--- Чем я могу ещё тебе помочь? --- просто спросил Хитрам.

--- Помоги одеться, --- сказал я.

Вскоре всё было готово к походу.
Я проверил в последний раз ремни на <<разбойничьих крыльях>> и потянулся за наручами;
кормилец тут же протянул мне свой, одинокий, который он зачем-то заткнул за пояс.
Обычный, отделанный кожей наруч, пожалуй, разве что чересчур тяжёлый...

Я почти выхватил его из широких рабочих рук.

--- Ликхэ в кустах возле колодца, я её одеялом накрыл, --- сообщил Хитрам и выбросил за подоконник любовно свёрнутую бухту.
\ml{$0$}
{--- Храни тебя лесные духи, дитя.}
{``May the Spirits bless you, child.}
\ml{$0$}
{Я закрою окно.}
{I will close the window.''}

\chapter{Ночь}

\section{Старые добрые времена}

Я притаился в колючих зарослях юкки, ожидая сигнала Сиртху-лехэ.

<<Чем менее удобно твоё укрытие, тем меньше шансов, что там будут искать, --- говорил Конфетка, заставляя меня часами сидеть в колючках или висеть на дереве.
--- Твоё тело должно быть готово к статическим нагрузкам и неприятным ощущениям>>.

Вскоре на этаже раздался тихий стук --- словно упало тело в доспехах.
Кто-то возбуждённо свистнул --- сигнал тревоги.
Тут же раздался ещё один стук, и из темноты три раза проплакала каменная жаба.
Можно было лезть.

Я осторожно закинул верёвку, завернулся в одеяло и полез, стараясь держаться в тени козырька.
Уже почти у края каменная жаба загулила снова, и я замер.
Наверху тихо шуршали рубахи --- поднятые по тревоге воины рассредоточились на втором этаже.

Я мимоходом восхитился старым охотником.
Его не было видно нигде.
Между тем прямо надо мной упал ещё кто-то, его товарищ лёг пластом и тихо выругался.

<<Опасайся стариков, --- вспомнил я слова Конфетки.
--- Тебе может показаться, что их мышцы вялы, что их глаза слабы, что их речь медленна и несуразна.
Но вспомни, через сколько сражений они прошли, прежде чем стать стариками>>.

Вдруг в темноте возле домов показался сгорбленный силуэт с луком, и о колонну щёлкнула шальная стрела.

<<В шестнадцати на восемь, --- прожужжал наруч.
--- Вы трое, спуститесь, атакующая схема четыре.
Ликхэ, мы тебя ждём, поддержка по сигналу>>.

<<Я уже на площади>>, --- чётко отбарабанил я.
Какое счастье, что Кхарас отказался от изобретения Трукхвала --- с тем большим наручем вождь вмиг заподозрил бы неладное.

Воины отправились исполнять приказание.
Один спрыгнул со второго этажа в каких-то пяти шагах от меня.
Воины уже почти дошли до цели, как силуэт исчез.

<<Куда он подевался?
Там же...>>

Вскрик --- и посланных на поимку воинов стало на одного меньше.

<<Все вниз, --- прочирикал вождь.
--- Он там>>.

<<Кхарас, нам не следует покидать храм.
Это отвлекающий манёвр>>, --- возмутился Ситрис.

<<Я его узнала, это старик из дома Люм>>, --- сказала Кхохо.

<<Значит, я прав.
Где-то должен быть Ликхмас, и не только он, --- предположил Ситрис.
--- Кхарас, труби тревогу, поднимай горожан>>.

<<Отставить, справимся сами.
Ликхмас и Сиртху дома.
Я оставил Ликхэ сторожить>>, --- ответил Кхарас.

<<А ты уверен, что Ликхэ по-прежнему сторожит, а не валяется без чувств в кустах?>> --- возразил ему Ситрис.

\ml{$0$}
{<<Со мной всё отлично>>, --- возмутился я.}
{\emph{``I'm all right,''} I said.}

<<И ты, похоже, настолько этому обрадовалась, что прохлопала старика, --- заметил Ситрис.
\ml{$0$}
{--- Внизу есть кто живой?>>}
{\emph{``Anybody alive down there?''}}

Ответом было молчание.

\ml{$0$}
{<<Кхарас, труби в ракушку, тебе сказали, --- прожужжала Кхохо.}
{\emph{``\Kcharas, blow the shell, you are told,''} \Kchoho\ buzzed.}
\ml{$0$}
{--- Ты оглох или сдурел?}
{\emph{``You deaf or mad?}}
\ml{$0$}
{У нас четверых с ходу положили!>>}
{\emph{We've got four down on the fly!''}}

<<Если это тот старик, то он был в прошлом одним из лучших охотников, --- сообщил Кхарас.
--- Говорят, что это он Живодёра\FM\ убил...>>
\FA{
Живодёр --- огромный ягуар-людоед, который в течение десяти дождей терроризировал окрестности Тхитрона.
Джунгли, бывшие его вотчиной, ныне носят его имя.
}

<<Благодарю, что напомнил, с кем мы имеем дело, --- я почти слышал привычный рассеянный сарказм в сообщении Ситриса.
--- Так мы трубим тревогу или нет?>>

\ml{$0$}
{<<Я уже всё сказал --- раковину не трогать.}
{\emph{``You have the order: don't touch the seashell.}}
По двое встаньте в дверных проёмах и следите за вторым этажом.
Двое у главного входа, трое встаньте на указанные точки, --- Кхарас методично назвал координаты позиций.
\ml{$0$}
{--- Ликхэ, ложись, где ты есть.}
{\emph{``\Likchoe, lie down where you are.}}
Держи колени\FM\ в направлении четыре-шесть, пока не будет точных координат или движения.
\FA{Имеется в виду стрельба на уровне ног.
Кхарас хотел взять старика живым.
}
\ml{$0$}
{По возможности избегайте боя, скоро уже всё закончится>>.}
{\emph{Avoid fighting as long as possible, it's almost done anyway.''}}

\ml{$0$}
{Храм перешёл к глухой обороне.}
{The Temple has circled the wagons.}
Именно этого и добивался Сиртху-лехэ.

Я подтянулся на руках и спрятался за колонной.
Судя по данным Кхараса, прямо за мной должно быть двое.

<<Он здесь, прямо под восточной стеной>>, --- сообщил кто-то.
Я выглянул из-за колонны... и встретился глазами с Кхохо.

<<Ящерицыно семя...>> --- подумал я и медленно взялся за нож.
Кхохо с живейшим интересом наблюдала за мной, словно я был котёнком, играющим с костяной погремушкой.

--- Ты меня убить собрался, Лисичка? --- осклабилась она.

--- Ликхмас, быстрее, знак кирпича был, жертвоприношение вот-вот начнётся! --- выпрыгнул из темноты Ситрис.
Воин улыбался во весь рот.
--- Как в старые добрые времена...
Я так и знал, что ты отсюда полезешь...

--- Да, время.
Пойдём.
Давай верёвку, --- Кхохо, не дожидаясь разрешения, вырвала у меня из рук бухту.

Ситрис подозвал меня к стене и подтянул тонкий, едва заметный в темноте шнурок, заранее пропущенный через колонну сверху.

--- Тебя ждут у лесов, потому что с них проще всего забраться на крышу, старайся там не светиться, --- скороговоркой зашептала Кхохо, яростно завязывая узлы.
--- На крыше пятеро, трое у алтаря, двое с краю, с этой стороны.
Держи ружьё, Ситрис, дай ему, в магазине пять синих стрел.
Смажь нож кураре, хотя нет, в бездну --- Чёрного плесни, только сам осторожно, не порежься.
Да не сейчас, глупый, как ты с отравленным ножом полезешь?
Бей со спины, как мы тебя учили, без самодеятельности --- шея, икры, плечи, старайся одним ударом захватить нескольких и придерживай тела, они громко падают.
Одного оставь для Безумного, а то городу придётся несладко.
Об остальных мы позаботимся чуть попозже.
Сиртху-лехэ молодец, но его сейчас убьют.
Выходи через главные ворота, там не будут ждать, Ситрис встанет на часы...

Громкий воинственный клич Сиртху-лехэ и лязг оружия подтвердил --- старик попался.

--- Кхохо, Кхохо, --- остановил я её, --- город не сможет без жрецов.
Как вы...

--- Хай, так уж и не сможет, --- ухмыльнулась женщина.
--- Ты думаешь, что мы не сможем полечить или принести жертву?
Я уже приносила один раз, совсем одна.
Это несложно.

--- Но...

--- Лис, если выбирать между шайкой разбойников и пустым храмом, я выберу второе.
Попросим помощи у горожан, в конце концов.
<<Я принесу жертву и в скорби, и в болезни, и при смерти...>>.
По мне, вот эта строчка из вашей клятвы --- полное дерьмо.

Я бросился к ней на шею.
Кхохо обняла меня с такой силой, что у меня хрустнули рёбра.

--- Держитесь.
Я приведу в город Бродячий Храм, если мне удастся его найти.
Надеюсь, они остались не только в легендах.

--- Лезь, Ликхмас, и спасай Чханэ, --- улыбнулся Ситрис.
--- Мы справимся.
Мы все --- один Храм.

\section{Переплёт}

\epigraph
{Бой нужен лишь для того, чтобы спустить пар чувств.
Всё прочее можно достичь миром.}
{Лусафейру Лёгкая Рука}

--- Мы все повязаны одной красной лентой.
Можете меня убить, но вам не спрятаться от грядущего дождя под моим телом.

Ситрис стоял в кольце наёмников.
По его смуглому лбу струился холодный пот.
Под рёбра южанина уткнулись восемь фаланг и одно копьё, готовых в любой момент поднять его в воздух.

--- Да перестаньте вы! --- бессильно воскликнула Ликхлам, всплеснув руками.

--- Ты нас использовал, --- прорычала Согхо.
Копьё больно укололо разбойника в бок.

--- Да, --- легко согласился Ситрис.
Его колени ходили ходуном, но голос лился чистым, звучным потоком, словно песня пропавшего менестреля Атриса.
--- Да, я вас использовал.
Но я вас использовал, потому что хотел жить, а не потому, что не посмел сказать слово Храму, которым командовал.
Кто хочет меня убить --- давайте, налетайте!
Кто хочет назвать меня трусом --- пожалуйста, прошу вас!
Но не смейте говорить, что сейчас я предал своего командира!

Наёмники застыли.
Острия неуверенно отодвинулись.

--- Вы сами виноваты в сложившейся ситуации, --- заключил Ситрис.
\ml{$0$}
{--- Кто из вас вступился за Митхэ, когда Король-жрец отказался выполнить обещание?}
{``Who among you supported \Mitchoe\ when Priest-king refused to honour his word?}
\ml{$0$}
{Разумеется.}
{Indeed.}
\ml{$0$}
{Вы все думали, что я заслужил презрение и боль, что я страдаю за дело.}
{You thought I deserved pain and scorn, you took my misery as just punishment.}
Вы все думали, что участь кутрапа --- это что-то далёкое, что никогда не коснётся вас.
\ml{$0$}
{Смотрите теперь, смотрите во все глаза.}
{Look now, all of you, keep your eyes open.}
\ml{$0$}
{Чем отличается участь вашего командира, которому вы дали клятву верности, от участи презренного труса Ситриса?}
{The fate of your head you swore allegiance to, the fate of despicable coward \Sitris---what's the difference between them?''}

Воины молчали.
Наконец клинки опустились.
Согхо вышла вперёд и прижала руку к груди.
Её примеру последовали остальные.

Разбойника охватило странное чувство лёгкости.
Отряд чести стал его крыльями, и Ситрис знал, что эти крылья способны нести его за пределы возможного.

--- А теперь найдите Митхэ, --- скомандовал он.
--- Мы идём в город на суд.

--- Не выйдет, --- из коридора показался запыхавшийся воин.
--- Ситрис, извини.
Жрецы всё видели и спустились в город.
Боюсь, после их рассказа нас никто не будет слушать.

В зале воцарилось молчание.

--- Мы в первеплот, --- констатировал Аурвелий.

--- Мы в переплёте с того самого момента, как Митхэ покрасила жрецом келью, сенвиор Амвросий, --- огрызнулся Ситрис.

--- Я тебвя не обвинвять, Ситриций, --- заметил Аурвелий.
--- Нет и нет, клянусь глваза моей matri, я с твобой unas capita ut'e manie\FM.
\FA{
Единая голова и руки (ноа-лингва) --- полная поддержка.
}
Я имветь, что не все отрвяд сейчас в храм.
Я совветую вспомнвить, что капита Миция назначвить тебвя старшина, и позаботвиться о её лвюди.

Ситрис испуганно выругался.

--- Ликхлам, Согхо, Кхотрис!

--- Да, старшина, --- воины вытянулись в струнку.

--- Быстро бегите в город и попытайтесь перехватить кого-то из жрецов.
Никакого насилия, просто объясните ситуацию и скажите, что почти все воины живы.

Ситрис махнул рукой и ринулся на крышу.
Он осторожно выбрался из мёртвой зоны и лёг на узкий карниз, стараясь держаться подальше от колонн и в то же время не выходить из тени.

Вскоре окрестности огласило тонкое переливчатое поскуливание.
Уиуиуи...
В ответ завыл ветер и эхо загуляло по полуденному Тхартхаахитру.

И снова --- уиуиуи...
Сигнал подхватили двое --- в квартале ткачей и у Молочного моста.

Сидевший возле складов воин --- Ситрис не мог понять, кто --- встал и быстро пошёл прочь с площади.
К нему присоединились ещё пятеро, из постоялого двора.

Уйти им не дали.
На улице Скромных Апельсинов из подворотни полетели стрелы.
Четверо воинов упали сразу, двое оставшихся --- судя по тактике, Скорпионы --- забились в дверной проём и прикрылись клешнями.

С десяток крестьян опрометчиво бросились к ним.
Их тела так и остались лежать, создав вокруг Скорпионов мёртвую баррикаду.

Исход битвы решили две рыбачки.
Они использовали простое орудие своего промысла, неожиданно превратившееся в грозное оружие --- перемёт.
Рыбачки бросили спутанный перемёт в воинов и выдернули их из укрытия.
Подбежавший кузнец два раза опустил молот, разбив прикрытые шлемами головы.

Ситрис выругался.
Он уже успел пожалеть, что взялся за оружие.

Нападавшие задержались ровно настолько, чтобы влить каждому мертвецу в рот хэситр, растащить тела по укромным местам и засыпать песком кровавые пятна.
Вскоре об отгремевшей схватке не напоминало ничто --- солнце по-прежнему заливало обманчиво сонную, пустынную улицу.

\section{Боевой нож, жертвенный нож}

Я взобрался на карниз и осторожно выглянул.
Наблюдения Кхохо подтвердились --- двое стояли прямо передо мной.
Я вздохнул и очистил разум, сделав его похожим на свежий лист пергамента --- Конфетка считал, что крадущегося врага выдают чересчур громкие мысли...

Тела жрецов обвисли на собственных поясах.
Я убил их на нижней точке выдоха, без шума --- Конфетка мог бы гордиться своим учеником.
Я аккуратно положил тела на камень, вложил нож в ножны, вытащил духовое ружьё и обошёл поверженных врагов.
Шнурок штанов зацепил висевший на поясе мертвеца амулет Сана, и бронзовый диск предательски звякнул.
Трое у алтаря обернулись ко мне.

<<Ящерицыно семя...>> --- только и успел обругать я себя.

Первый жрец улыбнулся и снова занялся жертвой:

--- Ликхмас ар'Люм.
Ну здравствуй.
Мы тебя ждали.

--- Лис, это ловушка... --- выдохнула Чханэ.
--- Беги... ммм... ээии!..

Первый воспользовался неосторожной фразой, чтобы засунуть девушке в рот расширитель.
Только сейчас я понял смысл этой детали.
Жертвам не давали шанса прикусить язык и облегчить таким образом свою участь.

Чханэ, прикрученная бронзовыми скобами к базальту, яростно чирикнула на мучителя.
Два адепта медленно вытащили ножи и направились ко мне.

Я поднял духовое ружьё, взял на прицел одного из жрецов, прикинул угол до другого.
Стоит кому-то из них двинуться, указать мне направление --- и мгновение спустя в каждом будет по две синих стрелы...

Жрец даже не сделал попытку уклониться.
Он широко ухмыльнулся, и в тот же самый момент ружьё треснуло вдоль.
У меня замерло сердце.

<<Лесные духи, почему это произошло именно сейчас?!>>

Запоздало пропел спусковой механизм.
Осколки ружья с дробным деревянным стуком упали на каменные плиты.
Я попятился и вытащил нож.

Новые жрецы неожиданно оказались мастерами ножевого боя.
Такую подготовку и безупречную тактику редко можно встретить даже среди воинов.
Я пожалел, что не взял с собой щит.

Атака, отражение.
Нож едва не вылетел из моей руки, по правому запястью проехал длинный клинок одного из адептов, окровавив желтоватую ткань рубахи.
Руку стегнуло болью.
Лицо жреца скрывала утиная маска, глаза над ней насмешливо щурились.
Отбросив мысль о возможном яде, я снова ринулся в атаку.
И снова два жреца отбили её лениво, словно я был назойливым насекомым.
На моей рубахе прямо под воротом осталось два аккуратных надреза.

<<Лесные духи, что это за нелюди?>>
Меня вдруг охватило давно забытое чувство --- чувство мальчишки, которого зажали в углу два матёрых разбойника.

--- Alle, finita fiesto\FM, --- донёсся со стороны алтаря ухмыляющийся голос Первого.
\FA{
Всё, хватит игр (сектум-лингва).
}
--- Conpeda et co h\r{a}'s\'{e}n-nai\FM.
\FA{
Связать и отвести в <<холодную комнату>> (смесь сектум-лингва и цатрон).
}

От сказанного меня на мгновение пробрал мороз.
Я слышал похожий язык --- ноа-лингва, на котором говорили приморские племена Кита.
Но сам выговор...
Я словно наткнулся на древнюю глубоководную рыбу --- сильную, огромную, оснащённую острыми плавниками и непробиваемыми костными щитками, неповторимо грубую и уродливую в своей давно отжившей древности.
И эта рыба жила и плавала прямо передо мной...

Мне пришлось уйти в глухую оборону.
Пальцы порезанной руки не желали сжиматься --- лезвие надсекло нерв, не задев крупные сосуды.
Жрецы не просто были мастерами --- они били с хирургической точностью.
Я перекинул нож в левую, изловчился и изо всех сил толкнул обоих врагов, сбив их с ног.
Они какое-то время лежали и хохотали во всё горло, словно происходящее было игрой.
Короткая передышка дала возможность вынуть фалангу, сорвать верёвку с пояса и зубами прикрутить оружие к непослушной руке на манер щита.

--- Да хватит вам уже, --- одёрнул их Первый, впрочем, не слишком убедительно.

Клинки свистели, как сиу-сиу после дождей, когда им приходит время выводить птенцов.
Я пытался прорваться к алтарю, но меня теснили ко входу в крипту, где, вероятно, сидели ещё несколько жрецов.

Первый жрец тем временем занимался своим делом.
Я бросил взгляд на алтарь как раз в тот момент, когда обсидиановый нож опустился в первый раз.

Мне показалось, что я слышу рёв смертельно раненного зверя, медленно переходящий в дикий женский крик.
Рука, её рука, которую я столько раз целовал долгими ночами, ещё живая, ещё тёплая, с глухим стуком упала на жёлтый камень.
Жрец снял с жаровни раскалённые щипцы.
Зашипела и затрещала горящая плоть, и крик усилился.
Чханэ не молила и не плакала, только кричала от боли.
Потому что нельзя было не кричать.

Я стиснул зубы, продолжая неравный бой.
Надежда ещё жила.
<<Он прижёг, значит, кровь остановилась.
Помогите мне, духи лесные, только позвольте мне её отсюда вытащить...>>

Осознание смертельной опасности творило чудеса.
<<Селезень>> совершил первую ошибку, ставшую для него последней.
Разрубленная маска улетела в сторону, тело грохнулось о колонну и сползло вниз, а я в десятый раз мысленно поблагодарил Конфетку за науку.

--- Aiax\FM! --- каркнул второй адепт.
\FA{
Клич Красного Картеля.
Вероятно, искаж. auxa (сектум-лингва) --- <<помощь>>.
}
Глумливость в его глазах вдруг исчезла, он впервые увидел во мне серьёзного противника.
Наружу выбежали оставшиеся жрецы и с воинственным кличем бросились в бой.
В это время в кровавом облаке жуткого крика с алтаря мягко упала вторая рука.

Шансов у меня не было.
Усталость камнем засела в груди, я рыскал из стороны в сторону, избегая ближнего боя.
Ножи, стрелы, какая-то утварь пролетали в чечевичном зерне от моего тела.
Чей-то ловкий клинок надрезал ещё один нерв --- меня стегнуло болью, и рука обвисла насовсем.
Болтающуюся мёртвым грузом фалангу пришлось срезать.
Добить, прекратить мучения любимого человека --- эта мысль засела в пульсирующей голове.
А там будь что будет...

Зарычав, забыв о боли, я одним безрассудным десятилоктевым прыжком пронёсся над головами адептов, целясь ножом в шею Чханэ.
Но мою ногу обвила невесть откуда взявшаяся верёвочная петля, и я под новый крик боли всем своим весом шлёпнулся на камень рядом с алтарём.
Нож не долетел до цели какую-то пядь.

Жрецы-помощники втроём навалились на меня, не давая даже вздохнуть.
В спину вонзилась игла --- тело медленно расслабилось, дыхание замедлилось, перед глазами понеслось звёздное небо.

<<Всё, отбегался.
Ситрис... Кхохо... Неужели вы меня...>>

\section{Предатель}

\epigraph{
\ml{$0$}
{Рим предателям не платит.}
{Rome does not pay traitors.}
}{Квинт Сервилий Цепион.
Ответ приближённым Вириата, умертвившим своего вождя в надежде на вознаграждение.
Культура Ромай, Древняя Земля.}

--- Они знали всё, --- Митхэ бушевала настолько, насколько позволяла ситуация.
--- Наши позывные, нашу тактику.
Всё до мелочей...

Удерживаемый наёмниками храм пал в одно мгновение, словно в его стенах забаррикадировались играющие дети.
Митхэ, Акхсар, Ситрис и Эрхэ пробирались по подземному ходу.
Ситрис уверял, что ход вёл к Ху'тресоааса, где раньше располагался сгоревший квартал кожевников.

--- Кусачка не сказал, где именно находится выход, но он уверен, что люди о нём не знают, --- пояснил разбойник.
--- Значит, как минимум он должен быть хорошо спрятан...

По лицу Митхэ текли злые слёзы.
Там, в тёмном зале остались лежать её люди.
Они прошли с ней половину обитаемого мира, делили с ней пищу и воду.
Она смеялась их шуткам, она помогала устраивать их зачатых в походах детей, а порой и сама делила с ними ложе.
Они были родными --- кормильцами, братьями, сёстрами, которых Митхэ не знала в детстве.
\ml{$0$}
{Их больше нет.}
{They've gone.}

<<Мои люди и один ни в чём не повинный травник>>.

\ml{$0$}
{Кусачка тоже погиб.}
{Biter was killed too.}
До этого момента Митхэ даже не подозревала, что травник умеет драться.
Его голова едва доставала до пояса врагов, но он прыгал, словно взбесившийся сверчок, расшвыривая людей и нанося им страшные раны когтями.
Он мог беспрепятственно уйти в город к своему народу, но предпочёл защищать отряд чести.

\ml{$0$}
{<<Ты Это сделал, Кусачка.}
{\textit{Thou hast done This, Biter.}}
\ml{$0$}
{Возвращайся, и ты снова Это сделаешь>>.}
{\textit{Thou shalt come back, thou shalt have done This again.}}

Ход круто повернул к северу и закончился завалом.
Воины оттащили несколько камней и скользнули наружу, в заросли колючей черноягоды.
Кусты радостно проросли прямо в полуразрушенном доме, почуяв пробившися сквозь остов крыши солнечный свет.
Кусты жили своей медленной растительной жизнью.
Если бы у них были глаза, они бы видели, как солнце летает вокруг Тра-Ренкхаля, подобно привлечённому собственным отражением индиго-светлячку, кружащему у погасшего фонаря.
Звёзды были бы для них кольцами, а не точками.
И смерть их всегда была неожиданной, как было неожиданным всё, что движется быстрее солнца.

--- И кто нас предал? --- Эрхэ озвучила давно мучивший всех вопрос.

\ml{$0$}
{Все посмотрели на Ситриса.}
{All as one glowered at \Sitris.}
Тот мгновенно оказался в пяти шагах от Митхэ.

\ml{$0$}
{--- Это не я!}
{``It wasn't me!''}

\ml{$0$}
{--- А чего тогда боишься? --- осведомилась Эрхэ.}
{``So why are you so scared?'' \Oerchoe\ asked.}
Костяшки пальцев, сжимающих саблю, побелели.

\ml{$0$}
{--- Если тебя в чём-то подозревают, бояться нормально, --- пояснил разбойник.}
{``If you are suspected of something, it's normal to be scared,'' the outlaw explained.}
\ml{$0$}
{--- Митхэ, это правда не я.}
{``\Mitchoe, it really wasn't me.''}

--- Это не он, --- подтвердил возникший из ниоткуда Аурвелий, невозмутимо стряхивая с рубахи сухие листья.
--- Я развведка.
Ликламия продать нас, как курвица.
Только это не помвочь --- тело уже в землвя.

--- Ты её убил, сенвиор Амвросий? --- спросил Акхсар.

--- Лвюди Король-жрец лвюбезно сделать это за меня, --- мрачно ответил Аурвелий.
--- Её убвить воин сразу после информация, и Король-жрец дваже не пожурить его.

Вдруг перед воинами возникла ещё одна фигура.
Митхэ среагировала молниеносно --- Легенда Серого Рассвета обагрила землю и безмятежно растущие кусты.
Молодая женщина упала, удивлённо простонав: <<Хай?>>

--- Оружейница из Серого Рассвета, --- буркнул Акхсар, разглядывая убитую.
--- Правила мне клинок.
Хэситр есть?

--- Возьми у меня в правом нижнем кармане, --- Эрхэ бросила другу свой мешок.
--- Митхэ, спрячь саблю, дело уже сделано.

Митхэ потрясённо стояла, выставив саблю перед собой;
Эрхэ, поколебавшись, аккуратно забрала оружие из рук подруги и вытерла подолом рубахи.

Где-то в ночи снова раздались крики и зазвенело оружие.
И опять тишина.

--- Они обвязательно провверить жилище, --- шёпотом заявил Аурвелий, пристально вглядываясь в темноту.
--- Они провверять все укрытие.
Надо уходить.

\ml{$0$}
{--- Доверимся реке, --- прошептал Ситрис.}
{``Let's trust in the river,'' \Sitris\ whispered.}
--- Она скроет наше тепло и запах.
Вон там стебли муравьиного зонта, через них хорошо дышать.
Найдите камни потяжелее и избавьтесь от всего, что может очень невовремя всплыть.

Митхэ продолжала стоять на месте;
её взгляд был направлен в бесконечность...

--- Митхэ, --- Эрхэ дёрнула остолбеневшую подругу за штанину, затем отвесила пощёчину.
--- Митхэ!
Нельзя стоять на месте!
Пошла!

\section{Преображение}

Новый крик, ужаснее прежних, разорвал ночь.
Но на этот раз я слышал его как-то по-другому.
Всё вокруг внезапно изменилось, словно открылся древний, невообразимо глубокий колодец знаний.
В голове вдруг появились воспоминания о вещах, которые я никогда не видел, слова языков, о которых я никогда не слышал.

Я ощутил невообразимо приятное, до боли знакомое чувство --- чувство открывающихся глаз.
И этим глазам предстало величественное зрелище --- другой спектр, другое пространство, бесконечное, многомерное, разрисованное невообразимо сложной квантовой вязью.
Я почувствовал подобие света, исходящего от меня...

И сгустки древнего, ненавистного, голодного <<антисвета>>, которые меня окружали.
Ну что ж, Картель, я снова проснулся.

<<Уничтожить>>.

Эффект внезапности сработал.
Мой удар был похож на взмах крыла, перпендикулярный пространству и времени.
Два жреца замерли с открытыми в ужасе ртами --- их демоны перестали существовать.
Первый жрец, окутанный эманациями умирающей девушки, неожиданно легко отклонил этим <<облаком>> мою атаку, чем привёл меня в сильное замешательство.
Однако затем стратег совершил ошибку --- попытался нанести ответный удар вместо того, чтобы бежать.
Отражение, контратака, обход слабой защиты, уничтожение.
Тем временем моё тело, взбодрившись от поступившего в кровь кураре-антидота, схватило нож и разделалось с телами врагов.
Один из них бросился бежать, и я просто охладил ствол его мозга до температуры замерзания глицерина.
Жрец шлёпнулся, как кожаный ремень.

Я склонился над Чханэ, провёл рукой по высокому, забрызганному кровью лбу.
Её тело было одной сплошной раной.
Трудно поверить, что прикрученный к чёрному алтарю человеческий обрубок ещё на закате был красивой молодой девушкой.
Остекленевшие оранжевые глаза остановились на моём лице, судорожно приоткрылся искривлённый мукой рот:

--- Никх\FM...
\FA{
Добей (сигнальная система Десять Песчинок).
}

--- Нет, Чханэ, --- ласково прошептал я.
--- Я тебя спасу.

<<Ликхмас, нет.
Энергия тебе нужна.
Это всего лишь человек>>.

<<Нет, Аркадиу Люпино.
Поступай так, как нужно мне.
Остальное потом>>.

И снова из того же колодца услужливо появились нужные знания.
Мозг... Да, вот здесь...

Глаза Чханэ расширились, маска страдания сменилась выражением крайнего удивления.
Боль ушла.

Я положил руку на развороченный жертвенным кинжалом живот.
В моей голове мелькали, сменяя друг друга, формулы, многомерные диаграммы, вычислительные данные.
Живот Чханэ засветился слабым светом.
Медленно начал восстанавливаться изрезанный кишечник, растворялись тромбы, исчезали случайно попавшие в рану микроорганизмы, сосуды сшивались, вставали на место и наполнялись кровью.

--- Лис, что ты делаешь?

Свет погас, и я ощупал гладкий, розовый и совершенно целый живот.
Снова вспыхнул свет --- и бронзовые скобы пробками вылетели из чёрного камня алтаря.
Я аккуратно взял в руки культю, поднял отрубленную ногу.
С этим посложнее будет, но тоже исправимо...

Вновь брызнул белый свет.
Срослись, как ни в чём не бывало, кость и мышцы, нашли друг друга сосуды и нервы, исчезли из тканей успевшие накопиться продукты метаболизма.

Не успели звёзды пройти и ноготь пути, как Чханэ, дрожа всем телом от пронизывающего ветра, лежала на алтаре живая и совершенно невредимая.

--- Лис...

--- Меня зовут Arcadiv Valerianv Luppino.

Древняя рыба снова всплыла на поверхность --- на этот раз внутри меня.
Я впервые увидел в глазах Чханэ ужас.

Тут словно рухнула плотина, перегораживающая реку.
Интеграция личностей завершилась.
В мою голову хлынула масса информации, ждавшей своего часа.
Жизни, обучение, задания, приказ, с которым нас отправили на эту планету.
Мои товарищи --- Анкарьяль Кровавый Шторм и Грейсвольд Каменный Молот.

Девушка, шатаясь, попыталась встать с алтаря и чуть не упала.
Я поддержал Чханэ и вбросил ей в ликвор чуть-чуть морфина-44.
Это помогло --- взгляд стал более ясным.

--- Пойдём, Чханэ, нужно спешить.

Как же интересно звучит родной язык сели.
На нём вдруг в разы сложнее стало говорить.
Родной?
Хай, ej-a secunda, на каком lingai mej думать-raj?..
Я помотал головой.
Да, кажется, такое случается при форсированной интеграции личностей.

--- Лис, в храме остались подростки из семей Разрушителей.
Они прошли Отбор, но их всё равно принесут в жертву!

--- Без жертв Хатрафелю город погибнет, --- возразил я.
--- Я собирался спасти только тебя.

--- Лис!
Я тебя не узнаю!
Представь, что это наши дети!
Ты бы их оставил?

--- Погибнут не они, так другие.
Какая разница?

Чханэ беспомощно уставилась на меня.
Я вдруг вспомнил, как впервые обнаружил её в лесу с одним кинжалом...

--- Хорошо.
Только найди себе одежду и оружие.

Чханэ методично содрала рубаху и штаны с мёртвого жреца, схватила мою фалангу.
Спустя минуту она уже стояла навытяжку, словно перед походом, как будто и не пытали её смертной мукой.
Поневоле зауважаешь.

Мы вошли в крипту, освещённую пылающей трескучей жаровней.
Комната была пуста, лишь на скамье сидело чьё-то уже остывшее тело с гротескно поднятыми руками.
Тень мертвеца плясала под дудочку пламени свой последний танец.
Со страшным чувством пустоты в голове я узнал Трукхвала.
Погиб ли мой учитель при попытке освободить Чханэ или по прихоти агентов Картеля, рассказать было уже некому.

Я аккуратно отвязал костлявые руки от факельных колец, опустил тело на скамью и с усилием разогнул окоченевшие ноги.
Разорванная рубаха разошлась, и за пазухой старика я вдруг заметил тугой свиток пергамента.

Я достал его, развернул и поднёс к свету:

\begin{quote}
<<Людям Золотой Пчелы от Трукхвала ар'Хэ, жреца Тхитрона.\\
Змея 4 дождя 12011, Год Церемонии 5.\\
~\\
Я провёл расследование недавних диверсий в Храме Тхитрона и получил доказательства, что эти диверсии совершены группой жрецов Винтового Храма с целью захвата власти.
Есть основания утверждать, что рекомые жрецы были в сговоре со жрецами Трёхэтажного Храма, из-за чего я вынужден обратиться за помощью к вам, а не к Тхартхаахитру.
Перечень доказательств, а также список людей, которые могут подтвердить эти данные, ниже.
Письмо доставит мой ученик, Ликхмас ар'Люм, жрец Тхитрона, вынужденный бежать из-за угрозы убийством.
Скорее всего, я умру в ближайшие дни и не смогу выступить на Деле Перекрёстка.
Поэтому прошу снять с письма две копии и оставить эти копии у надёжных людей>>.
\end{quote}

\section{Золотая дорога}

Золотая дорога стояла обманчиво пустой.
Разумеется, это была ловушка --- Машина засекла взлом и поставила ложную картинку.
Она уверилась в победе.
Она даже не сломала систему, разворачивающую защитное поле, а наоборот --- повысила мощность поля до максимума, чтобы мы не выбрались.
Это было тактической ошибкой.
Ценой огромных усилий и стараниями трёх техников --- меня, Заяц и прибежавшего с разведки Фонтанчика --- удалось переключить управление полем, и тюрьма стала защитой.
Между нами и кораблём лежали два километра пути, тси и роботы-техники, превращённые Машиной в солдат.

Разумеется, все знали, что задумал Комар.
Все до единого.
Я понял это, едва увидел его грустное и слегка насмешливое лицо.
Лицо, которое я так любил.
Хотя кого я обманываю.
Я буду помнить и любить его до самой смерти.
Я знал --- он и прочие члены особого отдела просто не смогут жить после того, что им придётся сделать.
Что они скажут себе --- что они выполняли мой приказ?
Этой отговорки хватит лишь до трапа корабля.

\ml{$0$}
{--- Небо, мы подадим сигнал с панели в коридоре, когда доделаем всё, --- бросил Комар.}
{``Sky, we'll call you from the control board in the corridor, when we've done,'' Mosquito said.}

\ml{$0$}
{--- \emph{Если}, --- тихо хрюкнул кто-то.}
{``\emph{If},'' somebody quietly grunted.}
Кажется, Кукловод.

\ml{$0$}
{--- \emph{Когда}, --- поправил его Комар.}
{``\emph{When},'' Mosquito corrected him.}
--- Кто хочет уйти --- уходите.
С каждым ушедшим наши шансы упадут в два раза.
Лиана, кажется, ты хотела?

--- В жизни и смерти я буду щитом, --- возмутилась где-то рядом с ним Лиана.
--- Командир, мы вообще-то ждём.
Или ты предлагаешь мне сказать за тебя то, что следует?

Я хорошо помню, что после этих слов силач Фонтанчик зарыдал.
Это было страшное зрелище --- он плакал и выл, не прекращая поддержку систем ни на секунду.
Подсолнух подбежала к нему, запрограммировала плечевой имплант на инфузию стимулятора и вытерла ладошкой слёзы, катящиеся по его морде.

--- Солнышко, держись, скоро всё закончится! --- шептала она и облизывала ему уши.

--- Это какое-то безумие...

--- Так-так, --- Фонтанчик титаническим усилием воли остановил рыдания и махнул рукой.
--- Держите Заяц, она горит!

Заяц медленно села на пол.
Хомяк и Пирожок бросились к ней и приняли управление на себя.

--- Я не верю, --- померкшим голосом проговорила Заяц и легла на пол, глядя расширенными глазами куда-то в бесконечность.
--- Я не верю, что это происходит на самом деле.
Этого просто не может быть.
Я сейчас проснусь и поеду на работу.
У меня завтра куча дел...

--- Ради жизни, --- раздался в телекоме сухой голос Комара, --- особому отделу реализовать протокол <<Тайфун>>, модуль <<Живая сталь>>.
Повторяю, особому отделу реализовать протокол <<Тайфун>>, модуль <<Живая сталь>>.

...Телеком прервал вещание спустя секунду.
Я не знаю, что было в том коридоре, в этих пятистах двадцати метрах кошмара.
Я никогда этого не узнаю.

Сигнала мы не дождались.
Мы взяли персональные поля и оружие, прекрасно понимая, что теперь пришла наша очередь убивать или быть убитыми.
Я ступил на золотой пол первым.
Телепатический голос возвестил:

\begin{quote}
<<Золотая дорога --- галерея достижений.
Узри, о тси, жизнь и процветание!>>
\end{quote}

Товарищи один за другим проходили сквозь ворота и ёжились, услышав этот голос.
Под нашими ногами хлюпала вода, окрашенная в красный, коричневый, зелёный...
Она медленно сочилась из разбитых смотровых лазеек, хлестала рекой из вывороченных с корнем водопроводных труб.
С водой безуспешно пытался справиться одинокий сливной насос, давясь обрывками ткани, осколками стекла, завывая то ли от боли, то ли от отчаяния...

Обессилевшего Фонтанчика нёс на себе мускулистый Мак, обливаясь потом и шепча что-то вроде <<Вот-вот, хорошо, почти дошли>>.
Не привыкший к такому двухметровый канин слабо протестовал.
Заяц шла сама, но мне пришлось вести её за руку;
женщина плохо понимала, где находится и что происходит.

--- Жизнь и процветание, --- Заяц улыбнулась жалкой полубезумной улыбкой, услышав телепатический голос.
--- Это просто был сон.
Конечно, вот же она, явь!
Жизнь и процветание!

--- Заяц, пожалуйста, не останавливайся, --- взмолился я, когда она явно вознамерилась поглядеть голографический фильм.
--- Верь во что хочешь, но иди куда следует!

Перед нами начали разворачиваться голограммы.
Улыбающиеся, радостные тси.
Весёлые празднества.
Великие стройки, дух оживления и взаимопомощи, который царил на них.
Простые будни рабочих на объектах, с их обыденными добрыми шутками вроде бегающей по сверхстерильному цеху голографической мышки.
Однако голограммы не могли скрыть груды искорёженного металла и плавающее в грязной воде кровавое месиво.
Зеленоватый свет полуденной Звезды, лившийся сквозь кристаллический потолок, радостно освещал клубы тяжёлого, остро пахнущего горелой плотью дыма.

Танковые доспехи апидов Ди лежали, раскорячив ноги, словно огромные мёртвые жуки-бронзовки.
Их последний поход, спустя почти сто тысяч лет после Тараканьей войны, закончился.
Доспех моей любви с гордым золотым комариком на борту принял на себя первый снаряд, не оставивший пилоту ни единого шанса.
Баночка закрыл мне рукой глаза, когда мы проходили мимо.
На меня пахнуло тошнотворным запахом горелого хитина.
По голосу Баночки я понял, что зрелище было ужасающим.

--- Не надо, пожалуйста, не смотри, --- со слезами бормотал он, когда я попытался отмахнуться.
--- Запомни его другим...

Из пехоты мы узнали только двоих --- Лиана, как ребёнка, держала полуоторванной рукой Кукловода, вернее, его верхнюю часть.
На зелёном миниатюрном личике Кукловода замерло удивление, на обугленном лице Лианы --- усталая улыбка.
Искорёженные, но не сломленные, они сидели у панели, замерев в последней попытке подать нам сигнал.

\section{Путь охотника}

Дорога на второй этаж оказалась пустынной --- воины, не подозревая о произошедшем на вершине, сидели у входов и ждали других лазутчиков.
Храм словно вымер.
Комната для ожидающих жертвоприношение тоже оказалась пустой.
Я подумал о Сиртху-лехэ --- погиб ли он в бою или его пленили, чтобы принести в жертву?

--- Как ты попал в храм, Лис?

--- Сиртху-лехэ отвлёк воинов, Ситрис и Кхохо помогли забраться на крышу.

--- Ситрис, --- Чханэ заворчала злобно, как мурзящийся кот.
--- Предатель, урод пресноводный.
Тебя на смерть послал, а сам, значит, отсидеться решил...

--- Нас никто не предавал, --- перебил я подругу.
--- План был хорош, но воины понятия не имели, с какими силами нам пришлось столкнуться.
Ситрис --- не исключение.
Где дети?

--- Там, --- она указала на ведущую в глубь храма лестницу.

Да, Храм определённо не был тем, что раньше.
Прежде в этих помещениях держали еду, изредка кур и других мелких животных на убой.
Недолго, разумеется --- сели избегали держать зверей взаперти.
Теперь же в холодной металлической клетке, пропахшей куриным навозом и шерстью, сидели люди.
Трое подростков испуганно жались друг к другу в углу --- девчонка постарше держала в объятиях совсем маленьких мальчишек, почти детей.
В противоположном углу полулежал истекающий кровью старик.

Чханэ повозилась с замком и забежала внутрь.

--- Митхэ, Зонтик, Веточка! --- шёпотом позвала она.

--- Чханэ?
Это ты?
Тебя не забрали боги? --- услышал я знакомый голос.

--- Нет, не забрали.
И вас не заберут, если будете вести себя тихо.
Идём.
Вставайте.

\ml{$0$}
{--- Учитель!}
{``Teacher!''}

Митхэ кинулась мне на шею.
\ml{$0$}
{Так я узнал, что незаслуженные объятья тоже могут причинить боль --- ведь я чуть не оставил её здесь.}
{It's how I found out that undeserved hugs may be painful---I was so close to leave her there.}
\ml{$0$}
{<<Кто же знал>>, --- попытался возразить я себе.}
{\textit{Who would know,} I tried to counter.}
\ml{$0$}
{Вышло неубедительно.}
{It didn't work out.}

--- Учитель Ликхмас, я знала, я знала, что ты придёшь, --- шептала Митхэ и осыпала моё лицо мокрыми поцелуями.

\ml{$0$}
{--- Прости меня, милая, --- сказал я.}
{``I'm sorry, love,'' I said.}
\ml{$0$}
{--- Я больше не твой учитель.}
{``I'm not your teacher now.}
\ml{$0$}
{Я даже не жрец.}
{I'm not even a priest.''}

\ml{$0$}
{--- Это навсегда, --- всхлипнула Митхэ.}
{``It can't be undone,'' \Mitchoe\ sobbed.}
\ml{$0$}
{--- Это же как кормилец или друг, да?}
{``It's like a nurse or friend, isn't it?''}

Мы подошли к Сиртху-лехэ.
Два удара тонкого клинка пришлись в печень и селезёнку.
Бледный как снег старик, зажимая окровавленной рукой живот и часто дыша, слабо улыбнулся серыми губами.

--- Ты смог...

--- Да, дедушка, --- я вдруг почувствовал слёзы на лице.

--- Хорошо, --- выдохнул он.
--- Бегите, скоро они вернутся...

Глаза старика вдруг остекленели, он больно схватил меня за запястье.

--- Ситрис у главного входа.
Ситрис --- друг, --- отчеканил он.
--- Кхарас и Эрликх ходят по маршруту зал --- западный коридор --- заднее крыло --- восточный коридор.
Кхарас и Эрликх --- враги.

--- Я понял и запомнил, Дом-с-Карамельными-Окнами, --- так же чётко прошептал я.

Ветер Духов улетел, и хватка Сиртху ослабла.

--- Хэситр, --- засуетилась Чханэ, оглядываясь по сторонам.
Я увидел, что по её лицу тоже текут слёзы.
--- Нужен хэситр...

--- Брось, милая, брось, --- улыбнулся старик.
--- Неужели ты веришь, что мне, всю жизнь ходившему лесными тропами, нужна какая-то жалкая чаша с водой, чтобы достичь пристанища духов?

--- Ты лучший дедушка, --- сказал я.
--- Я люблю тебя.

--- И я вас люблю, малыши, --- слабо усмехнулся Сиртху-лехэ.
--- Будьте счастливы.
А я прожил хорошую жизнь...

Голос его затих.
Чханэ, обмакнув палец в кровь, нарисовала на бледном лбу старика Обнимающего Сита.
Митхэ отпустила руки мальчиков, сняла с запястья сплетённый из юкки браслетик и положила в сухую, испачканную кровью ладонь усопшего.

\section{Нежданная помощь}

Митхэ сидела на берегу и перебирала в руках мокрые обломки муравьиного зонта.
Эрхэ с виноватым видом сидела рядом.

--- Золото, извини.
Я ж не знала, что она у тебя в кармане, да и потянула...

Митхэ не ответила.
Она повертела пальцами один кусочек, затем другой, потом сложила их вместе.
Облупившиеся полоски чёрной краски сложились в закрытый глаз Плачущего Митра.

--- Это его подарок был, да? --- сочувствующе спросил Ситрис.
--- Хрупкая вещь, не для воинской жизни.
Как такую не сломать...

--- Атрис говорил, что я тоже не для воинской жизни, --- заметила Митхэ.
\ml{$0$}
{--- Может, он и прав был.}
{``Perhaps he was right.''}

\ml{$0$}
{<<Был>>.}
{\textit{He was.}}
\ml{$0$}
{Митхэ яростно помотала головой.}
{\Mitchoe\ furiously shook her head.}
\ml{$0$}
{<<Я не смогу смириться с твоей смертью.}
{\textit{I won't deal with your death.}}
\ml{$0$}
{Просто не смогу.}
{\textit{I just can't.}}
\ml{$0$}
{Ты сделаешь мне новую, своими руками, в нашем собственном доме.}
{\textit{You shall make a new ones, by your own hands, in our own house.}}
\ml{$0$}
{Мне и каждому из наших детей, чтобы музыка всегда была с нами>>.}
{\textit{For me and each our child---so the music will always be with us.}}

Она вытряхнула остатки соли из кожаного мешочка, положила в него обломки флейты и повесила мешочек на шею.

\ml{$0$}
{--- Найдите место на ночь.}
{``Find a place for the night.}
Утром мы пройдём по окрестностям и оставим знаки.
Нужно собрать отряд, всех, кто...

--- Нельзя, --- отрезал Ситрис.
\ml{$0$}
{--- Даже не думай.}
{``Don't even think about it.''}

--- Не думал, что скажу это, Золото, но Ситрис прав, --- хмуро сказал Акхсар.
--- Люди, связанные с тобой клятвой, пойдут на верную смерть.
\ml{$0$}
{Давай пропадём, скроемся в сельве, умрём для всего мира.}
{Let's disappear, shelter in the Silva, die for the whole world.}
Поодиночке у каждого из них будет шанс устроить свою жизнь, свободную от бед и угрызений совести.

Эрхэ согласно кивнула.

\ml{$0$}
{--- По-твоему, это будет честно? --- обратилась Митхэ к подруге.}
{``Do you think it would be fair?'' \Mitchoe\ asked her friend.}

\ml{$0$}
{--- Это будет милосердно, Золотце.}
{``It would be merciful, Goldy.}
Ты не можешь освободить их от клятвы иначе, не подвергая смертельной опасности.

Митхэ задумалась.

--- Хай.
\ml{$0$}
{Моя репутация и так уже, так что...}
{My reputation has just---so---}
\ml{$0$}
{Как ты там обычно говоришь, Ситрис?}
{What did you use to say in such situations, \Sitris?''}

\ml{$0$}
{--- Сгорел амбар, гори и город! --- улыбнулся разбойник.}
{``The barn burned down, let's burn the city!'' the outlaw smiled.}

Митхэ кивнула.

\ml{$0$}
{--- Именно.}
{``This.}
\ml{$0$}
{Тогда... ваша служба окончена.}
{So, your service is over.''}

--- А мы-то тут при чём, --- проворчал Аурвелий.
--- Мы пока не терять в лес.

--- В общем и целом сенвиор Амвросий прав, --- согласился Акхсар.
--- О присутствующих здесь речи не шло.

Вдруг на берегу показался силуэт.
Он робко приближался к сидящим воинам;
все застыли, не в силах преодолеть потустороннюю мощь наваждения.
Копыта ступали по камням с лёгким стуком;
ветвистые рога словно лучились в слабом сиянии звёзд.
Первой опомнилась Митхэ.

--- Серебряный! --- воительница бросилась к оленю и обняла его.
Животное ласково пожевало её рубаху.
--- Хоть одна хорошая новость...
Как ты выбрался?
Кто тебя оседлал?
Я же оставила тебя на постоялом дворе!

\ml{$0$}
{Митхэ залезла в седельную сумку и с недоумением вытащила оттуда связку провианта.}
{\Mitchoe\ opened the saddle bag and was surprised by food pack inside.}
\ml{$0$}
{На банках стояли печати Трёхэтажного Храма.}
{All containers got Three-storey seal.}
\ml{$0$}
{В другой сумке оказались когти и возмущённое шипение;}
{Second bag was full of claws, hiss, and indignation;}
спустя секхар оттуда выпрыгнул Цапка-лехэ и тут же высокомерно ушёл в кусты с поднятым хвостом, показывая, что он вовсе не рад видеть хозяйку.

--- Цапка-лехэ!
Ну куда ты?
Серебряный!
Как же я рада...

\ml{$0$}
{--- Митхэ, нам придётся отпустить оленя, --- пробормотал Акхсар.}
{`\Mitchoe, we have to release the deer,'' \Akchsar\ mumbled.}
\ml{$0$}
{--- Он белый, его видно за кхене.}
{``He's white, he could be seen for a whole \kchenoel.''}

\ml{$0$}
{--- Я тоже белая, Снежок, --- вздохнула воительница.}
{``I'm white too, Snow,'' the warrior sighed.}
\ml{$0$}
{--- Белее некуда.}
{``I'm the whitest.''}

Она подхватила оленя под уздцы.
Эрхэ, Акхсар и Аурвелий, вздохнув, начали приводить в порядок амуницию.

--- Пойдём.
Ночь длинная, мы успеем пройти по тракту достаточно кхене.
Цапка-лехэ, дедуля, мы уходим.
Да оставь ты в покое эту мышь, я тебя покормлю.
Что <<мяу>>?
Я ждать не собираюсь.
Пойдём.
Ситрис, и ты тоже.
Ситрис?

Разбойник сидел на том же месте, где она оставила его.
Непроницаемо чёрные матовые глаза не отрываясь смотрели в воду.

--- Ситрис, ты со мной?

Разбойник вздохнул и встал.

--- Да пусть остаётся или идёт, куда хочет, --- плюнул Акхсар.
--- Он бросит нас в самый ответственный момент.

--- Тихо, --- остановила друга Митхэ.
--- Ситрис, ты пойдёшь со мной?

Ситрис снова вздохнул.
Эрхэ и Акхсар хмуро смотрели на него.

--- Хорошо, --- сказал наконец Ситрис.
--- Пойдём искать твоего мужчину.
По правде сказать, мне и идти-то больше некуда.
Я и не разбойник, и не честный человек.
За пятьдесят дождей я приобрёл много знаний и навыков, но так и не понял, чего хочу.
Пойдём, Митхэ.
Может, что-то из этого и выйдет.

В тот момент Митхэ осознала, что эти слова кажутся ей дороже всех клятв.
Как и Атрис, Ситрис не был связан с ней ни долгом, ни моральными ценностями.
Он шёл с ней, потому что хотел идти.
И он был честен.
Он был кристально честен в своей трусости, беспринципности и эгоизме.

\ml{$0$}
{<<Поэтому я ему и верю>>.}
{``\emph{That's why he's trusted.}''}

\section{Кораблю --- взлёт!}

\epigraph{
С тех пор Мерлин дремлет в своей башне,\\
надёжно хранимый оком Нинианы.\\
Легенда гласит,\\
что в тяжёлые для мира времена он проснётся,\\
чтобы снова служить на благо людям.
}{
\ml{$0$}
{Легенда о кинью\FM\ Артьюре, Древняя Земля}
{The legend of Kinju\FM\ Artjur, Mother Earth}
}
\FA{
Кинью (энглис) --- предводитель союза племён.
}

Одиннадцать тысяч триста пятидесятая капсула анабиоза откликнулась радостным сигналом: <<Живое существо успешно остановлено>>.
Эти капсулы были последним словом техники.
Потребовались десятилетия, чтобы исключить попадание внутрь капсул квантов извне --- скопившиеся в зоне пенумбры\FM\ фотоны могли сжечь находящегося в капсуле сапиента.
\FA{
Пенумбра --- область градиентного изменения скорости времени.
}
Когда технология только появилась, её пытались применить ко всему --- были даже хранилища для пищи.
Правда, создатели шутили, что мясо не только остаётся свежим, но и поджаривается при попытке его вытащить.
Временем управлять просто, потомок.
Сложнее избежать расплаты...

Водного лифта к Дракону не было, поэтому нам потребовалось ещё около двух часов, чтобы на руках перенести в корабль всех дельфинов.
Они не возражали, только молча плакали.
В дельфиньем отсеке капсул хватило не для всех;
дельфины договорились, что бодрствовать будут по очереди.
Малышей мы выманили сладостями и аккуратно положили в <<сухопутные>> капсулы --- конструкция это допускала.
Хвала и честь предусмотрительности древних инженеров.

Мак и Листик перед анабиозом навестили кольцевую теплицу корабля.
Она чувствовала себя превосходно.
Баночка не пожелал лезть в капсулу;
он вызвался наладить связь Стального Дракона, да так и заснул в своём кресле.
Мы не стали его будить --- налаживать было нечего и незачем.

Заяц потихоньку начала отходить от временного помрачения сознания, а мы потихоньку привыкли к её жуткому невидящему взгляду, направленному в бесконечность.
Костёр осмотрел больную, побеседовал с ней, поколдовал с препаратами и объявил, что Заяц выздоровеет, но настоятельно рекомендовал воздержаться от анабиоза до полного выздоровления.

--- Это реакция на быструю смену обстановки, --- объяснил он.
--- А анабиоз тоже в своём роде быстрая смена обстановки --- сегодня здесь, завтра уже там.

--- А что у неё с глазами? --- тихо спросил Фонтанчик врача.
--- Это же кошмар, я первый раз её такой вижу...

--- Ты в первый раз оказался в таком переплёте.
Это пройдёт, --- успокоил друга Костёр.
--- Выглядит пугающе, но ничего серьёзного.
Просто мозг работает в режиме энергосбережения.

Мы усадили подругу в кресло и дали в руки первую подвернувшуюся игрушку-антистресс.
Заяц погладила лягушку, и та вцепилась в пальцы женщины лапками-присосками.

--- Знаешь, о чём я думал всё это время? --- тихо сказал я.
--- Помнишь, мы хотели...
В тот вечер, когда...

--- Печеньки, какао и тёплая постелька, --- из невидящих глаз Заяц покатились слёзы.
--- Небо, всё ещё впереди.
Мы вырастим какао и злаки, если найдём хоть горсть почвы.
Мы построим дома даже на лодках или на морском дне.
Всё обязательно будет.

--- Мне ужасно стыдно, --- признался я.
--- Я думал о них все эти дни.
Когда тси умирали, когда Комар... когда мы решали, что надо улетать, я думал о какао и печенье.

--- Знаешь, Небо, --- заметил Фонтанчик, --- полагаю, что именно поэтому наши предки построили общество.
Может, война и была нужна для жизни, но для процветания важны простые радости --- домашний уют, вкусная пища, общение, секс, игры и любимое дело.
То, что по сути могут обеспечить твоя собственная голова, крохотный кусочек планеты и два-три живых существа рядом.

Я промолчал.

--- Небо, мы бы хотели, чтобы ты радовался.
Комар прожил долгую, полную радостей жизнь, какую и заслужил.
Он умер так, как многим и не снилось --- защищая народ от неминуемой гибели.
Не закрывай любовь в себе.
Твой вид немногочисленен, но среди выживших есть прекрасные апиды, которым она нужна как никогда.

Я посмотрел на измождённое доброе лицо Фонтанчика.
Он всегда находил силы подбадривать других, даже если самому приходилось несладко.

--- Дружище, иди-ка ты отдыхай, --- сказал я и погладил его приятную шерстистую шею.

--- Нет, мой отдых начнётся сейчас, --- осклабился он и, откинув упавшую на глаза седую гриву, сел за штурвал.
\ml{$0$}
{--- Небо, второй пилот.}
{``Sky, co-pilot.}
\ml{$0$}
{Заяц, сиди и гладь лягушонка, сейчас за окном будет красиво.}
{Hare, sit and pet your froggy-moggy, it soon will be a nice view beyond the window.}
\ml{$0$}
{Баночка... хм.}
{Flask ... um.}
\ml{$0$}
{Спи дальше, дружище.}
{Sleep well, old friend.''}

Да, я уже и забыл, что Фонтанчик страсть как любит водить транспорт.
Он предпочитал старые аналоговые системы с громоздкими штурвалами современному нейроинтерфейсу.
Мы часто вылетали с ним в открытый океан.
Этот бескрайний, бездонный океан смывает все невзгоды;
он превращает маленькие проблемы в марлинов, а большие --- в китов, и они, величественно взрезав гладь, навеки остаются в чистой, насыщенной жизненной радостью лазури.

Наша планета пала за три дня.
Я не мог представить кита крупнее;
должно быть, с ним мог сравниться разве что доисторический лебиапант, который, по преданию ушедших предков, до сих пор лежит где-то в водах Древней Земли.
Однако, глядя на Фонтанчика, который с просветлевшим взглядом настраивал космолёт, глядя на Заяц, которая с интересом следила за другом и повторяла его движения, глядя на мирно спящего Баночку, я поверил в океан.
Я поверил в то, что там найдётся достаточно места даже лебиапанту.

Курс был проложен --- три парсака, жёлтый карлик, тси-ди-подобная необитаемая планета, мельком упомянутая пленёнными демонами Ордена Преисподней.
Мы наудачу в последний раз разослали сигналы во все места, где могли остаться тси.
Молчание.
В живых осталась только она --- родная планета, которая мирно спала и не спешила отвечать существу, выросшему в её объятьях и прожившему на ней почти всю жизнь.

Старое название Тси-Ди --- Мерлин-Ниниана.
Когда я впервые узнал это название, я долго повторял его про себя, удивляясь чудесному ассонансу.
Историки утверждают, что это мужское и женское имена, распространённые во времена Эпохи Богов.
Возможно, так звали первооткрывателей.
Кем они были?
Любовниками?
Братом и сестрой?
Хорошими друзьями?
Всем вместе?
Время надёжно укрыло тайну своим покрывалом.

\ml{$0$}
{Корабль медленно набирал скорость.}
{The spaceship was slowly gathering speed.}
\ml{$0$}
{Небо становилось всё темнее и темнее, а океан засиял, словно грань живого лазерного сапфира.}
{The sky was darkening more and more, and the ocean was beaming like a face of animated laser sapphire.}
\ml{$0$}
{Вскоре Тси-Ди превратилась в две далёкие яркие точки.}
{Soon Qi-Di has turned into a couple of bright points.}
\ml{$0$}
{Спи, мой первый и последний дом.}
{Sleep, my first and last home.}
\ml{$0$}
{Машина уничтожила тси, но тебя она в обиду не даст.}
{Machine has wiped Qi out, but you won't be harmed.}
\ml{$0$}
{Это я знаю совершенно точно.}
{I know it for sure.}

\section{Время на сон}

\epigraph
{Истинное испытание на человечность лишено правил и границ.}
{Леам эб-Салах}

Завывал ветер, тёмным морем колыхалась в алеющем рассвете далёкая поверхность джунглей.
Где-то печально свистела на своей флейте согхо.
Мы впятером быстро и осторожно спускались со ступеней храма.
Один из мальчишек оступился и покатился вниз, почти беззвучно вскрикнув от боли.
Чханэ едва успела схватить его за руку.
Мы остановились.

--- Я ногу подвернул, --- тихо плакал мальчик, показывая на ободранную лодыжку.

--- Ничего, Веточка, терпи, ты же будущий воин, --- ласково шепнула Чханэ и, пригладив воину чёрную как смоль шевелюру, забросила мальчишку на спину.
--- Цепляйся крепче.
Лис тебя потом вылечит.
Меня вот вылечил...

Я обернулся, чтобы кивком поблагодарить Ситриса.
Пожилой воин едва заметно шевельнул пальцами руки.

--- Под стену, быстро, --- шикнула Чханэ детям, даже не успев осмыслить жестовый сигнал.

Вскоре мы уже переводили дух, прислонившись спинами к холодной влажной стене храма.

--- Светает... --- обречённо шепнула Митхэ.
--- Скоро прокричат за утро...

--- Не бойся, ученица, успеем, --- успокоил я её.
--- На окраине возле реки есть пустое жилище.
Спрячемся там на время.

В это время к Ситрису подошёл Кхарас.
Часовой защёлкал кресалом, и сверху потянуло удушливым, слегка сладковатым конопляным туманом.

\ml{$0$}
{--- Все хорошо, --- тихо сказал Кхарас.}
{``All right,'' \Kcharas\ quietly said.}
\ml{$0$}
{--- Старик не собирался убивать --- стрелы были смазаны кураре.}
{``The old man wasn't going to kill, arrows were dipped in curare.}
\ml{$0$}
{Скоро воины придут в себя.}
{Warriors will recover soon.}
\ml{$0$}
{Как дела, Ситрис?}
{How it goes, \Sitris?''}

\ml{$0$}
{--- Всё просто замечательно, Кхарас, как и должно быть, --- отозвался воин.}
{``Wonderful, \Kcharas, as it should go,'' the warrior responded.}
\ml{$0$}
{Такого счастья в его голосе я не слышал никогда.}
{I've never heard his voice so happy.}
\ml{$0$}
{--- Можешь скомандовать отбой --- Ликхмас и Чханэ уже далеко.}
{``You can command `lights out', \Likchmas\ and \Chhanei\ are far-far away.''}

\ml{$0$}
{Молчание.}
{Silence.}
\ml{$0$}
{Я почти слышал, как в голове боевого вождя крутится и скрипит мельничный жёрнов.}
{I could almost hear a millstone rolling and creaking inside warchief's head.}

\ml{$0$}
{--- Тху, на алтарь бы тебя, гнилая рыбина, --- беззлобно проворчал наконец Кхарас и, громко топая, вошёл в здание.}
{``\Tchu, may you lie on the altar, you rotten fish,'' \Kcharas\ muttered without malice, then went inside tromping loudly.}

% ----- INFO PAGE -----
\newpage\thispagestyle{plain}
{\centering

% ----- LICENSE SIGNS BLOCK -----
\includegraphics[width=2.5em]{cc.pdf}~\includegraphics[width=2.5em]{by.pdf}
% ----- LICENSE SIGNS BLOCK -----

\vspace{0.2em}

% ----- LICENSE BLOCK -----
{\ml{$0$}{Данная книга распространяется под лицензией \textbf{CC~BY~4.0}.}{This book is distributed under the \textbf{CC~BY~4.0} license.}\par}
{\ml{$0$}{Подробнее о лицензии:}{Details:} \href{https://creativecommons.org/licenses/by/4.0/deed.ru}{creativecommons.org/licenses/by/4.0}\par}
% ----- LICENSE BLOCK -----

\vfill

% ----- AUTHOR AND TITLE -----
{\large\bookauthor\par}
\vspace{0.5em}
{\Large\textbf{\booktitle}\par}
% ----- AUTHOR AND TITLE -----

\vfill

% ----- BOOK INFO -----
{\ml{$0$}{Начато:}{Started:} \textit{\bookstarted}\par}
{\ml{$0$}{Последняя редакция:}{Latest revision:} \textit{\bookfinished}\par}
{\ml{$0$}{Подробнее о книге:}{The book details:} \href{https://github.com/regnveig/tofa}{github.com/regnveig/tofa}\par}
% ----- BOOK INFO -----

\vfill

% ----- DOCUMENT INFO -----
{\ml{$0$}{Дизайн обложки: \textit{Э.\,Весна}}{The cover design by E.\,Viesn\'a}\par}
{\ml{$0$}{Компьютерная вёрстка: \textit{Э.\,Весна}}{The computer layout by E.\,Viesn\'a}\par}
{\ml{$0$}{Создано с помощью \XeLaTeX}{Created with \XeLaTeX}\par}
% ----- DOCUMENT INFO -----

\vspace{0.5em}

% ----- PUBLISHER INFO -----
{\textbf{\ml{$0$}{Свободное издательство <<Цунами>>}{\textsc{Tsunami}, an independent publisher}}\par}
{\ml{$0$}{Томск 2023}{Tomsk 2023}\par}
% ----- PUBLISHER INFO -----

\vfill

% ----- QR CODES -----
{\hfill
\begin{minipage}[t]{0.25\textwidth}
{\centering
{\includegraphics[width=0.8\textwidth]{qr_github_tofa.pdf}\par}
\vspace{0.2em}
{\small{\ml{$0$}{О книге и авторе}{About the book and the author}}\par}
}
\end{minipage}
\hspace{2em}
\begin{minipage}[t]{0.25\textwidth}
{\centering
{\includegraphics[width=0.8\textwidth]{qr_license_by40.pdf}\par}
\vspace{0.2em}
{\small{\ml{$0$}{Текст лицензии}{License full text}}\par}
}
\end{minipage}
\hfill}
% ----- QR CODES -----

% }
% % ----- INFO PAGE -----
%
% % ----- EMPTY PAGE -----
% \newpage\thispagestyle{plain}~
% % ----- EMPTY PAGE -----
%
% % ----- BACK COVER -----
% \includepdf[pages={1}]{cover_back.pdf}
% % ----- BACK COVER -----

\end{document}
